\protect \chapter {线性方程组}
\protect \section {第 I 节：求解线性方程组}
\protect \subsection {一.I.1: 高斯消元法}
\begin{ans}{一.I.1.17}
      \begin{exparts}
        \partsitem 高斯消元法
          \begin{equation*}
            \grstep{-(1/2)\rho_1+\rho_2}
            \begin{linsys}{2}
               2x  &+  &3y      &=  &13  \\
                   &-  &(5/2)y  &=  &-15/2
            \end{linsys}
          \end{equation*}
          得到解为$y=3$与$x=2$。
        \partsitem 这里高斯消元法
          \begin{equation*}
            \grstep[\rho_1+\rho_3]{-3\rho_1+\rho_2}
            \begin{linsys}{3}
              x   &  &  &-  &z  &=  &0  \\
                  &  &y &+  &3z &=  &1  \\
                  &  &y &   &   &=  &4
            \end{linsys}
            \grstep{-\rho_2+\rho_3}
            \begin{linsys}{3}
              x   &  &  &-  &z    &=  &0  \\
                  &  &y &+  &3z   &=  &1  \\
                  &  &  &   &-3z  &=  &3
            \end{linsys}
          \end{equation*}
          给出$x=-1$，$y=4$，$z=-1$。
      \end{exparts}
    
\end{ans}
\begin{ans}{一.I.1.18}
      如果方程组含有矛盾方程，那么它无解。
      否则，如果有变量不引领任何一行，那么它有无穷多解。
      最后一种情形是，既没有矛盾方程，而每个变量又都引领某一行，
      这时它有唯一解。
      \begin{exparts}
        \partsitem 唯一解
        \partsitem 无穷多解
        \partsitem 无穷多解
        \partsitem 无解
        \partsitem 无穷多解
        \partsitem 无穷多解
        \partsitem 无解
        \partsitem 无穷多解
        \partsitem 无解
        \partsitem 唯一解
      \end{exparts}
    
\end{ans}
\begin{ans}{一.I.1.19}
      \begin{exparts}
       \partsitem 高斯消元
        \begin{equation*}
          \grstep{-(1/2)\rho_1+\rho_2}
          \begin{linsys}{2}
             2x  &+  &2y  &=  &5  \\
                 &   &-5y &=  &-5/2
          \end{linsys}
        \end{equation*}
        表明\( y=1/2 \)与\( x=2 \)是唯一解。
      \partsitem 高斯消元法
        \begin{equation*}
          \grstep{\rho_1+\rho_2}
          \begin{linsys}{2}
             -x  &+  &y   &=  &1  \\
                 &   &2y  &=  &3
           \end{linsys}
        \end{equation*}
        给出\( y=3/2 \)与\( x=1/2 \)这个唯一解。
      \partsitem 行化简
        \begin{equation*}
            \grstep{-\rho_1+\rho_2}
            \begin{linsys}{3}
                x  &-  &3y  &+  &z  &=  &1  \\
                   &   &4y  &+  &z  &=  &13
             \end{linsys}
        \end{equation*}
        由于变量$z$不是任何一行的首变量，
        这表明方程组有无穷多解。
      \partsitem 行化简
        \begin{equation*}
          \grstep{-3\rho_1+\rho_2}
          \begin{linsys}{2}
             -x  &-  &y   &=  &1  \\
                 &   &0   &=  &-1
           \end{linsys}
        \end{equation*}
        表明方程组无解。
      \partsitem 高斯消元法
        \begin{align*}
            \grstep{\rho_1\leftrightarrow\rho_4}
            \begin{linsys}{3}
                x  &+  &y   &-  &z  &=  &10 \\
               2x  &-  &2y  &+  &z  &=  &0  \\
                x  &   &    &+  &z  &=  &5  \\
                   &   &4y  &+  &z  &=  &20
             \end{linsys}
            &\grstep[-\rho_1+\rho_3]{-2\rho_1+\rho_2}
            \begin{linsys}{3}
                x  &+  &y   &-  &z  &=  &10 \\
                   &   &-4y &+  &3z &=  &-20\\
                   &   &-y  &+  &2z &=  &-5 \\
                   &   &4y  &+  &z  &=  &20
             \end{linsys}                                      \\
            &\grstep[\rho_2+\rho_4]{-(1/4)\rho_2+\rho_3}
            \begin{linsys}{3}
                x  &+  &y   &-  &z      &=  &10 \\
                   &   &-4y &+  &3z     &=  &-20\\
                   &   &    &   &(5/4)z &=  &0  \\
                   &   &    &   &4z     &=  &0
             \end{linsys}
        \end{align*}
        给出唯一解\( (x,y,z)=(5,5,0) \)。
      \partsitem 这里高斯消元法给出
         \begin{align*}
            &\grstep[-2\rho_1+\rho_4]{-(3/2)\rho_1+\rho_3}
            \begin{linsys}{4}
               2x  &\spaceforemptycolumn   &   &+  &z       &+  &w       &=  &5  \\
                   &   &y  &   &        &-  &w       &=  &-1 \\
                   &   &   &-  &(5/2)z  &-  &(5/2)w  &=  &-15/2  \\
                   &   &y  &   &        &-  &w       &=  &-1
             \end{linsys}                                             \\
            &\grstep{-\rho_2+\rho_4}
            \begin{linsys}{4}
               2x  &\spaceforemptycolumn   &   &+  &z       &+  &w       &=  &5  \\
                   &   &y  &   &        &-  &w       &=  &-1 \\
                   &   &   &-  &(5/2)z  &-  &(5/2)w  &=  &-15/2  \\
                   &   &   &   &        &   &0       &=  &0
             \end{linsys}
         \end{align*}
         这表明方程组有无穷多解。
      \end{exparts}
    
\end{ans}
\begin{ans}{一.I.1.20}
      \begin{exparts}
       \partsitem
         % sage: M = matrix(RDF, [[1,1,1], [1,-1,0], [0,1,2]])
         % sage: v = vector(RDF, [5,0,7])
         % sage: M_prime = M.augment(v)
         % sage: gauss_method(M_prime)
         % [ 1.0  1.0  1.0  5.0]
         % [ 1.0 -1.0  0.0  0.0]
         % [ 0.0  1.0  2.0  7.0]
         %  take -1.0 times row 1 plus row 2
         % [ 1.0  1.0  1.0  5.0]
         % [ 0.0 -2.0 -1.0 -5.0]
         % [ 0.0  0.0  1.5  4.5]
         %  take 0.5 times row 2 plus row 3
         % [ 1.0  1.0  1.0  5.0]
         % [ 0.0 -2.0 -1.0 -5.0]
         % [ 0.0  0.0  1.5  4.5]
          高斯消元法
          \begin{align*}
           \begin{linsys}{3}
              x  &+  &y   &+  &z   &=  &5  \\
              x  &-  &y   &   &    &=  &0  \\
                 &   &y   &+  &2z  &=  &7
           \end{linsys}
           &\grstep{-\rho_1+\rho_2}
           \begin{linsys}{3}
              x  &+  &y   &+  &z   &=  &5  \\
              0  &   &-2y  &-   &z    &=  &-5  \\
                 &   &y   &+  &2z  &=  &7
           \end{linsys}                                     \\
           &\grstep{(1/2)\rho_2+\rho_3}
           \begin{linsys}{3}
              x  &+  &y   &+  &z   &=  &5  \\
              0  &   &-2y  &-   &1    &=  &-5  \\
                 &   &    &+  &(3/2)z  &=  &7/2
           \end{linsys}
          \end{align*}
          再经回代得到$x=1$，$y=1$，$z=3$。
       \partsitem
         % sage: M = matrix(QQ, [[3,0,1], [1,-1,3], [1,2,-5]])
         % sage: v = vector(RDF, [7,4,-1])
         % sage: v = vector(QQ, [7,4,-1])
         % sage: M_prime = M.augment(v)
         % sage: gauss_method(M_prime)
         % [ 3  0  1  7]
         % [ 1 -1  3  4]
         % [ 1  2 -5 -1]
         %  take -1/3 times row 1 plus row 2
         %  take -1/3 times row 1 plus row 3
         % [    3     0     1     7]
         % [    0    -1   8/3   5/3]
         % [    0     2 -16/3 -10/3]
         %  take 2 times row 2 plus row 3
         % [  3   0   1   7]
         % [  0  -1 8/3 5/3]
         % [  0   0   0   0]
          这里高斯消元法
          \begin{align*}
            \begin{linsys}{3}
                3x  &  &    &+  &z   &=  &7  \\
                 x  &-  &y   &+   &3z    &=  &4  \\
                 x   &+  &2y   &-  &5z  &=  &-1
            \end{linsys}
            &\grstep[-(1/3)\rho_1+\rho_3]{-(1/3)\rho_1+\rho_2}
            \begin{linsys}{3}
                3x  &  &    &+  &z        &=  &7  \\
                    &  &-y   &+ &(8/3)z   &=  &5/3  \\
                    &  &2y   &- &(16/3)z  &=  &-(10/3)
            \end{linsys}                                      \\
            &\grstep{2\rho_2+\rho_3}
            \begin{linsys}{3}
                3x  &  &    &+  &z        &=  &7  \\
                    &  &-y   &+ &(8/3)z   &=  &5/3  \\
                    &  &     &  &0        &=  &0
            \end{linsys}
        \end{align*}
        发现变量~$z$不引领任何一行。
        方程组有无穷多解。
        \partsitem
          % sage: M = matrix(QQ, [[1,3,1], [-1,-1,0], [-1,1,2]])
          % sage: v = vector(QQ, [0,2,8])
          % sage: M_prime = M.augment(v)
          % sage: gauss_method(M_prime)
          % [ 1  3  1  0]
          % [-1 -1  0  2]
          % [-1  1  2  8]
          %  take 1 times row 1 plus row 2
          %  take 1 times row 1 plus row 3
          % [1 3 1 0]
          % [0 2 1 2]
          % [0 4 3 8]
          %  take -2 times row 2 plus row 3
          % [1 3 1 0]
          % [0 2 1 2]
          % [0 0 1 4]
          下列步骤
          \begin{equation*}
            \begin{linsys}{3}
               x   &+  &3y   &+  &z   &=  &0  \\
               -x  &-  &y   &   &    &=  &2  \\
               -x   &+  &y   &+  &2z  &=  &8
            \end{linsys}
            \grstep[\rho_1+\rho_3]{\rho_1+\rho_2}
            \begin{linsys}{3}
               x   &+  &3y   &+  &z   &=  &0  \\
                   &   &2y   &+  &z   &=  &2  \\
                   &   &4y   &+  &3z  &=  &8
            \end{linsys}
            \grstep{-2\rho_2+\rho_3}
            \begin{linsys}{3}
               x   &+  &3y   &+  &z   &=  &0  \\
                   &   &2y   &+  &z   &=  &2  \\
                   &   &     &   &z   &=  &4
            \end{linsys}
          \end{equation*}
          给出$(x,y,z)=(-1,-1,4)$。
      \end{exparts}
    
\end{ans}
\begin{ans}{一.I.1.21}
      \begin{exparts}
        \partsitem 由$x=1-3y$得到$2(1-3y)+y=-3$，从而$y=1$。
        \partsitem 由$x=1-3y$得到$2(1-3y)+2y=0$，
           从而得出$y=1/2$的结论。
      \end{exparts}
      使用这个方法的人必须把每个可能的解代回所有方程加以检验。
    
\end{ans}
\begin{ans}{一.I.1.22}
      作消元
      \begin{equation*}
        \grstep{-3\rho_1+\rho_2}
        \begin{linsys}{2}
          x  &-  &y  &=  &1\hfill  \\
             &   &0  &=  &-3+k\hfill
        \end{linsys}
      \end{equation*}
      由此得出：若\( k\neq 3 \)，则该方程组无解；
      若\( k=3 \)，则它有无穷多解。
      它永远不会有唯一解。
    
\end{ans}
\begin{ans}{一.I.1.23}
      令\( x=\sin\alpha \)，\( y=\cos\beta \)，\( z=\tan\gamma \)：
      \begin{equation*}
        \begin{linsys}{3}
           2x  &-  &y  &+  &3z  &=  &3  \\
           4x  &+  &2y &-  &2z  &=  &10  \\
           6x  &-  &3y &+  &z   &=  &9
        \end{linsys}
         \grstep[-3\rho_1+\rho_3]{-2\rho_1+\rho_2}
         \begin{linsys}{3}
           2x  &-  &y  &+  &3z  &=  &3  \\
               &   &4y &-  &8z  &=  &4   \\
               &   &   &   &-8z &=  &0
         \end{linsys}
      \end{equation*}
      得到\( z=0 \)，\( y=1 \)，\( x=2 \)。
      注意，没有\( \alpha\in\R \)满足$\sin\alpha=2$。
      
\end{ans}
\begin{ans}{一.I.1.24}
      \begin{exparts}
       \partsitem 高斯消元法
         \begin{equation*}
           \grstep[-\rho_1+\rho_3 \\ -2\rho_1+\rho_4]{-3\rho_1+\rho_2}
           \begin{linsys}{2}
              x  &-  &3y  &=  &b_1\hfill \\
                 &   &10y &=  &-3b_1+b_2\hfill \\
                 &   &10y &=  &-b_1+b_3\hfill \\
                 &   &10y &=  &-2b_1+b_4\hfill
            \end{linsys}
           \grstep[-\rho_2+\rho_4]{-\rho_2+\rho_3}
           \begin{linsys}{2}
              x  &-  &3y  &=  &b_1\hfill \\
                 &   &10y &=  &-3b_1+b_2\hfill \\
                 &   &0   &=  &2b_1-b_2+b_3\hfill \\
                 &   &0   &=  &b_1-b_2+b_4\hfill
            \end{linsys}
         \end{equation*}
         表明该方程组相容当且仅当
         \( b_3=-2b_1+b_2 \)与\( b_4=-b_1+b_2 \)同时成立。
       \partsitem 消元
         \begin{align*}
            &\grstep[-\rho_1+\rho_3]{-2\rho_1+\rho_2}
            \begin{linsys}{3}
              x_1  &+  &2x_2  &+  &3x_3  &=  &b_1\hfill  \\
                   &   &x_2   &-  &3x_3  &=  &-2b_1+b_2\hfill  \\
                   &   &-2x_2 &+  &5x_3  &=  &-b_1+b_3\hfill
             \end{linsys}                                          \\
            &\grstep{2\rho_2+\rho_3}
            \begin{linsys}{3}
              x_1  &+  &2x_2  &+  &3x_3  &=  &b_1\hfill  \\
                   &   &x_2   &-  &3x_3  &=  &-2b_1+b_2\hfill  \\
                   &   &      &   &-x_3  &=  &-5b_1+2b_2+b_3\hfill
             \end{linsys}
         \end{align*}
         表明\( b_1 \)、\( b_2 \)与\( b_3 \)各自可以取任何
         实数\Dash 该方程组总有唯一解。
      \end{exparts}
      
\end{ans}
\begin{ans}{一.I.1.25}
      下面这个未知数多于方程个数的方程组
      \begin{equation*}
        \begin{linsys}{3}
          x  &+  &y  &+  &z  &=  &0  \\
          x  &+  &y  &+  &z  &=  &1
        \end{linsys}
      \end{equation*}
      无解。
      
\end{ans}
\begin{ans}{一.I.1.26}
      是的。
      例如，同一个反应可以在两个不同的烧瓶中进行，
      这一事实表明任何一个解的两倍是另一个不同的解
      （如果物理上确实发生反应，那么至少必须
      有一个非零解）。
    
\end{ans}
\begin{ans}{一.I.1.27}
      由\( f(1)=2 \)、\( f(-1)=6 \)以及\( f(2)=3 \)我们得到
      一个线性方程组。
      \begin{equation*}
        \begin{linsys}{3}
          1a  &+  &1b  &+  &c  &=  &2  \\
          1a  &-  &1b  &+  &c  &=  &6  \\
          4a  &+  &2b  &+  &c  &=  &3
         \end{linsys}
      \end{equation*}
      高斯消元法
      \begin{equation*}
         \grstep[-4\rho_1+\rho_3]{-\rho_1+\rho_2}
         \begin{linsys}{3}
            a  &+  &b  &+  &c  &=  &2  \\
               &   &-2b&   &   &=  &4  \\
               &   &-2b&-  &3c &=  &-5
          \end{linsys}
         \grstep{-\rho_2+\rho_3}
         \begin{linsys}{3}
            a  &+  &b  &+  &c  &=  &2  \\
               &   &-2b&   &   &=  &4  \\
               &   &   &   &-3c&=  &-9
           \end{linsys}
      \end{equation*}
      表明解为\( f(x)=1x^2-2x+3 \)。
      
\end{ans}
\begin{ans}{一.I.1.28}
      这里$S_0=\set{(1,1)}$
      \begin{equation*}
          \begin{linsys}{2}
             x  &+  &y  &=  &2  \\
             x  &-  &y  &=  &0
           \end{linsys}
          \grstep{0\rho_2}
          \begin{linsys}{2}
             x  &+  &y  &=  &2  \\
                &   &0  &=  &0
           \end{linsys}
      \end{equation*}
      而$S_1$是一个真超集，因为它
      至少包含两个点：$(1,1)$和~$(2,0)$。
      在这个例子中解集没有改变。
      \begin{equation*}
          \begin{linsys}{2}
             x  &+  &y  &=  &2  \\
            2x  &+  &2y &=  &4
           \end{linsys}
          \grstep{0\rho_2}
          \begin{linsys}{2}
             x  &+  &y  &=  &2  \\
                &   &0  &=  &0
           \end{linsys}
      \end{equation*}
    
\end{ans}
\begin{ans}{一.I.1.29}
       \begin{exparts}
        \partsitem 能。观察可知，所给方程来自
          \( -\rho_1+\rho_2 \)。
        \partsitem 不能。
          数对\( (1,1) \)满足所给方程。
          然而，这个数对
          不满足该方程组中的第一个方程。
        \partsitem 能。
          要判断所给行是否为\( c_1\rho_1+c_2\rho_2 \)，求解
          联系$x$、$y$、
          $z$的系数与常数的方程组：
          \begin{equation*}
            \begin{linsys}{2}
               2c_1  &+  &6c_2  &=  &6  \\
                c_1  &-  &3c_2  &=  &-9 \\
               -c_1  &+  &c_2   &=  &5  \\
               4c_1  &+  &5c_2  &=  &-2
             \end{linsys}
          \end{equation*}
          得到$c_1=-3$与$c_2=2$，所以所给行是
          \( -3\rho_1+2\rho_2 \)。
      \end{exparts}
     
\end{ans}
\begin{ans}{一.I.1.30}
      若\( a\neq 0 \)，则第一个方程的解集是
      下面这个。
      \begin{equation*}
        \set{(x,y)\in\Re^2\suchthat
             \text{$x=(c-by)/a=(c/a)-(b/a)\cdot y$}}
        \tag{$*$}
      \end{equation*}
      于是，给定~$y$就能算出相应的~$x$。
      取$y=0$得到解$(c/a,0)$；由于第二个
      方程$ax+dy=e$被假定有相同的解集，
      代入其中得到$a(c/a)+d\cdot 0=e$，所以$c=e$。
      在($*$)中取$y=1$得到$a((c-b)/a)+d\cdot 1=e$，
      从而$b=d$。
      因此它们是同一个方程。

      当\( a=0 \)时，两个方程可以不同却仍
      有相同的解集，例如
      \( 0x+3y=6 \)与\( 0x+6y=12 \)。
     
\end{ans}
\begin{ans}{一.I.1.31}
      我们分三种情形：$a\neq 0$；$a=0$而
      $c\neq 0$；以及$a=0$与$c=0$同时成立。

      第一种情形，我们假设\( a\neq 0 \)。
      那么消元
      \begin{equation*}
        \grstep{-(c/a)\rho_1+\rho_2}
        \begin{linsys}{2}
          ax  &+  &by                  &=  &j \hfill\hbox{} \\
              &   &(-(cb/a)+d)y  &=  &-(cj/a)+k \hfill
         \end{linsys}
      \end{equation*}
      表明该方程组有唯一解当且仅当
      \( -(cb/a)+d\neq 0   \)；记住\( a\neq 0 \)，因
      此回代给出唯一的\( x \)
      （顺带指出，\( j \)与\( k \)对“存在唯一解”这一
      结论没有任何影响，不过若存在唯一解，
      则它们会影响这个解的值）。
      但\( -(cb/a)+d = (ad-bc)/a \)，而一个分数等于\( 0 \)
      当且仅当它的分子等于\( 0 \)。
      于是，在第一种情形，存在唯一解当且仅当
      $ad-bc\neq 0$。

      在第二种情形，若\( a=0 \)但\( c\neq 0 \)，则我们对换
      \begin{equation*}
        \begin{linsys}{2}
          cx  &+  &dy  &=  &k  \\
              &   &by  &=  &j
        \end{linsys}
      \end{equation*}
      从而得出：该方程组有唯一解当且仅当
      \( b\neq 0 \)
      （在回代中我们利用情形假设\( c\neq 0 \)来得到唯一的
      \( x \)）。
      但\Dash 在\( a=0 \)而\( c\neq 0 \)时\Dash
      条件“\( b\neq 0 \)”
      等价于条件“\( ad-bc\neq 0 \)”。
      第二种情形到此完成。

      最后，第三种情形，
      若\( a \)与\( c \)都是\( 0 \)，则方程组
      \begin{equation*}
        \begin{linsys}{2}
          0x  &+  &by  &=  &j  \\
          0x  &+  &dy  &=  &k
        \end{linsys}
      \end{equation*}
      可能无解（如果第二个方程不是第一个方程的倍式），
      也可能有无穷多解（如果第二个方程是第一个方程的倍式，
      那么对每个满足两个方程的\( y \)，任何数对\( (x,y) \)都行），
      但它永远不会有唯一解。
      注意\( a=0 \)与\( c=0 \)给出\( ad-bc=0 \)。
    
\end{ans}
\begin{ans}{一.I.1.32}
      回忆一下，若两条直线有两个不同的公共点，则
      它们是同一条直线。
      这是因为两点确定一条直线，所以这两个点
      同时确定了这两条直线，
      因而它们是同一条直线。

      于是，两条直线可以有一个公共点（给出唯一解）、
      没有公共点（给出无解），
      或者至少有两个公共点（这使它们成为同一条直线）。
    
\end{ans}
\begin{ans}{一.I.1.33}
     对于用非零
     实数$k$乘$\rho_i$的消元操作，我们有：\( (s_1,\ldots,s_n) \)满足
     方程组
     \begin{equation*}
       \begin{linsys}{4}
         a_{1,1}x_1  &+  &a_{1,2}x_2 &+  &\cdots  &+  &a_{1,n}x_n  &=  &d_1  \\
                   &   &          &  &        &   &           &\vdotswithin{=}   \\
        ka_{i,1}x_1  &+  &ka_{i,2}x_2 &+  &\cdots  &+  &ka_{i,n}x_n &=  &kd_i  \\
                   &   &           &   &        &   &            &\vdotswithin{=}   \\
         a_{m,1}x_1  &+  &a_{m,2}x_2 &+  &\cdots  &+  &a_{m,n}x_n  &=  &d_m
       \end{linsys}
     \end{equation*}
     当且仅当
     % \begin{align*}
     %    a_{1,1}s_1+a_{1,2}s_2+\cdots+a_{1,n}s_n
     %    &=d_1                                              \\
     %    &\alignedvdots                                     \\
     %    \text{and\ } ka_{i,1}s_1+ka_{i,2}s_2+\cdots+ka_{i,n}s_n
     %    &=kd_i                                              \\
     %    &\alignedvdots                                      \\
     %    \text{and\ } a_{m,1}s_1+a_{m,2}s_2+\cdots+a_{m,n}s_n
     %    &=d_m
     % \end{align*}
     $a_{1,1}s_1+a_{1,2}s_2+\cdots+a_{1,n}s_n=d_1$
     且\ldots{} $ka_{i,1}s_1+ka_{i,2}s_2+\cdots+ka_{i,n}s_n=kd_i$
     且\ldots{}
     $a_{m,1}s_1+a_{m,2}s_2+\cdots+a_{m,n}s_n=d_m$，
     这由“满足”的定义得出。
     因为\( k\neq 0 \)，上式成立当且仅当
     % \begin{align*}
     %    a_{1,1}s_1+a_{1,2}s_2+\cdots+a_{1,n}s_n
     %    &=d_1                                              \\
     %    &\alignedvdots                                     \\
     %    \text{and\ } a_{i,1}s_1+a_{i,2}s_2+\cdots+a_{i,n}s_n
     %    &=d_i                                              \\
     %    &\alignedvdots                                      \\
     %    \text{and\ } a_{m,1}s_1+a_{m,2}s_2+\cdots+a_{m,n}s_n
     %    &=d_m
     % \end{align*}
     $a_{1,1}s_1+a_{1,2}s_2+\cdots+a_{1,n}s_n=d_1$
     且\ldots{}
     $a_{i,1}s_1+a_{i,2}s_2+\cdots+a_{i,n}s_n=d_i$
     且\ldots{}
     $a_{m,1}s_1+a_{m,2}s_2+\cdots+a_{m,n}s_n=d_m$
     （这只是把第$i$个方程两边直接消去$k$），
     这说明\( (s_1,\ldots,s_n) \)解方程组
     \begin{equation*}
       \begin{linsys}{4}
         a_{1,1}x_1  &+  &a_{1,2}x_2 &+  &\cdots  &+  &a_{1,n}x_n  &=  &d_1  \\
                     &   &           &   &        &   &            &\vdotswithin{=}   \\
         a_{i,1}x_1  &+  &a_{i,2}x_2 &+  &\cdots  &+  &a_{i,n}x_n  &=  &d_i  \\
                     &   &           &   &        &   &            &\vdotswithin{=}   \\
         a_{m,1}x_1  &+  &a_{m,2}x_2 &+  &\cdots  &+  &a_{m,n}x_n  &=
              &d_m
         \end{linsys}
     \end{equation*}
     即为所需。

     对于组合操作$k\rho_i+\rho_j$，元组
     \( (s_1,\ldots,s_n) \)满足
     \begin{equation*}
       \begin{linsys}{4}
         a_{1,1}x_1             &+  &\cdots  &+  &a_{1,n}x_n  &=  &d_1\hfill\hbox{} \\
                                &   &        &   &            &\vdotswithin{=}   \\
         a_{i,1}x_1             &+  &\cdots  &+  &a_{i,n}x_n  &=  &d_i\hfill\hbox{} \\
                                &   &        &   &            &\vdotswithin{=}   \\
         (ka_{i,1}+a_{j,1})x_1  &+  &\cdots  &+  &(ka_{i,n}+a_{j,n})x_n
               &=  &kd_i+d_j \hfill \\
                                &   &        &   &            &\vdotswithin{=}   \\
         a_{m,1}x_1             &+   &\cdots  &+  &a_{m,n}x_n  &=
          &d_m\hfill\hbox{}
        \end{linsys}
     \end{equation*}
     当且仅当
     $a_{1,1}s_1+\cdots+a_{1,n}s_n=d_1$
     且\ldots{}
     $a_{i,1}s_1+\cdots+a_{i,n}s_n=d_i$
     且\ldots{}
     $(ka_{i,1}+a_{j,1})s_1+\cdots+(ka_{i,n}+a_{j,n})s_n=kd_i+d_j$
     且\ldots{}
     $a_{m,1}s_1+a_{m,2}s_2+\cdots+a_{m,n}s_n=d_m$，
     同样由“满足”的定义得出。
     从方程~\( j \)中减去\( k \)倍的方程~\( i \)。
     （这正是需要\( i\neq j \)的地方；若\( i=j \)，则上面
     两个\( d_i \)不相等。）
     前面的复合命题成立当且仅当
     $a_{1,1}s_1+\cdots+a_{1,n}s_n=d_1$
     且\ldots{}
     $a_{i,1}s_1+\cdots+a_{i,n}s_n=d_i$
     且\ldots{} $
     (ka_{i,1}+a_{j,1})s_1+\cdots+(ka_{i,n}+a_{j,n})s_n
         -(ka_{i,1}s_1+\cdots+ka_{i,n}s_n)
         =kd_i+d_j-kd_i$
      且\ldots{} $a_{m,1}s_1+\cdots+a_{m,n}s_n=d_m$，
     消去之后，这说明\( (s_1,\ldots,s_n) \)解方程组
     \begin{equation*}
       \begin{linsys}{4}
         a_{1,1}x_1  &+   &\cdots  &+  &a_{1,n}x_n  &=  &d_1  \\
                     &    &        &   &            &\vdotswithin{=}   \\
         a_{i,1}x_1  &+   &\cdots  &+  &a_{i,n}x_n  &=  &d_i  \\
                     &    &        &   &            &\vdotswithin{=}   \\
         a_{j,1}x_1  &+  &\cdots  &+  &a_{j,n}x_n  &=  &d_j  \\
                     &   &        &   &            &\vdotswithin{=}   \\
         a_{m,1}x_1  &+  &\cdots  &+  &a_{m,n}x_n  &=
              &d_m\hfill\hbox{}
       \end{linsys}
     \end{equation*}
     即为所需。
   
\end{ans}
\begin{ans}{一.I.1.34}
      存在，下面这个只有一个方程的方程组：
      \begin{equation*}
         0x+0y=0
      \end{equation*}
      每个\( (x,y)\in\Re^2 \)都满足它。
    
\end{ans}
\begin{ans}{一.I.1.35}
      有。
      下面这串操作对换了行\( i \)与行\( j \)
      \begin{equation*}
         \grstep{\rho_i+\rho_j}
         \repeatedgrstep{-\rho_j+\rho_i}
         \repeatedgrstep{\rho_i+\rho_j}
         \repeatedgrstep{-1\rho_i}
      \end{equation*}
      所以在另外两种操作存在的情况下，对换操作是多余的。
     
\end{ans}
\begin{ans}{一.I.1.36}
      行对换$\rho_i\leftrightarrow\rho_j$可以用
      对换回去$\rho_j\leftrightarrow\rho_i$来逆转。
      $k\rho_i$这一步是用\( k\neq 0  \)乘
      某一行的两边，
      可以用除以~$(1/k)\rho_i$来逆转。

      倍加的情形是非平凡的。
      操作$k\rho_i+\rho_j$
      得到如下的第$j$行。
      \begin{equation*}
        k\cdot a_{i,1}+a_{j,1}+\cdots+k\cdot a_{i,n}+a_{j,n}=k\cdot d_i+d_j
      \end{equation*}
      由于$i\neq j$的限制，第$i$行不变。
      因为第$i$行不变，操作$-k\rho_i+\rho_j$
      把第$j$行变回原状。

      （注意$k\rho_i+\rho_j$中\( i=j \)的条件
       是必需的，否则可能出现
       \begin{equation*}
         \begin{linsys}{2}
           3x  &+  &2y  &=  &7
         \end{linsys}
         \grstep{2\rho_1+\rho_1}
         \begin{linsys}{2}
           9x  &+  &6y  &=  &21
          \end{linsys}
         \grstep{-2\rho_1+\rho_1}
         \begin{linsys}{2}
          -9x  &-  &6y  &=  &-21
         \end{linsys}
       \end{equation*}
       于是结论不成立。）
    
\end{ans}
\begin{ans}{一.I.1.37}
      设\( p \)、\( n \)、\( d \)分别是
      1美分、5美分和10美分硬币的个数。
      当变量取实数值时，方程组
      \begin{equation*}
         \begin{linsys}{3}
              p  &+ &n   &+  &d   &=  &13   \\
              p  &+ &5n  &+  &10d &=  &83
         \end{linsys}
         \grstep{-\rho_1+\rho_2}
         \begin{linsys}{3}
              p  &+ &n   &+  &d   &=  &13   \\
                 &  &4n  &+  &9d  &=  &70
          \end{linsys}
      \end{equation*}
      的解不止一个；事实上，它有无穷多个解。
      然而，\( p \)、\( n \)、
      \( d \)为非负整数的解只有有限多个。
      逐一检查\( d=0 \)，\ldots，\( d=8 \)可知
      \( (p,n,d)=(3,4,6) \)
      是唯一使用自然数的解。
    
\end{ans}
\begin{ans}{一.I.1.38}
      求解方程组
      \begin{equation*}
        \begin{linsys}{2}
        (1/3)(a+b+c)  &+  &d  &=  &29  \\
        (1/3)(b+c+d)  &+  &a  &=  &23  \\
        (1/3)(c+d+a)  &+  &b  &=  &21  \\
        (1/3)(d+a+b)  &+  &c  &=  &17
        \end{linsys}
      \end{equation*}
      我们得到$a=12$，$b=9$，$c=3$，$d=21$。
      因此第二项，21，是正确答案。
     
\end{ans}
\begin{ans}{一.I.1.39}
        \answerasgiven
        比较这个加法竖式的个位列与百位列可知，
        十位列必有进位。
        十位列于是告诉我们\( A<H \)，所以
        个位列与百位列都不会有进位。
        五个列于是给出下面五个方程。
        \begin{align*}
          A+E  &=  W  \\
          2H   &=  A+10  \\
          H    &=  W+1  \\
          H+T  &=  E+10  \\
          A+1  &=  T
        \end{align*}
        这五个未知数的五个线性方程若联立求解，
        给出唯一解：\( A=4 \)，\( T=5 \)，\( H=7 \)，
        \( W=6 \)与\( E=2 \)，于是原来的加法算式
        是\( 47474+5272=52746 \)。
      
\end{ans}
\begin{ans}{一.I.1.40}
       \answerasgiven{}
       \textit{其中一些补充材料取自\cite{joriki}。}
       有八名委员投票给$B$。
       为说明这一点，我们将利用已知信息来研究有多少投票者
       选择了$A$、$B$、$C$的每种次序。

       六种偏好次序是$ABC$、$ACB$、$BAC$、$BCA$、$CAB$、
       $CBA$；假设它们分别获得$a$、$b$、$c$、$d$、$e$、$f$票。
       我们知道
       \begin{equation*}
         \begin{linsys}{3}
           a  &+  &b  &+  &e  &=  &11  \\
           d  &+  &e  &+  &f  &=  &12  \\
           a  &+  &c  &+  &d  &=  &14
         \end{linsys}
       \end{equation*}
       它们分别来自偏好$A$胜过$B$的人数、偏好
       $C$胜过$A$的人数，以及偏好$B$胜过$C$的人数。
       因为总共投了20票，
       所以$a+b+\cdots+f=20$。
       从两倍的前一个方程中减去上面三个方程之和，
       得到$b+c+f=3$。
       我们已规定每种偏好次序至少获得
       一票，所以这意味着$b=c=f=1$。

       由上面三个方程，完整解为
       $a=6$，$b=1$，$c=1$，$d=7$，$e=4$以及~$f=1$，
       这可以用高斯消元法求出。
       把$B$列为第一选择的委员人数
       因此是$c+d=1+7=8$。

       \par\noindent {\em 评注。}
       即使我们只知道{\em 至少\/}有14名委员偏好$B$胜过$C$，
       这个问题的答案也会一样。

       委员们的偏好看似悖论的性质
       （$A$优于$B$，$B$优于$C$，而$C$又
       优于$A$）是“非传递支配”的一个例子，
       这种情形在汇总个人选择时很常见。
     
\end{ans}
\begin{ans}{一.I.1.41}
       \answerasgiven
       \textit{我们还没有用到“相依”一词；
       它在这里的意思是高斯
       消元法表明解不唯一。}
       若\( n\geq 3 \)，则方程组相依且解不
       唯一。
       因此\( n<3 \)。
       但“方程组”一词意味着\( n>1 \)。
       因此\( n=2 \)。
       若方程为
       \begin{equation*}
         \begin{linsys}{2}
              ax  &+ &(a+d)y  &=  &a+2d  \\
         (a+3d)x  &+ &(a+4d)y &=  &a+5d
         \end{linsys}
       \end{equation*}
       则\( x=-1 \)，\( y=2 \)。
    
\end{ans}
\protect \subsection {一.I.2: 解集的描述}
\begin{ans}{一.I.2.15}
      \begin{exparts*}
        \partsitem \( 2 \)
        \partsitem \( 3 \)
        \partsitem \(-1 \)
        \partsitem 无定义。
      \end{exparts*}
    
\end{ans}
\begin{ans}{一.I.2.16}
      \begin{exparts*}
        \partsitem \( \nbym{2}{3} \)
        \partsitem \( \nbym{3}{2} \)
        \partsitem \( \nbym{2}{2} \)
      \end{exparts*}
    
\end{ans}
\begin{ans}{一.I.2.17}
      \begin{exparts*}
        \partsitem \( \colvec[r]{5 \\ 1 \\ 5} \)
        \partsitem \( \colvec[r]{20 \\ -5} \)
        \partsitem \( \colvec[r]{-2 \\ 4 \\ 0} \)
        \partsitem \( \colvec[r]{41 \\ 52} \)
        \partsitem 无定义。
        \partsitem \( \colvec[r]{12 \\ 8 \\ 4} \)
      \end{exparts*}
     
\end{ans}
\begin{ans}{一.I.2.18}
      \begin{exparts}
        \partsitem 这一消元
          \begin{equation*}
            \begin{amat}[r]{2}
              3  &6  &18 \\
              1  &2  &6
            \end{amat}
            \grstep{(-1/3)\rho_1+\rho_2}
            \begin{amat}[r]{2}
              3  &6  &18 \\
              0  &0  &0
            \end{amat}
          \end{equation*}
          留下\( x \)为首变量而\( y \)为自由变量。
          取\( y \)为参数，得到\( x=6-2y \)以及这个解
          集。
          \begin{equation*}
            \set{\colvec[r]{6 \\ 0}+\colvec[r]{-2 \\ 1}y
              \suchthat y\in\Re}
          \end{equation*}
        \partsitem 一次消元
          \begin{equation*}
            \begin{amat}[r]{2}
              1  &1  &1  \\
              1  &-1 &-1
            \end{amat}
            \grstep{-\rho_1+\rho_2}
            \begin{amat}[r]{2}
              1  &1  &1  \\
              0  &-2 &-2
            \end{amat}
          \end{equation*}
          给出唯一解\( y=1 \)、\( x=0 \)。
          解集如下。
          \begin{equation*}
            \set{\colvec[r]{0 \\ 1} }
          \end{equation*}
        \partsitem 高斯消元法
          \begin{equation*}
            \begin{amat}[r]{3}
              1  &0  &1  &4  \\
              1  &-1 &2  &5  \\
              4  &-1 &5  &17
            \end{amat}
            \grstep[-4\rho_1+\rho_3]{-\rho_1+\rho_2}
            \begin{amat}[r]{3}
              1  &0  &1  &4  \\
              0  &-1 &1  &1  \\
              0  &-1 &1  &1
            \end{amat}
            \grstep{-\rho_2+\rho_3}
            \begin{amat}[r]{3}
              1  &0  &1  &4  \\
              0  &-1 &1  &1  \\
              0  &0  &0  &0
            \end{amat}
          \end{equation*}
          留下\( x_1 \)和\( x_2 \)为首变量，而\( x_3 \)为自由变量。
          解集如下。
          \begin{equation*}
            \set{\colvec[r]{4 \\ -1 \\ 0}+\colvec[r]{-1 \\ 1 \\ 1}x_3
              \suchthat x_3\in\Re}
          \end{equation*}
        \partsitem 这一消元
          \begin{align*}
            \begin{amat}[r]{3}
              2  &1  &-1 &2  \\
              2  &0  &1  &3  \\
              1  &-1 &0  &0
            \end{amat}
            &\grstep[-(1/2)\rho_1+\rho_3]{-\rho_1+\rho_2}
            \begin{amat}[r]{3}
              2  &1    &-1   &2  \\
              0  &-1   &2    &1  \\
              0  &-3/2 &1/2  &-1
            \end{amat}                                          \\
            &\grstep{(-3/2)\rho_2+\rho_3}
            \begin{amat}[r]{3}
              2  &1  &-1   &2  \\
              0  &-1 &2    &1  \\
              0  &0  &-5/2 &-5/2
            \end{amat}
          \end{align*}
          表明解集是单元素集。
          \begin{equation*}
            \set{\colvec[r]{1 \\ 1 \\ 1}}
          \end{equation*}
        \partsitem 这一消元很容易
          \begin{align*}
            \begin{amat}[r]{4}
              1  &2  &-1 &0  &3 \\
              2  &1  &0  &1  &4 \\
              1  &-1 &1  &1  &1
            \end{amat}
            &\grstep[-\rho_1+\rho_3]{-2\rho_1+\rho_2}
            \begin{amat}[r]{4}
              1  &2  &-1 &0  &3  \\
              0  &-3 &2  &1  &-2 \\
              0  &-3 &2  &1  &-2
            \end{amat}                                       \\
            &\grstep{-\rho_2+\rho_3}
            \begin{amat}[r]{4}
              1  &2  &-1 &0  &3  \\
              0  &-3 &2  &1  &-2 \\
              0  &0  &0  &0  &0
            \end{amat}
          \end{align*}
          结束时\( x \)和$y$为首变量，而\( z \)和\( w \)为
          自由变量。
          解出\( y \)得到\( y=(2+2z+w)/3 \)，代入表明
          \( x+2(2+2z+w)/3-z=3 \)，于是\( x=(5/3)-(1/3)z-(2/3)w \)，
          由此得到这个解集。
          \begin{equation*}
            \set{\colvec[r]{5/3 \\ 2/3 \\ 0 \\ 0}
                 +\colvec[r]{-1/3 \\ 2/3 \\ 1 \\ 0}z
                 +\colvec[r]{-2/3 \\ 1/3 \\ 0 \\ 1}w
                 \suchthat z,w\in\Re}
          \end{equation*}
        \partsitem 消元
          \begin{align*}
            \begin{amat}[r]{4}
              1  &0  &1  &1  &4 \\
              2  &1  &0  &-1 &2 \\
              3  &1  &1  &0  &7
            \end{amat}
            &\grstep[-3\rho_1+\rho_3]{-2\rho_1+\rho_2}
            \begin{amat}[r]{4}
              1  &0  &1  &1  &4 \\
              0  &1  &-2 &-3 &-6\\
              0  &1  &-2 &-3 &-5
            \end{amat}                                       \\
            &\grstep{-\rho_2+\rho_3}
            \begin{amat}[r]{4}
              1  &0  &1  &1  &4 \\
              0  &1  &-2 &-3 &-6\\
              0  &0  &0  &0  &1
            \end{amat}
          \end{align*}
          表明无解\Dash 解集为空。
      \end{exparts}
     
\end{ans}
\begin{ans}{一.I.2.19}
      \begin{exparts}
      \partsitem 这一消元
        \begin{equation*}
          \begin{amat}[r]{3}
            2  &1  &-1  &1  \\
            4  &-1 &0   &3
          \end{amat}
          \grstep{-2\rho_1+\rho_2}
          \begin{amat}[r]{3}
            2  &1  &-1  &1  \\
            0  &-3 &2   &1
          \end{amat}
        \end{equation*}
        结束时\( x \)和\( y \)为首变量，而\( z \)为自由变量。
        解出\( y \)得到\( y=(1-2z)/(-3) \)，然后代入
        \( 2x+(1-2z)/(-3)-z=1 \)表明\( x=((4/3)+(1/3)z)/2 \)。
        因此解集如下。
        \begin{equation*}
          \set{\colvec[r]{2/3 \\ -1/3 \\ 0}
               +\colvec[r]{1/6 \\ 2/3 \\ 1}z
              \suchthat z\in\Re}
        \end{equation*}
      \partsitem 这样应用高斯消元法
        \begin{align*}
          \begin{amat}[r]{4}
            1  &0  &-1  &0  &1 \\
            0  &1  &2   &-1 &3 \\
            1  &2  &3   &-1 &7
          \end{amat}
          &\grstep{-\rho_1+\rho_3}
          \begin{amat}[r]{4}
            1  &0  &-1  &0  &1 \\
            0  &1  &2   &-1 &3 \\
            0  &2  &4   &-1 &6
          \end{amat}                                   \\
          &\grstep{-2\rho_2+\rho_3}
          \begin{amat}[r]{4}
            1  &0  &-1  &0  &1 \\
            0  &1  &2   &-1 &3 \\
            0  &0  &0   &1  &0
          \end{amat}
        \end{align*}
        留下\( x \)、\( y \)和\( w \)为首变量。
        解集如下。
        \begin{equation*}
          \set{\colvec[r]{1 \\ 3 \\ 0 \\ 0}
               +\colvec[r]{1 \\ -2 \\ 1 \\ 0}z
              \suchthat z\in\Re}
        \end{equation*}
      \partsitem 这次行化简
        \begin{align*}
          \begin{amat}[r]{4}
            1  &-1 &1   &0  &0 \\
            0  &1  &0   &1  &0 \\
            3  &-2 &3   &1  &0 \\
            0  &-1 &0   &-1 &0
          \end{amat}
          &\grstep{-3\rho_1+\rho_3}
          \begin{amat}[r]{4}
            1  &-1 &1   &0  &0 \\
            0  &1  &0   &1  &0 \\
            0  &1  &0   &1  &0 \\
            0  &-1 &0   &-1 &0
          \end{amat}                                       \\
          &\grstep[\rho_2+\rho_4]{-\rho_2+\rho_3}
          \begin{amat}[r]{4}
            1  &-1 &1   &0  &0 \\
            0  &1  &0   &1  &0 \\
            0  &0  &0   &0  &0 \\
            0  &0  &0   &0  &0
          \end{amat}
        \end{align*}
        结束时\( z \)和\( w \)为自由变量。
        我们得到这个解集。
        \begin{equation*}
          \set{\colvec[r]{0 \\ 0 \\ 0 \\ 0}
               +\colvec[r]{-1 \\ 0 \\ 1 \\ 0}z
               +\colvec[r]{-1 \\ -1 \\ 0 \\ 1}w
              \suchthat z,w\in\Re}
        \end{equation*}
      \partsitem 按这种方式做高斯消元法
        \begin{equation*}
          \begin{amat}[r]{5}
            1  &2  &3   &1  &-1 &1  \\
            3  &-1 &1   &1  &1  &3
          \end{amat}
          \grstep{-3\rho_1+\rho_2}
          \begin{amat}[r]{5}
            1  &2  &3   &1  &-1 &1  \\
            0  &-7 &-8  &-2 &4  &0
          \end{amat}
        \end{equation*}
        结束时\( c \)、\( d \)和\( e \)为自由变量。
        解出\( b \)表明\( b=(8c+2d-4e)/(-7) \)，然后
        代入
        \( a+2(8c+2d-4e)/(-7)+3c+1d-1e=1 \)表明
        \( a=1-(5/7)c-(3/7)d-(1/7)e \)，于是我们得到解集。
        \begin{equation*}
          \set{\colvec[r]{1 \\ 0 \\ 0 \\ 0 \\ 0}
               +\colvec[r]{-5/7 \\ -8/7 \\ 1 \\ 0 \\ 0}c
               +\colvec[r]{-3/7 \\ -2/7 \\ 0 \\ 1 \\ 0}d
               +\colvec[r]{-1/7 \\ 4/7 \\ 0 \\ 0 \\ 1}e
              \suchthat c,d,e\in\Re}
        \end{equation*}
    \end{exparts}
   
\end{ans}
\begin{ans}{一.I.2.20}
      \begin{exparts}
        \partsitem 这一消元
          \begin{align*}
            \begin{amat}{3}
              3 &2  &1  &1 \\
              1 &-1 &1  &2 \\
              5 &5  &1  &0
            \end{amat}
            &\grstep[-(5/3)\rho_1+\rho_3]{-(1/3)\rho_1+\rho_2}
            \begin{amat}{3}
              3 &2    &1     &1 \\
              0 &-5/3 &2/3   &5/3 \\
              0 &5/3  &-2/3  &-5/3
            \end{amat}                                     \\
            &\grstep{\rho_2+\rho_3}
            \begin{amat}{3}
              3 &2    &1     &1 \\
              0 &-5/3 &2/3   &5/3 \\
              0 &0    &0     &0
            \end{amat}
          \end{align*}
          给出这个解集。
          \begin{equation*}
            \set{\colvec{x \\ y \\ z}=\colvec{1 \\ -1 \\ 0}
                                       +\colvec{-3/5 \\ 2/5 \\ 1}z
                             \suchthat z\in\Re}
          \end{equation*}
        \partsitem
          消元如下。
          \begin{align*}
            \begin{amat}{3}
              1 &1   &-2 &0 \\
              1 &-1  &0  &3  \\
              3 &-1  &-2 &-6 \\
              0 &2   &-2 &3
            \end{amat}
            &\grstep[-3\rho_1+\rho_3]{-\rho_1+\rho_2}
            \begin{amat}{3}
              1 &1   &-2 &0 \\
              0 &-2  &2  &-3  \\
              0 &-4  &4  &-6 \\
              0 &2   &-2 &3
            \end{amat}                                  \\
            &\grstep[\rho_2+\rho_4]{-2\rho_2+\rho_3}
            \begin{amat}{3}
              1 &1   &-2 &0 \\
              0 &-2  &2  &-3  \\
              0 &0   &0  &0 \\
              0 &0   &0  &0
            \end{amat}
          \end{align*}
          解集如下。
          \begin{equation*}
            \set{\colvec{-3/2 \\ 3/2 \\ 0}
                 +\colvec{1 \\ 1 \\ 1}z
                 \suchthat z\in\Re}
          \end{equation*}
      \partsitem
         高斯消元法
         \begin{equation*}
           \begin{amat}{4}
             2 &-1 &-1 &1 &4 \\
             1 &1  &1  &0 &-1
           \end{amat}
           \grstep{-(1/2)\rho_1+\rho_2}
           \begin{amat}{4}
             2 &-1   &-1   &1    &4 \\
             0 &3/2  &3/2  &-1/2 &-3
           \end{amat}
         \end{equation*}
         给出解集。
         \begin{equation*}
           \set{\colvec{1 \\ -2 \\ 0 \\ 0}
                  +\colvec{0 \\ -1 \\ 1 \\ 0}z
                   \colvec{-1/3 \\ 1/3 \\ 0 \\ 1}w
                 \suchthat z,w\in\Re}
         \end{equation*}
       \partsitem
          消元如下。
          \begin{equation*}
            \begin{amat}{3}
              1 &1  &-2 &0  \\
              1 &-1 &0  &-3 \\
              3 &-1 &-2 &0
           \end{amat}
           \grstep[-3\rho_1+\rho_3]{-\rho_1+\rho_2}
           \begin{amat}{3}
             1 &1  &-2 &0  \\
             0 &-2 &2  &-3 \\
             0 &-4  &4  &0
          \end{amat}
          \grstep{-2\rho_2+\rho_3}
          \begin{amat}{3}
            1 &1  &-2 &0  \\
            0 &-2 &2  &-3 \\
            0 &0  &0  &6
          \end{amat}
        \end{equation*}
        解集为空集~$\set{}$。
      \end{exparts}
    
\end{ans}
\begin{ans}{一.I.2.21}
      对每一小题，我们通过考察分量的等式得到一个
      线性方程组。
      \begin{exparts}
       \partsitem $k=5$
       \partsitem 第二个分量表明$i=2$，第三个
       分量表明$j=1$。
       \partsitem $m=-4$，$n=2$
      \end{exparts}
    
\end{ans}
\begin{ans}{一.I.2.22}
      对每一小题，我们通过考察分量的等式得到一个
      线性方程组。
      \begin{exparts}
        \partsitem 是；取$k=-1/2$。
        \partsitem 否；由方程$5=5\cdot j$和
            $4=-4\cdot j$组成的方程组无解。
        \partsitem 是；取$r=2$。
        \partsitem 否。
           第二个分量给出$k=0$。
           于是第三个分量给出$j=1$。
           但第一个分量对不上。
      \end{exparts}
     
\end{ans}
\begin{ans}{一.I.2.23}
      \begin{exparts}
        \item 设$c$为玉米的英亩数，$s$为
          大豆的英亩数，$a$为燕麦的英亩数。
          \begin{equation*}
            \begin{linsys}{3}
              c   &+   &s   &+   &a   &=   &1200 \\
            20c   &+   &50s &+   &12a &=   &40\,000
            \end{linsys}
            \grstep{-20\rho_1+\rho_2}
            \begin{linsys}{3}
              c   &+   &s   &+   &a   &=   &1200 \\
                  &    &30s &-   &8a  &=   &16\,000
            \end{linsys}
          \end{equation*}
          为描述解集，我们可以用$a$参数化。
          \begin{equation*}
            \set{\colvec{c \\ s \\ a}
                 =\colvec{20\,000/30 \\ 16\,000/30 \\ 0}
                  +\colvec{-38/30 \\ 8/30 \\ 1}a
                 \suchthat a\in\Re}
          \end{equation*}
        \item 这里有许多可能的答案。
          例如我们可以取$a=0$得到$c=20\,000/30\approx 666.66$和
          $s=16000/30\approx 533.33$。
          另一个例子是取$a=20\,000/38\approx 526.32$，给出
          $c=0$和$s=7360/38\approx 193.68$。
        \item 把前一题的答案代入
          $100c+300s+80a$。
      \end{exparts}
    
\end{ans}
\begin{ans}{一.I.2.24}
      这个方程组有一个方程。
      首变量是\( x_1 \)，其余变量是自由变量。
      \begin{equation*}
        \set{\colvec{-1 \\ 1 \\ \vdotswithin{-1} \\ 0}x_2
             +\cdots+
             \colvec{-1 \\ 0 \\ \vdotswithin{-1} \\ 1}x_n
             \suchthat x_2,\ldots,x_n\in\Re}
      \end{equation*}
     
\end{ans}
\begin{ans}{一.I.2.25}
      \begin{exparts}
        \partsitem 这里高斯消元法给出
          \begin{align*}
            \begin{amat}{4}
              1  &2  &0  &-1  &a  \\
              2  &0  &1  &0   &b  \\
              1  &1  &0  &2   &c
            \end{amat}
            &\grstep[-\rho_1+\rho_3]{-2\rho_1+\rho_2}
            \begin{amat}{4}
              1  &2  &0  &-1  &a  \\
              0  &-4 &1  &2   &-2a+b  \\
              0  &-1 &0  &3   &-a+c
            \end{amat}                                  \\
            &\grstep{-(1/4)\rho_2+\rho_3}
            \begin{amat}{4}
              1  &2  &0    &-1  &a  \\
              0  &-4 &1    &2   &-2a+b  \\
              0  &0  &-1/4 &5/2 &-(1/2)a-(1/4)b+c
            \end{amat}
          \end{align*}
          留下\( w \)为自由变量。
          求解：\( z=2a+b-4c+10w \)，
          且\( -4y=-2a+b-(2a+b-4c+10w)-2w \)于是
          \( y=a-c+3w \)，并且
          \( x=a-2(a-c+3w)+w=-a+2c-5w. \)
          因此解集如下。
          \begin{equation*}
             \set{\colvec{-a+2c \\ a-c \\ 2a+b-4c \\ 0}
                  +\colvec{-5 \\ 3 \\ 10 \\ 1}w
                  \suchthat w\in\Re}
          \end{equation*}
        \partsitem 代入\( a=3 \)、\( b=1 \)和\( c=-2 \)。
          \begin{equation*}
             \set{\colvec[r]{-7 \\ 5 \\ 15 \\ 0}
                  +\colvec[r]{-5 \\ 3 \\ 10 \\ 1}w
                  \suchthat w\in\Re}
          \end{equation*}
      \end{exparts}
     
\end{ans}
\begin{ans}{一.I.2.26}
       省略逗号，比如写成\( a_{123} \)，
       是有歧义的，因为它可能指$a_{1,23}$或$a_{12,3}$。
    
\end{ans}
\begin{ans}{一.I.2.27}
      \begin{exparts*}
        \partsitem \(
           \begin{mat}[r]
             2  &3  &4  &5  \\
             3  &4  &5  &6  \\
             4  &5  &6  &7  \\
             5  &6  &7  &8
           \end{mat} \)
        \partsitem \(
           \begin{mat}[r]
             1  &-1  &1   &-1  \\
            -1  &1   &-1  &1  \\
             1  &-1  &1   &-1  \\
            -1  &1   &-1  &1
           \end{mat} \)
      \end{exparts*}
    
\end{ans}
\begin{ans}{一.I.2.28}
      \begin{exparts*}
        \partsitem \( \begin{mat}[r]
                   1  &4  \\
                   2  &5  \\
                   3  &6
                 \end{mat}  \)
        \partsitem \( \begin{mat}[r]
                   2  &1  \\
                  -3  &1
                 \end{mat}  \)
        \partsitem \( \begin{mat}[r]
                   5  &10 \\
                  10  &5
                 \end{mat}  \)
        \partsitem \( \rowvec{1 &1 &0}  \)
      \end{exparts*}
     
\end{ans}
\begin{ans}{一.I.2.29}
      \begin{exparts}
        \partsitem 代入\( x=1 \)和\( x=-1 \)得到
          \begin{equation*}
            \begin{linsys}{3}
              a  &+  &b   &+  &c  &=  &2  \\
              a  &-  &b   &+  &c  &=  &6
            \end{linsys}
            \grstep{-\rho_1+\rho_2}
            \begin{linsys}{3}
              a  &+  &b   &+  &c  &=  &2  \\
                 &   &-2b &   &   &=  &4
              \end{linsys}
          \end{equation*}
          所以函数的集合是
          \( \set{f(x)=(4-c)x^2-2x+c\suchthat c\in\Re} \)。
        \partsitem 代入\( x=1 \)得到
          \begin{equation*}
            \begin{linsys}{3}
              a  &+  &b   &+  &c  &=  &2
            \end{linsys}
          \end{equation*}
          所以函数的集合是
          \( \set{f(x)=(2-b-c)x^2+bx+c\suchthat b,c\in\Re} \)。
      \end{exparts}
    
\end{ans}
\begin{ans}{一.I.2.30}
      把五对$(x,y)$代入，我们得到一个有
      五个方程、六个未知数$a$、\ldots、$f$的方程组。
      由于未知数比方程多，如果方程之间没有不相容，
      那么就有无穷多解
      （至少有一个变量最终会成为自由变量）。

      但不相容不可能存在，因为$a=0$、\ldots、$f=0$是
      一个解（我们只是用这个零解来表明方程组
      是相容的\Dash 前面一段表明
      存在非零解）。
    
\end{ans}
\begin{ans}{一.I.2.31}
      \begin{exparts}
      \partsitem 这是一个\Dash 第四个方程是多余的
        但仍然可以。
        \begin{equation*}
          \begin{linsys}{4}
             x  &+  &y  &-  &z  &+  &w  &=  &0  \\
                &   &y  &-  &z  &   &   &=  &0  \\
                &   &   &   &2z &+  &2w &=  &0  \\
                &   &   &   &z  &+  &w  &=  &0
          \end{linsys}
        \end{equation*}
      \partsitem 这是一个。
        \begin{equation*}
          \begin{linsys}{4}
             x  &+  &y  &-  &z  &+  &w  &=  &0  \\
                &   &   &   &   &   &w  &=  &0  \\
                &   &   &   &   &   &w  &=  &0  \\
                &   &   &   &   &   &w  &=  &0
          \end{linsys}
        \end{equation*}
      \partsitem 这是一个。
        \begin{equation*}
          \begin{linsys}{4}
             x  &+  &y  &-  &z  &+  &w  &=  &0  \\
             x  &+  &y  &-  &z  &+  &w  &=  &0  \\
             x  &+  &y  &-  &z  &+  &w  &=  &0  \\
             x  &+  &y  &-  &z  &+  &w  &=  &0
          \end{linsys}
        \end{equation*}
    \end{exparts}
   
\end{ans}
\begin{ans}{一.I.2.32}
        \answerasgiven
        我的解法是把arbuzoid的
        数量定义为$3$维向量，并把所有可能的
        基本转变也表示成这样的向量：
        \begin{center}
          \begin{tabular}{rr}
            R: &$13$  \\
            G: &$15$  \\
            B: &$17$
          \end{tabular}
          \qquad
          操作：
          $\colvec[r]{-1 \\ -1 \\ 2}$、
          $\colvec[r]{-1 \\ 2 \\ -1}$、
          以及
          $\colvec[r]{2 \\ -1 \\ -1}$
        \end{center}
        现在，只需检查下面某个线性方程组的解
        是否
        存在：
        \begin{equation*}
          \colvec[r]{13 \\ 15 \\ 17}
          +x\colvec[r]{-1 \\ -1  \\ 2}
          +y\colvec[r]{-1 \\ 2 \\ -1}
          +z\colvec[r]{2 \\ -1 \\ -1}
          =\colvec[r]{0 \\ 0 \\ 45}
          \qquad
          \text{（或 $\colvec[r]{0 \\ 45 \\ 0}$ 或 $\colvec[r]{45 \\ 0 \\ 0}$）}
        \end{equation*}
        求解
        \begin{equation*}
          \begin{amat}[r]{3}
           -1  &-1 &2  &-13  \\
           -1  &2  &-1 &-15  \\
            2  &-1 &-1 &28
          \end{amat}
          \grstep[2\rho_1+\rho_3]{-\rho_1+\rho_2}
          \repeatedgrstep{\rho_2+\rho_3}
          \begin{amat}[r]{3}
           -1  &-1 &2  &-13  \\
            0  &3  &-3 &-2  \\
            0  &0  &0  &0
          \end{amat}
        \end{equation*}
        给出$y+2/3=z$，所以如果变换次数$z$是一个整数
        那么$y$就不是。
        另外两个方程组给出类似的结论，因此无解。
     
\end{ans}
\begin{ans}{一.I.2.33}
       \answerasgiven
       \begin{exparts}
        \partsitem 对方程组作形式求解得到
          \begin{equation*}
            x=\frac{a^3-1}{a^2-1}
            \qquad
            y=\frac{-a^2+a}{a^2-1}.
          \end{equation*}
          如果$a+1\neq 0$且$a-1\neq 0$，那么方程组有唯一的
          解
          \begin{equation*}
            x=\frac{a^2+a+1}{a+1}
            \qquad
            y=\frac{-a}{a+1}.
          \end{equation*}
          如果$a=-1$，或者$a=+1$，那么这些公式没有意义；在
          第一种情形，我们得到方程组
          \begin{equation*}
            \left\{
            \begin{linsys}{2}
              -x &+  &y  &=  &1 \\
               x &-  &y  &=  &1
            \end{linsys}\right.
          \end{equation*}
          这是一个矛盾方程组。
          在第二种情形，我们有
          \begin{equation*}
            \left\{
            \begin{linsys}{2}
               x &+  &y  &=  &1 \\
               x &+  &y  &=  &1
            \end{linsys}\right.
          \end{equation*}
          它有无穷多个解（例如，$x$任取，
          $y=1-x$）。
        \partsitem 求解方程组得到
          \begin{equation*}
            x=\frac{a^4-1}{a^2-1}
            \qquad
            y=\frac{-a^3+a}{a^2-1}.
          \end{equation*}
          这里，若$a^2-1\neq 0$，则方程组有唯一解
          $x=a^2+1$，$y=-a$。
          对于$a=-1$和$a=1$，我们得到方程组
          \begin{equation*}
            \left\{
            \begin{linsys}{2}
              -x &+  &y  &=  &-1 \\
               x &-  &y  &=  &1
            \end{linsys}\right.
            \qquad
            \left\{
            \begin{linsys}{2}
               x &+  &y  &=  &1 \\
               x &+  &y  &=  &1
            \end{linsys}\right.
         \end{equation*}
         两者都有无穷多个解。
      \end{exparts}
    
\end{ans}
\begin{ans}{一.I.2.34}
      \answerasgiven
      设\( u \)、\( v \)、\( x \)、\( y \)、\( z \)分别为球中所含
      Al、Cu、Pb、Ag和Au的体积（单位
      \( {\rm cm}^3 \)），我们假设球不是空心的。
      由于在水中（比重\( 1.00 \)）的失重
      为
      \( 1000 \)克，球的体积为\( 1000\mbox{ cm}^3 \)。
      于是这些数据（其中一些是多余的，尽管是相容的）只导出
      \( 2 \)个独立的方程，一个联系体积，另一个联系重量。
      \begin{equation*}
        \begin{linsys}{5}
           u  &+  &v    &+  &x     &+  &y     &+  &z     &=  &1000  \\
        2.7u  &+  &8.9v &+  &11.3x &+  &10.5y &+  &19.3z &=  &7558
        \end{linsys}
      \end{equation*}
      显然，球必须含有一些铝，以使其平均比重
      低于其他所有金属的比重。
      问题的这一部分没有唯一的结果，因为三种金属的量
      可以任意选取，只要所选取的量
      不会使任何金属的量为负。

      如果球只含铝和金，那么有
      \( 294.5\mbox{ cm}^3 \)的金和\( 705.5\mbox{ cm}^3 \)的铝。
      另一种可能是Cu、Au、Pb和
      Ag各\( 124.7\mbox{ cm}^3 \)，以及\( 501.2\mbox{ cm}^3  \)的Al。
    
\end{ans}
\protect \subsection {一.I.3: 通解$=$特解$+$齐次解}
\begin{ans}{一.I.3.14}
     下面的化简求解了该方程组。
     \begin{align*}
        \begin{amat}{3}
          1 &1   &-2 &0 \\
          1 &-1  &0  &3  \\
          3 &-1  &-2 &-6 \\
          0 &2   &-2 &3
        \end{amat}
        &\grstep[-3\rho_1+\rho_3]{-\rho_1+\rho_2}
        \begin{amat}{3}
          1 &1   &-2 &0 \\
          0 &-2  &2  &-3  \\
          0 &-4  &4  &-6 \\
          0 &2   &-2 &3
        \end{amat}                                 \\
        &\grstep[\rho_2+\rho_4]{-2\rho_2+\rho_3}
        \begin{amat}{3}
          1 &1   &-2 &0 \\
          0 &-2  &2  &-3  \\
          0 &0   &0  &0 \\
          0 &0   &0  &0
        \end{amat}
      \end{align*}
      解集如下。
      \begin{equation*}
        \set{\colvec{-3/2 \\ 3/2 \\ 0}
             +\colvec{1 \\ 1 \\ 1}z
             \suchthat z\in\Re}
      \end{equation*}
      类似地，我们可以化简相应的齐次方程组
      \begin{align*}
        \begin{amat}{3}
          1 &1   &-2 &0 \\
          1 &-1  &0  &0  \\
          3 &-1  &-2 &0 \\
          0 &2   &-2 &0
        \end{amat}
        &\grstep[-3\rho_1+\rho_3]{-\rho_1+\rho_2}
        \begin{amat}{3}
          1 &1   &-2 &0 \\
          0 &-2  &2  &0  \\
          0 &-4  &4  &0 \\
          0 &2   &-2 &0
        \end{amat}                                         \\
        &\grstep[\rho_2+\rho_4]{-2\rho_2+\rho_3}
        \begin{amat}{3}
          1 &1   &-2 &0 \\
          0 &-2  &2  &0  \\
          0 &0   &0  &0 \\
          0 &0   &0  &0
        \end{amat}
      \end{align*}
      得到它的解集。
      \begin{equation*}
        \set{
             \colvec{1 \\ 1 \\ 1}z
             \suchthat z\in\Re}
      \end{equation*}
    
\end{ans}
\begin{ans}{一.I.3.15}
      这几题的具体计算参见前一小节的答案。

      \begin{exparts}
        \partsitem
          解集如下。
          \begin{equation*}
            S=\set{\colvec[r]{6 \\ 0}+\colvec[r]{-2 \\ 1}y
              \suchthat y\in\Re}
          \end{equation*}
          下面是特解以及相应的齐次方程组的解集。
          \begin{equation*}
            \colvec[r]{6 \\ 0}
              \quad\text{与}\quad
            \set{\colvec[r]{-2 \\ 1}y
              \suchthat y\in\Re}
          \end{equation*}
          \textit{注。}
          这里可能有两点混淆。
          首先，上面给出的集合 $S$ 等于这个集合
          \begin{equation*}
            T=\set{\colvec[r]{4 \\ 1}+\colvec[r]{-2 \\ 1}y
              \suchthat y\in\Re}
          \end{equation*}
          因为这两个集合含有相同的元素。
          对于“什么是特解？”这一问题，下面这些都是正确的答案：
          \begin{equation*}
            \colvec{6 \\ 0},\quad
            \colvec{4 \\ 1},\quad
            \colvec{2 \\ 2},\quad
            \colvec{1 \\ 2.5}
          \end{equation*}
          第二点混淆是，我们在集合中使用的字母无关紧要。
          这个集合也等于 $S$。
          \begin{equation*}
            U=\set{\colvec[r]{6 \\ 0}+\colvec[r]{-2 \\ 1}u
                   \suchthat u\in\Re}
          \end{equation*}
        \partsitem
          解集如下。
          \begin{equation*}
            \set{\colvec[r]{0 \\ 1} }
          \end{equation*}
          下面是一个特解，
          以及相应的齐次方程组的解集。
          \begin{equation*}
            \colvec[r]{0 \\ 1}
              \qquad
            \set{\colvec[r]{0 \\ 0} }
          \end{equation*}
        \partsitem
          解集是无穷的。
          \begin{equation*}
            \set{\colvec[r]{4 \\ -1 \\ 0}+\colvec[r]{-1 \\ 1 \\ 1}x_3
              \suchthat x_3\in\Re}
          \end{equation*}
          下面是一个特解以及相应的齐次方程组的解集。
          \begin{equation*}
            \colvec[r]{4 \\ -1 \\ 0}
              \qquad
            \set{\colvec[r]{-1 \\ 1 \\ 1}x_3
              \suchthat x_3\in\Re}
          \end{equation*}
        \partsitem
          解集是单元素集。
          \begin{equation*}
            \set{\colvec[r]{1 \\ 1 \\ 1}}
          \end{equation*}
          下面是一个特解以及相应的齐次方程组的解集。
          \begin{equation*}
            \colvec[r]{1 \\ 1 \\ 1}
              \qquad
            \set{\colvec[r]{0 \\ 0 \\ 0}}
          \end{equation*}
        \partsitem
          解集是无穷的。
          \begin{equation*}
            \set{\colvec[r]{5/3 \\ 2/3 \\ 0 \\ 0}
                 +\colvec[r]{-1/3 \\ 2/3 \\ 1 \\ 0}z
                 +\colvec[r]{-2/3 \\ 1/3 \\ 0 \\ 1}w
                 \suchthat z,w\in\Re}
          \end{equation*}
          下面是一个特解以及相应的齐次方程组的解集。
          \begin{equation*}
            \colvec[r]{5/3 \\ 2/3 \\ 0 \\ 0}
              \qquad
            \set{\colvec[r]{-1/3 \\ 2/3 \\ 1 \\ 0}z
                 +\colvec[r]{-2/3 \\ 1/3 \\ 0 \\ 1}w
                 \suchthat z,w\in\Re}
          \end{equation*}
        \partsitem 该方程组的解集为空。
          因此，不存在特解。
          相应的齐次方程组的解集如下。
          \begin{equation*}
            \set{\colvec[r]{-1 \\ 2 \\ 1 \\ 0}z
                 +\colvec[r]{-1 \\ 3 \\ 0 \\ 1}w
                 \suchthat z,w\in\Re}
          \end{equation*}
      \end{exparts}
    
\end{ans}
\begin{ans}{一.I.3.16}
    前一小节的答案给出了行变换。
    这里的每个答案只列出解集、特解和齐次解。
    \begin{exparts}
      \partsitem
        解集如下。
        \begin{equation*}
          \set{\colvec[r]{2/3 \\ -1/3 \\ 0}
               +\colvec[r]{1/6 \\ 2/3 \\ 1}z
              \suchthat z\in\Re}
        \end{equation*}
        下面是一个特解以及相应的齐次方程组的解集。
        \begin{equation*}
          \colvec[r]{2/3 \\ -1/3 \\ 0}
            \qquad
        \set{\colvec[r]{1/6 \\ 2/3 \\ 1}z
            \suchthat z\in\Re}
        \end{equation*}
      \partsitem
        解集是无穷的。
        \begin{equation*}
          \set{\colvec[r]{1 \\ 3 \\ 0 \\ 0}
               +\colvec[r]{1 \\-2 \\ 1 \\ 0}z
              \suchthat z\in\Re}
        \end{equation*}
        下面是
        一个特解以及相应的齐次方程组的解集。
        \begin{equation*}
          \colvec[r]{1 \\ 3 \\ 0 \\ 0}
            \qquad
          \set{\colvec[r]{1 \\ -2 \\ 1 \\ 0}z
              \suchthat z\in\Re}
        \end{equation*}
      \partsitem
        解集如下。
        \begin{equation*}
          \set{\colvec[r]{0 \\ 0 \\ 0 \\ 0}
               +\colvec[r]{-1 \\ 0 \\ 1 \\ 0}z
               +\colvec[r]{-1 \\ -1 \\ 0 \\ 1}w
              \suchthat z,w\in\Re}
        \end{equation*}
        下面是
        一个特解以及相应的齐次方程组的解集。
        \begin{equation*}
          \colvec[r]{0 \\ 0 \\ 0 \\ 0}
            \qquad
          \set{\colvec[r]{-1 \\ 0 \\ 1 \\ 0}z
               +\colvec[r]{-1 \\ -1 \\ 0 \\ 1}w
              \suchthat z,w\in\Re}
        \end{equation*}
      \partsitem
        解集如下。
        \begin{equation*}
          \set{\colvec[r]{1 \\ 0 \\ 0 \\ 0 \\ 0}
               +\colvec[r]{-5/7 \\ -8/7 \\ 1 \\ 0 \\ 0}c
               +\colvec[r]{-3/7 \\ -2/7 \\ 0 \\ 1 \\ 0}d
               +\colvec[r]{-1/7 \\ 4/7 \\ 0 \\ 0 \\ 1}e
              \suchthat c,d,e\in\Re}
        \end{equation*}
        下面是
        一个特解以及相应的齐次方程组的解集。
        \begin{equation*}
          \colvec[r]{1 \\ 0 \\ 0 \\ 0 \\ 0}
            \qquad
          \set{\colvec[r]{-5/7 \\ -8/7 \\ 1 \\ 0 \\ 0}c
               +\colvec[r]{-3/7 \\ -2/7 \\ 0 \\ 1 \\ 0}d
               +\colvec[r]{-1/7 \\ 4/7 \\ 0 \\ 0 \\ 1}e
              \suchthat c,d,e\in\Re}
        \end{equation*}
    \end{exparts}
   
\end{ans}
\begin{ans}{一.I.3.17}
      直接把它们代入，看它们是否满足全部三个方程。
      \begin{exparts}
        \partsitem 否。
        \partsitem 是。
        \partsitem 是。
      \end{exparts}
    
\end{ans}
\begin{ans}{一.I.3.18}
      对相应的齐次方程组施行高斯消元法
      \begin{align*}
        \begin{amat}[r]{4}
           1  &-1  &0  &1  &0  \\
           2  &3   &-1 &0  &0  \\
           0  &1   &1  &1  &0
        \end{amat}
        &\grstep{-2\rho_1+\rho_2}
        \begin{amat}[r]{4}
           1  &-1  &0  &1  &0  \\
           0  &5   &-1 &-2 &0  \\
           0  &1   &1  &1  &0
        \end{amat}                                \\
        &\grstep{-(1/5)\rho_2+\rho_3}
        \begin{amat}[r]{4}
           1  &-1  &0  &1  &0  \\
           0  &5   &-1 &-2 &0  \\
           0  &0   &6/5&7/5&0
        \end{amat}
      \end{align*}
      得到这个解是齐次问题的解。
      \begin{equation*}
        \set{\colvec[r]{-5/6 \\ 1/6 \\ -7/6 \\ 1}w\suchthat w\in\Re}
      \end{equation*}
      \begin{exparts}
        \partsitem 该向量确实是一个特解，因此所求的
          通解如下。
          \begin{equation*}
            \set{\colvec[r]{0 \\ 0 \\ 0 \\ 4}+
                 \colvec[r]{-5/6 \\ 1/6 \\ -7/6 \\ 1}w\suchthat w\in\Re}
          \end{equation*}
        \partsitem 该向量是一个特解，因此所求的
          通解如下。
          \begin{equation*}
            \set{\colvec[r]{-5 \\ 1 \\ -7 \\ 10}+
                 \colvec[r]{-5/6 \\ 1/6 \\ -7/6 \\ 1}w\suchthat w\in\Re}
          \end{equation*}
        \partsitem 该向量不是该方程组的解，因为
          它不满足第三个方程。
          这样的通解不存在。
      \end{exparts}
    
\end{ans}
\begin{ans}{一.I.3.19}
       第一个非奇异而第二个奇异。
       直接用高斯消元法，看阶梯形结果的
       对角线上每个位置的数是否都非 $0$。
     
\end{ans}
\begin{ans}{一.I.3.20}
      \begin{exparts}
      \partsitem 非奇异：
        \begin{equation*}
          \grstep{-\rho_1+\rho_2}
          \begin{mat}[r]
            1  &2  \\
            0  &1
          \end{mat}
        \end{equation*}
        最终每行都含有一个首主元。
      \partsitem 奇异：
        \begin{equation*}
          \grstep{3\rho_1+\rho_2}
          \begin{mat}[r]
            1  &2  \\
            0  &0
          \end{mat}
        \end{equation*}
        最终第 \( 2 \) 行没有首主元。
      \partsitem 都不是。
        矩阵必须是方阵，这两个词才适用。
      \partsitem 奇异。
      \partsitem 非奇异。
     \end{exparts}
    
\end{ans}
\begin{ans}{一.I.3.21}
        在每种情形下，我们都要判断该向量是否为集合中那些向量的
        线性组合。
        \begin{exparts}
          \partsitem 是。
            求解
            \begin{equation*}
              c_1\colvec[r]{1 \\ 4}+c_2\colvec[r]{1 \\ 5}=\colvec[r]{2 \\ 3}
            \end{equation*}
            结合
            \begin{equation*}
              \begin{amat}[r]{2}
                1  &1  &2  \\
                4  &5  &3
              \end{amat}
              \grstep{-4\rho_1+\rho_2}
              \begin{amat}[r]{2}
                1  &1  &2  \\
                0  &1  &-5
              \end{amat}
            \end{equation*}
            可以得出存在给出该组合的 $c_1$ 和 $c_2$。
          \partsitem 否。
            化简
            \begin{equation*}
              \begin{amat}[r]{2}
                2  &1  &-1 \\
                1  &0  &0  \\
                0  &1  &1
              \end{amat}
              \grstep{-(1/2)\rho_1+\rho_2}
              \begin{amat}[r]{2}
                2  &1     &-1 \\
                0  &-1/2  &1/2  \\
                0  &1     &1
              \end{amat}
              \grstep{2\rho_2+\rho_3}
              \begin{amat}[r]{2}
                2  &1     &-1 \\
                0  &-1/2  &1/2  \\
                0  &0     &2
              \end{amat}
            \end{equation*}
            表明
            \begin{equation*}
              c_1\colvec[r]{2 \\ 1 \\ 0}+c_2\colvec[r]{1 \\ 0 \\ 1}
                =\colvec[r]{-1 \\ 0 \\ 1}
            \end{equation*}
            无解。
          \partsitem 是。
            化简
            \begin{align*}
              \begin{amat}[r]{4}
                1  &2  &3  &4  &1  \\
                0  &1  &3  &2  &3  \\
                4  &5  &0  &1  &0
              \end{amat}
              &\grstep{-4\rho_1+\rho_3}
              \begin{amat}[r]{4}
                1  &2  &3  &4  &1  \\
                0  &1  &3  &2  &3  \\
                0  &-3 &-12&-15&-4
              \end{amat}                   \\
              &\grstep{3\rho_2+\rho_3}
              \begin{amat}[r]{4}
                1  &2  &3  &4  &1  \\
                0  &1  &3  &2  &3  \\
                0  &0  &-3 &-9 &5
              \end{amat}
            \end{align*}
            表明有无穷多种方式
            \begin{equation*}
              \set{\colvec[r]{c_1 \\ c_2 \\ c_3 \\ c_4}=
                   \colvec[r]{-10 \\ 8 \\ -5/3 \\ 0}+
                   \colvec[r]{-9 \\ 7 \\ -3 \\ 1}c_4
                    \suchthat c_4\in\Re}
            \end{equation*}
            来写出这个组合。
            \begin{equation*}
              \colvec[r]{1 \\ 3 \\ 0}=
              c_1\colvec[r]{1 \\ 0 \\ 4}+
              c_2\colvec[r]{2 \\ 1 \\ 5}+
              c_3\colvec[r]{3 \\ 3 \\ 0}+
              c_4\colvec[r]{4 \\ 2 \\ 1}
            \end{equation*}
          \partsitem 否。
            看第三个分量。
        \end{exparts}
      
\end{ans}
\begin{ans}{一.I.3.22}
        由于系数矩阵是非奇异的，高斯消元法
        结束时得到的阶梯形中每个变量都在某个方程中居于首位。
        回代给出唯一解。

      （另一种看出解唯一的方法是注意：
      当系数矩阵非奇异时，由定义，
      相应的齐次方程组有唯一解。
      由于通解是一个特解与
      每个齐次解的和，所以通解
      至多只有一个元素。）
     
\end{ans}
\begin{ans}{一.I.3.23}
      此时解集是整个 \( \Re^n \)，而且我们可以
      用所要求的形式表示它。
      \begin{equation*}
        \set{c_1\colvec[r]{1 \\ 0 \\ \vdotswithin{1} \\ 0}
             +c_2\colvec[r]{0 \\ 1 \\ \vdotswithin{0} \\ 0}
             +\cdots
             +c_n\colvec[r]{0 \\ 0 \\ \vdotswithin{0} \\ 1}
             \suchthat c_1,\ldots,c_n\in\Re}
      \end{equation*}
     
\end{ans}
\begin{ans}{一.I.3.24}
      设 \( \vec{s},\vec{t}\in\Re^n \)，并把它们写成如下形式。
      \begin{equation*}
        \vec{s}=\colvec{s_1 \\ \vdotswithin{s_1} \\ s_n}
          \qquad
        \vec{t}=\colvec{t_1 \\ \vdotswithin{t_1} \\ t_n}
      \end{equation*}
      再设 \( a_{i,1}x_1+\cdots+a_{i,n}x_n=0 \) 是齐次方程组中的
      第 \( i \) 个方程。
      \begin{exparts}
        \partsitem 验证很容易。
          \begin{multline*}
            a_{i,1}(s_1+t_1)+\cdots+a_{i,n}(s_n+t_n)      \\
            =
            (a_{i,1}s_1+\cdots+a_{i,n}s_n)
            +(a_{i,1}t_1+\cdots+a_{i,n}t_n)
            =
            0+0
          \end{multline*}
        \partsitem 这与前一个类似。
          \begin{equation*}
            a_{i,1}(3s_1)+\cdots+a_{i,n}(3s_n)
            =3(a_{i,1}s_1+\cdots+a_{i,n}s_n)
            =3\cdot 0=0
          \end{equation*}
        \partsitem 这个也不难多少。
          \begin{multline*}
            a_{i,1}(ks_1+mt_1)+\cdots+a_{i,n}(ks_n+mt_n)  \\
            =
            k(a_{i,1}s_1+\cdots+a_{i,n}s_n)
            +m(a_{i,1}t_1+\cdots+a_{i,n}t_n)
            =
            k\cdot 0+m\cdot 0
          \end{multline*}
      \end{exparts}
     这个论证的错误在于，只涉及
     零向量的任何线性组合都会得到零向量。
   
\end{ans}
\begin{ans}{一.I.3.25}
      先看证明。

      若方程组的系数全为有理数，则
      带回代的高斯消元法将只用到有理数（例如，
      \( -(m/n)\rho_i+\rho_j \)）。
      于是，我们可以只用有理数
      作为每个向量的分量来表示解集。
      因此，特解全为有理数。

      至于无穷多解的问题，
      有理向量解有无穷多个，当且仅当
      相应的齐次方程组有无穷多个
      实向量解。
      这是因为把任意参数取成有理数都会产生一个
      全为有理数的解。
   
\end{ans}
\protect \section {第 II 节：线性几何}
\protect \subsection {一.II.1: 空间中的向量}
\begin{ans}{一.II.1.1}
      \begin{exparts*}
        \partsitem \( \colvec[r]{2 \\ 1}  \)
        \partsitem \( \colvec[r]{-1 \\ 2}  \)
        \partsitem \( \colvec[r]{4 \\ 0 \\ -3}  \)
        \partsitem \( \colvec[r]{0 \\ 0 \\ 0}  \)
      \end{exparts*}
     
\end{ans}
\begin{ans}{一.II.1.2}
      \begin{exparts}
        \partsitem 不相等，它们的规范位置不同。
          \begin{equation*}
            \colvec[r]{1 \\ -1}
            \qquad
            \colvec[r]{0 \\ 3}
          \end{equation*}
        \partsitem 相等，它们的规范位置相同。
          \begin{equation*}
            \colvec[r]{1 \\ -1 \\ 3}
          \end{equation*}
      \end{exparts}
     
\end{ans}
\begin{ans}{一.II.1.3}
      这条直线是这个集合。
      \begin{equation*}
        \set{\colvec[r]{-2 \\ 1 \\ 1 \\ 0}
             +\colvec[r]{7 \\ 9 \\ -2 \\ 4}t \suchthat t\in\Re }
      \end{equation*}
      注意，方程组
      \begin{equation*}
        \begin{linsys}{2}
          -2  &+  &7t  &=  &1  \\
           1  &+  &9t  &=  &0  \\
           1  &-  &2t  &=  &2  \\
           0  &+  &4t  &=  &1
        \end{linsys}
      \end{equation*}
      无解。
      因此所给的点不在这条直线上。
    
\end{ans}
\begin{ans}{一.II.1.4}
      \begin{exparts}
        \partsitem 注意到
          \begin{equation*}
            \colvec[r]{2 \\ 2 \\ 2 \\ 0}
            -\colvec[r]{1 \\ 1 \\ 5 \\ -1}
            =\colvec[r]{1 \\ 1 \\ -3 \\ 1}
            \qquad
            \colvec[r]{3 \\ 1 \\ 0 \\ 4}
            -\colvec[r]{1 \\ 1 \\ 5 \\ -1}
            =\colvec[r]{2 \\ 0 \\ -5 \\ 5}
          \end{equation*}
          所以该平面是这个集合。
          \begin{equation*}
            \set{\colvec[r]{1 \\ 1 \\ 5 \\ -1}
                 +\colvec[r]{1 \\ 1 \\ -3 \\ 1}t
                 +\colvec[r]{2 \\ 0 \\ -5 \\ 5}s
                \suchthat t,s\in\Re}
          \end{equation*}
        \partsitem 不在；方程组
          \begin{equation*}
            \begin{linsys}{3}
              1  &+  &1t  &+  &2s  &=  &0  \\
              1  &+  &1t  &   &    &=  &0  \\
              5  &-  &3t  &-  &5s  &=  &0  \\
             -1  &+  &1t  &+  &5s  &=  &0
            \end{linsys}
          \end{equation*}
          无解。
      \end{exparts}
    
\end{ans}
\begin{ans}{一.II.1.5}
      \begin{exparts}
        \item 把 $x-2y+z=4$ 看作只含一个方程的线性方程组，
          以变量 $y$ 和~$z$ 为参数进行参数化，得到
          $x=4+2y-z$。
          由此得到这个平面的向量描述。
          \begin{equation*}
            \set{\colvec{x \\ y \\ z}=\colvec{4 \\ 0 \\ 0}
                                       +\colvec{2 \\ 1 \\ 0}\cdot y
                                       +\colvec{-1 \\ 0 \\ 1}\cdot z
                             \suchthat y,z\in\Re}
          \end{equation*}
        \item 参数化给出 $x=-(1/2)-(1/2)y-2z$，
          所以这就是向量描述。
          \begin{equation*}
            \colvec{x \\ y \\ z}=\colvec{-1/2 \\ 0 \\ 0}
                                       +\colvec{-1/2 \\ 1 \\ 0}\cdot y
                                       +\colvec{-2 \\ 0 \\ 1}\cdot z
          \end{equation*}
        \item 这里
          $x=10-y-z-w$，于是我们得到含三个参数的向量描述。
          \begin{equation*}
            \colvec{x \\ y \\ z \\ w}=\colvec{10 \\ 0 \\ 0 \\ 0}
                                       +\colvec{-1 \\ 1 \\ 0 \\ 0}\cdot y
                                       +\colvec{-1 \\ 0 \\ 1 \\ 0}\cdot z
                                       +\colvec{-1 \\ 0 \\ 0 \\ 1}\cdot w
          \end{equation*}
      \end{exparts}
    
\end{ans}
\begin{ans}{一.II.1.6}
      向量
      \begin{equation*}
        \colvec[r]{2 \\ 0 \\ 3}
      \end{equation*}
      不在这条直线上。
      由于
      \begin{equation*}
        \colvec[r]{2 \\ 0 \\ 3}
        -\colvec[r]{-1 \\ 0 \\ -4}
        =\colvec[r]{3 \\ 0 \\ 7}
      \end{equation*}
      我们可以这样描述该平面。
      \begin{equation*}
        \set{\colvec[r]{-1 \\ 0 \\ -4}
             +m\colvec[r]{1 \\ 1 \\ 2}
             +n\colvec[r]{3 \\ 0 \\ 7}
            \suchthat m,n\in\Re}
      \end{equation*}
   
\end{ans}
\begin{ans}{一.II.1.7}
      重合的点就是下面这个方程组的解。
      \begin{equation*}
        \begin{linsys}{2}
         t  &  &   &= &1+2m\hfill      \\
         t  &+ &s  &= &1+3k\hfill      \\
         t  &+ &3s &= &4m\hfill
        \end{linsys}
      \end{equation*}
      用高斯消元法
      \begin{align*}
        \begin{amat}{4}
          1  &0  &0  &-2  &1  \\
          1  &1  &-3 &0   &1  \\
          1  &3  &0  &-4  &0
        \end{amat}
        &\grstep[-\rho_1+\rho_3]{-\rho_1+\rho_2}
        \begin{amat}{4}
          1  &0  &0  &-2  &1  \\
          0  &1  &-3 &2   &0  \\
          0  &3  &0  &-2  &-1
        \end{amat}                       \\
        &\grstep{-3\rho_2+\rho_3}
        \begin{amat}{4}
          1  &0  &0  &-2  &1  \\
          0  &1  &-3 &2   &0  \\
          0  &0  &9  &-8  &-1
        \end{amat}
      \end{align*}
      得到 \( k=-(1/9)+(8/9)m \)，所以 \( s=-(1/3)+(2/3)m \) 且 \( t=1+2m \)。
      交是下面这个集合。
      \begin{equation*}
        \set{\colvec[r]{1 \\ 1 \\ 0}+
             \colvec[r]{0 \\ 3 \\ 0}(-\frac{1}{9}+\frac{8}{9}m)+
             \colvec[r]{2 \\ 0 \\ 4}m
             \suchthat m\in\Re}
        =\set{\colvec[r]{1 \\ 2/3 \\ 0}
             +\colvec[r]{2 \\ 8/3 \\ 4}m
             \suchthat m\in\Re}
      \end{equation*}
    
\end{ans}
\begin{ans}{一.II.1.8}
       \begin{exparts}
        \partsitem 方程组
          \begin{equation*}
            \begin{linsys}{1}
              1     &= &1\hfill         \\
              1+t   &= &3+s\hfill         \\
              2+t   &= &-2+2s\hfill
            \end{linsys}
          \end{equation*}
          给出 \( s=6 \) 与 \( t=8 \)，所以这就是解集。
          \begin{equation*}
            \set{\colvec[r]{1 \\ 9 \\ 10}     }
          \end{equation*}
        \partsitem 这个方程组
          \begin{equation*}
            \begin{linsys}{1}
            2+t   &=  &0 \hfill        \\
            t     &=  &s+4w\hfill        \\
            1-t   &=  &2s+w\hfill
            \end{linsys}
          \end{equation*}
          给出 \( t=-2 \)、\( w=-1 \) 与 \( s=2 \)，所以它们的交
          是这个点。
          \begin{equation*}
            \colvec[r]{0 \\ -2 \\ 3}
          \end{equation*}
      \end{exparts}
    
\end{ans}
\begin{ans}{一.II.1.9}
        我们将在后面把它定义成恰有一个元素的集合\Dash 一个
        “原点”。
      
\end{ans}
\begin{ans}{一.II.1.10}
      \answerasgiven %
      向量三角形如下图所示，因此 \( \vec{w}=3\sqrt{2} \)，
      风从西北方向吹来。
      \begin{center}
        \setlength{\unitlength}{4pt}      % wind triangles.
        \begin{picture}(20,12)(0,0)
           \put(0,10){\vector(1,0){20} }
           \put(0,10){\vector(1,-1){10} }
              \put(3.8,5){\makebox(0,0)[r]{\small \( \vec{w} \)} }
           \put(10,0){\vector(1,1){10} }
           \put(10,0){\line(0,1){10} }
        \end{picture}
      \end{center}
     
\end{ans}
\begin{ans}{一.II.1.11}
      欧几里得无疑是在设想 \( \Re^3 \) 之内的平面。
      然而，请注意 \( \Re^1 \) 与 \( \Re^2 \) 也都满足
      这个定义。
    
\end{ans}
\protect \subsection {一.II.2: 长度与夹角的度量}
\begin{ans}{一.II.2.11}
      \begin{exparts*}
        \partsitem \( \sqrt{3^2+1^2}=\sqrt{10}  \)
        \partsitem \( \sqrt{5}  \)
        \partsitem \( \sqrt{18}  \)
        \partsitem \( 0  \)
        \partsitem \( \sqrt{3}  \)
      \end{exparts*}
    
\end{ans}
\begin{ans}{一.II.2.12}
       \begin{exparts*}
         \partsitem \( \arccos(9/\sqrt{85})\approx 0.22\mbox{\ 弧度} \)
         \partsitem \( \arccos(8/\sqrt{85})\approx 0.52\mbox{\ 弧度} \)
         \partsitem 无定义。
      \end{exparts*}
     
\end{ans}
\begin{ans}{一.II.2.13}
      我们把每个位移表示为一个向量，舍入到一位
      小数，因为这就是题目陈述所具有的精度，
      然后相加得到总位移
      （忽略地球曲率）。
      \begin{equation*}
        \colvec[r]{0.0 \\ 1.2}
        +\colvec[r]{3.8 \\ -4.8}
        +\colvec[r]{4.0 \\ 0.1}
        +\colvec[r]{3.3 \\ 5.6}
        =\colvec[r]{11.1 \\ 2.1}
      \end{equation*}
      距离为 \( \sqrt{11.1^2+2.1^2}\approx 11.3 \)。
    
\end{ans}
\begin{ans}{一.II.2.14}
       解 \( (k)(4)+(1)(3)=0 \) 得 \( k=-3/4 \)。
    
\end{ans}
\begin{ans}{一.II.2.15}
      我们可以用参数把集合
      \begin{equation*}
         \set{\colvec{x \\ y \\ z}\suchthat 1x+3y-1z=0}
      \end{equation*}
      描述成这种形式：
      \begin{equation*}
         \set{\colvec[r]{-3 \\ 1 \\ 0}y+\colvec[r]{1 \\ 0 \\ 1}z
               \suchthat y,z\in\Re}
      \end{equation*}
    
\end{ans}
\begin{ans}{一.II.2.16}
      \begin{exparts}
        \partsitem 我们可以取 \( x \) 轴。
          \begin{equation*}
            \arccos (\frac{(1)(1)+(0)(1)}{\sqrt{1}\sqrt{2}})
            \approx 0.79~\text{弧度}
          \end{equation*}
        \partsitem 同样，取 \( x \) 轴。
          \begin{equation*}
            \arccos (\frac{(1)(1)+(0)(1)+(0)(1)}{\sqrt{1}\sqrt{3}})
            \approx 0.96~\text{弧度}
          \end{equation*}
        \partsitem \( x \) 轴前面很管用，这里它同样管用。
          \begin{equation*}
            \arccos (\frac{(1)(1)+\cdots+(0)(1)}{\sqrt{1}\sqrt{n}})
            =\arccos (\frac{1}{\sqrt{n}})
          \end{equation*}
        \partsitem 利用上一小题的公式，
            $\lim_{n\to\infty} \arccos(1/\sqrt{n})
              =\pi/2\text{\ 弧度}$。
      \end{exparts}
    
\end{ans}
\begin{ans}{一.II.2.17}
      显然，\( u_1u_1+\cdots+u_nu_n \) 为零当且仅当
      每个 \( u_i \) 都为零。
      所以只有 \( \zero\in\Re^n \) 与它自身垂直。
    
\end{ans}
\begin{ans}{一.II.2.18}
      在下面每一小题中，假设向量
      \( \vec{u},\vec{v},\vec{w}\in\Re^n \)
      的分量分别为
      \( u_1,\ldots,u_n,v_1,\ldots,w_n \)。
      \begin{exparts}
        \partsitem 点积对加法右分配。
           \begin{align*}
              (\vec{u}+\vec{v})\dotprod\vec{w}
              &=[\colvec{u_1 \\ \vdotswithin{u_1} \\ u_n}
                +\colvec{v_1 \\ \vdotswithin{v_1} \\ v_n}]\dotprod
                \colvec{w_1 \\ \vdotswithin{w_1} \\ w_n}               \\
              &=
              \colvec{u_1+v_1 \\ \vdotswithin{u_1+v_1} \\ u_n+v_n}\dotprod
                \colvec{w_1 \\ \vdotswithin{w_1} \\ w_n}               \\
              &=
              (u_1+v_1)w_1+\cdots+(u_n+v_n)w_n              \\
              &=
              (u_1w_1+\cdots+u_nw_n)+(v_1w_1+\cdots+v_nw_n)  \\
              &=
              \vec{u}\dotprod\vec{w}+\vec{v}\dotprod\vec{w}
           \end{align*}
        \partsitem 点积也左分配：
          $\vec{w}\dotprod(\vec{u}+\vec{v})=
              \vec{w}\dotprod\vec{u}+\vec{w}\dotprod\vec{v}$。
          证明与上一小题完全类似。
        \partsitem 点积可交换。
          \begin{equation*}
            \colvec{u_1 \\ \vdotswithin{u_1} \\ u_n}\dotprod
              \colvec{v_1 \\ \vdotswithin{v_1} \\ v_n}
            =u_1v_1+\cdots+u_nv_n
            =v_1u_1+\cdots+v_nu_n
            =\colvec{v_1 \\ \vdotswithin{v_1} \\ v_n}\dotprod
              \colvec{u_1 \\ \vdotswithin{u_1} \\ u_n}
          \end{equation*}
        \partsitem 由于 \( \vec{u}\dotprod\vec{v} \)
          是标量而不是向量，
          表达式 \( (\vec{u}\dotprod\vec{v})\dotprod\vec{w} \) 没有
          意义；标量与向量的点积没有定义。
        \partsitem 这是一个含糊的问题，所以它有多种答案。
          其中有
          (1)~\( k(\vec{u}\dotprod\vec{v})=(k\vec{u})\dotprod\vec{v} \)
          以及 \( k(\vec{u}\dotprod\vec{v})=\vec{u}\dotprod(k\vec{v}) \)，
          (2)~\( k(\vec{u}\dotprod\vec{v})\neq (k\vec{u})\dotprod(k\vec{v}) \)
          （一般而言；容易举出例子），以及
          (3)~\( \absval{k\vec{v}\,}=\lvert k\rvert\absval{\vec{v}\,} \)
          （长度与点积之间的联系在于：
          长度的平方就是向量
          与自身的点积）。
      \end{exparts}
    
\end{ans}
\begin{ans}{一.II.2.19}
       \begin{exparts}
         \partsitem 容易验证：对 \( k\in\Re \) 与
           \( \vec{x},\vec{y}\in\Re^n \)，
           \( (k\vec{x})\dotprod\vec{y}=k(\vec{x}\dotprod\vec{y})
           =\vec{x}\dotprod(k\vec{y}) \)。
           现在，对 \( k\in\Re \) 与 \( \vec{v},\vec{w}\in\Re^n \)，
           若 \( \vec{u}=k\vec{v} \)，则
           \( \vec{u}\dotprod\vec{v}=(k\vec{v})\dotprod\vec{v}
             =k(\vec{v}\dotprod\vec{v}) \)，
           它是 \( k \) 乘一个非负
           实数。

           \( \vec{v}=k\vec{u} \) 的情形类似（实际上，取本段中的
           \( k \) 为上文 \( k \) 的倒数即知，我们只需担心
           \( k=0 \) 的情形）。
         \partsitem 我们先考虑 $\vec{u}\dotprod\vec{v}\geq 0$ 的情形。
           由三角不等式可知，
           \( \vec{u}\dotprod\vec{v}=\absval{\vec{u}\,}\,\absval{\vec{v}\,} \)
           当且仅当其中一个向量是另一个向量的非负标量倍。
           但这就足够了，因为
           本题第一小题表明，
           在两个向量的点积为正的背景下，
           下面两个命题
           “其中一个向量是
           另一个向量的标量倍”与“其中一个向量是另一个向量的非负
           标量倍”
           等价。

           最后考虑 $\vec{u}\dotprod\vec{v}< 0$ 的情形。
           由于
           $0<\absval{\vec{u}\dotprod\vec{v}}=-(\vec{u}\dotprod\vec{v})
            =(-\vec{u})\dotprod\vec{v}$ 且
           $\absval{\vec{u}\,}\,\absval{\vec{v}\,}
              =\absval{-\vec{u}\,}\,\absval{\vec{v}\,}$，
           我们有
           $0<(-\vec{u})\dotprod\vec{v}=\absval{-\vec{u}\,}\,\absval{\vec{v}\,}$。
           于是上一段适用，给出：$-\vec{u}$ 与 $\vec{v}$ 这两个向量
           之一是另一个的标量倍。
           而这等价于断言 $\vec{u}$ 与 $\vec{v}$ 这两个向量
           之一是另一个的标量倍，
           即为所证。
      \end{exparts}
    
\end{ans}
\begin{ans}{一.II.2.20}
      不必。
      下面这些向量给出一个例子。
      \begin{equation*}
        \vec{u}=\colvec[r]{1 \\ 0}
        \quad
        \vec{v}=\colvec[r]{1 \\ 0}
        \quad
        \vec{w}=\colvec[r]{1 \\ 1}
      \end{equation*}
    
\end{ans}
\begin{ans}{一.II.2.21}
      我们证明：向量长度为零当且仅当它的
      所有分量都为零。

      设 \( \vec{u}\in\Re^n \) 的分量为 \( u_1,\ldots,u_n \)。
      回忆一下：任何实数的平方都大于或等于
      零，且仅当该实数为零时才取等号。
      于是
      \( \absval{\vec{u}\,}^2={u_1}^2+\cdots+{u_n}^2 \) 是若干个
      大于或等于零的数之和，因而它本身也大于或等于
      零，且当且仅当每个 \( u_i \) 都为零时取等号。
      因此 \( \absval{\vec{u}\,}=0 \) 当且仅当
      \( \vec{u} \) 的所有分量都为零。
    
\end{ans}
\begin{ans}{一.II.2.22}
      容易验证
      \begin{equation*}
        \bigl( \frac{x_1+x_2}{2},\frac{y_1+y_2}{2}  \bigr)
      \end{equation*}
      在连接这两点的直线上，并且到两点的距离相等。
      推广是显然的。
    
\end{ans}
\begin{ans}{一.II.2.23}
      设 \( \vec{v}\in\Re^n \) 的分量为 \( v_1,\ldots,v_n \)。
      若 \( \vec{v}\neq \zero \)，则有下面的计算。
      \begin{multline*}
        \sqrt{\left(\frac{v_1}{\sqrt{{v_1}^2+\cdots+{v_n}^2}}\right)^2+
              \dots+\left(\frac{v_n}{\sqrt{{v_1}^2+\cdots+{v_n}^2}}\right)^2}
        \\
        \begin{aligned}
          &=\sqrt{\left(\frac{{v_1}^2}{{v_1}^2+\cdots+{v_n}^2}\right)+
              \dots+\left(\frac{{v_n}^2}{{v_1}^2+\cdots+{v_n}^2}\right)}   \\
          &=1
        \end{aligned}
      \end{multline*}
      若 \( \vec{v}=\zero \)，则 \( \vec{v}/\absval{\vec{v}\,} \) 没有
      定义。
    
\end{ans}
\begin{ans}{一.II.2.24}
      对于第一个问题，设 \( \vec{v}\in\Re^n \) 且
      \( r\geq 0 \)，开方并提取公因子。
      \begin{equation*}
        \absval{r\vec{v}\,}
        =\sqrt{(rv_1)^2+\cdots+(rv_n)^2}
        =\sqrt{r^2({v_1}^2+\cdots+{v_n}^2}
        =r\absval{\vec{v}\,}
      \end{equation*}
      对于第二个问题，结果的长度是 \( r \) 倍，但它
      指向相反的方向，这体现在
      \( r\vec{v}+(-r)\vec{v}=\zero \)。
    
\end{ans}
\begin{ans}{一.II.2.25}
      设 \( \vec{u},\vec{v}\in\Re^n \) 的长度都是 \( 1 \)。
      应用 Cauchy–Schwarz：
      $\absval{\vec{u}\dotprod\vec{v}}
             \leq\absval{\vec{u}\,}\,\absval{\vec{v}\,}=1$。

      为看到“小于”能发生，在 \( \Re^2 \) 中取
      \begin{equation*}
        \vec{u}=\colvec[r]{1 \\ 0}
        \qquad
        \vec{v}=\colvec[r]{0 \\ 1}
      \end{equation*}
      并注意 \( \vec{u}\dotprod\vec{v}=0 \)。
      对于“等于”，注意 \( \vec{u}\dotprod\vec{u}=1 \)。
    
\end{ans}
\begin{ans}{一.II.2.26}
      \begin{exparts}
        \partsitem 图中所示的向量
          \begin{center}
            \includegraphics{gr/mp/ch1.32}
          \end{center}
          并不是将
          \begin{equation*}
            \colvec[r]{2 \\ 0 \\ 0}
              +\colvec[r]{-0.5 \\ 1 \\ 0}\cdot 1
            =\colvec[r]{1.5 \\ 1 \\ 0}
          \end{equation*}
          加倍的结果，而是将参数加倍的结果。
          \begin{equation*}
            \colvec[r]{2 \\ 0 \\ 0}
              +\colvec[r]{-0.5 \\ 1 \\ 0}\cdot 2
            =\colvec[r]{1 \\ 2 \\ 0}
          \end{equation*}
          下面比较长度。
          \begin{equation*}
            \sqrt{1^2+2^2+0^2}
            =\sqrt{(2\cdot 0.5)^2 + (2\cdot 1)^2+ (2\cdot 0)^2}
            =\sqrt{2^2\cdot(0.5^2+1^2+0^2)}
            =2\cdot\sqrt{0.5^2+1^2+0^2}
          \end{equation*}
        \partsitem 向量
          \begin{center}
            \includegraphics{gr/mp/ch1.19}
          \end{center}
          并不是下式相加的结果：
          \begin{equation*}
            (\colvec[r]{2 \\ 0 \\ 0}
              +\colvec[r]{-0.5 \\ 1 \\ 0}\cdot 1)
            +
            (\colvec[r]{2 \\ 0 \\ 0}
              +\colvec[r]{-0.5 \\ 0 \\ 1}\cdot 1)
          \end{equation*}
          而是
          \begin{equation*}
            \colvec[r]{2 \\ 0 \\ 0}
              +\colvec[r]{-0.5 \\ 1 \\ 0}\cdot 1
              +\colvec[r]{-0.5 \\ 0 \\ 1}\cdot 1
            =\colvec[r]{1 \\ 1 \\ 1}
          \end{equation*}
          后者将参数相加。
      \end{exparts}
    
\end{ans}
\begin{ans}{一.II.2.27}
      “若”这一半是直接的。
      若 \( b_1-a_1=d_1-c_1 \) 且 \( b_2-a_2=d_2-c_2 \)，则
      \begin{equation*}
        \sqrt{(b_1-a_1)^2+(b_2-a_2)^2}
        =\sqrt{(d_1-c_1)^2+(d_2-c_2)^2}
      \end{equation*}
      所以二者长度相同；斜率同样容易：
      \begin{equation*}
        \frac{b_2-a_2}{b_1-a_1}
        =\frac{d_2-c_2}{d_1-a_1}
      \end{equation*}
      （若分母为 \( 0 \)，则二者的斜率都无定义）。

     对于“仅当”，假设
     两条线段有相同的长度与斜率
     （斜率无定义的情形是容易的；我们处理两条
     线段都有斜率 \( m \) 的情形）。
     另外不失一般性，假设 $a_1<b_1$ 且
     $c_1<d_1$。
     第一条线段是
     \( \overline{(a_1,a_2)(b_1,b_2)}=
     \set{(x,y)\suchthat y=mx+n_1,\; x\in [a_1..b_1]} \)
     （$n_1$ 为某个截距）
     而第二条线段是 \( \overline{(c_1,c_2)(d_1,d_2)}=
     \set{(x,y)\suchthat y=mx+n_2,\; x\in [c_1..d_1]} \)
     （$n_2$ 为某个数）。
     于是这两条线段的长度为
     \begin{equation*}
     \sqrt{(b_1-a_1)^2+((mb_1+n_1)-(ma_1+n_1))^2}
       =\sqrt{(1+m^2)(b_1-a_1)^2}
     \end{equation*}
     以及，类似地，\( \sqrt{(1+m^2)(d_1-c_1)^2} \)。
     因此 \( \absval{b_1-a_1}=\absval{d_1-c_1} \)。
     于是，由于我们假定了 \( a_1<b_1 \) 与 \( c_1<d_1 \)，故有
     \( b_1-a_1=d_1-c_1 \)。

     另一个等式类似。
   
\end{ans}
\begin{ans}{一.II.2.28}
      成立；
      我们可以用归纳法证明它。

      设这些向量都在某个 \( \Re^k \) 中。
      显然，该论断对单个向量成立。
      三角不等式就是该论断用于两个向量的情形。
      归纳步骤假设该论断对 \( n \) 个或更少的
      向量成立。
      那么，由两个向量的三角不等式可得下式
      \begin{equation*}
         \absval{\vec{u}_1+\cdots+\vec{u}_n+\vec{u}_{n+1}}
         \leq
         \absval{\vec{u}_1+\cdots+\vec{u}_n}+\absval{\vec{u}_{n+1}}
      \end{equation*}
      再对右端第一个和项
      应用归纳假设，
      便知它小于等于
      \( \absval{\vec{u}_1}+\cdots+\absval{\vec{u}_n}+\absval{\vec{u}_{n+1}} \)。
    
\end{ans}
\begin{ans}{一.II.2.29}
      按定义
      \begin{equation*}
        \frac{\vec{u}\dotprod\vec{v}}{
            \absval{\vec{u}\,}\,\absval{\vec{v}\,}}=\cos\theta
      \end{equation*}
      其中 \( \theta \) 是两向量之间的夹角。
      于是该比值就是 \( \absval{\cos\theta} \)。
   
\end{ans}
\begin{ans}{一.II.2.30}
      这样，“向量正交当且仅当它们的
      点积为零”这一论断就没有例外了。
    
\end{ans}
\begin{ans}{一.II.2.31}
      我们可用下式
      来求 \( (a) \) 与 \( (b) \)（其中 \( a,b\neq 0 \)）的夹角
      \begin{equation*}
         \arccos(\frac{ab}{\sqrt{a^2}\sqrt{b^2}}).
      \end{equation*}
      若 \( a \) 或 \( b \) 为零，则夹角为 \( \pi/2 \) 弧度。
      否则，若 \( a \) 与 \( b \) 异号，则夹角为
      \( \pi \) 弧度，不然夹角为零~弧度。
   
\end{ans}
\begin{ans}{一.II.2.32}
       \( \vec{u} \) 与 \( \vec{v} \) 的夹角当
       \( \vec{u}\dotprod\vec{v}> 0 \) 时为锐角，当
       \( \vec{u}\dotprod\vec{v}=0 \) 时为直角，当
       \( \vec{u}\dotprod\vec{v}<0 \) 时为钝角。
       这是因为在夹角公式中，分母绝不会
       为负。
     
\end{ans}
\begin{ans}{一.II.2.33}
      设 \( \vec{u},\vec{v}\in\Re^n \)。
      若 \( \vec{u} \) 与 \( \vec{v} \) 垂直，则
      \begin{equation*}
        \absval{\vec{u}+\vec{v}\,}^2
        =(\vec{u}+\vec{v})\dotprod(\vec{u}+\vec{v})
        =\vec{u}\dotprod\vec{u}+2\,\vec{u}\dotprod\vec{v}
               +\vec{v}\dotprod\vec{v}
        =\vec{u}\dotprod\vec{u}+\vec{v}\dotprod\vec{v}
        =\absval{\vec{u}\,}^2+\absval{\vec{v}\,}^2
      \end{equation*}
      （第三个等号成立是因为 \( \vec{u}\dotprod\vec{v}=0 \)）。
    
\end{ans}
\begin{ans}{一.II.2.34}
      当 \( \vec{u},\vec{v}\in\Re^n \) 时，向量
      \( \vec{u}+\vec{v} \) 与 \( \vec{u}-\vec{v} \) 垂直当且
      仅当
      $0=(\vec{u}+\vec{v})\dotprod(\vec{u}-\vec{v})
         =\vec{u}\dotprod\vec{u}-\vec{v}\dotprod\vec{v}$，
      这表明这两个向量垂直当且仅当
      \( \vec{u}\dotprod\vec{u}=\vec{v}\dotprod\vec{v} \)。
      而这又当且仅当 \( \absval{\vec{u}\,}=\absval{\vec{v}\,} \) 时成立。
   
\end{ans}
\begin{ans}{一.II.2.35}
      设 \( \vec{u}\in\Re^n \) 与
      \( \vec{v}\in\Re^n \) 及 \( \vec{w}\in\Re^n \) 都垂直。
      那么，对任意 \( k,m\in\Re \) 我们有下式。
      \begin{equation*}
        \vec{u}\dotprod(k\vec{v}+m\vec{w})
        =k(\vec{u}\dotprod\vec{v})+m(\vec{u}\dotprod\vec{w})
        =k(0)+m(0)=0
      \end{equation*}
    
\end{ans}
\begin{ans}{一.II.2.36}
      我们来证明一个更一般的结论：~若
      \( \absval{\vec{z}_1}=\absval{\vec{z}_2} \)，其中
      \( \vec{z}_1,\vec{z}_2\in\Re^n \)，则 \( \vec{z}_1+\vec{z}_2 \)
      平分 \( \vec{z}_1 \) 与 \( \vec{z}_2 \) 之间的夹角
      \begin{center}
        \setlength{\unitlength}{4pt}      % sum of equal length vectors.
        \begin{picture}(37,12)(0,0)
          \put(0,0){\begin{picture}(12,12)(0,0)
                      \thicklines
                      \put(0,0){\vector(2,1){8} }
                      \put(0,0){\vector(1,2){4} }
                      \put(0,0){\vector(1,1){12} }
                      \thinlines
                      \put(8,4){\line(1,2){4} }
                      \put(4,8){\line(2,1){8} }
                    \end{picture} }
          \put(18.5,7){\makebox(0,0){\small 可得} }
          \put(25,0){\begin{picture}(12,12)(0,0)
                      \thinlines
                      \put(0,0){\line(2,1){8} }
                      \put(0,0){\line(1,2){4} }
                      \put(0,0){\line(1,1){12} }
                      \put(8,4){\line(1,2){4} }
                      \put(4,8){\line(2,1){8} }

                      \put(2,3.8){\( \prime \)}
                      \put(4.2,1.5){\( \prime \)}
                      \multiput(5.7,5.2)(0.5,0.5){2}{\( \prime \)}
                      \multiput(9.0,5.8)(0.5,1.0){3}{\( \prime \)}
                      \multiput(6.6,8.7)(0.6,0.3){3}{\( \prime \)}
                    \end{picture} }
        \end{picture}
      \end{center}
      （我们忽略 \( \vec{z}_1 \) 与 \( \vec{z}_2 \) 为
      零向量的情形）。

      \( \vec{z}_1+\vec{z}_2=\zero \) 的情形很容易。
      对其余情形，由夹角的定义，
      只需证明下式即告完成。
      \begin{equation*}
        \frac{\vec{z}_1\dotprod(\vec{z}_1+\vec{z}_2)}{
              \absval{\vec{z}_1}\,\absval{\vec{z}_1+\vec{z}_2} }
        =
        \frac{\vec{z}_2\dotprod(\vec{z}_1+\vec{z}_2)}{
              \absval{\vec{z}_2}\,\absval{\vec{z}_1+\vec{z}_2} }
      \end{equation*}
      而在每个表达式内部展开则得
      \begin{equation*}
        \frac{\vec{z}_1\dotprod\vec{z}_1+\vec{z}_1\dotprod\vec{z}_2}{
              \absval{\vec{z}_1}\,\absval{\vec{z}_1+\vec{z}_2} }
        \qquad
        \frac{\vec{z}_2\dotprod\vec{z}_1+\vec{z}_2\dotprod\vec{z}_2}{
              \absval{\vec{z}_2}\,\absval{\vec{z}_1+\vec{z}_2} }
      \end{equation*}
      且 \( \vec{z}_1\dotprod\vec{z}_1=\absval{\vec{z}_1}^2
              =\absval{\vec{z}_2}^2=\vec{z}_2\dotprod\vec{z}_2 \)，故
      两者相等。
    
\end{ans}
\begin{ans}{一.II.2.37}
      我们可以把这两个论断一并证明。
      设 \( \vec{u}, \vec{v}\in\Re^n \)，写
      \begin{equation*}
         \vec{u}=\colvec{u_1 \\ \vdotswithin{u_1} \\ u_n}
         \qquad
         \vec{v}=\colvec{v_1 \\ \vdotswithin{v_1} \\ v_n}
      \end{equation*}
      并计算。
      \begin{equation*}
        \cos\theta=
        \frac{ku_1v_1+\cdots+ku_nv_n}{
             \sqrt{{(ku_1)}^2+\cdots+{(ku_n)}^2}\sqrt{{b_1}^2+\cdots+{b_n}^2} }
        =\frac{k}{\absval{k}}
        \frac{\vec{u}\cdot\vec{v}}{\absval{\vec{u}\,}\,\absval{\vec{v}\,} }
        =\pm
        \frac{\vec{u}\dotprod\vec{v}}{\absval{\vec{u}\,}\,\absval{\vec{v}\,} }
      \end{equation*}
    
\end{ans}
\begin{ans}{一.II.2.38}
      设
      \begin{equation*}
        \vec{u}=\colvec{u_1 \\ \vdotswithin{u_1} \\ u_n},
        \quad
        \vec{v}=\colvec{v_1 \\ \vdotswithin{v_1} \\ v_n}
        \quad
        \vec{w}=\colvec{w_1 \\ \vdotswithin{w_1} \\ w_n}
      \end{equation*}
      于是
      \begin{align*}
        \vec{u}\dotprod\bigl(k\vec{v}+m\vec{w}\bigr)
        &=\colvec{u_1 \\ \vdotswithin{u_1} \\ u_n}\dotprod
         \bigl( \colvec{kv_1 \\ \vdotswithin{kv_1} \\ kv_n}
               +\colvec{mw_1 \\ \vdotswithin{mw_1} \\ mw_n} \bigr)   \\
        &=\colvec{u_1 \\ \vdotswithin{u_1} \\ u_n}\dotprod
         \colvec{kv_1+mw_1 \\ \vdotswithin{kv_1+mw_1} \\ kv_n+mw_n}    \\
        &=u_1(kv_1+mw_1)+\cdots+u_n(kv_n+mw_n)    \\
        &=ku_1v_1+mu_1w_1+\cdots+ku_nv_n+mu_nw_n    \\
        &=(ku_1v_1+\cdots+ku_nv_n)+(mu_1w_1+\cdots+mu_nw_n)    \\
        &=k(\vec{u}\dotprod\vec{v})+m(\vec{u}\dotprod\vec{w})
      \end{align*}
      这正是所要证的。
    
\end{ans}
\begin{ans}{一.II.2.39}
      \begin{exparts}
        \item 例如，生日为10~月$12$日时得到下式。
          \begin{equation*}
            \theta
            =
            \arccos(\,\frac{\colvec[r]{7 \\ 12}\dotprod\colvec[r]{10 \\ 12}}{
                  \absval{\colvec[r]{7 \\ 12}\,}\cdot\absval{\colvec[r]{10 \\ 12}\,} }\,)
            =\arccos(\frac{214}{\sqrt{244}\sqrt{193}})
            \approx \text{$0.17$~弧度}
          \end{equation*}
        \item 把同一公式用于 $\rowvec{9 &19}$，得
           约 $0.09$~弧度。
        \item 若对方生于
           同一天，则夹角为 $0$~弧度。
           若一个生日是另一个生日的标量倍数，夹角也会为 $0$。
           例如，生于3~月$6$日的人
           与生于2~月$4$日的人关系和睦。

          给定一个生日，我们可以让 \Sage{}
          绘出其他日期所对应的夹角。
          本例展示了所有日期与7~月12日的关系。
\begin{lstlisting}
  sage: plot3d(lambda x, y: math.acos((x*7+y*12)/(math.sqrt(7**2+12**2)*math.sqrt(x**2+y**2))),
                                                                 (1,12),(1,31))
\end{lstlisting}
          % The result looks like this.
          \begin{center}
            \includegraphics[height=2in]{gr/bin/rcbday.jpg}
          \end{center}
        \item 我们希望最大化下式。
          \begin{equation*}
            \theta
            =
            \arccos(\,\frac{\colvec[r]{7 \\ 12}\dotprod\colvec{m \\ d}}{
                  \absval{\colvec[r]{7 \\ 12}\,}\cdot\absval{\colvec{m \\ d}\,} }\,)
          \end{equation*}
          当然，$m$ 与 $d$ 都不能取负值，因此我们
          无法得到与给定向量正交的向量。
          下面的 Python 脚本用暴力法求出最大夹角。
\begin{lstlisting}
  import math
  days={1:31,  # Jan
        2:29, 3:31, 4:30, 5:31, 6:30, 7:31, 8:31, 9:30, 10:31, 11:30, 12:31}
  BDAY=(7,12)
  max_res=0
  max_res_date=(-1,-1)
  for month in range(1,13):
      for day in range(1,days[month]+1):
          num=BDAY[0]*month+BDAY[1]*day
          denom=math.sqrt(BDAY[0]**2+BDAY[1]**2)*math.sqrt(month**2+day**2)
          if denom>0:
              res=math.acos(min(num*1.0/denom,1))
              print "day:",str(month),str(day)," angle:",str(res)
              if res>max_res:
                  max_res=res
                  max_res_date=(month,day)
  print "For ",str(BDAY),"worst case",str(max_res),"rads, date",str(max_res_date)
  print "  That is ",180*max_res/math.pi,"degrees"
\end{lstlisting}
          结果为
\begin{lstlisting}
  For  (7, 12) worst case 0.95958064648 rads, date (12, 1)
    That is  54.9799211457 degrees
\end{lstlisting}

          更具概念性的方法是考察所有点
          $(\text{月},\text{日})$ 与点 $(7,12)$ 的关系。
          下图
          清楚地表明答案是12~月$1$日或1~月$31$日，取决于
          哪个离生日更远。
          虚线平分从原点到
          12~月$1$日的直线与从原点到1~月$31$日的直线之间的夹角。
          直线以上的生日离12~月$1$日最远，
          直线以下的生日离1~月$31$日最远。
          \begin{center}
            \includegraphics{gr/mp/ch1.54}
          \end{center}
      \end{exparts}
    
\end{ans}
\begin{ans}{一.II.2.40}
      \answerasgiven %
      风的实际速度 \( \vec{v} \) 等于
      船的速度与风的视速度之和。
      不失一般性，可设 \( \vec{a} \) 与
      \( \vec{b} \) 为单位向量，并可写成
      \begin{equation*}
         \vec{v}=\vec{v}_1+s\vec{a}=\vec{v}_2+t\vec{b}
      \end{equation*}
      其中 \( s \) 与 \( t \) 为待定标量。
      先与 \( \vec{a} \) 作点积、再与 \( \vec{b} \) 作点积
      以得到
      \begin{align*}
         s-t\vec{a}\dotprod\vec{b}
         &=\vec{a}\dotprod(\vec{v}_2-\vec{v}_1)    \\
         s\vec{a}\dotprod\vec{b}-t
         &=\vec{b}\dotprod(\vec{v}_2-\vec{v}_1)
      \end{align*}
      将第二式乘以 \( \vec{a}\dotprod\vec{b} \)，
      再从第一式中减去所得结果，便求得
      \begin{equation*}
         s=
         \frac{[\vec{a}-(\vec{a}\dotprod\vec{b})\vec{b}]
                   \dotprod(\vec{v}_2-\vec{v}_1)
              }{1-(\vec{a}\dotprod\vec{b})^2}.
      \end{equation*}
      代回原来的方程，便得
      \begin{equation*}
         \vec{v}=\vec{v}_1+
         \frac{[\vec{a}-(\vec{a}\dotprod\vec{b})\vec{b}]
                 \dotprod(\vec{v}_2-\vec{v}_1)
             \vec{a}}{1-(\vec{a}\dotprod\vec{b})^2}.
      \end{equation*}
    
\end{ans}
\begin{ans}{一.II.2.41}
       我们对 \( n \) 使用归纳法。

       在 \( n=1 \) 的基底情形，该恒等式化为
       \begin{equation*}
          (a_1b_1)^2=({a_1}^2)({b_1}^2)-0
       \end{equation*}
       显然成立。

       归纳步骤假设该公式对
       \( 0 \), \ldots, \( n \) 各情形成立。
       我们将证明它在 \( n+1 \) 情形也成立。
       从右端出发
       \begin{multline*}
         \bigl( \sum_{1\leq j\leq n+1}{a_j}^2\bigr)
         \bigl( \sum_{1\leq j\leq n+1}{b_j}^2\bigr)
         -
         \sum_{1\leq k<j\leq n+1}\bigl(a_kb_j-a_jb_k\bigr)^2  \\
       \begin{aligned}
         &=
         \bigl[ (\sum_{1\leq j\leq n}{a_j}^2)+{a_{n+1}}^2\bigr]
         \bigl[ (\sum_{1\leq j\leq n}{b_j}^2)+{b_{n+1}}^2\bigr]   \\
         &\hbox{}\quad -
         \bigl[\sum_{1\leq k<j\leq n}\bigl(a_kb_j-a_jb_k\bigr)^2+
         \sum_{1\leq k\leq n}\bigl(a_kb_{n+1}-a_{n+1}b_k\bigr)^2  \bigr] \\
         &=
         \bigl( \sum_{1\leq j\leq n}{a_j}^2\bigr)
         \bigl( \sum_{1\leq j\leq n}{b_j}^2\bigr)
         +
         \sum_{1\leq j\leq n}{b_j}^2{a_{n+1}}^2
         +
         \sum_{1\leq j\leq n}{a_j}^2{b_{n+1}}^2
         +
         {a_{n+1}}^2{b_{n+1}}^2                                 \\
         &\hbox{}\qquad\hbox{} -
         \bigl[\sum_{1\leq k<j\leq n}\bigl(a_kb_j-a_jb_k\bigr)^2+
         \sum_{1\leq k\leq n}\bigl(a_kb_{n+1}-a_{n+1}b_k\bigr)^2  \bigr] \\
         &=
         \bigl( \sum_{1\leq j\leq n}{a_j}^2\bigr)
         \bigl( \sum_{1\leq j\leq n}{b_j}^2\bigr)
         -\sum_{1\leq k<j\leq n}\bigl(a_kb_j-a_jb_k\bigr)^2   \\
         &\hbox{}\quad +
         \sum_{1\leq j\leq n}{b_j}^2{a_{n+1}}^2
         +
         \sum_{1\leq j\leq n}{a_j}^2{b_{n+1}}^2
         +
         {a_{n+1}}^2{b_{n+1}}^2                                 \\
         &\hbox{}\qquad\hbox{} -
         \sum_{1\leq k\leq n}\bigl(a_kb_{n+1}-a_{n+1}b_k\bigr)^2
        \end{aligned}
       \end{multline*}
       再应用归纳假设。
       \begin{align*}
         &=
         \bigl( \sum_{1\leq j\leq n}a_jb_j\bigr)^2
         +
         \sum_{1\leq j\leq n}{b_j}^2{a_{n+1}}^2
         +
         \sum_{1\leq j\leq n}{a_j}^2{b_{n+1}}^2
         +
         {a_{n+1}}^2{b_{n+1}}^2                          \\
         &\hbox{}\qquad\hbox{}-
         \bigl[\sum_{1\leq k\leq n}{a_k}^2{b_{n+1}}^2
           -2\sum_{1\leq k\leq n}a_kb_{n+1}a_{n+1}b_k
           +\sum_{1\leq k\leq n}{a_{n+1}}^2{b_k}^2\bigr]        \\
         &=
         \bigl( \sum_{1\leq j\leq n}a_jb_j\bigr)^2
         +2\bigl(\sum_{1\leq k\leq n}a_kb_{n+1}a_{n+1}b_k\bigr)
         +{a_{n+1}}^2{b_{n+1}}^2                                 \\
         &=
         \bigl[\bigl(\sum_{1\leq j\leq n}a_jb_j\bigr)+a_{n+1}b_{n+1}\bigr]^2
       \end{align*}
       即得左端。
   
\end{ans}
\protect \section {第 III 节：简化阶梯形}
\protect \subsection {一.III.1: 高斯–若尔当消元}
\begin{ans}{一.III.1.8}
       这些解答只给出了高斯–若尔当消元的过程。
       有了它，描述解集就很容易了。
       \begin{exparts}
         \partsitem 解集只含一个元素。
           \begin{multline*}
             \begin{amat}[r]{2}
               1  &1  &2  \\
               1  &-1 &0
             \end{amat}
             \grstep{-\rho_1+\rho_2}
             \begin{amat}[r]{2}
               1  &1  &2  \\
               0  &-2 &-2
             \end{amat}                          \\
             \grstep{-(1/2)\rho_2}
             \begin{amat}[r]{2}
               1  &1  &2  \\
               0  &1  &1
             \end{amat}
             \grstep{-\rho_2+\rho_1}
             \begin{amat}[r]{2}
               1  &0  &1  \\
               0  &1  &1
             \end{amat}
           \end{multline*}
         \partsitem 解集含有一个参数。
            \begin{equation*}
             \begin{amat}[r]{3}
               1  &0  &-1  &4  \\
               2  &2  &0   &1
             \end{amat}
             \grstep{-2\rho_1+\rho_2}
             \begin{amat}[r]{3}
               1  &0  &-1  &4  \\
               0  &2  &2   &-7
             \end{amat}
             \grstep{(1/2)\rho_2}
             \begin{amat}[r]{3}
               1  &0  &-1  &4  \\
               0  &1  &1   &-7/2
             \end{amat}
           \end{equation*}
         \partsitem 有唯一解。
           \begin{multline*}
             \begin{amat}[r]{2}
               3  &-2  &1  \\
               6  &1   &1/2
             \end{amat}
             \grstep{-2\rho_1+\rho_2}
             \begin{amat}[r]{2}
               3  &-2  &1  \\
               0  &5   &-3/2
             \end{amat}                            \\
             \grstep[(1/5)\rho_2]{(1/3)\rho_1}
             \begin{amat}[r]{2}
               1  &-2/3&1/3 \\
               0  &1   &-3/10
             \end{amat}
             \grstep{(2/3)\rho_2+\rho_1}
             \begin{amat}[r]{2}
               1  &0   &2/15 \\
               0  &1   &-3/10
             \end{amat}
          \end{multline*}
         \partsitem 第二步作一次对换两行的变换，会使算术更简单。
          \begin{multline*}
           \begin{amat}[r]{3}
             2  &-1  &0  &-1  \\
             1  &3   &-1 &5   \\
             0  &1   &2  &5
           \end{amat}
           \grstep{-(1/2)\rho_1+\rho_2}
           \begin{amat}[r]{3}
             2  &-1  &0  &-1   \\
             0  &7/2 &-1 &11/2 \\
             0  &1   &2  &5
           \end{amat}                                \\
           \begin{aligned}
           &\grstep{\rho_2\leftrightarrow\rho_3}
           \begin{amat}[r]{3}
             2  &-1  &0  &-1   \\
             0  &1   &2  &5    \\
             0  &7/2 &-1 &11/2
           \end{amat}
             \grstep{-(7/2)\rho_2+\rho_3}
             \begin{amat}[r]{3}
               2  &-1  &0  &-1   \\
               0  &1   &2  &5    \\
               0  &0   &-8 &-12
             \end{amat}                               \\
             &\grstep[-(1/8)\rho_2]{(1/2)\rho_1}
             \begin{amat}[r]{3}
               1  &-1/2&0  &-1/2 \\
               0  &1   &2  &5    \\
               0  &0   &1  &3/2
             \end{amat}
             \grstep{-2\rho_3+\rho_2}
             \begin{amat}[r]{3}
               1  &-1/2&0  &-1/2 \\
               0  &1   &0  &2    \\
               0  &0   &1  &3/2
             \end{amat}                                    \\
             &\grstep{(1/2)\rho_2+\rho_1}
             \begin{amat}[r]{3}
               1  &0   &0  &1/2  \\
               0  &1   &0  &2    \\
               0  &0   &1  &3/2
             \end{amat}
           \end{aligned}
         \end{multline*}
       \end{exparts}
     
\end{ans}
\begin{ans}{一.III.1.9}
       \begin{exparts}
         \partsitem
           \begin{multline*}
             \begin{amat}{3}
               1 &1  &-1 &3 \\
               2 &-1 &-1 &1 \\
               3 &1  &2  &0
             \end{amat}
             \grstep[-3\rho_1+\rho_3]{-2\rho_1+\rho_2}
             \begin{amat}{3}
               1 &1  &-1 &3 \\
               0 &-3 &1  &-5 \\
               0 &-2  &5  &-9
             \end{amat}                             \\
           \begin{aligned}
             &\grstep{-(2/3)\rho_2+\rho_3}
             \begin{amat}{3}
               1 &1  &-1    &3 \\
               0 &-3 &1     &-5 \\
               0 &0  &13/3  &-17/3
             \end{amat}
             \grstep[(3/13)\rho_3]{-(1/3)\rho_2}
             \begin{amat}{3}
               1 &1  &-1    &3 \\
               0 &1  &-1/3    &5/3 \\
               0 &0  &1      &-17/13
             \end{amat}                                \\
             &\grstep[(1/3)\rho_3+\rho_2]{\rho_3+\rho_1}
             \begin{amat}{3}
               1 &1  &0    &22/13 \\
               0 &1  &0    &16/13 \\
               0 &0  &1      &-17/13
             \end{amat}
             \grstep{-\rho_2+\rho_1}
             \begin{amat}{3}
               1 &0  &0    &6/13 \\
               0 &1  &0    &16/13 \\
               0 &0  &1      &-17/13
             \end{amat}
           \end{aligned}
           \end{multline*}
         \partsitem
           \begin{multline*}
             \begin{amat}{3}
               1 &1  &2 &0 \\
               2 &-1 &1 &1 \\
               4 &1  &5 &1
             \end{amat}
             \grstep[-4\rho_1+\rho_3]{-2\rho_1+\rho_2}
             \begin{amat}{3}
               1 &1  &2  &0 \\
               0 &-3 &-3 &1 \\
               0 &-3 &-3 &1
             \end{amat}
             \grstep{-\rho_2+\rho_3}
             \begin{amat}{3}
               1 &1  &2  &0 \\
               0 &-3 &-3 &1 \\
               0 &0  &0  &0
             \end{amat}                              \\
             \grstep{-(1/3)\rho_2}
             \begin{amat}{3}
               1 &1  &2  &0 \\
               0 &1  &1 &-1/3 \\
               0 &0  &0  &0
             \end{amat}
             \grstep{-\rho_2+\rho_1}
             \begin{amat}{3}
               1 &0  &1  &1/3 \\
               0 &1  &1 &-1/3 \\
               0 &0  &0  &0
             \end{amat}
           \end{multline*}
       \end{exparts}
     
\end{ans}
\begin{ans}{一.III.1.10}
      使用高斯–若尔当消元。
      \begin{exparts}
        \partsitem 简化阶梯形除了由1构成的对角线以外全为0。
          \begin{equation*}
            \grstep{-(1/2)\rho_1+\rho_2}
            \begin{mat}[r]
              2  &1  \\
              0  &5/2
            \end{mat}
            \grstep[(2/5)\rho_2]{(1/2)\rho_1}
            \begin{mat}[r]
              1  &1/2\\
              0  &1
            \end{mat}
            \grstep{-(1/2)\rho_2+\rho_1}
            \begin{mat}[r]
              1  &0  \\
              0  &1
            \end{mat}
          \end{equation*}
        \partsitem 与上一题一样，简化阶梯形
          除了由1构成的对角线以外全是0。
          \begin{multline*}
            \grstep[\rho_1+\rho_3]{-2\rho_1+\rho_2}
            \begin{mat}[r]
              1  &3  &1  \\
              0  &-6 &2  \\
              0  &0  &-2
            \end{mat}
            \grstep[-(1/2)\rho_3]{-(1/6)\rho_2}
            \begin{mat}[r]
              1  &3  &1     \\
              0  &1  &-1/3  \\
              0  &0  &1
            \end{mat}                                      \\
            \grstep[-\rho_3+\rho_1]{(1/3)\rho_3+\rho_2}
            \begin{mat}[r]
              1  &3  &0     \\
              0  &1  &0     \\
              0  &0  &1
            \end{mat}
            \grstep{-3\rho_2+\rho_1}
            \begin{mat}[r]
              1  &0  &0     \\
              0  &1  &0     \\
              0  &0  &1
            \end{mat}
           \end{multline*}
        \partsitem 列数比行数多，所以我们得到的
          不只是由1构成的一条对角线。
          \begin{multline*}
            \grstep[-3\rho_1+\rho_3]{-\rho_1+\rho_2}
            \begin{mat}[r]
              1  &0  &3  &1  &2  \\
              0  &4  &-1 &0  &3  \\
              0  &4  &-1 &-2 &-4
            \end{mat}
            \grstep{-\rho_2+\rho_3}
            \begin{mat}[r]
              1  &0  &3  &1  &2  \\
              0  &4  &-1 &0  &3  \\
              0  &0  &0  &-2 &-7
            \end{mat}                           \\
            \grstep[-(1/2)\rho_3]{(1/4)\rho_2}
            \begin{mat}[r]
              1  &0  &3    &1  &2  \\
              0  &1  &-1/4 &0  &3/4  \\
              0  &0  &0    &1  &7/2
            \end{mat}
            \grstep{-\rho_3+\rho_1}
            \begin{mat}[r]
              1  &0  &3    &0  &-3/2  \\
              0  &1  &-1/4 &0  &3/4     \\
              0  &0  &0    &1  &7/2
            \end{mat}
          \end{multline*}
        \partsitem 与上一小题一样，这不是一个方阵。
          \begin{multline*}
            \grstep{\rho_1\leftrightarrow\rho_3}
            \begin{mat}[r]
              1  &5  &1  &5  \\
              0  &0  &5  &6  \\
              0  &1  &3  &2
            \end{mat}
            \grstep{\rho_2\leftrightarrow\rho_3}
            \begin{mat}[r]
              1  &5  &1  &5  \\
              0  &1  &3  &2  \\
              0  &0  &5  &6
            \end{mat}                    \\
            \begin{aligned}
              &\grstep{(1/5)\rho_3}
              \begin{mat}[r]
                1  &5  &1  &5  \\
                0  &1  &3  &2  \\
                0  &0  &1  &6/5
              \end{mat}
              \grstep[-\rho_3+\rho_1]{-3\rho_3+\rho_2}
              \begin{mat}[r]
                1  &5  &0  &19/5  \\
                0  &1  &0  &-8/5  \\
                0  &0  &1  &6/5
              \end{mat}                   \\
              &\grstep{-5\rho_2+\rho_1}
              \begin{mat}[r]
                1  &0  &0  &59/5  \\
                0  &1  &0  &-8/5  \\
                0  &0  &1  &6/5
              \end{mat}
            \end{aligned}
          \end{multline*}
      \end{exparts}
    
\end{ans}
\begin{ans}{一.III.1.11}
      \begin{exparts}
        \partsitem 先作对换。
          \begin{equation*}
            \grstep{\rho_1\leftrightarrow\rho_2}
            \grstep{\rho_1+\rho_3}
            \grstep{\rho_2+\rho_3}
            \grstep[(1/2)\rho_2 \\ (1/2)\rho_3]{(1/2)\rho_1}
            \grstep[(-1/2)\rho_3+\rho_2]{(-1/2)\rho_3+\rho_1}
            \grstep{(1/2)\rho_2+\rho_1}
            \begin{mat}
              1 &0 &0 \\
              0 &1 &0 \\
              0 &0 &1
            \end{mat}
          \end{equation*}
        \partsitem 这里对换在中间进行。
        \begin{equation*}
          \grstep[\rho_1+\rho_3]{-2\rho_1+\rho_2}
          \grstep{\rho_2\leftrightarrow\rho_3}
          \grstep{(1/3)\rho_2}
          \grstep{-3\rho_2+\rho_1}
          \begin{mat}
            1 &0 &0   \\
            0 &1 &1/3 \\
            0 &0 &0
          \end{mat}
        \end{equation*}
      \end{exparts}
    
\end{ans}
\begin{ans}{一.III.1.12}
      至于高斯消元的那一半，
      参见第一章第~I.2 节习题
      \nearbyexercise{exer:SlvMatNot} 的解答。
      \begin{exparts}
      \partsitem “若尔当”那一半如下进行。
        \begin{equation*}
          \grstep[-(1/3)\rho_2]{(1/2)\rho_1}
          \begin{amat}[r]{3}
            1  &1/2 &-1/2 &1/2  \\
            0  &1   &-2/3 &-1/3
          \end{amat}
          \grstep{-(1/2)\rho_2+\rho_1}
          \begin{amat}[r]{3}
            1  &0   &-1/6 &2/3  \\
            0  &1   &-2/3 &-1/3
          \end{amat}
        \end{equation*}
        解集为下面这个集合
        \begin{equation*}
          \set{\colvec[r]{2/3 \\ -1/3 \\ 0}
               +\colvec[r]{1/6 \\ 2/3 \\ 1}z
              \suchthat z\in\Re}
        \end{equation*}
      \partsitem 后一半是
        \begin{equation*}
          \grstep{\rho_3+\rho_2}
          \begin{amat}[r]{4}
            1  &0  &-1  &0  &1 \\
            0  &1  &2   &0  &3 \\
            0  &0  &0   &1  &0
          \end{amat}
        \end{equation*}
        于是解如下。
        \begin{equation*}
          \set{\colvec[r]{1 \\ 3 \\ 0 \\ 0}
               +\colvec[r]{1 \\ -2 \\ 1 \\ 0}z
              \suchthat z\in\Re}
        \end{equation*}
      \partsitem 这个若尔当一半
        \begin{equation*}
          \grstep{\rho_2+\rho_1}
          \begin{amat}[r]{4}
            1  &0  &1   &1  &0 \\
            0  &1  &0   &1  &0 \\
            0  &0  &0   &0  &0 \\
            0  &0  &0   &0  &0
          \end{amat}
        \end{equation*}
        给出
        \begin{equation*}
          \set{\colvec[r]{0 \\ 0 \\ 0 \\ 0}
               +\colvec[r]{-1 \\ 0 \\ 1 \\ 0}z
               +\colvec[r]{-1 \\ -1 \\ 0 \\ 1}w
              \suchthat z,w\in\Re}
        \end{equation*}
        （当然，我们可以在描述中省去零向量）。
      \partsitem “若尔当”那一半
        \begin{align*}
          &\grstep{-(1/7)\rho_2}
          \begin{amat}[r]{5}
            1  &2  &3   &1   &-1   &1  \\
            0  &1  &8/7 &2/7 &-4/7 &0
          \end{amat}                   \\
          &\grstep{-2\rho_2+\rho_1}
          \begin{amat}[r]{5}
            1  &0  &5/7 &3/7 &1/7  &1  \\
            0  &1  &8/7 &2/7 &-4/7 &0
          \end{amat}
        \end{align*}
        以这个解集收尾。
        \begin{equation*}
          \set{\colvec[r]{1 \\ 0 \\ 0 \\ 0 \\ 0}
               +\colvec[r]{-5/7 \\ -8/7 \\ 1 \\ 0 \\ 0}c
               +\colvec[r]{-3/7 \\ -2/7 \\ 0 \\ 1 \\ 0}d
               +\colvec[r]{-1/7 \\ 4/7 \\ 0 \\ 0 \\ 1}e
              \suchthat c,d,e\in\Re}
        \end{equation*}
    \end{exparts}
   
\end{ans}
\begin{ans}{一.III.1.13}
      按常规作高斯消元法便得到一个：
      \begin{equation*}
        \grstep[-(1/2)\rho_1+\rho_3]{-3\rho_1+\rho_2}
        \begin{mat}[r]
          2  &1  &1  &3  \\
          0  &1  &-2 &-7 \\
          0  &9/2&1/2&7/2
        \end{mat}
        \grstep{-(9/2)\rho_2+\rho_3}
        \begin{mat}[r]
          2  &1  &1    &3  \\
          0  &1  &-2   &-7 \\
          0  &0  &19/2 &35
        \end{mat}
      \end{equation*}
      而任何外观上的改动，比如把最下面一行乘以 \( 2 \)，
      \begin{equation*}
        \begin{mat}[r]
          2  &1  &1    &3  \\
          0  &1  &-2   &-7 \\
          0  &0  &19   &70
        \end{mat}
      \end{equation*}
      便给出另一个。
    
\end{ans}
\begin{ans}{一.III.1.14}
      在下面列出的各情形中，我们取 $a,b\in\Re$。
      因此，下面列出的某些标准形
      实际上包含了无穷多种情形。
      特别地，它们包括了 $a=0$ 与 $b=0$ 的情形。
      \begin{exparts}
        \partsitem
          $\begin{mat}[r]
            0  &0  \\
            0  &0
          \end{mat}$,
          $\begin{mat}[r]
            1  &a  \\
            0  &0
          \end{mat}$,
          $\begin{mat}[r]
            0  &1  \\
            0  &0
          \end{mat}$,
          $\begin{mat}[r]
            1  &0  \\
            0  &1
          \end{mat}$
        \partsitem
          $\begin{mat}[r]
               0  &0  &0  \\
               0  &0  &0
             \end{mat}$,
          $\begin{mat}[r]
               1  &a  &b  \\
               0  &0  &0
             \end{mat}$,
          $\begin{mat}[r]
               0  &1  &a  \\
               0  &0  &0
             \end{mat}$,
          $\begin{mat}[r]
               0  &0  &1  \\
               0  &0  &0
             \end{mat}$,
          $\begin{mat}[r]
               1  &0  &a  \\
               0  &1  &b
             \end{mat}$,
          $\begin{mat}[r]
               1  &a  &0  \\
               0  &0  &1
             \end{mat}$,
          $\begin{mat}[r]
               0  &1  &0  \\
               0  &0  &1
             \end{mat}$
        \partsitem
          $\begin{mat}[r]
               0  &0  \\
               0  &0  \\
               0  &0
             \end{mat}$,
          $\begin{mat}[r]
               1  &a  \\
               0  &0  \\
               0  &0
             \end{mat}$,
          $\begin{mat}[r]
               0  &1  \\
               0  &0  \\
               0  &0
             \end{mat}$,
          $\begin{mat}[r]
               1  &0  \\
               0  &1  \\
               0  &0
             \end{mat}$
        \partsitem
          $\begin{mat}[r]
               0  &0  &0  \\
               0  &0  &0  \\
               0  &0  &0
             \end{mat}$,
          $\begin{mat}[r]
               1  &a  &b  \\
               0  &0  &0  \\
               0  &0  &0
             \end{mat}$,
          $\begin{mat}[r]
               0  &1  &a  \\
               0  &0  &0  \\
               0  &0  &0
             \end{mat}$,
          $\begin{mat}[r]
               0  &1  &0  \\
               0  &0  &1  \\
               0  &0  &0
             \end{mat}$,
          $\begin{mat}[r]
               0  &0  &1  \\
               0  &0  &0  \\
               0  &0  &0
             \end{mat}$,
          $\begin{mat}[r]
               1  &0  &a  \\
               0  &1  &b  \\
               0  &0  &0
             \end{mat}$,
          $\begin{mat}[r]
               1  &a  &0  \\
               0  &0  &1  \\
               0  &0  &0
             \end{mat}$,
          $\begin{mat}[r]
               1  &0  &0  \\
               0  &1  &0  \\
               0  &0  &1
             \end{mat}$
      \end{exparts}
    
\end{ans}
\begin{ans}{一.III.1.15}
      非奇异的齐次线性方程组有唯一解。
      所以非奇异矩阵必定能化简成一个（方）
      矩阵，它除了从左上到右下的对角线上是 \( 1 \) 之外
      全为 \( 0 \)，
      比如下面这些。
      \begin{equation*}
         \begin{mat}[r]
           1  &0  \\
           0  &1  \\
         \end{mat}
         \qquad
         \begin{mat}[r]
           1  &0  &0  \\
           0  &1  &0  \\
           0  &0  &1
         \end{mat}
      \end{equation*}
    
\end{ans}
\begin{ans}{一.III.1.16}
     \begin{exparts}
       \partsitem  这是一个等价关系。

          若两个矩阵相关，我们可以写成 $M_1\sim M_2$。
          $\sim$~关系是自反的，因为任何矩阵与它自身
          有相同的 $1,1$ 元素。
          该关系是对称的，因为如果 $M_1$ 与~$M_2$ 有相同的 $1,1$
          元素，那么显然 $M_2$ 与~$M_1$ 也有相同的
          $1,1$ 元素。
          最后，该关系是传递的，因为如果 $M_1\sim M_2$，即两者
          有相同的 $1,1$ 元素，
          且 $M_2\sim M_3$，即两者
          也有相同的元素，那么三者都有相同的~$1,1$ 元素，
          从而 $M_1\sim M_3$。
       % \partsitem  This is an equivalence.
       %    Write $M_1\sim M_2$ if they are related.

       %    This relation is reflexive because any matrix has the
       %    same sum of entries as itself.
       %    The $\sim$~relation is symmetric because if $M_1$ has the same
       %    sum of
       %    entries as~$M_2$, then also $M_2$ has the same sum of entries
       %    as~$M_1$.
       %    Finally, the relation is transitive because if $M_1\sim M_2$ so their
       %    entry sum is the same,
       %    and $M_2\sim M_3$ so they have the
       %    same entry sum as each other, then all three have the same entry sum,
       %    and $M_1\sim M_3$.
       \partsitem
          这不是等价关系，因为它不是传递的。
          下面第一个矩阵与第二个矩阵因 $1,1$ 元素而相关，
          第二个与第三个因 $2,2$~元素而相关。
          但第一个与第三个不相关。
          \begin{equation*}
            \begin{mat}
              1 &0 \\
              0 &0
            \end{mat}
            \quad
            \begin{mat}
              1 &0 \\
              0 &-1
            \end{mat}
            \quad
            \begin{mat}
              0 &0 \\
              0 &-1
            \end{mat}
          \end{equation*}
     \end{exparts}
   
\end{ans}
\begin{ans}{一.III.1.17}
      它是等价关系。
      为证明这一点，我们必须验证该关系
      是自反的、对称的和传递的。

      设所有矩阵都是 $\nbyn{2}$ 的。
      自反性：我们注意到矩阵与它自身
      的元素之和相同。
      对称性：设 $A$ 与~$B$ 的元素之和相同，
      那么显然 $B$ 与~$A$ 的元素之和也相同。
      传递性同样不难\Dash 如果 $A$ 与 $B$ 的元素之和相同，
      且 $B$ 与 $C$ 的元素之和相同，
      那么 $A$ 与 $C$ 的元素之和相同。
    
\end{ans}
\begin{ans}{一.III.1.18}
    \begin{exparts}
      \partsitem 例如，
        \begin{equation*}
          \begin{mat}[r]
            1  &2  \\
            3  &4
          \end{mat}
          \grstep{-\rho_1+\rho_1}
          \begin{mat}[r]
            0  &0  \\
            3  &4
          \end{mat}
          \grstep{\rho_1+\rho_1}
          \begin{mat}[r]
            0  &0  \\
            3  &4
          \end{mat}
        \end{equation*}
        矩阵被改变了。
      \partsitem 这个操作
        \begin{equation*}
          \begin{mat}
            \vdotswithin{a_{i,1}}                     \\
            a_{i,1}  &\cdots  &a_{i,n}  \\
            \vdotswithin{a_{i,1}}                     \\
            a_{j,1}  &\cdots  &a_{j,n}  \\
            \vdotswithin{a_{i,1}}
          \end{mat}
          \grstep{k\rho_i+\rho_j}
          \begin{mat}
            \vdotswithin{a_{i,1}}                                      \\
            a_{i,1}           &\cdots  &a_{i,n}          \\
            \vdotswithin{a_{i,1}}                                      \\
            ka_{i,1}+a_{j,1}  &\cdots  &ka_{i,n}+a_{j,n}  \\
            \vdotswithin{a_{i,1}}
          \end{mat}
        \end{equation*}
        由于 $i\neq j$ 的限制而使第$i$行保持不变。
        因为第$i$行未变，所以这个操作
        \begin{equation*}
          \grstep{-k\rho_i+\rho_j}
          \begin{mat}
            \vdotswithin{a_{i,1}}                                      \\
            a_{i,1}           &\cdots  &a_{i,n}          \\
            \vdotswithin{a_{i,1}}                                      \\
            -ka_{i,1}+ka_{i,1}+a_{j,1}  &\cdots &-ka_{i,n}+ka_{i,n}+a_{j,n} \\
            \vdotswithin{a_{i,1}}
          \end{mat}
        \end{equation*}
        把第$j$行还原到它原来的状态。
        % (If $i=j$ then the third matrix would have entries of the
        % form $-k(ka_{i,j}+a_{i,j})+ka_{i,j}+a_{i,j}$.)
    \end{exparts}
   
\end{ans}
\begin{ans}{一.III.1.19}
    要成为等价关系，每个关系都必须是自反的、对称的和
    传递的。
    \begin{exparts}
      \item 这个关系
        不是对称的，因为如果 $x$ 修了 $4$~门课而 $y$
        修了 $3$ 门，那么 $x$ 与 $y$ 相关，但 $y$ 与 $x$
        不相关。
      \item 这是自反的，因为 $x$ 的名字与 $x$ 自己的名字
        以同一个字母开头。
        它是对称的，因为如果 $x$ 的名字与 $y$ 的名字
        以同一个字母开头，那么 $y$ 的名字也与~$x$ 的以同一个字母开头。
        它还是传递的，因为如果 $x$ 的名字与~$y$ 的以同一个字母开头，
        且 $y$ 的名字与 $z$ 的
        以同一个字母开头，那么 $x$ 的名字与 $z$ 的以同一个字母开头。
        所以它是等价关系。
    \end{exparts}
  
\end{ans}
\begin{ans}{一.III.1.20}
    对每一个我们都必须验证自反性、对称性和
    传递性这三个条件。
    \begin{exparts}
      \item 显然，任何矩阵沿对角线的乘积都与它自身相同，
        所以该关系是自反的。
        该关系是对称的，因为如果 $A$ 沿对角线的乘积
        与~$B$ 的相同，
        即 $a_{1,1}\cdot a_{2,2}=b_{1,1}\cdot b_{2,2}$，
        那么 $B$ 的乘积也与~$A$ 的相同。

        传递性类似：设 $A$ 的乘积为~$r$，
        且等于 $B$ 的乘积。
        又设 $B$ 的乘积等于 $C$ 的乘积。
        那么三者的乘积都是~$r$，即 $A$ 的乘积等于~$C$ 的乘积。

        每个实数对应一个等价类，也就是说，
        该等价类包含所有沿对角线
        乘积为这个实数的 $\nbyn{2}$ 矩阵。
      \item 自反性：如果矩阵 $A$ 有一个元素为~$1$，那么它与
        自身相关；如果没有，那么它同样与
        自身相关。
        对称性也分两种情形：设 $A$ 与~$B$ 相关。
        如果 $A$ 有一个元素为~$1$，那么~$B$ 也有，于是 $B$ 与~$A$ 相关。
        如果 $A$ 没有~$1$，那么 $B$ 也没有，同样 $B$ 与
        $A$ 相关。

        传递性：设 $A$ 与~$B$ 相关，$B$ 与~$C$ 相关。
        如果 $A$ 有一个元素为~$1$，那么~$B$ 也有；又因 $B$ 与~$C$ 相关，
        所以~$C$ 也有，
        从而 $A$ 与~$C$ 相关。
        同样，如果 $A$ 没有~$1$，那么~$B$ 也没有，因而
        ~$C$ 也没有，结论同样是 $A$ 与~$C$ 相关。

        恰有两个等价类：一个包含任何
        至少有一个元素为~$1$ 的 $\nbyn{2}$ 矩阵，
        另一个包含所有不含 $1$ 的矩阵。
    \end{exparts}
  
\end{ans}
\begin{ans}{一.III.1.21}
      \begin{exparts}
        \item 这个关系不是自反的。
          例如，任何左上角元素为~$1$ 的矩阵
          都与自身不相关。
        \item 这个关系不是传递的。
          对下面这三个矩阵，$A$ 与~$B$ 相关，$B$ 与~$C$ 相关，
          但 $A$ 与~$C$ 不相关。
          \begin{equation*}
            A=\begin{mat}
              0 &0 \\
              0 &0
            \end{mat},\quad
            B=\begin{mat}
              4 &0 \\
              0 &0
            \end{mat},\quad
            C=\begin{mat}
              8 &0 \\
              0 &0
            \end{mat},\quad
          \end{equation*}
      \end{exparts}
    
\end{ans}
\protect \subsection {一.III.2: 线性组合引理}
\begin{ans}{一.III.2.11}
      把每个矩阵都化成简化阶梯形，然后进行比较。
      \begin{exparts}
        \partsitem 第一个给出
          \begin{equation*}
            \grstep{-4\rho_1+\rho_2}
            \begin{mat}[r]
              1  &2  \\
              0  &0
            \end{mat}
          \end{equation*}
          而第二个给出
          \begin{equation*}
            \grstep{\rho_1\leftrightarrow\rho_2}
            \begin{mat}[r]
              1  &2  \\
              0  &1
            \end{mat}
            \grstep{-2\rho_2+\rho_1}
            \begin{mat}[r]
              1  &0  \\
              0  &1
            \end{mat}
          \end{equation*}
          两个简化阶梯形矩阵并不相同，因此
          原来的矩阵不行等价。
        \partsitem 第一个如下。
          \begin{equation*}
            \grstep[-5\rho_1+\rho_3]{-3\rho_1+\rho_2}
            \begin{mat}[r]
              1  &0  &2  \\
              0  &-1 &-5 \\
              0  &-1 &-5
            \end{mat}
            \grstep{-\rho_2+\rho_3}
            \begin{mat}[r]
              1  &0  &2  \\
              0  &-1 &-5 \\
              0  &0  &0
            \end{mat}
            \grstep{-\rho_2}
            \begin{mat}[r]
              1  &0  &2  \\
              0  &1  &5  \\
              0  &0  &0
            \end{mat}
          \end{equation*}
          第二个如下。
          \begin{equation*}
            \grstep{-2\rho_1+\rho_3}
            \begin{mat}[r]
              1  &0  &2  \\
              0  &2  &10 \\
              0  &0  &0
            \end{mat}
            \grstep{(1/2)\rho_2}
            \begin{mat}[r]
              1  &0  &2  \\
              0  &1  &5  \\
              0  &0  &0
            \end{mat}
          \end{equation*}
          这两个矩阵行等价。
        \partsitem 这两个矩阵不行等价，因为它们的
          尺寸不同。
        \partsitem 第一个，
          \begin{equation*}
            \grstep{\rho_1+\rho_2}
            \begin{mat}[r]
              1  &1  &1  \\
              0  &3  &3
            \end{mat}
            \grstep{(1/3)\rho_2}
            \begin{mat}[r]
              1  &1  &1  \\
              0  &1  &1
            \end{mat}
            \grstep{-\rho_2+\rho_1}
            \begin{mat}[r]
              1  &0  &0  \\
              0  &1  &1
            \end{mat}
          \end{equation*}
          而第二个，
          \begin{equation*}
            \grstep{\rho_1\leftrightarrow\rho_2}
            \begin{mat}[r]
              2  &2  &5  \\
              0  &3  &-1
            \end{mat}
            \grstep[(1/3)\rho_2]{(1/2)\rho_1}
            \begin{mat}[r]
              1  &1  &5/2 \\
              0  &1  &-1/3
            \end{mat}
            \grstep{-\rho_2+\rho_1}
            \begin{mat}[r]
              1  &0  &17/6 \\
              0  &1  &-1/3
            \end{mat}
          \end{equation*}
          这两个矩阵不行等价。
        \partsitem 这里第一个是
          \begin{equation*}
            \grstep{(1/3)\rho_2}
            \begin{mat}[r]
              1  &1  &1  \\
              0  &0  &1
            \end{mat}
            \grstep{-\rho_2+\rho_1}
            \begin{mat}[r]
              1  &1  &0  \\
              0  &0  &1
            \end{mat}
          \end{equation*}
          而这是第二个。
          \begin{equation*}
            \grstep{\rho_1\leftrightarrow\rho_2}
            \begin{mat}[r]
              1  &-1 &1  \\
              0  &1  &2
            \end{mat}
            \grstep{\rho_2+\rho_1}
            \begin{mat}[r]
              1  &0  &3  \\
              0  &1  &2
            \end{mat}
          \end{equation*}
          这两个矩阵不行等价。
       \end{exparts}
     
\end{ans}
\begin{ans}{一.III.2.12}
      对每一个作高斯–若尔当消元。
      两个矩阵行等价当且仅当它们具有相同的
      简化阶梯形。
      下面是每个矩阵的简化形式。
      % sage: M = matrix(QQ, [[1,3], [2,4]])
      % sage: gauss_jordan(M)
      % [1 3]
      % [2 4]
      %  take -2 times row 1 plus row 2
      % [ 1  3]
      % [ 0 -2]
      %  take -1/2 times row 2
      % [1 3]
      % [0 1]
      %  take -3 times row 2 plus row 1
      % [1 0]
      % [0 1]
      % sage: M = matrix(QQ, [[1,5], [2,10]])
      % sage: gauss_jordan(M)
      % [ 1  5]
      % [ 2 10]
      %  take -2 times row 1 plus row 2
      % [1 5]
      % [0 0]
      % sage: M = matrix(QQ, [[1,-1], [3,0]])
      % sage: gauss_jordan(M)
      % [ 1 -1]
      % [ 3  0]
      %  take -3 times row 1 plus row 2
      % [ 1 -1]
      % [ 0  3]
      %  take 1/3 times row 2
      % [ 1 -1]
      % [ 0  1]
      %  take 1 times row 2 plus row 1
      % [1 0]
      % [0 1]
      % sage: M = matrix(QQ, [[2,6], [4,10]])
      % sage: gauss_jordan(M)
      % [ 2  6]
      % [ 4 10]
      %  take -2 times row 1 plus row 2
      % [ 2  6]
      % [ 0 -2]
      %  take 1/2 times row 1
      %  take -1/2 times row 2
      % [1 3]
      % [0 1]
      %  take -3 times row 2 plus row 1
      % [1 0]
      % [0 1]
      % sage: M = matrix(QQ, [[0,1], [-1,0]])
      % sage: gauss_jordan(M)
      % [ 0  1]
      % [-1  0]
      %  swap row 1 with row 2
      % [-1  0]
      % [ 0  1]
      %  take -1 times row 1
      % [1 0]
      % [0 1]
      % sage: M = matrix(QQ, [[3,3], [2,2]])
      % sage: gauss_jordan(M)
      % [3 3]
      % [2 2]
      %  take -2/3 times row 1 plus row 2
      % [3 3]
      % [0 0]
      %  take 1/3 times row 1
      % [1 1]
      % [0 0]
      \begin{exparts*}
        \partsitem
          \(\begin{mat}
              1  &0  \\
              0  &1
            \end{mat} \)
        \partsitem
          \(\begin{mat}
              1  &5  \\
              0  &0
            \end{mat} \)
        \partsitem
          \(\begin{mat}
             1   &0  \\
             0   &1
            \end{mat} \)
        \partsitem
          \(\begin{mat}
             1   &0  \\
             0   &1
            \end{mat} \)
        \partsitem
          \(\begin{mat}
             1   &0  \\
             0   &1
            \end{mat} \)
        \partsitem
          \(\begin{mat}
             1   &1  \\
             0   &0
            \end{mat} \)
      \end{exparts*}
    
\end{ans}
\begin{ans}{一.III.2.13}
     对每个小题，只需对起始矩阵作一些行
     变换即可。
      \begin{exparts}
        \partsitem
          把第一行乘以~$3$，得到下式。
          \begin{equation*}
            \begin{mat}
              3 &9  \\
              4 &-1
            \end{mat}
          \end{equation*}
          （这个行变换的选择并没有什么讲究；它
          只是脑海中想到的第一个而已。）
          另外两个行变换是行对换 $\rho_1\leftrightarrow\rho_2$
          和倍加 $\rho_1+\rho_2$。
          \begin{equation*}
            \begin{mat}
              4 &-1  \\
              1 &3
            \end{mat}
            \quad
            \begin{mat}
              1 &3  \\
              5 &2
            \end{mat}
          \end{equation*}
        \partsitem 作同样的三个任意行变换，得到
          这三个。
          \begin{equation*}
           \begin{mat}
            0 &3 &6 \\
            1 &1 &1  \\
            2 &3 &4
           \end{mat}
           \quad
           \begin{mat}
            1 &1 &1  \\
            0 &1 &2 \\
            2 &3 &4
           \end{mat}
           \quad
           \begin{mat}
            0 &1 &2 \\
            1 &2 &3  \\
            2 &3 &4
           \end{mat}
           \quad
          \end{equation*}
      \end{exparts}
    
\end{ans}
\begin{ans}{一.III.2.14}
      高斯化简是常规计算。
      % sage: load("gauss_method.sage")
      % sage: M = matrix (QQ, [[1,2,1], [3,-1,0], [0,4,0]])
      % sage: gauss_method(M)
      % [ 1  2  1]
      % [ 3 -1  0]
      % [ 0  4  0]
      %  take -3 times row 1 plus row 2
      % [ 1  2  1]
      % [ 0 -7 -3]
      % [ 0  4  0]
      %  take 4/7 times row 2 plus row 3
      % [    1     2     1]
      % [    0    -7    -3]
      % [    0     0 -12/7]
      \begin{equation*}
        \begin{mat}
          1 &2   &1 \\
          3 &-1  &0 \\
          0 &4   &0
        \end{mat}
        \grstep{-3\rho_1+\rho_2}
        \begin{mat}
          1 &2   &1 \\
          0 &-7  &-3 \\
          0 &4   &0
        \end{mat}
        \grstep{(4/7)\rho_2+\rho_3}
        \begin{mat}
          1 &2   &1 \\
          0 &-7  &-3 \\
          0 &0   &-12/7
        \end{mat}
      \end{equation*}
      把这三个矩阵分别记为 $A$、$D$ 和~$B$，则有下式。
      \begin{align*}
        \begin{mat}
          \alpha_1 \\
          \alpha_2 \\
          \alpha_3
        \end{mat}
        &\grstep{-3\rho_1+\rho_2}
        \begin{mat}
          \delta_1=\alpha_1 \\
          \delta_2=-3\alpha_1+\alpha_2 \\
          \delta_3=\alpha_3
        \end{mat}                                   \\
        & \grstep{(4/7)\rho_2+\rho_3}
        \begin{mat}
          \beta_1=\alpha_1  \\
          \beta_2=-3\alpha_1+\alpha_2 \\
          \beta_3=-(12/7)\alpha_1+(4/7)\alpha_2+\alpha_3
        \end{mat}
      \end{align*}
    
\end{ans}
\begin{ans}{一.III.2.15}
       首先，与全
       \( 0 \) 矩阵行等价的矩阵只有它自己（因为行变换对它没有效果）。

       其次，能化简为
       \begin{equation*}
         \begin{mat}
           1  &a  \\
           0  &0
         \end{mat}
       \end{equation*}
       的矩阵形如
       \begin{equation*}
         \begin{mat}
           b  &ba \\
           c  &ca
         \end{mat}
       \end{equation*}
       （其中 \( a,b,c\in\Re \)，且 \(b\) 与 \(c\) 不全为零）。

       再次，能化简为
       \begin{equation*}
         \begin{mat}[r]
           0  &1  \\
           0  &0
         \end{mat}
       \end{equation*}
       的矩阵形如
       \begin{equation*}
         \begin{mat}
           0  &a \\
           0  &b
         \end{mat}
       \end{equation*}
       （其中 \( a,b\in\Re \)，且不全为零）。

       最后，能化简为
       \begin{equation*}
         \begin{mat}[r]
           1  &0  \\
           0  &1
         \end{mat}
       \end{equation*}
       的矩阵是非奇异的矩阵。
       这是因为以这个矩阵为系数矩阵的线性方程组将有唯一解，而这正是
       非奇异的定义。
       （换一种说法：它们不属于上面的任何一个等价类。）
     
\end{ans}
\begin{ans}{一.III.2.16}
      \begin{exparts}
        \partsitem 它们形如
          \begin{equation*}
            \begin{mat}
              a  &0  \\
              b  &0
            \end{mat}
          \end{equation*}
          其中 \( a,b\in\Re \) 至少有一个不为零。
        \partsitem 它们具有如下形式
          \begin{equation*}
            \begin{mat}
             a  &2a \\
             b  &2b
            \end{mat}
          \end{equation*}
          其中 \( a,b\in\Re \) 至少有一个不为零。
        \partsitem 所给矩阵是非奇异的。
          所以该行等价类由所有非奇异的 $\nbyn{2}$
          矩阵组成。
          （给出公式的话，它们形如
          \begin{equation*}
            \begin{mat}
              a  &b  \\
              c  &d
            \end{mat}
          \end{equation*}
          其中 \( a,b,c,d\in\Re \)）且 \( ad-bc\neq 0 \)。
          我们将在第四章看到，这个公式
          判定了 \( \nbyn{2} \) 矩阵
          何时非奇异。）
      \end{exparts}
    
\end{ans}
\begin{ans}{一.III.2.17}
       有无穷多个。
       例如，在
       \begin{equation*}
         \begin{mat}
           1  &k  \\
           0  &0
         \end{mat}
       \end{equation*}
       中，每个 $k\in\Re$ 给出一个不同的等价类。
    
\end{ans}
\begin{ans}{一.III.2.18}
      不能。
      行变换不改变矩阵的尺寸。
    
\end{ans}
\begin{ans}{一.III.2.19}
     \begin{exparts}
      \partsitem 对全零矩阵作行变换没有效果。
        因此每个这样的矩阵在它的行等价类中只有自己。
      \partsitem 没有。
        任何非零元都可以倍乘。
     \end{exparts}
    
\end{ans}
\begin{ans}{一.III.2.20}
      这里有两个。
      \begin{equation*}
        \begin{mat}[r]
          1  &1  &0  \\
          0  &0  &1
        \end{mat}
        \quad\text{和}\quad
        \begin{mat}[r]
          1  &0  &0  \\
          0  &0  &1
        \end{mat}
      \end{equation*}
     
\end{ans}
\begin{ans}{一.III.2.21}
      任意两个 \( \nbyn{n} \) 非奇异矩阵具有
      相同的简化阶梯形，即除对角线上
      为 \( 1 \) 之外全为 \( 0 \) 的那个矩阵。
      \begin{equation*}
        \begin{mat}
          1  &0  &       &0  \\
          0  &1  &       &0  \\
             &   &\ddots &   \\
          0  &0  &       &1
        \end{mat}
      \end{equation*}

      两个同尺寸的奇异矩阵未必行等价。
      例如，这两个 \( \nbyn{2} \) 奇异矩阵
      不行等价。
      \begin{equation*}
        \begin{mat}[r]
          1  &1  \\
          0  &0
        \end{mat}
        \quad\text{和}\quad
        \begin{mat}[r]
          1  &0  \\
          0  &0
        \end{mat}
      \end{equation*}
    
\end{ans}
\begin{ans}{一.III.2.22}
      由于每个等价类中有且仅有一个简化阶梯形矩阵，
      我们只需列出所有可能的简化阶梯形矩阵。

      这份清单参见 \nearbyexercise{exer:PossRedEchFrms} 的答案。
    
\end{ans}
\begin{ans}{一.III.2.23}
          \begin{exparts}
           \partsitem 如果存在 $c_0$ 不为零的线性关系
             那么我们可以从两边减去 $c_0\vec{\beta}_0$ 再除以 $-c_0$，
             把 $\vec{\beta}_0$ 表示成其余向量的线性
             组合。
             （注：
             如果集合中没有其他向量\Dash 即关系式
             是 $\zero=3\cdot\zero$\Dash 那么该命题仍然成立，因为
             按定义零向量是空向量集
             的和。）

             反之，如果 $\vec{\beta}_0$ 是其余向量的组合
             $\vec{\beta}_0=c_1\vec{\beta}_1+\dots+c_n\vec{\beta}_n$
             那么从两边减去
             $\vec{\beta}_0$ 就给出一个系数
             不全为零的线性关系；即
             $\vec{\beta}_0$ 前面的 $-1$。
           \partsitem 第一行不是其余各行的线性组合，
             理由
             如证明中所述：~在包含第一行首主元的那一列
             的分量等式中，唯一非零的元素
             是来自第一行的首主元，因此
             它的系数必为零。
             于是，由本题前面的小题，第一行与其
             余各行没有线性关系。

             因此，在考虑第二行能否与其余各行有线性
             关系时，
             我们可以不考虑第一行。
             但这时刚才对第一行的论证也将适用于
             第二行。
             （也就是说，我们这里是用归纳法论证的。）
         \end{exparts}
      
\end{ans}
\begin{ans}{一.III.2.24}
      我们知道 $4s+c+10d=8.45$ 与 $3s+c+7d=6.30$，而且想知道
      $s+c+d$ 是多少。
      幸运的是，$s+c+d$ 是 $4s+c+10d$ 与 $3s+c+7d$ 的线性组合。
      把未知的价格记为 $p$，我们有如下化简过程。
      \begin{align*}
        \begin{amat}{3}
          4  &1  &10  &8.45 \\
          3  &1  &7   &6.30 \\
          1  &1  &1   &p
        \end{amat}
        &\grstep[-(1/4)\rho_1+\rho_3]{-(3/4)\rho_1+\rho_2}
        \begin{amat}{3}
          4  &1    &10     &8.45      \\
          0  &1/4  &-1/2   &-0.037\,5 \\
          0  &3/4  &-3/2   &p-2.112\,5
        \end{amat}                                      \\
        &\grstep{-3\rho_2+\rho_3}
        \begin{amat}{3}
          4  &1    &10     &8.45      \\
          0  &1/4  &-1/2   &-0.037\,5 \\
          0  &0    &0      &p-2.00
        \end{amat}
     \end{align*}
     所付的价钱是 \$$2.00$。
   
\end{ans}
\begin{ans}{一.III.2.25}
     \begin{enumerate}
        \item 一个容易的答案是这个。
          \begin{equation*}
            0=3
          \end{equation*}
          想要一个不那么抬杠的答案，求解该方程组：
          \begin{equation*}
            \begin{amat}[r]{2}
              3  &4  &8  \\
              2  &1   &3
            \end{amat}
            \grstep{-(2/3)\rho_1+\rho_2}
            \begin{amat}[r]{2}
              3  &4    &8    \\
              0  &-5/3 &-7/3
            \end{amat}
          \end{equation*}
          得到 \( y=7/5 \) 与 \( x=4/5 \)。
          于是任何不被 \( (7/5,4/5) \) 满足的方程都可以，
          例如 \( 5x+5y=10 \)。
        \item 从无解的方程组可以导出每一个方程。
          例如，下面演示如何从
          \( 0=5 \) 导出 \( 3x+2y=4 \)。
          首先，
          \begin{equation*}
            0=5
            \grstep{(3/5)\rho_1}
            0=3
            \grstep{x\rho_1}
            0=3x
          \end{equation*}
          （\( x=0 \) 情形的有效性需要另外论证，但很显然）。
          类似地，\( 0=2y \)。
          对 \( 0=4 \) 亦然。
          而现在，\( 0+0=0 \) 便给出 \( 3x+2y=4 \)。
     \end{enumerate}
    
\end{ans}
\begin{ans}{一.III.2.26}
      定义：若它们的增广
      矩阵行等价，则称线性方程组等价。
      等价方程组有相同解集的证明很容易。
    
\end{ans}
\begin{ans}{一.III.2.27}
      \begin{exparts}
        \partsitem 三种可能的行对换很容易，
          三种可能的倍乘也是如此。
          六种可能的倍加之一是 \( k\rho_1+\rho_2 \)：
          \begin{equation*}
            \begin{mat}
              1           &2           &3  \\
              k\cdot 1+3  &k\cdot 2+0  &k\cdot 3+3  \\
              1           &4           &5
            \end{mat}
          \end{equation*}
          此时第一列与第二列之和仍等于第三列。
          其余五种倍加类似。
        \partsitem 显然的猜想是：行变换不改变
          列之间的线性关系。
        \partsitem 逐情形的
          证明按照第一小题给出的思路进行。
      \end{exparts}
   
\end{ans}
\protect \topic {计算机代数系统}
\begin{ans}{1}
      \begin{exparts}
        \partsitem
\Sage{} 求解如下。
\begin{lstlisting}
sage: var('h,c')
(h, c)
sage: statics = [40*h + 15*c == 100,
....:            25*c == 50 + 50*h]
sage: solve(statics, h,c)
[[h == 1, c == 4]]
\end{lstlisting}
其他计算机代数系统也有类似的命令。
下面这些 \Maple{} 命令
\begin{lstlisting}
> A:=array( [[40,15],
             [-50,25]] );
> u:=array([100,50]);
> linsolve(A,u);
\end{lstlisting}
           给出答案 \lstinline![1,4]!。
        \partsitem 这里有一个自由变量。
\Sage{} 给出如下结果。
\begin{lstlisting}
sage: var('h,i,j,k')
(h, i, j, k)
sage: chemistry = [7*h == 7*j,
....:              8*h + 1*i == 5*j + 2*k,
....:              1*i == 3*j,
....:              3*i == 6*j + 1*k]
sage: solve(chemistry, h,i,j,k)
[[h == 1/3*r1, i == r1, j == 1/3*r1, k == r1]]
\end{lstlisting}
类似地，下面这个 \Maple{} 会话
\begin{lstlisting}
> A:=array( [[7,0,-7,0],
             [8,1,-5,2],
             [0,1,-3,0],
             [0,3,-6,-1]] );
> u:=array([0,0,0,0]);
> linsolve(A,u);
\end{lstlisting}
         给出同样的回答（但参数为 $t_1$）。
      \end{exparts}
    
\end{ans}
\begin{ans}{2}
      \begin{exparts}
        \partsitem 答案为 \( x=2 \) 与 \( y=1/2 \)。
          用 \Sage{} 会话求解如下。
\begin{lstlisting}
sage: var('x,y')
(x, y)
sage: system = [2*x + 2*y == 5,
....:              x - 4*y == 0]
sage: solve(system, x,y)
[[x == 2, y == (1/2)]]
\end{lstlisting}
          下面这个 \Maple{} 会话
\begin{lstlisting}
> A:=array( [[2,2],
             [1,-4]] );
> u:=array([5,0]);
> linsolve(A,u);
\end{lstlisting}
          当然得到相同的答案。
        \partsitem 答案为 \( x=1/2 \) 与 \( y=3/2 \)。
\begin{lstlisting}
sage: var('x,y')
(x, y)
sage: system = [-1*x + y == 1,
....:              x + y == 2]
sage: solve(system, x,y)
[[x == (1/2), y == (3/2)]]
\end{lstlisting}
        \partsitem 该方程组有无穷多解。
           在第一小节中，取 $z$ 为参数，
           我们得到 $x=(43-7z)/4$ 与 $y=(13-z)/4$。
           \Sage{} 得到相同的结果。
\begin{lstlisting}
sage: var('x,y,z')
(x, y, z)
sage: system = [x - 3*y + z == 1,
....:           x + y + 2*z == 14]
sage: solve(system, x,y)
[[x == -7/4*z + 43/4, y == -1/4*z + 13/4]]
\end{lstlisting}
           \Maple{} 回答 $(-12+7t_1, t_1, 13-4t_1)$，
           它更愿意取 $y$ 作参数。
        \partsitem 该方程组无解。
           \Sage{} 给出空列表。
\begin{lstlisting}
sage: var('x,y')
(x, y)
sage: system = [-1*x - y == 1,
....:           -3*x - 3*y == 2]
sage: solve(system, x,y)
[]
\end{lstlisting}
           类似地，
           当把数组 $A$ 与向量 $u$ 交给 \Maple{}
           并要求它执行 \lstinline!linsolve(A,u)! 时，
           它完全不返回任何结果；也就是说，它以
           无解作答。
        \partsitem \Sage{} 求得
\begin{lstlisting}
sage: var('x,y,z')
(x, y, z)
sage: system = [       4*y + z == 20,
....:            2*x - 2*y + z == 0,
....:             x +        z == 5,
....:             x +    y - z == 10]
sage: solve(system, x,y,z)
[[x == 5, y == 5, z == 0]]
\end{lstlisting}
           即解为 \( (x,y,z)=(5,5,0) \)。
        \partsitem 解有无穷多个。
           \Sage{} 求解如下。
\begin{lstlisting}
sage: var('x,y,z,w')
(x, y, z, w)
sage: system = [ 2*x +        z + w == 5,
....:                   y -       w == -1,
....:            3*x -        z - w == 0,
....:            4*x +  y + 2*z + w == 9]
sage: solve(system, x,y,z,w)
[[x == 1, y == r2 - 1, z == -r2 + 3, w == r2]]
\end{lstlisting}
           \Maple{} 给出 $(1, -1+t_1, 3-t_1, t_1)$。
      \end{exparts}
    
\end{ans}
\begin{ans}{3}
      \begin{exparts}
        \partsitem 该方程组有无穷多解。
              在第二小节中我们把解集写成
              \begin{equation*}
              \set{\colvec{6 \\ 0}+\colvec{-2 \\ 1}y
                      \suchthat y\in\Re}
              \end{equation*}
              \Sage{} 求得
\begin{lstlisting}
sage: var('x,y')
(x, y)
sage: system = [ 3*x + 6*y == 18,
....:              x + 2*y == 6]
sage: solve(system, x,y)
[[x == -2*r3 + 6, y == r3]]
\end{lstlisting}
              而 \Maple{} 回答 $(6-2t_1,t_1)$。
        \partsitem 解集只有一个元素。
          \begin{equation*}
             \set{\colvec{0 \\ 1} }
          \end{equation*}
          \Sage{} 给出如下结果。
\begin{lstlisting}
sage: var('x,y')
(x, y)
sage: system = [ x + y == 1,
....:            x - y == -1]
sage: solve(system, x,y)
[[x == 0, y == 1]]
\end{lstlisting}
        \partsitem 该方程组的解集是无限的
          \begin{equation*}
            \set{\colvec{4 \\ -1 \\ 0}+\colvec{-1 \\ 1 \\ 1}x_3
                             \suchthat x_3\in\Re}
          \end{equation*}
          \Sage{} 给出如下结果
\begin{lstlisting}
sage: var('x1,x2,x3')
(x1, x2, x3)
sage: system = [   x1 +        x3 == 4,
....:              x1 - x2 + 2*x3 == 5,
....:            4*x1 - x2 + 5*x3 == 17]
sage: solve(system, x1,x2,x3)
[[x1 == -r4 + 4, x2 == r4 - 1, x3 == r4]]
\end{lstlisting}
          而 Maple 给出 $(t_1,-t_1+3,-t_1+4)$。
        \partsitem 存在唯一解
           \begin{equation*}
             \set{\colvec{1 \\ 1 \\ 1}}
           \end{equation*}
           \Sage{} 也求出了它。
\begin{lstlisting}
sage: var('a,b,c')
(a, b, c)
sage: system = [ 2*a +  b - c == 2,
....:            2*a +      c == 3,
....:              a - b      == 0]
sage: solve(system, a,b,c)
[[a == 1, b == 1, c == 1]]
\end{lstlisting}
        \partsitem 该方程组有无穷多解；在
           第二小节中我们用两个参数描述了
           该解集。
           \begin{equation*}
             \set{\colvec{5/3 \\ 2/3 \\ 0 \\ 0}
                  +\colvec{-1/3 \\ 2/3 \\ 1 \\ 0}z
                  +\colvec{-2/3 \\ 1/3 \\ 0 \\ 1}w
                  \suchthat z,w\in\Re}
           \end{equation*}
           \Sage{} 也给出同样的结果。
\begin{lstlisting}
sage: var('x,y,z,w')
(x, y, z, w)
sage: system = [   x + 2*y - z     == 3,
....:            2*x +   y     + w == 4,
....:              x -   y + z + w == 1]
sage: solve(system, x,y,z,w)
[[x == r6, y == -r5 - 2*r6 + 4, z == -2*r5 - 3*r6 + 5, w == r5]]
\end{lstlisting}
           \Maple{} 也是如此：$(3-2t_1+t_2,t_1,t_2,-2+3t_1-2t_2)$。
        \partsitem 解集为空。
\begin{lstlisting}
sage: var('x,y,z,w')
(x, y, z, w)
sage: system = [   x +      z + w == 4,
....:            2*x +  y     - w == 2,
....:            3*x +  y + z     == 7]
sage: solve(system, x,y,z,w)
[]
\end{lstlisting}
      \end{exparts}
    
\end{ans}
\begin{ans}{4}
      \Sage{} 求解如下。
\begin{lstlisting}
sage: var('x,y,a,b,c,d,p,q')
(x, y, a, b, c, d, p, q)
sage: system = [ a*x +  c*y == p,
....:            b*x +  d*y == q]
sage: solve(system, x,y)
[[x == -(d*p - c*q)/(b*c - a*d), y == (b*p - a*q)/(b*c - a*d)]]
\end{lstlisting}
       对于如下输入
\begin{lstlisting}
> A:=array( [[a,c],
             [b,d]] );
> u:=array([p,q]);
> linsolve(A,u);
\end{lstlisting}
      \Maple{} 给出如下回答。
      \begin{equation*}
        \bigl[-\frac{-d\,p+q\,c}{-b\,c+a\,d},
          \frac{-b\,p+a\,q}{-b\,c+a\,d}\bigr]
      \end{equation*}
    
\end{ans}
\protect \topic {投入产出分析}
\begin{ans}{1}
      这些答案来自\textit{Octave}。
      \begin{exparts}
        \partsitem \Sage{} 得到$s\approx 25\,952$与$a\approx 30\,312$。
\begin{lstlisting}
sage: var('s,a')
(s, a)
sage: system = [(20053/25448)*s -  (2664/30346)*a == 17789,
....:             -(48/25448)*s + (21316/30346)*a == 21243]
sage: solve(system, s,a)
[[s == (2772443022712/106830469), a == (6476439541923/213660938)]]
sage: n(2772443022712/106830469), n(6476439541923/213660938)
(25951.8005365305, 30311.7621898814)
\end{lstlisting}
        \partsitem \Sage{} 得到$s\approx 25\,857$与$a\approx 30\,596$。
\begin{lstlisting}
sage: system = [(20053/25448)*s -  (2664/30346)*a == 17689,
....:             -(48/25448)*s + (21316/30346)*a == 21443]
sage: solve(system, s,a)
[[s == (2762271457112/106830469), a == (6537219545323/213660938)]]
sage: n(2762271457112/106830469), n(6537219545323/213660938)
(25856.5883213711, 30596.2316112410)
\end{lstlisting}
        \partsitem \Sage{} 得到$s\approx 25\,984$与$a\approx 30\,597$。
\begin{lstlisting}
sage: system = [(20053/25448)*s -  (2664/30346)*a == 17789,
....:             -(48/25448)*s + (21316/30346)*a == 21443]
sage: solve(system, s,a)
[[s == (2775832696312/106830469), a == (6537292375723/213660938)]]
sage: n(2775832696312/106830469), n(6537292375723/213660938)
(25983.5300012771, 30596.5724802865)
\end{lstlisting}
      \end{exparts}
    
\end{ans}
\begin{ans}{2}
      下面的\textit{Sage}会话
\begin{lstlisting}
sage: var('a,s')
(a, s)
sage: eqns=[(20053/25448)*s - (2664/30346)*a == 17589,
....:       (-48/25448)*s + (21316/30346)*a == 21243]
sage: solve(eqns, s, a)
[[s == (2745320544312/106830469), a == (6476293881123/213660938)]]
sage: s_by_s = (5395/25448)*2745320544312/106830469
sage: s_by_s
582010544505/106830469
sage: n(s_by_s)
5447.98267716114
sage: s_by_a = (2664/30346)*6476293881123/213660938
sage: n(s_by_a)
2660.93449955743
sage: a_by_s = (48/25448)*2745320544312/106830469
sage: n(a_by_s)
48.4713936058822
sage: a_by_a = (9030/30346)*6476293881123/213660938
sage: n(a_by_a)
9019.60905818451
\end{lstlisting}
      给出此表。
      \begin{center}
        \begin{tabular}{r|rrrr}
         \multicolumn{1}{c}{\ }  %put the | in the right place
         &\multicolumn{1}{r}{\begin{tabular}[b]{@{}c@{}}
                                \textit{使用方} \\[-.65ex] \textit{钢铁业}
                             \end{tabular}}
         &\multicolumn{1}{r}{\begin{tabular}[b]{@{}c@{}}
                                \textit{使用方} \\[-.65ex] \textit{汽车业}
                             \end{tabular}}
         &\multicolumn{1}{r}{\begin{tabular}[b]{@{}c@{}}
                                \textit{使用方} \\[-.65ex] \textit{其他部门}
                              \end{tabular}}
         &\textit{总计}                                                \\
         \cline{2-5}
         \begin{tabular}{r} \textit{钢材} \\[-.65ex] \textit{的价值} \end{tabular}
           &$5\,448$  &$2\,661$  &$17\,589$     &$25\,698$                          \\
         \begin{tabular}{r} \textit{汽车} \\[-.65ex] \textit{的价值} \end{tabular}
           &$48$      &$9\,020$  &$21\,243$     &$30\,311$
        \end{tabular}
      \end{center}
      作为比较，这里是原来的表。
      \begin{center}
        \begin{tabular}{r|rrrr}
         \multicolumn{1}{c}{\ }  %put the | in the right place
         &\multicolumn{1}{r}{\begin{tabular}[b]{@{}c@{}}
                                \textit{使用方} \\[-.65ex] \textit{钢铁业}
                             \end{tabular}}
         &\multicolumn{1}{r}{\begin{tabular}[b]{@{}c@{}}
                                \textit{使用方} \\[-.65ex] \textit{汽车业}
                             \end{tabular}}
         &\multicolumn{1}{r}{\begin{tabular}[b]{@{}c@{}}
                                \textit{使用方} \\[-.65ex] \textit{其他部门}
                              \end{tabular}}
         &\textit{总计}                                                \\
         \cline{2-5}
         \begin{tabular}{r} \textit{钢材} \\[-.65ex] \textit{的价值} \end{tabular}
           &$5\,395$  &$2\,664$  &$17\,389$     &$25\,448$                          \\
         \begin{tabular}{r} \textit{汽车} \\[-.65ex] \textit{的价值} \end{tabular}
           &$48$      &$9\,030$  &$21\,268$     &$30\,346$
        \end{tabular}
      \end{center}
    
\end{ans}
\begin{ans}{3}
      \begin{exparts}
        \partsitem \Sage{} 给出$s=24,244$、$a=30,307$。
\begin{lstlisting}
 sage: var('s,a')
(s, a)
sage: system = [(20053/25448)*s -  0.0500*a == 17589,
....:             -(48/25448)*s + (21316/30346)*a == 21243]
sage: solve(system, s,a)
[[s == (12951731858499/534221147), a == (32381469405615/1068442294)]]
sage: n(12951731858499/534221147), n(32381469405615/1068442294)
(24244.1392131918, 30307.1767071166)
\end{lstlisting}
        \partsitem Sage 给出$s=24,267$、$a=30,673$。
\begin{lstlisting}
sage: system = [(20053/25448)*s -  0.0500*a == 17589,
....:             -(48/25448)*s + (21316/30346)*a == 21500]
sage: solve(system, s,a)
[[s == (12964136043940/534221147), a == (16386224431390/534221147)]]
sage: n(12964136043940/534221147), n(16386224431390/534221147)
(24267.3584090448, 30673.1107957993)
\end{lstlisting}
      \end{exparts}
    
\end{ans}
\begin{ans}{4}
      \begin{exparts}
        \partsitem
          方程如下。
^^I  \begin{equation*}
^^I    \begin{linsys}{2}
^^I       (11.79/18.69)s   &-   &(1.28/14.27)a &= &10.51 \\
^^I       -(0/18.69)s   &+   &(9.87/14.27)a &= &9.87
^^I    \end{linsys}
^^I  \end{equation*}
          \textit{Sage}给出如下结果（含一个初始检验）。
\begin{lstlisting}
sage: var('a,s')
(a, s)
sage: eqns=[(11.79/18.69)*s - (1.28/14.27)*a == 10.51,
....:       (-0)*s         + (9.87/14.27)*a == 9.87]
sage: solve(eqns,s,a)
[[s == (1869/100), a == (1427/100)]]
sage: n(1869/100)
18.6900000000000
sage: n(1427/100)
14.2700000000000
sage: eqns=[(11.79/18.69)*s - (1.28/14.27)*a == 11.561,
....:       (-0)*s         + (9.87/14.27)*a == 11.3505]
sage: solve(eqns,s,a)
[[s == (8119559/393000), a == (32821/2000)]]
sage: n(8119559/393000)
20.6604554707379
sage: n(32821/2000)
16.4105000000000
\end{lstlisting}
^^I\partsitem
          下面是1947年的比值，取三位小数。
^^I  \begin{center}
            \begin{tabular}{r|cc}
               1947   &\textit{钢铁业使用}  &\textit{汽车业使用}  \\ \hline
               \textit{钢材的使用}  &$(6.90/18.69)=0.369$ &$(1.28/14.27)=0.090$  \\
               \textit{汽车的使用}  &$(0/18.69)=0.00$ &$(4.40/14.27)=0.308$
            \end{tabular}
          \end{center}
          下面是1958年的比值。
          \begin{center}
            \begin{tabular}{r|cc}
               1958           &\textit{钢铁业使用}  &\textit{汽车业使用}  \\ \hline
               \textit{钢材的使用}  &$(5395/25448)=0.212$ &$(2664/30346)=0.088$  \\
               \textit{汽车的使用}  &$(48/25448)=0.002$ &$(9030/30346)=0.298$
            \end{tabular}
          \end{center}
^^I\partsitem
          1958年的外部需求为$17.589$与~$21.243$，
          单位是1958年的十亿美元。
          按1948年的美元计，它们为$(17.589/1.3)=13.53$与
          $(21.243/1.3)=16.34$。
          \textit{Sage}给出这些数字。
\begin{lstlisting}
sage: eqns=[(11.79/18.69)*s - (1.28/14.27)*a == 13.53,
....:       (-0)*s         + (9.87/14.27)*a == 16.34]
sage: solve(eqns,s,a)
[[s == (137466107/5541300), a == (1165859/49350)]]
sage: n(137466107/5541300)
24.8075554472777
sage: n(1165859/49350)
23.6242958459980
\end{lstlisting}
      乘以$1.3$换回1958年的美元，得到
      $32.25$与~$30.71$。
      \end{exparts}
    
\end{ans}
\protect \topic {计算的精度}
\begin{ans}{1}
      科学记数法便于表达两位的限制：
      $.25\times 10^{2}+.67\times 10^{0}=.25\times 10^{2}$。
      注意加上 $2/3$ 对总和没有影响。
    
\end{ans}
\begin{ans}{2}
      消元
      \begin{equation*}
        \grstep{-3\rho_1+\rho_2}
        \begin{linsys}{2}
          x  &+  &2y   &=  &3  \\
             &   &-8y  &=  &-7.992
        \end{linsys}
      \end{equation*}
      给出 \( (x,y)=(1.002,0.999) \)。
      所以对该方程组而言，常数的微小变化只产生解的微小变化。
    
\end{ans}
\begin{ans}{3}
      \begin{exparts}
        \partsitem 完全精确的解是 $x=10$ 与 $y=1$。
        \partsitem 四位数字的消元
          \begin{equation*}
            \grstep{-(.3454/.0003)\rho_1+\rho_2}
            \begin{amat}{2}
              .0003  &1.556  &1.569  \\
              0      &-1794   &-1805
            \end{amat}
          \end{equation*}
          给出结论~$y=1.006$，这还不错，
          以及 $x=12.21$。
          当然，这与正确答案相差百分之二十。
      \end{exparts}
    
\end{ans}
\begin{ans}{4}
      \begin{exparts}
        \partsitem 对于第一个，首先 $(2/3)-(1/3)$ 为
          $.666\,666\,67-.333\,333\,33=.333\,333\,34$，
          于是
          $(2/3)+((2/3)-(1/3))=.666\,666\,67+.333\,333\,34=1.000\,000\,0$。

          对于另一个，首先
          $((2/3)+(2/3))=.666\,666\,67+.666\,666\,67=1.333\,333\,3$，
          于是
          $((2/3)+(2/3))-(1/3)=1.333\,333\,3-.333\,333\,33=.999\,999\,97$。
        \partsitem 第一个方程是
          $.333\,333\,33\cdot x+1.000\,000\,0\cdot y=0$，
          而第二个方程是
          $.666\,666\,67\cdot x+2.000\,000\,0\cdot y=0$。
      \end{exparts}
    
\end{ans}
\begin{ans}{5}
      \begin{exparts}
        \partsitem 这个计算
          \begin{align*}
            &\grstep[-(1/3)\rho_1+\rho_3]{-(2/3)\rho_1+\rho_2}
            \begin{amat}{3}
              3  &2                   &1                   &6               \\
              0  &-(4/3)+2\varepsilon &-(2/3)+2\varepsilon &-2+4\varepsilon \\
              0  &-(2/3)+2\varepsilon &-(1/3)-\varepsilon  &-1+\varepsilon
            \end{amat}                                                    \\
            &\grstep{-(1/2)\rho_2+\rho_3}
            \begin{amat}{3}
              3  &2                   &1                   &6               \\
              0  &-(4/3)+2\varepsilon &-(2/3)+2\varepsilon &-2+4\varepsilon \\
              0  &\varepsilon         &-2\varepsilon       &-\varepsilon
            \end{amat}
          \end{align*}
          给出第三个方程 $y-2z=-1$。
          代入第二个方程给出
          $((-10/3)+6\varepsilon)\cdot z=(-10/3)+6\varepsilon$，
          所以 $z=1$，从而 $y=1$。
          有了这些，第一个方程给出 $x=1$。
        \partsitem
          与上面一样，科学记数法便于表达
          对位数的限制。

          保留两位数字的解
          为 $z=2.1$，$y=2.6$，$x=-.43$。
          \begin{multline*}
            \begin{amat}{3}
              .30\times 10^{1}  &.20\times 10^{1}  &.10\times 10^{1}
                 &.60\times 10^{1}        \\
              .10\times 10^{1}  &.20\times 10^{-3} &.20\times 10^{-3}
                 &.20\times 10^{1}        \\
              .30\times 10^{1}  &.20\times 10^{-3}  &-.10\times 10^{-3}
                 &.10\times 10^{1}
            \end{amat}                                           \\
            \begin{aligned}
            &\grstep[-(1/3)\rho_1+\rho_3]{-(2/3)\rho_1+\rho_2}
            \begin{amat}{3}
              .30\times 10^{1}  &.20\times 10^{1}  &.10\times 10^{1}
                 &.60\times 10^{1}        \\
              0                 &-.13\times 10^{1} &-.67\times 10^{0}
                 &-.20\times 10^{1}        \\
              0                 &-.67\times 10^{0}  &-.33\times 10^{0}
                 &-.10\times 10^{1}
            \end{amat}                                          \\
            &\grstep{-(.67/1.3)\rho_2+\rho_3}
            \begin{amat}{3}
              .30\times 10^{1}  &.20\times 10^{1}  &.10\times 10^{1}
                 &.60\times 10^{1}        \\
              0                 &-.13\times 10^{1} &-.67\times 10^{0}
                 &-.20\times 10^{1}        \\
              0                 &0                  &.15\times 10^{-2}
                 &.31\times 10^{-2}
            \end{amat}
            \end{aligned}
          \end{multline*}
      \end{exparts}
    
\end{ans}
\protect \topic {网络分析}
\begin{ans}{1}
      \begin{exparts}
        \partsitem 总电阻为$7$~欧姆。
          在$9$~伏特的电势下，流量将为$9/7$~安培。
          顺带一提，此时各电压降将为：~$3$~欧姆电阻器上$27/7$~伏特，
          而两个$2$~欧姆电阻器上各为$18/7$~伏特。
        \partsitem 处理这个网络的一种方法是注意到左侧的$2$~欧姆
          电阻器上有$9$~伏特的电压降
          （因而通过它的流量为$9/2$~安培），而
          右侧剩余部分的电压降也为
          $9$~伏特，于是我们可以像上一小题那样分析它。
          我们也可以使用线性方程组。
          \begin{center}
            \includegraphics{gr/mp/ch1.48}
          \end{center}
          使用图中的变量，我们得到线性方程组
          \begin{equation*}
            \begin{linsys}{4}
              i_0  &- &i_1  &- &i_2  &  &    &= &0  \\
                   &  &i_1  &+ &i_2  &- &i_3 &= &0  \\
                   &  &2i_1 &  &     &  &    &= &9  \\
                   &  &     &  &7i_2 &  &    &= &9
            \end{linsys}
          \end{equation*}
          它给出唯一解$i_0=81/14$、$i_1=9/2$、$i_2=9/7$
          与$i_3=81/14$。

          当然，第一段与第二段给出相同的答案。
          本质上，在第一段中我们用一种不如高斯消元法系统的方法
          求解了该线性方程组：先解出一些变量，再代入。
        \partsitem
          使用这些变量
          \begin{center}
            \includegraphics{gr/mp/ch1.49}
          \end{center}
          足以给出唯一解的一个线性方程组如下。
          \begin{equation*}
            \begin{linsys}{7}
              i_0  &- &i_1  &- &i_2  &  &    &  &    & &    & &    &= &0  \\
                   &  &     &  &i_2  &- &i_3 &- &i_4 & &    & &    &= &0  \\
                   &  &     &  &     &  &i_3 &+ &i_4 &-&i_5 & &    &= &0  \\
                   &  &i_1  &  &     &  &    &  &    &+&i_5 &-&i_6 &= &0  \\
                   &  &3i_1 &  &     &  &    &  &    & &    & &    &= &9  \\
                   &  &     &  &3i_2 &  &    &+ &2i_4&+&2i_5& &    &= &9  \\
                   &  &     &  &3i_2 &+ &9i_3&  &    &+&2i_5& &    &= &9
            \end{linsys}
          \end{equation*}
          （后三个方程来自包含
            $i_0$–$i_1$–$i_6$的回路、
            包含$i_0$–$i_2$–$i_4$–$i_5$–$i_6$的回路、
            以及包含$i_0$–$i_2$–$i_3$–$i_5$–$i_6$的回路。）
           Octave给出
            $i_0=4.35616$、$i_1=3.00000$、$i_2=1.35616$、
            $i_3=0.24658$、$i_4=1.10959$、$i_5=1.35616$、$i_6=4.35616$。
      \end{exparts}
    
\end{ans}
\begin{ans}{2}
      \begin{exparts}
        \partsitem
          使用先前分析中的变量，
          \begin{displaymath}
            \begin{linsys}{3}
              i_0&- &i_1    &-  &i_2   &=  &0 \\
             -i_0&+ &i_1    &+  &i_2   &=  &0  \\
                 &  &5i_1   &   &      &=  &20  \\
                 &  &       &   &8i_2  &=  &20  \\
                 &  &-5i_1  &+  &8i_2  &=  &0
            \end{linsys}
          \end{displaymath}
          此时各支路中流过的电流
          为$i_2=20/8=2.5$、$i_1=20/5=4$与$i_0=13/2=6.5$，单位均为安培。
          于是并联部分表现得像一个
          阻值为$20/(13/2)\approx 3.08$~欧姆的单一电阻器。
        \partsitem
          类似的分析给出
          $i_2=i_1=20/8=2.5$与$i_0=40/8=5$~ 安培。
          等效电阻为$20/5=4$~欧姆。
        \partsitem
          再一次类似的分析给出
          $i_2=20/r_2$、$i_1=20/r_1$、
          与$i_0=20(r_1+r_2)/(r_1r_2)$，单位均为安培。
          于是并联部分表现得像一个
          阻值为$20/i_0=r_1r_2/(r_1+r_2)$~欧姆的单一电阻器。
          （这个等式常被表述为：~等效
          电阻~$r$满足$1/r=(1/r_1)+(1/r_2)$。）
      \end{exparts}
    
\end{ans}
\begin{ans}{3}
     将 Kirchhoff 电流定律用于
     $r_1$、$r_2$与~$r_g$相汇的节点，以及
     $r_3$、$r_4$与~$r_g$相汇的节点，得到下列各式。
     \begin{equation*}
       \begin{linsys}{3}
         i_1 &- &i_2 &- &i_g &= &0  \\
         i_3 &- &i_4 &+ &i_g &= &0
       \end{linsys}
     \end{equation*}
     将 Kirchhoff 电压定律用于右上方的回路
     以及右下方的回路，得到下列各式。
     \begin{equation*}
       \begin{linsys}{3}
         i_3r_3 &- &i_gr_g &- &i_1r_1 &= &0  \\
         i_4r_4 &- &i_2r_2 &+ &i_gr_g &= &0
       \end{linsys}
     \end{equation*}
     假定$i_g$为零，即得$i_1=i_2$、$i_3=i_4$、
     $i_1r_1=i_3r_3$以及$i_2r_2=i_4r_4$。
     然后整理最后一个等式
     \begin{equation*}
       r_4=\frac{i_2r_2}{i_4}\cdot\frac{i_3r_3}{i_1r_1}
     \end{equation*}
     再约去各个$i$，便得到所要的结论。
   
\end{ans}
\begin{ans}{4}
      \begin{exparts}
        \partsitem
           一种改造是：~在任一路口，流入的流量等于
           流出的流量。
           在这种情形下它看起来确实合理，
           除非车辆长时间滞留在某个路口。
        \partsitem
           由于有$50$辆车经 Main 驶出而只有$25$~辆车驶入，
           故$i_1-25=i_2$。
           类似地，Pier 的进出平衡意味着$i_2=i_3$，
           而 North 给出$i_3+25=i_1$。
           我们得到如下方程组。
           \begin{equation*}
             \begin{linsys}{3}
               i_1  &-  &i_2  &   &     &=  &25  \\
                    &   &i_2  &-  &i_3  &=  &0   \\
              -i_1  &   &     &+  &i_3  &=  &-25
              \end{linsys}
           \end{equation*}
        \partsitem
           行变换$\rho_1+\rho_2$与$\rho_2+\rho_3$
           导出结论：有无穷多解。
           以$i_3$为参数，解集如下。
           \begin{equation*}
             \set{\colvec[c]{25+i_3  \\ i_3  \\i_3} \suchthat i_3\in\Re}
           \end{equation*}
           由于问题是以车辆数来陈述的，
           我们可以把$i_3$限制为自然数。
        \partsitem
           如果我们想象一个最初为空、具有给定输入/输出行为的
           环形交叉口，那么可以叠加$i_3$辆无限绕行的
           汽车，得到一个新的解。
        \partsitem
           一个合适的重述可能是：~驶入
           环形交叉口的车辆数必须等于驶出的车辆数。
           这一条的合理性就不那么清楚了。
           在五分钟的时间段内，我们可能发现
           驶入的车辆比驶出的多出五六辆，
           尽管题中的进出表
           确实满足这一性质。
           无论如何，它对获得唯一解没有帮助，
           因为要做到这一点，我们需要知道无限绕行的车辆数。
      \end{exparts}
    
\end{ans}
\begin{ans}{5}
      \begin{exparts}
        \partsitem 下面是每个未知街区的一个变量；每个已知
          街区的流量如图所示。
          \begin{center}
            \includegraphics{gr/mp/ch1.51}
          \end{center}
          我们应用 Kirchhoff 原理——流入 Willow 与 Shelburne
          交叉口的车流必须等于流出的——得到
          $i_1+25=i_2+125$。
          按从右到左、从上到下的顺序处理各交叉口，
          得到下列方程。
          \begin{equation*}
            \begin{linsys}{7}
              i_1 &- &i_2 &  &    &  &    &  &    &  &    &  &    &= &10  \\
             -i_1 &  &    &+ &i_3 &  &    &  &    &  &    &  &    &= &15   \\
                  &  &i_2 &  &    &+ &i_4 &- &i_5 &  &    &  &    &= &5   \\
                  &  &    &  &-i_3&- &i_4 &  &    &+ &i_6 &  &    &= &-50 \\
                  &  &    &  &    &  &    &  &i_5 &  &    &- &i_7 &= &-10 \\
                  &  &    &  &    &  &    &  &    &  &-i_6&+ &i_7 &= &30
            \end{linsys}
          \end{equation*}
          行变换$\rho_1+\rho_2$，接着$\rho_2+\rho_3$，
          再$\rho_3+\rho_4$与$\rho_4+\rho_5$，最后$\rho_5+\rho_6$，
          得到如下方程组。
          % runfile("gauss_method.sage")
          % M = matrix(QQ, [[1,-1,0,0,0,0,0], [-1,0,1,0,0,0,0], [0,1,0,1,-1,0,0], [0,0,-1,-1,0,1,0], [0,0,0,0,1,0,-1], [0,0,0,0,0,-1,1]])
          % v = vector(QQ, [10,15,5,-50,-10,30])
          % M_prime = M.augment(v, subdivide=True)
          % gauss_method(M_prime)
          \begin{equation*}
            \begin{linsys}{7}
              i_1 &- &i_2 &  &    &  &    &  &    &  &    &  &    &= &10  \\
                  &  &-i_2&+ &i_3 &  &    &  &    &  &    &  &    &= &25   \\
                  &  &    &  &i_3 &+ &i_4 &- &i_5 &  &    &  &    &= &30  \\
                  &  &    &  &    &  &    &  &-i_5&+ &i_6 &  &    &= &-20  \\
                  &  &    &  &    &  &    &  &    &  &i_6 &- &i_7 &= &-30 \\
                  &  &    &  &    &  &    &  &    &  &    &  &0   &= &0
            \end{linsys}
          \end{equation*}
          由于自由变量是$i_4$与$i_7$，我们取它们为
          参数。
          \begin{equation*}
          \begin{split}
            i_6  &=  i_7-30  \\
            i_5  &=  i_6+20=(i_7-30)+20=i_7-10 \\
            i_3  &=  -i_4+i_5+30=-i_4+(i_7-10)+30=-i_4+i_7+20 \\
            i_2  &=  i_3-25=(-i_4+i_7+20)-25=-i_4+i_7-5 \\
            i_1  &=  i_2+10=(-i_4+i_7-5)+10=-i_4+i_7+5
          \end{split}
          \end{equation*}
          显然$i_4$与$i_7$必须为正，而且
          第一个方程表明$i_7$至少为$30$。
          若从$i_7$出发，则$i_2$~方程表明
          $0\leq i_4\leq i_7-5$。
        \partsitem 我们不能取$i_7$为零，否则$i_6$将
          为负（这意味着汽车在单行道 Jay 上
          逆行）。
          不过，我们可以取$i_7$小至$30$，此时
          有许多合适的$i_4$。
          例如，取$i_4=0$即得到解
          \begin{equation*}
            (i_1,i_2,i_3,i_4,i_5,i_6,i_7)
            =
            (35,25,50,0,20,0,30)
          \end{equation*}
      \end{exparts}
    
\end{ans}
\protect \chapter {向量空间}
\protect \section {第 I 节：向量空间的定义}
\protect \subsection {二.I.1: 定义与例题}
\begin{ans}{二.I.1.17}
      \begin{exparts}
        \partsitem \( 0+0x+0x^2+0x^3 \)
        \partsitem \( \begin{mat}[r]
                   0  &0  &0  &0  \\
                   0  &0  &0  &0
                 \end{mat} \)
        \partsitem 常值函数 \( f(x)=0 \)
        \partsitem 常值函数 \( f(n)=0 \)
      \end{exparts}
    
\end{ans}
\begin{ans}{二.I.1.18}
      \begin{exparts*}
        \partsitem \( 3+2x-x^2 \)
        \partsitem \( \begin{mat}[r]
                   -1  &+1  \\
                    0  &-3
                 \end{mat} \)
        \partsitem \( -3e^x+2e^{-x} \)
      \end{exparts*}
    
\end{ans}
\begin{ans}{二.I.1.19}
      \begin{exparts}
        \partsitem
          三个元素是：
          $1+2x$、$2-1x$ 与 $x$。
          （当然，可以有许多不同的答案。）

          验证过程与 \nearbyexample{ex:RealVecSpaces} 完全类似。
          我们先处理
          \nearbydefinition{def:VecSpace} 中与加法有关的
          条件 $1$-$5$。
          关于对加法封闭，即条件~(1)，
          注意当 $a+bx,c+dx\in\polyspace_1$ 时，
          $(a+bx)+(c+dx)=(a+c)+(b+d)x$ 是系数为实数的一次多项式，
          因而是
          $\polyspace_1$ 的元素。
          条件~(2) 的验证如下：当 $a+bx,c+dx\in\polyspace_1$ 时，
          $(a+bx)+(c+dx)=(a+c)+(b+d)x$，而按另一顺序相加则得
          $(c+dx)+(a+bx)=(c+a)+(d+b)x$，由于这些都是实数，
          故 $a+c=c+a$
          且 $b+d=d+b$。
          条件~(3) 类似：设
          $a+bx,c+dx,e+fx\in\polyspace$，则
          $((a+bx)+(c+dx))+(e+fx)=(a+c+e)+(b+d+f)x$
          而
          $(a+bx)+((c+dx)+(e+fx))=(a+c+e)+(b+d+f)x$，两者相等
          （也就是说，实数加法满足结合律，故 $(a+c)+e=a+(c+e)$
          且 $(b+d)+f=b+(d+f)$）。
          对于条件~(4)，注意一次多项式
          $0+0x\in\polyspace_1$ 满足
          $(a+bx)+(0+0x)=a+bx$ 与
          $(0+0x)+(a+bx)=a+bx$。
          对于这一部分的最后一个条件，即条件~(5)，
          注意对任意 $a+bx\in\polyspace_1$，其加法逆元
          是 $-a-bx\in\polyspace_1$，因为
          $(a+bx)+(-a-bx)=(-a-bx)+(a+bx)=0+0x$。

          接下来我们还要检查与
          标量乘法有关的条件~(6)-(10)。
          对于~(6)，即空间对标量乘法封闭这一条件，
          设 $r$ 是实数，$a+bx$ 是
          $\polyspace_1$ 的元素，则 $r(a+bx)=(ra)+(rb)x$ 是
          $\polyspace_1$ 的元素，因为它是系数为
          实数的一次多项式。
          条件~(7) 成立，因为
          $(r+s)(a+bx)=r(a+bx)+s(a+bx)$ 可由实数乘法的分配律
          得到。
          条件~(8) 类似：
          $r((a+bx)+(c+dx))=r((a+c)+(b+d)x)=r(a+c)+r(b+d)x=(ra+rc)+(rb+rd)x
          =r(a+bx)+r(c+dx)$。
          对于~(9)，我们有
          $(rs)(a+bx)=(rsa)+(rsb)x=r(sa+sbx)=r(s(a+bx))$。
          最后，条件~(10) 是
          $1(a+bx)=(1a)+(1b)x=a+bx$。
        \partsitem
          将该集合记为 $P$。
          本练习的前一小题对系数没有限制，
          而这里我们只关注满足 $a_0-2a_1=0$ 的一次多项式，
          即常数项减去一次项系数的两倍为零的一次多项式。
          于是，$P$ 的三个典型元素是
          $2+1x$、$6+3x$ 与 $-4-2x$。

          对于条件~(1)，我们必须证明：
          若两个满足该限制的一次多项式相加，
          则所得仍是一次满足该限制的一次多项式：这里的论证是，
          若 $a+bx,c+dx\in P$，则
          $(a+bx)+(c+dx)=(a+c)+(b+d)x$ 是 $P$ 的元素，因为
          $(a+c)-2(b+d)=(a-2b)+(c-2d)=0+0=0$。
          条件~(2) 可以这样验证：
          当 $a+bx,c+dx\in\polyspace_1$ 时，
          $(a+bx)+(c+dx)=(a+c)+(b+d)x$，而按另一顺序相加则得
          $(c+dx)+(a+bx)=(c+a)+(d+b)x$，由于这些都是实数，
          故 $a+c=c+a$
          且 $b+d=d+b$。
          （也就是说，这一条件不受该限制的影响，
          其验证与本练习第一小题中的验证相同。）
          条件~(3) 同样不受这一附加限制的影响：
          设
          $a+bx,c+dx,e+fx\in\polyspace$，则
          $((a+bx)+(c+dx))+(e+fx)=(a+c+e)+(b+d+f)x$
          而
          $(a+bx)+((c+dx)+(e+fx))=(a+c+e)+(b+d+f)x$，两者相等。
          对于条件~(4)，注意一次多项式
          $0+0x\in P$ 满足该限制，因为
          其常数项减去其一次项系数的两倍
          为零，于是本小题中的
          验证即可适用：
          $0+0x\in\polyspace_1$ 满足
          $(a+bx)+(0+0x)=a+bx$ 与
          $(0+0x)+(a+bx)=a+bx$。
          为检查条件~(5)，
          注意对任意 $a+bx\in P$，其加法逆元
          是 $-a-bx$，因为
          它是 $P$ 的元素（由 $a+bx\in P$ 知
          $a-2b=0$，两边同乘 $-1$ 得 $-a+2b=0$），
          并且如第一小题那样，它起到加法逆元的作用：
          $(a+bx)+(-a-bx)=(-a-bx)+(a+bx)=0+0x$。

          我们还必须检查
          标量乘法的条件~(6)-(10)。
          对于~(6)，即空间对标量乘法封闭这一条件，
          设 $r$ 是实数且 $a+bx\in P$（于是 $a-2b=0$），
          则 $r(a+bx)=(ra)+(rb)x$ 是
          $P$ 的元素，因为它是系数为实数、
          且满足 $(ra)-2(rb)=r(a-2b)=0$ 的一次多项式。
          条件~(7) 成立的理由与本练习第一小题中
          相同，因为
          $(r+s)(a+bx)=r(a+bx)+s(a+bx)$ 可由实数乘法的分配律
          得到。
          条件~(8) 也与第一小题中相同：
          $r((a+bx)+(c+dx))=r((a+c)+(b+d)x)=r(a+c)+r(b+d)x=(ra+rc)+(rb+rd)x
          =r(a+bx)+r(c+dx)$。
          对~(9) 也是如此：
          $(rs)(a+bx)=(rsa)+(rsb)x=r(sa+sbx)=r(s(a+bx))$。
          最后，条件~(10) 也是如此：
          $1(a+bx)=(1a)+(1b)x=a+bx$。
       \end{exparts}
     
\end{ans}
\begin{ans}{二.I.1.20}
      可参照
      \nearbyexample{ex:RealVecSpaces} 作为指导。
      （\textit{评注。}
      由于许多条件都很容易检查，
      有时人们会觉得一定是漏掉了什么。
      请记住：容易做或常规，与不必做是两回事。）
      \begin{exparts}
        \partsitem
          以下是三个元素。
          \begin{equation*}
            \begin{mat}
              1  &2  \\
              3  &4
            \end{mat},\,
            \begin{mat}
              -1  &-2  \\
              -3  &-4
            \end{mat},\,
            \begin{mat}
              0  &0  \\
              0  &0
            \end{mat}
          \end{equation*}

          对于~(1)，$\nbyn{2}$ 实矩阵的和是 $\nbyn{2}$
          实矩阵。
          对于~(2)，我们考察两个矩阵的和
          \begin{equation*}
            \begin{mat}
              a  &b  \\
              c  &d
            \end{mat}+
            \begin{mat}
              e  &f  \\
              g  &h
            \end{mat}
            =
            \begin{mat}
              a+e  &b+f  \\
              c+g  &d+h
            \end{mat}
          \end{equation*}
          并利用实数加法的交换律
          \begin{equation*}
            =\begin{mat}
              e+a  &f+b  \\
              g+c  &h+d
            \end{mat}
            =
            \begin{mat}
              e  &f  \\
              g  &h
            \end{mat}
            +
            \begin{mat}
              a  &b  \\
              c  &d
            \end{mat}
          \end{equation*}
          来验证矩阵加法的交换性。
          条件~(3) 的验证，即矩阵加法的结合律，
          与前面的验证类似：
          \begin{equation*}
            \big(
              \begin{mat}
                a  &b  \\
                c  &d
              \end{mat}
              +\begin{mat}
                e  &f  \\
                g  &h
              \end{mat}
            \big)
            +
            \begin{mat}
              i  &j  \\
              k  &l
            \end{mat}
            =
            \begin{mat}
              (a+e)+i  &(b+f)+j  \\
              (c+g)+k  &(d+h)+l
            \end{mat}
          \end{equation*}
          而
          \begin{equation*}
            \begin{mat}
              a  &b  \\
              c  &d
            \end{mat}
            +
              \big(
              \begin{mat}
                e  &f  \\
                g  &h
              \end{mat}
              +
              \begin{mat}
                i  &j  \\
                k  &l
              \end{mat}
            \big)
            =
            \begin{mat}
              a+(e+i)  &b+(f+j)  \\
              c+(g+k)  &d+(h+l)
            \end{mat}
          \end{equation*}
          且由于实数加法满足结合律，两者逐项相同。
          对于~(4)，这个空间的零元素是全为零的 $\nbyn{2}$
          矩阵。
          条件~(5) 成立，因为对任意 $\nbyn{2}$ 矩阵~$A$，
          其加法逆元是各元素为 $A$ 相应元素的相反数的矩阵，
          即矩阵 $-1\cdot A$。

          条件~($6$)
          成立，因为 $\nbyn{2}$ 矩阵的标量倍
          仍是 $\nbyn{2}$ 矩阵。
          对于条件~(7)，我们有：
          \begin{multline*}
            (r+s)
            \begin{mat}
              a  &b  \\
              c  &d
            \end{mat}
            =
            \begin{mat}
              (r+s)a  &(r+s)b  \\
              (r+s)c  &(r+s)d
            \end{mat}               \\
            =
            \begin{mat}
              ra+sa  &rb+sb  \\
              rc+sc  &rd+sd
            \end{mat}
            =
            r
            \begin{mat}
              a  &b  \\
              c  &d
            \end{mat}
            +
            s
            \begin{mat}
              a  &b  \\
              c  &d
            \end{mat}
          \end{multline*}
          条件~(8) 的处理方式相同。
          \begin{multline*}
            r
            \big(
              \begin{mat}
                a  &b  \\
                c  &d
              \end{mat}
              +
              \begin{mat}
                e  &f  \\
                g  &h
              \end{mat}
            \big)
            =
            r
            \begin{mat}
              a+e  &b+f  \\
              c+g  &d+h
            \end{mat}
            =
            \begin{mat}
              ra+re  &rb+rf  \\
              rc+rg  &rd+rh
            \end{mat}              \\
            =
            r
            \begin{mat}
              a  &b  \\
              c  &d
            \end{mat}
            +
            r
            \begin{mat}
              e  &f  \\
              g  &h
            \end{mat}
            =
            r
            \big(
            \begin{mat}
              a  &b  \\
              c  &d
            \end{mat}
            +
            \begin{mat}
              e  &f  \\
              g  &h
            \end{mat}
            \big)
          \end{multline*}
          对于~(9)，我们有：
          \begin{equation*}
            (rs)
            \begin{mat}
              a  &b  \\
              c  &d
            \end{mat}
            =
            \begin{mat}
              rsa  &rsb  \\
              rsc  &rsd
            \end{mat}
            =
            r
            \begin{mat}
              sa  &sb  \\
              sc  &sd
            \end{mat}
            =
            r\big( s
            \begin{mat}
              a  &b  \\
              c  &d
            \end{mat}
            \big)
          \end{equation*}
          条件~(10) 同样简单。
          \begin{equation*}
            1
            \begin{mat}
              a  &b  \\
              c  &d
            \end{mat}
            =
            \begin{mat}
              1\cdot a  &1\cdot b  \\
              1\cdot c  &1\cdot d
            \end{mat}
            =
            \begin{mat}
              sa  &sb  \\
              sc  &sd
            \end{mat}
          \end{equation*}
        \partsitem
          本小题与本练习前一小题的差别仅在于：
          我们限定于第二行第一列
          为零的矩阵集合 $T$。
          以下是 $T$ 的三个元素。
          \begin{equation*}
            \begin{mat}
              1  &2  \\
              0  &4
            \end{mat},\,
            \begin{mat}
              -1  &-2  \\
              0  &-4
            \end{mat},\,
            \begin{mat}
              0  &0  \\
              0  &0
            \end{mat}
          \end{equation*}
          本小题的某些验证与本练习第一小题相同，
          下面我们只做那些不同的验证。

          对于~(1)，$2,1$ 元为零的 $\nbyn{2}$ 实矩阵之和
          也是 $2,1$ 元为零的 $\nbyn{2}$ 实矩阵。
          \begin{equation*}
            \begin{mat}
              a  &b  \\
              0  &d
            \end{mat}
            +
            \begin{mat}
              e  &f  \\
              0  &h
            \end{mat}
            \begin{mat}
              a+e  &b+f  \\
              0    &d+h
            \end{mat}
          \end{equation*}
          前一小题中给出的条件~(2) 的验证
          在本小题同样适用。
          条件~(3) 也是如此。
          对于~(4)，注意全为零的 $\nbyn{2}$
          矩阵是 $T$ 的元素。
          条件~(5) 成立，因为对任意 $\nbyn{2}$ 矩阵~$A$，
          其加法逆元是矩阵
          $-1\cdot A$，因此 $2,1$ 元为零的矩阵的加法逆元
          也是 $2,1$ 元为零的矩阵。

          条件~$6$ 成立，因为 $2,1$ 元为零的 $\nbyn{2}$ 矩阵的标量倍
          仍是 $2,1$ 元为零的 $\nbyn{2}$ 矩阵。
          条件~(7) 的验证与前一小题相同。
          条件~(8)、(9) 与~(10) 的验证也是如此。
      \end{exparts}
    
\end{ans}
\begin{ans}{二.I.1.21}
      % Most of the conditions are easy to check; use
      % \nearbyexample{ex:RealVecSpaces} as a guide.
      \begin{exparts}
        \partsitem
          三个元素是 $\rowvec{1  &2  &3}$、
          $\rowvec{2  &1  &3}$ 与 $\rowvec{0  &0  &0}$。

          我们必须检查 \nearbydefinition{def:VecSpace} 中的条件 (1)-(10)。
          条件 (1)-(5) 与加法有关。
          对于条件~(1)，回顾两个三分量行向量的和
          \begin{equation*}
            \rowvec{a  &b  &c}
            +\rowvec{d  &e  &f}
            =\rowvec{a+d  &b+e  &c+f}
          \end{equation*}
          也是三分量行向量
         （字母 $a,\ldots,f$ 全都表示实数）。
         对~(2) 的验证是常规的
          \begin{multline*}
            \rowvec{a  &b  &c}
            +\rowvec{d  &e  &f}
            =\rowvec{a+d  &b+e  &c+f}   \\
            =\rowvec{d+a  &e+b  &f+c}
            =\rowvec{d  &e  &f}
            +\rowvec{a  &b  &c}
          \end{multline*}
          （第二个等号成立是因为三个分量都是实数，
          而实数加法可交换）。
         条件~(3) 的验证类似。
          \begin{multline*}
            \big(\rowvec{a  &b  &c}
              +\rowvec{d  &e  &f}
            \big)
            +\rowvec{g  &h  &i}
            =\rowvec{(a+d)+g  &(b+e)+h  &(c+f)+i}    \\
            =\rowvec{a+(d+g)  &b+(e+h)  &c+(f+i)}
            =\rowvec{a  &b  &c}
            +\big(\rowvec{d  &e  &f}
               +\rowvec{g  &h  &i}
             \big)
          \end{multline*}
          对于~(4)，注意三分量行向量
          $\rowvec{0  &0  &0}$ 是加法单位元：
          $\rowvec{a  &b  &c}+\rowvec{0  &0  &0}=\rowvec{a  &b  &c}$。
          为验证条件~(5)，设给定集合中的
          元素 $\rowvec{a  &b  &c}$，并注意
          $\rowvec{-a  &-b  &-c}$ 也在此集合中且具有
          所需性质：
          $\rowvec{a  &b  &c}+\rowvec{-a  &-b  &-c}=\rowvec{0  &0  &0}$。

          条件 (6)-(10) 涉及标量乘法。
          为验证~(6)，即空间对所给的
          标量乘法运算封闭，
          注意 $r\rowvec{a  &b  &c}=\rowvec{ra  &rb  &rc}$ 是
          实分量的三分量行向量。
          对于~(7)，我们计算如下。
          \begin{multline*}
            (r+s)\rowvec{a  &b  &c}=\rowvec{(r+s)a  &(r+s)b  &(r+s)c}
            =\rowvec{ra+sa  &rb+sb  &rc+sc}                             \\
            =\rowvec{ra  &rb  &rc}+\rowvec{sa  &sb  &sc}
            =r\rowvec{a  &b  &c}+s\rowvec{a  &b  &c}
          \end{multline*}
          条件~(8) 非常类似。
          \begin{multline*}
            r\big(\rowvec{a  &b  &c}+\rowvec{d  &e  &f}\big)  \\
            \begin{aligned}
            &=r\rowvec{a+d  &b+e  &c+f}
            =\rowvec{r(a+d)  &r(b+e)  &r(c+f)}                     \\
            &=\rowvec{ra+rd  &rb+re  &rc+rf}
            =\rowvec{ra  &rb  &rc}+\rowvec{rd  &re  &rf}          \\
            &=r\rowvec{a  &b  &c}+r\rowvec{d  &e  &f}
            \end{aligned}
          \end{multline*}
          条件~(9) 的计算也是如此。
          \begin{equation*}
            (rs)\rowvec{a  &b  &c}
             =\rowvec{rsa  &rsb  &rsc}
             =r\rowvec{sa  &sb  &sc}
             =r\big(s\rowvec{a  &b  &c}\big)
          \end{equation*}
          条件~(10) 同样是常规的：
          $1\rowvec{a  &b  &c}
             =\rowvec{1\cdot a  &1\cdot b  &1\cdot c}
             =\rowvec{a  &b  &c}$。
        \partsitem
          将该集合记为 $L$。
          加法的封闭性，即条件~(1)，需要检查：若各加数都是~$L$ 的成员，则
          和
          \begin{equation*}
            \colvec{a \\ b \\ c \\ d}
            +\colvec{e \\ f \\ g \\ h}
            =
            \colvec{a+e \\ b+f \\ c+g \\ d+h}
          \end{equation*}
          也是 \( L \) 的成员；这是成立的，因为它满足
          $L$ 的成员资格判据：
          \( (a+e)+(b+f)-(c+g)+(d+h)
          =(a+b-c+d)+(e+f-g+h)=0+0 \)。
          条件~(2)、(3) 与~(5) 的验证与本练习第一部分中的
          验证类似。
          对于条件~(4)，注意零向量是~$L$ 的成员，
          因为其第一分量加第二分量、减第三分量、
          再加第四分量，总和为零。

          条件~(6)，即标量乘法的封闭性，类似：
          当该向量是~$L$ 的元素时，
          \begin{equation*}
            r\colvec{a  \\ b \\ c \\ d}
            =\colvec{ra  \\  rb  \\  rc  \\  rd}
          \end{equation*}
          也是~$L$ 的元素，因为 $ra+rb-rc+rd=r(a+b-c+d)=r\cdot 0=0$。
          条件~(7)、(8)、(9) 与~(10) 的验证与本练习
          前一小题中相同。
     \end{exparts}
     
\end{ans}
\begin{ans}{二.I.1.22}
      每一小题中都将该集合记为 \( Q \)。
      对其中某些小题，还有其他正确方法可证明 $Q$ 不是
      向量空间。
      \begin{exparts}
        \partsitem 它对加法不封闭；它不满足
          条件~(1)。
          \begin{equation*}
            \colvec[r]{1 \\ 0 \\ 0},
            \colvec[r]{0 \\ 1 \\ 0}\in Q
            \qquad
            \colvec[r]{1 \\ 1 \\ 0}\not\in Q
          \end{equation*}
        \partsitem 它对加法不封闭。
          \begin{equation*}
            \colvec[r]{1 \\ 0 \\ 0},
            \colvec[r]{0 \\ 1 \\ 0}\in Q
            \qquad
            \colvec[r]{1 \\ 1 \\ 0}\not\in Q
          \end{equation*}
        \partsitem 它对加法不封闭。
          \begin{equation*}
            \begin{mat}[r]
              0  &1  \\
              0  &0
            \end{mat},
            \,
            \begin{mat}[r]
              1  &1  \\
              0  &0
            \end{mat}\in Q
            \qquad
            \begin{mat}[r]
              1  &2  \\
              0  &0
            \end{mat}\not\in Q
          \end{equation*}
        \partsitem 它对标量乘法不封闭。
          \begin{equation*}
            1+1x+1x^2\in Q
            \qquad
            -1\cdot(1+1x+1x^2)\not\in Q
          \end{equation*}
        \item 它是空集，违反条件~(4)。
      \end{exparts}
    
\end{ans}
\begin{ans}{二.I.1.23}
      通常的运算
      \( (v_0+v_1i)+(w_0+w_1i)=(v_0+w_0)+(v_1+w_1)i \) 与
      \( r(v_0+v_1i)=(rv_0)+(rv_1)i \) 即可满足要求。
      检查很容易。
    
\end{ans}
\begin{ans}{二.I.1.24}
       不是，它对标量乘法不封闭，因为例如
       \( \pi\cdot (1) \) 不是有理数。
    
\end{ans}
\begin{ans}{二.I.1.25}
      自然运算是
      \( (v_1x+v_2y+v_3z)+(w_1x+w_2y+w_3z)=(v_1+w_1)x+(v_2+w_2)y+(v_3+w_3)z \)
      与 \( r\cdot(v_1x+v_2y+v_3z)=(rv_1)x+(rv_2)y+(rv_3)z \)。
      检查它是向量空间很容易；以
      \nearbyexample{ex:RealVecSpaces} 作为指导。
    
\end{ans}
\begin{ans}{二.I.1.26}
      “\( + \)” 运算不可交换（即条件~(2) 不
      满足）；很容易找出该集合的两个成员来
      证实这一断言。
    
\end{ans}
\begin{ans}{二.I.1.27}
      \begin{exparts}
        \partsitem 它不是向量空间。
          \begin{equation*}
            (1+1)\cdot\colvec[r]{1 \\ 0 \\ 0}\neq
            \colvec[r]{1 \\ 0 \\ 0}
            +\colvec[r]{1 \\ 0 \\ 0}
          \end{equation*}
        \partsitem 它不是向量空间。
          \begin{equation*}
            1\cdot\colvec[r]{1 \\ 0 \\ 0}\neq\colvec[r]{1 \\ 0 \\ 0}
          \end{equation*}
      \end{exparts}
    
\end{ans}
\begin{ans}{二.I.1.28}
      对每个“是”的答案，只需在该空间中检查几个
      线性组合即可。
      对每个“否”的答案，给出其中一个
      条件不成立的具体例子。
      \begin{exparts}
        \partsitem 是。
          随手取的一个线性组合如下。
          \begin{equation*}
            2\cdot\begin{mat}
              3  &0 \\
              0  &4
            \end{mat}
            +5\cdot\begin{mat}
              6  &0 \\
              0  &7
            \end{mat}
          \end{equation*}
          这个集合成为向量空间的关键在于：
          结果是对角矩阵。
          \begin{equation*}
            =\begin{mat}
              36 &0 \\
              0  &43
            \end{mat}
          \end{equation*}
        \partsitem 是。
          线性组合如下。
          \begin{equation*}
            10\cdot\begin{mat}
              1  &3  \\
              3  &2
            \end{mat}
            +11\cdot\begin{mat}
              3  &7  \\
              7  &4
            \end{mat}
          \end{equation*}
          我们看到结果保持了如下形式：左上
          与右下元素相加得到右上元素，
          同时也得到左下元素。
          \begin{equation*}
            =\begin{mat}
               43  &107  \\
              107  &64
            \end{mat}
          \end{equation*}
        \partsitem 否，这个集合对自然加法
          运算不封闭。
          全部分量都是 $1/4$ 的向量是这个集合的成员，
          但它自加所得的结果，即
          全部分量都是 $1/2$ 的向量，不是成员。
        \partsitem 是。
        \partsitem 否，\( f(x)=e^{-2x}+(1/2) \) 在集合中，
           但 \( 2\cdot f \) 不在（即条件~(6) 不成立）。
      \end{exparts}
    
\end{ans}
\begin{ans}{二.I.1.29}
      它是向量空间。
      向量空间定义中的大多数条件的验证都是常规的；这里我们
      只检查封闭性。
      对于加法，
      \( (f_1+f_2)\,(7)=f_1(7)+f_2(7)=0+0=0 \)。
      对于标量乘法，
      \( (r\cdot f)\,(7)=rf(7)=r0=0 \)。
    
\end{ans}
\begin{ans}{二.I.1.30}
      我们逐条检查 \nearbydefinition{def:VecSpace}。

      首先，对`\( + \)'封闭
      成立，因为两个正实数之积是
      正实数。
      第二个条件满足，因为实数乘法可交换。
      类似地，由于实数乘法可结合，第三条也成立。
      对于第四个条件，注意用
      \( 1\in\Re^+ \)中的数去乘一个数不会改变它。
      第五，任何正实数都有一个倒数，它是正实数。

      第六条，对`\( \cdot \)'封闭，
      成立是因为正实数的任意次幂仍是
      正实数。
      第七个条件就是\( v^{r+s} \)等于
      \( v^r \)与\( v^s \)之积这条法则。
      第八个条件说\( (vw)^r=v^rw^r \)。
      第九个条件断言\( (v^r)^s=v^{rs} \)。
      最后一个条件说\( v^1=v \)。
    
\end{ans}
\begin{ans}{二.I.1.31}
        \begin{exparts}
           \partsitem 不是：\( 1\cdot(0,1)+1\cdot(0,1)\neq (1+1)\cdot(0,1) \)。
           \partsitem 不是；与前一答案相同的计算表明，向量空间定义中的
              一个条件被违反了。
              违反向空间条件的另一个例子是
              \( 1\cdot (0,1)\neq (0,1) \)。
        \end{exparts}
      
\end{ans}
\begin{ans}{二.I.1.32}
      它不是向量空间，因为它对加法不封闭，例如
      \( (x^2)+(1+x-x^2) \)不在此集合中。
    
\end{ans}
\begin{ans}{二.I.1.33}
      \begin{exparts}
        \partsitem \( 6 \)
        \partsitem \( nm \)
        \partsitem \( 3 \)
        \partsitem 为看出答案是\( 2 \)，把它改写为
        \begin{equation*}
          \set{\begin{mat}
            a  &0  \\
            b  &-a-b
          \end{mat} \suchthat a,b\in\Re}
        \end{equation*}
        可见其中有两个参数。
      \end{exparts}
     
\end{ans}
\begin{ans}{二.I.1.34}
      {
      \def\plus{\mathbin{\vec{+}}}
      \def\tim{\mathbin{\vec{\cdot}}}
        \definend{向量空间}\index{vector space!definition}
        （\( \Re \)上的）由一个集合\( V \)连同两个运算`\( \plus \)'与
        `\( \tim \)'组成，满足以下条件。
        当\( \vec{v},\vec{w}\in V \)时，
        (1)~它们的\definend{向量和}
          \( \vec{v}\plus\vec{w} \)是\( V \)中的一个元素。
        若\( \vec{u},\vec{v},\vec{w}\in V \)，则
        (2)~\( \vec{v}\plus\vec{w}=\vec{w}\plus\vec{v} \)且
        (3)~\( (\vec{v}\plus\vec{w})\plus\vec{u}
                 =\vec{v}\plus(\vec{w}\plus\vec{u}) \)。
        (4)~存在一个\definend{零向量}
          \( \zero\in V \)，使得
          对所有\( \vec{v}\in V\)都有\( \vec{v}\plus\zero=\vec{v}\, \)。
        (5)~每个\( \vec{v}\in V \)都有一个
          \definend{加法逆元}
          \( \vec{w}\in V \)，使得\( \vec{w}\plus\vec{v}=\zero \)。
        若\( r,s \)是\definend{标量\/}
        （即\( \Re \)的成员），
        且\( \vec{v},\vec{w}\in V \)，则
        (6)~每个
           \definend{标量倍数}
           \( r\cdot\vec{v} \)都在\( V \)中。
        若\( r,s\in\Re \)且\( \vec{v},\vec{w}\in V \)，则
        (7)~\( (r+ s)\cdot\vec{v}=r\cdot\vec{v}\plus s\cdot\vec{v} \)，
        以及(8)~\( r\tim (\vec{v}+\vec{w})
           =r\tim\vec{v}+r\tim\vec{w} \)，
        以及(9)~\( (rs)\tim \vec{v} =r\tim (s\tim\vec{v}) \)，
        以及(10)~\( 1\tim \vec{v}=\vec{v} \)。
     }
    
\end{ans}
\begin{ans}{二.I.1.35}
      \begin{exparts}
        \partsitem 设\( V \)是向量空间，
          设\( \vec{v}\in V \)，并
          假设\( \vec{w}\in V \)是$\vec{v}$的加法逆元，
          于是\( \vec{w}+\vec{v}=\zero \)。
          因为加法可交换，
          \( \zero=\vec{w}+\vec{v}=\vec{v}+\vec{w} \)，
          所以\( \vec{v} \)也是
          \( \vec{w} \)的加法逆元。
        \partsitem 设\( V \)是向量空间且设
          \( \vec{v},\vec{s},\vec{t}\in V \)。
          \( \vec{v} \)的加法逆元是\( -\vec{v} \)，于是
          \( \vec{v}+\vec{s}=\vec{v}+\vec{t} \)给出
          \( -\vec{v}+\vec{v}+\vec{s}=-\vec{v}+\vec{v}+\vec{t} \)，
          即\( \zero+\vec{s}=\zero+\vec{t} \)，从而
          \( \vec{s}=\vec{t} \)。
      \end{exparts}
     
\end{ans}
\begin{ans}{二.I.1.36}
      加法可交换，所以在任何向量空间中，
      对任意向量\( \vec{v} \)我们有
      \( \vec{v}=\vec{v}+\zero=\zero+\vec{v} \)。
    
\end{ans}
\begin{ans}{二.I.1.37}
      它不是向量空间，因为尺寸不同的两个矩阵相加没有定义，
      因而该集合不满足封闭性
      条件。
    
\end{ans}
\begin{ans}{二.I.1.38}
      向量空间中每个元素有且只有一个加法
      逆元。

      设\( V \)是向量空间且设\( \vec{v}\in V \)。
      若\( \vec{w}_1,\vec{w}_2\in V \)都是
      \( \vec{v} \)的加法逆元，则考虑\( \vec{w}_1+\vec{v}+\vec{w}_2 \)。
      一方面，我们有它等于$\vec{w}_1+(\vec{v}+\vec{w}_2)=
      \vec{w}_1+\zero=\vec{w}_1$。
      另一方面，我们有它等于$(\vec{w}_1+\vec{v})+\vec{w}_2=
      \zero+\vec{w}_2=\vec{w}_2$。
      因此$\vec{w}_1=\vec{w}_2$。
    
\end{ans}
\begin{ans}{二.I.1.39}
     \begin{exparts}
      \partsitem 每个这样的集合都具有形式
        \( \set{r\cdot\vec{v}+s\cdot\vec{w}\suchthat r,s\in\Re} \)，
        其中\( \vec{v},\vec{w} \)之一或两者可以是\( \zero \)。
        在继承的运算下，加法的封闭性
        \( (r_1\vec{v}+s_1\vec{w})+(r_2\vec{v}+s_2\vec{w})
           =(r_1+r_2)\vec{v}+(s_1+s_2)\vec{w} \)
        与标量乘法的封闭性
        \( c(r\vec{v}+s\vec{w})=(cr)\vec{v}+(cs)\vec{w} \)
        是容易的。
        其他条件也是常规的。
      \partsitem 这样的集合在继承的
        运算下不能是向量空间，因为它没有零元素。
     \end{exparts}
    
\end{ans}
\begin{ans}{二.I.1.40}
      设\( \vec{v}\in V \)不是\( \zero \)。
      \begin{exparts}
        \partsitem 当且仅当的一个方向是清楚的：~若$r=0$，
          则$r\cdot\vec{v}=\zero$。
          另一个方向，设\( r \)是非零标量。
          若\( r\vec{v}=\zero \)，则
          \( (1/r)\cdot r\vec{v}=(1/r)\cdot \zero \)表明
          $\vec{v}=\zero$，与假设矛盾。
        \partsitem 当\( r_1,r_2 \)是标量时，
          \( r_1\vec{v}=r_2\vec{v}\, \)
          成立当且仅当\( (r_1-r_2)\vec{v}=\zero \)。
          由前一项，于是\( r_1-r_2=0 \)。
        \partsitem 非平凡空间有向量
          \( \vec{v}\neq\zero \)。
          考虑集合\( \set{k\cdot\vec{v}\suchthat k\in\Re} \)。
          由前一项，这个集合是无限的。
        \partsitem 解集要么是平凡的，要么是非平凡的。
          在第二种情形，它是无限的。
     \end{exparts}
    
\end{ans}
\begin{ans}{二.I.1.41}
      是。
      第一学期微积分的一个定理说，可微函数之和可微，且
      \( (f+g)^\prime=f^\prime+g^\prime \)；还说
      可微函数的倍数
      可微，且\( (r\cdot f)^\prime=r\,f^\prime \)。
    
\end{ans}
\begin{ans}{二.I.1.42}
      这个验证是常规的。
      注意`\( 1 \)'是\( 1+0i \)，零元素是下面这些。
      \begin{exparts}
        \partsitem \( (0+0i)+(0+0i)x+(0+0i)x^2 \)
        \partsitem \( \begin{mat}
                   0+0i  &0+0i  \\
                   0+0i  &0+0i
                 \end{mat} \)
      \end{exparts}
    
\end{ans}
\begin{ans}{二.I.1.43}
      向量空间的定义中显著缺失的
      是距离的度量。
    
\end{ans}
\begin{ans}{二.I.1.44}
      \begin{exparts}
        \partsitem 小小的重新排列即可奏效。
          \begin{align*}
            (\vec{v}_1+(\vec{v}_2+\vec{v}_3))+\vec{v}_4
            &=((\vec{v}_1+\vec{v}_2)+\vec{v}_3)+\vec{v}_4  \\
            &=(\vec{v}_1+\vec{v}_2)+(\vec{v}_3+\vec{v}_4)  \\
            &=\vec{v}_1+(\vec{v}_2+(\vec{v}_3+\vec{v}_4))  \\
            &=\vec{v}_1+((\vec{v}_2+\vec{v}_3)+\vec{v}_4)
          \end{align*}
          上面每个等号都来自三个向量的结合性，
          它作为一条条件写在向量空间的定义中。
          例如，第二个`$=$'应用法则
          $(\vec{w}_1+\vec{w}_2)+\vec{w}_3=\vec{w}_1+(\vec{w}_2+\vec{w}_3)$，
          取$\vec{w}_1$为$\vec{v}_1+\vec{v}_2$，
          取$\vec{w}_2$为$\vec{v}_3$，
          取$\vec{w}_3$为$\vec{v}_4$。
        \partsitem 归纳的基础是三个向量的情形。
          这一情形
          \( \vec{v}_1+(\vec{v}_2+\vec{v}_3)
          =(\vec{v}_1+\vec{v}_2)+\vec{v}_3 \)
          是向量空间定义中的条件之一。

          归纳步骤：假设三个向量的任意两个和、
          四个向量的任意两个和、
          \ldots、$k$个向量的任意两个和，
          无论我们如何给这些和加括号都相等。
          我们将证明任意\( k+1 \)个向量的和都等于这一个
          \( ((\cdots((\vec{v}_1+\vec{v}_2)+\vec{v}_3)+\cdots)+\vec{v}_k)
              +\vec{v}_{k+1} \)。

          任何加了括号的和都有一个最外层的`\( + \)'。
          设它位于\( \vec{v}_m \)与\( \vec{v}_{m+1} \)之间，
          于是这个和看起来像这样。
          \begin{equation*}
            (\cdots\,\vec{v}_1\cdots\vec{v}_m\,\cdots)
           +(\cdots\,\vec{v}_{m+1}\cdots\vec{v}_{k+1}\,\cdots)
          \end{equation*}
          后一半涉及的加法少于$k+1$次，所以
          由归纳假设我们可以重新加括号，
          使得它从里到外自左向右读，特别地，
          使得其最外层的`$+$'恰好出现在
          $\vec{v}_{k+1}$之前。
          \begin{equation*}
            =(\cdots\,\vec{v}_1\,\cdots\,\vec{v}_m\,\cdots)
             +((\cdots(\vec{v}_{m+1}+\vec{v}_{m+2})+\cdots+\vec{v}_{k})
                 +\vec{v}_{k+1})
          \end{equation*}
          应用三个对象之和的结合性
          \begin{equation*}
            =((\,\cdots\, \vec{v}_1\,\cdots\,\vec{v}_m\,\cdots\,)
             +(\,\cdots\,(\vec{v}_{m+1}+\vec{v}_{m+2})+\cdots\,\vec{v}_k))
             +\vec{v}_{k+1}
          \end{equation*}
          最后在这些最外层括号的内部应用归纳假设，即完成证明。
      \end{exparts}
    
\end{ans}
\begin{ans}{二.I.1.45}
     设$\vec{v}$是$\Re^2$的成员，分量为$v_1$与$v_2$。
     两个分量同号或为~$0$这一条件可以缩写
     为$v_1v_2\geq 0$。

     为证这个集合对标量乘法封闭，注意
     $r\vec{v}$的分量满足$(rv_1)(rv_2)=r^2(v_1v_2)$，
     且$r^2\geq 0$，所以$r^2v_1v_2\geq 0$。

     为证这个集合对加法不封闭，我们只需给出一个
     例子。
     分量为$-1$与~$0$的向量加到
     分量为$0$与~$1$的向量上，得到一个分量为$-1$与~$1$、
     符号混合的向量。
   
\end{ans}
\protect \subsection {二.I.2: 子空间与张成集}
\begin{ans}{二.I.2.20}
      根据 \nearbylemma{th:SubspIffClosed}，要判断 $\matspace_{\nbyn{2}}$
      的每个子集是否为子空间，只需检查它是否非空且封闭。
      \begin{exparts}
        \partsitem 是，容易验证它非空且封闭。
          这就是它的一个参数化描述。
          \begin{equation*}
            \set{a\begin{mat}[r]
                    1  &0  \\
                    0  &0
                  \end{mat}
                 +b\begin{mat}[r]
                    0  &0  \\
                    0  &1
                   \end{mat}
                 \suchthat a,b\in\Re}
          \end{equation*}
          顺便指出，这个参数化描述本身也表明它是一个子空间，
          因为它被表示为由两个矩阵构成的集合张成，
          而任何张成都是子空间。
        \partsitem 是；容易验证它非空且封闭。
         或者，如前一个答案所述，参数化描述的存在性
         表明它是一个子空间。
         要得到参数化描述，
         可将条件 $a+b=0$ 改写为 $a=-b$。
         于是有下式。
          \begin{equation*}
            \set{\begin{mat}
                    -b  &0  \\
                    0   &b
                  \end{mat}
                 \suchthat b\in\Re}
            =\set{b\begin{mat}[r]
                    -1  &0  \\
                    0   &1
                  \end{mat}
                 \suchthat b\in\Re}
          \end{equation*}
        \partsitem 不是。
          它对加法不封闭。
          例如，
          \begin{equation*}
            \begin{mat}[r]
              5  &0  \\
              0  &0
            \end{mat}
            +\begin{mat}[r]
              5  &0  \\
              0  &0
            \end{mat}
            =\begin{mat}[r]
              10  &0  \\
              0  &0
            \end{mat}
          \end{equation*}
          不在该集合中。
          （这个集合对标量乘法也不封闭，
          例如，它不包含零矩阵。）
        \partsitem 是。
          \begin{equation*}
            \set{b\begin{mat}[r]
                    -1  &0  \\
                    0   &1
                  \end{mat}
                 +c\begin{mat}[r]
                    0  &1  \\
                    0  &0
                   \end{mat}
                 \suchthat b,c\in\Re}
          \end{equation*}
      \end{exparts}
     
\end{ans}
\begin{ans}{二.I.2.21}
      不是，它不封闭。
      具体来说，它对标量乘法不封闭，因为它
      不包含零多项式。
    
\end{ans}
\begin{ans}{二.I.2.22}
      方程
      \begin{equation*}
        \colvec{1 \\ 0 \\ 3} =
                 c_1\colvec{2 \\ 1 \\ -1}
                 +c_2\colvec{1 \\ -1 \\ 1}
      \end{equation*}
      给出一个线性方程组
      \begin{equation*}
        \begin{amat}{2}
          2  &1  &1  \\
          1  &-1 &0  \\
          -1 &1  &3
        \end{amat}
        \grstep[(1/2)\rho_1+\rho_3]{(-1/2)\rho_1+\rho_2}
        \begin{amat}{2}
          2  &1    &1  \\
          0  &-3/2 &-1/2  \\
          0  &0  &3
        \end{amat}
      \end{equation*}
      该方程组无解，所以该向量不在张成之中。
    
\end{ans}
\begin{ans}{二.I.2.23}
      \begin{exparts}
         \partsitem 是，由
          \begin{equation*}
            r_1\colvec[r]{1 \\ 0 \\ 0}+r_2\colvec[r]{0 \\ 0 \\ 1}
              =\colvec[r]{2 \\ 0 \\ 1}
          \end{equation*}
          得到的线性方程组的解给出 \( r_1=2 \) 与 \( r_2=1 \)。
        \partsitem 是；由
          \( r_1(x^2)+r_2(2x+x^2)+r_3(x+x^3)=x-x^3 \)
          \begin{equation*}
            \begin{linsys}{3}
                  &  &2r_2 &+ &r_3 &= &1  \\
              r_1 &+ &r_2  &  &    &= &0  \\
                  &  &     &  &r_3 &= &-1
            \end{linsys}
          \end{equation*}
          可得 \( -1(x^2)+1(2x+x^2)-1(x+x^3)=x-x^3 \)。
       \partsitem 不是；所给两个矩阵的任意组合，其右上角元素都为零。
      \end{exparts}
    
\end{ans}
\begin{ans}{二.I.2.24}
      \begin{exparts}
        \partsitem 关系式
          \begin{equation*}
            \colvec{-5 \\ 7}=a\colvec{3 \\ 1}+b\colvec{1 \\ 2}
          \end{equation*}
          给出如下方程组。
          \begin{equation*}
            \begin{linsys}{2}
              3a &+ &b  &= &-5 \\
               a &+ &2b &= &7
            \end{linsys}
            \grstep{-(1/3)\rho_1+\rho_2}
            \begin{linsys}{2}
              3a &+ &b      &= &-5 \\
                 &  &(5/3)b &= &26/3
            \end{linsys}
          \end{equation*}
          所以 $b=26/5$。
          代入 $3a+(26/5)=-5$ 得 $a=-17/5$。
          用本小节的术语来说，向量 $\binom{-5}{7}$
          在另外两个向量构成的集合的张成之中。
        \partsitem 没有藏身之处。
          下面的高斯消元法化简
          \begin{equation*}
            \begin{linsys}{2}
              3a &+ &b  &= &x \\
               a &+ &2b &= &y
            \end{linsys}
            \grstep{-(1/3)\rho_1+\rho_2}
            \begin{linsys}{2}
              3a &+ &b      &= &x\hfill\hbox{} \\
                 &  &(5/3)b &= &-(1/3)x+y
            \end{linsys}
          \end{equation*}
          表明，对任意 $x,y\in\R$ 都存在一对 $a,b\in\R$。
          用本小节的术语来说，这两个向量的张成
          是整个平面。
        \partsitem 没有藏身之处。
          考虑这个方程组。
          \begin{equation*}
            a\colvec{-1 \\ 0 \\ 3}+b\colvec{2 \\ 1 \\ 0}+c\colvec{5 \\ 4 \\ -9}
                = \colvec{x \\ y \\ z}
          \end{equation*}
          化简后
          \begin{align*}
            \begin{linsys}{3}
              -a &+ &2b &+ &5c &= &x \\
                 &  &b  &+ &4c &= &y \\
              3a &  &   &- &9c &= &z \\
            \end{linsys}
            & \grstep{3\rho_1+\rho_3}
            \begin{linsys}{3}
              -a &+ &2b  &+ &5c &= &x\hfill\hbox{} \\
                 &  &b   &+ &4c &= &y\hfill\hbox{} \\
                 &  &6b  &+ &6c &= &3x+z \\
            \end{linsys}                                  \\
            & \grstep{-6\rho_2+\rho_3}
            \begin{linsys}{3}
              -a &+ &2b  &+ &5c  &= &x\hfill\hbox{} \\
                 &  &b   &+ &4c  &= &y\hfill\hbox{} \\
                 &  &    & &-18c &= &3x+-6y+z \\
            \end{linsys}
          \end{align*}
          表明，对任意 $x,y,z\in\R$，
          所需的 $a$、$b$ 与~$c$ 都存在。
          超级英雄的装置张成了整个空间。
      \end{exparts}
    
\end{ans}
\begin{ans}{二.I.2.25}
      \begin{exparts}
        \partsitem 是；它在那个张成之中，因为
          \( 1\cdot\cos^2x+1\cdot\sin^2x=f(x) \)。
        \partsitem 不是，因为 \( r_1\cos^2x+r_2\sin^2x=3+x^2 \) 没有对
          所有 \( x \) 都成立的标量解。
          例如，取 $x$ 为 $0$ 和 $\pi$ 得到两个方程 $r_1\cdot 1+r_2\cdot 0=3$ 与
          $r_1\cdot 1+r_2\cdot 0=3+\pi^2$，二者
          相互矛盾。
        \partsitem 不是；考虑取 $x$ 为 $\pi/2$ 与
          $3\pi/2$ 时的结果。
        \partsitem 是，\( \cos (2x)=1\cdot\cos^2(x)-1\cdot\sin^2(x) \)。
      \end{exparts}
     
\end{ans}
\begin{ans}{二.I.2.26}
      \begin{exparts}
         \partsitem 是，对任意 \( x,y,z\in\Re \)，方程
          \begin{equation*}
             r_1\colvec[r]{1 \\ 0 \\ 0}
             +r_2\colvec[r]{0 \\ 2 \\ 0}
             +r_3\colvec[r]{0 \\ 0 \\ 3}
             =\colvec{x \\ y \\ z}
          \end{equation*}
          有解 \( r_1=x \)，\( r_2=y/2 \) 与
          \( r_3=z/3 \)。
        \partsitem 是，方程
          \begin{equation*}
            r_1\colvec[r]{2 \\ 0 \\ 1}
            +r_2\colvec[r]{1 \\ 1 \\ 0}
            +r_3\colvec[r]{0 \\ 0 \\ 1}
            =\colvec{x \\ y \\ z}
          \end{equation*}
          给出如下
          \begin{equation*}
            \begin{linsys}{3}
              2r_1 &+  &r_2  &  &    &=  &x \\
                   &   &r_2  &  &    &=  &y \\
               r_1 &   &     &+ &r_3 &=  &z \\
            \end{linsys}
            \grstep{-(1/2)\rho_1+\rho_3}\repeatedgrstep{(1/2)\rho_2+\rho_3}
            \begin{linsys}{3}
              2r_1 &+  &r_2  &  &    &=  &x\hfill\hbox{} \\
                   &   &r_2  &  &    &=  &y\hfill\hbox{} \\
                   &   &     &  &r_3 &=  &-(1/2)x+(1/2)y+z \\
            \end{linsys}
          \end{equation*}
          于是，给定任意 $x$、$y$ 和 $z$，我们可以算出
          \( r_3=(-1/2)x+(1/2)y+z \)、\( r_2=y \) 与
          \( r_1=(1/2)x-(1/2)y \)。
        \partsitem 不是。
          特别地，我们无法得到向量
          \begin{equation*}
            \colvec[r]{0 \\ 0 \\ 1}
          \end{equation*}
          的线性组合，因为所给的两个
          向量的第三个分量都为零。
       \partsitem 是。
         方程
         \begin{equation*}
           r_1\colvec[r]{1 \\ 0 \\ 1}
           +r_2\colvec[r]{3 \\ 1 \\ 0}
           +r_3\colvec[r]{-1\\ 0 \\ 0}
           +r_4\colvec[r]{2 \\ 1 \\ 5}
           =\colvec{x \\ y \\ z}
         \end{equation*}
         导出如下化简。
         \begin{equation*}
           \begin{amat}{4}
             1  &3  &-1  &2  &x  \\
             0  &1  &0   &1  &y  \\
             1  &0  &0   &5  &z
           \end{amat}
           \grstep{-\rho_1+\rho_3}\repeatedgrstep{3\rho_2+\rho_3}
           \begin{amat}{4}
             1  &3  &-1  &2  &x \\
             0  &1  &0   &1  &y  \\
             0  &0  &1   &6  &-x+3y+z
           \end{amat}
         \end{equation*}
         它有无穷多解。
         例如，我们可以取 $r_4$ 为零，然后按通常的
         回代方法，用 $x$、$y$ 和 $z$ 解出
         $r_3$、$r_2$ 和 $r_1$。
       \partsitem 不是。
         方程
         \begin{equation*}
           r_1\colvec[r]{2 \\ 1 \\ 1}
           +r_2\colvec[r]{3 \\ 0 \\ 1}
           +r_3\colvec[r]{5 \\ 1 \\ 2}
           +r_4\colvec[r]{6 \\ 0 \\ 2}
           =\colvec{x \\ y \\ z}
         \end{equation*}
         导出如下化简。
         \begin{multline*}
           \begin{amat}{4}
             2  &3  &5   &6  &x  \\
             1  &0  &1   &0  &y  \\
             1  &1  &2   &2  &z
           \end{amat}                                             \\
           \grstep[-(1/2)\rho_1+\rho_3]{-(1/2)\rho_1+\rho_2}
           \repeatedgrstep{-(1/3)\rho_2+\rho_3}
           \begin{amat}{4}
             2  &3     &5     &6  &x \\
             0  &-3/2  &-3/2  &-3 &-(1/2)x+y  \\
             0  &0     &0     &0  &-(1/3)x-(1/3)y+z
           \end{amat}
         \end{multline*}
         这表明并非每个三元向量都能如此表示。
         只有满足约束 $-(1/3)x-(1/3)y+z=0$ 的向量才在张成之中。
         （要看出任何这样的向量确实可以表示，
         可取 $r_3$ 与 $r_4$
         为零，然后用回代法以 $x$、$y$ 和
         $z$ 解出 $r_1$ 与 $r_2$。）
      \end{exparts}
    
\end{ans}
\begin{ans}{二.I.2.27}
      \begin{exparts}
        \partsitem \( \set{\rowvec{c &b &c}\suchthat b,c\in\Re}
                      =\set{b\rowvec{0 &1 &0}+c\rowvec{1 &0 &1}
                        \suchthat b,c\in\Re} \)
           张成集合的显然选择是
           $\set{\rowvec{0 &1 &0},\rowvec{1 &0 &1}}$。
        \partsitem \( \set{\begin{mat}
                       -d &b  \\
                       c  &d
                      \end{mat} \suchthat b,c,d\in\Re}
                     =\set{b\begin{mat}[r]
                       0  &1  \\
                       0  &0
                      \end{mat}
                     +c\begin{mat}[r]
                       0  &0  \\
                       1  &0
                      \end{mat}
                     +d\begin{mat}[r]
                       -1  &0  \\
                       0  &1
                      \end{mat}  \suchthat b,c,d\in\Re} \)
            张成这个空间的一个集合就由这三个矩阵组成。
        \partsitem 方程组
          \begin{equation*}
            \begin{linsys}{4}
              a  &+  &3b  &   &   &  &  &=  &0  \\
             2a  &   &    &   &-c &- &d &=  &0
            \end{linsys}
          \end{equation*}
          给出 \( b=-(c+d)/6 \) 与 \( a=(c+d)/2 \)。
          于是一种描述如下。
          \begin{equation*}
            \set{c\begin{mat}[r]
                       1/2  &-1/6  \\
                       1    &0
                      \end{mat}
                     +d\begin{mat}[r]
                       1/2  &-1/6  \\
                       0    &1
                      \end{mat}  \suchthat c,d\in\Re}
           \end{equation*}
          这表明张成该子空间的一个集合由那
          两个矩阵组成。
        \partsitem 由 $a=2b-c$ 可知，集合
           \( \set{(2b-c)+bx+cx^3 \suchthat b,c\in\Re} \)
           等于集合
           \( \set{b(2+x)+c(-1+x^3) \suchthat b,c\in\Re}  \)。
           所以该子空间是集合 $\set{2+x, -1+x^3}$ 的张成。
        \partsitem 集合
          \( \set{a+bx+cx^2\suchthat a+7b+49c=0} \)
          可以参数化为
          \begin{equation*}
            \set{b(-7+x)+c(-49+x^2)\suchthat b,c\in\Re}
          \end{equation*}
          因此
          有张成集 $\set{-7+x,-49+x^2}$。
      \end{exparts}
    
\end{ans}
\begin{ans}{二.I.2.28}
      \begin{exparts}
        \partsitem 我们可以这样参数化
          \begin{equation*}
             \set{\colvec{x \\ 0 \\ z}\suchthat x,z\in\Re}
             =\set{x\colvec[r]{1 \\ 0 \\ 0}
                  +z\colvec[r]{0 \\ 0 \\ 1}\suchthat x,z\in\Re}
          \end{equation*}
          由此得到张成集如下。
          \begin{equation*}
             \set{\colvec[r]{1 \\ 0 \\ 0},\colvec[r]{0 \\ 0 \\ 1}}
          \end{equation*}
        \item 这是参数化描述，以及相应的张成集。
          \begin{equation*}
             \set{y\colvec[r]{-2/3 \\ 1 \\ 0}+z\colvec[r]{-1/3 \\ 0 \\ 1}
                 \suchthat y,z\in\Re }
             \qquad
             \set{\colvec[r]{-2/3 \\ 1 \\ 0},\colvec[r]{-1/3 \\ 0 \\ 1} }
          \end{equation*}
        \partsitem \( \set{\colvec[r]{1 \\ -2 \\ 1 \\ 0},
                      \colvec[r]{-1/2 \\ 0 \\ 0 \\ 1} } \)
        \partsitem 把描述参数化为
          \( \set{-a_1+a_1x+a_3x^2+a_3x^3\suchthat a_1,a_3\in\Re } \)
          得到 \( \set{-1+x,x^2+x^3}. \)
        \partsitem \( \set{1,x,x^2,x^3,x^4} \)
        \partsitem \( \set{ \begin{mat}[r]
                   1  &0  \\
                   0  &0
                 \end{mat},
                 \begin{mat}[r]
                   0  &1  \\
                   0  &0
                 \end{mat},
                 \begin{mat}[r]
                   0  &0  \\
                   1  &0
                 \end{mat},
                 \begin{mat}[r]
                   0  &0  \\
                   0  &1
                 \end{mat} } \)
      \end{exparts}
    
\end{ans}
\begin{ans}{二.I.2.29}
      严格地说，不是。
      \( \Re^3 \) 的子空间都是三元向量构成的集合，而
      \( \Re^2 \) 是二元向量构成的集合。
      不过显然，\( \Re^2 \) 与 \( \Re^3 \) 的这个子空间“几乎
      一样”。
      \begin{equation*}
        \set{\colvec{x \\ y \\ 0}\suchthat x,y\in\Re}
      \end{equation*}
    
\end{ans}
\begin{ans}{二.I.2.30}
      当然，加法与标量乘法运算就是
      从外围空间继承的那两个运算。
      \begin{exparts}
        \partsitem 这是一个子空间。
          它不是空集，因为它至少包含上面给出的两个示例函数。
          它是封闭的，因为若 \( f_1,f_2 \) 是偶函数且
          \( c_1,c_2 \) 是标量，则有下式。
          \begin{equation*}
            (c_1f_1+c_2f_2)\,(-x)
            =c_1\,f_1(-x)+c_2\,f_2(-x)
            =c_1\,f_1(x)+c_2\,f_2(x)
            =(c_1f_1+c_2f_2)\,(x)
          \end{equation*}
        \partsitem 这也是一个子空间；验证方法与
          前一个类似。
      \end{exparts}
    
\end{ans}
\begin{ans}{二.I.2.31}
      它可以不是真子空间。
      例如，
      若向量空间是 \( \Re^1 \) 且 \( \vec{v}\neq\zero\, \)，
      则子空间 $\set{r\cdot\vec{v}\suchthat r\in\Re}$
      就是整个 $\Re^1$。
    
\end{ans}
\begin{ans}{二.I.2.32}
      不构成，这样的集合不封闭。
      首先，它不包含零向量。
    
\end{ans}
\begin{ans}{二.I.2.33}
     \begin{exparts}
       \item $\matspace_{\nbyn{2}}$ 的这个非空子集不是子空间。
         \begin{equation*}
           A=\set{
             \begin{mat}[r]
               1 &2 \\
               3 &4
             \end{mat},
             \begin{mat}[r]
               5 &6 \\
               7 &8
             \end{mat}}
         \end{equation*}
         它不是 $\matspace_{\nbyn{2}}$ 的子空间的一个理由是
         它不包含零矩阵。
         （另一个理由是它对加法不封闭，因为两者的和
         不是 $A$ 的元素。
         它对标量乘法也不封闭。）
        \item 这个由两个向量构成的集合不张成 $\Re^2$。
          \begin{equation*}
            \set{\colvec[r]{1 \\ 1},\colvec[r]{3 \\ 3}}
          \end{equation*}
          这两个向量的任何线性组合都无法给出
          第二个分量不等于第一个分量的向量。
      \end{exparts}
   
\end{ans}
\begin{ans}{二.I.2.34}
      不是。
      \( \Re^1 \) 仅有的子空间是它本身与其
      平凡子空间。
      $\Re$ 的任何包含非零成员 $\vec{v}$ 的子空间 $S$
      必须包含其全部标量倍数构成的集合
      $\set{r\cdot\vec{v}\suchthat r\in\Re}$。
      但这个集合就是整个 $\Re$。
    
\end{ans}
\begin{ans}{二.I.2.35}
      第~(1) 项已在正文中验证。

      第~(2) 条有五个条件。
      第一条是封闭性：若 \( c\in\Re \) 且 \( \vec{s}\in S \)，则
      \( c\cdot\vec{s}\in S \)，因为
      \( c\cdot\vec{s}=c\cdot\vec{s}+0\cdot\zero \)。
      第二，因为 \( S \) 中的运算继承自 \( V \)，所以对 \( c,d\in\Re \) 与 \( \vec{s}\in S \)，
      \( S \) 中的标量乘积 \( (c+d)\cdot\vec{s}\, \) 等于 \( V \) 中的乘积
      \( (c+d)\cdot\vec{s}\, \)，而它等于 \( V \) 中的
      \( c\cdot\vec{s}+d\cdot\vec{s}\, \)，进而等于 \( S \) 中的
      \( c\cdot\vec{s}+d\cdot\vec{s}\, \)。

      第三、第四和第五个条件的验证与刚才给出的第二个条件的验证类似。
    
\end{ans}
\begin{ans}{二.I.2.36}
      前一小节的一道练习证明了每个向量空间只有一个零向量（也就是说，只有
      一个向量是该空间的加法单位元）。
      而平凡空间只有一个元素，且该元素必是这个（唯一的）零向量。
    
\end{ans}
\begin{ans}{二.I.2.37}
      如提示所言，根本的原因是第一章的线性组合引理。
      为给出完整证明，我们将证明这两个集合互相包含。

      第一个包含关系 \( \spanof{ \spanof{S} }\supseteq\spanof{S}  \)
      是一个更一般且显然的事实的特例：对向量空间的任意子集
      \( T \)，都有 \( \spanof{T}\supseteq T \)。

      对于另一个包含关系，
      即 \( \spanof{ \spanof{S} }\subseteq\spanof{S} \)，
      从 \( \spanof{S} \) 中取 $m$ 个向量，即
      \( c_{1,1}\vec{s}_{1,1}+\cdots+c_{1,n_1}\vec{s}_{1,n_1} \)、\ldots、
      \( c_{1,m}\vec{s}_{1,m}+\cdots+c_{1,n_m}\vec{s}_{1,n_m} \)，
      并注意它们的任意线性组合
      \begin{equation*}
         r_1(c_{1,1}\vec{s}_{1,1}+\cdots+c_{1,n_1}\vec{s}_{1,n_1})+\cdots
        +r_m(c_{1,m}\vec{s}_{1,m}+\cdots+c_{1,n_m}\vec{s}_{1,n_m})
      \end{equation*}
      是 \( S \) 中元素的线性组合
      \begin{equation*}
        = (r_1c_{1,1})\vec{s}_{1,1}+\cdots+(r_1c_{1,n_1})\vec{s}_{1,n_1}+\cdots
        +(r_mc_{1,m})\vec{s}_{1,m}+\cdots+(r_mc_{1,n_m})\vec{s}_{1,n_m}
      \end{equation*}
      因而在 \( \spanof{S} \) 中。
      也就是说，只需回想：线性组合的线性组合（$S$ 中成员的）
      仍是线性组合（同样还是 $S$ 中成员的）。
    
\end{ans}
\begin{ans}{二.I.2.38}
      \begin{exparts}
         \partsitem 它不是子空间，因为这些运算不是继承来的运算。
           例如，在这个空间中，
           \begin{equation*}
             0\cdot\colvec{x \\ y \\ z}=\colvec[r]{1 \\ 0 \\ 0}
           \end{equation*}
           而这当然在 $\Re^3$ 中不成立。
         \partsitem 我们可以把证明对加法封闭的论证
           与证明对标量乘法封闭的论证合并成
           一个证明对两个向量的线性组合封闭的单一论证。
           若 $r_1,r_2,x_1,x_2,y_1,y_2,z_1,z_2$ 都属于 $\Re$，则
           \begin{multline*}
              r_1\colvec{x_1 \\ y_1 \\ z_1}
              +r_2\colvec{x_2 \\ y_2 \\ z_2}
              =\colvec{r_1x_1-r_1+1 \\ r_1y_1 \\ r_1z_1}
               +\colvec{r_2x_2-r_2+1 \\ r_2y_2 \\ r_2z_2}                 \\
            =\colvec{r_1x_1-r_1+r_2x_2-r_2+1 \\ r_1y_1+r_2y_2 \\ r_1z_1+r_2z_2}
           \end{multline*}
           （注意这个空间中加法的定义是：第一个分量按
           $(r_1x_1-r_1+1)+(r_2x_2-r_2+1)-1$ 的方式合并，
           所以最后一个向量的第一个分量并非表示
           `$\hbox{}+2$'）。
           将最后一个向量的三个分量相加得
           $r_1(x_1-1+y_1+z_1)+r_2(x_2-1+y_2+z_2)+1=r_1\cdot0+r_2\cdot0+1=1$。

           其余大部分条件的验证都很容易（尽管这些运算的古怪之处
           使它们无法按常规进行）。
           加法的交换律如下。
           \begin{equation*}
             \colvec{x_1 \\ y_1 \\ z_1}+\colvec{x_2 \\ y_2 \\ z_2}
             =\colvec{x_1+x_2-1 \\ y_1+y_2 \\ z_1+z_2}
             =\colvec{x_2+x_1-1 \\ y_2+y_1 \\ z_2+z_1}
             =\colvec{x_2 \\ y_2 \\ z_2}+\colvec{x_1 \\ y_1 \\ z_1}
           \end{equation*}
           加法的结合律则是
           \begin{equation*}
             (\colvec{x_1 \\ y_1 \\ z_1}+\colvec{x_2 \\ y_2 \\ z_2})
              +\colvec{x_3 \\ y_3 \\ z_3}
             =\colvec{(x_1+x_2-1)+x_3-1 \\ (y_1+y_2)+y_3 \\ (z_1+z_2)+z_3}
           \end{equation*}
           而
           \begin{equation*}
             \colvec{x_1 \\ y_1 \\ z_1}
             +(\colvec{x_2 \\ y_2 \\ z_2}+\colvec{x_3 \\ y_3 \\ z_3})
             =\colvec{x_1+(x_2+x_3-1)-1 \\ y_1+(y_2+y_3) \\ z_1+(z_2+z_3)}
           \end{equation*}
           且两者相等。
           关于这个加法运算的单位元
           如下
           \begin{equation*}
             \colvec{x \\ y \\ z}+\colvec[r]{1 \\ 0 \\ 0}
             =\colvec{x+1-1 \\ y+0 \\ z+0}
             =\colvec{x \\ y \\ z}
           \end{equation*}
           加法逆元类似。
           \begin{equation*}
             \colvec{x \\ y \\ z}+\colvec{-x+2 \\ -y \\ -z}
             =\colvec{x+(-x+2)-1 \\ y-y \\ z-z}
             =\colvec[r]{1 \\ 0 \\ 0}
           \end{equation*}

           关于标量乘法的那些条件也容易验证。
           对第一个条件，
           \begin{equation*}
             (r+s)\colvec{x \\ y \\ z}
             =\colvec{(r+s)x-(r+s)+1 \\ (r+s)y \\ (r+s)z}
           \end{equation*}
           而
           \begin{multline*}
             r\colvec{x \\ y \\ z}+s\colvec{x \\ y \\ z}
             =\colvec{rx-r+1 \\ ry \\ rz}+\colvec{sx-s+1 \\ sy \\ sz}  \\
             =\colvec{(rx-r+1)+(sx-s+1)-1 \\ ry+sy \\ rz+sz}
           \end{multline*}
           两者相等。
           第二个条件比较
           \begin{equation*}
             r\cdot(\colvec{x_1 \\ y_1 \\ z_1}+\colvec{x_2 \\ y_2 \\ z_2})
             =r\cdot\colvec{x_1+x_2-1 \\ y_1+y_2 \\ z_1+z_2}
             =\colvec{r(x_1+x_2-1)-r+1 \\ r(y_1+y_2) \\ r(z_1+z_2)}
           \end{equation*}
           与
           \begin{multline*}
             r\colvec{x_1 \\ y_1 \\ z_1}+r\colvec{x_2 \\ y_2 \\ z_2}
             =\colvec{rx_1-r+1 \\ ry_1 \\ rz_1}
                     +\colvec{rx_2-r+1 \\ ry_2 \\ rz_2}                  \\
             =\colvec{(rx_1-r+1)+(rx_2-r+1)-1 \\ ry_1+ry_2 \\ rz_1+rz_2}
           \end{multline*}
           且它们相等。
           对第三个条件，
           \begin{equation*}
             (rs)\colvec{x \\ y \\ z}
             =\colvec{rsx-rs+1 \\ rsy \\ rsz}
           \end{equation*}
           而
           \begin{equation*}
             r(s\colvec{x \\ y \\ z})
             =r(\colvec{sx-s+1 \\ sy \\ sz})
             =\colvec{r(sx-s+1)-r+1 \\ rsy \\ rsz}
           \end{equation*}
           两者相等。
           对于被 $1$ 的标量乘法，我们有下式。
           \begin{equation*}
             1\cdot\colvec{x \\ y \\ z}
             =\colvec{1x-1+1 \\ 1y \\ 1z}
             =\colvec{x \\ y \\ z}
           \end{equation*}
           这样，这两个运算满足向量空间的所有条件。

           \textit{注记。}
           理解这个向量空间的一种方式是把它看作 $\Re^3$ 中的平面
           \begin{equation*}
             P=\set{\colvec{x \\ y \\ z}\suchthat x+y+z=0}
           \end{equation*}
           沿 $x$ 轴离开原点平移 $1$ 所得的平面。
           于是加法变成：~要加上这个空间的两个成员
           \begin{equation*}
             \colvec{x_1 \\ y_1 \\ z_1},\;\colvec{x_2 \\ y_2 \\ z_2}
           \end{equation*}
           （其中 $x_1+y_1+z_1=1$ 且 $x_2+y_2+z_2=1$）
           先把它们移回 $1$ 放入 $P$，再按通常方式相加，
           \begin{equation*}
             \colvec{x_1-1 \\ y_1 \\ z_1}+\colvec{x_2-1 \\ y_2 \\ z_2}
             =\colvec{x_1+x_2-2 \\ y_1+y_2 \\ z_1+z_2}
             \qquad\text{（在 $P$ 中）}
           \end{equation*}
           然后把结果沿 $x$ 轴再移出 $1$。
           \begin{equation*}
             \colvec{x_1+x_2-1 \\ y_1+y_2 \\ z_1+z_2}.
           \end{equation*}
           标量乘法类似。
         \partsitem 为使该子空间对继承的标量
           乘法封闭，设 $\vec{v}$ 是该子空间的一个成员，
           \begin{equation*}
             0\cdot\vec{v}=\colvec[r]{0 \\ 0 \\ 0}
           \end{equation*}
           它也必须是成员。

           逆命题不成立。
           下面是 $\Re^3$ 的一个包含原点的子集
           \begin{equation*}
             \set{\colvec[r]{0 \\ 0 \\ 0},\colvec[r]{1 \\ 0 \\ 0}}
           \end{equation*}
           （这个子集只有两个元素）但它不是子空间。
       \end{exparts}
     
\end{ans}
\begin{ans}{二.I.2.39}
      \begin{exparts}
        \item $(\vec{v}_1+\vec{v}_2+\vec{v}_3)-(\vec{v}_1+\vec{v}_2)
                 =\vec{v}_3$
        \item $(\vec{v}_1+\vec{v}_2)-(\vec{v}_1)
                 =\vec{v}_2$
        \item 显然，$\vec{v}_1$。
        \item 取那个只含一项的和再相减，得
          （$\vec{v}_1)-\vec{v}_1=\zero$。
      \end{exparts}
    
\end{ans}
\begin{ans}{二.I.2.40}
      是的；任何空间都是它自身的子空间，所以每个空间都包含另一个。
    
\end{ans}
\begin{ans}{二.I.2.41}
      \begin{exparts}
        \partsitem \( \Re^2 \) 中 \( x \) 轴与 \( y \) 轴的并就是一例。
        \partsitem 整数集，作为 \( \Re^1 \) 的子集，就是一例。
        \partsitem \( \Re^2 \) 的子集 \( \set{\vec{v}} \) 就是一例，
          其中 $\vec{v}$ 是任意非零向量。
      \end{exparts}
     
\end{ans}
\begin{ans}{二.I.2.42}
      由于向量空间加法是交换的，重新排列被加项不会改变线性组合。
    
\end{ans}
\begin{ans}{二.I.2.43}
      我们总是在一个包围空间的背景下考虑这个张成。
    
\end{ans}
\begin{ans}{二.I.2.44}
      “当”与“仅当”两个方向都成立。

      对于“当”的方向，
      设 \( S \) 是向量空间 \( V \) 的子集，并假设
      \( \vec{v}\in S \) 满足
      \( \vec{v}=c_1\vec{s}_1+\dots+c_n\vec{s}_n \)，其中
      \( c_1,\ldots,c_n \) 是标量，且
      \( \vec{s}_1,\ldots,\vec{s}_n\in S \)。
      我们必须证明 \( \spanof{S\union\set{\vec{v}} }=\spanof{S} \)。

      一侧的包含关系，
      \( \spanof{S}\subseteq\spanof{S\union\set{\vec{v}} } \) 是显然的。
      对于另一个方向，
      \( \spanof{S\union\set{\vec{v}} }\subseteq\spanof{S} \)，注意若一个
      向量属于左边的集合，则它形如
      \( d_0\vec{v}+d_1\vec{t}_1+\dots+d_m\vec{t}_m \)，其中诸 \( d \) 是
      标量，诸 \( \vec{t}\, \) 属于 \( S \)。
      把它改写为
      \( d_0(c_1\vec{s}_1+\dots+c_n\vec{s}_n)
      +d_1\vec{t}_1+\cdots+d_m\vec{t}_m \)，并注意
      该结果是 \( S \) 的张成中的一个成员。

      “仅当”显然成立\Dash 加入 \( \vec{v} \)
      会把张成扩大到至少包含 \( \vec{v} \)。
    
\end{ans}
\begin{ans}{二.I.2.45}
      \begin{exparts}
        \partsitem 总是。

          设 \( A,B \) 是 \( V \) 的子空间。
          注意
          它们的交不空，因为两者都包含零向量。
          若 \( \vec{w},\vec{s}\in A\intersection B \) 且 \( r,s \) 是
          标量，则 \( r\vec{v}+s\vec{w}\in A \)，因为
          每个向量都属于 \( A \)，所以线性组合属于 \( A \)，
          出于同样的理由 \(r\vec{v}+s\vec{w}\in B \)。
          于是这个交是封闭的。
          这时 \nearbylemma{th:SubspIffClosed} 适用。
        \partsitem 有时（更准确地说，仅当 \( A\subseteq B \) 或
          \( B\subseteq A \)）。

          为看出答案不是“总是”，取 \( V \) 为 \( \Re^3 \)，
          取 \( A \) 为 $x$ 轴，\( B \) 为
          \( y \) 轴。
          注意
          \begin{equation*}
            \colvec[r]{1 \\ 0}\in A \text{ 且 }\colvec[r]{0 \\ 1}\in B
            \quad\text{而}\quad
            \colvec[r]{1 \\ 0}+\colvec[r]{0 \\ 1}\not\in A\union B
          \end{equation*}
          因为该和既不在 \( A \) 中也不在 \( B \) 中。

          答案不是“从不”，因为若 \( A\subseteq B \) 或
          \( B\subseteq A \)，则显然 \( A\union B \) 是子空间。

          为证明 \( A\union B \) 是子空间仅当一个子空间
          包含另一个，我们假设 \( A\not\subseteq B \)
          且 \( B\not\subseteq A \)，并证明
          该并不是子空间。
          \( A \) 不是 \( B \) 的子集这一假设意味着存在
          \( \vec{a}\in A \) 且 \( \vec{a}\not\in B \)。
          另一假设给出一个 \( \vec{b}\in B \)，满足
          \( \vec{b}\not\in A \)。
          考虑 \( \vec{a}+\vec{b} \)。
          注意该和不是 \( A \) 的元素，否则
          \( (\vec{a}+\vec{b})-\vec{a} \) 将属于 \( A \)，但它不属于。
          类似地，该和也不是 \( B \) 的元素。
          因此该和不是 \( A\union B \) 的元素，所以该并
          不是子空间。
        \partsitem 从不。
          由于 \( A \) 是子空间，它包含零向量，因此
          作为 $A$ 的补集的集合 $V-A$ 不包含它。
          没有零向量，这个补集不可能是向量空间。
      \end{exparts}
    
\end{ans}
\begin{ans}{二.I.2.46}
      一个集合的张成不依赖于包围它的空间。
      \( S \) 中向量的线性组合给出同样的和，
      无论我们把运算视为 \( W \) 的还是 \( V \) 的，
      因为 \( W \) 的运算继承自 \(  V \)。
    
\end{ans}
\begin{ans}{二.I.2.47}
      是传递的；
      应用 \nearbylemma{th:SubspIffClosed}。
      （你必须考虑下面这一点。
      设 \( B \) 是向量空间 \( V \) 的子空间，再设
      \( A\subseteq B\subseteq V \) 是一个子空间。
      \( A \) 从哪个空间继承它的运算？
      答案是这无关紧要\Dash 两种情形下 \( A \) 继承的
      都是相同的运算。）
    
\end{ans}
\begin{ans}{二.I.2.48}
      \begin{exparts}
         \partsitem 总是；
           若 \( S\subseteq T \)，则 \( S \) 中元素的线性组合
           也是 \( T \) 中元素的线性组合。
         \partsitem 有时（更准确地说，当且仅当
           \( S\subseteq T \) 或 \( T\subseteq S \)）。

           答案不是“总是”，\( \Re^3 \) 中的这个例子
           表明了这一点
           \begin{equation*}
             S=\set{\colvec[r]{1 \\ 0 \\ 0},\colvec[r]{0 \\ 1 \\ 0}},\quad
             T=\set{\colvec[r]{1 \\ 0 \\ 0},\colvec[r]{0 \\ 0 \\ 1}}
           \end{equation*}
           因为有下式。
           \begin{equation*}
             \colvec[r]{1 \\ 1 \\ 1}\in\spanof{S\union T}
             \qquad
             \colvec[r]{1 \\ 1 \\ 1}\not\in\spanof{S}\union \spanof{T}
           \end{equation*}

           答案不是“从不”，因为若其中一个集合包含另一个，
           则等式显然成立。
           我们可以通过假设 \( S\not\subseteq T \)（意味着存在
           向量 \( \vec{s}\in S \) 满足
           \( \vec{s}\not\in T \)）
           且 \( T\not\subseteq S \)（给出一个 \( \vec{t}\in T \)，满足
           \( \vec{t}\not\in S \)），并注意
           \( \vec{s}+\vec{t}\in\spanof{S\union T} \)，
           再证明
           \( \vec{s}+\vec{t}\not\in\spanof{S}\union\spanof{T} \)，
           来刻画等式仅在其中一个集合包含另一个时才发生。
         \partsitem 有时。

           显然
           \( \spanof{S\intersection T}
             \subseteq\spanof{S}\intersection\spanof{T} \)
           因为来自
           \( S\intersection T \)
           的向量的任何线性组合既是 \( S \) 中向量的组合，也是
           \( T \) 中向量的组合。

           反方向的包含并不总成立。
           例如，在 \( \Re^2 \) 中，取
           \begin{equation*}
             S=\set{\colvec[r]{1 \\ 0},\colvec[r]{0 \\ 1}},\quad
             T=\set{\colvec[r]{2 \\ 0}}
           \end{equation*}
           于是 \( \spanof{S}\intersection\spanof{T} \) 是 \( x \) 轴，
           但 \( \spanof{S\intersection T}  \) 是平凡子空间。

           要精确刻画等式何时成立并不容易。
           显然若其中一个集合包含另一个则等式成立，但由 \( \Re^3 \) 中的这个例子
           可知那不是“仅当”的。
           \begin{equation*}
             S=\set{\colvec[r]{1 \\ 0 \\ 0},\colvec[r]{0 \\ 1 \\ 0}},
             \quad
             T=\set{\colvec[r]{1 \\ 0 \\ 0},\colvec[r]{0 \\ 0 \\ 1}}
           \end{equation*}
         \partsitem 从不，因为补集的张成是子空间，而
           张成的补不是（它不包含零
           向量）。
      \end{exparts}
     
\end{ans}
\begin{ans}{二.I.2.49}
      要出现这种情况，保证加法与标量乘法运算合理性的条件中
      必须有一个被违反。
      考虑带这些运算的 \( \Re^2 \)。
      \begin{equation*}
        \colvec{x_1 \\ y_1}
        +\colvec{x_2 \\ y_2}
        =\colvec[r]{0 \\ 0}
        \qquad
        r\colvec{x \\ y}
        =\colvec[r]{0 \\ 0}
      \end{equation*}
      集合 $\Re^2$ 对这些运算封闭。
      但它不是向量空间。
      \begin{equation*}
        1\cdot\colvec[r]{1 \\ 1}
        \neq\colvec[r]{1 \\ 1}
      \end{equation*}
    
\end{ans}
\protect \section {第 II 节：线性无关}
\protect \subsection {二.II.1: 定义与例题}
\begin{ans}{二.II.1.21}
      对这里的每一小题，当子集是线性无关的时候你必须证明这一点，而当子集
      是线性相关的时候你必须给出一个相关关系的例子。
      \begin{exparts}
        \partsitem 它是线性相关的。
          考虑
          \begin{equation*}
             c_1\colvec{1 \\ -3 \\ 5}
             +c_2\colvec{2 \\ 2 \\ 4}
             +c_3\colvec{4 \\ -4 \\ 14}
             =\colvec{0 \\ 0 \\ 0}
          \end{equation*}
          就得到这个线性方程组。
          \begin{equation*}
            \begin{linsys}{3}
              c_1  &+  &2c_2  &+  &4c_3  &=  &0  \\
              -3c_1&+  &2c_2  &-  &4c_3  &=  &0  \\
              5c_1 &+  &4c_2  &+  &14c_3 &=  &0
            \end{linsys}
          \end{equation*}
          用高斯消元法
          \begin{equation*}
            \begin{amat}[r]{3}
              1  &2  &4  &0  \\
              -3 &2  &-4 &0  \\
              5  &4  &14 &0
            \end{amat}
            \grstep[-5\rho_1+\rho_3]{3\rho_1+\rho_2}
            \repeatedgrstep{(3/4)\rho_2+\rho_3}
            \begin{amat}[r]{3}
              1  &2  &4  &0  \\
              0  &8  &8  &0  \\
              0  &0  &0  &0
            \end{amat}
          \end{equation*}
          会得到一个自由变量，所以有无穷多解。
          作为相关关系的一个具体例子，我们可以取 $c_3$ 为，比如说，$1$。
          于是得到 \( c_2=-1 \) 和 \( c_1=-2 \)。
        \partsitem 它是线性相关的。
          这里出现的线性方程组
          \begin{equation*}
            \begin{amat}[r]{3}
              1  &2  &3  &0  \\
              7  &7  &7  &0  \\
              7  &7  &7  &0
            \end{amat}
            \grstep[-7\rho_1+\rho_3]{-7\rho_1+\rho_2}
            \repeatedgrstep{-\rho_2+\rho_3}
            \begin{amat}[r]{3}
              1  &2  &3   &0  \\
              0  &-7 &-14 &0  \\
              0  &0  &0   &0
            \end{amat}
          \end{equation*}
          有无穷多解。
          我们可以取 $c_3$ 为，比如说，$1$，再回代求出相应的 $c_2$ 和 $c_1$，
          从而得到一个特解。
        \partsitem 它是线性无关的。
          方程组
          \begin{equation*}
            \begin{amat}[r]{2}
              0  &1  &0  \\
              0  &0  &0  \\
              -1 &4  &0
            \end{amat}
            \grstep{\rho_1\leftrightarrow\rho_2}
            \repeatedgrstep{\rho_3\leftrightarrow\rho_1}
            \begin{amat}{2}
              -1 &4  &0  \\
              0  &1  &0  \\
              0  &0  &0
            \end{amat}
          \end{equation*}
          只有解 $c_1=0$ 与 $c_2=0$。
          （我们本来也可以凭观察得出答案\Dash 第二个向量显然不是第一个
          向量的倍数，反之亦然。）
        \partsitem 它是线性相关的。
          线性方程组
          \begin{equation*}
            \begin{amat}[r]{4}
              9  &2  &3  &12  &0  \\
              9  &0  &5  &12  &0  \\
              0  &1  &-4 &-1  &0
            \end{amat}
          \end{equation*}
          中未知量的个数多于方程的个数，因此用高斯消元法化简后最后必定
          至少有一个自由变量（不可能出现矛盾方程，因为该方程组是齐次的，
          从而至少有全零解）。
          为了展示出一个组合，我们可以作如下化简
          \begin{equation*}
            \grstep{-\rho_1+\rho_2}
            \grstep{(1/2)\rho_2+\rho_3}
            \begin{amat}[r]{4}
              9  &2  &3  &12  &0  \\
              0  &-2 &2  &0   &0  \\
              0  &0  &-3 &-1  &0
            \end{amat}
          \end{equation*}
          然后取，比如说，$c_4=1$。
          于是有 $c_3=-1/3$、$c_2=-1/3$ 以及 $c_1=-31/27$。
      \end{exparts}
    
\end{ans}
\begin{ans}{二.II.1.22}
      在线性无关的情形，你必须证明它确实是线性无关的。否则，你必须展示出
      一个相关关系。（这里我们给出一个具体的相关关系，但也可能有别的。）
      \begin{exparts}
        \partsitem 这个集合是线性无关的。
          建立关系式
          \( c_1(3-x+9x^2)+c_2(5-6x+3x^2)+c_3(1+1x-5x^2)=0+0x+0x^2 \)
          就得到线性方程组
          \begin{equation*}
            \begin{amat}[r]{3}
              3  &5  &1  &0  \\
              -1 &-6 &1  &0  \\
              9  &3  &-5 &0
            \end{amat}
            \grstep[-3\rho_1+\rho_3]{(1/3)\rho_1+\rho_2}
            \repeatedgrstep{3\rho_2}
            \repeatedgrstep{-(12/13)\rho_2+\rho_3}
            \begin{amat}[r]{3}
              3  &5   &1        &0  \\
              0  &-13 &4        &0  \\
              0  &0   &-128/13  &0
            \end{amat}
          \end{equation*}
          只有唯一解：\( c_1=0 \)、\( c_2=0 \) 以及 \( c_3=0 \)。
        \partsitem 这个集合是线性无关的。
           我们可以直接从线性无关的定义出发，凭观察看出这一点。
           显然两者中任何一个都不是另一个的倍数。
        \partsitem 这个集合是线性无关的。
           线性方程组按如下方式化简
           \begin{equation*}
             \begin{amat}[r]{3}
               2  &3  &4  &0  \\
               1  &-1 &0  &0  \\
               7  &2  &-3 &0
             \end{amat}
             \grstep[-(7/2)\rho_1+\rho_3]{-(1/2)\rho_1+\rho_2}
             \repeatedgrstep{-(17/5)\rho_2+\rho_3}
             \begin{amat}[r]{3}
               2  &3    &4      &0  \\
               0  &-5/2 &-2     &0  \\
               0  &0    &-51/5  &0
             \end{amat}
           \end{equation*}
           由此看出只有解 $c_1=0$、$c_2=0$ 以及 $c_3=0$。
        \partsitem 这个集合是线性相关的。
           线性方程组
           \begin{equation*}
             \begin{amat}[r]{4}
               8  &0  &2  &8  &0  \\
               3  &1  &2  &-2 &0  \\
               3  &2  &2  &5  &0
             \end{amat}
           \end{equation*}
           在化简后必定最后至少有一个自由变量（变量的个数多于方程的
           个数，而且由于该方程组是齐次的，不可能出现矛盾方程）。
           我们可以把自由变量取作参数来描述解集。然后令参数取一个
           非零值，就得到一个非平凡的线性关系。
       \end{exparts}
      
\end{ans}
\begin{ans}{二.II.1.23}
      \begin{exparts}
        \partsitem   自然向量空间是 $\Re^3$。
          建立方程
          \begin{equation*}
            \colvec{0 \\ 0 \\ 0}=c_1\colvec{1 \\ 2 \\ 0}
                                 +c_2\colvec{-1 \\ 1 \\ 0}
          \end{equation*}
          并考虑由此得到的齐次方程组。
          \begin{equation*}
             \begin{amat}{2}
               1  &-1 &0  \\
               2  &1  &0  \\
               0  &0  &0
             \end{amat}
             \grstep{-2\rho_1+\rho_2}
             \begin{amat}{2}
              1  &-1 &0  \\
              0  &3  &0  \\
              0  &0  &0
            \end{amat}
          \end{equation*}
          它有唯一解，即 $c_1=0$、$c_2=0$。
          所以它是线性无关的。
        \partsitem   自然向量空间是三个分量的行向量全体。
          方程
          \begin{equation*}
            \rowvec{0 &0 &0}=c_1\rowvec{1 &3 &1}
                            +c_2\rowvec{-1 &4 &3}
                            +c_3\rowvec{-1 &11 &7}
          \end{equation*}
          引出线性方程组
          \begin{align*}
            \begin{amat}{3}
              1  &-1  &-1  &0  \\
              3  &4   &11  &0  \\
              1  &3   &7   &0
           \end{amat}
           &\grstep[-\rho_1+\rho_3]{-3\rho_1+\rho_2}
           \begin{amat}{3}
             1  &-1  &-1  &0  \\
             0  &7   &14  &0  \\
             0  &4   &8   &0
          \end{amat}                                   \\
          &\grstep{(4/7)\rho_2+\rho_3}
          \begin{amat}{3}
            1  &-1  &-1  &0  \\
            0  &7   &14  &0  \\
            0  &0   &0   &0
         \end{amat}
       \end{align*}
       它有无穷多解，也就是说，不只是平凡解。
       \begin{equation*}
         \set{\colvec{c_1 \\ c_2 \\ c_3}
              =\colvec{-1 \\ -2 \\ 1}c_3
              \suchthat c_3\in\Re}
       \end{equation*}
       所以这个集合是线性相关的。
       一个相关关系来自取 $c_3=2$，此时得到 $c_1=-2$ 和 $c_2=-4$。
      \partsitem   不必建立方程组我们就能看出，集合中第二个元素是
        第一个元素的倍数（即第一个元素的 $0$ 倍）。
      \end{exparts}
    
\end{ans}
\begin{ans}{二.II.1.24}
      设 $\map{Z}{\Re^+}{\Re}$ 为零函数 $Z(x)=0$，它是所讨论的向量空间中
      的加法单位元。
      \begin{exparts}
        \partsitem 这个集合是线性无关的。
          考虑 \( c_1\cdot f(x)+c_2\cdot g(x)=Z(x) \)。
          代入 \( x=1 \) 和 \( x=2 \) 就得到线性方程组
          \begin{equation*}
            \begin{linsys}{2}
              c_1\cdot 1  &+  &c_2\cdot 1     &=  &0  \\
              c_1\cdot 2  &+  &c_2\cdot (1/2) &=  &0
            \end{linsys}
          \end{equation*}
          它有唯一解 \( c_1=0 \)、\( c_2=0 \)。
        \partsitem 这个集合是线性无关的。
          考虑 \( c_1\cdot f(x)+c_2\cdot g(x)=Z(x) \)，
          代入 \( x=\pi \) 和 \( x=\pi/2 \) 得到
          \begin{equation*}
            \begin{linsys}{2}
              c_1\cdot (-1)  &+  &c_2\cdot 0     &=  &0  \\
              c_1\cdot 0     &+  &c_2\cdot 1     &=  &0
            \end{linsys}
          \end{equation*}
          它显然给出 \( c_1=0 \)、\( c_2=0 \)。
        \partsitem 这个集合也是线性无关的。
          考虑 \( c_1\cdot f(x)+c_2\cdot g(x)=Z(x) \)
          并代入 \( x=1 \) 和 \( x=e \)
          \begin{equation*}
            \begin{linsys}{2}
              c_1\cdot e    &+  &c_2\cdot 0     &=  &0  \\
              c_1\cdot e^e  &+  &c_2\cdot 1     &=  &0
            \end{linsys}
          \end{equation*}
          得到 \( c_1=0 \) 与 \( c_2=0 \)。
      \end{exparts}
     
\end{ans}
\begin{ans}{二.II.1.25}
      在每一种情形，如果集合是线性无关的，那么你必须证明这一点；如果它是
      线性相关的，那么你必须展示出一个相关关系。
      \begin{exparts}
        \partsitem 这个集合是线性相关的。
          熟知的关系式 $\sin^2(x)+\cos^2(x)=1$ 表明，
          $2=c_1\cdot(4\sin^2(x))+c_2\cdot(\cos^2(x))$ 在
          $c_1=1/2$ 和 $c_2=2$ 时成立。
        \partsitem 这个集合是线性无关的。
          考虑关系式
          $c_1\cdot 1+c_2\cdot\sin(x)+c_3\cdot\sin(2x)=0$
          （其中 `$0$' 指零函数）。
          取三个适当的点，比如 $x=\pi$、$x=\pi/2$、$x=\pi/4$，
          就得到方程组
          \begin{equation*}
            \begin{linsys}{3}
               c_1  &   &                &   &      &=  &0  \\
               c_1  &+  &c_2             &   &      &=  &0  \\
               c_1  &+  &(\sqrt{2}/2)c_2 &+  &c_3   &=  &0
            \end{linsys}
          \end{equation*}
          其唯一解是
          $c_1=0$、$c_2=0$ 以及 $c_3=0$。
        \partsitem 凭观察可知，这个集合是线性无关的。
          任何形如 $\cos(x)=c\cdot x$ 的相关关系都不可能成立，因为余弦
          函数不是恒等函数的倍数
          （我们在这里应用 \nearbycorollary{cor:LDMeansLC}）。
        \partsitem 凭观察我们发现存在一个相关关系。
          因为 $(1+x)^2=x^2+2x+1$，所以
          $c_1\cdot(1+x)^2+c_2\cdot(x^2+2x)=3$ 在
          $c_1=3$ 和 $c_2=-3$ 时成立。
        \partsitem 这个集合是线性相关的。
          看出这一点最简单的方法是想起三角关系式
          $\cos^2(x)-\sin^2(x)=\cos(2x)$。
          （\textit{注}：
          想不起来这个关系式的人，若试着代入一些 $x$ 值，只会始终得不到
          通向唯一解的方程组，也就永远无法得出该集合线性无关的结论。
          当然，这个人可能会怀疑是不是恰好没试对那组 $x$ 值，但试过
          几次之后，大多数人就会转而去寻找相关关系。）
        \partsitem 这个集合是线性相关的，因为它包含该向量空间中的
          零对象，即零多项式。
      \end{exparts}
     
\end{ans}
\begin{ans}{二.II.1.26}
      不是，那个方程并不是一个线性关系。
      事实上这个集合是线性无关的，取 \( x \) 为 \( 0 \)、\( \pi/6 \)
      和 \( \pi/4 \) 所得到的方程组即可说明这一点。
    
\end{ans}
\begin{ans}{二.II.1.27}
      不是。
      该平面中有两个元素，第二个是第一个的倍数。
      \begin{equation*}
        \colvec{1  \\ 0 \\ 0},\,\colvec{2 \\ 0 \\ 0}
      \end{equation*}
      （答案为“不是”的另一个理由是，零向量是该平面的元素，而包含零
      向量的集合都不可能是线性无关的。）
    
\end{ans}
\begin{ans}{二.II.1.28}
      我们已经证明过这一点：线性组合引理及其推论指出，在阶梯形矩阵中，
      任何非零行都不是其余各行的线性组合。
    
\end{ans}
\begin{ans}{二.II.1.29}
       \begin{exparts}
         \partsitem 假设
           \( \set{\vec{u},\vec{v},\vec{w}} \) 是
           线性无关的，于是任何关系式
           $d_0\vec{u}+d_1\vec{v}+d_2\vec{w}=\zero$ 都导致
           $d_0=0$、$d_1=0$ 以及 $d_2=0$ 的结论。

           考虑关系式
           \( c_1(\vec{u})+c_2(\vec{u}+\vec{v})+c_3(\vec{u}+\vec{v}+\vec{w})
           =\zero \)。
           将其改写，得到
           \( (c_1+c_2+c_3)\vec{u}+(c_2+c_3)\vec{v}+(c_3)\vec{w}=\zero \)。
           取 $d_0$ 为 $c_1+c_2+c_3$，取 $d_1$ 为 $c_2+c_3$，
           取 $d_2$ 为 $c_3$，就得到这个方程组。
           \begin{equation*}
             \begin{linsys}{3}
               c_1  &+  &c_2  &+  &c_3  &=  &0  \\
                    &   &c_2  &+  &c_3  &=  &0  \\
                    &   &     &   &c_3  &=  &0
              \end{linsys}
            \end{equation*}
            结论：~各 $c$ 全为零，所以该集合是线性
            无关的。
         \partsitem 第二个集合是线性相关的
            \begin{equation*}
              1\cdot(\vec{u}-\vec{v})
              +1\cdot(\vec{v}-\vec{w})
              +1\cdot(\vec{w}-\vec{u})
              =\zero
            \end{equation*}
            无论第一个集合是否线性无关。
       \end{exparts}
     
\end{ans}
\begin{ans}{二.II.1.30}
      \begin{exparts}
        \partsitem 单元素集 $\set{\vec{v}}$ 是线性无关的
          当且仅当 $\vec{v}\neq\zero$。
          对于“若”的方向，当 $\vec{v}\neq\zero$ 时，考虑
          关系式
          \( c\cdot\vec{v}=\zero \) 并注意唯一解
          是平凡的：~$c=0$，就可以应用
          \nearbylemma{le:LDIffANonTrivLinRel}。
          对于“仅当”的方向，只需回忆
          \nearbyexample{ex:SetWithZeroVecLD}
          表明 $\set{\zero}$ 是线性相关的，因此如果集合
          $\set{\vec{v}}$ 是线性无关的，那么 $\vec{v}\neq\zero$。

          （\textit{注}：
          另一种回答是说，这是
          \nearbylemma{lm:AddVecLiSetIsLiIffVecNotInSpan}
          在 \( S=\emptyset \) 时的特殊情形。）
        \partsitem 两个元素的集合是线性无关的
          当且仅当其中任何一个元素都不是另一个的倍数
          （注意，如果其中一个是零向量，那么它是另一个的倍数）。
          等价的说法：~一个集合是线性相关的当且仅当
          其中一个元素是另一个的倍数。

          证明很容易。
          集合 $\set{\vec{v}_1,\vec{v}_2}$ 是线性相关的当且仅当
          存在关系式 $c_1\vec{v}_1+c_2\vec{v}_2=\zero$，
          其中 $c_1\neq 0$ 或 $c_2\neq 0$（或两者都成立）。
          这当且仅当 $\vec{v}_1=(-c_2/c_1)\vec{v}_2$
          或 $\vec{v}_2=(-c_1/c_2)\vec{v}_1$（或两者都成立）时成立。
       \end{exparts}
     
\end{ans}
\begin{ans}{二.II.1.31}
      这个集合是线性相关的，因为它包含零向量。
    
\end{ans}
\begin{ans}{二.II.1.32}
      \nearbylemma{le:SubsetPreserveLI} 给出了“若”的那一半。
      逆命题（“仅当”的陈述）不成立。
      举个例子：考虑向量空间 \( \Re^2 \) 和
      这些向量。
      \begin{equation*}
         \vec{x}=\colvec{1 \\ 0},\quad
         \vec{y}=\colvec{0 \\ 1},\quad
         \vec{z}=\colvec{1 \\ 1}
      \end{equation*}
    
\end{ans}
\begin{ans}{二.II.1.33}
      \begin{exparts}
        \partsitem 由
          \begin{equation*}
            c_1\colvec{1 \\ 1 \\ 0}
            +c_2\colvec{-1 \\ 2 \\ 0}
            =\colvec{0 \\ 0 \\ 0}
          \end{equation*}
          得到的线性方程组有唯一解 \( c_1=0 \) 和 \( c_2=0 \)。
        \partsitem 由
          \begin{equation*}
            c_1\colvec{1 \\ 1 \\ 0}
            +c_2\colvec{-1 \\ 2 \\ 0}
            =\colvec{3 \\ 2 \\ 0}
          \end{equation*}
          得到的线性方程组有唯一解 \( c_1=8/3 \) 和 \( c_2=-1/3 \)。
        \partsitem 假设 \( S \) 是线性无关的。
          再假设我们同时有 $\vec{v}=c_1\vec{s}_1+\dots+c_n\vec{s}_n$
          和 $\vec{v}=d_1\vec{t}_1+\dots+d_m\vec{t}_m$
          （其中这些向量都是 $S$ 的元素）。
          于是，
          \begin{equation*}
            c_1\vec{s}_1+\dots+c_n\vec{s}_n
            =\vec{v}
            =d_1\vec{t}_1+\dots+d_m\vec{t}_m
          \end{equation*}
          可以如下改写。
          \begin{equation*}
            c_1\vec{s}_1+\dots+c_n\vec{s}_n
            -d_1\vec{t}_1-\dots-d_m\vec{t}_m
            =\zero
          \end{equation*}
          可能有些 $\vec{s}$ 与某些 $\vec{t}$ 相等；
          我们可以把相应的系数合并
          （即，如果 $\vec{s}_i=\vec{t}_j$，那么
          $\cdots+c_i\vec{s}_i+\dots-d_j\vec{t}_j-\cdots$ 可以改写
          为 $\cdots+(c_i-d_j)\vec{s}_i+\cdots$）。
          这个等式是集合 $S$ 的（在合并完成之后的）不同元素之间的
          一个线性关系。
          我们已假设 $S$ 是线性无关的，所以所有系数都为零。
          如果 $i$ 使得 $\vec{s}_i$ 不等于任何 $\vec{t}_j$，
          那么 $c_i$ 为零。
          如果 $j$ 使得 $\vec{t}_j$ 不等于任何 $\vec{s}_i$，
          那么 $d_j$ 为零。
          在最后一种情形，我们有 $c_i-d_j=0$，从而 $c_i=d_j$。

          因此，原来的两个和是相同的，也许除去一些
          $0\cdot\vec{s}_i$ 或 $0\cdot\vec{t}_j$ 项，这些项可以忽略。
        \partsitem
          这个集合不是线性无关的：
          \begin{equation*}
            S=\set{\colvec{1 \\ 0},\colvec{2 \\ 0}}\subset\Re^2
          \end{equation*}
          而这两个线性组合给出同样的结果
          \begin{equation*}
            \colvec{0 \\ 0}=2\cdot\colvec{1 \\ 0}-1\cdot\colvec{2 \\ 0}
                            =4\cdot\colvec{1 \\ 0}-2\cdot\colvec{2 \\ 0}
          \end{equation*}
          因此，一个线性相关的集合的和可能不是互不相同的。

          事实上，一个更强的论断成立：~如果一个集合是线性相关的，那么它
          必定具有这样的性质：存在两个不同的线性组合求和得到同一个向量。
          简言之，若有 \( c_1\vec{s}_1+\dots+c_n\vec{s}_n=\zero \)，
          则将关系式两边乘以二就得到另一个关系式。
          如果第一个关系式是非平凡的，那么第二个也是。
      \end{exparts}
    
\end{ans}
\begin{ans}{二.II.1.34}
       在这个“当且仅当”的陈述中，“若”的那一半是清楚的\Dash 如果
       多项式是零多项式，那么由该多项式的作用导出的函数必定是零
       函数 $x\mapsto 0$。
       对于“仅当”，我们写 $p(x)=c_nx^n+\dots+c_0$。
       代入零，$p(0)=0$，得到 $c_0=0$。
       取导数再代入零，$p^\prime(0)=0$，得到
       $c_1=0$。
       类似地，我们得到每个 $c_i$ 都为零，于是 $p$ 是零
       多项式。
     
\end{ans}
\begin{ans}{二.II.1.35}
      本节的工作提示我们，应当把 \( n \) 维
      非退化线性曲面定义为 \( n \) 个向量组成的线性无关集合的
      张成。
    
\end{ans}
\begin{ans}{二.II.1.36}
      \begin{exparts}
        \partsitem 对任意的 $a_{1,1}$、\ldots、$a_{2,4}$，
          \begin{equation*}
            c_1\colvec{a_{1,1} \\ a_{2,1}}
            +c_2\colvec{a_{1,2} \\ a_{2,2}}
            +c_3\colvec{a_{1,3} \\ a_{2,3}}
            +c_4\colvec{a_{1,4} \\ a_{2,4}}
            =\colvec{0 \\ 0}
          \end{equation*}
          给出线性方程组
          \begin{equation*}
             \begin{linsys}{4}
              a_{1,1}c_1 &+ &a_{1,2}c_2 &+ &a_{1,3}c_3 &+ &a_{1,4}c_4 &= &0  \\
              a_{2,1}c_1 &+ &a_{2,2}c_2 &+ &a_{2,3}c_3 &+ &a_{2,4}c_4 &= &0
             \end{linsys}
          \end{equation*}
        它有无穷多解（高斯消元法至少留下
        两个自由变量）。
        因此，$\Re^2$ 中给定的这些元素之间存在非平凡的线性关系。
      \partsitem 任意五个向量的集合是某个四向量集合的超集，
        所以是线性相关的。

        对于来自 $\Re^2$ 的三个向量，上一小项的论证仍然适用，只是略有变化：高斯消元法现在只留下至少一个自由变量（但这仍然给出无穷多解）。
      \partsitem 上一小项表明，$\Re^2$ 中任何三元素子集都不是线性无关的。
        我们知道，$\Re^2$ 中存在线性无关的二元素子集\Dash 例如
        \begin{equation*}
          \set{\colvec{1  \\ 0},\colvec{0  \\ 1}}
        \end{equation*}
        因此答案是二。
     \end{exparts}
    
\end{ans}
\begin{ans}{二.II.1.37}
      存在；下面就是一例。
      \begin{equation*}
         \set{\colvec{1 \\ 0 \\ 0},
              \colvec{0 \\ 1 \\ 0},
              \colvec{0 \\ 0 \\ 1},
              \colvec{1 \\ 1 \\ 1} }
      \end{equation*}
    
\end{ans}
\begin{ans}{二.II.1.38}
      是的。
      任取一个向量空间~$V$，并设 $S$ 是~$V$ 的一个线性相关子集。
      至于~$S$ 的线性相关子集，取~$S$ 本身即可。
      至于~$S$ 的线性无关子集，可以取空集。
    
\end{ans}
\begin{ans}{二.II.1.39}
      在 \( \Re^4 \) 中，最大的线性无关集合含有四个向量。
      这样的集合有很多例子，下面是一个。
      \begin{equation*}
        \set{\colvec{1 \\ 0 \\ 0 \\ 0},
             \colvec{0 \\ 1 \\ 0 \\ 0},
             \colvec{0 \\ 0 \\ 1 \\ 0},
             \colvec{0 \\ 0 \\ 0 \\ 1}  }
      \end{equation*}
      为了看出任何含五个或更多向量的集合都不可能线性无关，列出
      \begin{equation*}
             c_1\colvec{a_{1,1} \\ a_{2,1} \\ a_{3,1} \\ a_{4,1}}
            +c_2\colvec{a_{1,2} \\ a_{2,2} \\ a_{3,2} \\ a_{4,2}}
            +c_3\colvec{a_{1,3} \\ a_{2,3} \\ a_{3,3} \\ a_{4,3}}
            +c_4\colvec{a_{1,4} \\ a_{2,4} \\ a_{3,4} \\ a_{4,4}}
            +c_5\colvec{a_{1,5} \\ a_{2,5} \\ a_{3,5} \\ a_{4,5}}
             =\colvec{0 \\ 0 \\ 0 \\ 0}
      \end{equation*}
      并注意由此得到的线性方程组
      \begin{equation*}
        \begin{linsys}{5}
          a_{1,1}c_1 &+ &a_{1,2}c_2 &+ &a_{1,3}c_3
              &+ &a_{1,4}c_4 &+ &a_{1,5}c_5 &=  &0   \\
          a_{2,1}c_1 &+ &a_{2,2}c_2 &+ &a_{2,3}c_3
              &+ &a_{2,4}c_4 &+ &a_{2,5}c_5 &=  &0   \\
          a_{3,1}c_1 &+ &a_{3,2}c_2 &+ &a_{3,3}c_3
              &+ &a_{3,4}c_4 &+ &a_{3,5}c_5 &=  &0   \\
          a_{4,1}c_1 &+ &a_{4,2}c_2 &+ &a_{4,3}c_3
              &+ &a_{4,4}c_4 &+ &a_{4,5}c_5 &=  &0
        \end{linsys}
      \end{equation*}
      有四个方程与五个未知数，
      所以高斯消元法结束时至少有一个 \( c \) 变量是自由变量，
      故有无穷多解，
      从而上述四个分量的向量之间的线性关系，其解不只是平凡解。

      最小的线性无关集合是空集。

      最大的线性相关集合是 \( \Re^4 \)。
      最小的是 \( \set{\zero} \)。
    
\end{ans}
\begin{ans}{二.II.1.40}
      \begin{exparts}
        \partsitem 两个线性无关集合的交 $S\intersection T$ 必定线性无关，
          因为它是线性无关集合 $S$ 的子集
          （当然，也是线性无关集合 $T$ 的子集）。
        \partsitem 线性无关集合的补是线性相关的，
          因为它包含零向量。
        \partsitem \( \Re^2 \) 中的一个简单例子是下面这两个集合。
          \begin{equation*}
            S=\set{\colvec{1 \\ 0}}
            \qquad
            T=\set{\colvec{0 \\ 1}}
          \end{equation*}
          一个稍微精细些的例子，同样取自 \( \Re^2 \)，是下面这两个。
          \begin{equation*}
            S=\set{\colvec{1 \\ 0}}
            \qquad
            T=\set{\colvec{1 \\ 0}, \colvec{0 \\ 1}}
          \end{equation*}
        \partsitem 我们必须构造一个例子。
          \( \Re^2 \) 中的一个例子是
          \begin{equation*}
            S=\set{\colvec{1 \\ 0}}
            \qquad
            T=\set{\colvec{2 \\ 0}}
          \end{equation*}
          因为 \( S_1\union S_2 \) 的线性相关性显而易见。
      \end{exparts}
    
\end{ans}
\begin{ans}{二.II.1.41}
      \begin{exparts}
        \partsitem \nearbylemma{le:LDIffANonTrivLinRel}
          要求诸向量
          $\vec{s}_1,\ldots,\vec{s}_n,\vec{t}_1,\ldots,\vec{t}_m$
          互不相同。
          但可能出现这种情况：并 $S\cup T$ 线性无关，而某个 $\vec{s}_i$ 等于某个 $\vec{t}_j$。
        \partsitem \( \Re^2 \) 中的一个例子是下面这两个。
          \begin{equation*}
            S=\set{\colvec{1 \\ 0}}
            \qquad
            T=\set{\colvec{1 \\ 0}, \colvec{0 \\ 1}}
          \end{equation*}
        \partsitem
          取自 \( \Re^2 \) 的一个例子是这些集合。
          \begin{equation*}
            S=\set{\colvec{1 \\ 0}, \colvec{0 \\ 1}}
            \qquad
            T=\set{\colvec{1 \\ 0}, \colvec{1 \\ 1}}
          \end{equation*}
        \partsitem 两个线性无关集合的并 $S\union T$
          线性无关，当且仅当 \( S \) 的张成
          与 \( T-(S\cap T)\) 的张成之交是平凡的，
          即 $\spanof{S}\intersection \spanof{T-(S\cap T)}=\set{\zero}$。
          为证明这一点，设 \( S \) 与 \( T \) 是某个向量空间中的线性
          无关子集。

          对于“仅当”方向，设张成之交是平凡的
          \( \spanof{S}\intersection \spanof{T-(S\cap T)}=\set{\zero} \)。
          考虑集合 $S\union (T-(S\cap T))=S\union T$，
          并考虑线性关系
          $c_1\vec{s}_1+\dots+c_n\vec{s}_n
            +d_1\vec{t}_1+\dots+d_m\vec{t}_m=\zero$。
          移项得
          $c_1\vec{s}_1+\dots+c_n\vec{s}_n=
            -d_1\vec{t}_1-\dots-d_m\vec{t}_m$。
          该等式左边求和得到 $\spanof{S}$ 中的一个向量，
          而右边是 $\spanof{T-(S\cap T)}$ 中的一个向量。
          因此，由于张成之交是平凡的，两边都等于零向量。
          因为 $S$ 线性无关，所以所有 $c$ 都是零。
          因为 $T$ 线性无关，所以 $T-(S\cap T)$ 也线性无关，
          从而所有 $d$ 都是零。
          于是，$S\union T$ 中各向量之间的这个线性关系
          只有在所有系数都为零时才成立。
          因此，$S\union T$ 线性无关。

          对于“若”这一半，我们可以把同样的论证反着进行。
          设并 $S\union T$ 线性无关。
          考虑 $S$ 与 $T-(S\cap T)$ 中向量之间的一个线性关系。
          $c_1\vec{s}_1+\cdots+c_n\vec{s}_n
            +d_1\vec{t}_1+\cdots+d_m\vec{t}_m
            =\zero$
          注意没有哪个 $\vec{s}_i$ 等于某个 $\vec{t}_j$，
          所以这是互不相同的向量的组合，正如
          \nearbylemma{le:LDIffANonTrivLinRel}
          所要求的那样。
          于是唯一解
          是平凡解 $c_1=0$，\ldots，$d_m=0$。
          由于张成交集
          $\spanof{S}\intersection \spanof{T-(S\cap T)}$
          中的任意向量 $\vec{v}$
          都可以写成
          $\vec{v}=c_1\vec{s}_1+\cdots+c_n\vec{s}_n
            =-d_1\vec{t}_1-\cdots-d_m\vec{t}_m$，
          而每个标量都是零，故它必是零向量。
      \end{exparts}
    
\end{ans}
\begin{ans}{二.II.1.42}
      \begin{exparts}
         \partsitem 我们对有限集
           \( S \) 中向量的个数作归纳。

           基础情形是 $S$ 没有元素。
           此时 $S$ 线性无关，无需检验任何东西\Dash
           与 $S$ 张成相同的 $S$ 的子集就是 $S$
           本身。

           归纳步骤：假设该定理对大小为 $n=0$、$n=1$、\ldots、$n=k$
           的所有集合都成立，以此证明当 \( S \) 有 $n=k+1$ 个元素时它也成立。
           若 $k+1$ 元集合 \( S=\set{\vec{s}_0,\dots,\vec{s}_{k}} \)
           线性无关，则定理平凡成立，
           故设它是线性相关的。
           由 \nearbycorollary{cor:LDMeansLC}，存在某个 \( \vec{s}_i \)，
           它是 \( S \) 中其他向量的线性组合。
           令 \( S_1=S-\set{\vec{s}_i} \)，并注意
           由 \nearbycorollary{cor:OmitVecNotShrinkSpanIffVecIsDep}，
           \( S_1 \) 与 \( S \) 有相同的张成。
           集合 \( S_1 \) 有 \( k \) 个元素，
           故归纳假设适用，于是它有一个张成相同的线性无关子集。
           \( S_1 \) 的这个子集就是所要求的 \( S \) 的子集。
         \partsitem 下面是论证的概要。
           我们略去了归纳论证的细节。

           若有限集 \( S \) 为空，则无需证明。
           若 \( S=\set{\zero} \)，则取空子集即可。

           否则，取某个非零向量 \( \vec{s}_1\in S \)，
           并令 \( S_1=\set{\vec{s}_1} \)。
           若 \( \spanof{S_1}=\spanof{S} \)，则
           注意到 \( S_1 \) 线性无关，本证明即告完成。

           若不然，则存在非零
           向量 \( \vec{s}_2\in S-\spanof{S_1} \)
           （因为若每个 \( \vec{s}\in S \) 都属于 \( \spanof{S_1} \)，则
           \( \spanof{S_1}=\spanof{S} \)）。
           令 \( S_2=S_1\union\set{\vec{s}_2} \)。
           若 \( \spanof{S_2}=\spanof{S} \)，则
           可利用 \nearbytheorem{cor:LDMeansLC}
           证明 \( S_2 \) 线性无关，证明即告完成。

           将上一段重复进行，直到出现一个张成足够大的集合。
           这最终必定发生，因为 \( S \) 是有限的，并且
           最迟在我们用尽
           \( S \) 中的全部向量时就会达到 \( \spanof{S} \)。
      \end{exparts}
    
\end{ans}
\begin{ans}{二.II.1.43}
      \begin{exparts}
        \partsitem
          先设 \( a\neq 0 \)，
          \begin{equation*}
            x\colvec{a \\ c}
            +y\colvec{b \\ d}
            =\colvec{0 \\ 0}
          \end{equation*}
          给出
          \begin{equation*}
            \begin{linsys}{2}
               ax  &+  &by &=  &0  \\
               cx  &+  &dy &=  &0
            \end{linsys}
            \grstep{-(c/a)\rho_1+\rho_2}
            \begin{linsys}{2}
               ax  &+  &by           &=  &0  \\
                   &   &(-(c/a)b+d)y &=  &0
             \end{linsys}
          \end{equation*}
          它有解当且仅当
          \( 0\neq-(c/a)b+d=(-cb+ad)/d \)
          （这一情形我们已设 \( a\neq 0 \)，
          故回代给出唯一解）。

          \( a=0 \) 的情形也不难\Dash
          把它分成 \( c\neq 0 \) 与 \( c=0 \) 两种子情形，
          并注意在这些情形 \( ad-bc=0\cdot d-bc \)。

          \textit{评注。}
          前面的一道练习已证明：由两个向量组成的集合线性相关，
          当且仅当其中一个向量是另一个的标量倍。
          我们也可以利用那个结论来进行这里的计算。
        \partsitem 方程
          \begin{equation*}
            c_1\colvec{a \\ d \\ g}
            +c_2\colvec{b \\ e \\ h}
            +c_3\colvec{c \\ f \\ i}
            =\colvec{0 \\ 0 \\ 0}
          \end{equation*}
         表示一个齐次线性方程组。
         我们将它写成矩阵形式并应用高斯消元法。

         我们先把矩阵化为上三角形式。
         设 \( a\neq 0 \)。
         有了这一点，我们就可以对第一列向下消元。
         \begin{equation*}
           \grstep{(1/a)\rho_1}
           \begin{amat}{3}
              1   &b/a   &c/a  &0 \\
              d   &e     &f    &0 \\
              g   &h     &i    &0
            \end{amat}
           \grstep[-g\rho_1+\rho_3]{-d\rho_1+\rho_2}
           \begin{amat}{3}
              1   &b/a           &c/a        &0   \\
              0   &(ae-bd)/a     &(af-cd)/a  &0   \\
              0   &(ah-bg)/a     &(ai-cg)/a  &0
            \end{amat}
         \end{equation*}
         然后我们把第二行、第二列处的元素变成 $1$。
         （为执行这一行化简步骤，暂时设 \( ae-bd\neq 0 \)。）
         \begin{equation*}
           \grstep{(a/(ae-bd))\rho_2}
           \begin{amat}{3}
              1   &b/a           &c/a             &0  \\
              0   &1             &(af-cd)/(ae-bd) &0  \\
              0   &(ah-bg)/a     &(ai-cg)/a       &0
            \end{amat}
         \end{equation*}
         然后，在上述假设下，我们施行行变换 $((ah-bg)/a)\rho_2+\rho_3$，
         得到下式。
         \begin{equation*}
           % \grstep{((ah-bg)/a)\rho_2+\rho_3}
           \begin{amat}{3}
              1   &b/a   &c/a                              &0 \\
              0   &1     &(af-cd)/(ae-bd)                   &0 \\
              0   &0     &(aei+bgf+cdh-hfa-idb-gec)/(ae-bd) &0
            \end{amat}
         \end{equation*}
         因此，原方程组是非奇异的
         当且仅当上面的 \( 3,3 \) 元非零
         （由于 \( ae-bd\neq 0 \) 的假设，这个分式有定义）。
         它等于零当且仅当分子为零。

         接下来我们处理这些假设。
         首先，若 \( a\neq 0 \) 但 \( ae-bd=0 \)，则我们作对换
         \begin{multline*}
           \begin{amat}{3}
              1   &b/a           &c/a        &0   \\
              0   &0             &(af-cd)/a  &0   \\
              0   &(ah-bg)/a     &(ai-cg)/a  &0
            \end{amat}                                    \\
           \grstep{\rho_2\leftrightarrow\rho_3}
           \begin{amat}{3}
              1   &b/a           &c/a        &0   \\
              0   &(ah-bg)/a     &(ai-cg)/a  &0   \\
              0   &0             &(af-cd)/a  &0
            \end{amat}
         \end{multline*}
         并得出结论：方程组非奇异当且仅当
         \( ah-bg=0 \) 或 \( af-cd=0 \)。
         这等价于问它们的乘积是否为零：
         \begin{align*}
            ahaf-ahcd-bgaf+bgcd
            &=0                   \\
            ahaf-ahcd-bgaf+aegc
            &=0                   \\
            a(haf-hcd-bgf+egc)
            &=0
         \end{align*}
         （从第一行到第二行，我们利用了情形假设 $ae-bd=0$，
         用 $ae$ 代替 $bd$）。
         由于我们已设 \( a\neq 0 \)，
         我们便有 \( haf-hcd-bgf+egc=0 \)。
         利用 $ae-bd=0$，我们可以把它改写成所需的形式：~
         在这个 \( a\neq 0 \) 且 \( ae-bd=0 \) 的情形中，
         当
         \( haf-hcd-bgf+egc-i(ae-bd)=0 \) 时，
         给定方程组非奇异，正如所要求的。

         其余的情形具有同样的特点。
         对于 \( a=0 \) 但 \( d\neq 0 \) 的情形，以及 \( a=0 \) 且
         \( d=0 \) 但 \( g\neq 0 \) 的情形，
         先对换行，然后照上面那样继续进行。
         \( a=0 \)、\( d=0 \)、\( g=0 \) 的情形很简单\Dash
         含零向量的集合线性相关，而公式算出来恰好为零。
       \partsitem 它线性相关当且仅当其中一个向量是另一个的倍数。
          也就是说，它不独立当且仅当
          \begin{equation*}
            \colvec{a \\ d \\ g}=r\cdot\colvec{b \\ e \\ h}
            \quad\text{或}\quad
            \colvec{b \\ e \\ h}=s\cdot\colvec{a \\ d \\ g}
          \end{equation*}
          （或两者同时成立），其中 $r$ 与 $s$ 是标量。
          消去 $r$ 与 $s$，以便只用给定的字母 $a$、$b$、$d$、$e$、$g$、$h$
          来重新陈述这个条件，我们得到：它不独立\Dash
          即线性相关\Dash 当且仅当
          \( ae-bd=ah-gb=dh-ge \)。
        \partsitem 相关还是无关是下标的函数，所以确实存在一个公式
          （尽管乍看之下人们可能以为该公式要分情况讨论：
          “如果第一个向量的第一个分量为零，那么\ldots”，
          但这个猜测并不正确）。
       \end{exparts}
     
\end{ans}
\begin{ans}{二.II.1.44}
      回想一下，\( \Re^n \) 中的两个向量垂直当且仅当它们的点积为零。
      \begin{exparts}
         \partsitem 设 \( \vec{v} \) 与 \( \vec{w} \) 是
           $\Re^n$ 中互相垂直的非零向量，其中 \( n>1 \)。
           对于线性关系 \( c\vec{v}+d\vec{w}=\zero \)，
           在两边点乘 \( \vec{v} \)，得到 \( c\cdot\norm{\vec{v}}^2+d\cdot 0=0 \)。
           因为 \( \vec{v}\neq\zero \)，所以 \( c=0 \)。
           类似地用 \( \vec{w} \) 点乘可知 \( d=0 \)。
         \partsitem \( \Re^1 \) 中的两个向量垂直当且仅当
           其中至少有一个是零向量。

           我们把 \( \Re^0 \) 定义为平凡空间，
           因此 $\vec{v}$ 与 $\vec{w}$ 都是零向量。
         \partsitem 正确的推广是考察由
           \definend{两两正交}
           （也称为 \definend{两两垂直}）的向量组成的集合
           \( \set{\vec{v}_1,\dots,\vec{v}_n}\subseteq\Re^k \)：~若
           \( i\neq j \)，则 \( \vec{v}_i \) 垂直于
           \( \vec{v}_j \)。
           模仿上面第一小项的证明可知，这样的非零向量集合是线性无关的。
      \end{exparts}
    
\end{ans}
\begin{ans}{二.II.1.45}
      \begin{exparts}
        \partsitem 这个检验是常规的。
        \partsitem 该求和是无穷的（有无穷多项）。
          而线性组合的定义只涉及有限和。
        \partsitem \( \set{g,f_0,f_1,\ldots} \) 中成员的任何非平凡有限和
           都不会加成零对象：~设
           \begin{equation*}
              c_0\cdot (1/(1-x))+c_1\cdot 1+\dots+c_n\cdot x^n=0
           \end{equation*}
           （任何有限和都有一个最高次幂，这里是 \( n \)）。
           两边乘以 \( 1-x \) 可知每个系数都为零，
           因为仅当多项式是零多项式时，它才描述零函数。
      \end{exparts}
     
\end{ans}
\begin{ans}{二.II.1.46}
      “若”与“仅当”两个方向都成立。

      设 \( T \) 是向量空间 \( V \) 的子空间 \( S \) 的一个子集。
      断言“\( T \) 中成员之间的任何线性关系
      $c_1\vec{t}_1+\dots+c_n\vec{t}_n=\zero$
      必定是平凡关系 $c_1=0$，\ldots，$c_n=0$”
      这一命题在 \( S \) 中成立当且仅当它在 \( V \) 中成立，
      因为子空间 \( S \) 的加法与
      标量乘法运算继承自 \( V \)。
    
\end{ans}
\protect \section {第 III 节：基与维数}
\protect \subsection {二.III.1: 基}
\begin{ans}{二.III.1.20}
      \begin{exparts}
        \partsitem 这是 $\polyspace_2$ 的一个基。
           为证明它张成该空间，考虑一个一般的
           $a_2x^2+a_1x+a_0\in\polyspace_2$，寻找
           满足
           $a_2x^2+a_1x+a_0=c_1\cdot(x^2-x+1) +c_2\cdot(2x+1) +c_3(2x-1)$
           的标量 $c_1, c_2, c_3\in\Re$。
           由此得到的线性方程组如下。
            \begin{equation*}
              \begin{linsys}{3}
                c_1  &\spaceforemptycolumn  &     &  &     &=  &a_2 \\
                     &  &2c_2 &+ &2c_3 &=  &a_1 \\
                     &  &c_2  &- &c_3   &=  &a_0
              \end{linsys}
              \grstep{(-1/2)\rho_2+\rho_3}
              \begin{linsys}{3}
                c_1  &\spaceforemptycolumn  &     &  &     &=  &a_2\hfill\hbox{} \\
                     &  &2c_2 &+ &2c_3    &=  &a_1\hfill\hbox{} \\
                     &  &     &  &-2c_3   &=  &a_0-(1/2)a_1
              \end{linsys}
            \end{equation*}
           于是，给定任意一组 $a_i$，我们都能据此算出 $c_j$：
           $c_1=a_2$，$c_2=(1/4)a_1+(1/2)a_0$，以及 $c_3=(1/4)a_1-(1/2)a_0$。
           因此 $\polyspace_2$ 的每个元素都是给定的
           $x^2-x+1$、$2x+1$ 与~$2x-1$ 的组合。

           为证明所给三个元素构成的集合线性无关，可以建立方程
           $0x^2+0x+0=c_1\cdot(x^2-x+1) +c_2\cdot(2x+1) +c_3(2x-1)$
           并求解，它将给出 $c_1=0$、$c_2=0$、$c_3=0$。
           或者，我们也可以换一种方式，
           注意到上一段中的解是唯一的，
           并引用 \nearbytheorem{th:BasisIffUniqueRepWRT}。
        \partsitem 这不是基。
           它不张成该空间，因为 $c_1\cdot (x+x^2) +c_2\cdot (x-x^2)$
           这两个元素的任何组合都不能合成多项式
           $3\in\polyspace_2$。
      \end{exparts}
    
\end{ans}
\begin{ans}{二.III.1.21}
      由 \nearbytheorem{th:BasisIffUniqueRepWRT}，每一个是基当且仅当
      我们能以唯一的方式把空间中的每个向量表示为
      所给向量的线性组合。
      \begin{exparts}
        \partsitem 是的，这是一个基。
          关系式
          \begin{equation*}
            c_1\colvec[r]{1 \\ 2 \\ 3}
            +c_2\colvec[r]{3 \\ 2 \\ 1}
            +c_3\colvec[r]{0 \\ 0 \\ 1}
            =\colvec{x \\ y \\ z}
          \end{equation*}
          给出
          \begin{equation*}
            \begin{amat}{3}
              1  &3  &0  &x  \\
              2  &2  &0  &y  \\
              3  &1  &1  &z
            \end{amat}
            \grstep[-3\rho_1+\rho_3]{-2\rho_1+\rho_2}
            \repeatedgrstep{-2\rho_2+\rho_3}
            \begin{amat}{3}
              1  &3  &0  &x \hfill\hbox{}  \\
              0  &-4 &0  &-2x+y \hfill\hbox{}  \\
              0  &0  &1  &x-2y+z
            \end{amat}
          \end{equation*}
          它有唯一的解
          \( c_3=x-2y+z \)、\( c_2=x/2-y/4 \)
          以及 \( c_1=-x/2+3y/4 \)。
        \partsitem 这不是基。
          按照上一小题的方式建立关系式
          \begin{equation*}
            c_1\colvec[r]{1 \\ 2 \\ 3}
            +c_2\colvec[r]{3 \\ 2 \\ 1}
            =\colvec{x \\ y \\ z}
          \end{equation*}
          得到一个线性方程组，其求解过程为
          \begin{equation*}
            \begin{amat}{2}
              1  &3  &x  \\
              2  &2  &y  \\
              3  &1  &z
            \end{amat}
            \grstep[-3\rho_1+\rho_3]{-2\rho_1+\rho_2}
            \repeatedgrstep{-2\rho_2+\rho_3}
            \begin{amat}{2}
              1  &3  &x\hfill\hbox{}  \\
              0  &-4 &-2x+y\hfill\hbox{}  \\
              0  &0  &x-2y+z
            \end{amat}
          \end{equation*}
          当且仅当这个三元向量的分量
          $x$、$y$、$z$ 满足 $x-2y+z=0$ 时，求解才可能。
          例如，当 $x=1$、$y=1$、$z=1$ 时，
          我们能找到起作用的系数 $c_1$ 与 $c_2$。
          然而，对于 $x=1$、$y=1$、$z=2$，
          则不存在起作用的 $c$。
          因此这不是基；它不张成该空间。
        \partsitem 是的，这是一个基。
         建立关系式后得到如下化简
          \begin{equation*}
            \begin{amat}{3}
              0  &1  &2  &x \hfill\hbox{} \\
              2  &1  &5  &y \hfill\hbox{} \\
             -1  &1  &0  &z
            \end{amat}
            \grstep{\rho_1\swap\rho_3}
            \repeatedgrstep{2\rho_1+\rho_2}
            \repeatedgrstep{-(1/3)\rho_2+\rho_3}
            \begin{amat}{3}
             -1  &1  &0   &z  \hfill\hbox{} \\
              0  &3  &5   &y+2z \hfill\hbox{} \\
              0  &0  &1/3 &x-y/3-2z/3
            \end{amat}
          \end{equation*}
          对于每一组分量
          $x$、$y$、$z$，它都有唯一的解。
        \partsitem 不，这不是基。
          化简过程
          \begin{equation*}
            \begin{amat}{3}
              0  &1  &1  &x   \\
              2  &1  &3  &y   \\
             -1  &1  &0  &z
            \end{amat}
            \grstep{\rho_1\swap\rho_3}
            \repeatedgrstep{2\rho_1+\rho_2}
            \repeatedgrstep{(-1/3)\rho_2+\rho_3}
            \begin{amat}{3}
             -1  &1  &0  &z  \hfill\hbox{} \\
              0  &3  &3  &y+2z  \hfill\hbox{} \\
              0  &0  &0  &x-y/3-2z/3
            \end{amat}
          \end{equation*}
          它并非对每一组三元 $x$、$y$、$z$ 都有解。
          恰恰相反，所给集合的张成
          只包含那些满足 \( x=y/3+2z/3 \) 的三元向量。
      \end{exparts}
     
\end{ans}
\begin{ans}{二.III.1.22}
      \begin{exparts}
        \partsitem 我们求解
          \begin{equation*}
            c_1\colvec[r]{1 \\ 1}
            +c_2\colvec[r]{-1 \\ 1}
            =\colvec[r]{1 \\ 2}
          \end{equation*}
          其过程为
          \begin{equation*}
             \begin{amat}[r]{2}
               1  &-1  &1  \\
               1  &1   &2
             \end{amat}
             \grstep{-\rho_1+\rho_2}
             \begin{amat}[r]{2}
               1  &-1  &1  \\
               0  &2   &1
             \end{amat}
          \end{equation*}
          并由此得出 \( c_2=1/2 \)，从而 \( c_1=3/2 \)。
          于是，表示如下。
          \begin{equation*}
            \rep{\colvec[r]{1 \\ 2}}{B}=\colvec[r]{3/2 \\ 1/2}_B
          \end{equation*}
        \partsitem 关系式
           $c_1\cdot(1)+c_2\cdot(1+x)+c_3\cdot(1+x+x^2)+c_4\cdot(1+x+x^2+x^3)
             =x^2+x^3$
           一眼即可解出：$c_4=1$，$c_3=0$，$c_2=-1$，
           以及 $c_1=0$。
            \begin{equation*}
               \rep{x^2+x^3}{D}=\colvec[r]{0 \\ -1 \\ 0 \\ 1}_D
            \end{equation*}
        \partsitem \( \rep{\colvec[r]{0 \\ -1 \\ 0 \\ 1}}{\stdbasis_4}
                     =\colvec[r]{0 \\ -1 \\ 0 \\ 1}_{\stdbasis_4} \)
      \end{exparts}
    
\end{ans}
\begin{ans}{二.III.1.23}
       求解
      \begin{equation*}
        \colvec{3 \\ -1}=\colvec{1 \\ -1}\cdot c_1
                         +\colvec{1 \\ 1}\cdot c_2
      \end{equation*}
      给出 $c_1=2$ 与~$c_2=1$。
      \begin{equation*}
        \rep{\colvec{3 \\ -1}}{B_1}
           =\colvec{2 \\ 1}_{B_1}
      \end{equation*}
      类似地，求解
      \begin{equation*}
        \colvec{3 \\ -1}=\colvec{1 \\ 2}\cdot c_1
                         +\colvec{1 \\ 3}\cdot c_2
      \end{equation*}
      给出下式。
      \begin{equation*}
        \rep{\colvec{3 \\ -1}}{B_2}
           =\colvec{10 \\ -7}_{B_2}
      \end{equation*}
    
\end{ans}
\begin{ans}{二.III.1.24}
      一个自然的基是 \( \sequence{1,x,x^2} \)。
      $\polyspace_2$ 的有些基不含有任何
      一次或零次多项式。
      例如 \( \sequence{1+x+x^2,x+x^2,x^2} \)。
      （不过，每个基都至少含有一个二次多项式。）
    
\end{ans}
\begin{ans}{二.III.1.25}
      化简
      \begin{equation*}
        \begin{amat}[r]{4}
          1  &-4  &3  &-1  &0  \\
          2  &-8  &6  &-2  &0
        \end{amat}
        \grstep{-2\rho_1+\rho_2}
        \begin{amat}[r]{4}
          1  &-4  &3  &-1  &0  \\
          0  &0   &0  &0   &0
        \end{amat}
      \end{equation*}
      给出：唯一的条件是
      $x_1=4x_2-3x_3+x_4$。
      解集为
      \begin{multline*}
        \set{\colvec{4x_2-3x_3+x_4 \\ x_2 \\ x_3 \\ x_4}
               \suchthat x_2,x_3,x_4\in\Re}                          \\
        =\set{x_2\colvec[r]{4 \\ 1 \\ 0 \\ 0}
             +x_3\colvec[r]{-3 \\ 0 \\ 1 \\ 0}
             +x_4\colvec[r]{1 \\ 0 \\ 0 \\ 1} \suchthat x_2,x_3,x_4\in\Re}
      \end{multline*}
      因此，基的显然候选就是下面这个。
      \begin{equation*}
       \sequence{\colvec[r]{4 \\ 1 \\ 0 \\ 0},
                   \colvec[r]{-3 \\ 0 \\ 1 \\ 0},
                   \colvec[r]{1 \\ 0 \\ 0 \\ 1}  }
      \end{equation*}
      我们已经证明它张成该空间，而证明它还线性无关
      则是例行的。
    
\end{ans}
\begin{ans}{二.III.1.26}
      基有很多。
      下面是一个自然的基。
      \begin{equation*}
        \sequence{
           \begin{mat}[r]
             1  &0  \\
             0  &0
           \end{mat},
           \begin{mat}[r]
             0  &1  \\
             0  &0
           \end{mat},
           \begin{mat}[r]
             0  &0  \\
             1  &0
           \end{mat},
           \begin{mat}[r]
             0  &0  \\
             0  &1
           \end{mat}  }
      \end{equation*}
    
\end{ans}
\begin{ans}{二.III.1.27}
        对每一小题，都有多种可能的答案。
        \begin{exparts}
          \partsitem 一种做法是参数化：把 $a_2$ 表示为另外两个的组合，
            即 $a_2=2a_1+a_0$。
            于是
            $a_2x^2+a_1x+a_0$ 即为 $(2a_1+a_0)x^2+a_1x+a_0$，而
            \begin{multline*}
              \set{(2a_1+a_0)x^2+a_1x+a_0\suchthat a_1,a_0\in\Re}          \\
              =\set{a_1\cdot(2x^2+x)+a_0\cdot(x^2+1)\suchthat a_1,a_0\in\Re}
            \end{multline*}
            提示了
            $\sequence{2x^2+x,x^2+1}$。
            这只表明它张成该空间，但验证它线性无关
            是例行的。
          \partsitem 对 $\set{\rowvec{a  &b  &c}\suchthat a+b=0}$ 参数化
            得到 $\set{\rowvec{-b &b &c}\suchthat b,c\in\Re}$，这提示我们
            使用序列 $\sequence{\rowvec{-1 &1 &0},\rowvec{0 &0 &1}}$。
            我们已证明它张成该空间，而验证它线性无关
            也容易。
          \partsitem 改写
            \begin{equation*}
              \set{\begin{mat}
                     a  &b  \\
                     0  &2b
                   \end{mat}\suchthat a,b\in\Re}
              =\set{a\cdot\begin{mat}[r]
                     1  &0  \\
                     0  &0
                   \end{mat}
                  +b\cdot\begin{mat}[r]
                           0  &1  \\
                           0  &2
                          \end{mat} \suchthat a,b\in\Re}
            \end{equation*}
            提示了下面这个基。
            \begin{equation*}
              \sequence{\begin{mat}[r]
                          1  &0  \\
                          0  &0
                        \end{mat},
                        \begin{mat}[r]
                          0  &1  \\
                          0  &2
                        \end{mat}}
            \end{equation*}
         \end{exparts}
      
\end{ans}
\begin{ans}{二.III.1.28}
      \begin{exparts}
        \partsitem
          将 $a-2b+c-d=0$ 参数化为 $a=2b-c+d$，$b=b$，$c=c$，以及~$d=d$，得到
          把 $M$ 描述为一个三元集合的张成。
          \begin{equation*}
            M=\set{
             (2+x)\cdot b+(-1+x^2)\cdot c+(1+x^3)\cdot d
             \suchthat b,c,d\in\Re
             }
          \end{equation*}
          为证明这个三元集合是基，剩下只需
          验证它线性无关。
          \begin{equation*}
             0+0x+0x^2+0x^3=(2+x)\cdot c_1+(-1+x^2)\cdot c_2+(1+x^3)\cdot c_3
          \end{equation*}
          由 $x$~项可知 $c_1=0$。
          由 $x^2$~项可知 $c_2=0$。
          $x^3$~项给出 $c_3=0$。
       \partsitem
         先把该描述参数化
         （注意，$b$ 与~$d$ 未在 $W$ 的
         描述中被提到
         并不意味着
         它们为零或不存在，而是意味着它们不受限制）。
         \begin{equation*}
           W=\set{
             \begin{mat}
               0 &1 \\
               0 &0
             \end{mat}\cdot b
             +
             \begin{mat}
               1 &0 \\
               1 &0
             \end{mat}\cdot c
             +
             \begin{mat}
               0 &0 \\
               0 &1
             \end{mat}\cdot d
             \suchthat b,c,d\in\Re}
         \end{equation*}
         这就把 $W$ 给成了一个三元集合的张成。
         若再证明该集合线性无关，我们就完成了。
         \begin{equation*}
             \begin{mat}
               0 &0 \\
               0 &0
             \end{mat}
             =
             \begin{mat}
               0 &1 \\
               0 &0
             \end{mat}\cdot c_1
             +
             \begin{mat}
               1 &0 \\
               1 &0
             \end{mat}\cdot c_2
             +
             \begin{mat}
               0 &0 \\
               0 &1
             \end{mat}\cdot c_3
         \end{equation*}
         用右上角的元素可知 $c_1=0$。
         左上角的元素给出 $c_2=0$，而左下角的元素
         表明 $c_3=0$。
     \end{exparts}
   
\end{ans}
\begin{ans}{二.III.1.29}
      我们将证明第二个是基；第一个类似。
      我们将直接从基的定义出发来证明，
      因为这个例子出现在
      \nearbytheorem{th:BasisIffUniqueRepWRT} 之前。

      为看出它线性无关，
      我们建立
      \( c_1\cdot(\cos\theta-\sin\theta)+c_2\cdot(2\cos\theta+3\sin\theta)=
        0\cos\theta+0\sin\theta \)。
      取 \( \theta=0 \) 与 \( \theta=\pi/2 \) 给出这个方程组
      \begin{equation*}
        \begin{linsys}{2}
          c_1\cdot 1    &+  &c_2\cdot 2   &=  &0  \\
          c_1\cdot (-1) &+  &c_2\cdot 3   &=  &0
        \end{linsys}
        \grstep{\rho_1+\rho_2}
        \begin{linsys}{2}
          c_1  &+  &2c_2   &=  &0  \\
               &+  &5c_2   &=  &0
        \end{linsys}
      \end{equation*}
      它表明 $c_1=0$ 与 $c_2=0$。

      张成的计算也不难；对任意 $x,y\in\Re$，
      我们有
      \( c_1\cdot(\cos\theta-\sin\theta)+c_2\cdot(2\cos\theta+3\sin\theta)=
        x\cos\theta+y\sin\theta \)
      给出 \( c_2=x/5+y/5 \) 与 \( c_1=3x/5-2y/5 \)，
      因此张成就是整个空间。
    
\end{ans}
\begin{ans}{二.III.1.30}
      \begin{exparts}
         \partsitem 问哪些 $a_0+a_1x+a_2x^2$ 可以表示为
           $c_1\cdot (1+x)+c_2\cdot (1+2x)$，
           就产生三个线性方程，
           分别比较 $x^2$ 的系数、$x$ 的系数与常数项。
           \begin{equation*}
             \begin{linsys}{2}
               c_1 &+ &c_2  &= &a_0 \\
               c_1 &+ &2c_2 &= &a_1 \\
                   &  &0    &= &a_2
             \end{linsys}
           \end{equation*}
           带回代的高斯消元法表明，
           只要 $a_2=0$，就有 $c_2=-a_0+a_1$
           与 $c_1=2a_0-a_1$。
           于是，当 $a_2=0$ 时，对任意 $a_0$ 和 $a_1$，
           我们都能算出合适的
           $c_1$ 与 $c_2$。
           所以张成就是全体一次多项式
           $\set{a_0+a_1x\suchthat a_0,a_1\in\Re}$。
           将该集合参数化为
           $\set{a_0\cdot 1+a_1\cdot x\suchthat a_0,a_1\in\Re}$
           提示了一个基 $\sequence{1,x}$
           （我们已证明它张成；验证线性无关也容易）。
        \partsitem 利用
          \begin{equation*}
            a_0+a_1x+a_2x^2
            =c_1\cdot(2-2x)+c_2\cdot(3+4x^2)
            =(2c_1+3c_2)+(-2c_1)x+(4c_2)x^2
          \end{equation*}
          我们得到这个方程组。
           \begin{equation*}
             \begin{linsys}{2}
               2c_1  &+ &3c_2  &= &a_0 \\
               -2c_1 &  &      &= &a_1 \\
                     &  &4c_2  &= &a_2
              \end{linsys}
             \grstep{\rho_1+\rho_2}
             \repeatedgrstep{(-4/3)\rho_2+\rho_3}
             \begin{linsys}{2}
               2c_1  &+ &3c_2  &= &a_0\hfill\hbox{} \\
                     &  &3c_2  &= &a_0+a_1\hfill\hbox{} \\
                     &  &0     &= &(-4/3)a_0-(4/3)a_1+a_2
              \end{linsys}
           \end{equation*}
           于是，有相应 $c$ 的二次多项式 $a_0+a_1x+a_2x^2$
           只有那些满足
           $0=(-4/3)a_0-(4/3)a_1+a_2$ 的。
           因此张成如下。
           \begin{equation*}
             \set{(-a_1+(3/4)a_2)+a_1x+a_2x^2\suchthat a_1,a_2\in\Re}
           \end{equation*}
           参数化给出
           $\set{a_1\cdot (-1+x)+a_2\cdot ((3/4)+x^2)\suchthat a_1,a_2\in\Re}$，
           这提示了 $\sequence{-1+x,(3/4)+x^2}$
           （验证它线性无关是例行的）。
      \end{exparts}
    
\end{ans}
\begin{ans}{二.III.1.31}
       \begin{exparts}
        \partsitem 该子空间如下。
          \begin{equation*}
             \set{a_0+a_1x+a_2x^2+a_3x^3 \suchthat a_0+7a_1+49a_2+343a_3=0 }
          \end{equation*}
          改写 $a_0=-7a_1-49a_2-343a_3$ 给出下式。
          \begin{equation*}
             \set{(-7a_1-49a_2-343a_3)+a_1x+a_2x^2+a_3x^3
               \suchthat a_1,a_2,a_3\in\Re }
          \end{equation*}
          分离出
          各参数后，这提示 \( \sequence{-7+x,-49+x^2,-343+x^3} \)
          可作基（容易验证）。
        \partsitem 所给子空间是三次多项式
          $p(x)=a_0+a_1x+a_2x^2+a_3x^3$ 的全体，其中 $a_0+7a_1+49a_2+343a_3=0$
          且 $a_0+5a_1+25a_2+125a_3=0$。
          高斯消元法
          \begin{equation*}
            \begin{linsys}{4}
              a_0  &+  &7a_1  &+  &49a_2  &+  &343a_3  &=  &0  \\
              a_0  &+  &5a_1  &+  &25a_2  &+  &125a_3  &=  &0
            \end{linsys}
            \grstep{-\rho_1+\rho_2}
            \begin{linsys}{4}
              a_0  &+  &7a_1  &+  &49a_2  &+  &343a_3  &=  &0  \\
                   &   &-2a_1 &-  &24a_2  &-  &218a_3  &=  &0
            \end{linsys}
          \end{equation*}
          给出 $a_1=-12a_2-109a_3$ 与 $a_0=35a_2+420a_3$。
          把 $(35a_2+420a_3)+(-12a_2-109a_3)x+a_2x^2+a_3x^3$
          改写为 $a_2\cdot(35-12x+x^2)+a_3\cdot(420-109x+x^3)$
          提示了基 $\sequence{35-12x+x^2,420-109x+x^3}$。
          上面已证明它张成该空间。
          验证它线性无关是例行的。
          （\textit{评注。}
          一个值得做的检验是验证基中的两个多项式
          都以七和五为根。）
        \partsitem 这里对三次多项式有三个条件：
          $a_0+7a_1+49a_2+343a_3=0$，$a_0+5a_1+25a_2+125a_3=0$，
          以及 $a_0+3a_1+9a_2+27a_3=0$。
          高斯消元法
          \begin{equation*}
            \begin{linsys}{4}
              a_0  &+  &7a_1  &+  &49a_2  &+  &343a_3  &=  &0  \\
              a_0  &+  &5a_1  &+  &25a_2  &+  &125a_3  &=  &0  \\
              a_0  &+  &3a_1  &+  &9a_2   &+  &27a_3   &=  &0
            \end{linsys}
            \grstep[-\rho_1+\rho_3]{-\rho_1+\rho_2}
            \repeatedgrstep{-2\rho_2+\rho_3}
            \begin{linsys}{4}
              a_0  &+  &7a_1  &+  &49a_2  &+  &343a_3  &=  &0  \\
                   &   &-2a_1 &-  &24a_2  &-  &218a_3  &=  &0  \\
                   &   &      &   &8a_2   &+  &120a_3  &=  &0
            \end{linsys}
          \end{equation*}
          给出唯一的自由变量 $a_3$，其中
          $a_2=-15a_3$，$a_1=71a_3$，$a_0=-105a_3$。
          参数化如下。
          \begin{multline*}
            \set{(-105a_3)+(71a_3)x+(-15a_3)x^2+(a_3)x^3\suchthat a_3\in\Re} \\
            =
            \set{a_3\cdot(-105+71x-15x^2+x^3)\suchthat a_3\in\Re}
          \end{multline*}
          因此，基的一个自然候选是
          $\sequence{-105+71x-15x^2+x^3}$。
          由上面的推导可知它张成该空间。
          它显然线性无关，因为它是单元素
          集合（且那个唯一的元素不是该空间的零对象）。
          于是，过三点 $(7,0)$、$(5,0)$ 和
          $(3,0)$ 的任何三次多项式都是它的倍数。
          （\textit{评注。}
          与前一题一样，
          一个值得做的检验是验证把七、五、三
          代入这个多项式每次都得到零。）
        \partsitem 这是 $\polyspace_3$ 的平凡子空间。
          因此，基为空 $\sequence{}$。
      \end{exparts}
      \noindent\textit{注。}
      另外，我们也可以通过展开 $(x-7)(x-5)(x-3)$
      来导出第三小题中的多项式。
     
\end{ans}
\begin{ans}{二.III.1.32}
      是的。
      线性无关与张成都不会因重排而改变。
    
\end{ans}
\begin{ans}{二.III.1.33}
      线性无关的集合都不包含零向量。
    
\end{ans}
\begin{ans}{二.III.1.34}
      \begin{exparts}
         \partsitem 为证明它线性无关，注意如果
           \( d_1(c_1\vec{\beta}_1)+d_2(c_2\vec{\beta}_2)+d_3(c_3\vec{\beta}_3)
               =\zero \)，
           那么
           \( (d_1c_1)\vec{\beta}_1+(d_2c_2)\vec{\beta}_2+(d_3c_3)\vec{\beta}_3
               =\zero \)，而这转而
           蕴含每个 \( d_ic_i \) 都是零。
           但由 \( c_i\neq 0 \)，这意味着每个 \( d_i \) 都是零。
           证明它张成该空间也大体相同；
           由于 $\sequence{\vec{\beta}_1,\vec{\beta}_2,\vec{\beta}_3}$
           是基，从而张成该空间，故对任意 $\vec{v}$ 我们可写成
           \( \vec{v}=d_1\vec{\beta}_1+d_2\vec{\beta}_2+d_3\vec{\beta}_3 \)，
           进而
           \( \vec{v}=(d_1/c_1)(c_1\vec{\beta}_1)
               +(d_2/c_2)(c_2\vec{\beta}_2)+(d_3/c_3)(c_3\vec{\beta}_3) \)。

           如果这些标量中有零，那么结果不是基，
           因为它不线性无关。
         \partsitem 证明
           $\sequence{2\vec{\beta}_1,\vec{\beta}_1+\vec{\beta}_2,
             \vec{\beta}_1+\vec{\beta}_3}$ 线性无关是容易的。
           为证明它张成该空间，设
           \( \vec{v}=d_1\vec{\beta}_1+d_2\vec{\beta}_2+d_3\vec{\beta}_3 \)。
           那么，我们可以把同一个 \( \vec{v} \) 关于
           \( \sequence{2\vec{\beta}_1,\vec{\beta}_1+\vec{\beta}_2,
                         \vec{\beta}_1+\vec{\beta}_3} \)
           这样表示：
           $\vec{v}=(1/2)(d_1-d_2-d_3)(2\vec{\beta}_1)
           +d_2(\vec{\beta}_1+\vec{\beta}_2)+d_3(\vec{\beta}_1+\vec{\beta}_3)$。
      \end{exparts}
    
\end{ans}
\begin{ans}{二.III.1.35}
      如果我们略去 $\vec{v}$，每一个都构成线性无关的集合。
      为保持线性无关，我们必须扩大每一个的张成。
      也就是说，我们必须确定每一个的张成（把 $\vec{v}$ 排除在外），
      然后选取一个位于该张成之外的 $\vec{v}$。
      最后，我们还必须检验结果张成整个所给的
      空间。
      这些检验都是例行的。
      \begin{exparts}
        \partsitem 任何不是所给向量的倍数的向量，
          即任何不在直线 $y=x$ 上的向量，在这里都可以。
          例如 $\vec{v}=\vec{e}_1$。
        \partsitem 通过观察，我们注意到向量 $\vec{e}_3$
          不在所给两个向量构成的集合的张成中。
          检验所得集合是 $\Re^3$ 的基
          是例行的。
        \partsitem 对张成
          $\set{c_1\cdot(x)+c_2\cdot(1+x^2)\suchthat c_1,c_2\in\Re}$ 中的任何元素，
          $x^2$ 的系数都等于常数项。
          所以，如果我们添加一个不具此性质的二次多项式，
          比如 $\vec{v}=1-x^2$，就能扩大张成。
          检验结果是 $\polyspace_2$ 的基是容易的。
      \end{exparts}
    
\end{ans}
\begin{ans}{二.III.1.36}
    \begin{exparts}
      \partsitem $1\cdot(2+4x^2)-3\cdot(1+3x^2)+1\cdot(1+5x^2)=0+0x+0x^2$
      \partsitem 简单的计算给出下式。
        \begin{equation*}
          \rep{2+4x^2}{B}=\colvec{1 \\ 1 \\ 4}
          \quad
          \rep{1+3x^2}{B}=\colvec{1/2 \\ 1/2 \\ 3}
          \quad
          \rep{1+5x^2}{B}=\colvec{1/2 \\ 1/2 \\ 5}
        \end{equation*}
      \partsitem 检验是直截了当的。
        \begin{equation*}
          1\cdot\colvec{1 \\ 1 \\ 4}
          -3\cdot\colvec{1/2 \\ 1/2 \\ 3}
          +1\cdot\colvec{1/2 \\ 1/2 \\ 5}
          =\colvec{0 \\ 0 \\ 0}
        \end{equation*}
    \end{exparts}
    
\end{ans}
\begin{ans}{二.III.1.37}
      为证明每个标量都是零，只需作差：
      \( c_1\vec{\beta}_1+\dots+c_k\vec{\beta}_k
          -c_{k+1}\vec{\beta}_{k+1}-\dots-c_n\vec{\beta}_n=\zero \)。
      显然的推广是：在任何只含
      \( \vec{\beta} \) 且其中每个 \( \vec{\beta} \) 只出现
      一次的方程中，每个标量都是零。
      例如，一个方程右端是偶数编号基向量的组合
      （即 $\vec{\beta}_2$、$\vec{\beta}_4$ 等），左端是
      奇数编号基向量的组合，也同样给出
      所有系数都是零的结论。
    
\end{ans}
\begin{ans}{二.III.1.38}
      不能；线性无关的集合都不包含零向量。
    
\end{ans}
\begin{ans}{二.III.1.39}
      这是 $\Re^2$ 的一个不是基的子集，以及它的元素的
      两个不同的线性组合，二者都合成为同一个向量。
      \begin{equation*}
        \set{\colvec[r]{1 \\ 2},\colvec[r]{2 \\ 4}}
        \qquad
        2\cdot\colvec[r]{1 \\ 2}+0\cdot\colvec[r]{2 \\ 4}
        =0\cdot\colvec[r]{1 \\ 2}+1\cdot\colvec[r]{2 \\ 4}
      \end{equation*}
      于是，当子集不是基时，可能出现
      其线性组合不唯一的情形。

      但子集不是基并不意味着
      它的组合必定不唯一。
      例如，这个集合
      \begin{equation*}
        \set{\colvec[r]{1 \\ 2}}
      \end{equation*}
      确实具有如下性质：
      \begin{equation*}
        c_1\cdot\colvec[r]{1 \\ 2}
        =
        c_2\cdot\colvec[r]{1 \\ 2}
      \end{equation*}
      蕴含 $c_1=c_2$。
      这里的要点是，这个子集之所以不是基，是因为它
      不张成该空间；该定理的证明建立了
      线性组合唯一当且仅当该子集线性
      无关。
     
\end{ans}
\begin{ans}{二.III.1.40}
      \begin{exparts}
        \partsitem 把该向量空间描述为
          \begin{equation*}
             \set{\begin{mat}
                     a  &b  \\
                     b  &c
                  \end{mat}  \suchthat a,b,c\in\Re}
          \end{equation*}
          提示了下面这个基。
          \begin{equation*}
            \sequence{
              \begin{mat}[r]
                1  &0  \\
                0  &0
              \end{mat},
              \begin{mat}[r]
                0  &0  \\
                0  &1
              \end{mat},
              \begin{mat}[r]
                0  &1  \\
                1  &0
              \end{mat}  }
          \end{equation*}
          容易验证。
        \partsitem 这是一个可能的基。
          \begin{equation*}
            \sequence{
              \begin{mat}[r]
                1  &0  &0  \\
                0  &0  &0  \\
                0  &0  &0
              \end{mat},
              \begin{mat}[r]
                0  &0  &0  \\
                0  &1  &0  \\
                0  &0  &0
              \end{mat},
              \begin{mat}[r]
                0  &0  &0  \\
                0  &0  &0  \\
                0  &0  &1
              \end{mat},
              \begin{mat}[r]
                0  &1  &0  \\
                1  &0  &0  \\
                0  &0  &0
              \end{mat},
              \begin{mat}[r]
                0  &0  &1  \\
                0  &0  &0  \\
                1  &0  &0
              \end{mat},
              \begin{mat}[r]
                0  &0  &0  \\
                0  &0  &1  \\
                0  &1  &0
              \end{mat}  }
          \end{equation*}
         \partsitem 与前两题一样，我们可以用两类
           矩阵构成一个基。
           第一类是只有一个对角元为 1、
           其余元素全为零的矩阵（这样的矩阵有 \( n \) 个）。
           第二类是两个相对的非对角元素为 1、
           其余元素全为零的矩阵。
           （也就是说，$M$ 中除 $m_{i,j}$ 与 $m_{j,i}$ 为 1 外，
           其余元素全为零。）
      \end{exparts}
    
\end{ans}
\begin{ans}{二.III.1.41}
       \begin{exparts}
         \partsitem 来自 $\Re^3$ 的任意四个向量都是线性相关的，因为
           向量方程
           \begin{equation*}
             c_1\colvec{x_1 \\ y_1 \\ z_1}
             +c_2\colvec{x_2 \\ y_2 \\ z_2}
             +c_3\colvec{x_3 \\ y_3 \\ z_3}
             +c_4\colvec{x_4 \\ y_4 \\ z_4}
             =\colvec[r]{0 \\ 0 \\ 0}
           \end{equation*}
           产生一个线性方程组
           \begin{equation*}
             \begin{linsys}{4}
                x_1c_1  &+  &x_2c_2  &+  &x_3c_3  &+  &x_4c_4  &=  &0 \\
                y_1c_1  &+  &y_2c_2  &+  &y_3c_3  &+  &y_4c_4  &=  &0 \\
                z_1c_1  &+  &z_2c_2  &+  &z_3c_3  &+  &z_4c_4  &=  &0
             \end{linsys}
           \end{equation*}
           它是齐次的（因而有解），且有
           四个未知数却只有三个方程，
           因而有非平凡解。
           （当然，这个论证适用于 $\Re^3$ 的任何
           含有四个或更多向量的子集。）
         \partsitem 我们只做两个向量的情形。
           给定 $x_1$、\ldots、$z_2$，
           \begin{equation*}
             S=\set{\colvec{x_1 \\ y_1 \\ z_1},
             \colvec{x_2 \\ y_2 \\ z_2}}
           \end{equation*}
           为判定哪些向量
           \begin{equation*}
             \colvec{x \\ y \\ z}
           \end{equation*}
           在 $S$ 的张成中，建立
           \begin{equation*}
             c_1\colvec{x_1 \\ y_1 \\ z_1}
             +c_2\colvec{x_2 \\ y_2 \\ z_2}
             =\colvec{x \\ y \\ z}
           \end{equation*}
           并对所得方程组作行化简。
           \begin{equation*}
             \begin{linsys}{2}
                x_1c_1  &+  &x_2c_2   &=  &x \\
                y_1c_1  &+  &y_2c_2   &=  &y \\
                z_1c_1  &+  &z_2c_2   &=  &z
             \end{linsys}
           \end{equation*}
           只有两个
           变量 $c_1$ 和 $c_2$，却有三个方程，因此
           当高斯消元法结束时，最后一
           行上会有某个形如
           $0=m_1x+m_2y+m_3z$ 的关系式。
           因此，二元集合 $S$ 的张成中的向量
           必须满足某种限制。
           因此该张成不是整个 $\Re^3$。
       \end{exparts}
    
\end{ans}
\begin{ans}{二.III.1.42}
      我们有（小心地使用这些古怪的运算）
      \begin{multline*}
        \set{\colvec{1-y-z \\ y \\ z}\suchthat y,z\in\Re}
         =\set{\colvec{-y+1 \\ y \\ 0}
               +\colvec{-z+1 \\ 0 \\ z}
              \suchthat y,z\in\Re}                           \\
         =\set{y\cdot\colvec{0 \\ 1 \\ 0}
               +z\cdot\colvec{0 \\ 0 \\ 1}
              \suchthat y,z\in\Re}
      \end{multline*}
      因此基的一个自然候选是下面这个。
      \begin{equation*}
        \sequence{\colvec[r]{0 \\ 1 \\ 0},\colvec[r]{0 \\ 0 \\ 1}}
      \end{equation*}
      为检验线性无关，我们建立
      \begin{equation*}
         c_1\colvec[r]{0 \\ 1 \\ 0}+c_2\colvec[r]{0 \\ 0 \\ 1}
         =\colvec[r]{1 \\ 0 \\ 0}
      \end{equation*}
      （右边的向量是这个空间中的零对象）。
      这给出线性方程组
      \begin{equation*}
        \begin{linsys}{3}
          (-c_1+1)  &+  &(-c_2+1)  &-  &1  &=  &1  \\
             c_1    &   &          &   &   &=  &0  \\
                    &   &c_2       &   &   &=  &0
        \end{linsys}
      \end{equation*}
      它仅有解 $c_1=0$ 与 $c_2=0$。
      检验张成也是类似的。
   
\end{ans}
\protect \subsection {二.III.2: 维数}
\begin{ans}{二.III.2.15}
      一个基是 \( \sequence{1,x,x^2} \)，于是
      维数是三。
    
\end{ans}
\begin{ans}{二.III.2.16}
      解集为
      \begin{equation*}
        \set{\colvec{4x_2-3x_3+x_4 \\ x_2 \\ x_3 \\ x_4}
               \suchthat x_2,x_3,x_4\in\Re}
      \end{equation*}
      于是一个自然的基是这个
      \begin{equation*}
       \sequence{\colvec[r]{4 \\ 1 \\ 0 \\ 0},
                   \colvec[r]{-3 \\ 0 \\ 1 \\ 0},
                   \colvec[r]{1 \\ 0 \\ 0 \\ 1}  }
      \end{equation*}
      （检验线性无关是容易的）。
      于是维数是三。
    
\end{ans}
\begin{ans}{二.III.2.17}
      \begin{exparts}
        \partsitem
          参数化，得到该空间的如下描述。
          \begin{equation*}
            \set{\colvec{w-z \\ y \\ z \\ w}
                 =\colvec{0 \\ 1 \\ 0 \\ 0}y+
                 \colvec{-1 \\ 0 \\ 1 \\ 0}z+
                 \colvec{1 \\ 0 \\ 0 \\ 1}w
                 \suchthat y,z,w\in\Re}
          \end{equation*}
          这把该空间表示为那个由三个向量组成的集合的张成空间。
          为证明这个三向量集构成一个基，我们检验它是
          线性无关的。
          \begin{equation*}
             \colvec{0 \\ 0 \\ 0 \\ 0}
                 =\colvec{0 \\ 1 \\ 0 \\ 0}c_1+
                 \colvec{-1 \\ 0 \\ 1 \\ 0}c_2+
                 \colvec{1 \\ 0 \\ 0 \\ 1}c_3
          \end{equation*}
          第二个分量给出 $c_1=0$，第三、第四个
          分量给出 $c_2=0$ 和~$c_3=0$。
          于是一个基是这个。
          \begin{equation*}
              \sequence{
                 \colvec{0 \\ 1 \\ 0 \\ 0},
                 \colvec{-1 \\ 0 \\ 1 \\ 0},
                 \colvec{1 \\ 0 \\ 0 \\ 1}
              }
          \end{equation*}
          维数是基中向量的个数：$3$。
        \partsitem
          自然的参数化是这个。
          \begin{multline*}
            \set{
              \begin{mat}
                a &0 &0 &0 &0 \\
                0 &b &0 &0 &0 \\
                0 &0 &c &0 &0 \\
                0 &0 &0 &d &0 \\
                0 &0 &0 &0 &e
              \end{mat}
              \suchthat a,\ldots,e\in\Re}   \\
            \set{
              \begin{mat}
                1 &0 &0 &0 &0 \\
                0 &0 &0 &0 &0 \\
                0 &0 &0 &0 &0 \\
                0 &0 &0 &0 &0 \\
                0 &0 &0 &0 &0
              \end{mat}\cdot a
              +\begin{mat}
                0 &0 &0 &0 &0 \\
                0 &1 &0 &0 &0 \\
                0 &0 &0 &0 &0 \\
                0 &0 &0 &0 &0 \\
                0 &0 &0 &0 &0
              \end{mat}\cdot b
              +\cdots
              \suchthat a,\ldots,e\in\Re}
          \end{multline*}
          检验这个五元集线性无关是平凡的。
          所以这是一个基；维数是 $5$。
          \begin{equation*}
            \sequence{
              \begin{mat}
                1 &0 &0 &0 &0 \\
                0 &0 &0 &0 &0 \\
                0 &0 &0 &0 &0 \\
                0 &0 &0 &0 &0 \\
                0 &0 &0 &0 &0
              \end{mat},
              \begin{mat}
                0 &0 &0 &0 &0 \\
                0 &1 &0 &0 &0 \\
                0 &0 &0 &0 &0 \\
                0 &0 &0 &0 &0 \\
                0 &0 &0 &0 &0
              \end{mat},
              \,\ldots\,,
              \begin{mat}
                0 &0 &0 &0 &0 \\
                0 &0 &0 &0 &0 \\
                0 &0 &0 &0 &0 \\
                0 &0 &0 &0 &0 \\
                0 &0 &0 &0 &1
              \end{mat}
          }
          \end{equation*}
        \partsitem
          这些约束条件构成一个两个方程、四个未知数的线性方程组。
          将该方程组参数化，用自由变量表示
          首变量，得到 $a_0=-a_1$、$a_1=a_1$、$a_2=2a_3$ 以及~$a_3=a_3$。
          \begin{multline*}
            \set{-a_1+a_1x+2a_3x^2+a_3x^3\suchthat a_1,a_3\in\Re}         \\
            =\set{(-1+x)\cdot a_1+(2x^2+x^3)\cdot a_3\suchthat a_1,a_3\in\Re}
          \end{multline*}
          该描述表明这个空间是二元集
          $\set{-1+x,x^2+x^3}$ 的张成空间。
          只要证明该集合线性无关即完成证明。
          关系式
          \begin{equation*}
            0+0x+0x^2+0x^3=(-1+x)\cdot c_1+(2x^2+x^3)\cdot c_2
          \end{equation*}
          由常数项给出 $c_1=0$，由
          三次项给出 $c_2=0$。
          该空间的一个基是 $\sequence{-1+x,2x^2+x^3}$。
          这是一个二维空间。

      \end{exparts}
    
\end{ans}
\begin{ans}{二.III.2.18}
      对这个空间
      \begin{multline*}
        \set{\begin{mat}
               a  &b  \\
               c  &d
             \end{mat} \suchthat a,b,c,d\in\Re}            \\
        =\set{a\cdot\begin{mat}[r]
             1  &0  \\
             0  &0
           \end{mat}
           +\dots+
           d\cdot\begin{mat}[r]
             0  &0  \\
             0  &1
           \end{mat} \suchthat a,b,c,d\in\Re}
      \end{multline*}
      这是一个自然的基。
      \begin{equation*}
        \sequence{
           \begin{mat}[r]
             1  &0  \\
             0  &0
           \end{mat},
           \begin{mat}[r]
             0  &1  \\
             0  &0
           \end{mat},
           \begin{mat}[r]
             0  &0  \\
             1  &0
           \end{mat},
           \begin{mat}[r]
             0  &0  \\
             0  &1
           \end{mat}  }
      \end{equation*}
      维数是四。
    
\end{ans}
\begin{ans}{二.III.2.19}
     \begin{exparts}
      \partsitem 与前面的练习一样，不加限制的矩阵空间 $\matspace_{\nbyn{2}}$
        有这个基
        \begin{equation*}
         \sequence{
           \begin{mat}[r]
             1  &0  \\
             0  &0
           \end{mat},
           \begin{mat}[r]
             0  &1  \\
             0  &0
           \end{mat},
           \begin{mat}[r]
             0  &0  \\
             1  &0
           \end{mat},
           \begin{mat}[r]
             0  &0  \\
             0  &1
           \end{mat}  }
        \end{equation*}
        于是维数是四。
      \partsitem 对这个空间
        \begin{multline*}
         \set{\begin{mat}
               a  &b  \\
               c  &d
             \end{mat} \suchthat \text{$a=b-2c$ 且 $d\in\Re$}}     \\
         =\set{b\cdot\begin{mat}[r]
             1  &1  \\
             0  &0
           \end{mat}
           +c\cdot\begin{mat}[r]
             -2  &0  \\
              1  &0
           \end{mat}
           +d\cdot\begin{mat}[r]
             0  &0  \\
             0  &1
           \end{mat} \suchthat b,c,d\in\Re}
        \end{multline*}
        这是一个自然的基。
        \begin{equation*}
          \sequence{
            \begin{mat}[r]
              1  &1  \\
              0  &0
            \end{mat},
            \begin{mat}[r]
              -2  &0  \\
               1  &0
            \end{mat},
            \begin{mat}[r]
              0  &0  \\
              0  &1
            \end{mat}  }
        \end{equation*}
        维数是三。
      \partsitem 对该两个方程的线性方程组应用高斯消元法，得到
        $c=0$ 以及 $a=-b$。
        于是，我们有这个描述
        \begin{equation*}
         \set{\begin{mat}
               -b  &b  \\
                0  &d
             \end{mat} \suchthat b,d\in\Re}
         =\set{b\cdot\begin{mat}[r]
             -1  &1  \\
             0  &0
           \end{mat}
           +d\cdot\begin{mat}[r]
             0  &0  \\
             0  &1
           \end{mat} \suchthat b,d\in\Re}
        \end{equation*}
        因此这是一个自然的基。
        \begin{equation*}
          \sequence{
            \begin{mat}[r]
              -1  &1  \\
               0  &0
            \end{mat},
            \begin{mat}[r]
              0  &0  \\
              0  &1
            \end{mat}  }
        \end{equation*}
        维数是二。
     \end{exparts}
    
\end{ans}
\begin{ans}{二.III.2.20}
      我们不能简单地数参数的个数。
      也就是说，答案不是~$3$。
      相反，注意我们可以把每个成员
      $\binom{x}{y}\in\Re^2$ 表示成
      \begin{equation*}
        \colvec{x \\ y}=\colvec{a+b \\ a+c}
      \end{equation*}
      的形式，其中取 $a=0$、$b=x$ 和 $c=y$
      （其他取法也是可能的）。
      所以 $S$ 就是集合 $S=\Re^2$。
      它的维数是~$2$。
    
\end{ans}
\begin{ans}{二.III.2.21}
      这些空间的基已
      在前一小节的答案集中导出。
      \begin{exparts}
        \partsitem 一个基是 \( \sequence{-7+x,-49+x^2,-343+x^3} \)。
          维数是三。
        \partsitem 一个基是 $\sequence{35-12x+x^2,420-109x+x^3}$，于是
          维数是二。
        \partsitem 一个基是 $\set{-105+71x-15x^2+x^3}$。
          维数是一。
        \partsitem 这是 $\polyspace_3$ 的平凡子空间，因此
          基是空的。
          维数是零。
      \end{exparts}
    
\end{ans}
\begin{ans}{二.III.2.22}
       先回忆 $\cos2\theta=\cos^2\theta-\sin^2\theta$，所以
       从该集合中去掉 $\cos2\theta$ 不会改变张成空间。
       剩下的，即集合
       $\set{\cos^2\theta,\sin^2\theta,\sin2\theta}$，线性无关
       （考虑关系式
       $c_1\cos^2\theta+c_2\sin^2\theta+c_3\sin2\theta=Z(\theta)$，
       其中 $Z$ 是零函数，然后取
       $\theta=0$、$\theta=\pi/4$ 和 $\theta=\pi/2$，
       即可断定每个 $c$ 都是零）。
       因此它是其张成空间的一个基。
       这表明该张成空间是一个维数为三的向量空间。
     
\end{ans}
\begin{ans}{二.III.2.23}
      这是一个基
      \begin{equation*}
        \sequence{(1+0i,0+0i,\dots,0+0i),\,
                  (0+1i,0+0i,\dots,0+0i),(0+0i,1+0i,\dots,0+0i),\ldots }
      \end{equation*}
      于是维数是 \( 2\cdot 47=94 \)。
    
\end{ans}
\begin{ans}{二.III.2.24}
      一个基是
      \begin{equation*}
        \sequence{
          \begin{mat}[r]
            1  &0  &0  &0  &0  \\
            0  &0  &0  &0  &0  \\
            0  &0  &0  &0  &0
          \end{mat},
          \begin{mat}[r]
            0  &1  &0  &0  &0  \\
            0  &0  &0  &0  &0  \\
            0  &0  &0  &0  &0
          \end{mat},
          \ldots,
          \begin{mat}[r]
            0  &0  &0  &0  &0  \\
            0  &0  &0  &0  &0  \\
            0  &0  &0  &0  &1
          \end{mat}  }
      \end{equation*}
      于是维数是 \( 3\cdot 5=15 \)。
    
\end{ans}
\begin{ans}{二.III.2.25}
       在四维空间中，由四个向量组成的集合线性
       无关当且仅当它张成该空间。
       这些向量的形式使线性无关容易证明
       （先看第四个分量的方程，再看第三个分量的
       方程，等等）。
    
\end{ans}
\begin{ans}{二.III.2.26}
      由本节的结果，因为 $\polyspace_2$ 的维数是~$3$，
      要证明一个线性无关集是一个基，我们只需
      观察到它有三个成员。
      要证明一个集合不是基，我们
      只需观察到它没有三个成员（此时
      我们不必担心线性无关）。
      \begin{exparts}
        \partsitem  这是一个基；一眼即可看出它线性无关
           （第一个元素没有二次项或一次项，第二个有
           二次项但没有一次项，第三个有一次项），
           并且它有三个元素。
        \partsitem 这不是基，因为它只有两个元素。
        \partsitem 这个三元素集是一个基。
        \partsitem 这不是基，因为它有四个元素。
      \end{exparts}
    
\end{ans}
\begin{ans}{二.III.2.27}
      \begin{exparts}
       \partsitem $\polyspace_2$ 的图有四层。
          最顶层有唯一的三维子空间，
          即 $\polyspace_2$ 本身。
          下一层包含二维子空间
          （\emph{并非}只有一次多项式；任何二维
          子空间都行，比如形如 $ax^2+b$ 的多项式）。
          再下面是一维子空间。
          最后，当然是唯一的零维子空间，
          即平凡子空间。
        \partsitem 对 $\matspace_{\nbyn{2}}$，图有五层，
          包括维数从四到零的子空间。
     \end{exparts}
    
\end{ans}
\begin{ans}{二.III.2.28}
      \begin{exparts*}
        \partsitem 一
        \partsitem 二
        \partsitem \( n \)
      \end{exparts*}
    
\end{ans}
\begin{ans}{二.III.2.29}
      我们只需给出一个无限的线性无关集。
      一个这样的序列是 \( \sequence{f_1,f_2,\ldots} \)，其中
      \( \map{f_i}{\Re}{\Re} \)
      为
      \begin{equation*}
         f_i(x)=\begin{cases}
                   1  &\text{若 \( x=i \)}  \\
                   0  &\text{其他情形}
                \end{cases}
      \end{equation*}
      即仅在 $x=i$ 处取值 $1$ 的函数。
    
\end{ans}
\begin{ans}{二.III.2.30}
      函数是有序对
      $(x,f(x))$ 组成的集合。
      所以以空集为定义域的函数只有一个，即空集本身。
      只含一个元素的向量空间是平凡向量空间，
      其维数是零。
    
\end{ans}
\begin{ans}{二.III.2.31}
      应用\nearbycorollary{cor:NoLiSetGreatDim}。
    
\end{ans}
\begin{ans}{二.III.2.32}
      平面具有
      $\set{\vec{p}+t_1\vec{v}_1+t_2\vec{v}_2\suchthat t_1,t_2\in\Re}$ 的形式。
      （第一章也把它称为“$2$-平坦集”，其中
      还讨论了为什么这等价于微积分中常采用的描述，即
      满足形如 $ax+by+cz=d$ 的条件的点 $(x,y,z)$
      组成的集合）。
      当平面过原点时，我们可以取特向量
      $\vec{p}$ 为 $\zero$。
      于是，用本章建立的语言来说，过原点的平面是
      由两个向量组成的集合的张成空间。

      下面给出陈述。
      断言这三个向量不共面，就等同于断言没有向量落在另外两个向量的张成之中\Dash 没有向量是另两个向量的线性组合。
      这恰恰就是断言这个三元集合是线性无关的。
       由\nearbycorollary{cor:NVecsRNSpanIffLI}可知，这等价于断言该集合是 $\Re^3$ 的一个基（更精确地说，由该集合的元素构成的任何序列都是一个基）。
    
\end{ans}
\begin{ans}{二.III.2.33}
      设空间 $V$ 是有限维的，且 $S$ 是 $V$ 的一个子空间。

      若 $S$ 不是有限维的，那么它有一个无穷的线性无关集（从空集出发，不断添加不与该集合线性相关的向量；这一过程可以持续无穷多步，否则 $S$ 就会是有限维的）。
      但 $S$ 的任何线性无关子集也是 $V$ 的线性无关子集，这与\nearbycorollary{cor:NoLiSetGreatDim}矛盾
      % \begin{exparts}
      %   \partsitem The empty set is a linearly independent subset of $S$.
      %     By \nearbycorollary{cor:LIExpBas}, it can be expanded to a basis
      %     for the vector space $S$.
      %   \partsitem Any basis for the subspace $S$ is a linearly independent
      %     set in the superspace $V$.
      %     Hence it can be expanded to a basis for the superspace, which is
      %     finite dimensional.
      %     Therefore it has only finitely many members.
      % \end{exparts}
    
\end{ans}
\begin{ans}{二.III.2.34}
      它保证我们能穷尽 \( \vec{\beta} \) 中的各个向量。也就是说，它为最后一段的第一句话提供了依据。
    
\end{ans}
\begin{ans}{二.III.2.35}
       设 \( B_U \) 是 \( U \) 的一个基，\( B_W \) 是 \( W \) 的一个基。
       考虑将这两个基的序列首尾相连。
       若有重复出现的元素，则交 \( U\intersection W \) 是非平凡的。
       否则，
       集合 \( B_U\union B_W \) 是线性相关的，因为它是五维空间 \( \Re^5 \) 的一个六元子集。
       无论哪种情形，\( B_W \) 中总有某个成员落在 \( B_U \) 的张成之中，于是 \( U\intersection W \) 就不只是平凡空间 \( \set{\zero\,} \)。

       推广：
       若 \( U,W \) 是维数为 \( n \) 的向量空间的子空间，且 \( \dim(U)+\dim(W)>n \)，则它们的交是非平凡的。
    
\end{ans}
\begin{ans}{二.III.2.36}
       首先注意，一个集合是某个空间的基当且仅当它是线性无关的，因为此时它是自身张成空间的基。
       \begin{exparts}
         \partsitem 第二段中问题的答案是“是的”
           （这意味着第一段中两个问题的答案也都是“是的”）。
           若 \( B_U \) 是 \( U \) 的一个基，则 \( B_U \) 是 \( W \) 的一个线性无关子集。
           用\nearbycorollary{cor:LIExpBas}把它扩充为 \( W \) 的一个基。
           那就是所求的 \( B_W \)。

           第三段中问题的答案是“不是”，这意味着第四段问题的答案是“不是”。
           下面举例说明：超空间的一个基，可以没有任何子基构成子空间的基——在 \( W=\Re^2 \) 中，考虑标准基 \( \stdbasis_2 \)。
           $\stdbasis_2$ 的任何子基都不能构成 $\Re^2$ 中直线 \( y=x \) 这个子空间 \( U \) 的基。
         \partsitem 它是一个基（是其张成空间的基），因为线性无关集的交是线性无关的（交是其中每一个线性无关集的子集）。

           然而，它并不是这些空间的交的基。
           例如，下面是 \( \Re^2 \) 的基：
            \begin{equation*}
              B_1=\sequence{\colvec[r]{1 \\ 0},\colvec[r]{0 \\ 1}}
              \quad\text{和}\quad
              B_2=\sequence[r]{\colvec{2 \\ 0},\colvec[r]{0 \\ 2}}
            \end{equation*}
            而 \( \Re^2\intersection\Re^2=\Re^2 \)，但 \( B_1\intersection B_2 \) 为空。
            我们只能说，基的 $\cap$ 是这些空间的交的某个子集的基。
         \partsitem 基的 $\cup$ 未必是基：在 \( \Re^2 \) 中
            \begin{equation*}
              B_1=\sequence{\colvec[r]{1 \\ 0},\colvec[r]{1 \\ 1}}
              \quad\text{和}\quad
              B_2=\sequence{\colvec[r]{1 \\ 0},\colvec[r]{0 \\ 2}}
            \end{equation*}
            \( B_1\union B_2 \) 不是线性无关的。
            两个基的 $\cup$ 为基的一个充要条件是
            \begin{equation*}
              B_1\union B_2 \text{ 线性无关 }
              \quad\iff\quad
              \spanof{B_1\intersection B_2}=\spanof{B_1}\intersection
                                             \spanof{B_2}
            \end{equation*}
            证明它足够容易（但应用起来或许很难）。
          \partsitem 一个基的补不可能是一个基，因为它包含零向量。
        \end{exparts}
      
\end{ans}
\begin{ans}{二.III.2.37}
       \begin{exparts}
          \partsitem \( U \) 的一个基是 \( W \) 中的线性无关集，因此可以用\nearbycorollary{cor:LIExpBas}把它扩充为 \( W \) 的一个基。
            第二个基的成员至少与第一个基一样多。
          \partsitem 一个方向是清楚的：若 \( V=W \)，则二者维数相同。
            反之，设 \( B_U \) 是 \( U \) 的一个基。
            它是 \( W \) 的一个线性无关子集，因此可以扩充为 \( W \) 的一个基。
            若 \( \dim(U)=\dim(W) \)，则 \( W \) 的这个基的成员不会比 \( B_U \) 多，于是它就等于 \( B_U \)。
            由于 \( U \) 与 \( W \) 有相同的基，所以它们相等。
          \partsitem 设 \( W \) 是次数有限的多项式构成的空间，\( U \) 是只含偶数次幂项的多项式构成的子空间。
             \begin{equation*}
                U=\set{a_0+a_1x^2+a_2x^4+\dots+a_nx^{2n}\suchthat
                           a_0,\ldots,a_n\in\Re}
             \end{equation*}
            两个空间都是无限维的，但 \( U \) 是真子空间。
        \end{exparts}
      
\end{ans}
\begin{ans}{二.III.2.38}
      \begin{exparts}
        \partsitem
          $\colvec[r]{1 \\ 2 \\ 0}
             =(-1)\cdot\colvec[r]{1 \\ 0  \\ 0}
              +2\colvec[r]{1 \\ 1 \\ 0}
              +0\cdot\colvec[r]{0 \\ 0 \\ 2}$
         \partsitem 有两个基向量带有非零系数。
            比如我们可以取第一个。
            \begin{equation*}
              \hat{B}=\sequence{\colvec[r]{1 \\ 2  \\ 0},
                \colvec[r]{1 \\ 1 \\ 0},
                \colvec[r]{0 \\ 0 \\ 2}}
            \end{equation*}
            验证它是 \( \Re^3 \) 的另一个基是常规的。
         \partsitem
          称交换的结果为
          \( \hat{B}=\sequence{\vec{\beta}_1,\dots,\vec{\beta}_{i-1},\vec{v},
            \vec{\beta}_{i+1},\dots,\vec{\beta}_n}  \)。
          我们先证明 $\hat{B}$ 是线性无关的。
          $\hat{B}$ 的成员之间的任一关系
          \( d_1\vec{\beta}_1+\dots+d_i\vec{v}+\dots+d_n\vec{\beta}_n=\zero \)，
          在代入 $\vec{v}$ 之后，
          \begin{equation*}
          d_1\vec{\beta}_1+\dots
          +d_i\cdot(c_1\vec{\beta}_1+\dots+c_i\vec{\beta}_i+\dots+c_n\vec{\beta}_n)
          +\dots+d_n\vec{\beta}_n
          =\zero
          \tag*{($*$)}\end{equation*}
          就给出 $B$ 的成员之间的一个线性关系。
          基 $B$ 是线性无关的，所以 $\vec{\beta}_i$ 的系数 $d_ic_i$ 为零。
          由于我们假设了 $c_i$ 非零，故 $d_i=0$。
          把它用于方程~$(*)$ 可知，其余的 $d$ 也全为零。
          因此 $\hat{B}$ 是线性无关的。

          最后我们证明 $\hat{B}$ 与 $B$ 有相同的张成。
          这个论证的一半，即
          $\spanof{\smash{\hat{B}}}\subseteq\spanof{B}$，
          是容易的；我们可以把 $\spanof{\smash{\hat{B}}}$ 的任一成员
          $d_1\vec{\beta}_1+\dots+d_i\vec{v}+\dots+d_n\vec{\beta}_n$
          写成
          $
          d_1\vec{\beta}_1+\dots
          +d_i\cdot(c_1\vec{\beta}_1+\dots+c_n\vec{\beta}_n)
          +\dots+d_n\vec{\beta}_n
          $，
          它是 $B$ 的成员的线性组合的线性组合，从而属于 $\spanof{B}$。
          对于论证的另一半 $\spanof{B}\subseteq\spanof{\smash{\hat{B}}}$，回忆一下：若
          $\vec{v}=c_1\vec{\beta}_1+\dots+c_n\vec{\beta}_n$ 且 $c_i\neq 0$，
          则可以把该等式改写为
          $\vec{\beta}_i=(-c_1/c_i)\vec{\beta}_1+\dots+(1/c_i)\vec{v}+\dots
          +(-c_n/c_i)\vec{\beta}_n$。
          现在，考虑 $\spanof{B}$ 的任一成员
          $d_1\vec{\beta}_1+\dots+d_i\vec{\beta}_i+\dots+d_n\vec{\beta}_n$，
          把 $\vec{\beta}_i$ 替换为它关于 $\hat{B}$ 的成员的线性组合表达式，并像本论证的前一半那样看出，
          结果是 $\hat{B}$ 的成员的线性组合的线性组合，
          从而属于 $\spanof{\smash{\hat{B}}}$。
         \partsitem 取定一个至少有一个有限基的向量空间。
           从这个空间的所有基中，选出一个规模最小的
           \( B=\sequence{\vec{\beta}_1,\dots,\vec{\beta}_n} \)。
           我们将证明任何其他的基
           \( D=\smash{\sequence{\vec{\delta}_1,\vec{\delta}_2,\ldots}} \)
           也有相同的成员数 $n$。
           因为 \( B \) 规模最小，\( D \) 有不少于 \( n \)~个向量。
           我们将论证它不可能有多于 \( n \) 个向量。

           基 \( B \) 张成该空间，而 \( \vec{\delta}_1 \)
           在该空间中，
           所以 \( \vec{\delta}_1 \) 是 \( B \) 的元素的一个非平凡线性组合。
           由交换引理，可以把 \( B \) 中的一个向量换成
           \( \vec{\delta}_1 \)，得到基 \( B_1 \)，
           其中一个元素是
           \( \vec{\delta}_1 \)，
           而其余 \( n-1 \) 个元素都是 \( \vec{\beta} \)。

           上一段构成归纳论证的基础步骤。
           归纳步骤从一个基 \( B_k \)（其中 \( 1\leq k<n \)）开始，
           它包含 \( D \) 的 \( k \) 个成员与 \( B \) 的 \( n-k \) 个成员。
           我们知道 \( D \) 至少有 \( n \) 个成员，所以存在
           \( \vec{\delta}_{k+1} \)。
           把它表示为 \( B_k \) 的元素的线性组合。
           关键点：在那个表示中，
           至少有一个非零标量必须与某个 \( \vec{\beta}_i \) 相关联，
           否则该表示就会是线性无关集 \( D \) 的元素之间的一个
           非平凡线性关系。
           用 \( \vec{\delta}_{k+1} \) 换掉 \( \vec{\beta}_i \)，
           得到新的基
           \( B_{k+1} \)，与前面的基 \( B_k \) 相比，它多一个 \( \vec{\delta} \)、少一个
           \( \vec{\beta} \)。

           重复这一步，直到没有 \( \vec{\beta} \) 剩下，
           于是 \( B_n \) 包含
           $\vec{\delta}_1,\dots,\vec{\delta}_n$。
           现在，\( D \) 不可能有多于这 \( n \) 个的向量，因为
           剩下的任何 \( \vec{\delta}_{n+1} \) 都会落在
           \( B_n \)（因为它是一个基）的张成中，
           从而是其余 $\vec{\delta}$ 的线性组合，
           这与 $D$ 线性无关矛盾。
      \end{exparts}
      
\end{ans}
\begin{ans}{二.III.2.39}
      （这个解答来自网站上 Yuzhou Gu 的一条评论。）
      给每个人指定一个向量 $\vec{v}_i\in\R^n$，
      若此人读过该书，则相应分量为 $1$，否则为 $0$。
      由于这些向量线性相关，我们将得到一个形如 $\sum a_i v_i = 0$ 的等式。
      那么，使 $a_i>0$ 的那些 $i$ 构成的集合与使
      $a_i<0$ 的那些 $i$ 构成的集合，就是两个不相交的集合，且读过的书的并集相同。
    
\end{ans}
\begin{ans}{二.III.2.40}
      $V$ 的维数的可能是 $0$、$1$、$n-1$
      和 $n$。

      为看到这一点，先考虑 $\vec{v}$ 的所有坐标都相等的情形。
      \begin{equation*}
         \vec{v}=\colvec{z \\ z \\ \vdots \\ z}
      \end{equation*}
      那么对每个置换 $\sigma$ 都有 $\sigma(\vec{v})=\vec{v}$，
      所以 $V$ 就是 $\vec{v}$ 的张成，
      依 $\vec{v}$ 是否为 $\zero$ 而定，其维数为 $0$ 或 $1$。

      现在假设 $\vec{v}$ 的坐标不全相等；设 $x$
      与 $y$（$x\neq y$）是 $\vec{v}$ 的坐标中的两个。
      那么我们可以找到置换 $\sigma_1$ 与 $\sigma_2$，使得
      \begin{equation*}
        \sigma_1(\vec{v})=\colvec{x \\ y \\ a_3 \\ \vdots \\ a_n}
        \quad\text{和}\quad
        \sigma_2(\vec{v})=\colvec{y \\ x \\ a_3 \\ \vdots \\ a_n}
      \end{equation*}
      其中 $a_3,\ldots,a_n\in\Re$。
      因此，
      \begin{equation*}
        \frac{1}{y-x}\bigl(\sigma_1(\vec{v})-\sigma_2(\vec{v})\bigr)=
        \colvec[r]{-1 \\ 1 \\ 0 \\ \vdotswithin{-1} \\ 0}
      \end{equation*}
      属于 $V$。
      也就是说，$\vec{e}_2-\vec{e}_1\in V$，其中 $\vec{e}_1$、$\vec{e}_2$、
      \ldots、$\vec{e}_n$ 是 $\Re^n$ 的标准基。
      类似地，$\vec{e}_3-\vec{e}_2$、\ldots、$\vec{e}_n-\vec{e}_1$
      都在 $V$ 中。
      容易看出，向量
      $\vec{e}_2-\vec{e}_1$、$\vec{e}_3-\vec{e}_2$、\ldots、
      $\vec{e}_n-\vec{e}_1$ 是线性无关的（也就是说，构成一个线性无关集），
      所以 $\dim V\geq n-1$。

      最后，我们可以写出
      \begin{align*}
        \vec{v} &=x_1\vec{e}_1+x_2\vec{e}_2+\dots+x_n\vec{e}_n  \\
                &=(x_1+x_2+\dots+x_n)\vec{e}_1+x_2(\vec{e}_2-\vec{e}_1)+\dots
                   +x_n(\vec{e}_n-\vec{e}_1)
      \end{align*}
      这表明，若 $x_1+x_2+\dots+x_n=0$，则 $\vec{v}$ 落在
      $\vec{e}_2-\vec{e}_1$、\ldots、$\vec{e_n}-\vec{e}_1$ 的张成中
      （也就是说，落在这些向量构成的集合的张成中）；
      类似地，每个 $\sigma(\vec{v})$
      也在这个张成中，所以 $V$ 就等于这个张成，且 $\dim V=n-1$。
      另一方面，若 $x_1+x_2+\cdots+x_n\neq 0$，则上面的等式
      表明 $\vec{e}_1\in V$，从而 $\vec{e}_1,\dots,\vec{e}_n\in V$，
      所以 $V=\Re^n$ 且 $\dim V=n$。
    
\end{ans}
\protect \subsection {二.III.3: 向量空间与线性系统}
\begin{ans}{二.III.3.16}
      \begin{exparts*}
        \partsitem \( \begin{mat}[r]
                   2  &3  \\
                   1  &1
                 \end{mat}  \)
        \partsitem \( \begin{mat}[r]
                   2  &1  \\
                   1  &3
                 \end{mat}  \)
        \partsitem \( \begin{mat}[r]
                   1  &6  \\
                   4  &7  \\
                   3  &8
                 \end{mat}  \)
        \partsitem \( \rowvec{0 &0 &0} \)
        \partsitem \( \colvec[r]{-1 \\ -2} \)
      \end{exparts*}
    
\end{ans}
\begin{ans}{二.III.3.17}
       \begin{exparts}
         \partsitem 是的。
           为判断是否存在 $c_1$ 与 $c_2$ 使得
           \( c_1\cdot\rowvec{2 &1}+c_2\cdot\rowvec{3 &1}=\rowvec{1 &0} \)，
           我们求解
           \begin{equation*}
              \begin{linsys}{2}
                 2c_1  &+  &3c_2  &=  &1  \\
                  c_1  &+  &c_2   &=  &0
              \end{linsys}
           \end{equation*}
           解得 \( c_1=-1 \) 与 \( c_2=1 \)。
           因此该向量在行空间中。
         \partsitem 否。
           方程
           \( c_1\rowvec{0 &1 &3}
              +c_2\rowvec{-1 &0 &1}
              +c_3\rowvec{-1 &2 &7}
              =\rowvec{1 &1 &1} \)
           无解。
           \begin{equation*}
             \begin{amat}[r]{3}
               0  &-1  &-1  &1  \\
               1  &0   &2   &1  \\
               3  &1   &7   &1
             \end{amat}
             \grstep{\rho_1\leftrightarrow\rho_2}
             \grstep{-3\rho_1+\rho_3}
             \repeatedgrstep{\rho_2+\rho_3}
             \begin{amat}[r]{3}
               1  &0   &2   &1  \\
               0  &-1  &-1  &1  \\
               0  &0   &0   &-1
             \end{amat}
           \end{equation*}
           因此，该向量不在行空间中。
      \end{exparts}
    
\end{ans}
\begin{ans}{二.III.3.18}
       \begin{exparts}
          \partsitem 否。
            为判断是否存在 \( c_1,c_2\in\Re \) 使得
            \begin{equation*}
              c_1\colvec[r]{1 \\ 1}
              +c_2\colvec[r]{1 \\ 1}
              =\colvec[r]{1 \\ -3}
            \end{equation*}
            我们可以对由此得到的线性方程组使用高斯消元法。
            \begin{equation*}
              \begin{linsys}{2}
                 c_1  &+  &c_2  &=  &1  \\
                 c_1  &+  &c_2  &=  &3
              \end{linsys}
              \grstep{-\rho_1+\rho_2}
              \begin{linsys}{2}
                 c_1  &+  &c_2  &=  &1  \\
                      &   &0    &=  &2
              \end{linsys}
            \end{equation*}
            该方程组无解，因此该向量不在列空间中。
          \partsitem 是的。
            由关系式
            \begin{equation*}
              c_1\colvec[r]{1 \\ 2 \\ 1}
              +c_2\colvec[r]{3 \\ 0 \\ -3}
              +c_3\colvec[r]{1 \\ 4 \\ -3}
              =\colvec[r]{1 \\ 0 \\ 0}
            \end{equation*}
            我们得到一个线性方程组，对其使用高斯消元法，
            \begin{equation*}
              \begin{amat}[r]{3}
                1  &3  &1  &1  \\
                2  &0  &4  &0  \\
                1  &-3 &-3 &0
              \end{amat}
              \grstep[-\rho_1+\rho_3]{-2\rho_1+\rho_2}
              \repeatedgrstep{-\rho_2+\rho_3}
              \begin{amat}[r]{3}
                1  &3  &1  &1  \\
                0  &-6 &2  &-2 \\
                0  &0  &-6 &1
              \end{amat}
            \end{equation*}
            即得一个解。
            因此，该向量在列空间中。
% sage: load('gauss_method.sage')
% sage: M = matrix(QQ, [[1,3, 1, 1], [2, 0, 4, 0], [1, -3, -3, 0]])
% sage: M
% [ 1  3  1  1]
% [ 2  0  4  0]
% [ 1 -3 -3  0]
% sage: gauss_method(M)
% [ 1  3  1  1]
% [ 2  0  4  0]
% [ 1 -3 -3  0]
% (' take', -2, 'times row', 1, 'plus row', 2)
% (' take', -1, 'times row', 1, 'plus row', 3)
% [ 1  3  1  1]
% [ 0 -6  2 -2]
% [ 0 -6 -4 -1]
% (' take', -1, 'times row', 2, 'plus row', 3)
% [ 1  3  1  1]
% [ 0 -6  2 -2]
% [ 0  0 -6  1]
      \end{exparts}
    
\end{ans}
\begin{ans}{二.III.3.19}
      \begin{exparts}
        \partsitem 是的；
          我们问的是是否存在标量 \( c_1 \) 与 \( c_2 \) 使得
          \begin{equation*}
            c_1\colvec[r]{2 \\ 2}+c_2\colvec[r]{1 \\ 5}=\colvec[r]{1 \\ -3}
          \end{equation*}
          由此得到线性方程组
          \begin{equation*}
            \begin{linsys}{2}
              2c_1  &+  &c_2  &=  &1  \\
              2c_1  &+  &5c_2 &=  &-3
            \end{linsys}
            \grstep{-\rho_1+\rho_2}
            \begin{linsys}{2}
              2c_1  &+  &c_2  &=  &1  \\
                    &   &4c_2 &=  &-4
            \end{linsys}
          \end{equation*}
          高斯消元法给出 \( c_2=-1 \) 与 \( c_1=1 \)。
          也就是说，确实存在这样一对标量，因此该向量
          确实在矩阵的列空间中。
        \partsitem 否；
          我们问的是是否存在标量 $c_1$ 与 $c_2$
          使得
          \begin{equation*}
            c_1\colvec[r]{4 \\ 2}+c_2\colvec[r]{-8 \\ -4}=\colvec[r]{0 \\ 1}
          \end{equation*}
          一种做法是考虑由此得到的线性方程组
          \begin{equation*}
            \begin{linsys}{2}
              4c_1  &-  &8c_2  &=  &0  \\
              2c_1  &-  &4c_2  &=  &1
            \end{linsys}
          \end{equation*}
          容易看出它无解。
          另一种做法是注意到
          左边各列的任意线性组合
          其第二个分量都是第一个分量的一半，
          而右边的向量不满足这一条件。
        \partsitem 是的；我们可以直接观察到该向量
          是第一列减去第二列。
          或者，若一时看不出，可以写出各列之间的关系
          \begin{equation*}
            c_1\colvec[r]{1 \\ 1 \\ -1}
             +c_2\colvec[r]{-1 \\ 1 \\ -1}
             +c_3\colvec[r]{1 \\ -1 \\ 1}
             =\colvec[r]{2 \\ 0 \\ 0}
          \end{equation*}
          并考虑由此得到的线性方程组
          \begin{equation*}
            \begin{linsys}{3}
              c_1  &-  &c_2  &+  &c_3  &=  &2  \\
              c_1  &+  &c_2  &-  &c_3  &=  &0  \\
             -c_1  &-  &c_2  &+  &c_3  &=  &0
            \end{linsys}
            \grstep[\rho_1+\rho_3]{-\rho_1+\rho_2}
            \begin{linsys}{3}
              c_1  &-  &c_2  &+  &c_3  &=  &2  \\
                   &   &2c_2 &-  &2c_3 &=  &-2  \\
                   &   &-2c_2&+  &2c_3 &=  &2
            \end{linsys}
            \grstep{\rho_2+\rho_3}
            \begin{linsys}{3}
              c_1  &-  &c_2  &+  &c_3  &=  &2  \\
                   &   &2c_2 &-  &2c_3 &=  &-2  \\
                   &   &     &   &0    &=  &0
            \end{linsys}
          \end{equation*}
          给出了额外的信息（除了解至少存在一个之外）：有无穷多解。
          参数化给出 $c_2=-1+c_3$ 与 $c_1=1$，取 $c_3$ 为
          零即得一个特解 $c_1=1$，$c_2=-1$，
          $c_3=0$（这当然就是开头所作的观察）。
      \end{exparts}
    
\end{ans}
\begin{ans}{二.III.3.20}
     % Using Sage:
     % M = matrix(QQ, [[2,0,3,4], [0,1,1,-1], [3,1,0,2], [1,0,-4,-1]])
     % M1 = M.with_added_multiple_of_row(2,0,-3/2)
     % M2 = M1.with_added_multiple_of_row(3,0,-1/2)
     % M3 = M2.with_added_multiple_of_row(2,1,-1)
     % M4 = M3.with_added_multiple_of_row(3,2,-1)
     常规的高斯消去
     \begin{equation*}
       \begin{mat}[r]
         2  &0  &3 &4   \\
         0  &1  &1  &-1 \\
         3  &1  &0  &2  \\
         1  &0  &-4  &-1
       \end{mat}
       \grstep[-(1/2)\rho_1+\rho_4]{-(3/2)\rho_1+\rho_3}
       \grstep{-\rho_2+\rho_3}
       \repeatedgrstep{-\rho_3+\rho_4}
       \begin{mat}[r]
         2  &0  &3      &4   \\
         0  &1  &1      &-1 \\
         0  &0  &-11/2  &-3  \\
         0  &0  &0      &0
       \end{mat}
     \end{equation*}
     提示了这个基
     $\sequence{\rowvec{2 &0 &3 &4},
         \rowvec{0 &1 &1 &-1},\rowvec{0 &0 &-11/2 &-3}}$。

      另一个或许更方便的做法是先对换两行，
      % From Sage
      % M = matrix(QQ, [[2,0,3,4], [0,1,1,-1], [3,1,0,2], [1,0,-4,-1]])
      % M1 = M.with_swapped_rows(0,3)
      % M2 = M1.with_added_multiple_of_row(2,0,-3)
      % M3 = M2.with_added_multiple_of_row(3,0,-2)
      % M4 = M3.with_added_multiple_of_row(2,1,-1)
      % M5 = M4.with_added_multiple_of_row(3,2,-1)
      \begin{equation*}
        \grstep{\rho_1\swap\rho_4}
        \repeatedgrstep[-2\rho_1+\rho_4]{-3\rho_1+\rho_3}
        \repeatedgrstep{-\rho_2+\rho_3}
        \repeatedgrstep{-\rho_3+\rho_4}
        \begin{mat}[r]
           1  &0  &-4 &-1 \\
           0  &1  &1  &-1 \\
           0  &0  &11 &6  \\
           0  &0  &0  &0
         \end{mat}
      \end{equation*}
      从而得到基
      \( \sequence{\rowvec{1 &0 &-4 &-1},\rowvec{0 &1 &1 &-1},
                   \rowvec{0 &0 &11 &6}} \)。
    
\end{ans}
\begin{ans}{二.III.3.21}
      \begin{exparts}
         \partsitem 这一化简
           \begin{equation*}
             \mbox{}
             \grstep[-(1/2)\rho_1+\rho_3]{-(1/2)\rho_1+\rho_2}
             \repeatedgrstep{-(1/3)\rho_2+\rho_3}
             \begin{mat}[r]
               2  &1     &3     \\
               0  &-3/2  &1/2   \\
               0  &0     &4/3
             \end{mat}
           \end{equation*}
           表明行秩（从而秩）为三。
         \partsitem 观察各列可见其余各列都是第一列的倍数（观察各行也可看出同样的结论）。
           因此秩为一。

           或者，化简
           \begin{equation*}
             \begin{mat}[r]
               1  &-1  &2  \\
               3  &-3  &6  \\
               -2 &2   &-4
             \end{mat}
             \grstep[2\rho_1+\rho_3]{-3\rho_1+\rho_2}
             \begin{mat}[r]
               1  &-1  &2  \\
               0  &0   &0  \\
               0  &0   &0
             \end{mat}
           \end{equation*}
           也表明同样的事实。
         \partsitem 这一计算
           \begin{equation*}
             \begin{mat}[r]
               1  &3  &2  \\
               5  &1  &1  \\
               6  &4  &3
             \end{mat}
             \grstep[-6\rho_1+\rho_3]{-5\rho_1+\rho_2}
             \repeatedgrstep{-\rho_2+\rho_3}
             \begin{mat}[r]
               1  &3   &2  \\
               0  &-14 &-9 \\
               0  &0   &0
             \end{mat}
           \end{equation*}
           表明秩为二。
         \partsitem 秩为零。
       \end{exparts}
     
\end{ans}
\begin{ans}{二.III.3.22}
    我们要求这个张成空间的一个基。
    \begin{equation*}
      \spanof{\colvec{1 \\ 2 \\ 0},
              \colvec{3 \\ 1 \\ 1},
              \colvec{-1 \\ 1 \\ 1},
              \colvec{2 \\ 0 \\ 4}}\subseteq\Re^3
    \end{equation*}
    最直接的做法
    是把这些列转置成行，用高斯消元法求出这些行的张成空间的一个基，
    然后再把它们转置回
    列。
    \begin{equation*}
      \begin{mat}
        1 &2 &0 \\
        3 &1 &1 \\
       -1 &1 &1 \\
        2 &0 &4
      \end{mat}
      \grstep[\rho_1+\rho_3 \\ -2\rho_1+\rho_4]{-3\rho_1+\rho_2}
      \grstep[-(4/5)\rho_2+\rho_4]{(3/5)\rho_2+\rho_3}
      \grstep{-2\rho_3+\rho_4}
      \begin{mat}
        1 &2  &0 \\
        0 &-5 &1 \\
        0 &0  &8/5 \\
        0 &0  &0
      \end{mat}
    \end{equation*}
    舍去零向量（它表明初始向量之间存在冗余），得到列空间的这个基。
    \begin{equation*}
      \sequence{
        \colvec{1 \\ 2 \\ 0},
        \colvec{0 \\ -5 \\ 1},
        \colvec{0 \\ 0 \\ 8/5}
        }
    \end{equation*}
    矩阵的秩是其列空间的维数，所以为三。
    （它也等于其行空间的维数。）
  
\end{ans}
\begin{ans}{二.III.3.23}
      \begin{exparts}
        \partsitem 这一化简
          \begin{equation*}
             \begin{mat}[r]
               1  &3  \\
              -1  &3  \\
               1  &4  \\
               2  &1
             \end{mat}
             \grstep[-\rho_1+\rho_3 \\ -2\rho_1+\rho_4]{\rho_1+\rho_2}
             \grstep[(5/6)\rho_2+\rho_4]{-(1/6)\rho_2+\rho_3}
             \begin{mat}[r]
               1  &3  \\
               0  &6  \\
               0  &0  \\
               0  &0
             \end{mat}
          \end{equation*}
          给出 \( \sequence{\rowvec{1 &3},\rowvec{0 &6}} \)。
        \partsitem 转置并化简
          \begin{equation*}
             \begin{mat}[r]
               1  &2  &1  \\
               3  &1  &-1 \\
               1  &-3 &-3
             \end{mat}
             \grstep[-\rho_1+\rho_3]{-3\rho_1+\rho_2}
             \begin{mat}[r]
               1  &2  &1  \\
               0  &-5 &-4 \\
               0  &-5 &-4
             \end{mat}
             \grstep{-\rho_2+\rho_3}
             \begin{mat}[r]
               1  &2  &1  \\
               0  &-5 &-4 \\
               0  &0  &0
             \end{mat}
          \end{equation*}
          再转置回去，得到这个基。
          \begin{equation*}
            \sequence{\colvec[r]{1 \\ 2 \\ 1},
              \colvec[r]{0 \\ -5 \\ -4}   }
          \end{equation*}
        \partsitem 首先注意外围空间是
          $\polyspace_3$，而不是 $\polyspace_2$。
          然后，把第一个多项式 $1+1\cdot x+0\cdot x^2+0\cdot x^3$
          取作与行向量 $\rowvec{1 &1 &0 &0}$“相同”，等等，
          就得到
          \begin{equation*}
            \begin{mat}[r]
              1  &1  &0  &0 \\
              1  &0  &-1 &0 \\
              3  &2  &-1 &0
            \end{mat}
            \grstep[-3\rho_1+\rho_3]{-\rho_1+\rho_2}
            \repeatedgrstep{-\rho_2+\rho_3}
            \begin{mat}[r]
              1  &1  &0  &0 \\
              0  &-1 &-1 &0 \\
              0  &0  &0  &0
            \end{mat}
          \end{equation*}
          由此得到基 \( \sequence{1+x,-x-x^2} \)。
        \partsitem 这里“相同”的做法给出
          \begin{equation*}
            \begin{mat}[r]
              1  &0  &1  &3  &1  &-1  \\
              1  &0  &3  &2  &1  &4   \\
             -1  &0  &-5  &-1 &-1 &-9
            \end{mat}
            \grstep[\rho_1+\rho_3]{-\rho_1+\rho_2}
            \repeatedgrstep{2\rho_2+\rho_3}
            \begin{mat}[r]
              1  &0  &1  &3  &1  &-1  \\
              0  &0  &2  &-1 &0  &5   \\
              0  &0  &0  &0  &0  &0
            \end{mat}
          \end{equation*}
          从而得到这个基。
          \begin{equation*}
            \sequence{
            \begin{mat}[r]
              1  &0  &1  \\
              3  &1  &-1
            \end{mat},
            \begin{mat}[r]
              0  &0  &2  \\
             -1  &0  &5
            \end{mat}  }
          \end{equation*}
      \end{exparts}
     
\end{ans}
\begin{ans}{二.III.3.24}
      \begin{exparts}
        \partsitem
          把这些列转置成行，化为阶梯形
          （然后去掉零行），再转置回列。
          \begin{equation*}
            \begin{mat}
              1 &1  &3 \\
             -1 &2  &0 \\
              0 &12 &6
            \end{mat}
            \grstep{\rho_1+\rho_2}
            \grstep{-4\rho_2+\rho_3}
            \begin{mat}
              1 &1  &3 \\
              0 &3  &3 \\
              0 &0  &-6
            \end{mat}
          \end{equation*}
          该张成空间的一个基如下。
          \begin{equation*}
            \sequence{
              \colvec{1 \\ 1 \\ 3},
              \colvec{0 \\ 3 \\ 3},
              \colvec{0 \\ 0 \\ -6}
          }
          \end{equation*}
       \partsitem
          与前一小区一样，我们把它们看作行，
          以利用高斯消元法已完成的工作。
          \begin{multline*}
            \begin{mat}
              0 &1  &1 \\
              2 &-2 &0 \\
              7 &0 &0 \\
              4 &3  &2
            \end{mat}
            \grstep{\rho_1\leftrightarrow\rho_2}
            \grstep[2\rho_1+\rho_4]{-(7/2)\rho_1+\rho_3}
            \grstep[-7\rho_2+\rho_4]{-7\rho_2+\rho_3}          \\
            \grstep{-(5/7)\rho_3+\rho_4}
            \begin{mat}
              2 &-2  &0 \\
              0 &1 &1 \\
              0 &0  &-7 \\
              0 &0 &0
            \end{mat}
          \end{multline*}
          该集合的张成空间的一个基是
          $\sequence{2-2x, x+x^2, -5x^2}$。
      \end{exparts}
    
\end{ans}
\begin{ans}{二.III.3.25}
      只有零矩阵的秩为零。
      秩为一的矩阵仅有如下形式
      \begin{equation*}
        \begin{mat}
           k_1\cdot \rho  \\
            \vdots  \\
           k_m\cdot \rho
        \end{mat}
      \end{equation*}
      其中 \( \rho \) 是某个非零行向量，且诸 \( k_i \) 不全为零。
      （\textit{注}：
      我们不能简单地说所有行都是第一行的倍数，因为第一行可能是零行。
      \textit{再注}：
      上述结论在把“行”换成“列”后同样适用。）
    
\end{ans}
\begin{ans}{二.III.3.26}
      若 \( a\neq 0 \)，则取 \( d=(c/a)b \) 可使第二行是第一行的倍数，
      具体地说，是第一行的 \( c/a \) 倍。
      若 \( a=0 \) 且 \( b=0 \)，则 \( d \) 取任意非 \( 0 \) 值都能保证
      第二行非零。
      若 \( a=0 \) 且 \( b\neq 0 \) 且 \( c=0 \)，则 \( d \) 取任何值都行，
      因为该矩阵自动具有秩一（即使取 $d=0$）。
      最后，若 \( a=0 \) 且 \( b\neq 0 \) 且 \( c\neq 0 \)，则 \( d \) 取
      任何值都不够，因为该矩阵必定有秩二。
    
\end{ans}
\begin{ans}{二.III.3.27}
      列秩是二。
      一种看法是直接观察\Dash 列空间由高度为二的列组成，因而其维数
      至少可以是二，而且我们容易找到两列，它们合起来组成一个
      线性无关的集合（例如第四列与第五列）。
      另一种看法是回忆列秩等于行秩，然后施行高斯消元法，
      结果剩下两个非零行。
    
\end{ans}
\begin{ans}{二.III.3.28}
      我们应用\nearbytheorem{th:RankVsSoltnSp}。
      线性方程组的系数矩阵 $A$ 的列数等于未知数的个数 $n$。
      一个至少有一个解的线性方程组至多有一个解当且仅当相应的
      齐次方程组的解空间维数为零（回忆：~在
      “$\text{通解}=\text{特解}+\text{齐次解}$”这一方程
       $\vec{v}=\vec{p}+\vec{h}$ 中，只要这样的 $\vec{p}$ 存在，
       解 $\vec{v}$ 唯一当且仅当向量 $\vec{h}$ 唯一，即 $\vec{h}=\zero$）。
       而由该定理，这意味着 $n=r$。
    
\end{ans}
\begin{ans}{二.III.3.29}
      列的集合必定是线性相关的，因为该矩阵的秩至多是五，而列却有
      九个。
    
\end{ans}
\begin{ans}{二.III.3.30}
      它们几乎没有相等的可能，因为行空间是行向量的集合，而列空间是
      列向量的集合（除非该矩阵是 \( \nbyn{1} \) 的，此时两个空间
      必定相等）。

      \textit{注记。}
      考察
      \begin{equation*}
        A=\begin{mat}[r]
            1  &3  \\
            2  &6
          \end{mat}
      \end{equation*}
      并注意行空间是 \( \rowvec{1 &3} \) 的所有倍数组成的集合，
      而列空间由以下列向量的倍数组成
      \begin{equation*}
        \colvec[r]{1 \\ 2}
      \end{equation*}
      所以我们也不能论证说这两个空间必定恰好互为转置。
    
\end{ans}
\begin{ans}{二.III.3.31}
      首先，该向量空间是在自然运算下的实数四元组的集合。
      尽管这不是四宽的行向量的集合，但差别不大\Dash 它与那个集合
      是“相同的”。
      所以我们把四元组当作四宽的向量来处理。

      有了这一点，要看出 \( (1,0,1,0) \) 不在第一个集合的张成中，
      一种办法是注意这个化简
      \begin{equation*}
        \begin{mat}[r]
          1  &-1  &2  &-3  \\
          1  &1   &2  &0   \\
          3  &-1  &6  &-6
        \end{mat}
        \grstep[-3\rho_1+\rho_3]{-\rho_1+\rho_2}
        \repeatedgrstep{-\rho_2+\rho_3}
        \begin{mat}[r]
          1  &-1  &2  &-3  \\
          0  &2   &0  &3   \\
          0  &0   &0  &0
        \end{mat}
      \end{equation*}
      以及这个化简
      \begin{equation*}
        \begin{mat}[r]
          1  &-1  &2  &-3  \\
          1  &1   &2  &0   \\
          3  &-1  &6  &-6  \\
          1  &0   &1  &0
        \end{mat}
        \grstep[-3\rho_1+\rho_3 \\ -\rho_1+\rho_4]{-\rho_1+\rho_2}
        \repeatedgrstep[-(1/2)\rho_2+\rho_4]{-\rho_2+\rho_3}
        \repeatedgrstep{\rho_3\leftrightarrow\rho_4}
        \begin{mat}[r]
          1  &-1  &2  &-3  \\
          0  &2   &0  &3   \\
          0  &0   &-1 &3/2 \\
          0  &0   &0  &0
        \end{mat}
      \end{equation*}
      得到的矩阵在秩上不同。
      这意味着把 $(1,0,1,0)$ 添加到前三个四元组的集合中会增大该
      集合的秩，从而增大其张成。
      因此 $(1,0,1,0)$ 原本就不在这个张成中。
    
\end{ans}
\begin{ans}{二.III.3.32}
      它是子空间，因为它是系数矩阵
      \begin{equation*}
        \begin{mat}[r]
          3  &2  &4  \\
          1  &0  &-1 \\
          2  &2  &5
        \end{mat}
      \end{equation*}
      的列空间。
      为了求列空间的一个基，
      \begin{equation*}
        \set{c_1\colvec[r]{3 \\ 1 \\ 2}
             +c_2\colvec[r]{2 \\ 0 \\ 2}
             +c_3\colvec[r]{4 \\ -1 \\ 5}
             \suchthat c_1,c_2,c_3\in\Re}
      \end{equation*}
      我们通过转置、化简
      \begin{equation*}
         \begin{mat}[r]
           3  &1  &2  \\
           2  &0  &2  \\
           4  &-1 &5
         \end{mat}
         \grstep[-(4/3)\rho_1+\rho_3]{-(2/3)\rho_1+\rho_2}
         \repeatedgrstep{-(7/2)\rho_2+\rho_3}
         \begin{mat}[r]
           3  &1     &2  \\
           0  &-2/3  &2/3  \\
           0  &0     &0
         \end{mat}
      \end{equation*}
      并略去零行，再转置回去，来消去张成集中三个列向量之间的
      线性关系。
      \begin{equation*}
        \sequence{ \colvec[r]{3 \\ 1 \\ 2},
                   \colvec[r]{0 \\ -2/3 \\ 2/3}  }
      \end{equation*}
    
\end{ans}
\begin{ans}{二.III.3.33}
      我们可以通过直接的计算来完成。
      \begin{align*}
        \trans{(rA+sB)}
        &=\trans{\begin{mat}
             ra_{1,1}+sb_{1,1}  &\ldots &ra_{1,n}+sb_{1,n} \\
                          &\vdots            \\
             ra_{m,1}+sb_{m,1}  &\ldots &ra_{m,n}+sb_{m,n}
                \end{mat}  }  \\
        &=\begin{mat}
           ra_{1,1}+sb_{1,1}  &\ldots &ra_{m,1}+sb_{m,1}  \\
                            &\vdots                           \\
           ra_{1,n}+sb_{1,n}  &\ldots &ra_{m,n}+sb_{m,n}
         \end{mat}           \\
        &=\begin{mat}
           ra_{1,1}  &\ldots &ra_{m,1}  \\
                    &\vdots                       \\
           ra_{1,n}  &\ldots &ra_{m,n}
         \end{mat}
        +\begin{mat}
           sb_{1,1}  &\ldots &sb_{m,1}  \\
                    &\vdots                       \\
           sb_{1,n}  &\ldots &sb_{m,n}
         \end{mat}           \\
        &=r\trans{A}+s\trans{B}
      \end{align*}
    
\end{ans}
\begin{ans}{二.III.3.34}
      \begin{exparts}
        \partsitem 这些化简给出不同的基。
          \begin{equation*}
            \begin{mat}[r]
              1  &2  &0  \\
              1  &2  &1
            \end{mat}
            \grstep{-\rho_1+\rho_2}
            \begin{mat}[r]
              1  &2  &0  \\
              0  &0  &1
            \end{mat}
            \qquad
            \begin{mat}[r]
              1  &2  &0  \\
              1  &2  &1
            \end{mat}
            \grstep{-\rho_1+\rho_2}
            \repeatedgrstep{2\rho_2}
            \begin{mat}[r]
              1  &2  &0  \\
              0  &0  &2
            \end{mat}
          \end{equation*}
        \partsitem 一个简单的例子是这个。
          \begin{equation*}
            \begin{mat}[r]
              1  &2  &1  \\
              3  &1  &4
            \end{mat}
            \qquad
            \begin{mat}[r]
              1  &2  &1  \\
              3  &1  &4  \\
              0  &0  &0
            \end{mat}
          \end{equation*}
          这是一个稍微不那么简单的例子。
          \begin{equation*}
            \begin{mat}[r]
              1  &2  &1  \\
              3  &1  &4
            \end{mat}
            \qquad
            \begin{mat}[r]
              1  &2  &1  \\
              3  &1  &4  \\
              2  &4  &2  \\
              4  &3  &5
            \end{mat}
          \end{equation*}
        \partsitem
          % We give two proofs.
          % The first is a bit lower-level but more constructive.
          % The second is higher level.

          % \textit{First proof.}
          % Assume that \( A \) and \( B \) are matrices with equal row spaces.
          % Construct a matrix \( C \) with the rows of \( A \) above the rows
          % of \( B \), and another matrix \( D \) with the rows of \( B \)
          % above the rows of \( A \).
          % \begin{equation*}
          %   C=\begin{mat}
          %        A \\ B
          %     \end{mat}
          %   \qquad
          %   D=\begin{mat}
          %        B \\ A
          %     \end{mat}
          % \end{equation*}
          % Observe that \( C \) and \( D \) are row-equivalent (via a sequence
          % of row-swaps) and so Gauss-Jordan reduce to the same reduced
          % echelon form matrix.

          % Because the row spaces are equal, the rows of \( B \)
          % are linear combinations of the rows of \( A \) so Gauss-Jordan
          % reduction on \( C \) simply turns the rows of \( B \) to zero rows
          % and thus the nonzero rows of $C$ are just the nonzero rows obtained
          % by Gauss-Jordan reducing \( A \).
          % The same can be said for the matrix \( D \)\Dash Gauss-Jordan
          % reduction on \( D \) gives the same non-zero rows as are produced
          % by reduction on \( B \) alone.
          % Therefore,
          % \( A \) yields the same nonzero rows as \( C  \), which yields
          % the same nonzero rows as \( D \), which yields the same nonzero
          % rows as \( B \).

          % \textit{Second proof (N~Schenker).}
          由于 \( A \) 与 \( B \) 的行空间相等，所以二者行等价。
          由于每个行等价类都包含唯一的简化阶梯形
          （定理 One.III.\ref{th:ReducedEchelonFormIsUnique}），
          \( A \) 的简化阶梯形必定等于 \( B \) 的简化阶梯形。
      \end{exparts}
    
\end{ans}
\begin{ans}{二.III.3.35}
      它不可能更大。
    
\end{ans}
\begin{ans}{二.III.3.36}
      极大线性无关集合中行的个数不会超过行的个数。
      一个更好的界（一般情形下这是最佳可能的界）是 \( m \) 与 \( n \)
      中的最小者，因为行秩等于列秩。
     
\end{ans}
\begin{ans}{二.III.3.37}
      由于矩阵 \( A \) 的行是 \( \trans{A} \) 的列，所以 \( A \) 的
      行空间的维数等于 \( \trans{A} \) 的列空间的维数。
      而 \( A \) 的行空间的维数就是 \( A \) 的秩，并且 \( \trans{A} \) 的
      列空间的维数就是 \( \trans{A} \) 的秩。
      于是这两个秩相等。
    
\end{ans}
\begin{ans}{二.III.3.38}
      假的。
      前者是列的集合，而后者是行的集合。

      然而这个例子
      \begin{equation*}
         A=\begin{mat}[r]
             1  &2  &3  \\
             4  &5  &6
           \end{mat},
         \qquad
         \trans{A}=\begin{mat}[r]
             1  &4  \\
             2  &5  \\
             3  &6
           \end{mat}
      \end{equation*}
      表明，一旦我们对“相同”有了形式化的含义，就可以在这里应用它：
      \begin{equation*}
        \colspace{A}=\spanof{\set{
                        \colvec[r]{1 \\ 4},
                        \colvec[r]{2 \\ 5},
                        \colvec[r]{3 \\ 6} }  }
      \end{equation*}
      而
      \begin{equation*}
         \rowspace{\trans{A}}=\spanof{\set{
                                 \rowvec{1 &4},
                                 \rowvec{2 &5},
                                 \rowvec{3 &6} }  }
      \end{equation*}
      彼此是“相同的”。
    
\end{ans}
\begin{ans}{二.III.3.39}
      不会。
      这里，高斯消元法没有改变列空间。
      \begin{equation*}
        \begin{mat}[r]
          1  &0  \\
          3  &1
        \end{mat}
        \grstep{-3\rho_1+\rho_2}
        \begin{mat}[r]
          1  &0  \\
          0  &1
        \end{mat}
      \end{equation*}
    
\end{ans}
\begin{ans}{二.III.3.40}
      线性方程组
      \begin{equation*}
        c_1\vec{a}_1+\dots+c_n\vec{a}_n=\vec{d}
      \end{equation*}
      有解当且仅当 \( \vec{d} \) 在集合 \( \set{\vec{a}_1,\dots,\vec{a}_n} \)
      的张成中。
      这成立当且仅当增广矩阵的列秩等于系数矩阵的列秩。
      由于秩等于列秩，方程组有解当且仅当其增广矩阵的秩等于其
      系数矩阵的秩。
    
\end{ans}
\begin{ans}{二.III.3.41}
      \begin{exparts}
        \partsitem 行秩等于列秩，所以二者各自至多是行数与列数的
          最小者。
          因此只有当行数等于列数时二者才能都满。
          （当然，逆命题不成立：~方阵不必行满秩或列满秩。）
        \partsitem 若 \( A \) 行满秩，则无论右侧是什么，对增广矩阵施行
          高斯消元法结束时每一行都有一个首主元一，并且这些首主元一
          没有一个位于最右边的列（“增广”列）。
          回代随即给出一个解。

          另一方面，如果该线性方程组对某个右侧没有解，只可能是因为
          高斯消元法留下某行，使得它在“增广”竖线左边全为零，
          而在右边有非零元。
          于是，如果 $A$ 对某些右侧没有解，那么 \( A \) 不具有
          行满秩，因为它的某些行已被消去。
        \partsitem 矩阵 \( A \) 列满秩当且
          仅当其列
          组成一个线性无关的集合。
          这等价于诸列之间只存在平凡的线性关系，于是方程组唯一的解是
          每个变量都取 $0$ 的情形。
        \partsitem 矩阵 \( A \) 列满秩当且
          仅当其列的集合是线性无关的，从而构成其张成的一个基。
          这等价于该张成中的所有向量都有唯一的线性表示。
          这就证明了它，
          因为张成中一个向量的任何线性表示都是
          该线性方程组的解。
      \end{exparts}
    
\end{ans}
\begin{ans}{二.III.3.42}
      行空间不再相同，$B$ 的行空间将是 $A$ 的行空间的一个子空间
      （有可能相等）。
    
\end{ans}
\begin{ans}{二.III.3.43}
      显然 \( \rank(A)=\rank(-A) \)，因为高斯消元法允许我们把一个
      矩阵的所有行乘以 \( -1 \)。
      同样地，当 \( k\neq 0 \) 时我们有 \( \rank(A)=\rank(kA) \)。

      加法更有意思。
      和的秩可以比各加数的秩小。
      \begin{equation*}
        \begin{mat}[r]
          1  &2  \\
          3  &4
        \end{mat}
        +
        \begin{mat}[r]
         -1  &-2  \\
         -3  &-4
        \end{mat}
        =
        \begin{mat}[r]
         0   &0   \\
         0   &0
        \end{mat}
      \end{equation*}
      和的秩也可以比各加数的秩大。
      \begin{equation*}
        \begin{mat}[r]
          1  &2  \\
          0  &0
        \end{mat}
        +
        \begin{mat}[r]
          0  &0   \\
          3  &4
        \end{mat}
        =
        \begin{mat}[r]
         1   &2   \\
         3   &4
        \end{mat}
      \end{equation*}
      但存在一个上界（除矩阵尺寸之外的上界）。
      一般地，\( \rank(A+B)\leq\rank(A)+\rank(B) \)。

      为了证明它，注意我们可以用两种方式之一对 \( A+B \) 实施
      高斯消去法：
      可以先把 \( A \) 加到 \( B \) 上，然后施行适当的化简步骤序列
      \begin{equation*}
        (A+B)
        \grstep{\text{步骤}_1}
        \cdots
        \grstep{\text{步骤}_k}
        \text{阶梯形}
      \end{equation*}
      或者也可以分别在 \( A \) 与 \( B \) 上实施 \( \text{步骤}_1 \) 到
      \( \text{步骤}_k \)，然后再相加，得到同样的结果。
      在第二种方式中我们最终能得到的最大秩显然是两个秩之和。
      （上面的矩阵给出了两种可能性，
       $\rank(A+B)<\rank(A)+\rank(B)$ 与 $\rank(A+B)=\rank(A)+\rank(B)$，
       都发生的例子。）
    
\end{ans}
\protect \subsection {二.III.4: 子空间的组合}
\begin{ans}{二.III.4.20}
       对这里的每一项我们都可以应用\nearbylemma{le:DirectSumTwoSp}。
       \begin{exparts}
         \partsitem 是的。
           该平面是这个 $W_1$ 与 $W_2$ 的和，因为对任意标量 $a$ 与 $b$，
           \begin{equation*}
             \colvec{a \\ b}
             =\colvec{a-b \\ 0}
             +\colvec{b \\ b}
           \end{equation*}
           表明任一向量都是来自这两部分的向量之和。
           而且，这两个子空间是（不同的）过原点的直线，
           因此它们的交是平凡的。
         \partsitem 是的。
           为看出平面中任一向量都是这些部分中向量的组合，
           考察下面这个关系式。
           \begin{equation*}
             \colvec{a \\ b}=c_1\colvec[r]{1 \\ 1}+c_2\colvec[r]{1 \\ 1.1}
           \end{equation*}
           我们现在只需注意，集合
           \begin{equation*}
             \set{\colvec[r]{1 \\ 1},\colvec[r]{1 \\ 1.1}}
           \end{equation*}
           是该空间的一组基（因为它显然线性无关，且在 $\Re^2$ 中含有两个元素），
           因此上述方程有且仅有一解，这意味着所有分解都是唯一的。
           另一种做法是求解
           \begin{equation*}
             \begin{linsys}{2}
               c_1  &+  &c_2    &=  &a  \\
               c_1  &+  &1.1c_2 &=  &b
             \end{linsys}
             \grstep{-\rho_1+\rho_2}
             \begin{linsys}{2}
               c_1  &+  &c_2    &=  &a  \\
                    &   &0.1c_2 &=  &-a+b
             \end{linsys}
           \end{equation*}
           得到 $c_2=10(-a+b)$ 与 $c_1=11a-10b$，于是我们有
           \begin{equation*}
             \colvec{a \\ b}=\colvec{11a-10b \\ 11a-10b}
                             +\colvec{-10a+10b \\ 1.1\cdot(-10a+10b)}
           \end{equation*}
           正如所求。
           与前一个答案一样，这两个子空间各自都是过原点的直线，
           它们的交是平凡的。
         \partsitem 是的。
           平面中的每个向量都可如下分解为和
           \begin{equation*}
             \colvec{x \\ y}=\colvec{x \\ y}+\colvec[r]{0 \\ 0}
           \end{equation*}
           且这两个子空间的交是平凡的。
         \partsitem 不是。
           交不是平凡的。
         \partsitem 不是。
           它们不是子空间。
      \end{exparts}
    
\end{ans}
\begin{ans}{二.III.4.21}
      对这里的每一项我们都可以使用\nearbylemma{le:DirectSumTwoSp}。
      \begin{exparts}
        \partsitem \( \Re^3 \) 中的任意向量都可以分解为这个和。
          \begin{equation*}
            \colvec{x \\ y \\ z}
            =\colvec{x \\ y \\ 0}
            +\colvec{0 \\ 0 \\ z}
          \end{equation*}
          而且，$xy$ 平面与 $z$ 轴的交是平凡子空间。
        \partsitem \( \Re^3 \) 中的任意向量都可以分解为
          \begin{equation*}
            \colvec{x \\ y \\ z}
            =\colvec{x-z \\ y-z \\ 0}
            +\colvec{z \\ z \\ z}
          \end{equation*}
          且这两个空间的交是平凡的。
      \end{exparts}
    
\end{ans}
\begin{ans}{二.III.4.22}
      是。
      证明这两个集合是子空间是常规的。
      为看出整个空间是这两个的直和，只需注意
      \( \polyspace_2 \) 中每个元素都有唯一的分解
      \( m+nx+px^2=(m+px^2)+(nx) \)。
     
\end{ans}
\begin{ans}{二.III.4.23}
      证明它们是子空间是常规的。
      我们将利用\nearbylemma{le:DirectSumTwoSp}来论证它们互补。
      交 \( \mathcal{E}\intersection\mathcal{O} \) 是平凡的，
      因为同时满足条件 $p(-x)=p(x)$ 与
      $p(-x)=-p(x)$ 的多项式只有零多项式。
      为看出整个空间是这些子空间的和
      \( \mathcal{E}+\mathcal{O}=\polyspace_n \)，
      注意多项式
      \( p_0(x)=1 \)、\( p_2(x)=x^2 \)、\( p_4(x)=x^4 \) 等都在
      \( \mathcal{E} \) 中，同时注意多项式
      \( p_1(x)=x \)、\( p_3(x)=x^3 \) 等都在 \( \mathcal{O} \) 中。
      于是 \( \polyspace_n \) 的任何元素都是
      \( \mathcal{E} \) 与 \( \mathcal{O} \) 中元素的组合。
    
\end{ans}
\begin{ans}{二.III.4.24}
     这里的每一个都是 $\Re^3$。
     \begin{exparts}
      \partsitem
        为便于阅读，这里把它们分成若干行列出。
        \begin{center}
          \begin{tabular}{l}
            $W_1+W_2+W_3$,
              $W_1+W_2+W_3+W_4$,
              $W_1+W_2+W_3+W_5$,  \\ \quad %line too long
              $W_1+W_2+W_3+W_4+W_5$,
              $W_1+W_2+W_4$,
              $W_1+W_2+W_4+W_5$,     \\ \quad %line too long
              $W_1+W_2+W_5$,
              $W_1+W_3+W_4$,
              $W_1+W_3+W_5$,
              $W_1+W_3+W_4+W_5$,     \\ \quad
              $W_1+W_4$,
              $W_1+W_4+W_5$,
              $W_1+W_5$,             \\
              $W_2+W_3+W_4$,
              $W_2+W_3+W_4+W_5$,
              $W_2+W_4$,
              $W_2+W_4+W_5$,         \\
              $W_3+W_4$,
              $W_3+W_4+W_5$,         \\
              $W_4+W_5$
          \end{tabular}
         \end{center}
     \partsitem
       $W_1\directsum W_2\directsum W_3$,
       $W_1\directsum W_4$,
       $W_1\directsum W_5$,
       $W_2\directsum W_4$,
       $W_3\directsum W_4$
     \end{exparts}
    
\end{ans}
\begin{ans}{二.III.4.25}
      显然每一个都是子空间。
      这些子空间的基 \( B_i=\sequence{x^i} \) 依次连接起来，
      就构成整个空间的一组基。
    
\end{ans}
\begin{ans}{二.III.4.26}
       是 \( W_2 \)。
    
\end{ans}
\begin{ans}{二.III.4.27}
      由\nearbylemma{le:UniqDecIffBasisDec}知其为真。
    
\end{ans}
\begin{ans}{二.III.4.28}
      \( \Re^4 \) 的两个不同的直和分解很容易找到。
      其一是
      \( W_1=\spanof{\set{\vec{e}_1,\vec{e}_2}} \) 与
      \( W_2=\spanof{\set{\vec{e}_3,\vec{e}_4}} \)，
      其二是
      \( U_1=\spanof{\set{\vec{e}_1}} \) 与
      \( U_2=\spanof{\set{\vec{e}_2,\vec{e}_3,\vec{e}_4}} \)。
      （还有很多种可能，
      例如 \( \Re^4 \) 与它的平凡子空间。）

      相比之下，对 \( \Re^1 \) 的单向量基的任何划分，
      都会给出一组没有元素的基和另一组只有一个元素的基。
      因此任何分解都涉及 \( \Re^1 \) 与它的平凡
      子空间。
    
\end{ans}
\begin{ans}{二.III.4.29}
        一个方向的包含关系是容易的：
        \( \set{\vec{w}_1+\dots+\vec{w}_k\suchthat \vec{w}_i\in W_i} \)
        是 \( \spanof{W_1\union \dots \union W_k}  \) 的子集，
        因为每个 \( \vec{w}_1+\dots+\vec{w}_k \) 都是这个并集中向量的和。

        对于另一个方向的包含，对来自该并集的向量的任一线性组合，
        利用向量加法的交换性，
        把来自 \( W_1 \) 的向量放在最前面，
        接着是来自 \( W_2 \) 的向量，等等。
        把来自 \( W_1 \) 的向量相加得到一个 \( \vec{w}_1\in W_1 \)，
        把来自 \( W_2 \) 的向量相加得到一个 \( \vec{w}_2\in W_2 \)，等等。
        所得结果具有所要求的形式。
      
\end{ans}
\begin{ans}{二.III.4.30}
      一个例子是取空间为 $\Re^3$，
      并取子空间为 $xy$ 平面、$xz$ 平面与 $yz$ 平面。
    
\end{ans}
\begin{ans}{二.III.4.31}
      当然，零向量在所有子空间中，所以
      交至少包含这一个向量。
      由直和的定义，集合
      \( \set{W_1,\dots,W_k} \) 是无关的，因此当
      \( i\neq j \) 时，\( W_i \) 中没有非零向量
      是 \( W_j \) 中某个元素的倍数。
      特别是，\( W_i \) 中没有非零向量等于
      \( W_j \) 中的元素。
    
\end{ans}
\begin{ans}{二.III.4.32}
      它可以包含一个平凡子空间；
      \( \Re^3 \) 的下面这个子空间集合是无关的：
      \( \set{\set{\zero},\text{\( x \)-轴}} \)。
      平凡空间 \( \set{\zero} \) 中没有非零向量
      是 \( x \) 轴中某个
      向量的倍数，这仅仅是因为平凡空间
      没有可以作为这种倍数候选者的非零向量（而且
      \( x \) 轴中也没有非零向量是平凡子空间中
      零向量的倍数）。
     
\end{ans}
\begin{ans}{二.III.4.33}
      有。
      对向量空间的任何子空间，我们可以取该
      子空间的任意一组基 \( \sequence{\vec{\omega}_1,\dots,\vec{\omega}_k} \)，
      并把它扩充成整个
      空间的一组基
      \( \sequence{\vec{\omega}_1,\dots,\vec{\omega}_k,
                     \vec{\beta}_{k+1},\dots,\vec{\beta}_n} \)。
      于是原子空间的补具有这组基
      \( \sequence{\vec{\beta}_{k+1},\dots,\vec{\beta}_n} \)。
     
\end{ans}
\begin{ans}{二.III.4.34}
      \begin{exparts}
        \partsitem 它必定张成。
          我们可以把 \( W_1+W_2 \) 的任何成员写成 \( \vec{w}_1+\vec{w}_2 \)，
          其中 \( \vec{w}_1\in W_1 \) 且 \( \vec{w}_2\in W_2 \)。
          由于 \( S_1 \) 张成 \( W_1 \)，向量 \( \vec{w}_1 \) 是
          \( S_1 \) 中成员的组合。
          类似地，\( \vec{w}_2 \) 是 \( S_2 \) 中成员的组合。
        \partsitem 看出它可以是线性无关的一个简单方法
          是取两者都为空集。
          另一方面，在空间 $\Re^1$ 中，
          若 \( W_1=\Re^1 \) 且 \( W_2=\Re^1 \)，而 $S_1=\set{1}$，
          $S_2=\set{2}$，则它们的并 $S_1\union S_2$ 不是无关的。
      \end{exparts}
     
\end{ans}
\begin{ans}{二.III.4.35}
      \begin{exparts}
        \partsitem 交与和是
          \begin{equation*}
             \set{\begin{mat}
                    0  &0  \\
                    c  &0
                  \end{mat} \suchthat c\in\Re  }
             \qquad
             \set{\begin{mat}
                    0  &b  \\
                    c  &d
                   \end{mat} \suchthat b,c,d\in\Re  }
          \end{equation*}
         其维数分别为一与三。
      \partsitem 我们把 $U\intersection W$ 的基记为
        $B_{U\intersection W}$，
        把 $U$ 的基记为 $B_U$，
        把 $W$ 的基记为 $B_W$，
        而把正在考虑的基记为 $B_{U+W}$。

        为看出 $B_{U+W}$ 张成 $U+W$，注意我们可以把
        $U+W$ 中的任一向量 $c\vec{u}+d\vec{w}$ 写成
        $B_{U+W}$ 中向量的线性
        组合，
        只需把 $\vec{u}$ 用
        $B_U$ 表示、把 $\vec{w}$ 用 $B_W$ 表示即可。

         我们最后证明 $B_{U+W}$ 线性无关。
         考察
         \begin{equation*}
           c_1\vec{\mu}_1+\dots+c_{j+1}\vec{\beta}_1+\dots
              +c_{j+k+p}\vec{\omega}_p=\zero
         \end{equation*}
         它可以如下改写。
         \begin{equation*}
           c_1\vec{\mu}_1+\dots+c_j\vec{\mu}_j=
             -c_{j+1}\vec{\beta}_1-\dots-c_{j+k+p}\vec{\omega}_p
         \end{equation*}
         注意左边求和得到 \( U \) 中的向量，而
         右边求和得到 \( W \) 中的向量，于是两边求和都得到
         $U\intersection W$ 的一个成员。
         由于左边是 $U\intersection W$ 的成员，
         它可以用 $B_{U\intersection W}$ 的成员表示，
         这给出：上式左边 $\vec{\mu}$ 的组合
         等于 $\vec{\beta}$ 的一个组合。
         但是，基 $B_U$ 线性无关这一事实表明
         任何这样的组合
         都是平凡的，特别地，上面左边的系数 $c_1$, \ldots,
         $c_j$ 全为零。
         类似地，\( \vec{\omega} \) 的系数也全为零。
         这使上面的方程成为
         \( \vec{\beta} \) 之间的一个线性关系，
         但 $B_{U\intersection W}$ 线性无关，
         因此 $\vec{\beta}$ 的所有系数也都
         为零。
        \partsitem 只需数一数上一项中的
          基向量：~$\dim(U+W)=j+k+p$，而 $\dim(U)=j+k$，$\dim(W)=k+p$，
          以及 $\dim(U\intersection W)=k$。
        \partsitem  我们知道
          \( \dim(W_1+W_2)=\dim(W_1)+\dim(W_2)-\dim(W_1\intersection W_2) \)。
          因为 \( W_1\subseteq W_1+W_2 \)，
          我们知道 \( W_1+W_2 \) 的维数必定大于
          $W_1$ 的维数，也就是说，其维数必定是八、九或十。
          代入得到三种可能：
          $8=8+8-\dim(W_1\intersection W_2)$、
          $9=8+8-\dim(W_1\intersection W_2)$ 或
          $10=8+8-\dim(W_1\intersection W_2)$。
          因此 \( \dim(W_1\intersection W_2) \) 必定是
          八、七或六。
          （要举例说明这三种情形中的每一种都可能
          出现是容易的，例如在 \( \Re^{10} \) 中。）
      \end{exparts}
    
\end{ans}
\begin{ans}{二.III.4.36}
      把每个 \( S_i \) 扩充成 \( W_i \) 的一组基 $B_i$。
      这些基的连接 $\cat{\cat{B_1}{\cdots}}{B_k}$
      是 \( V \) 的一组基，因此其中的成员构成一个线性无关
      集合。
      但并集 $S_1\union\cdots\union S_k$ 是这个线性无关集合的
      子集，因此它自身
      也线性无关。
     
\end{ans}
\begin{ans}{二.III.4.37}
      \begin{exparts}
        \partsitem 这样的两个矩阵如下。
          \begin{equation*}
            \begin{mat}[r]
              1  &2  \\
              2  &3
            \end{mat}
            \qquad
            \begin{mat}[r]
              0  &1  \\
              -1 &0
            \end{mat}
          \end{equation*}
          对于反对称的那个，对角线上的元素必须为零。
        \partsitem 对称方阵等于它的转置。
          反对称方阵等于它的转置的负值。
        \partsitem 证明这两个集合是子空间是容易的。
         设 \( A\in\matspace_{\nbyn{n}} \)。
         为把 $A$ 表示成一个对称矩阵与一个反对称矩阵的和，
         我们注意到
         \begin{equation*}
            A=(1/2)(A+\trans{A}) + (1/2)(A-\trans{A})
         \end{equation*}
         并注意第一个加项是对称的，而第二个是
         反对称的。
         因此 \( \matspace_{\nbyn{n}} \) 是这两个子空间的和。
         为证明这个和是直和，
         设矩阵 \( A \) 既是对称的
         \( A=\trans{A} \) 又是反对称的 \( A=-\trans{A} \)。
         则 \( A=-A \)，于是 \( A \) 的所有元素都是零。
      \end{exparts}
     
\end{ans}
\begin{ans}{二.III.4.38}
      设
      \( \vec{v}\in (W_1\intersection W_2)+(W_1\intersection W_3) \)。
      则 \( \vec{v}=\vec{w}_2+\vec{w}_3 \)，其中
      \( \vec{w}_2\in W_1\intersection W_2 \) 且
      \( \vec{w}_3\in W_1\intersection W_3 \)。
      注意 \( \vec{w}_2,\vec{w}_3\in W_1 \)，又子空间对加法
      封闭，故 \( \vec{w}_2+\vec{w}_3\in W_1 \)。
      于是 \( \vec{v}=\vec{w}_2+\vec{w}_3\in W_1\intersection(W_2+W_3) \)。

      这个例子证明包含可能是严格的：
      在 \( \Re^2 \) 中取 \( W_1 \) 为 \( x \) 轴，取 \( W_2 \)
      为 \( y \) 轴，取 \( W_3 \) 为直线 \( y=x \)。
      则 \( W_1\intersection W_2 \) 与 \( W_1\intersection W_3 \)
      都是平凡的，因此它们的和是平凡的。
      但 \( W_2+W_3 \) 是整个 \( \Re^2 \)，所以
      \( W_1\intersection(W_2+W_3) \) 是 \( x \) 轴。
     
\end{ans}
\begin{ans}{二.III.4.39}
      当 \( W_1,W_2 \) 中至少有一个是平凡的时，这会发生。
      但这是它可能发生的唯一方式。

      为证明这一点，设两者都是非平凡的，
      从各自中取非零向量 \( \vec{w}_1, \vec{w}_2 \)，
      并考察 \( \vec{w}_1+\vec{w}_2 \)。
      这个和不属于 \( W_1 \)，因为
      \( \vec{w}_1+\vec{w}_2=\vec{v}\in W_1 \) 将蕴含
      \( \vec{w}_2=\vec{v}-\vec{w}_1 \) 属于 \( W_1 \)，
      这与子空间无关性的假设矛盾。
      类似地，\( \vec{w}_1+\vec{w}_2 \) 不属于 \( W_2 \)。
      于是 \( V \) 中有一个元素不属于 \( W_1\union W_2 \)。
     
\end{ans}
\begin{ans}{二.III.4.40}
      是。
      从左到右的蕴含就是
      \nearbycorollary{cor:DirSumDimsAdd}。
      对于另一个方向，
      设 \( \dim(V)=\dim(W_1)+\dots+\dim(W_k) \)。
      取 \( B_1,\dots, B_k \) 为 \( W_1,\dots, W_k \) 的基。
      由于 $V$ 是这些子空间的和，我们可以把
      任一 \( \vec{v}\in V \) 写成
      \( \vec{v}=\vec{w}_1+\cdots+\vec{w}_k \)，而把每个
      \( \vec{w}_i \) 表示成相应
      基 \( B_i \) 中向量的组合，即表明连接
      \( \cat{B_1}{\cat{\cdots}{B_k}}  \) 张成 \( V \)。
      现在，这个连接有
      \( \dim(W_1)+\dots+\dim(W_k) \) 个成员，因此它是一个大小为
      \( \dim(V) \) 的张成集。
      所以这个连接是 \( V \) 的一组基。
      于是 \( V \) 是直和。
     
\end{ans}
\begin{ans}{二.III.4.41}
      不能。
      \( \Re^2 \) 的标准基不能拆分成互补子空间，即直线 \( x=y \) 与直线
      \( x=-y \) 的基。
    
\end{ans}
\begin{ans}{二.III.4.42}
      \begin{exparts}
        \partsitem  是的，对一切子空间 \( W_1,W_2 \)
          都有 \( W_1+W_2=W_2+W_1 \)，
          因为每一边都是 \( W_1\union W_2=W_2\union W_1 \) 的张成。
        \partsitem 这一条与前一条类似\Dash 该等式的每一边
          都是
          \( (W_1\union W_2)\union W_3=W_1\union (W_2\union W_3) \)
          的张成。
        \partsitem 由于这是集合之间的等式，我们可以用相互包含
          来证明它成立。
          显然 \( W\subseteq W+W \)。
          对于 \( W+W\subseteq W \)，只需记得每个子集都对加法封闭，
          所以形如 \( \vec{w}_1+\vec{w}_2 \) 的任何和都在
          \( W \) 中。
        \partsitem 在每个向量空间中，关于子空间加法的
          单位元是平凡子空间。
        \partsitem   左消去与右消去都不必成立。
          举个例子，在 \( \Re^3 \) 中取 \( W_1 \) 为
          \( xy \) 平面，取 \( W_2 \) 为 \( x \) 轴，
          取 \( W_3 \) 为 \( y \) 轴。
      \end{exparts}
     
\end{ans}
\protect \topic {域}
\begin{ans}{1}
      逐条检查这五个条件可知，它们都是初等数学中我们熟悉的内容。
    
\end{ans}
\begin{ans}{2}
      与前面的问题一样，逐条检查这五个条件
      可知，
      对这两个结构而言，
      这些性质都是我们熟悉的。
    
\end{ans}
\begin{ans}{3}
       整数不满足条件~(5)。
       例如，$2$ 没有乘法逆元\Dash 虽然
       $2$ 是整数，$1/2$ 却不是。
     
\end{ans}
\begin{ans}{4}
       我们可以通过列举所有可能的情况来做这些验证。
       例如，为验证条件~(2)的前半部分，我们必须检查
       该结构对加法封闭，并且加法
       满足交换律 $a+b=b+a$，
       由于这样的数对
       只有四对，我们可以对所有可能的数对 $a$ 与~$b$ 检验这两点。
       类似地，对于结合律，三元组 $a$、$b$、
       $c$ 只有八种，因此检验并不太长。
       （还有其他做这些验证的方法；特别地，你可能会
       认出这些运算就是模~$2$ 的算术。
       但穷举检验也并不繁重）
     
\end{ans}
\begin{ans}{5}
       下面这些运算就可以。
       \begin{center}
          \begin{tabular}{c|ccc}
             \( + \) &\( 0 \) &\( 1 \) &$2$ \\
             \hline
             \( 0 \) &\( 0 \) &\( 1 \) &$2$ \\
             \( 1 \) &\( 1 \) &\( 2 \) &$0$ \\
             $2$     &$2$     &$0$     &$1$
          \end{tabular}
          \qquad
          \begin{tabular}{c|ccc}
            \( \cdot \) &\( 0 \) &\( 1 \) &$2$ \\
            \hline
            \( 0 \)  &\( 0 \) &\( 0 \) &$0$ \\
            \( 1 \)  &\( 0 \) &\( 1 \) &$2$ \\
            $2$      &$0$     &$2$     &$1$
          \end{tabular}
       \end{center}
       与前面的题目一样，我们可以
       通过列举所有情况来
       验证它们满足这些条件。
     
\end{ans}
\protect \topic {晶体}
\begin{ans}{1}
     每个基本单元是 $3.34\times 10^{-10}$~cm，所以这样的单元约有
     $0.1/(3.34\times 10^{-10})$ 个。
     这给出 $2.99\times 10^{8}$，因此差不多有
     $300,000,000$（三亿）个区域。
   
\end{ans}
\begin{ans}{2}
      \begin{exparts}
        \partsitem 我们求解
          \begin{equation*}
            c_1\colvec{1.42 \\ 0}+c_2\colvec{0.71 \\ 1.23}
              =\colvec{5.67 \\ 3.14}
            \quad\implies\quad
            \begin{linsys}{2}
              1.42c_1  &+  &0.71c_2  &=  &5.67  \\
                       &   &1.23c_2  &=  &3.14
            \end{linsys}
          \tag*{}\end{equation*}
          解得 $c_1\approx 2.74$ 与 $c_2\approx 2.55$。
        \partsitem 下面是这个点在晶格中的定位。
          在左图中，原胞上叠加了
          两个基向量 $\vec{\beta}_1$ 与 $\vec{\beta}_2$，
          以及一个标出偏移量
          $2.55\vec{\beta}_1+2.72\vec{\beta}_2$ 的方框。
          右图显示的是它在晶体晶格内部出现的位置，这里取
          左下角那个六边形的左下顶点作为原点。
          \begin{center}  %graphite
            \includegraphics{vs/mp/ch2.17}
            \qquad
            \includegraphics{vs/mp/ch2.18}
       \end{center}
       所以这个点向右移过两列六边形，
       向上移过一行六边形。
      \partsitem 这第二个基
          \begin{equation*}
            \sequence{\colvec{1.42 \\ 0},\colvec{0 \\ 1.42}}
          \tag*{}\end{equation*}
          使计算更容易
          \begin{equation*}
            c_1\colvec{1.42 \\ 0}+c_2\colvec{0 \\ 1.42}
              =\colvec{5.67 \\ 3.14}
            \quad\implies\quad
            \begin{linsys}{2}
              1.42c_1  &   &         &=  &5.67  \\
                       &   &1.42c_2  &=  &3.14
            \end{linsys}
          \tag*{}\end{equation*}
          （我们得到 $c_2\approx 2.21$ 与 $c_1\approx 3.99$），但它似乎
          与我们正在研究的物理结构没有多大关系。
      \end{exparts}
    
\end{ans}
\begin{ans}{3}
      用基表示时，角上原子的位置是
      $(0,0,0)$、$(1,0,0)$、\ldots{}、$(1,1,1)$。
      面上原子的位置是 $(0.5,0.5,1)$、$(1,0.5,0.5)$、
      $(0.5,1,0.5)$、$(0,0.5,0.5)$、$(0.5,0,0.5)$ 以及 $(0.5,0.5,0)$。
      自顶面向下四分之一处的原子的位置
      是 $(0.75,0.75,0.75)$ 与 $(0.25,0.25,0.25)$。
      自底面向上四分之一处的原子
      位于 $(0.75,0.25,0.25)$ 与 $(0.25,0.75,0.25)$。
      换算成埃很容易。
    
\end{ans}
\begin{ans}{4}
      \begin{exparts}
        \partsitem $195.08/6.02\times 10^{23}=3.239\times 10^{-22}$
        \partsitem 每个面中央的铂原子
          由两个立方体分摊，所以到这里共有 $6/2=3$ 个原子。
          每个角上的原子由八个立方体分摊，所以
          又增加 $8/8=1$~个原子，因此总数是 $4$。
        \partsitem $4\cdot 3.239\times 10^{-22}=1.296\times 10^{-21}$
        \partsitem $1.296\times 10^{-21}/21.45=6.042\times 10^{-23}$ 立方
           厘米
        \partsitem $3.924\times 10^{-8}$ 厘米。
        \partsitem
            $\sequence{\colvec{3.924\times 10^{-8} \\ 0 \\ 0},
                       \colvec{0 \\ 3.924\times 10^{-8} \\ 0},
                       \colvec{0 \\ 0 \\ 3.924\times 10^{-8}}}$
      \end{exparts}
    
\end{ans}
\protect \topic {投票悖论}
\begin{ans}{1}
      这个例子给出单个选民的一个非理性偏好顺序。
      \smallskip
      \begin{center}
        \begin{tabular}{l|c|c|c|}
             \multicolumn{1}{c}{\ }
                  &\multicolumn{1}{c}{\textit{品格}}
                  &\multicolumn{1}{c}{\textit{经验}}
                  &\multicolumn{1}{c}{\textit{政策}}   \\
          \cline{2-4}
         \textit{民主党}   &最偏好    &中间  &最不偏好   \\
         \textit{共和党} &中间  &最不偏好   &最偏好    \\
         \textit{第三党}      &最不偏好   &最偏好    &中间  \\
          \cline{2-4}
        \end{tabular}
      \end{center}
      \smallskip
      在品格和经验上，民主党优于共和党。
      在品格和政策上，共和党胜过第三党。
      而在经验和政策上，第三党击败民主党。
    
\end{ans}
\begin{ans}{2}
      首先，把本专题中讲过的 $D>R>T$ 分解
      \begin{equation*}
        \votepreflist{-1}{1}{1}
          =\frac{1}{3}\cdot\votepreflist{1}{1}{1}
           +\frac{2}{3}\cdot\votepreflist{-1}{1}{0}
           +\frac{2}{3}\cdot\votepreflist{-1}{0}{1}
      \end{equation*}
      与相反的 $T>R>D$ 选民的分解
      \begin{equation*}
        \votepreflist{1}{-1}{-1}
          =d_1\cdot\votepreflist{1}{1}{1}
           +d_2\cdot\votepreflist{-1}{1}{0}
           +d_3\cdot\votepreflist{-1}{0}{1}
      \tag*{}\end{equation*}
      相比较。
      显然，第二个是第一个的相反向量，因此
      $d_1=-1/3$、$d_2=-2/3$、$d_3=-2/3$。
      这一原则对任何一对相反的选民都成立，所以我们只需
      对来自第二行的一位选民与来自第三行的
      一位选民作计算。
      对于第二行中具有正旋向的选民，
      \begin{equation*}
        \begin{linsys}{3}
          c_1  &-  &c_2  &-  &c_3  &=  &1 \\
          c_1  &+  &c_2  &   &     &=  &1  \\
          c_1  &   &     &+  &c_3  &=  &-1
        \end{linsys}
        \grstep[-\rho_1+\rho_3]{-\rho_1+\rho_2}
        \repeatedgrstep{(-1/2)\rho_2+\rho_3}
        \begin{linsys}{3}
          c_1  &-  &c_2  &-  &c_3      &=  &1 \\
               &   &2c_2 &+  &c_3      &=  &0  \\
               &   &     &   &(3/2)c_3 &=  &-2
        \end{linsys}
      \tag*{}\end{equation*}
      得 $c_3=-4/3$、$c_2=2/3$、$c_1=1/3$。
      对于第三行中具有正旋向的选民，
      \begin{equation*}
        \begin{linsys}{3}
          c_1  &-  &c_2  &-  &c_3  &=  &1 \\
          c_1  &+  &c_2  &   &     &=  &-1  \\
          c_1  &   &     &+  &c_3  &=  &1
        \end{linsys}
        \grstep[-\rho_1+\rho_3]{-\rho_1+\rho_2}
        \repeatedgrstep{(-1/2)\rho_2+\rho_3}
        \begin{linsys}{3}
          c_1  &-  &c_2  &-  &c_3      &=  &1 \\
               &   &2c_2 &+  &c_3      &=  &-2  \\
               &   &     &   &(3/2)c_3 &=  &1
        \end{linsys}
        \tag*{}\end{equation*}
        得 $c_3=2/3$、$c_2=-4/3$、$c_1=1/3$。
      
\end{ans}
\begin{ans}{3}
      模拟选举按第一个表所示的方式对应于
      第~\pageref{table:Voting}~页的表格，
      抵消之后的结果是第二个表。
     \smallskip
      \begin{center}
         % \renewcommand{\arraystretch}{1.2}
         \begin{tabular}{|c|c|} \hline
           \textit{正旋向}
           &\textit{负旋向}  \\ \hline
           \begin{tabular}{l}
             $D>R>T$        \\
             $5$ 名选民
           \end{tabular}
          &\begin{tabular}{l}
             $T>R>D$          \\
             $2$ 名选民
           \end{tabular}                  \\ \hline
          \begin{tabular}{l}
             $R>T>D$           \\
             $8$ 名选民
          \end{tabular}
         &\begin{tabular}{l}
             $D>T>R$             \\
             $4$ 名选民
           \end{tabular}                  \\ \hline
          \begin{tabular}{l}
            $T>D>R$             \\
            $8$ 名选民
          \end{tabular}
         &\begin{tabular}{l}
            $R>D>T$            \\
            $2$ 名选民
          \end{tabular}                  \\ \hline
        \end{tabular}
        \qquad
         \begin{tabular}{|c|c|} \hline
           \textit{正旋向}
           &\textit{负旋向}  \\ \hline
           \begin{tabular}{l}
             $D>R>T$        \\
             $3$ 名选民
           \end{tabular}
          &\begin{tabular}{l}
             $T>R>D$          \\
             \textit{-已抵消-}
           \end{tabular}                  \\ \hline
          \begin{tabular}{l}
             $R>T>D$           \\
             $4$ 名选民
          \end{tabular}
         &\begin{tabular}{l}
             $D>T>R$             \\
             \textit{-已抵消-}
           \end{tabular}                  \\ \hline
          \begin{tabular}{l}
            $T>D>R$             \\
            $6$ 名选民
          \end{tabular}
         &\begin{tabular}{l}
            $R>D>T$            \\
            \textit{-已抵消-}
          \end{tabular}                  \\ \hline
        \end{tabular}
       \end{center}
       \smallskip
       三组全部来自表格的同一侧（左侧），
       正如本专题的结果所指出的必然情形。
       现在我们可以统计，利用抵消后的数字
      \begin{equation*}
        % \psset{xunit=4pt,yunit=4pt,runit=4pt}
        3\cdot \votinggraphic{31}
%          \pspicture[.4](-10.5,-4.2)(10,4) %\psgrid(-10,-10)(10,10)
%            \SpecialCoor
%            \rput(3;90){\small $D$}
%            \psarc{->}(0,0){3}{10}{70} %-rd->
%               \rput[lb](3.3;20){\scriptsize $1$ voter}
%            \rput[B](3;210){\small $T$}
%            \psarc{->}(0,0){3}{220}{320} %-tr->
%               \rput[t](3.1;270){\scriptsize $1$ voter}
%            \rput[B](3;330){\small $R$}
%            \psarc{->}(0,0){3}{110}{175} %-dt->
%               \rput[rb](3.3;160){\scriptsize $-1$ voter}
%          \endpspicture
        +4\cdot\votinggraphic{32}
%          \pspicture[.4](-10,-4.2)(10,4) %\psgrid(-10,-10)(10,10)
%            \SpecialCoor
%            \rput(3;90){\small $D$}
%            \psarc{->}(0,0){3}{10}{70} %-rd->
%               \rput[lb](3.3;20){\scriptsize $-1$ voter}
%            \rput[B](3;210){\small $T$}
%            \psarc{->}(0,0){3}{220}{320} %-tr->
%               \rput[t](3.1;270){\scriptsize $1$ voter}
%            \rput[B](3;330){\small $R$}
%            \psarc{->}(0,0){3}{110}{175} %-dt->
%               \rput[rb](3.3;160){\scriptsize $1$ voter}
%          \endpspicture
        +6\cdot\votinggraphic{33}
%          \pspicture[.4](-9,-4.2)(10,4) %\psgrid(-10,-10)(10,10)
%            \SpecialCoor
%            \rput(3;90){\small $D$}
%            \psarc{->}(0,0){3}{10}{70} %-rd->
%               \rput[lb](3.3;20){\scriptsize $1$ voter}
%            \rput[B](3;210){\small $T$}
%            \psarc{->}(0,0){3}{220}{320} %-tr->
%               \rput[t](3.1;270){\scriptsize $-1$ voter}
%            \rput[B](3;330){\small $R$}
%            \psarc{->}(0,0){3}{110}{175} %-dt->
%               \rput[rb](3.3;160){\scriptsize $1$ voter}
%          \endpspicture
        = \votinggraphic{34}
%          \pspicture[.4](-9,-4.2)(10,4) %\psgrid(-10,-10)(10,10)
%            \SpecialCoor
%            \rput(3;90){\small $D$}
%            \psarc{->}(0,0){3}{10}{70} %-rd->
%               \rput[lb](3.3;20){\scriptsize $5$ voters}
%            \rput[B](3;210){\small $T$}
%            \psarc{->}(0,0){3}{220}{320} %-tr->
%               \rput[t](3.1;270){\scriptsize $1$ voter}
%            \rput[B](3;330){\small $R$}
%            \psarc{->}(0,0){3}{110}{175} %-dt->
%               \rput[rb](3.3;160){\scriptsize $7$ voters}
%          \endpspicture
      \tag*{}\end{equation*}
       得到同样的结果。
    
\end{ans}
\begin{ans}{4}
      \begin{exparts}
        % \partsitem We can rewrite the two as $-c\leq a-b$ and $-c\leq b-a$.
        %   Either $a-b$ or $b-a$ is nonpositive and so $-c\leq-\absval{a-b}$,
        %   as required.
        % \partsitem This is immediate from the supposition that $0\leq a+b-c$.
        \partsitem 一个平凡的例子从零选民的选举开始，
          它有一个平凡的多数循环，然后添加任意一名选民。
          一个更有趣的例子是取政治学课堂的模拟
          选举并添加两名 $T>D>R$ 选民
          （为了满足题目中
          「再添加一名选民」的要求，
          我们可以一次添加一名）。
          新选民的旋向为正，这正是原模拟选举中
          抵消后剩下的选票的旋向。
          下面是所得的选民表，旁边是
          抵消的结果。
      \smallskip
      \begin{center}
         % \renewcommand{\arraystretch}{1.3}
         \begin{tabular}{|c|c|} \hline
           \textit{正旋向}
           &\textit{负旋向}  \\ \hline
           \begin{tabular}{l}
             $D>R>T$        \\
             $5$ 名选民
           \end{tabular}
          &\begin{tabular}{l}
             $T>R>D$          \\
             $2$ 名选民
           \end{tabular}                  \\ \hline
          \begin{tabular}{l}
             $R>T>D$           \\
             $8$ 名选民
          \end{tabular}
         &\begin{tabular}{l}
             $D>T>R$             \\
             $4$ 名选民
           \end{tabular}                  \\ \hline
          \begin{tabular}{l}
            $T>D>R$             \\
            $10$ 名选民
          \end{tabular}
         &\begin{tabular}{l}
            $R>D>T$            \\
            $2$ 名选民
          \end{tabular}                  \\ \hline
        \end{tabular}
        \qquad
         \begin{tabular}{|c|c|} \hline
           \textit{正旋向}
           &\textit{负旋向}  \\ \hline
           \begin{tabular}{l}
             $D>R>T$        \\
             $3$ 名选民
           \end{tabular}
          &\begin{tabular}{l}
             $T>R>D$          \\
             \textit{-已抵消}
           \end{tabular}                  \\ \hline
          \begin{tabular}{l}
             $R>T>D$           \\
             $4$ 名选民
          \end{tabular}
         &\begin{tabular}{l}
             $D>T>R$             \\
             \textit{-已抵消-}
           \end{tabular}                  \\ \hline
          \begin{tabular}{l}
            $T>D>R$             \\
            $8$ 名选民
          \end{tabular}
         &\begin{tabular}{l}
            $R>D>T$            \\
            \textit{-已抵消-}
          \end{tabular}                  \\ \hline
        \end{tabular}
       \end{center}
       \smallskip
       这是利用抵消后的数字所作的选举。
       \begin{equation*}
         % \psset{xunit=4pt,yunit=4pt,runit=4pt}
         3\cdot\votinggraphic{35}
%          \pspicture[.4](-10.5,-4.2)(10,4) %\psgrid(-10,-10)(10,10)
%            \SpecialCoor
%            \rput(3;90){\small $D$}
%            \psarc{->}(0,0){3}{10}{70} %-rd->
%               \rput[lb](3.3;20){\scriptsize $1$ voter}
%            \rput[B](3;210){\small $T$}
%            \psarc{->}(0,0){3}{220}{320} %-tr->
%               \rput[t](3.1;270){\scriptsize $1$ voter}
%            \rput[B](3;330){\small $R$}
%            \psarc{->}(0,0){3}{110}{175} %-dt->
%               \rput[rb](3.3;160){\scriptsize $-1$ voter}
%          \endpspicture
         +4\cdot\votinggraphic{36}
%          \pspicture[.4](-10,-4.2)(10,4) %\psgrid(-10,-10)(10,10)
%            \SpecialCoor
%            \rput(3;90){\small $D$}
%            \psarc{->}(0,0){3}{10}{70} %-rd->
%               \rput[lb](3.3;20){\scriptsize $-1$ voter}
%            \rput[B](3;210){\small $T$}
%            \psarc{->}(0,0){3}{220}{320} %-tr->
%               \rput[t](3.1;270){\scriptsize $1$ voter}
%            \rput[B](3;330){\small $R$}
%            \psarc{->}(0,0){3}{110}{175} %-dt->
%               \rput[rb](3.3;160){\scriptsize $1$ voter}
%          \endpspicture
         +8\cdot\votinggraphic{37}
%          \pspicture[.4](-9,-4.2)(10,4) %\psgrid(-10,-10)(10,10)
%            \SpecialCoor
%            \rput(3;90){\small $D$}
%            \psarc{->}(0,0){3}{10}{70} %-rd->
%               \rput[lb](3.3;20){\scriptsize $1$ voter}
%            \rput[B](3;210){\small $T$}
%            \psarc{->}(0,0){3}{220}{320} %-tr->
%               \rput[t](3.1;270){\scriptsize $-1$ voter}
%            \rput[B](3;330){\small $R$}
%            \psarc{->}(0,0){3}{110}{175} %-dt->
%               \rput[rb](3.3;160){\scriptsize $1$ voter}
%          \endpspicture
         =\votinggraphic{38}
%          \pspicture[.4](-9,-4.2)(10,4) %\psgrid(-10,-10)(10,10)
%            \SpecialCoor
%            \rput(3;90){\small $D$}
%            \psarc{->}(0,0){3}{10}{70} %-rd->
%               \rput[lb](3.3;20){\scriptsize $7$ voters}
%            \rput[B](3;210){\small $T$}
%            \psarc{->}(0,0){3}{220}{320} %-tr->
%               \rput[t](3.1;270){\scriptsize $-1$ voter}
%            \rput[B](3;330){\small $R$}
%            \psarc{->}(0,0){3}{110}{175} %-dt->
%               \rput[rb](3.3;160){\scriptsize $9$ voters}
%          \endpspicture
      \tag*{}\end{equation*}
       多数循环确实消失了。
      \partsitem 把前面的练习反过来。
         也就是说，在前一个练习的结果上添加一名选民，
         把使循环消失的那位选民抵消掉。
      \partsitem 一个这样的条件是：抵消之后，
         三个数全部非负或全部非正，
         并且：~$\absval{c}<\absval{a+b}$、$\absval{b}<\absval{a+c}$
         与 $\absval{a}<\absval{b+c}$。
         这由下图得出。
         \begin{equation*}
           \votinggraphic{27}  %\votecycle{a}{a}{-a}{4}
           +
           \votinggraphic{28}  %\votecycle{-b}{b}{b}{4}
           +
           \votinggraphic{29}  %\votecycle{c}{-c}{c}{4}
           =\hspace{.30in}
           \votinggraphic{30}  %\votecycle{a-b+c}{a+b-c}{-a+b+c}{4}
           \hspace{.30in}\hbox{}
         \end{equation*}
%          \begin{center}
%            \votecycle{a}{a}{-a}{4}
%            +
%            \votecycle{-b}{b}{b}{4}
%            +
%            \votecycle{c}{-c}{c}{4}
%            =\hspace{.30in}
%            \votecycle{a-b+c}{a+b-c}{-a+b+c}{4}\hspace{.30in}\hbox{}
%          \end{center}
      \end{exparts}
    
\end{ans}
\begin{ans}{5}
      \begin{exparts}
        \partsitem 两名选民的选举可能以两种方式出现多数循环。
          第一，两名选民可能恰好相反，
          抵消后得到平凡的选举（多数循环全为零）。
          第二，两名选民可能具有相同旋向但来自
          不同的行，如下所示。
          \begin{equation*}
            % \psset{xunit=4pt,yunit=4pt,runit=4pt}
            1\cdot \votinggraphic{39}
%             \pspicture[.4](-10.5,-4.2)(10,4) %\psgrid(-10,-10)(10,10)
%               \SpecialCoor
%               \rput(3;90){\small $D$}
%               \psarc{->}(0,0){3}{10}{70} %-rd->
%                  \rput[lb](3.3;20){\scriptsize $1$ voter}
%               \rput[B](3;210){\small $T$}
%               \psarc{->}(0,0){3}{220}{320} %-tr->
%                  \rput[t](3.1;270){\scriptsize $1$ voter}
%               \rput[B](3;330){\small $R$}
%               \psarc{->}(0,0){3}{110}{175} %-dt->
%                  \rput[rb](3.3;160){\scriptsize $-1$ voter}
%             \endpspicture
            +1\cdot\votinggraphic{40}
%             \pspicture[.4](-10,-4.2)(10,4) %\psgrid(-10,-10)(10,10)
%               \SpecialCoor
%               \rput(3;90){\small $D$}
%               \psarc{->}(0,0){3}{10}{70} %-rd->
%                  \rput[lb](3.3;20){\scriptsize $-1$ voter}
%               \rput[B](3;210){\small $T$}
%               \psarc{->}(0,0){3}{220}{320} %-tr->
%                  \rput[t](3.1;270){\scriptsize $1$ voter}
%               \rput[B](3;330){\small $R$}
%               \psarc{->}(0,0){3}{110}{175} %-dt->
%                  \rput[rb](3.3;160){\scriptsize $1$ voter}
%             \endpspicture
            +0\cdot\votinggraphic{41}
%             \pspicture[.4](-9,-4.2)(10,4) %\psgrid(-10,-10)(10,10)
%               \SpecialCoor
%               \rput(3;90){\small $D$}
%               \psarc{->}(0,0){3}{10}{70} %-rd->
%                  \rput[lb](3.3;20){\scriptsize $1$ voter}
%               \rput[B](3;210){\small $T$}
%               \psarc{->}(0,0){3}{220}{320} %-tr->
%                  \rput[t](3.1;270){\scriptsize $-1$ voter}
%               \rput[B](3;330){\small $R$}
%               \psarc{->}(0,0){3}{110}{175} %-dt->
%                  \rput[rb](3.3;160){\scriptsize $1$ voter}
%             \endpspicture
            =\votinggraphic{42}
%             \pspicture[.4](-9,-4.2)(10,4) %\psgrid(-10,-10)(10,10)
%               \SpecialCoor
%               \rput(3;90){\small $D$}
%               \psarc{->}(0,0){3}{10}{70} %-rd->
%                  \rput[lb](3.3;20){\scriptsize $0$ voters}
%               \rput[B](3;210){\small $T$}
%               \psarc{->}(0,0){3}{220}{320} %-tr->
%                  \rput[t](3.1;270){\scriptsize $2$ voters}
%               \rput[B](3;330){\small $R$}
%               \psarc{->}(0,0){3}{110}{175} %-dt->
%                  \rput[rb](3.3;160){\scriptsize $0$ voters}
%             \endpspicture
          \tag*{}\end{equation*}
        \partsitem 有两种情形。
          偶数个选民可以平分成相反的两半，
          例如一半选民是 $D>R>T$，一半是 $T>R>D$。
          那么抵消给出平凡的选举。
          如果选民数大于一且为奇数（形如
          $2k+1$，其中 $k>0$），那么利用证明中的循环图，
          \begin{equation*}
            \votinggraphic{27}  %\votecycle{a}{a}{-a}{4}
            +
            \votinggraphic{28}  %\votecycle{-b}{b}{b}{4}
            +
            \votinggraphic{29}  %\votecycle{c}{-c}{c}{4}
            =\hspace{.30in}
            \votinggraphic{30}  %\votecycle{a-b+c}{a+b-c}{-a+b+c}{4}
            \hspace{.30in}\hbox{}
          \end{equation*}
%          \begin{center}
%            \votecycle{a}{a}{-a}{4}
%            +
%            \votecycle{-b}{b}{b}{4}
%            +
%            \votecycle{c}{-c}{c}{4}
%            =\hspace{.30in}
%            \votecycle{a-b+c}{a+b-c}{-a+b+c}{4}\hspace{.30in}\hbox{}
%          \end{center}
          我们可以取 $a=k$、$b=k$、$c=1$。
          因为 $k>0$，这就是一个多数循环。
       \end{exparts}
    
\end{ans}
\begin{ans}{6}
      它非空，因为它包含零向量。
      为证它对其两个成员的线性组合封闭，
      设 $\vec{v}_1$ 与 $\vec{v}_2$ 属于 $U^\perp$，
      考虑 $c_1\vec{v}_1+c_2\vec{v}_2$。
      对任意 $\vec{u}\in U$，
      \begin{equation*}
        (c_1\vec{v}_1+c_2\vec{v}_2)\dotprod\vec{u}
        =c_1(\vec{v}_1\dotprod\vec{u})+c_2(\vec{v}_2\dotprod\vec{u})
        =c_1\cdot 0+c_2\cdot 0=0
      \tag*{}\end{equation*}
      所以 $c_1\vec{v}_1+c_2\vec{v}_2\in U^\perp$。

      至于 $U$ 是一个子集但不是子空间时是否仍然成立，
      答案是肯定的。
    
\end{ans}
\protect \topic {量纲分析}
\begin{ans}{1}
      \begin{exparts}
        \partsitem 假设这个
          \begin{equation*}
            (L^1M^0T^0)^{p_1}(L^1M^0T^0)^{p_2}(L^1M^0T^{-1})^{p_3}
             (L^0M^0T^0)^{p_4}(L^1M^0T^{-2})^{p_5}(L^0M^0T^1)^{p_6}
          \end{equation*}
          等于 $L^0M^0T^0$
          便产生这个线性方程组
          \begin{equation*}
            \begin{linsys}{5}
              p_1  &+  &p_2  &+  &p_3  &+  &p_5  &   &    &=  &0  \\
              %      &   &     &   &     &   &     &   &0   &=  &0  \\
                   &   &     &   &-p_3 &-  &2p_5 &+  &p_6 &=  &0
            \end{linsys}
          \end{equation*}
          （对 $p_4$ 没有限制）。
          自然的参数化利用自由变量给出
          $p_3=-2p_5+p_6$ 与 $p_1=-p_2+p_5-p_6$。
          所得解集的描述
          \begin{equation*}
            \set{\colvec{p_1 \\ p_2 \\ p_3 \\ p_4 \\ p_5 \\ p_6}
                =p_2\colvec[r]{-1 \\ 1 \\ 0 \\ 0  \\ 0 \\ 0}
                +p_4\colvec[r]{0 \\ 0 \\ 0 \\ 1 \\ 0 \\ 0}
                +p_5\colvec[r]{1 \\ 0 \\ -2 \\ 0 \\ 1 \\ 0}
                +p_6\colvec[r]{-1 \\ 0 \\ 1 \\ 0 \\ 0 \\ 1}
                \suchthat p_2, p_4, p_5, p_6\in\Re  }
          \end{equation*}
          给出 $\set{y/x,\theta,xt/{v_0}^2,v_0t/x}$
          作为一个完备的无量纲乘积集合
          （回忆一下，这里的「完备」并不意味着
          不存在其他无量纲乘积；
          它只是说这个集合是一组基）。
          然而，这并不是
          题目所要求的那个无量纲乘积集合。

          有两条路可走。
          第一条是调整参数的选择，希望能
          碰上正确的集合。
          为此，我们可以把前一段倒过来做。
          把给定的无量纲乘积 $gt/v_0$、
          $gx/v_0^2$、$gy/v_0^2$ 与 $\theta$
          化成向量，得到这个描述（注意参数将要放置
          的位置上的 \textit{?}）。
          \begin{equation*}
            \set{\colvec{p_1 \\ p_2 \\ p_3 \\ p_4 \\ p_5 \\ p_6}
                =\underline{\ \text{\textit{?}}\ }
                      \colvec[r]{0  \\ 0 \\ -1 \\ 0  \\ 1 \\ 1}
                +\underline{\ \text{\textit{?}}\ }
                   \colvec[r]{1 \\ 0 \\ -2 \\ 0 \\ 1 \\ 0}
                +\underline{\ \text{\textit{?}}\ }
                   \colvec[r]{0 \\ 1 \\ -2 \\ 0 \\ 1 \\ 0}
                +p_4\colvec[r]{0  \\ 0 \\ 0 \\ 1 \\ 0 \\ 0}
                \suchthat p_2, p_4, p_5, p_6\in\Re  }
          \end{equation*}
          $p_4$ 已经就位。
          检查各行可知，我们还能把 $p_6$、$p_1$
          与 $p_2$ 也放到相应位置上。

          按照提示的第二条路，
          是注意到给定的集合在四维向量空间中
          大小为四，
          因此我们只需证明它线性无关。
          通过观察向量第六、第一、第二
          与第四个分量即可轻松完成。
        \item 第一个方程可以改写为
          \begin{equation*}
            \frac{gx}{{v_0}^2}=\frac{gt}{v_0}\cos\theta
          \end{equation*}
          于是 Buckingham 函数为
          $f_1(\Pi_1,\Pi_2,\Pi_3,\Pi_4)=\Pi_2-\Pi_1\cos(\Pi_4)$。
          第二个方程可以改写为
          \begin{equation*}
            \frac{gy}{{v_0}^2}=\frac{gt}{v_0}\sin\theta
              -\frac{1}{2}\left(\frac{gt}{v_0}\right)^2
          \end{equation*}
          这里的 Buckingham 函数为
          $f_2(\Pi_1,\Pi_2,\Pi_3,\Pi_4)
             =\Pi_3-\Pi_1\sin(\Pi_4)+(1/2){\Pi_1}^2$。
      \end{exparts}
    
\end{ans}
\begin{ans}{2}
      考虑
      \begin{equation*}
        (L^0M^0T^{-1})^{p_1}(L^1M^{-1}T^2)^{p_2}
          (L^{-3}M^0T^0)^{p_3}(L^0M^1T^0)^{p_4}=(L^0M^0T^0)
      \end{equation*}
      它给出这些幂之间的关系。
      \begin{equation*}
        \begin{linsys}{4}
               &    &p_2   &-  &3p_3  &   &    &=  &0  \\
               &    &-p_2  &   &      &+  &p_4 &=  &0  \\
         -p_1  &+   &2p_2  &   &      &   &    &=  &0
        \end{linsys}
        \grstep{\rho_1\leftrightarrow\rho_3}
        \repeatedgrstep{\rho_2+\rho_3}
        \begin{linsys}{4}
         -p_1  &+   &2p_2  &   &      &   &    &=  &0  \\
               &    &-p_2  &   &      &+  &p_4 &=  &0  \\
               &    &      &\  &-3p_3 &+  &p_4 &=  &0
        \end{linsys}
      \end{equation*}
      这就是解空间
      （因为我们希望把 $k$ 表示成
      其他量的函数，
      所以取 $p_2$ 作为参数）。
      \begin{equation*}
        \set{\colvec[r]{2 \\ 1 \\ 1/3 \\ 1}p_2
             \suchthat p_2\in\Re}
      \end{equation*}
      因此，$\Pi_1=\nu^2kN^{1/3}m$ 是这个无量纲组合，
      并且 $k$ 等于
      $\nu^{-2}N^{-1/3}m^{-1}$ 乘以一个常数
      （函数 $\hat{f}$ 是常数，因为它没有自变量）。
    
\end{ans}
\begin{ans}{3}
      \begin{exparts}
        \partsitem 令
          \begin{equation*}
            (L^2M^1T^{-2})^{p_1}(L^0M^0T^{-1})^{p_2}(L^{3}M^0T^0)^{p_3}
              =(L^0M^0T^0)
          \end{equation*}
          得到这个
          \begin{equation*}
            \begin{linsys}{3}
              2p_1  &   &    &+  &3p_3  &=  &0  \\
              p_1   &   &    &   &      &=  &0  \\
              -2p_1 &-  &p_2 &   &      &=  &0
            \end{linsys}
          \end{equation*}
          它蕴含 $p_1=p_2=p_3=0$。
          也就是说，在具有这些量纲公式的量中，唯一
          的无量纲乘积是平凡的那个。
        \partsitem 令
          \begin{equation*}
            (L^2M^1T^{-2})^{p_1}(L^0M^0T^{-1})^{p_2}(L^{3}M^0T^0)^{p_3}
              (L^{-3}M^1T^0)^{p_4}=(L^0M^0T^0)
          \end{equation*}
          得到这个。
          \begin{multline*}
            \begin{linsys}{4}
              2p_1  &   &    &+  &3p_3  &-  &3p_4 &=  &0  \\
              p_1   &   &    &   &      &+  &p_4  &=  &0  \\
              -2p_1 &-  &p_2 &   &      &   &     &=  &0
            \end{linsys}                                         \\
            \grstep[\rho_1+\rho_3]{(-1/2)\rho_1+\rho_2}
            \repeatedgrstep{\rho_2\leftrightarrow\rho_3}
            \begin{linsys}{4}
              2p_1  &   &      &+  &3p_3       &-   &3p_4      &=  &0  \\
                    &   &-p_2  &+  &3p_3       &-   &3p_4      &=  &0  \\
                    &  &      &   &(-3/2)p_3  &+   &(5/2)p_4  &=  &0
            \end{linsys}
          \end{multline*}
          取 $p_1$ 作为参数来表示转矩，
          得到解集的这个描述。
          \begin{equation*}
            \set{\colvec[r]{1 \\ -2 \\ -5/3 \\ -1}p_1
                 \suchthat p_1\in\Re}
          \end{equation*}
          记转矩为 $\tau$、转速为 $r$、空气体积
          为 $V$、空气密度为 $d$，我们有
          $\Pi_1=\tau r^{-2} V^{-5/3} d^{-1}$，于是
          转矩是 $r^2V^{5/3}d$ 乘以一个常数。
      \end{exparts}
    
\end{ans}
\begin{ans}{4}
          \begin{exparts}
            \partsitem 这些是量纲公式。
              \begin{center}
                \begin{tabular}{r|l}
                  \multicolumn{1}{r}{\textit{量}}
                    &\multicolumn{1}{l}{\begin{tabular}[b]{@{}l}
                                           \textit{量纲} \\[-.3ex]
                                           \textit{公式}
                                       \end{tabular}} \\ \hline
                        波速 $v$               &$L^1M^0T^{-1}$ \\
                        骨牌间距 $d$      &$L^1M^0T^0$   \\
                        骨牌高度 $h$          &$L^1M^0T^0$ \\
                        重力加速度 $g$    &$L^1M^0T^{-2}$
                \end{tabular}
              \end{center}
            \partsitem 关系式
              \begin{equation*}
                (L^1M^0T^{-1})^{p_1}(L^1M^0T^0)^{p_2}(L^1M^0T^0)^{p_3}
                  (L^1M^0T^{-2})^{p_4}=(L^0M^0T^0)
              \end{equation*}
              给出这个线性方程组。
              \begin{equation*}
                \begin{linsys}{4}
                  p_1  &+  &p_2  &+  &p_3  &+  &p_4  &=  &0  \\
                       &   &     &   &     &   &0    &=  &0  \\
                 -p_1  &   &     &   &     &-  &2p_4 &=  &0
                \end{linsys}
                \grstep{\rho_1+\rho_4}
                \begin{linsys}{4}
                  p_1  &+  &p_2  &+  &p_3  &+  &p_4  &=  &0  \\
                       &   &p_2  &+  &p_3  &-  &p_4  &=  &0
                \end{linsys}
              \end{equation*}
              取 $p_3$ 与 $p_4$ 作为参数，我们可以这样描述
              解集。
              \begin{equation*}
                \set{\colvec[r]{0 \\ -1 \\ 1 \\ 0}p_3
                     +\colvec[r]{-2 \\ 1 \\ 0 \\ 1}p_4
                     \suchthat p_3,p_4\in\Re}
              \end{equation*}
              这给出 $\set{\Pi_1=h/d,\Pi_2=dg/v^2}$ 作为一个完备集合。
            \partsitem Buckingham 定理说
              $v^2=dg\cdot\hat{f}(h/d)$，于是由于 $g$ 是常数，
              如果 $h/d$ 固定，那么 $v$ 与 $\sqrt{d\,}$ 成正比。
          \end{exparts}
        
\end{ans}
\begin{ans}{5}
        逐条检验向量空间定义中的条件是例行的。
      
\end{ans}
\begin{ans}{6}
      \begin{exparts}
        \partsitem 周长的量纲公式是 $L$，
          即 $L^1M^0T^0$。
          面积的量纲公式是 $L^2$。
        \partsitem 一个是 $C+A=2\pi r + \pi r^2$。
        \partsitem 一个例子是把角所对的弧长
           与半径及以弧度计的角
           联系起来的这个公式：~$\ell-r\theta=0$。
           该公式中两项的量纲公式都是
           $L^1$。
           这一关系对某些单位制
           （例如英寸与弧度）成立，
           但并非对所有单位制都成立
           （例如英寸与度）。
      \end{exparts}
    
\end{ans}
\protect \chapter {空间之间的映射}
\protect \section {第 I 节：同构}
\protect \subsection {三.I.1: 定义与例子}
\begin{ans}{三.I.1.12}
     \begin{exparts}
     \partsitem 把这个映射记为 $f$。
       \begin{equation*}
         \rowvec{a &b} \mapsunder{f} \colvec{a \\ b}
       \end{equation*}
       它是单射，因为若 $f$ 把定义域中的两个元素映到同一个像，即若
       $f\left(\rowvec{a &b}\right)=f\left(\rowvec{c &d}\right)$，
       则由 $f$ 的定义可得
       \begin{equation*}
          \colvec{a \\ b}=\colvec{c \\ d}
       \end{equation*}
       由于相等的列向量具有相等的对应分量，
       我们有 $a=c$ 且 $b=d$。
       因此，若 $f$ 把定义域中的两个行向量映到同一个列向量，
       则这两个行向量相等：
       $\rowvec{a &b}=\rowvec{c &d}$。

       要证明 $f$ 是满射，就必须证明陪域 $\Re^2$ 中的每个元素
       都是某个行向量在 $f$ 下的像。
       这很容易；
       \begin{equation*}
          \colvec{x \\ y}
       \end{equation*}
       就是 $f\left(\rowvec{x &y}\right)$。

       验证保持加法的计算如下。
       \begin{multline*}
         f\left(\rowvec{a &b}+\rowvec{c &d}\right)
           = f\left(\rowvec{a+c &b+d}\right)
           = \colvec{a+c \\ b+d}                               \\
           = \colvec{a \\ b}+\colvec{c \\ d}
           = f\left(\rowvec{a &b}\right)+f\left(\rowvec{c &d}\right)
       \end{multline*}
       验证保持标量乘法的计算与此类似。
       \begin{equation*}
         f\left(r\cdot\rowvec{a &b}\right)
           = f\left(\rowvec{ra &rb}\right)
           = \colvec{ra \\ rb}
           = r\cdot\colvec{a \\ b}
           = r\cdot f\left(\rowvec{a &b}\right)
       \end{equation*}
     \partsitem 把\nearbyexample{exam:PolyTwoIsoRThree} 中的映射记为
      $f$。
      要证明它是单射，
      假设 \( f(a_0+a_1x+a_2x^2)=f(b_0+b_1x+b_2x^2) \)。
      则由函数的定义，
      \begin{equation*}
        \colvec{a_0 \\ a_1 \\ a_2}
        =\colvec{b_0 \\ b_1 \\ b_2}
      \end{equation*}
      因而 \( a_0=b_0 \)、\( a_1=b_1 \) 且 \( a_2=b_2 \)。
      所以 \( a_0+a_1x+a_2x^2=b_0+b_1x+b_2x^2 \)，从而 \( f \)
      是单射。

      函数 $f$ 是满射，因为对于
      \begin{equation*}
        \colvec{a \\ b \\ c}
      \end{equation*}
      总有多项式 \( a+bx+cx^2 \) 经 \( f \) 映到它。

      至于结构，
      下面的计算表明 $f$ 保持加法
      \begin{multline*}
        f\left(\,(a_0+a_1x+a_2x^2)+(b_0+b_1x+b_2x^2)\,\right)           \\
        \begin{aligned}
        &=f\left(\,(a_0+b_0)+(a_1+b_1)x+(a_2+b_2)x^2\,\right)       \\
        &=\colvec{a_0+b_0 \\ a_1+b_1 \\ a_2+b_2}                    \\
        &=\colvec{a_0 \\ a_1 \\ a_2} +\colvec{b_0 \\ b_1 \\ b_2}   \\
        &=f(a_0+a_1x+a_2x^2)+f(b_0+b_1x+b_2x^2)
        \end{aligned}
      \end{multline*}
      而下面的计算
      \begin{align*}
        f(\,r(a_0+a_1x+a_2x^2)\,)
        &=f(\,(ra_0)+(ra_1)x+(ra_2)x^2\,)    \\
        &=\colvec{ra_0 \\ ra_1 \\ ra_2}     \\
        &=r\cdot\colvec{a_0 \\ a_1 \\ a_2}     \\
        &=r\,f(a_0+a_1x+a_2x^2)
      \end{align*}
      表明它保持标量乘法。
    \end{exparts}
   
\end{ans}
\begin{ans}{三.I.1.13}
       各元素的像如下。
       \begin{exparts*}
        \partsitem \( \colvec[r]{5 \\ -2} \)
        \partsitem \( \colvec[r]{0 \\ 2} \)
        \partsitem \( \colvec[r]{-1 \\ 1} \)
      \end{exparts*}

      要证明 $f$ 是单射，假设它把两个次数不超过一的多项式
      映到同一个像，即 $f(a_1+b_1x)=f(a_2+b_2x)$。
      则
      \begin{equation*}
        \colvec{a_1-b_1 \\ b_1}=\colvec{a_2-b_2 \\ b_2}
      \end{equation*}
      由于相等的列向量具有相等的对应分量，
      可得 $b_1=b_2$ 且 $a_1=a_2$。
      这说明两个次数不超过一的多项式相等，因此 $f$ 是单射。

      要证明 $f$ 是满射，只需注意陪域中的元素
      \begin{equation*}
         \colvec{s \\ t}
      \end{equation*}
      就是定义域中元素 $(s+t)+tx$ 的像。

      要验证 $f$ 保持结构，
      可以使用\nearbylemma{le:PresStructIffPresCombos} 的第~(2) 项。
      \begin{align*}
        f\left(c_1\cdot (a_1+b_1x)+c_2\cdot (a_2+b_2x)\right)
        &=f\left((c_1a_1+c_2a_2)+(c_1b_1+c_2b_2)x\right)                \\
        &=\colvec{(c_1a_1+c_2a_2)-(c_1b_1+c_2b_2) \\ c_1b_1+c_2b_2}       \\
        &=c_1\cdot\colvec{a_1-b_1 \\ b_1}+c_2\cdot\colvec{a_2-b_2 \\ b_2}  \\
        &=c_1\cdot f(a_1+b_1x)+c_2\cdot f(a_2+b_2x)
      \end{align*}
     
\end{ans}
\begin{ans}{三.I.1.14}
      要验证它是单射，假设
       $f_1(c_1x+c_2y+c_3z)=f_1(d_1x+d_2y+d_3z)$。
      由 $f_1$ 的定义，得到 $c_1+c_2x+c_3x^2=d_1+d_2x+d_3x^2$。
      $\polyspace_2$ 中的元素相等时，对应系数必定相等，
      因而 $c_1=d_1$、$c_2=d_2$ 且 $c_3=d_3$。
      因此 $f_1(c_1x+c_2y+c_3z)=f_1(d_1x+d_2y+d_3z)$ 蕴含
      $c_1x+c_2y+c_3z=d_1x+d_2y+d_3z$，所以 $f_1$ 是单射。

      要验证它是满射，考虑陪域中的任意元素 $a_1+a_2x+a_3x^2$，
      注意它确实是定义域中某个元素的像，
      即它等于 $f_1(a_1x+a_2y+a_3z)$。
      （例如，$0+3x+6x^2=f_1(0x+3y+6z)$。）

      验证 $f_1$ 保持加法的计算如下。
      \begin{multline*}
        f_1\left(\,(c_1x+c_2y+c_3z)+(d_1x+d_2y+d_3z)\,\right)            \\
          \begin{aligned}
          &=f_1\left(\,(c_1+d_1)x+(c_2+d_2)y+(c_3+d_3)z\,\right)  \\
          &=(c_1+d_1)+(c_2+d_2)x+(c_3+d_3)x^2  \\
          &=(c_1+c_2x+c_3x^2)+(d_1+d_2x+d_3x^2)  \\
          &=f_1(c_1x+c_2y+c_3z)+f_1(d_1x+d_2y+d_3z)
          \end{aligned}
      \end{multline*}

      验证 $f_1$ 保持标量乘法的计算如下。
      \begin{align*}
        f_1(\,r\cdot (c_1x+c_2y+c_3z)\,)
          &=f_1(\,(rc_1)x+(rc_2)y+(rc_3)z\,)  \\
          &=(rc_1)+(rc_2)x+(rc_3)x^2  \\
          &=r\cdot (c_1+c_2x+c_3x^2)  \\
          &=r\cdot f_1(c_1x+c_2y+c_3z)
      \end{align*}
    
\end{ans}
\begin{ans}{三.I.1.15}
      要证明这个映射是单射，假设 $t(\vec{v}_1)=t(\vec{v}_2)$，
      目标是推出 $\vec{v}_1=\vec{v}_2$。
      也就是说，$t(a_1x^2+b_1x+c_1)=t(a_2x^2+b_2x+c_2)$。
      则 $b_1x^2-(a_1+c_1)x+a_1=b_2x^2-(a_2+c_2)x+a_2$。
      次数不超过二的多项式相等时，二次项、常数项和一次项的系数分别相等，
      因而 $b_1=b_2$、$a_1=a_2$，并由此得到 $c_1=c_2$。
      所以 $a_1x^2+b_1x+c_1=a_2x^2+b_2x+c_2$，这个函数是单射。

      要证明这个映射是满射，假设给定陪域中的一个元素~$\vec{w}$，
      我们来找定义域中映到它的元素~$\vec{v}$。
      设这个陪域元素为~$\vec{w}=px^2+qx+r$。
      注意，当 $\vec{v}=rx^2+px+(-q-r)$ 时，$t(\vec{v})=\vec{w}$。
      因此~$t$ 是满射。

      要证明这个映射是同态，我们验证它保持两个元素的线性组合。
      根据\nearbylemma{le:PresStructIffPresCombos}，这就说明该映射保持运算。
      \begin{multline*}
        t(r_1(a_1x^2+b_1x+c_1)+r_2(a_2x^2+b_2x+c_2))              \\
        \begin{split} \quad
        &=t((r_1a_1+r_2a_2)x^2+(r_1b_1+r_2b_2)x+(r_1c_1+r_2c_2))   \\
        &=(r_1b_1+r_2b_2)x^2-((r_1a_1+r_2a_2)+(r_1c_1+r_2c_2))x+(r_1a_1+r_2a_2)  \\
        &=(r_1b_1)x^2-(r_1a_1+r_1c_1)x+r_1a_1
           +(r_2b_2)x^2-(r_2a_2+r_2c_2)x+r_2a_2                       \\
        &=r_1t(a_1x^2+b_1x+c_1)+r_2t(a_2x^2+b_2x+c_2)
        \end{split}
      \end{multline*}
    
\end{ans}
\begin{ans}{三.I.1.16}
      先验证 $h$ 是单射。
      为此，我们将证明 $h(\vec{v}_1)=h(\vec{v}_2)$ 蕴含
      $\vec{v}_1=\vec{v}_2$。
      假设
      \begin{equation*}
          h(\vec{v}_1)
          =
          h(\colvec{a_1 \\ b_1 \\ c_1 \\ d_1})
          =
          h(\colvec{a_2 \\ b_2 \\ c_2 \\ d_2})
          =h(\vec{v}_2)
      \end{equation*}
      则有
      \begin{equation*}
          \begin{mat}
            c_1  &a_1+d_1 \\
            b_1  &d_1
          \end{mat}
          =
          \begin{mat}
            c_2  &a_2+d_2 \\
            b_2  &d_2
          \end{mat}
      \end{equation*}
      由此可得
      $c_1=c_2$（比较左上角的元素），
      $b_1=b_2$（比较左下角的元素），
      $d_1=d_2$（比较右下角的元素），
      再结合最后一个等式得到 $a_1=a_2$（比较右上角的元素）。
      因此 $\vec{v}_1=\vec{v}_2$。

      接下来证明这个映射是满射，即陪域 $\matspace_{\nbyn{2}}$ 中的每个元素
      都是定义域中某个四维列向量的像。
      给定
      \begin{equation*}
          \vec{w}=
          \begin{mat}
            m  &n \\
            p  &q
          \end{mat}
          \in\matspace_{\nbyn{2}}
      \end{equation*}
      注意，它是定义域中下列向量的像。
      \begin{equation*}
        \vec{v}=
        \colvec{n-q \\ p \\ m \\ q}
      \end{equation*}

      最后验证这个映射保持线性组合。
      根据\nearbylemma{le:PresStructIffPresCombos}，这就说明该映射保持运算。
      \begin{align*}
          h(r_1\cdot\colvec{a_1 \\ b_1 \\ c_1 \\ d_1}
            +
            r_2\cdot\colvec{a_2 \\ b_2 \\ c_2 \\ d_2})
          &=h(\colvec{r_1a_1+r_2a_2 \\ r_1b_1+r_2b_2 \\ r_1c_1+r_2c_2 \\ r_1d_1+r_2d_2})  \\
          &=
          \begin{mat}
             r_1c_1+r_2c_2 &(r_1a_1+r_2a_2)+(r_1d_1+r_2d_2)  \\
             r_1b_1+r_2b_2 &r_1d_1+r_2d_2
          \end{mat}                               \\
          &=
          r_1\begin{mat}
             c_1 &a_1+d_1  \\
             b_1 &d_1
          \end{mat}
          +
          r_2\begin{mat}
             c_2 &a_2+d_2  \\
             b_2 &d_2
          \end{mat}                               \\
         &=r_1\cdot h(\colvec{a_1 \\ b_1 \\ c_1 \\ d_1})
            +
            r_2\cdot h(\colvec{a_2 \\ b_2 \\ c_2 \\ d_2})
      \end{align*}
    
\end{ans}
\begin{ans}{三.I.1.17}
    \begin{exparts}
      \partsitem 不是；这个映射不是单射。
        例如，元素全为零的矩阵与元素全为一的矩阵具有相同的像。
      \partsitem 是，这个映射是同构。

        它是单射：
        \begin{align*}
          &\text{若 }
          f(\begin{mat}
              a_1  &b_1  \\
              c_1  &d_1
            \end{mat})
         =f(\begin{mat}
              a_2  &b_2  \\
              c_2  &d_2
            \end{mat})                                     \\
          &\qquad\text{ 则 }
          \colvec{a_1+b_1+c_1+d_1 \\ a_1+b_1+c_1 \\ a_1+b_1 \\ a_1}
          =\colvec{a_2+b_2+c_2+d_2 \\ a_2+b_2+c_2 \\ a_2+b_2 \\ a_2}
        \end{align*}
        给出 \( a_1=a_2 \)、\( b_1=b_2 \)、
        \( c_1=c_2 \) 且 \( d_1=d_2 \)。

        它是满射，因为下式
        \begin{equation*}
          \colvec{x \\ y \\ z \\ w}
         =f(\begin{mat}
              w    &z-w  \\
              y-z  &x-y
            \end{mat})
        \end{equation*}
        表明每个四维列向量都是某个 $\nbyn{2}$~矩阵的像。

        最后，它保持线性组合
        \begin{multline*}
          f(\,r_1\cdot\begin{mat}
              a_1  &b_1  \\
              c_1  &d_1
            \end{mat}
            +r_2\cdot\begin{mat}
              a_2  &b_2  \\
              c_2  &d_2
            \end{mat}\,)                                   \\
          \begin{aligned}
          &=f(\begin{mat}
              r_1a_1+r_2a_2  &r_1b_1+r_2b_2  \\
              r_1c_1+r_2c_2  &r_1d_1+r_2d_2
            \end{mat})                            \\
          &=\colvec{r_1a_1+\dots+r_2d_2 \\ r_1a_1+\dots+r_2c_2 \\
                       r_1a_1+\dots+r_2b_2 \\ r_1a_1+r_2a_2}          \\
          &=r_1\cdot\colvec{a_1+\dots+d_1 \\ a_1+\dots+c_1 \\
                       a_1+b_1 \\ a_1}
          +r_2\cdot\colvec{a_2+\dots+d_2 \\ a_2+\dots+c_2 \\
                       a_2+b_2 \\ a_2}                         \\
          &=r_1\cdot f(\begin{mat}
                  a_1  &b_1  \\
                  c_1  &d_1
            \end{mat})
          +r_2\cdot f(\begin{mat}
                  a_2  &b_2  \\
                  c_2  &d_2
            \end{mat})
         \end{aligned}
        \end{multline*}
        因此，根据\nearbylemma{le:PresStructIffPresCombos} 的第~(2) 项，
        它保持结构。
      \partsitem 是，这个映射是同构。

        要证明它是单射，假设定义域中的两个元素在 $f$ 下具有相同的像。
        \begin{equation*}
          f(\begin{mat}
              a_1  &b_1  \\
              c_1  &d_1
            \end{mat})
          =f(\begin{mat}
              a_2  &b_2  \\
              c_2  &d_2
            \end{mat})
        \end{equation*}
        根据 $f$ 的定义，可得
        $c_1+(d_1+c_1)x+(b_1+a_1)x^2+a_1x^3
          =c_2+(d_2+c_2)x+(b_2+a_2)x^2+a_2x^3$。
        再根据相等多项式的对应系数相等，得到线性方程组
        \begin{align*}
          c_1      &=  c_2      \\
          d_1+c_1  &=  d_2+c_2  \\
          b_1+a_1  &=  b_2+a_2  \\
          a_1      &=  a_2
        \end{align*}
        它只有一个解：$a_1=a_2$、$b_1=b_2$、$c_1=c_2$ 且
        $d_1=d_2$。

        要证明 $f$ 是满射，只需注意 $p+qx+rx^2+sx^3$
        是下面这个矩阵在 $f$ 下的像。
        \begin{equation*}
          \begin{mat}
            s  &r-s   \\
            p  &q-p
          \end{mat}
        \end{equation*}

        可以利用\nearbylemma{le:PresStructIffPresCombos} 的第~(2) 项，
        验证 $f$ 保持结构。
        \begin{multline*}
          f(r_1\cdot\begin{mat}
              a_1  &b_1  \\
              c_1  &d_1
            \end{mat}
            +r_2\cdot\begin{mat}
              a_2  &b_2  \\
              c_2  &d_2
            \end{mat})                                      \\
           \begin{aligned}
           &=
           f(\begin{mat}
               r_1a_1+r_2a_2  &r_1b_1+r_2b_2  \\
               r_1c_1+r_2c_2  &r_1d_1+r_2d_2
             \end{mat})                         \\
           &=\begin{aligned}[t]
               &(r_1c_1+r_2c_2)
               +(r_1d_1+r_2d_2+r_1c_1+r_2c_2)x                 \\
               &\mbox{}\quad +(r_1b_1+r_2b_2+r_1a_1+r_2a_2)x^2
                 +(r_1a_1+r_2a_2)x^3
             \end{aligned}                   \\
           &=\begin{aligned}[t]
                &r_1\cdot\left(c_1+(d_1+c_1)x+(b_1+a_1)x^2+a_1x^3\right) \\
                &\mbox{}\quad +r_2\cdot\left(c_2+(d_2+c_2)x
                            +(b_2+a_2)x^2+a_2x^3\right)
              \end{aligned} \\
           &=r_1\cdot f(\begin{mat}
              a_1  &b_1  \\
              c_1  &d_1
            \end{mat})
            +r_2\cdot f(\begin{mat}
              a_2  &b_2  \\
              c_2  &d_2
            \end{mat})
            \end{aligned}
        \end{multline*}
      \partsitem 不是，这个映射不保持结构。
        例如，它没有把元素全为零的矩阵映到零多项式。
    \end{exparts}
   
\end{ans}
\begin{ans}{三.I.1.18}
      它既是单射又是满射，即是一一对应，
      因为它有逆映射（即 \( f^{-1}(x)=\sqrt[3]{x} \)）。
      不过，它不是同构。
      例如，\( f(1)+f(1)\neq f(1+1) \)。
    
\end{ans}
\begin{ans}{三.I.1.19}
     这样的映射有很多。
     下面给出两个。
     \begin{equation*}
       \rowvec{a &b}\mapsto\colvec{b \\ a}
       \quad\text{且}\quad
       \rowvec{a &b}\mapsto\colvec{2a \\ b}
     \end{equation*}
     参照前面的验证稍作调整，即可完成验证。
    
\end{ans}
\begin{ans}{三.I.1.20}
      下面给出两个。
      \begin{equation*}
        a_0+a_1x+a_2x^2 \mapsto \colvec{a_1 \\ a_0 \\ a_2}
        \quad\text{且}\quad
        a_0+a_1x+a_2x^2 \mapsto \colvec{a_0+a_1 \\ a_1 \\ a_2}
      \end{equation*}
      验证很直接（对于第二个映射，要证明它是满射，只需注意
      \begin{equation*}
        \colvec{s \\ t \\ u}
      \end{equation*}
      是 $(s-t)+tx+ux^2$ 的像）。
     
\end{ans}
\begin{ans}{三.I.1.21}
       空间 $\Re^2$ 不是 $\Re^3$ 的子空间，因为它不是 $\Re^3$ 的子集。
       $\Re^2$ 中的二维列向量不是 $\Re^3$ 的元素。

       自然映射 \( \map{\iota}{\Re^2}{\Re^3} \)
       （称为\definend{嵌入}映射）如下；它给出到该平面的同构。
       \begin{equation*}
         \colvec{x \\ y}
           \mapsunder{\iota}
         \colvec{x \\ y \\ 0}
       \end{equation*}

       这个映射是单射，因为
       \begin{equation*}
         f(\colvec{x_1 \\ y_1})=f(\colvec{x_2 \\ y_2})
         \quad\text{蕴含}\quad
         \colvec{x_1 \\ y_1 \\ 0}=\colvec{x_2 \\ y_2 \\ 0}
       \end{equation*}
       进而推出 $x_1=x_2$ 且 $y_1=y_2$，
       因此最初的两个二维列向量相等。

       由于
       \begin{equation*}
         \colvec{x \\ y \\ 0}
         =f(\colvec{x \\ y})
       \end{equation*}
       这个映射到 $xy$ 平面是满射。

       要证明这个映射保持结构，利用\nearbylemma{le:PresStructIffPresCombos}
       的第~(2) 项，通过下式
       \begin{multline*}
         f(c_1\cdot \colvec{x_1 \\ y_1}+c_2\cdot \colvec{x_2 \\ y_2})
          = f(\colvec{c_1x_1+c_2x_2 \\ c_1y_1+c_2y_2})
          =\colvec{c_1x_1+c_2x_2 \\ c_1y_1+c_2y_2 \\ 0}                     \\
          =c_1\cdot \colvec{x_1 \\ y_1 \\ 0}+c_2\cdot \colvec{x_2 \\ y_2 \\ 0}
          =c_1\cdot f(\colvec{x_1 \\ y_1})+c_2\cdot f(\colvec{x_2 \\ y_2})
       \end{multline*}
       验证它保持两个向量的线性组合。
     
\end{ans}
\begin{ans}{三.I.1.22}
        下面给出两个：
        \begin{equation*}
           \colvec{r_1 \\ r_2 \\ \vdots \\ r_{16}}
              \mapsto
           \begin{mat}
              r_1  &r_2  &\ldots  \\
                   &              \\
                   &     &\ldots  &r_{16}
           \end{mat}
           \quad\text{且}\quad
           \colvec{r_1 \\ r_2 \\ \vdots \\ r_{16}}
              \mapsto
           \begin{mat}
              r_1    &                    \\
              r_2    &                    \\
              \vdots &   &        &\vdots \\
                   &     &        &r_{16}
           \end{mat}
        \end{equation*}
        很容易验证它们都是同构。
    
\end{ans}
\begin{ans}{三.I.1.23}
        当 $k$ 为乘积 \( k=mn \) 时，下面给出一个同构。
        \begin{equation*}
           \begin{mat}
              r_1  &r_2  &\ldots  \\
                   &\vdots        \\
                   &     &\ldots  &r_{m\cdot n}
           \end{mat}
           \mapsto
           \colvec{r_1 \\ r_2 \\ \vdots \\ r_{m\cdot n}}
        \end{equation*}
        很容易验证它是同构。
      
\end{ans}
\begin{ans}{三.I.1.24}
        若 \( n\geq 1 \)，则 \( \polyspace_{n-1}\isomorphicto\Re^n \)。
        （若把 \( \polyspace_{-1} \) 和 \( \Re^0 \) 视为平凡向量空间，
        则这一关系还可以向下延伸一个维数。）
        它们之间的自然同构如下。
        \begin{equation*}
          a_0+a_1x+\dots+a_{n-1}x^{n-1}
          \mapsto\colvec{a_0 \\ a_1 \\ \vdots \\ a_{n-1}}
        \end{equation*}
        验证它是同构很直接。
     
\end{ans}
\begin{ans}{三.I.1.25}
      把这个映射展开，得到
      \begin{multline*}
        f(a_0+a_1x+a_2x^2+a_3x^3+a_4x^4+a_5x^5)                        \\
        \begin{aligned}
        &=a_0+a_1(x-1)+a_2(x-1)^2+a_3(x-1)^3               \\
        &\mbox{}\quad +a_4(x-1)^4+a_5(x-1)^5               \\
        &=a_0+a_1(x-1)+a_2(x^2-2x+1)         \\
        &\mbox{}\quad +a_3(x^3-3x^2+3x-1)   \\
        &\mbox{}\quad +a_4(x^4-4x^3+6x^2-4x+1) \\
        &\mbox{}\quad +a_5(x^5-5x^4+10x^3-10x^2+5x-1)       \\
        &=(a_0-a_1+a_2-a_3+a_4-a_5)  \\
        &\mbox{}\quad +(a_1-2a_2+3a_3-4a_4+5a_5)x  \\
        &\mbox{}\quad +(a_2-3a_3+6a_4-10a_5)x^2  \\
        &\quad+(a_3-4a_4+10a_5)x^3         \\
        &\mbox{}\quad +(a_4-5a_5)x^4
                        +a_5x^5
        \end{aligned}
      \end{multline*}
      这个映射是一一对应，因为它有逆映射，
      即 \( p(x)\mapsto p(x+1) \)。

      最后，利用\nearbylemma{le:PresStructIffPresCombos} 的第~(2) 项，
      验证它保持两个多项式的线性组合，便可证明它是同构。
      简言之，
      $f(c\cdot (a_0+a_1x+\dots +a_5x^5)+d\cdot (b_0+b_1x+\dots +b_5x^5))$
      等于
      \begin{multline*}
         (ca_0-ca_1+ca_2-ca_3+ca_4-ca_5+db_0-db_1+db_2-db_3+db_4-db_5)  \\
          +\dots+(ca_5+db_5)x^5
      \end{multline*}
      而这又等于
      $c\cdot f(a_0+a_1x+\dots +a_5x^5)+d\cdot f(b_0+b_1x+\dots +b_5x^5)$。
    
\end{ans}
\begin{ans}{三.I.1.26}
       向量空间的底层集合不可能是空集。
       可以取 $\vec{v}$ 为零向量。
     
\end{ans}
\begin{ans}{三.I.1.27}
        是的；若两个空间分别为 \( \set{\vec{a}} \) 和
        \( \set{\vec{b}} \)，
        则把 \( \vec{a} \) 映到 \( \vec{b} \) 的映射显然既是单射又是满射，
        而且也保持了其中仅有的那些结构。
      
\end{ans}
\begin{ans}{三.I.1.28}
       $n=0$ 个向量的线性组合等于零向量，
       因此由\nearbylemma{le:IsoSendsZeroToZero} 可知，
       这三个命题在此情形下也等价。
    
\end{ans}
\begin{ans}{三.I.1.29}
      考虑 \( \polyspace_0 \) 的基 \( \sequence{1} \)，
      把 \( f(1)\in\Re \) 记为 \( k \)。
      对任意 \( a\in\polyspace_0 \)，都有
      \( f(a)=f(a\cdot 1)=af(1)=ak \)，
      所以 \( f \) 的作用就是乘以 \( k \)。
      注意 \( k\neq 0 \)，否则这个映射不是单射。
      （顺便说一句，任何这样的映射 $a\mapsto ka$ 都是同构，这很容易验证。）
     
\end{ans}
\begin{ans}{三.I.1.30}
      在每一题中，我们都依据\nearbylemma{le:PresStructIffPresCombos}
      的第~(2) 项，通过验证映射保持定义域中两个元素的线性组合，
      来证明它保持结构。
      \begin{exparts}
        \partsitem
          恒等映射显然既是单射又是满射。
          保持线性组合也很容易验证。
          \begin{equation*}
            \mbox{id}(c_1\cdot \vec{v}_1+c_2\cdot \vec{v}_2)
             =c_1\vec{v}_1+c_2\vec{v}_2
             =c_1\cdot \mbox{id}(\vec{v}_1)+c_2\cdot \mbox{id}(\vec{v}_2)
          \end{equation*}
        \partsitem 一一对应的逆映射仍是一一对应（见附录），
          因此只需验证逆映射保持线性组合。
          假设 \( \vec{w}_1=f(\vec{v}_1) \)
          （因而 \( f^{-1}(\vec{w}_1)=\vec{v}_1 \)），
          并假设 \( \vec{w}_2=f(\vec{v}_2) \)。
          \begin{align*}
            f^{-1}(c_1\cdot\vec{w}_1+c_2\cdot\vec{w}_2)
              &=f^{-1}\bigl(\,c_1\cdot f(\vec{v}_1)
                             +c_2\cdot f(\vec{v}_2)\,\bigr) \\
              &=f^{-1}(\,f\bigl(c_1\vec{v}_1+c_2\vec{v}_2)\,\bigr)    \\
              &=c_1\vec{v}_1+c_2\vec{v}_2                             \\
              &=c_1\cdot f^{-1}(\vec{w}_1)+c_2\cdot f^{-1}(\vec{w}_2)
          \end{align*}
        \partsitem 两个一一对应的复合仍是一一对应（见附录），
          因此只需验证复合映射保持线性组合。
          \begin{align*}
           \composed{g}{f}\,\bigl(c_1\cdot\vec{v}_1+c_2\cdot\vec{v}_2\bigr)
             &=g\bigl(\,f(c_1\vec{v}_1+c_2\vec{v}_2)\,\bigr)        \\
             &=g\bigl(\,c_1\cdot f(\vec{v}_1)+c_2\cdot f(\vec{v}_2)\,\bigr) \\
             &=c_1\cdot g\bigl(f(\vec{v}_1))+c_2\cdot g(f(\vec{v}_2)\bigr)  \\
             &=c_1\cdot\composed{g}{f}\,(\vec{v}_1)
                  +c_2\cdot\composed{g}{f}\,(\vec{v}_2)
          \end{align*}
      \end{exparts}
     
\end{ans}
\begin{ans}{三.I.1.31}
      一个方向很容易：~根据定义，若 \( f \) 是单射，
      则对任意 \( \vec{w}\in W \)，至多有一个 \( \vec{v}\in V \) 满足
      \( f(\vec{v}\,)=\vec{w} \)；特别地，
      \( V \) 中至多有一个元素被映到 \( \zero_W \)。
      \nearbylemma{le:IsoSendsZeroToZero} 的证明没有用到映射是一一对应这一条件，
      因而它实际上表明，任何保持结构的映射
      \( f \) 都把 \( \zero_V \) 映到 \( \zero_W \)。

      反过来，假设 \( V \) 中被映到 \( \zero_W \) 的唯一元素是 \( \zero_V \)。
      要证明 \( f \) 是单射，假设 \( f(\vec{v}_1)=f(\vec{v}_2) \)。
      则 \( f(\vec{v}_1)-f(\vec{v}_2)=\zero_W \)，从而
      \( f(\vec{v}_1-\vec{v}_2)=\zero_W \)。
      于是 \( \vec{v}_1-\vec{v}_2=\zero_V \)，所以
      $\vec{v}_1=\vec{v}_2$，即 \( f \) 是单射。
    
\end{ans}
\begin{ans}{三.I.1.32}
      我们将证明一个更强的结论\Dash 同构不仅保持线性相关关系的存在性，
      而且保持每一个具体的线性相关关系，即
      \begin{multline*}
        \vec{v}_i=c_1\vec{v}_1+\cdots+c_{i-1}\vec{v}_{i-1}
                   +c_{i+1}\vec{v}_{i+1}+\cdots+c_k\vec{v}_k   \\
        \iff
        f(\vec{v}_i)=c_1f(\vec{v}_1)+\cdots+c_{i-1}f(\vec{v}_{i-1})
                     +c_{i+1}f(\vec{v}_{i+1})+\cdots+c_kf(\vec{v}_k).
      \end{multline*}
      这个命题的 \( \implies \) 方向由\nearbylemma{le:PresStructIffPresCombos}
      的第~(3) 项可得。
      对于 \( \isimpliedby \) 方向，先将各项合并
      \begin{align*}
         f(\vec{v}_i)
           &=c_1f(\vec{v}_1)+\dots+c_{i-1}f(\vec{v}_{i-1})
                     +c_{i+1}f(\vec{v}_{i+1})+\dots+c_kf(\vec{v}_k)  \\
           &=f(c_1\vec{v}_1+\dots+c_{i-1}\vec{v}_{i-1}
                     +c_{i+1}\vec{v}_{i+1}+\dots+c_k \vec{v}_k)
      \end{align*}
      再利用 $f$ 是单射这一事实：两个向量 $\vec{v}_i$ 与
      $c_1\vec{v}_1+\dots+c_{i-1}\vec{v}_{i-1}
              +c_{i+1}f\vec{v}_{i+1}+\dots+c_kf(\vec{v}_k$
      若被 $f$ 映到同一个像，就必定相等。
    
\end{ans}
\begin{ans}{三.I.1.33}
      \begin{exparts}
        \partsitem
          这个映射是单射，因为若 \( d_s(\vec{v}_1)=d_s(\vec{v}_2) \)，
          则由定义有 \( s\cdot\vec{v}_1=s\cdot\vec{v}_2 \)，
          而 \( s \) 非零，所以 \( \vec{v}_1=\vec{v}_2 \)。
          这个映射是满射，因为任意 \( \vec{w}\in\Re^2 \) 都是
          \( \vec{v}=(1/s)\cdot\vec{w} \) 的像（这里同样用到 $s$ 非零）。
          （也可以通过逆映射看出它是一一对应：
          \( d_s \) 的逆映射是 \( d_{1/s} \)。）

          最后，注意这个映射保持线性组合
          \begin{equation*}
            d_s(c_1\cdot\vec{v}_1+c_2\cdot\vec{v}_2)
            =s(c_1\vec{v}_1+c_2\vec{v}_2)
            =c_1s\vec{v}_1+c_2s\vec{v}_2
            =c_1\cdot d_s(\vec{v}_1)+c_2\cdot d_s(\vec{v}_2)
          \end{equation*}
          因此它是同构。
        \partsitem 与前一题一样，注意到映射 $t_\theta$ 有逆映射 $t_{-\theta}$，
          就能说明它是一一对应。

          从几何上很容易看出这个映射保持结构。
          例如，先将两个向量相加再旋转，与先旋转再相加，效果相同。
          为给出代数论证，考虑极坐标：
          映射 \( t_\theta \) 把终点为 \( (r,\phi) \) 的向量
          映到终点为 \( (r,\phi+\theta) \) 的向量。
          利用熟悉的三角公式
          $\cos(\phi+\theta)=\cos\phi\,\cos\theta-\sin\phi\,\sin\theta$
          和
          $\sin(\phi+\theta)=\sin\phi\,\cos\theta+\cos\phi\,\sin\theta$，
          可以把这个映射的作用写成通常的直角坐标形式。
          \begin{equation*}
            \colvec{x \\ y}=\colvec{r\cos\phi \\ r\sin\phi}
              \mapsunder{t_\theta}
            \colvec{r\cos(\phi+\theta) \\ r\sin(\phi+\theta)}
            =\colvec{x\cos\theta-y\sin\theta \\ x\sin\theta+y\cos\theta}
          \end{equation*}
          现在，验证保持加法只需直接计算。
          \begin{multline*}
            \colvec{x_1+x_2 \\ y_1+y_2}
              \mapsunder{t_\theta}
            \colvec{(x_1+x_2)\cos\theta-(y_1+y_2)\sin\theta \\
                       (x_1+x_2)\sin\theta+(y_1+y_2)\cos\theta}    \\
            =\colvec{x_1\cos\theta-y_1\sin\theta \\
                       x_1\sin\theta+y_1\cos\theta}
            +\colvec{x_2\cos\theta-y_2\sin\theta \\
                       x_2\sin\theta+y_2\cos\theta}
          \end{multline*}
          验证保持标量乘法的计算与此类似。
        \partsitem
          这个映射是一一对应，因为它有逆映射（就是它自身）。

          与前一题一样，从几何上很容易看出反射映射保持结构：~例如，
          先将向量相加再反射，与先反射再相加，结果相同。
          为给出代数证明，假设直线 $\ell$ 的斜率为 $k$
          （斜率不存在的情形可以单独处理，也很容易）。
          按照提示，使用极坐标：~若直线 \( \ell \) 与 \( x \) 轴的夹角为
          \( \phi \)，则 \( f_\ell \) 的作用是
          把终点为 \( (r\cos\theta,r\sin\theta) \) 的向量映到
          终点为 \( (r\cos(2\phi-\theta),r\sin(2\phi-\theta)) \) 的向量。
          \begin{center}  \small
            \includegraphics{map/mp/ch3.72}
          \end{center}
          与前一题一样，我们使用三角公式转为直角坐标。
          首先，注意 $\cos\phi$ 和 $\sin\phi$ 可由直线的斜率 $k$ 确定。
          下图
          \begin{center}  \small
            \includegraphics{map/mp/ch3.73}
          \end{center}
          给出 $\cos\phi=1/\sqrt{1+k^2}$ 和 $\sin\phi=k/\sqrt{1+k^2}$。
          现在，
          \begin{align*}
            \cos(2\phi-\theta)
              &=\cos(2\phi)\,\cos\theta+\sin(2\phi)\,\sin\theta        \\
              &=\left(\cos^2\phi-\sin^2\phi\right)\,\cos\theta
                 +\left(2\sin\phi\cos\phi\right)\,\sin\theta        \\
              &=\left((\frac{1}{\sqrt{1+k^2}})^2
                      -(\frac{k}{\sqrt{1+k^2}})^2\right)
                   \,\cos\theta
                 +\left(2\frac{k}{\sqrt{1+k^2}}\frac{1}{\sqrt{1+k^2}}\right)
                   \,\sin\theta        \\
              &=\left(\frac{1-k^2}{1+k^2}\right)\,\cos\theta
                 +\left(\frac{2k}{1+k^2}\right)\,\sin\theta
          \end{align*}
          因此像向量的第一个分量如下。
          \begin{equation*}
            r\cdot\cos(2\phi-\theta)=\frac{1-k^2}{1+k^2}\cdot x
                                     +\frac{2k}{1+k^2}\cdot y
          \end{equation*}
          类似的计算表明，像向量的第二个分量如下。
          \begin{equation*}
            r\cdot\sin(2\phi-\theta)=\frac{2k}{1+k^2}\cdot x
                                     -\frac{1-k^2}{1+k^2}\cdot y
          \end{equation*}
          有了 $f_\ell$ 作用的这一代数表达式
          \begin{equation*}
            \colvec{x \\ y}
             \mapsunder{f_\ell}
            \colvec{(1-k^2/1+k^2)\cdot x
                                     +(2k/1+k^2)\cdot y \\
                   (2k/1+k^2)\cdot x
                                     -(1-k^2/1+k^2)\cdot y}
          \end{equation*}
          就很容易验证它保持结构。
      \end{exparts}
    
\end{ans}
\begin{ans}{三.I.1.34}
       首先，映射 $p(x)\mapsto p(x+k)$ 不符合要求，
       因为它也是 $p(x)\mapsto p(x-k)$ 这种形式的映射。
       下面给出一个正确答案（还有许多其他正确答案）：
       \( a_0+a_1x+a_2x^2 \mapsto a_2+a_0x+a_1x^2 \)。
       验证它是同构很直接。
    
\end{ans}
\begin{ans}{三.I.1.35}
      \begin{exparts}
        \partsitem 先证必要性。设 \( \map{f}{\Re^1}{\Re^1} \) 是同构。
          考虑基 \( \sequence{1}\subseteq\Re^1 \)，
          把 \( f(1) \) 记为 \( k \)。
          则对任意 \( x \)，都有
          \( f(x)=f(x\cdot 1)=x\cdot f(1)=xk \)，
          所以 $f$ 的作用就是乘以 $k$。
          最后，注意 \( k\neq 0 \)，否则 $f$ 不是单射。

          再证充分性，只需验证当 $k\neq 0$ 时，这样的映射是同构。
          要验证它是单射，假设 \( f(x_1)=f(x_2) \)，
          则 \( kx_1=kx_2 \)，两边除以非零因子 \( k \)，得到 $x_1=x_2$。
          要验证它是满射，注意任意 \( y\in\Re^1 \)
          都是 \( x=y/k \) 的像（这里同样用到 $k\neq 0$）。
          最后，下面的计算验证了这个映射保持定义域中两个元素的线性组合。
          \begin{equation*}
            f(c_1x_1+c_2x_2)
            =k(c_1x_1+c_2x_2)
            =c_1kx_1+c_2kx_2
            =c_1f(x_1)+c_2f(x_2)
          \end{equation*}
        \partsitem 由前一题可知，$f$ 的作用是 \( x\mapsto (7/3)x \)。
          因此 \( f(-2)=-14/3 \)。
        \partsitem 先证必要性，假设 \( \map{f}{\Re^2}{\Re^2} \) 是自同构。
          考虑 \( \Re^2 \) 的标准基 \( \stdbasis_2 \)。
          设
          \begin{equation*}
            f(\vec{e}_1)=\colvec{a \\ c}
            \quad\text{且}\quad
            f(\vec{e}_2)=\colvec{b \\ d}.
          \end{equation*}
          则 $f$ 对任意向量的作用，都由它对这两个基向量的作用决定。
          \begin{equation*}
            f(\colvec{x \\ y})
            =f(x\cdot\vec{e}_1+y\cdot\vec{e}_2)
            =x\cdot f(\vec{e}_1)+y\cdot f(\vec{e}_2)
            =x\cdot\colvec{a \\ c}+y\cdot\colvec{b \\ d}
            =\colvec{ax+by \\ cx+dy}
          \end{equation*}
          最后，注意若 \( ad-bc=0 \)，即
          \( f(\vec{e}_2) \) 是 \( f(\vec{e}_1) \) 的倍数，
          则 \( f \) 不是单射。

          再证充分性，需要在 $ad-bc\neq 0$ 的条件下验证这个映射是同构。
          保持结构很容易验证；这里我们来证明 \( f \) 是一一对应。
          要证明这个映射是单射，假设
          \begin{equation*}
            f(\colvec{x_1 \\ y_1})
            =f(\colvec{x_2 \\ y_2})
            \quad\text{因此}\quad
            \colvec{ax_1+by_1 \\ cx_1+dy_1}
            =\colvec{ax_2+by_2 \\ cx_2+dy_2}
          \end{equation*}
          由于 \( ad-bc\neq 0 \)，得到的方程组
          \begin{equation*}
            \begin{linsys}{2}
              a(x_1-x_2)  &+  &b(y_1-y_2)  &=  &0  \\
              c(x_1-x_2)  &+  &d(y_1-y_2)  &=  &0
            \end{linsys}
          \end{equation*}
          有唯一解，即平凡解
          $x_1-x_2=0$ 且 $y_1-y_2=0$
          （这可由提示推出）。

          证明这个映射是满射的思路与此密切相关\Dash 方程组
          \begin{equation*}
            \begin{linsys}{2}
              ax_1  &+  &by_1  &=  &x  \\
              cx_1  &+  &dy_1  &=  &y
            \end{linsys}
          \end{equation*}
          对任意 \( x \) 和 \( y \) 都有解，当且仅当集合
          \begin{equation*}
            \set{\colvec{a \\ c},
                 \colvec{b \\ d} }
          \end{equation*}
          张成 \( \Re^2 \)，也即当且仅当这个集合是一组基
          （因为它是 $\Re^2$ 的一个二元子集），
          也即当且仅当 \( ad-bc\neq 0 \)。
        \partsitem
          \begin{equation*}
            f(\colvec[r]{0 \\ -1})=f(\colvec[r]{1 \\ 3}-\colvec[r]{1 \\ 4})
            =f(\colvec[r]{1 \\ 3})-f(\colvec[r]{1 \\ 4})
            =\colvec[r]{2 \\ -1}-\colvec[r]{0 \\ 1}
            =\colvec[r]{2 \\ -2}
          \end{equation*}
      \end{exparts}
    
\end{ans}
\begin{ans}{三.I.1.36}
       答案有很多，其中两个是线性无关性和子空间。

       先证明：若集合 $\set{\vec{v}_1,\dots,\vec{v}_n}$ 线性无关，
       则它的像集 $\set{f(\vec{v}_1),\dots,f(\vec{v}_n)}$ 也线性无关。
       考虑像集中的元素之间的一个线性关系。
       \begin{equation*}
         0=c_1f(\vec{v}_1)+\dots+c_nf(\vec{v_n})
          =f(c_1\vec{v}_1)+\dots+f(c_n\vec{v_n})
          =f(c_1\vec{v}_1+\dots+c_n\vec{v_n})
       \end{equation*}
       因为这个映射是同构，所以它是单射。
       因此 $f$ 只把定义域中的一个向量映到值域中的零向量，
       也就是说，$c_1\vec{v}_1+\dots+c_n\vec{v}_n$ 等于零向量
       （当然，是定义域中的零向量）。
       但若 $\set{\vec{v}_1,\dots,\vec{v}_n}$ 线性无关，
       则所有 $c$ 都为零，
       因而 $\set{f(\vec{v}_1),\dots,f(\vec{v}_n)}$ 也线性无关。
       （\textit{注。}
       这一论证中有个细节需要说明。
       集合中的重复元素只计一次；严格地说，
       $\set{\vec{v},\vec{v}}$ 是一个单元素集合，
       因为列出的两个元素其实是同一个。
       不过，请注意上面论证中下标~$n$ 的使用。
       从定义域中的集合 $\set{\vec{v}_1,\dots,\vec{v}_n}$
       转到像集 $\set{f(\vec{v}_1),\dots,f(\vec{v}_n)}$ 时，
       元素个数不会因重复而减少，
       因为同构 $f$ 是单射，像集中没有重复元素。）

       设 $\map{f}{V}{W}$ 是同构，$U$ 是定义域 $V$ 的子空间。
       要证明像向量的集合
       $f(U)=\set{\vec{w}\in W
                 \suchthat\text{存在 $\vec{u}\in U$ 使 $\vec{w}=f(\vec{u})$}}$
       是 $W$ 的子空间，只需证明它对其中两个元素的线性组合封闭
       （它包含零向量的像，所以非空）。
       我们有
       \begin{equation*}
         c_1\cdot f(\vec{u}_1)+c_2\cdot f(\vec{u}_2)
        =f(c_1\vec{u}_1)+f(c_2\vec{u}_2)
        =f(c_1\vec{u}_1+c_2\vec{u}_2)
       \end{equation*}
       而 $c_1\vec{u}_1+c_2\vec{u}_2$ 属于 $U$，
       因为子空间对线性组合封闭。
       因此，$f(\vec{u}_1)$ 与 $f(\vec{u}_2)$ 的这个线性组合属于 $f(U)$。
      
\end{ans}
\begin{ans}{三.I.1.37}
      \begin{exparts}
        \partsitem 对应关系
          \begin{equation*}
            \vec{p}=c_1\vec{\beta}_1+c_2\vec{\beta}_2+c_3\vec{\beta}_3
            \mapsunder{\rep{\cdot}{B}}
            \colvec{c_1 \\ c_2 \\ c_3}
          \end{equation*}
          是函数，只需定义域中的每个元素 $\vec{p}$ 都至少对应陪域中的一个元素，
          并且定义域中的每个元素 $\vec{p}$ 都至多对应陪域中的一个元素。
          第一个条件成立，因为基 $B$ 张成定义域\Dash 每个 $\vec{p}$
          至少可以写成一种由 $\vec{\beta}$ 组成的线性组合。
          第二个条件成立，因为基 $B$ 线性无关\Dash 定义域中的每个元素 $\vec{p}$
          至多可以写成一种由 $\vec{\beta}$ 组成的线性组合。
        \partsitem 为证明它是单射，若 $\rep{\vec{p}}{B}=\rep{\vec{q}}{B}$，即若
          \( \rep{p_1\vec{\beta}_1+p_2\vec{\beta}_2+p_3\vec{\beta}_3}{B}
             =\rep{q_1\vec{\beta}_1+q_2\vec{\beta}_2+q_3\vec{\beta}_3}{B} \)，
          则
          \begin{equation*}
            \colvec{p_1 \\ p_2 \\ p_3}
            =\colvec{q_1 \\ q_2 \\ q_3}
          \end{equation*}
          因而 \( p_1=q_1 \)、\( p_2=q_2 \) 且 \( p_3=q_3 \)，
          由此得到 $\vec{p}=\vec{q}$。
          因此这个映射是单射。

          至于满射，只需注意
          \begin{equation*}
            \colvec{a \\ b \\ c}
          \end{equation*}
          等于 $\rep{a\vec{\beta}_1+b\vec{\beta}_2+c\vec{\beta}_3}{B}$，
          所以陪域 $\Re^3$ 中的每个元素都是定义域 $\polyspace_2$ 中某个元素的像。
        \partsitem 这个映射保持加法和标量乘法，
           因为它保持定义域中两个元素的线性组合
           （这里使用\nearbylemma{le:PresStructIffPresCombos} 的第~(2) 项）：
           若 $\vec{p}=p_1\vec{\beta}_1+p_2\vec{\beta}_2+p_3\vec{\beta}_3$
           且 $\vec{q}=q_1\vec{\beta}_1+q_2\vec{\beta}_2+q_3\vec{\beta}_3$，
           则有
           \begin{align*}
             \rep{c\cdot\vec{p}+d\cdot\vec{q}}{B}
             &=\rep{\,(cp_1+dq_1)\vec{\beta}_1+(cp_2+dq_2)\vec{\beta}_2
                  +(cp_3+dq_3)\vec{\beta}_3\,}{B}  \\
             &=\colvec{cp_1+dq_1 \\ cp_2+dq_2 \\ cp_3+dq_3}  \\
             &=c\cdot\colvec{p_1 \\ p_2 \\ p_3}
               +d\cdot\colvec{q_1 \\ q_2 \\ q_3}  \\
             &=\rep{\vec{p}}{B}+\rep{\vec{q}}{B}
           \end{align*}
        \partsitem 取 $\polyspace_2$ 的任意一组基 \( B \)，
          使其前两个元素为 \( x+x^2 \) 与 \( 1-x \)，
          例如 \( B=\sequence{x+x^2,1-x,1} \)。
      \end{exparts}
    
\end{ans}
\begin{ans}{三.I.1.38}
      见下一小节。
    
\end{ans}
\begin{ans}{三.I.1.39}
      \begin{exparts}
        \partsitem 向量空间定义中的大多数条件都很容易验证。
          这里简要验证定义的第~(1) 部分。

          对于 $U\times W$ 的封闭性，注意 \( U \) 与 \( W \) 都封闭，
          所以 $\vec{u}_1+\vec{u}_2\in U$ 且 $\vec{w}_1+\vec{w}_2\in W$，
          从而 \( (\vec{u}_1+\vec{u}_2,\vec{w}_1+\vec{w}_2)\in U\times W \)。
          \( U\times W \) 中加法的交换律可由
          \( U \) 与 \( W \) 中加法的交换律推出。
          \begin{equation*}
             (\vec{u}_1,\vec{w}_1)+(\vec{u}_2,\vec{w}_2)=
                 (\vec{u}_1+\vec{u}_2,\vec{w}_1+\vec{w}_2)
                 =(\vec{u}_2+\vec{u}_1,\vec{w}_2+\vec{w}_1)
             =(\vec{u}_2,\vec{w}_2)+(\vec{u}_1,\vec{w}_1)
          \end{equation*}
          加法结合律的验证与此类似。
          零元素为 \( (\zero_U,\zero_W)\in U\times W \)，
          而 \( (\vec{u},\vec{w}) \) 的加法逆元为 \( (-\vec{u},-\vec{w}) \)。

          向量空间定义的第二部分也很容易验证。
        \partsitem 下面是一组基
          \begin{equation*}
            \sequence{\, (1,\colvec[r]{0 \\ 0}), (x,\colvec[r]{0 \\ 0}),
              (x^2,\colvec[r]{0 \\ 0}), (1,\colvec[r]{1 \\ 0}),
              (1,\colvec[r]{0 \\ 1}) \,}
          \end{equation*}
          因为 $\polyspace_2\times \Re^2$ 中的任意元素
          关于这个集合都有且只有一种表示；下面给出一个例子。
          \begin{equation*}
            (3+2x+x^2,\colvec[r]{5 \\ 4})
             =
              3\cdot(1,\colvec[r]{0 \\ 0})
              +2\cdot(x,\colvec[r]{0 \\ 0})
              +(x^2,\colvec[r]{0 \\ 0})
              +5\cdot(1,\colvec[r]{1 \\ 0})
              +4\cdot(1,\colvec[r]{0 \\ 1})
          \end{equation*}
          这个空间的维数为五。
        \partsitem 我们有 \( \dim(U\times W)=\dim(U)+\dim(W) \)，
          因为下面是一组基。
          \begin{equation*}
            \sequence{ (\vec{\mu}_1,\zero_W), \dots,
              (\vec{\mu}_{\dim(U)},\zero_W), (\zero_U,\vec{\omega}_1),
              \ldots,(\zero_U,\vec{\omega}_{\dim(W)}) }
          \end{equation*}
        \partsitem 我们知道，若 \( V=U\directsum W \)，则每个
          \( \vec{v}\in V \) 都有且只有一种写法 \( \vec{v}=\vec{u}+\vec{w} \)。
          这正是证明所给函数为同构所需要的性质。

          首先，要证明 \( f \) 是单射，可以证明：若
          \( f\left((\vec{u}_1,\vec{w}_1)\right)
                 =\left((\vec{u}_2,\vec{w}_2)\right) \)，即若
          \( \vec{u}_1+\vec{w}_1=\vec{u}_2+\vec{w}_2 \)，则
          \( \vec{u}_1=\vec{u}_2 \) 且 \( \vec{w}_1=\vec{w}_2 \)。
          而“每个 $\vec{v}$ 都只有一种这样的和式表示”恰好给出这一结论。
          类似地，“每个 $\vec{v}$ 都至少有一种这样的和式表示”
          就证明了 \( f \) 是满射。

          这个映射也保持线性组合
          \begin{align*}
            f(\,c_1\cdot(\vec{u}_1,\vec{w}_1)+c_2\cdot (\vec{u}_2,\vec{w}_2)\,)
            &=f(\,(c_1\vec{u}_1+c_2\vec{u}_2,c_1\vec{w}_1+c_2\vec{w}_2)\,) \\
            &=c_1\vec{u}_1+c_2\vec{u}_2+c_1\vec{w}_1+c_2\vec{w}_2    \\
            &=c_1\vec{u}_1+c_1\vec{w}_1+c_2\vec{u}_2+c_2\vec{w}_2    \\
            &=c_1\cdot f(\,(\vec{u}_1,\vec{w}_1)\,)
               +c_2\cdot f(\,(\vec{u}_2,\vec{w}_2)\,)
          \end{align*}
          因此它是同构。
      \end{exparts}
    
\end{ans}
\protect \subsection {三.I.2: 维数刻画同构}
\begin{ans}{三.I.2.10}
       每一对空间同构，当且仅当它们的维数相同。
       同构时，可以给出一个具体的映射，但严格说来没有必要。
       \begin{exparts}
         \partsitem 不同构，维数不同。
         \partsitem 不同构，维数不同。
         \partsitem 同构，维数相同。
           下面给出一个同构。
           \begin{equation*}
             \begin{mat}
               a  &b  &c  \\
               d  &e  &f
             \end{mat}
             \mapsto
             \colvec{a \\ \vdots \\ f}
           \end{equation*}
         \partsitem 同构，维数相同。
           下面是一个同构。
           \begin{equation*}
             a+bx+\cdots+fx^5
             \mapsto
             \begin{mat}
               a  &b  &c  \\
               d  &e  &f
             \end{mat}
           \end{equation*}
         \partsitem 同构，两者的维数都是 \( 2k \)。
      \end{exparts}
    
\end{ans}
\begin{ans}{三.I.2.11}
      只需求出每个空间的维数，例如可以通过找一组基来确定维数，
      然后维数相同的空间就同构。
      各空间的维数如下。
      \begin{exparts*}
        \partsitem $3$
        \partsitem $4$
        \partsitem $4$
        \partsitem $4$
        \partsitem $3$
      \end{exparts*}
    
\end{ans}
\begin{ans}{三.I.2.12}
      \begin{exparts*}
        \partsitem \( \rep{3-2x}{B}=\colvec[r]{5 \\ -2} \)
        \partsitem \( \colvec[r]{0 \\ 2} \)
        \partsitem \( \colvec[r]{-1 \\ 1} \)
      \end{exparts*}
     
\end{ans}
\begin{ans}{三.I.2.13}
      对于每一题，最简单的方法都是找一组基来确定空间的维数。
      对于下面给出的各组基，我们省略验证它确为基的过程。
      \begin{exparts}
        \partsitem 它同构于 $\Re^5$。
          $\polyspace_4$ 的一组基是 $\set{x^4,x^3,x^2,x,1}$，
          因此该空间的维数为~$5$。
        \partsitem 它同构于 $\Re^2$，因为
          空间 $\polyspace_1=\set{a+bx\suchthat a,b\in\Re}$ 的一组基是 $\set{1,x}$。
        \partsitem 它同构于 $\Re^6$。
          下面六个矩阵构成一组基。
          \begin{equation*}
            \begin{mat}
              1 &0 &0 \\
              0 &0 &0
            \end{mat},\;
            \begin{mat}
              0 &1 &0 \\
              0 &0 &0
            \end{mat},\;
            \cdots\;
            \begin{mat}
              0 &0 &0 \\
              0 &0 &1
            \end{mat}
          \end{equation*}
        \partsitem 它是平面，因此同构于~$\Re^2$。
          更具体地，对这个平面作参数化，可得如下向量描述
          \begin{equation*}
            \set{\colvec{x \\ y \\ z}=\colvec{1/2 \\ 1 \\ 0}y
                                      +\colvec{-1/2 \\ 0 \\ 1}z
                   \suchthat  y,z\in\Re}
          \end{equation*}
          因此，其中的两个向量构成一组基。
        \partsitem 它同构于 $\Re^3$。
          一组基是线性组合的集合 $\set{x,y,z}$，
          即 $\set{x+0y+0z, 0x+y+0z,0x+0y+z}$。
      \end{exparts}
    
\end{ans}
\begin{ans}{三.I.2.14}
      它们的维数不同。
     
\end{ans}
\begin{ans}{三.I.2.15}
      是的，两者都是 \( mn \) 维的。
    
\end{ans}
\begin{ans}{三.I.2.16}
      是的，任意两个（非退化的）平面都是二维向量空间。
    
\end{ans}
\begin{ans}{三.I.2.17}
        答案有很多，其中一个是各个 \( \polyspace_k \) 组成的集合
        （把 \( \polyspace_{-1} \) 视为平凡向量空间）。
     
\end{ans}
\begin{ans}{三.I.2.18}
      错误（\( n=0 \) 时除外）。
      例如，若 \( \map{f}{V}{\Re^n} \) 是同构，
      则将它的输出乘以一个不等于一的非零标量，就会得到另一个不同的同构。
      （平凡空间之间的同构是唯一的；唯一可能的映射是 $\zero_V\mapsto 0_W$。）
    
\end{ans}
\begin{ans}{三.I.2.19}
      不能。
      真子空间的维数严格小于它所在空间的维数；
      若 $U$ 是 $V$ 的真子空间，则 $U$ 的任意线性无关子集
      都必须少于 $\dim(V)$ 个元素，否则这个集合就是 $V$ 的一组基，
      $U$ 就不是真子空间了。
    
\end{ans}
\begin{ans}{三.I.2.20}
      若 \( B=\sequence{\vec{\beta}_1,\ldots,\vec{\beta}_n} \)，
      则逆映射如下。
      \begin{equation*}
        \colvec{c_1 \\ \vdots \\ c_n}
        \mapsto c_1\vec{\beta}_1+\cdots+c_n\vec{\beta}_n
      \end{equation*}
    
\end{ans}
\begin{ans}{三.I.2.21}
      这三个空间的维数都等于矩阵的秩。
    
\end{ans}
\begin{ans}{三.I.2.22}
        我们必须证明：若 \( \vec{a}=\vec{b} \)，则
        \( f(\vec{a})=f(\vec{b}) \)。
        假设
        $a_1\vec{\beta}_1+\dots+a_n\vec{\beta}_n
           =b_1\vec{\beta}_1+\dots+b_n\vec{\beta}_n$。
        向量空间（这里是定义域空间）中的每个向量
        都有唯一的基向量线性组合表示，因此可以推出 \( a_1=b_1 \)，\ldots，
        \(a_n=b_n \)。
        于是，
        \begin{equation*}
          f(\vec{a})
          =\colvec{a_1 \\ \vdots \\ a_n}=\colvec{b_1 \\ \vdots \\ b_n}
          =f(\vec{b})
        \end{equation*}
        所以这个函数是良定义的。
     
\end{ans}
\begin{ans}{三.I.2.23}
      成立，因为零维空间就是平凡空间。
     
\end{ans}
\begin{ans}{三.I.2.24}
      \begin{exparts}
        \partsitem 不是，这个集合中没有奇数维空间。
        \partsitem 是，因为 $\polyspace_{k}\isomorphicto\Re^{k+1}$。
        \partsitem 不是，例如
          \( \matspace_{\nbym{2}{3}}\isomorphicto\matspace_{\nbym{3}{2}} \)。
      \end{exparts}
    
\end{ans}
\begin{ans}{三.I.2.25}
       一个方向很容易：~若两个空间通过 \( f \) 同构，
       则对任意基 \( B\subseteq V \)，
       集合 \( D=f(B) \) 也是基
       （\nearbylemma{lem:IsoImpliesSameDim} 已经证明了这一点）。
       对应向量具有相同的坐标，这很容易验证：
       \( f(c_1\vec{\beta}_1+\dots+c_n\vec{\beta}_n)
          =c_1f(\vec{\beta}_1)+\dots+c_nf(\vec{\beta}_n)
          =c_1\vec{\delta}_1+\dots+c_n\vec{\delta}_n   \)。

       反过来，假设存在这样的基，使对应向量关于这两组基具有相同的坐标。
       因为 \( f \) 是一一对应，要证明它是同构，只需证明它保持结构。
       由于 \( \rep{\vec{v}\,}{B}=\rep{f(\vec{v}\,)}{D} \)，
       映射 \( f \) 保持结构，当且仅当表示运算保持加法：
       \( \rep{\vec{v}_1+\vec{v}_2}{B}=\rep{\vec{v}_1}{B}+\rep{\vec{v}_2}{B} \)
       和标量乘法：
       \( \rep{r\cdot\vec{v}\,}{B}=r\cdot\rep{\vec{v}\,}{B} \)。
       加法的计算如下：
       \( (c_1+d_1)\vec{\beta}_1+\dots+(c_n+d_n)\vec{\beta}_n
          =c_1\vec{\beta}_1+\dots+c_n\vec{\beta}_n
          +d_1\vec{\beta}_1+\dots+d_n\vec{\beta}_n \)，
       标量乘法的计算与此类似。
     
\end{ans}
\begin{ans}{三.I.2.26}
       \begin{exparts}
        \partsitem 把定义从 \( \Re^4 \) 拉回到 \( \polyspace_3 \)，
          可得 \( a_0+a_1x+a_2x^2+a_3x^3 \)
          与 \( b_0+b_1x+b_2x^2+b_3x^3 \) 正交，当且仅当
          \( a_0b_0+a_1b_1+a_2b_2+a_3b_3=0 \)。
        \partsitem 一个自然的定义如下。
          \begin{equation*}
            D(\colvec{a_0 \\ a_1 \\ a_2 \\ a_3})=
              \colvec{a_1 \\ 2a_2 \\ 3a_3 \\ 0}
          \end{equation*}
      \end{exparts}
    
\end{ans}
\begin{ans}{三.I.2.27}
      是的。

      首先，\( \hat{f} \) 是良定义的，
      因为 \( V \) 中的每个元素都有且只有一种表示为 \( B \) 中元素的线性组合的方式。

      其次，需要证明 $\hat{f}$ 既是单射又是满射。
      它是单射，因为 $D$ 是基，\( W \) 中的每个元素
      都只能以一种方式表示为 \( D  \) 中元素的线性组合。
      \( \hat{f} \) 是满射，因为 \( W \) 中的每个元素
      都至少有一种表示为 \( D \) 中元素的线性组合的方式。

      最后，保持结构很容易验证。
      例如，保持加法的计算如下。
      \begin{multline*}
        \hat{f}(\,(c_1\vec{\beta}_1+\dots+c_n\vec{\beta}_n)+
                  (d_1\vec{\beta}_1+\dots+d_n\vec{\beta}_n)\,)    \\
       \begin{aligned}
        &=\hat{f}(\,(c_1+d_1)\vec{\beta}_1+\dots+(c_n+d_n)\vec{\beta}_n\,) \\
        &=(c_1+d_1)f(\vec{\beta}_1)+\dots+(c_n+d_n)f(\vec{\beta}_n) \\
        &=c_1f(\vec{\beta}_1)+\dots+c_nf(\vec{\beta}_n)
          +d_1f(\vec{\beta}_1)+\dots+d_nf(\vec{\beta}_n) \\
        &=\hat{f}(c_1\vec{\beta}_1+\dots+c_n\vec{\beta}_n)+
          +\hat{f}(d_1\vec{\beta}_1+\dots+d_n\vec{\beta}_n)
      \end{aligned}
      \end{multline*}
      （第二个等号使用了 $\hat{f}$ 的定义。）
      保持标量乘法的验证与此类似。
    
\end{ans}
\begin{ans}{三.I.2.28}
      因为 \( V_1\intersection V_2=\set{\zero_V} \)，且 \( f \) 是单射，
      所以 \( f(V_1)\intersection f(V_2)=\set{\zero_U} \)。
      最后计算维数：
      \( \dim(U)=\dim(V)=\dim(V_1)+\dim(V_2)=\dim(f(V_1))+\dim(f(V_2)) \)，
      即得所需结论。
    
\end{ans}
\begin{ans}{三.I.2.29}
      有理数有多种表示，例如 \( 1/2=3/6 \)，
      不同表示的分子可能不同。
    
\end{ans}
\protect \section {第 II 节：同态}
\protect \subsection {三.II.1: 定义}
\begin{ans}{三.II.1.18}
      \begin{exparts}
        \partsitem 是。
          验证很直接。
          \begin{align*}
            h( c_1\cdot\colvec{x_1 \\ y_1 \\ z_1}
               +c_2\cdot\colvec{x_2 \\ y_2 \\ z_2} )
            &=h( \colvec{c_1x_1+c_2x_2 \\ c_1y_1+c_2y_2 \\ c_1z_1+c_2z_2} ) \\
            &=\colvec{c_1x_1+c_2x_2 \\
                     c_1x_1+c_2x_2+c_1y_1+c_2y_2+c_1z_1+c_2z_2}          \\
            &=c_1\cdot\colvec{x_1 \\ x_1+y_1+z_1}
               +c_2\cdot\colvec{x_2 \\ c_2+y_2+z_2}                       \\
            &=c_1\cdot h(\colvec{x_1 \\ y_1 \\ z_1})
               +c_2\cdot h(\colvec{x_2 \\ y_2 \\ z_2})
          \end{align*}
        \partsitem 是。
          很容易验证。
          \begin{align*}
            h(c_1\cdot\colvec{x_1 \\ y_1 \\ z_1}
              +c_2\cdot\colvec{x_2 \\ y_2 \\ z_2})
            &=h( \colvec{c_1x_1+c_2x_2 \\ c_1y_1+c_2y_2 \\ c_1z_1+c_2z_2} ) \\
            &=\colvec[r]{0 \\ 0}                                      \\
            &=c_1\cdot h(\colvec{x_1 \\ y_1 \\ z_1})
               +c_2\cdot h(\colvec{x_2 \\ y_2 \\ z_2})
          \end{align*}
        \partsitem 不是。
          下面给出一个不保持加法的例子。
          \begin{equation*}
            h(\colvec[r]{0 \\ 0 \\ 0}+\colvec[r]{0 \\ 0 \\ 0})
            =\colvec[r]{1 \\ 1}
            \neq h(\colvec[r]{0 \\ 0 \\ 0})+h(\colvec[r]{0 \\ 0 \\ 0})
          \end{equation*}
        \partsitem 是。
           验证很直接。
          \begin{align*}
            h( c_1\cdot\colvec{x_1 \\ y_1 \\ z_1}
               +c_2\cdot\colvec{x_2 \\ y_2 \\ z_2} )
            &=h( \colvec{c_1x_1+c_2x_2 \\ c_1y_1+c_2y_2 \\ c_1z_1+c_2z_2} ) \\
            &=\colvec{2(c_1x_1+c_2x_2)+(c_1y_1+c_2y_2) \\
                      3(c_1y_1+c_2y_2)-4(c_1z_1+c_2z_2)}          \\
            &=c_1\cdot\colvec{2x_1+y_1 \\ 3y_1-4z_1}
               +c_2\cdot\colvec{2x_2+y_2 \\ 3y_2-4z_2}            \\
            &=c_1\cdot h(\colvec{x_1 \\ y_1 \\ z_1})
               +c_2\cdot h(\colvec{x_2 \\ y_2 \\ z_2})
          \end{align*}
      \end{exparts}
    
\end{ans}
\begin{ans}{三.II.1.19}
      对每一题，要么验证映射保持线性组合，
      要么给出一个它不保持的线性组合。
      \begin{exparts*}
         \partsitem 是。
           很容易验证它保持线性组合。
           \begin{align*}
             h(r_1\cdot\begin{mat}
                 a_1  &b_1  \\
                 c_1  &d_1
               \end{mat}
              +r_2\cdot\begin{mat}
                 a_2  &b_2  \\
                 c_2  &d_2
               \end{mat})
              &=h(\begin{mat}
                 r_1a_1+r_2a_2  &r_1b_1+r_2b_2  \\
                 r_1c_1+r_2c_2  &r_1d_1+r_2d_2
               \end{mat})                    \\
              &=(r_1a_1+r_2a_2)+(r_1d_1+r_2d_2)    \\
              &=r_1(a_1+d_1)+r_2(a_2+d_2)          \\
              &=r_1\cdot h(\begin{mat}
                             a_1  &b_1  \\
                             c_1  &d_1
                            \end{mat})
                +r_2\cdot h(\begin{mat}
                             a_2  &b_2  \\
                             c_2  &d_2
                            \end{mat})
           \end{align*}
        \partsitem 不是。
          例如，它不保持乘以标量 $2$ 的运算。
          \begin{equation*}
            h(2\cdot\begin{mat}[r]
                1  &0  \\
                0  &1
              \end{mat})
            =h(\begin{mat}[r]
                2  &0  \\
                0  &2
              \end{mat})
            =4
            \quad\text{而}\quad
            2\cdot h(\begin{mat}[r]
                1  &0  \\
                0  &1
              \end{mat})
            =2\cdot 1=2
          \end{equation*}
        \partsitem 是。
           下面验证它保持定义域中两个元素的线性组合。
           \begin{multline*}
             h(r_1\cdot\begin{mat} a_1 &b_1 \\ c_1 &d_1 \end{mat}
               +r_2\cdot\begin{mat} a_2 &b_2 \\ c_2 &d_2 \end{mat})     \\
             \begin{aligned}
             &=h(\begin{mat}
                   r_1a_1+r_2a_2 &r_1b_1+r_2b_2   \\
                   r_1c_1+r_2c_2 &r_1d_1+r_2d_2
                 \end{mat})                                 \\
             &=2(r_1a_1+r_2a_2)+3(r_1b_1+r_2b_2)
                    +(r_1c_1+r_2c_2)-(r_1d_1+r_2d_2)              \\
             &=r_1(2a_1+3b_1+c_1-d_1)
               +r_2(2a_2+3b_2+c_2-d_2)              \\
             &=r_1\cdot h(\begin{mat} a_1 &b_1 \\ c_1 &d_1 \end{mat}
               +r_2\cdot h(\begin{mat} a_2 &b_2 \\ c_2 &d_2 \end{mat})
             \end{aligned}
           \end{multline*}
        \partsitem 不是。
          下面给出一个不保持线性组合的例子：
          按一种次序计算得到
          \begin{equation*}
            h(\begin{mat}[r]
              1  &0  \\
              0  &0
            \end{mat}
            +\begin{mat}[r]
              1  &0  \\
              0  &0
            \end{mat})
           =h(\begin{mat}[r]
              2  &0  \\
              0  &0
            \end{mat})
           =4
           \end{equation*}
          而按另一种次序计算得到
           \begin{equation*}
            h(\begin{mat}[r]
              1  &0  \\
              0  &0
            \end{mat})
            +h(\begin{mat}[r]
              1  &0  \\
              0  &0
            \end{mat})
            =1+1
            =2
          \end{equation*}
      \end{exparts*}
    
\end{ans}
\begin{ans}{三.II.1.20}
      很容易验证它们都是同态。
      下面验证求导映射。
      \begin{multline*}
        \frac{d}{dx}(r\cdot (a_0+a_1x+a_2x^2+a_3x^3)
                    +s\cdot (b_0+b_1x+b_2x^2+b_3x^3))                 \\
        \begin{aligned}
           &=\frac{d}{dx}((ra_0+sb_0)+(ra_1+sb_1)x+(ra_2+sb_2)x^2
                                                    +(ra_3+sb_3)x^3)  \\
           &=(ra_1+sb_1)+2(ra_2+sb_2)x+3(ra_3+sb_3)x^2         \\
           &=r\cdot (a_1+2a_2x+3a_3x^2)+s\cdot (b_1+2b_2x+3b_3x^2)         \\
           &=r\cdot \frac{d}{dx}(a_0+a_1x+a_2x^2+a_3x^3)
                    +s\cdot \frac{d}{dx} (b_0+b_1x+b_2x^2+b_3x^3)
        \end{aligned}
      \end{multline*}
      （另一种证明只需指出，这是微积分中熟悉的求导性质。）

      这两个映射不互为逆映射，因为下面的复合
      对这个定义域元素的作用不同于恒等映射。
      \begin{equation*}
        1\in\polyspace_3\;\mapsunder{d/dx}\;
        0\in\polyspace_2\;\mapsunder{\int}\;
        0\in\polyspace_3
      \end{equation*}
     
\end{ans}
\begin{ans}{三.II.1.21}
      这些投影都是同态。
      到 $xz$ 平面和 $yz$ 平面的投影分别为
      \begin{equation*}
         \colvec{x \\ y \\ z}\mapsto\colvec{x \\ 0 \\ z}
           \qquad
         \colvec{x \\ y \\ z}\mapsto\colvec{0 \\ y \\ z}
      \end{equation*}
      到 $x$ 轴、$y$ 轴和 $z$ 轴的投影分别为
      \begin{equation*}
         \colvec{x \\ y \\ z}\mapsto\colvec{x \\ 0 \\ 0}
           \qquad
         \colvec{x \\ y \\ z}\mapsto\colvec{0 \\ y \\ 0}
           \qquad
         \colvec{x \\ y \\ z}\mapsto\colvec{0 \\ 0 \\ z}
      \end{equation*}
      到原点的投影为
      \begin{equation*}
         \colvec{x \\ y \\ z}\mapsto\colvec[r]{0 \\ 0 \\ 0}
      \end{equation*}
      很容易验证它们都是同态。
      （当然，最后一个是 $\Re^3$ 上的零变换。）
     
\end{ans}
\begin{ans}{三.II.1.22}
      \begin{exparts}
        \partsitem
          下式验证这个映射保持线性组合。
          根据\nearbylemma{le:HomoPreserveLinCombo}，这足以说明它是同态。
          \begin{multline*}
            h(\,d_1(a_1x^2+b_1x+c_1)+d_2(a_2x^2+b_2x+c_2)\,)      \\
             \begin{aligned}
             &=h((d_1a_1+d_2a_2)x^2+(d_1b_1+d_2b_2)x+(d_1c_1+d_2c_2))  \\
             &=\colvec{(d_1a_1+d_2a_2)+(d_1b_1+d_2b_2) \\
             (d_1a_1+d_2a_2)+(d_1c_1+d_2c_2)} \\
             &=\colvec{d_1a_1+d_1b_1 \\  d_1a_1+d_1c_1}+
                 \colvec{d_2a_2+d_2b_2 \\  d_2a_2+d_2c_2}  \\
             &=d_1\colvec{a_1+b_1 \\  a_1+c_1}+
                 d_2\colvec{a_2+b_2 \\  a_2+c_2}  \\
            &=d_1\cdot h(a_1x^2+b_1x+c_1)+d_2\cdot h(a_2x^2+b_2x+c_2)
            \end{aligned}
          \end{multline*}
        \partsitem
          它保持线性组合。
          \begin{align*}
            f(\,a_1\colvec{x_1 \\ y_1}+a_2\colvec{x_2 \\ y_2}\,)
              &=f(\,\colvec{a_1x_1+a_2x_2 \\ a_1y_1+a_2y_2}\,)  \\
              &=\colvec{0 \\ (a_1x_1+a_2x_2)-(a_1y_1+a_2y_2) \\
              3(a_1y_1+a_2y_2) }  \\
              &=a_1\colvec{0 \\ x_1-y_1 \\ 3y_1}
                 +a_2\colvec{0 \\ x_2-y_2 \\ 3y_2}  \\
              &=a_1f(\colvec{x_1 \\ y_1})+a_2f(\colvec{x_2 \\ y_2})
          \end{align*}
      \end{exparts}
    
\end{ans}
\begin{ans}{三.II.1.23}
      第一个不是满射；例如，没有多项式被映到常数多项式 $p(x)=1$。
      第二个不是单射；定义域中的下列两个元素
      \begin{equation*}
        \begin{mat}[r]
          1  &0  \\
          0  &0
        \end{mat}
        \quad\text{且}\quad
        \begin{mat}[r]
          0  &0  \\
          0  &1
        \end{mat}
      \end{equation*}
      都映到陪域中的同一个元素 $1\in\Re$。
     
\end{ans}
\begin{ans}{三.II.1.24}
      是；~在任意空间中，\( \text{id}(c\cdot \vec{v}+d\cdot \vec{w})
      = c\cdot \vec{v}+d\cdot \vec{w}
      = c\cdot\text{id}(\vec{v})+d\cdot\text{id}(\vec{w}) \)。
    
\end{ans}
\begin{ans}{三.II.1.25}
      \begin{exparts}
         \partsitem 这个映射不保持结构，因为
           \( f(1+1)=3 \)，而 \( f(1)+f(1)=2 \)。
         \partsitem 验证很直接。
           \begin{align*}
             f(r_1\cdot\colvec{x_1 \\ y_1}+r_2\cdot\colvec{x_2 \\ y_2})
             &=f(\colvec{r_1x_1+r_2x_2 \\ r_1y_1+r_2y_2})               \\
             &=(r_1x_1+r_2x_2)+2(r_1y_1+r_2y_2)                          \\
             &=r_1\cdot (x_1+2y_1)+r_2\cdot (x_2+2y_2)                   \\
             &=r_1\cdot f(\colvec{x_1 \\ y_1})+r_2\cdot f(\colvec{x_2 \\ y_2})
           \end{align*}
      \end{exparts}
     
\end{ans}
\begin{ans}{三.II.1.26}
      保持。
      若 \( \map{h}{V}{W} \) 是线性的，则
      \( h(\vec{u}-\vec{v})
         =h(\vec{u}+(-1)\cdot\vec{v})
         =h(\vec{u})+(-1)\cdot h(\vec{v})
         =h(\vec{u})-h(\vec{v}) \)。
    
\end{ans}
\begin{ans}{三.II.1.27}
      \begin{exparts}
       \partsitem 设 \( \vec{v}\in V \) 关于这组基的表示为
         \( \vec{v}=c_1\vec{\beta}_1+\dots+c_n\vec{\beta}_n \)。
         则 \( h(\vec{v})=h(c_1\vec{\beta}_1+\dots+c_n\vec{\beta}_n)
                  =c_1h(\vec{\beta}_1)+\dots+c_nh(\vec{\beta}_n)
                  =c_1\cdot\zero+\dots+c_n\cdot\zero
                  =\zero \)。
       \partsitem 论证与前一题类似。
         设 \( \vec{v}\in V \) 关于这组基的表示为
         \( \vec{v}=c_1\vec{\beta}_1+\dots+c_n\vec{\beta}_n \)。
         则 \( h(c_1\vec{\beta}_1+\dots+c_n\vec{\beta}_n)
            =c_1h(\vec{\beta}_1)+\dots+c_nh(\vec{\beta}_n)
            =c_1\vec{\beta}_1+\dots+c_n\vec{\beta}_n
            =\vec{v} \)。
       \partsitem 与上面一样，只是此时
         \( c_1h(\vec{\beta}_1)+\dots+c_nh(\vec{\beta}_n)
            =c_1r\vec{\beta}_1+\dots+c_nr\vec{\beta}_n
            =r(c_1\vec{\beta}_1+\dots+c_n\vec{\beta}_n)
            =r\vec{v} \)。
      \end{exparts}
     
\end{ans}
\begin{ans}{三.II.1.28}
      它是同态，可由熟悉的对数运算法则推出：
      乘积的对数等于对数之和 $\ln(ab)=\ln(a)+\ln(b)$，
      幂的对数等于对数的倍数 $\ln(a^r)=r\ln(a)$。
      这个映射是同构，因为它有逆映射，即指数映射，
      所以它是一一对应，从而是同构。
     
\end{ans}
\begin{ans}{三.II.1.29}
       令 \( \hat{x}=x/2 \) 且 \( \hat{y}=y/3 \)，则像集为
       \begin{equation*}
         \set{\colvec{\hat{x} \\ \hat{y}} \suchthat
           \frac{\displaystyle (2\hat{x})^2}{\displaystyle 4}
           +\frac{\displaystyle (3\hat{y})^2}{\displaystyle 9}=1}
         =\set{\colvec{\hat{x} \\ \hat{y}} \suchthat
            \hat{x}^2+\hat{y}^2=1}
       \end{equation*}
       即 \( \hat{x}\hat{y} \) 平面上的单位圆。
     
\end{ans}
\begin{ans}{三.II.1.30}
      圆周长函数 $r\mapsto 2\pi r$ 是线性的。
      因此有
      $2\pi\cdot (r_{\text{地球}}+6)-
         2\pi\cdot (r_{\text{地球}})=12\pi$。
      注意，把紧绕篮球的绳圈抬高到离篮球表面六英尺处，
      与把紧绕地球的绳圈抬高到离地面六英尺处，所需增加的绳长相同。
     
\end{ans}
\begin{ans}{三.II.1.31}
      很容易验证它是线性的。
      \begin{align*}
        h(c_1\cdot \colvec{x_1 \\ y_1 \\ z_1}
          +c_2\cdot \colvec{x_2 \\ y_2 \\ z_2})
        &=h(\colvec{c_1x_1+c_2x_2 \\ c_1y_1+c_2y_2 \\ c_1z_1+c_2z_2})   \\
        &=3(c_1x_1+c_2x_2)-(c_1y_1+c_2y_2)-(c_1z_1+c_2z_2)             \\
        &=c_1\cdot (3x_1-y_1-z_1)+c_2\cdot (3x_2-y_2-z_2)             \\
        &=c_1\cdot h(\colvec{x_1 \\ y_1 \\ z_1})
          +c_2\cdot h(\colvec{x_2 \\ y_2 \\ z_2})
      \end{align*}
      一个自然的推广猜想是：对任意固定的 \( \vec{k}\in\Re^3 \)，
      映射 \( \vec{v}\mapsto\vec{v}\dotprod\vec{k} \) 是线性的。
      这个命题成立。
      它可由前面见过的点积性质推出：
      \( (\vec{v}+\vec{u})\dotprod\vec{k}=\vec{v}\dotprod\vec{k}+
         \vec{u}\dotprod\vec{k} \)
      和 \( (r\vec{v})\dotprod\vec{k}=r(\vec{v}\dotprod\vec{k}) \)。
      （进一步自然地猜想：从 \( \Re^n \) 到 \( \Re \) 的映射，
      若其作用是把输入与固定向量 \( \vec{k}\in\Re^n \) 作点积，
      则它是线性的。这个结论也成立。）
    
\end{ans}
\begin{ans}{三.II.1.32}
      设 \( \map{h}{\Re^1}{\Re^1} \) 是线性的。
      线性映射由它在基上的作用确定，因此固定 \( \Re^1 \) 的基 \( \sequence{1} \)。
      对任意 \( r\in\Re^1 \)，有 \( h(r)=h(r\cdot 1)=r\cdot h(1) \)，
      所以 \( h \) 对任意输入 $r$ 的作用都是乘以常数 \( h(1) \)。
      若 \( h(1) \) 不为零，则这个映射是一一对应\Dash 其逆映射是除以
      \( h(1) \)\Dash 所以 $\Re^1$ 上任意非平凡变换都是同构。

      下面的投影映射说明，当 \( n>1 \) 时（包括 $n=2$），
      并非 \( \Re^n \) 上的每个变换都是乘以常数。
      \begin{equation*}
        \colvec{x_1 \\ x_2 \\ \vdots \\ x_n}
          \mapsto\colvec{x_1 \\ 0 \\ \vdots \\ 0}
      \end{equation*}
    
\end{ans}
\begin{ans}{三.II.1.33}
      设 \( c \) 和 \( d \) 为标量，则有
          \begin{multline*}
            h(c\cdot \colvec{x_1 \\ \vdots \\ x_n}
              +d\cdot \colvec{y_1 \\ \vdots \\ y_n})                     \\
            \begin{aligned}
            &=h(\colvec{cx_1+dy_1 \\ \vdots \\ cx_n+dy_n})        \\
            &=\colvec{a_{1,1}(cx_1+dy_1)+\dots+a_{1,n}(cx_n+dy_n) \\
                         \vdots                                   \\
                         a_{m,1}(cx_1+dy_1)+\dots+a_{m,n}(cx_n+dy_n)} \\
            &=c\cdot\colvec{a_{1,1}x_1+\dots+a_{1,n}x_n \\
                         \vdots                     \\
                         a_{m,1}x_1+\dots+a_{m,n}x_n}
            +d\cdot\colvec{a_{1,1}y_1+\dots+a_{1,n}y_n \\
                         \vdots \\
                         a_{m,1}y_1+\dots+a_{m,n}y_n} \\
            &=c\cdot h(\colvec{x_1 \\ \vdots \\ x_n})
              +d\cdot h(\colvec{y_1 \\ \vdots \\ y_n})
            \end{aligned}
          \end{multline*}
     
\end{ans}
\begin{ans}{三.II.1.34}
        \begin{exparts}
          \partsitem 求导算子的每个幂 $i$ 都是线性的，
            这是因为微积分中熟悉的下列法则成立。
            \begin{equation*}
              \frac{d^i}{dx^i}(\,f(x)+g(x)\,)=\frac{d^i}{dx^i}f(x)
                                         +\frac{d^i}{dx^i}g(x)
              \qquad
              \frac{d^i}{dx^i}\,r\cdot f(x)=r\cdot\frac{d^i}{dx^i}f(x)
            \end{equation*}
         \partsitem
          线性映射的任意线性组合仍是线性映射。
          因此，所给映射是 \( \polyspace_n \) 上的线性变换。
        \end{exparts}
      
\end{ans}
\begin{ans}{三.II.1.35}
      （这个论证已经在同构是等价关系的证明中出现过。）
      设 $\map{f}{U}{V}$ 与 $\map{g}{V}{W}$ 是线性的。
      复合映射保持线性组合
      \begin{multline*}
        \composed{g}{f}(c_1\vec{u}_1+c_2\vec{u}_2)
        =g(\,f(c_1\vec{u}_1+c_2\vec{u}_2)\,)
        =g(\,c_1f(\vec{u}_1)+c_2f(\vec{u}_2)\,)              \\
        =c_1\cdot g(f(\vec{u}_1))+c_2\cdot g(f(\vec{u}_2))
        =c_1\cdot \composed{g}{f}(\vec{u}_1)
           +c_2\cdot \composed{g}{f}(\vec{u}_2)
      \end{multline*}
      其中 $\vec{u}_1,\vec{u}_2\in U$，$c_1,c_2$ 为标量。
    
\end{ans}
\begin{ans}{三.II.1.36}
      \begin{exparts}
        \partsitem 是。
          若各个 $\vec{v}\,$ 线性相关，则各个 $\vec{w}\,$ 不可能线性无关，
          因为定义域中的任意非平凡关系
          \( \zero_V=c_1\vec{v}_1+\dots+c_n\vec{v}_n \)
          都会给出值域中的一个非平凡关系
          \( f(\zero_V)=\zero_W=f(c_1\vec{v}_1+\dots+c_n\vec{v}_n)
             =c_1f(\vec{v}_1)+\dots+c_nf(\vec{v}_n)
             =c_1\vec{w}+\dots+c_n\vec{w}_n \)。
        \partsitem 不一定。
          例如，\( \Re^2 \) 上的变换定义为
          \begin{equation*}
            \colvec{x \\ y} \mapsto \colvec{x+y \\ x+y}
          \end{equation*}
          它把定义域中下面这个线性无关集映成线性相关集。
          \begin{equation*}
            \set{\vec{v}_1,\vec{v}_2}=\set{\colvec[r]{1 \\ 0},\colvec[r]{1 \\ 1}}
            \;\mapsto\;
            \set{\colvec[r]{1 \\ 1},\colvec[r]{2 \\ 2}}=\set{\vec{w}_1,\vec{w}_2}
          \end{equation*}
        \partsitem 不一定。
          例如投影映射 \( \map{\pi}{\Re^3}{\Re^2} \)
          \begin{equation*}
            \colvec{x \\ y \\ z}\mapsto\colvec{x \\ y}
          \end{equation*}
          将下面这个不张成定义域的集合，映成张成陪域的集合。
          \begin{equation*}
            \set{\colvec[r]{1 \\ 0 \\ 0},\colvec[r]{0 \\ 1 \\ 0}}
            \mapsunder{\pi}\set{\colvec[r]{1 \\ 0},\colvec[r]{0 \\ 1}}
          \end{equation*}
        \partsitem 不一定。
          例如，嵌入映射 $\map{\iota}{\Re^2}{\Re^3}$ 把定义域的标准基
          $\stdbasis_2$ 映成一个不张成陪域的集合。
          （\textit{注。}
          不过，各个 $\vec{w}$ 组成的集合确实张成值域。
          这很容易证明。）
      \end{exparts}
     
\end{ans}
\begin{ans}{三.II.1.37}
      回顾一下，$M$ 的转置中第~$i$ 行第~$j$ 列的元素，
      是 $M$ 中第~$j$ 行第~$i$ 列的元素 $m_{j,i}$。
      现在很容易验证。
      从线性组合的转置开始。
      \begin{equation*}
        \trans{[r\cdot\begin{mat}
                \    &\vdots          \\
             \cdots &a_{i,j} &\cdots \\
                    &\vdots
           \end{mat}
         +s\cdot\begin{mat}
                \    &\vdots          \\
             \cdots &b_{i,j} &\cdots \\
                    &\vdots
           \end{mat}]}
      \end{equation*}
      先作线性组合，再转置。
      \begin{equation*}
        =\trans{\begin{mat}
                 \       &\vdots                    \\
                  \cdots &ra_{i,j}+sb_{i,j} &\cdots \\
                         &\vdots
                \end{mat}}
        =\begin{mat}
              \     &\vdots                    \\
            \cdots &ra_{j,i}+sb_{j,i} &\cdots \\
                   &\vdots
          \end{mat}
      \end{equation*}
      然后提出标量，并将各项写回转置形式。
      \begin{multline*}
        =r\cdot\begin{mat}
             \     &\vdots          \\
            \cdots &a_{j,i} &\cdots \\
                   &\vdots
          \end{mat}
         +s\cdot\begin{mat}
             \      &\vdots          \\
            \cdots &b_{j,i} &\cdots \\
                   &\vdots
          \end{mat}                             \\
        =r\cdot\trans{\begin{mat}
             \      &\vdots          \\
            \cdots &a_{j,i} &\cdots \\
                   &\vdots
          \end{mat} }
         +s\cdot\trans{\begin{mat}
              \     &\vdots          \\
            \cdots &b_{j,i} &\cdots \\
                   &\vdots
          \end{mat} }
      \end{multline*}
      定义域是 \( \matspace_{\nbym{m}{n}} \)，
      陪域是 \( \matspace_{\nbym{n}{m}} \)。
     
\end{ans}
\begin{ans}{三.II.1.38}
      \begin{exparts}
        \partsitem 对任意同态 \( \map{h}{\Re^n}{\Re^m} \)，有
          \begin{equation*}
            h(\ell)
            =\set{h(t\cdot\vec{u}+(1-t)\cdot\vec{v})\suchthat t\in [0..1]}
            =\set{t\cdot h(\vec{u})+(1-t)\cdot h(\vec{v})\suchthat t\in [0..1]}
          \end{equation*}
          这就是从 $h(\vec{u})$ 到 $h(\vec{v})$ 的线段。
        \partsitem 需要证明：定义域中的凸集的像，作为值域的子集，也是凸集。
          设 \( C\subseteq \Re^n \) 是凸集，考虑它的像 $h(C)$。
          要证明 $h(C)$ 是凸集，必须证明它的任意两个元素
          $\vec{d}_1$ 与 $\vec{d}_2$ 之间的线段
          \begin{equation*}
            \ell=\set{t\cdot\vec{d}_1+(1-t)\cdot\vec{d}_2\suchthat t\in [0..1]}
          \end{equation*}
          是 $h(C)$ 的子集。

          任取线段上的元素 $\hat{t}\cdot\vec{d}_1+(1-\hat{t})\cdot\vec{d}_2$。
          因为 $\ell$ 的端点在 $C$ 的像中，所以 $C$ 中存在映到它们的元素，
          设 $h(\vec{c}_1)=\vec{d}_1$ 且 $h(\vec{c}_2)=\vec{d}_2$。
          对于本段开头固定的标量 $\hat{t}$，注意
          $h(\hat{t}\cdot\vec{c}_1+(1-\hat{t})\cdot\vec{c}_2)
          =\hat{t}\cdot h(\vec{c}_1)+(1-\hat{t})\cdot h(\vec{c}_2)
          =\hat{t}\cdot\vec{d}_1+(1-\hat{t})\cdot\vec{d}_2$。
          因此 $\ell$ 中的每个元素都属于 $h(C)$，所以 $h(C)$ 是凸集。
      \end{exparts}
    
\end{ans}
\begin{ans}{三.II.1.39}
      \begin{exparts}
        \partsitem 对 \( \vec{v}_0,\vec{v}_1\in\Re^n \)，
          经过 \( \vec{v}_0 \)、方向为 \( \vec{v}_1 \) 的直线是集合
          $\set{\vec{v}_0+t\cdot \vec{v}_1\suchthat t\in\Re}$。
          这条直线在 $h$ 下的像
          $\set{h(\vec{v}_0+t\cdot \vec{v}_1)\suchthat t\in\Re}
             =\set{h(\vec{v}_0)+t\cdot h(\vec{v}_1)\suchthat t\in\Re}$
          是经过 $h(\vec{v}_0)$、方向为 $h(\vec{v}_1)$ 的直线。
          若 \( h(\vec{v}_1) \) 是零向量，则这条直线退化。
        \partsitem \( \Re^n \) 中的 \( k \) 维线性曲面映到
          \( \Re^m \) 中的 \( k \) 维线性曲面（可能退化）。
          证明与直线的情形一样。
      \end{exparts}
     
\end{ans}
\begin{ans}{三.II.1.40}
      设 \( \map{h}{V}{W} \) 是同态，\( S \) 是 \( V \) 的子空间。
      考虑由 \( \hat{h}(\vec{s})=h(\vec{s}) \)
      定义的映射 \( \map{\hat{h}}{S}{W} \)。
      （$\hat{h}$ 与 $h$ 的唯一区别在于定义域。）
      这个新映射是线性的：
      \( \hat{h}(c_1\cdot\vec{s}_1+c_2\cdot\vec{s}_2)=
              h(c_1\vec{s}_1+c_2\vec{s}_2)=c_1h(\vec{s}_1)+c_2h(\vec{s}_2)=
              c_1\cdot\hat{h}(\vec{s}_1)+c_2\cdot\hat{h}(\vec{s}_2) \)。
    
\end{ans}
\begin{ans}{三.II.1.41}
 下一小节将把这一结果作为引理。
      \begin{exparts}
        \partsitem 值域非空，因为 \( V \) 非空。
          最后只需证明它对线性组合封闭。
          值域向量的线性组合具有如下形式，其中 \( \vec{v}_1,\dots,\vec{v}_n\in V \)，
          \begin{equation*}
            c_1\cdot h(\vec{v}_1)+\dots+c_n\cdot h(\vec{v}_n)
            =
            h(c_1\vec{v}_1)+\dots+h(c_n\vec{v}_n)
            =
            h(c_1\cdot \vec{v}_1+\dots+c_n\cdot \vec{v}_n),
          \end{equation*}
          它本身也在值域中，因为 \( c_1\cdot \vec{v}_1+\dots+c_n\cdot \vec{v}_n \)
          属于定义域 \( V \)。
          因此值域是子空间。
        \partsitem 零空间包含 $\zero_V$，因此非空，
          因为 \( \zero_V \) 映到 \( \zero_W \)。
          它对线性组合封闭，因为若 \( \vec{v}_1,\dots,\vec{v}_n\in V \)
          属于原像 \( \set{\vec{v}\in V\suchthat h(\vec{v})=\zero_W} \)，
          则对 \( c_1,\ldots,c_n\in\Re \)，有
          \begin{equation*}
            \zero_W=c_1\cdot h(\vec{v}_1)+\dots+c_n\cdot h(\vec{v}_n)
            =h(c_1\cdot \vec{v}_1+\dots+c_n\cdot \vec{v}_n)
          \end{equation*}
          所以 \( c_1\cdot \vec{v}_1+\dots+c_n\cdot \vec{v}_n \)
          也在 \( \zero_W \) 的原像中。
        \partsitem \( U \) 的像非空，因为 \( U \) 非空。
          关于对线性组合的封闭性，若 \( \vec{u}_1,\ldots,\vec{u}_n\in U \)，则
          \begin{equation*}
            c_1\cdot h(\vec{u}_1)+\dots+c_n\cdot h(\vec{u}_n)
            =
            h(c_1\cdot \vec{u}_1)+\dots+h(c_n\cdot \vec{u}_n)
            =
            h(c_1\cdot \vec{u}_1+\dots+c_n\cdot \vec{u}_n)
          \end{equation*}
          它本身也属于 \( h(U) \)，因为
          \( c_1\cdot \vec{u}_1+\dots+c_n\cdot \vec{u}_n \) 属于 \( U \)。
          因此这个集合是子空间。
        \partsitem 自然的推广是：子空间的原像仍是子空间。

          设 \( X \) 是 \( W \) 的子空间。
          注意 \( \zero_W\in X \)，所以集合
          \( \set{\vec{v}\in V \suchthat h(\vec{v})\in X} \) 非空。
          要证明它对线性组合封闭，取 \( V \) 中的元素 \( \vec{v}_1,\dots,\vec{v}_n \)，
          使得 \( h(\vec{v}_1)=\vec{x}_1 \)，\ldots，
          \( h(\vec{v}_n)=\vec{x}_n \)，并注意
          \begin{equation*}
            h(c_1\cdot \vec{v}_1+\dots+c_n\cdot \vec{v}_n)
            =c_1\cdot h(\vec{v}_1)+\dots+c_n\cdot h(\vec{v}_n)
            =c_1\cdot \vec{x}_1+\dots+c_n\cdot \vec{x}_n
          \end{equation*}
          所以 \( h^{-1}(X) \) 中元素的线性组合仍属于 \( h^{-1}(X) \)。
      \end{exparts}
     
\end{ans}
\begin{ans}{三.II.1.42}
      不是；同构的集合不包含零映射（除非空间是平凡的）。
    
\end{ans}
\begin{ans}{三.II.1.43}
     若 $\sequence{\vec{\beta}_1,\ldots,\vec{\beta}_n}$ 不张成整个空间，
     则映射不一定唯一。
     例如，若只规定 $\vec{e}_1$ 映到自身，来定义 $\Re^2$ 到自身的映射，
     则可能的同态不止一个；恒等映射和投影到第一个分量的映射都满足这个条件。

     若去掉 $\sequence{\vec{\beta}_1,\ldots,\vec{\beta}_n}$ 线性无关的条件，
     则给定的要求可能互相矛盾（即这样的映射可能不存在）。
     例如，考虑 $\sequence{\vec{e}_2,\vec{e}_1,2\vec{e}_1}$，
     尝试定义 $\Re^2$ 到自身的映射，使 $\vec{e}_2$ 映到自身，
     而 $\vec{e}_1$ 和 $2\vec{e}_1$ 都映到 $\vec{e}_1$。
     不存在满足这三个条件的同态。
   
\end{ans}
\begin{ans}{三.II.1.44}
      \begin{exparts}
        \partsitem 线性的验证简述如下。
          \begin{multline*}
            F(r_1\cdot \vec{v}_1+r_2\cdot \vec{v}_2)
            =\colvec{f_1(r_1\vec{v}_1+r_2\vec{v}_2) \\
                        f_2(r_1\vec{v}_1+r_2\vec{v}_2)}           \\
            =r_1\colvec{f_1(\vec{v}_1) \\ f_2(\vec{v}_1)}
            +r_2\colvec{f_1(\vec{v}_2) \\ f_2(\vec{v}_2)}
            =r_1\cdot F(\vec{v}_1)+r_2\cdot F(\vec{v}_2)
          \end{multline*}
        \partsitem 成立。
          设 \( \map{\pi_1}{\Re^2}{\Re^1} \) 和
          \( \map{\pi_2}{\Re^2}{\Re^1} \) 是投影
          \begin{equation*}
            \colvec{x \\ y}\mapsunder{\pi_1} x
              \quad\text{且}\quad
            \colvec{x \\ y}\mapsunder{\pi_2} y
          \end{equation*}
          分别投到两个坐标轴上。
          令 \( f_1(\vec{v})=\pi_1(F(\vec{v})) \) 且
          \( f_2(\vec{v})=\pi_2(F(\vec{v})) \)，
          就得到所需的分量函数。
          \begin{equation*}
            F(\vec{v})=
            \colvec{f_1(\vec{v}) \\ f_2(\vec{v})}
          \end{equation*}
          它们是线性的，因为它们是线性函数的复合，
          而线性函数的复合仍是线性的，
          这一点已在同构是等价关系的证明中说明
          （也可以直接验证它们是线性的）。
        \partsitem 一般地，从向量空间 \( V \) 到 \( \Re^n \) 的映射是线性的，
          当且仅当它的每个分量函数都是线性的。
          验证与前一题相同。
      \end{exparts}
     
\end{ans}
\protect \subsection {三.II.2: 值域空间与零空间}
\begin{ans}{三.II.2.21}
      首先，要回答多项式是否在零空间中，
      必须把它看作定义域 $\polyspace_3$ 的元素。
      要回答它是否在值域空间中，则要把它看作陪域 $\polyspace_4$ 的元素。
      也就是说，对于 $p(x)=x^4$，问它是否在值域空间中是有意义的，
      但问它是否在零空间中则没有意义，因为它甚至不在定义域中。
      \begin{exparts}
        \partsitem 多项式 $x^3\in\polyspace_3$ 不在零空间中，
            因为 $h(x^3)=x^4$ 不是 $\polyspace_4$ 中的零多项式。
            多项式 $x^3\in\polyspace_4$ 在值域空间中，
            因为 $h$ 把 $x^2\in\polyspace_3$ 映到 $x^3$。
        \partsitem 两个问题的答案都是“在，因为 $h(0)=0$”。
            多项式 $0\in\polyspace_3$ 在零空间中，
            因为它被 $h$ 映到 $\polyspace_4$ 中的零多项式。
            多项式 $0\in\polyspace_4$ 在值域空间中，
            因为它是 $0\in\polyspace_3$ 在 $h$ 下的像。
        \partsitem 多项式 $7\in\polyspace_3$ 不在零空间中，
            因为 $h(7)=7x$ 不是 $\polyspace_4$ 中的零多项式。
            多项式 $7\in\polyspace_4$ 不在值域空间中，
            因为定义域中没有元素乘以 $x$ 后能得到常数多项式 $p(x)=7$。
        \partsitem 多项式 $12x-0.5x^3\in\polyspace_3$ 不在零空间中，
            因为 $h(12x-0.5x^3)=12x^2-0.5x^4$。
            多项式 $12x-0.5x^3\in\polyspace_4$ 在值域空间中，
            因为它是 $12-0.5x^2$ 的像。
        \partsitem 多项式 $1+3x^2-x^3\in\polyspace_3$ 不在零空间中，
            因为 $h(1+3x^2-x^3)=x+3x^3-x^4$。
            多项式 $1+3x^2-x^3\in\polyspace_4$ 不在值域空间中，
            因为它有非零常数项。
      \end{exparts}
    
\end{ans}
\begin{ans}{三.II.2.22}
      \begin{exparts}
        \partsitem
          $h$ 的值域是整个陪域~$\Re^2$，因为给定
          \begin{equation*}
            \colvec{x \\ y}\in\Re^2
          \end{equation*}
          它是定义域向量 $0x^2+bx+c$ 在~$h$ 下的像。
          所以~$h$ 的秩为~$2$。
        \partsitem
          值域是 $yz$~平面。任意
          \begin{equation*}
            \colvec{0 \\ a \\ b}
          \end{equation*}
          都是定义域中下列向量在~$f$ 下的像。
          \begin{equation*}
            \colvec{a+b/3  \\ b/3}
          \end{equation*}
          所以这个映射的秩为~$2$。
      \end{exparts}
    
\end{ans}
\begin{ans}{三.II.2.23}
      \begin{exparts}
        \partsitem 值域空间为
          \begin{equation*}
            \rangespace{h}
            =\set{a+ax+ax^2\in\polyspace_3\suchthat a,b\in\Re}
            =\set{a\cdot(1+x+x^2)\suchthat a\in\Re}
          \end{equation*}
          因此秩为一。
        \partsitem 值域空间
          \begin{equation*}
            \rangespace{h}=
             \set{a+d\suchthat a,b,c,d\in\Re}
          \end{equation*}
          是整个 $\Re$（令 $d$ 为 $0$，$a$ 为所需的数，就能得到任意实数）。
          因此秩为一。
        \partsitem 值域空间为
          $\rangespace{h}=\set{r+sx^2\suchthat r,s\in\Re}$。
          秩为二。
        \partsitem 值域空间是 \( \Re^4 \) 的平凡子空间，因此秩为零。
      \end{exparts}
    
\end{ans}
\begin{ans}{三.II.2.24}
      \begin{exparts}
        \partsitem 零空间为
          \begin{align*}
            \nullspace{h}
              &=\set{\colvec{a \\ b}\in\Re^2\suchthat
                                    a+ax+ax^2+0x^3=0+0x+0x^2+0x^3}  \\
              &=\set{\colvec{0 \\ b}\suchthat b\in\Re}
          \end{align*}
          因此零度为一。
        \partsitem 零空间如下。
          \begin{equation*}
            \nullspace{h}
             =\set{\begin{mat}
                     a  &b  \\
                     c  &d
                   \end{mat} \suchthat a+d=0}
             =\set{\begin{mat}
                     -d  &b  \\
                      c  &d
                   \end{mat} \suchthat b,c,d\in\Re }
          \end{equation*}
          因此零度为三。
        \partsitem 零空间
          \begin{equation*}
            \nullspace{h}
             =\set{\begin{mat}
                     a  &b  \\
                     c  &d
                   \end{mat} \suchthat a+b+c=0,\,d=0}
             =\set{\begin{mat}
                     -b-c  &b  \\
                      c  &0
                   \end{mat} \suchthat b,c\in\Re }
          \end{equation*}
          是维数为~$2$ 的空间，因此零度为二。
        \partsitem 定义域中的每个向量都被映到零向量，
          所以零空间为 $\nullspace{h}=\Re^3$。
      \end{exparts}
    
\end{ans}
\begin{ans}{三.II.2.25}
      每一题都利用秩加零度等于定义域维数这一结果。
      \begin{exparts*}
        \partsitem $0$
        \partsitem $3$
        \partsitem $3$
        \partsitem $0$
      \end{exparts*}
    
\end{ans}
\begin{ans}{三.II.2.26}
      因为
      \begin{equation*}
        \frac{d}{dx}\,(a_0+a_1x+\cdots+a_nx^n)
        =a_1+2a_2x+3a_3x^2+\cdots+na_nx^{n-1}
      \end{equation*}
      所以有
      \begin{align*}
        \nullspace{\frac{d}{dx}}
        &=\set{a_0+\cdots+a_nx^n\suchthat a_1+2a_2x+\cdots+na_nx^{n-1}
                                         =0+0x+\cdots+0x^{n-1}}       \\
        &=\set{a_0+\cdots+a_nx^n\suchthat
               \text{$a_1=0$，且 $a_2=0$，\ldots，$a_n=0$}}  \\
        &=\set{a_0+0x+0x^2+\cdots+0x^n\suchthat a_0\in\Re}
      \end{align*}
      同理，
      \begin{equation*}
        \nullspace{\frac{d^k}{dx^k}}
        =\set{a_0+a_1x+\cdots+a_{k-1}x^{k-1}\suchthat a_0,\ldots,a_{k-1}\in\Re}
      \end{equation*}
      其中 \( k\leq n \)。
    
\end{ans}
\begin{ans}{三.II.2.27}
        要证明值域空间是整个~$\Re^2$，只需注意：
        对任意 $u,v\in\Re$，下列方程组都有解。
        \begin{equation*}
          \begin{linsys}{3}
            x &+ &y &  &  &= &u  \\
            x &  &  &+ &z &= &v
          \end{linsys}
          \grstep{-\rho_1+\rho_2}
          \begin{linsys}{3}
            x &+ &y  &  &  &= &u  \\
              &  &-y &+ &z &= &-u+v
          \end{linsys}
        \end{equation*}
        （事实上，因为有自由变量 $z$，它有无穷多个解。）
        因此秩，即值域空间的维数，为~$2$。

        由于秩加零度等于定义域的维数，零度为~$1$。
        求出零空间可以验证这一点：
        \begin{equation*}
          \begin{linsys}{3}
            x &+ &y &  &  &= &0  \\
            x &  &  &+ &z &= &0
          \end{linsys}
          \grstep{-\rho_1+\rho_2}
          \begin{linsys}{3}
            x &+ &y  &  &  &= &0  \\
              &  &-y &+ &z &= &0
          \end{linsys}
        \end{equation*}
        零空间是集合
        \begin{equation*}
          \set{\colvec{-1 \\ 1 \\ 1}\cdot z\suchthat z \in \Re}
        \end{equation*}
        它是一维的。
      
\end{ans}
\begin{ans}{三.II.2.28}
      标量倍数的影子，等于影子的相应标量倍数。
    
\end{ans}
\begin{ans}{三.II.2.29}
      \begin{exparts}
        \partsitem 令 $a_0+(a_0+a_1)x+(a_2+a_3)x^3=0+0x+0x^2+0x^3$，
           得到 $a_0=0$、$a_0+a_1=0$ 且 $a_2+a_3=0$，因此零空间为
           \( \set{-a_3x^2+a_3x^3\suchthat a_3\in\Re} \)。
        \partsitem 令 $a_0+(a_0+a_1)x+(a_2+a_3)x^3=2+0x+0x^2-x^3$，
           得到 $a_0=2$、$a_1=-2$ 且 $a_2+a_3=-1$。
           取 $a_3$ 为参数，并改记为 $a_3=a$，得到集合描述
           \( \set{2-2x+(-1-a)x^2+ax^3\suchthat a\in\Re}
                =\set{(2-2x-x^2)+a\cdot(-x^2+x^3)\suchthat a\in\Re} \)。
        \partsitem 这个集合为空，因为 $h$ 的值域只包含二次项为 $0x^2$ 的多项式。
      \end{exparts}
    
\end{ans}
\begin{ans}{三.II.2.30}
      所有原像都是斜率为 $-2$ 的直线。
      \begin{center}
        \includegraphics{map/mp/ch3.74}
     \end{center}
    
\end{ans}
\begin{ans}{三.II.2.31}
      逆映射如下。
      \begin{exparts}
        \partsitem \( a_0+a_1x+a_2x^2+a_3x^3\mapsto
                        a_0+a_1x+(a_2/2)x^2+(a_3/3)x^3 \)
        \partsitem \( a_0+a_1x+a_2x^2+a_3x^3\mapsto a_0+a_2x+a_1x^2+a_3x^3 \)
        \partsitem \( a_0+a_1x+a_2x^2+a_3x^3\mapsto a_3+a_0x+a_1x^2+a_2x^3 \)
        \partsitem \( a_0+a_1x+a_2x^2+a_3x^3\mapsto
                        a_0+(a_1-a_0)x+(a_2-a_1)x^2+(a_3-a_2)x^3 \)
      \end{exparts}
      例如，对第二题，题目中的映射把
      $0+1x+2x^2+3x^3\mapsto 0+2x+1x^2+3x^3$，
      而上面的逆映射再把 $0+2x+1x^2+3x^3\mapsto 0+1x+2x^2+3x^3$。
      所以这个映射实际上与自身互逆。
     
\end{ans}
\begin{ans}{三.II.2.32}
       对任意向量空间 $V$，零空间
       \begin{equation*}
         \set{\vec{v}\in V\suchthat 2\vec{v}=\zero}
       \end{equation*}
       是平凡的，而值域空间
       \begin{equation*}
         \set{\vec{w}\in V\suchthat
                 \text{存在 $\vec{v}\in V$ 使 $\vec{w}=2\vec{v}$}}
       \end{equation*}
       是整个 \( V \)，因为每个向量 $\vec{w}$ 都是另一个向量的两倍，
       具体地说，是 $(1/2)\vec{w}$ 的两倍。
       （因此，这个变换实际上是自同构。）
    
\end{ans}
\begin{ans}{三.II.2.33}
      由于秩加零度等于定义域的维数（这里为五），且秩至多为三，
      所以可能的数对为：~\( (3,2) \)、\( (2,3) \)、
      \( (1,4) \) 和~\( (0,5) \)。
      很容易给出线性映射，说明每个数对确实都可能出现。
    
\end{ans}
\begin{ans}{三.II.2.34}
      没有（除非 \( \polyspace_n \) 是平凡空间），因为两个多项式
      \( f_0(x)=0 \) 和 \( f_1(x)=1 \) 的导数相同；
      一个映射必须是单射，才有逆映射。
    
\end{ans}
\begin{ans}{三.II.2.35}
      零空间如下。
      \begin{multline*}
        \set{a_0+a_1x+\dots+a_nx^n\suchthat
                   a_0(1)+\frac{\displaystyle a_1}{\displaystyle 2}(1^2)
                   +\dots
                   +\frac{\displaystyle a_n}{\displaystyle n+1}(1^{n+1})=0} \\
        =\set{a_0+a_1x+\dots+a_nx^n\suchthat
                   a_0+(a_1/2)+\dots+(a_{n}/n+1)=0}
      \end{multline*}
      因此零度为 \( n \)。
     
\end{ans}
\begin{ans}{三.II.2.36}
      \begin{exparts}
        \partsitem 一个方向显然成立：~若同态是满射，则值域就是陪域，
          因此秩等于陪域的维数。
          反过来，假设映射的秩等于陪域的维数。
          则值域是陪域的子空间，且维数等于陪域的维数。
          因此值域必定等于陪域，映射是满射。
          （这里“因此”的理由是：值域中存在一个线性无关子集，
          其元素个数等于陪域的维数；
          陪域中这样的线性无关子集必定是陪域的基，所以值域等于陪域。）
        \partsitem 根据\nearbytheorem{th:OOHomoEquivalence}，
          同态是单射，当且仅当零度为零。
          由于秩加零度等于定义域的维数，
          同态是单射，当且仅当秩等于定义域的维数。
          但这里定义域与陪域同维，因此映射是单射，当且仅当它是满射。
      \end{exparts}
   
\end{ans}
\begin{ans}{三.II.2.37}
      要证明的是：\( \map{h}{V}{W} \) 是单射，当且仅当
      对 \( V \) 的每个线性无关子集 \( S \)，
      \( W \) 的子集 \( h(S)=\set{h(\vec{s})\suchthat \vec{s}\in S} \)
      都线性无关。

      一个方向很容易\Dash 根据\nearbytheorem{th:OOHomoEquivalence}，
      若 \( h \) 不是单射，则零空间非平凡，即其中不仅有零向量。
      因此，取零空间中的 \( \vec{v}\neq\zero_V \)，
      单元素集合 \( \set{\vec{v\,}} \) 线性无关，
      而其像 \( \set{h(\vec{v})}=\set{\zero_W} \) 不是线性无关的。

      反过来，假设 \( h \) 是单射，
      则根据\nearbytheorem{th:OOHomoEquivalence}，其零空间平凡。
      于是，对任意 \( \vec{v}_1,\dots,\vec{v}_n\in V \)，关系
      \begin{equation*}
        \zero_W=
        c_1\cdot h(\vec{v}_1)+\dots+c_n\cdot h(\vec{v}_n)
        =h(c_1\cdot \vec{v}_1+\dots+c_n\cdot \vec{v}_n)
      \end{equation*}
      蕴含 \( c_1\cdot \vec{v}_1+\dots+c_n\cdot \vec{v}_n=\zero_V \)。
      因此，若 \( V \) 的子集线性无关，则它在 \( W \) 中的像也线性无关。

      \textit{注。}
      这里说的是：线性映射是单射，当且仅当它对\emph{所有}集合都保持线性无关性
      （即一个集合线性无关时，它的像也线性无关）。
      不是单射的映射完全可能保持某些集合的线性无关性。
      例如，下面这个从 $\Re^3$ 到 $\Re^2$ 的映射。
      \begin{equation*}
        \colvec{x \\ y \\ z}\mapsto\colvec{x+y+z \\ 0}
      \end{equation*}
      它保持下面这个集合的线性无关性
      \begin{equation*}
        \set{\colvec[r]{1 \\ 0 \\ 0}}\mapsto\set{\colvec[r]{1 \\ 0}}
      \end{equation*}
      也保持下面这个集合的线性无关性（这个例子稍微巧妙一些）
      \begin{equation*}
        \set{\colvec[r]{1 \\ 0 \\ 0},\colvec[r]{0 \\ 1 \\ 0}}
        \mapsto\set{\colvec[r]{1 \\ 0}}
      \end{equation*}
      （回顾一下，集合中的重复元素不计两次）。
      不过，确实存在一些集合，其线性无关性不被这个映射保持
      \begin{equation*}
        \set{\colvec[r]{1 \\ 0 \\ 0},\colvec[r]{0 \\ 2 \\ 0}}
        \mapsto\set{\colvec[r]{1 \\ 0},\colvec[r]{2 \\ 0}}
      \end{equation*}
      所以它并不保持所有集合的线性无关性。
    
\end{ans}
\begin{ans}{三.II.2.38}
      （沿用\nearbytheorem{th:HomoDetActOnBasis} 中的记号。）
      固定 $V$ 的基 $\sequence{\vec{\beta}_1,\dots,\vec{\beta}_n}$
      和 $W$ 的基 $\sequence{\vec{w}_1,\ldots,\vec{w}_k}$。
      若 $W$ 的维数 $k$ 不大于 $V$ 的维数 $n$，
      则由定理得到如下确定的从 $V$ 到 $W$ 的线性映射。
      \begin{equation*}
        \vec{\beta}_1\mapsto\vec{w}_1,\,\dots,\,\vec{\beta}_k\mapsto\vec{w}_k
        \quad\text{且}\quad
        \vec{\beta}_{k+1}\mapsto\vec{w}_k,
           \,\dots,\,\vec{\beta}_n\mapsto\vec{w}_k
      \end{equation*}
      只需验证它是满射。

      $W$ 的任意元素都可以写成基向量的线性组合
      $c_1\cdot \vec{w}_1+\dots+c_k\cdot \vec{w}_k$。
      这个向量是
      $c_1\cdot \vec{\beta}_1+\dots+c_k\cdot \vec{\beta}_k
      +0\cdot \vec{\beta}_{k+1}\dots+0\cdot \vec{\beta}_n$
      在上述映射下的像。
      因此映射是满射。
    
\end{ans}
\begin{ans}{三.II.2.39}
      存在。
      对于 \( \Re^2 \) 上的变换
      \begin{equation*}
        \colvec{x \\ y}\mapsunder{h}\colvec{0 \\ x}
      \end{equation*}
      有
      \begin{equation*}
        \nullspace{h}=\set{\colvec{0 \\ y}\suchthat y\in\Re}
        =\rangespace{h}
      \end{equation*}
      \textit{注。}
      第五章将进一步讨论这种情形。
    
\end{ans}
\begin{ans}{三.II.2.40}
      只需简单计算。
          \begin{align*}
            h(\spanof{S})
            &=\set{h(c_1\vec{s}_1+\dots+c_n\vec{s}_n)
                 \suchthat \text{\( c_1,\dots,c_n\in\Re \)
                                 且 \( \vec{s}_1,\dots,\vec{s}_n\in S \)} } \\
            &=\set{c_1h(\vec{s}_1)+\dots+c_nh(\vec{s}_n)
                 \suchthat \text{\( c_1,\dots,c_n\in\Re \)
                                 且 \( \vec{s}_1,\dots,\vec{s}_n\in S \)} } \\
            &=\spanof{h(S)}
          \end{align*}
     
\end{ans}
\begin{ans}{三.II.2.41}
      \begin{exparts}
        \partsitem
          通过互相包含来证明集合相等
          $h^{-1}(\vec{w})
            =\set{\vec{v}+\vec{n}\suchthat \vec{n}\in\nullspace{h} }$。
          对于
          $\set{\vec{v}+\vec{n}\suchthat \vec{n}\in\nullspace{h} }
            \subseteq h^{-1}(\vec{w})$，
          只需注意 $h(\vec{v}+\vec{n})=h(\vec{v})+h(\vec{n})$ 等于 $\vec{w}$，
          因此前一个集合的任意元素都属于后一个集合。
          对于
          $h^{-1}(\vec{w})
            \subseteq\set{\vec{v}+\vec{n}\suchthat \vec{n}\in\nullspace{h} }$，
          考虑 $\vec{u}\in h^{-1}(\vec{w})$。
          由于 \( h \) 是线性的，\( h(\vec{u})=h(\vec{v}) \)
          蕴含 \( h(\vec{u}-\vec{v})=\zero \)。
          把 \( \vec{u}-\vec{v} \) 记为 \( \vec{n} \)，
          则 $\vec{u}\in\set{\vec{v}+\vec{n}\suchthat \vec{n}\in\nullspace{h} }$，
          因为 $\vec{u}=\vec{v}+(\vec{u}-\vec{v})$，正如所需。
        \partsitem 验证很直接。
        \partsitem 结论立即成立。
        \partsitem 简述线性的验证：设 \( c,d \) 是标量，
          \( \vec{x},\vec{y}\in\Re^n \) 的分量分别为
          \( x_1,\dots,x_n \) 和 \( y_1,\dots,y_n \)，则有
          \begin{align*}
            h(c\cdot \vec{x}+d\cdot \vec{y})
            &=\colvec{a_{1,1}(cx_1+dy_1)+\dots+a_{1,n}(cx_n+dy_n) \\
                         \vdots \\
                         a_{m,1}(cx_1+dy_1)+\dots+a_{m,n}(cx_n+dy_n)}    \\
            &=\colvec{a_{1,1}cx_1+\dots+a_{1,n}cx_n \\
                         \vdots \\
                         a_{m,1}cx_1+\dots+a_{m,n}cx_n}
            +\colvec{a_{1,1}dy_1+\dots+a_{1,n}dy_n \\
                         \vdots \\
                         a_{m,1}dy_1+\dots+a_{m,n}dy_n}        \\
            &=c\cdot h(\vec{x})+d\cdot h(\vec{y})
          \end{align*}
          相应的结论是
          \( \text{通解}=\text{特解}+\text{齐次解} \)。
        \partsitem 求导算子的每个幂都是线性的，因为微积分中的法则
          \begin{equation*}
            \frac{d^k}{dx^k}(f(x)+g(x))=\frac{d^k}{dx^k}f(x)
                                         +\frac{d^k}{dx^k}g(x)
            \quad\text{且}\quad
            \frac{d^k}{dx^k}rf(x)=r\frac{d^k}{dx^k}f(x)
          \end{equation*}
          成立。
          因此，给定映射是这个空间上的线性变换，
          因为根据\nearbylemma{le:SpLinFcns}，线性映射的任意线性组合仍是线性映射。
          相应的结论为
          \( \text{通解}=\text{特解}+\text{齐次解} \)，
          其中对应的齐次微分方程右端的常数为 \( 0 \)。
      \end{exparts}
     
\end{ans}
\begin{ans}{三.II.2.42}
      由于 \( t \) 的秩为一，\( t \) 的值域空间是一维的。
      取 $\sequence{h(\vec{v})}$ 为基（其中 $\vec{v}$ 选取得当），
      则对任意 $\vec{w}\in V$，像 $h(\vec{w})\in V$ 都是这个基向量的倍数\Dash
      每个 $\vec{w}$ 都对应一个标量 \( c_{\vec{w}} \)，
      使得 \( t(\vec{w})=c_{\vec{w}}t(\vec{v}) \)。
      对等式两边应用 \( t \)，并取 \( r \) 为 \( c_{t(\vec{v})} \)，
      \begin{equation*}
        \composed{t}{t}(\vec{w})
        =t(c_{\vec{w}}\cdot t(\vec{v}))
        =c_{\vec{w}}\cdot\composed{t}{t}(\vec{v})
        =c_{\vec{w}}\cdot c_{t(\vec{v})}\cdot t(\vec{v})
        =c_{\vec{w}}\cdot r\cdot t(\vec{v})
        =r\cdot c_{\vec{w}}\cdot t(\vec{v})
        =r\cdot t(\vec{w})
      \end{equation*}
      就得到所需结论。
    
\end{ans}
\begin{ans}{三.II.2.43}
      根据假设，\( h \) 不是零映射，所以存在不在零空间中的向量 \( \vec{v}\in V \)。
      注意 \( \sequence{h(\vec{v})} \) 是 \( \Re \) 的一组基，
      因为它是 \( \Re \) 的单元素线性无关子集。
      因而 \( h \) 是满射，因为对任意 \( r\in\Re \)，
      存在标量 \( c \) 使 \( r=c\cdot h(\vec{v}) \)，所以 \( r=h(c\vec{v}) \)。

      因此 \( h \) 的秩为一。
      由于零度为 $n$，\( h \) 的定义域，即向量空间 \( V \)，维数为 \( n+1 \)。
      最后只需证明 \( \set{\vec{v},\vec{\beta}_1,\dots,\vec{\beta}_n} \) 线性无关，
      因为它是维数为~$n+1$ 的空间中含~$n+1$ 个元素的子集。
      由于 \( \set{\vec{\beta}_1,\dots,\vec{\beta}_n} \) 线性无关，
      只需证明 \( \vec{v} \) 不是其余向量的线性组合。
      但若 \( c_1\vec{\beta}_1+\dots+c_n\vec{\beta}_n=\vec{v} \)，
      则 \( -\vec{v}+c_1\vec{\beta}_1+\dots+c_n\vec{\beta}_n=\zero \)，
      对两边应用 \( h \) 就会产生矛盾。
    
\end{ans}
\begin{ans}{三.II.2.44}
      固定 \( V \) 的一组基 \( \sequence{\vec{\beta}_1,\dots,\vec{\beta}_n} \)。
      我们将证明下列映射
      \begin{equation*}
        h\mapsunder{\Phi}
        \colvec{h(\vec{\beta}_1) \\ \vdots \\ h(\vec{\beta}_n)}
      \end{equation*}
      是从 \( V^\ast \) 到 \( \Re^n \) 的同构。

      要证明 $\Phi$ 是单射，设 $h_1$ 和 $h_2$ 属于 $V^\ast$，
      且 $\Phi(h_1)=\Phi(h_2)$。
      则
      \begin{equation*}
        \colvec{h_1(\vec{\beta}_1) \\ \vdots \\ h_1(\vec{\beta}_n)}
        =\colvec{h_2(\vec{\beta}_1) \\ \vdots \\ h_2(\vec{\beta}_n)}
      \end{equation*}
      因而 $h_1(\vec{\beta}_1)=h_2(\vec{\beta}_1)$，依此类推。
      但同态由它在基上的作用确定，所以 \( h_1=h_2 \)，因此 \( \Phi \) 是单射。

      要证明 $\Phi$ 是满射，考虑
      \begin{equation*}
        \colvec{x_1 \\ \vdots \\ x_n}
      \end{equation*}
      其中 \( x_1,\ldots,x_n\in\Re \)。
      下面这个从 $V$ 到 $\Re$ 的函数 $h$
      \begin{equation*}
        c_1\vec{\beta}_1+\dots+c_n\vec{\beta}_n
        \mapsunder{h} c_1x_1+\dots+c_nx_n
      \end{equation*}
      是线性的，且 $\Phi$ 将它映到 $\Re^n$ 中给定的向量，
      所以 \( \Phi \) 是满射。

      映射 \( \Phi \) 也保持结构：~若
      \begin{align*}
        c_1\vec{\beta}_1+\dots+c_n\vec{\beta}_n
        &\mapsunder{h_1}
        c_1h_1(\vec{\beta}_1)+\dots+c_nh_1(\vec{\beta}_n)    \\
        c_1\vec{\beta}_1+\dots+c_n\vec{\beta}_n
        &\mapsunder{h_2}
        c_1h_2(\vec{\beta}_1)+\dots+c_nh_2(\vec{\beta}_n)
      \end{align*}
      则有
      \begin{multline*}
        (r_1h_1+r_2h_2)(c_1\vec{\beta}_1+\dots+c_n\vec{\beta}_n)       \\
        \begin{aligned}
        &=
        c_1(r_1h_1(\vec{\beta}_1)+r_2h_2(\vec{\beta}_1))
          +\dots
          +c_n(r_1h_1(\vec{\beta}_n)+r_2h_2(\vec{\beta}_n))    \\
        &=
        r_1(c_1h_1(\vec{\beta}_1)+\dots+c_nh_1(\vec{\beta}_n))
          + r_2(c_1h_2(\vec{\beta}_1)+\dots+c_nh_2(\vec{\beta}_n))
        \end{aligned}
      \end{multline*}
      所以 \( \Phi(r_1h_1+r_2h_2)=r_1\Phi(h_1)+r_2\Phi(h_2) \)。
    
\end{ans}
\begin{ans}{三.II.2.45}
      设 \( \map{h}{V}{W} \) 是线性的，固定 \( V \) 的一组基
      \( \sequence{\vec{\beta}_1,\dots,\vec{\beta}_n} \)。
      考虑下面 \( n \) 个从 \( V \) 到 \( W \) 的映射
      \begin{equation*}
        h_1(\vec{v})=c_1\cdot h(\vec{\beta}_1),
        \quad
        h_2(\vec{v})=c_2\cdot h(\vec{\beta}_2),
        \quad\ldots\quad,
        h_n(\vec{v})=c_n\cdot h(\vec{\beta}_n)
      \end{equation*}
      其中 \( \vec{v}=c_1\vec{\beta}_1+\dots+c_n\vec{\beta}_n \) 为任意向量。
      显然，\( h \) 是各个 \( h_i \) 之和。
      只需验证每个 \( h_i \) 都是线性的：~若
      \( \vec{u}=d_1\vec{\beta}_1+\dots+d_n\vec{\beta}_n \)，则
      \( h_i(r\vec{v}+s\vec{u})=rc_i+sd_i=rh_i(\vec{v})+sh_i(\vec{u}) \)。
    
\end{ans}
\begin{ans}{三.II.2.46}
     可以是（显然成立），也可以不是（几乎显然）。

     如果把 \( V \) “同态于” \( W \) 理解为存在从 \( V \) 到
     \( W \) 的同态（不一定是满射），
     那么每个空间都同态于其他任意空间，因为零映射总是存在。

     如果把 \( V \) “同态于” \( W \) 理解为存在从 \( V \) 到 \( W \) 的满射同态，
     那么这个关系就不是等价关系。
     例如，存在从 \( \Re^3 \) 到 \( \Re^2 \) 的满射同态
     （投影就是一个），但根据\nearbycorollary{cor:RankDecreases}，
     不存在从 \( \Re^2 \) 满射到 \( \Re^3 \) 的同态，
     所以这个关系不具有自反性。\appendrefs{equivalence relations}
   
\end{ans}
\begin{ans}{三.II.2.47}
      它们构成上述链是显然的。
      对于其余结论，这里证明
      \( \rangespace{t^{j+1}}=\rangespace{t^{j}} \)
      蕴含 \( \rangespace{t^{j+2}}=\rangespace{t^{j+1}} \)，
      然后应用数学归纳法即可。

      假设 \( \rangespace{t^{j+1}}=\rangespace{t^{j}} \)。
      则 \( \map{t}{\rangespace{t^{j+1}}}{\rangespace{t^{j+2}}} \)
      与 \( \map{t}{\rangespace{t^{j}}}{\rangespace{t^{j+1}}} \)
      是同一个映射，定义域也相同。
      因此值域相同：\( \rangespace{t^{j+2}}=\rangespace{t^{j+1}} \)。
    
\end{ans}
\protect \section {第 III 节：线性映射的计算}
\protect \subsection {三.III.1: 用矩阵表示线性映射}
\begin{ans}{三.III.1.13}
      \begin{exparts*}
\partsitem \(
           \colvec{1\cdot 2+3\cdot 1+1\cdot 0         \\
                        0\cdot 2+(-1)\cdot 1+2\cdot 0 \\
                        1\cdot 2+1\cdot 1+0\cdot 0}
           =\colvec[r]{5 \\ -1 \\ 3}   \)
        \partsitem 没有定义。
        \partsitem \(  \colvec[r]{0 \\ 0 \\ 0}  \)
      \end{exparts*}
    
\end{ans}
\begin{ans}{三.III.1.14}
      \begin{exparts}
\partsitem 没有定义。
\partsitem 有定义。
          \begin{equation*}
            \begin{mat}
              3  &1  \\
              2  &4
            \end{mat}
            \colvec{0 \\ -1}
            =
            \colvec{-1 \\ -4}
          \end{equation*}
\partsitem 没有定义。
      \end{exparts}
    
\end{ans}
\begin{ans}{三.III.1.15}
      \begin{exparts*}
\partsitem $\colvec{2\cdot 4 +1\cdot 2 \\
                            3\cdot 4-(1/2)\cdot 2}
                   =\colvec[r]{10 \\ 11}$
        \partsitem $\colvec[r]{4 \\ 1}$
        \partsitem 没有定义。
      \end{exparts*}
    
\end{ans}
\begin{ans}{三.III.1.16}
矩阵与向量相乘得到一个线性方程组。
      \begin{equation*}
        \begin{linsys}{3}
          2x  &+  &y  &+  &z  &=  &8  \\
              &   &y  &+  &3z &=  &4  \\
           x  &-  &y  &+  &2z &=  &4
        \end{linsys}
      \end{equation*}
高斯消元给出 \( z=1 \)、\( y=1 \) 和 \( x=3 \)。
    
\end{ans}
\begin{ans}{三.III.1.17}
下面给出两种求解方法。

首先，显然 $1-3x+2x^2=1\cdot 1-3\cdot x+2\cdot x^2$，利用保持线性组合的性质得
$h(1-3x+2x^2)
       =h(1\cdot 1-3\cdot x+2\cdot x^2)
       =1\cdot h(1)-3\cdot h(x)+2\cdot h(x^2)
       =1\cdot (1+x)-3\cdot (1+2x)+2\cdot (x-x^3)
       =-2-3x-2x^3$。

另一种方法使用本小节建立的计算方案。
由于知道这些空间元素的像，我们选取 \( B=\sequence{1,x,x^2} \) 为定义域的基。
陪域的基可任取为 \( D=\sequence{1,x,x^2,x^3} \)。在这些选择下，
      \begin{equation*}
        \rep{h}{B,D}
        =\begin{mat}[r]
           1   &1  &0  \\
           1   &2  &1  \\
           0   &0  &0  \\
           0   &0  &-1
         \end{mat}_{B,D}
      \end{equation*}
并且，由于
      \begin{equation*}
        \rep{1-3x+2x^2}{B}=\colvec[r]{1 \\ -3 \\ 2}_B
      \end{equation*}
矩阵与向量的乘法计算给出
      \begin{equation*}
        \rep{h(1-3x+2x^2)}{D}=
         \begin{mat}[r]
           1   &1  &0  \\
           1   &2  &1  \\
           0   &0  &0  \\
           0   &0  &-1
         \end{mat}_{B,D}
         \colvec[r]{1 \\ -3 \\ 2}_B
         =\colvec[r]{-2 \\ -3 \\ 0 \\ -2}_D
      \end{equation*}
因此，\( h(1-3x+2x^2)
              =-2\cdot 1-3\cdot x+0\cdot x^2-2\cdot x^3
              =-2-3x-2x^3 \)，与上面相同。
    
\end{ans}
\begin{ans}{三.III.1.18}
为 $\matspace_{\nbyn{2}}$ 固定以下自然基。
     \begin{equation*}
       D=\sequence{
       \begin{mat}
          1 &0  \\
          0 &0
       \end{mat},
       \begin{mat}
          0 &1  \\
          0 &0
       \end{mat},
       \begin{mat}
          0 &0  \\
          1 &0
       \end{mat},
       \begin{mat}
          0 &0  \\
          0 &1
       \end{mat}}
     \end{equation*}
映射~$h$ 相对于 $\stdbasis_2,D$ 的表示如下。
     \begin{equation*}
       \rep{h}{\stdbasis_2,D}=
       \begin{mat}
         1  &0  \\
         2  &-1  \\
         0  &1  \\
         1  &0
       \end{mat}
     \end{equation*}
由于一般向量相对于~$\stdbasis_2$ 的表示就是它自身，同样，每个矩阵相对于~$D$ 的表示由它自身的元素给出，映射的作用如下。
     \begin{equation*}
       \rep{h(\colvec{x \\ y})}{D}
       =
       \begin{mat}
         1  &0  \\
         2  &-1  \\
         0  &1  \\
         1  &0
       \end{mat}
       \colvec{x \\ y}
       =
       \colvec{x \\ 2x-y \\ y \\ x}
     \end{equation*}
   
\end{ans}
\begin{ans}{三.III.1.19}
再次回顾本小节所述：相对于 $\stdbasis_i$，列向量的表示就是它自身。
      \begin{exparts}
\partsitem 要表示 \( h \) 相对于 \( \stdbasis_2,\stdbasis_3 \) 的矩阵，先求定义域基向量的像，再相对于陪域的基表示这些像。第一个为
          \begin{equation*}
            \rep{\,h(\vec{e}_1)\,}{\stdbasis_3}
            =\rep{\colvec[r]{2 \\ 2 \\ 0}}{\stdbasis_3}
            =\colvec[r]{2 \\ 2 \\ 0}
          \end{equation*}
第二个为
          \begin{equation*}
            \rep{\,h(\vec{e}_2)\,}{\stdbasis_3}
            =\rep{\colvec[r]{0 \\ 1 \\ -1}}{\stdbasis_3}
            =\colvec[r]{0 \\ 1 \\ -1}
          \end{equation*}
将它们并排拼接成矩阵。
          \begin{equation*}
            \rep{h}{\stdbasis_2,\stdbasis_3}=
            \begin{mat}[r]
              2  &0  \\
              2  &1  \\
              0  &-1
            \end{mat}
          \end{equation*}
\partsitem 对定义域 \( \Re^2 \) 中的任意 \( \vec{v} \)，
          \begin{equation*}
            \rep{\vec{v}}{\stdbasis_2}
            =\rep{\colvec{v_1 \\ v_2}}{\stdbasis_2}
            =\colvec{v_1 \\ v_2}
          \end{equation*}
所以
          \begin{equation*}
            \rep{\,h(\vec{v})\,}{\stdbasis_3}
            =\begin{mat}[r]
              2  &0  \\
              2  &1  \\
              0  &-1
            \end{mat}
            \colvec{v_1 \\ v_2}
            =\colvec{2v_1 \\ 2v_1+v_2 \\ -v_2}
          \end{equation*}
就是所求的表示。
      \end{exparts}
    
\end{ans}
\begin{ans}{三.III.1.20}
该映射在定义域基向量上的作用如下。
      \begin{equation*}
          \colvec{1 \\ 1 \\ 1}\mapsto\colvec{2 \\ 2}
          \quad
          \colvec{1 \\ 1 \\ 0}\mapsto\colvec{2 \\ 1}
          \quad
          \colvec{1 \\ 0 \\ 0}\mapsto\colvec{1 \\ 1}
      \end{equation*}
将这些像相对于陪域的基表示出来。
      \begin{equation*}
        \rep{\colvec{2 \\ 2}}{D}=\colvec{2 \\ 1}_D
        \quad
        \rep{\colvec{2 \\ 1}}{D}=\colvec{2 \\ 1/2}_D
        \quad
        \rep{\colvec{1 \\ 1}}{D}=\colvec{1 \\ 1/2}_D
      \end{equation*}
将它们拼接成矩阵。
      \begin{equation*}
        \rep{h}{B,D}=
        \begin{mat}
          2 &2   &1 \\
          1 &1/2 &1/2
        \end{mat}
      \end{equation*}
    
\end{ans}
\begin{ans}{三.III.1.21}
      \begin{exparts*}
        \partsitem
先求定义域中每个基向量的像，再将它相对于陪域的基表示出来。
        \begin{multline*}
          \rep{\frac{d\,1}{dx}}{B}=\colvec[r]{0 \\ 0 \\ 0 \\ 0}
          \quad
          \rep{\frac{d\,x}{dx}}{B}=\colvec[r]{1 \\ 0 \\ 0 \\ 0}
          \quad
          \rep{\frac{d\,x^2}{dx}}{B}=\colvec[r]{0 \\ 2 \\ 0 \\ 0}
          \\
          \rep{\frac{d\,x^3}{dx}}{B}=\colvec[r]{0 \\ 0 \\ 3 \\ 0}
        \end{multline*}
将这些表示并排拼接，得到映射的表示矩阵。
        \begin{equation*}
          \rep{\frac{d}{dx}}{B,B}=
          \begin{mat}[r]
            0  &1  &0  &0  \\
            0  &0  &2  &0  \\
            0  &0  &0  &3  \\
            0  &0  &0  &0
          \end{mat}
        \end{equation*}
\partsitem 与上一问一样，先表示定义域基向量的像，
        \begin{multline*}
          \rep{\frac{d\,1}{dx}}{D}=\colvec[r]{0 \\ 0 \\ 0 \\ 0}
          \quad
          \rep{\frac{d\,x}{dx}}{D}=\colvec[r]{1 \\ 0 \\ 0 \\ 0}
          \quad
          \rep{\frac{d\,x^2}{dx}}{D}=\colvec[r]{0 \\ 1 \\ 0 \\ 0}
          \\
          \rep{\frac{d\,x^3}{dx}}{D}=\colvec[r]{0 \\ 0 \\ 1 \\ 0}
        \end{multline*}
再并排拼接成矩阵。
        \begin{equation*}
          \rep{\frac{d}{dx}}{B,D}=
          \begin{mat}[r]
            0  &1  &0  &0  \\
            0  &0  &1  &0  \\
            0  &0  &0  &1  \\
            0  &0  &0  &0
          \end{mat}
        \end{equation*}
      \end{exparts*}
    
\end{ans}
\begin{ans}{三.III.1.22}
每一问都要先求定义域各个基向量的像，再将这些像相对于陪域的基表示出来，最后将这些表示并排拼接成矩阵。
      \begin{exparts}
\partsitem 定义域基向量的像为
          \begin{equation*}
            1\mapsto 0
            \quad x\mapsto 1
            \quad x^2\mapsto 2x
            \quad \ldots
          \end{equation*}
这些像相对于陪域的基的表示为
          \begin{multline*}
            \rep{0}{B}=\colvec{0 \\ 0 \\ 0 \\ \vdots \\ \  \\  \ }
            \quad
            \rep{1}{B}=\colvec{1 \\ 0 \\ 0 \\ \vdots \\ \  \\ \ }
            \quad
            \rep{2x}{B}=\colvec{0 \\ 2 \\ 0 \\ \vdots \\ \  \\ \ }        \\
            \ldots\quad
            \rep{nx^{n-1}}{B}=\colvec{0 \\ 0 \\ 0 \\ \vdots \\ n \\ 0}
          \end{multline*}
矩阵
          \begin{equation*}
            \rep{\frac{d}{dx}}{B,B}
            =\begin{mat}
              0  &1  &0  &\ldots  &0  \\
              0  &0  &2  &\ldots  &0  \\
                 &\vdots             \\
              0  &0  &0  &\ldots  &n  \\
              0  &0  &0  &\ldots  &0
            \end{mat}
          \end{equation*}
有 $n+1$ 行和 $n+1$ 列。
\partsitem 求出定义域基向量在这个映射下的像后，
          \begin{equation*}
            1\mapsto x
            \quad x\mapsto x^2/2
            \quad x^2\mapsto x^3/3
            \quad \ldots
          \end{equation*}
将它们相对于陪域的基表示出来，
          \begin{multline*}
            \rep{x}{B_{n+1}}=\colvec{0 \\ 1 \\ 0 \\ \vdots \\ \ }
            \quad
            \rep{x^2/2}{B_{n+1}}=\colvec{0 \\ 0 \\ 1/2 \\ \vdots \\ \ }   \\
            \ldots\quad
            \rep{x^{n+1}/(n+1)}{B_{n+1}}
                     =\colvec{0 \\ 0 \\ 0 \\ \vdots \\ 1/(n+1)}
          \end{multline*}
再拼接成矩阵。
          \begin{equation*}
            \rep{\int}{B_{n},B_{n+1}}
            =\begin{mat}
              0  &0  &\ldots  &0      &0  \\
              1  &0  &\ldots  &0      &0  \\
              0  &1/2&\ldots  &0      &0  \\
                 &\vdots                  \\
              0  &0  &\ldots  &0      &1/(n+1)
            \end{mat}
          \end{equation*}
\partsitem 定义域基向量的像为
          \begin{equation*}
            1\mapsto 1
            \quad x\mapsto 1/2
            \quad x^2\mapsto 1/3
            \quad \ldots
          \end{equation*}
它们相对于陪域的基的表示为
          \begin{equation*}
            \rep{1}{\stdbasis_1}=1
            \quad \rep{1/2}{\stdbasis_1}=1/2
            \quad \ldots
          \end{equation*}
所以矩阵为
          \begin{equation*}
            \rep{\int}{B,\stdbasis_1}
            =\begin{mat}
              1  &1/2 &\cdots  &1/n    &1/(n+1)
            \end{mat}
          \end{equation*}
（这是一个 $\nbym{1}{(n+1)}$ 矩阵）。
\partsitem 定义域基向量的像为
          \begin{equation*}
            1\mapsto 1
            \quad x\mapsto 3
            \quad x^2\mapsto 9
            \quad \ldots
          \end{equation*}
它们在陪域中的表示为
          \begin{equation*}
            \rep{1}{\stdbasis_1}=1
            \quad\rep{3}{\stdbasis_1}=3
            \quad\rep{9}{\stdbasis_1}=9
            \quad \ldots
          \end{equation*}
因此矩阵如下。
          \begin{equation*}
            \rep{\text{eval}_3}{B,\stdbasis_1}
            =\begin{mat}[r]
              1  &3   &9  &\cdots  &3^n
            \end{mat}
          \end{equation*}
\partsitem 定义域基向量的像为
$1\mapsto 1$、$x\mapsto x+1=1+x$、$x^2\mapsto (x+1)^2=1+2x+x^2$、
$x^3\mapsto (x+1)^3=1+3x+3x^2+x^3$，依此类推。它们的表示如下。
          \begin{equation*}
            \rep{1}{B}=\colvec[r]{1 \\ 0 \\ 0 \\ 0 \\ \vdotswithin{0} \\ 0}
            \quad
            \rep{1+x}{B}=\colvec{1 \\ 1 \\ 0 \\ 0 \\ \vdotswithin{0} \\ 0}
            \quad
            \rep{1+2x+x^2}{B}=\colvec{1 \\ 2 \\ 1 \\ 0 \\ \vdotswithin{0} \\ 0}
            \quad\ldots
          \end{equation*}
得到的矩阵
          \begin{equation*}
            \renewcommand{\arraystretch}{1.2}
            \rep{\text{slide}_{-1}}{B,B}
            =\begin{mat}
              1  &1   &1  &1  &\ldots  &1           \\
              0  &1   &2  &3  &\ldots  &\binom{n}{1}  \\
              0  &0   &1  &3  &\ldots  &\binom{n}{2}  \\
                 &\vdots                        \\
              0  &0   &0  &   &\ldots  &1
            \end{mat}
          \end{equation*}
就是\definend{帕斯卡三角形}\index{Pascal's triangle}
（回顾，$\binom{n}{r}$ 表示从含 $n$ 个元素的集合中不计顺序、不重复地选取 $r$ 个元素的方法数）。
      \end{exparts}
    
\end{ans}
\begin{ans}{三.III.1.23}
若空间的维数为 \( n \)，则
      \begin{equation*}
        \rep{\text{id}}{B,B}=
        \begin{mat}
          1  &0  \ldots  &0  \\
          0  &1  \ldots  &0  \\
             &\vdots         \\
          0  &0  \ldots  &1
        \end{mat}_{B,B}
      \end{equation*}
是 $\nbyn{n}$ 单位矩阵。
    
\end{ans}
\begin{ans}{三.III.1.24}
取以下自然基，
      \begin{equation*}
        B=\sequence{\vec{\beta}_1,\vec{\beta}_2,\vec{\beta}_3,\vec{\beta}_4}
         =\sequence{
            \begin{mat}[r]
              1  &0  \\
              0  &0
            \end{mat},
            \begin{mat}[r]
              0  &1  \\
              0  &0
            \end{mat},
            \begin{mat}[r]
              0  &0  \\
              1  &0
            \end{mat},
            \begin{mat}[r]
              0  &0  \\
              0  &1
            \end{mat}   }
      \end{equation*}
转置映射的作用为
      \begin{equation*}
        \vec{\beta}_1\mapsto\vec{\beta}_1
        \quad \vec{\beta}_2\mapsto\vec{\beta}_3
        \quad \vec{\beta}_3\mapsto\vec{\beta}_2
        \quad \vec{\beta}_4\mapsto\vec{\beta}_4
      \end{equation*}
将这些像相对于陪域的基表示出来，再将所得列向量并排拼接，得到
      \begin{equation*}
        \rep{\text{trans}}{B,B}=
        \begin{mat}[r]
          1  &0  &0  &0  \\
          0  &0  &1  &0  \\
          0  &1  &0  &0  \\
          0  &0  &0  &1
        \end{mat}_{B,B}
      \end{equation*}
     
\end{ans}
\begin{ans}{三.III.1.25}
       \begin{exparts}
\partsitem 定义域基向量的像相对于陪域的基的表示为
           \begin{equation*}
             \rep{\vec{\beta}_2}{B}=\colvec[r]{0 \\ 1 \\ 0 \\ 0}
             \quad
             \rep{\vec{\beta}_3}{B}=\colvec[r]{0 \\ 0 \\ 1 \\ 0}
             \quad
             \rep{\vec{\beta}_4}{B}=\colvec[r]{0 \\ 0 \\ 0 \\ 1}
             \quad
             \rep{\zero}{B}=\colvec[r]{0 \\ 0 \\ 0 \\ 0}
           \end{equation*}
所以，该变换的表示矩阵如下。
           \begin{equation*}
             \begin{mat}[r]
                   0  &0  &0  &0  \\
                   1  &0  &0  &0  \\
                   0  &1  &0  &0  \\
                   0  &0  &1  &0
                 \end{mat}
         \end{equation*}
         \partsitem
              $\begin{mat}[r]
                   0  &0  &0  &0  \\
                   1  &0  &0  &0  \\
                   0  &0  &0  &0  \\
                   0  &0  &1  &0
                 \end{mat}$
         \partsitem
               $\begin{mat}[r]
                   0  &0  &0  &0  \\
                   1  &0  &0  &0  \\
                   0  &1  &0  &0  \\
                   0  &0  &0  &0
                 \end{mat}$
       \end{exparts}
     
\end{ans}
\begin{ans}{三.III.1.26}
      \begin{exparts}
        \partsitem The picture of $\map{d_s}{\Re^2}{\Re^2}$ is this.
          \begin{center}  \small
            \includegraphics{map/mp/ch3.75}
          \end{center}
这个映射在定义域标准基向量上的作用为
         \begin{equation*}
            \colvec[r]{1 \\ 0}\mapsunder{d_s}\colvec{s \\ 0}
            \qquad
            \colvec[r]{0 \\ 1}\mapsunder{d_s}\colvec{0 \\ s}
         \end{equation*}
这些像相对于陪域的基（仍为标准基）的表示就是它们自身。
         \begin{equation*}
           \rep{\colvec{s \\ 0}}{\stdbasis_2}=\colvec{s \\ 0}
           \qquad
           \rep{\colvec{0 \\ s}}{\stdbasis_2}=\colvec{0 \\ s}
         \end{equation*}
因此，伸缩映射的表示如下。
         \begin{equation*}
           \rep{d_s}{\stdbasis_2,\stdbasis_2}
           =\begin{mat}
              s  &0  \\
              0  &s
           \end{mat}
        \end{equation*}
      \partsitem The picture of $\map{f_\ell}{\Re^2}{\Re^2}$ is this.
         \begin{center}  \small
           \includegraphics{map/mp/ch3.76}
        \end{center}
计算可知（参见练习~I.\ref{exer:RigidPlaneMapsAutos}），当直线斜率为 $k$ 时，
       \begin{equation*}
         \colvec[r]{1 \\ 0}
             \mapsunder{f_\ell}\colvec{(1-k^2)/(1+k^2) \\ 2k/(1+k^2)}
         \qquad
         \colvec[r]{0 \\ 1}
             \mapsunder{f_\ell}\colvec{2k/(1+k^2) \\ -(1-k^2)/(1+k^2)}
       \end{equation*}
（斜率不存在的情形需另行处理，但很简单），所以反射的表示矩阵如下。
       \begin{equation*}
         \rep{f_\ell}{\stdbasis_2,\stdbasis_2}
         =\frac{1}{1+k^2}\cdot\begin{mat}
            1-k^2  &2k  \\
            2k     &-(1-k^2)
         \end{mat}
       \end{equation*}
      \end{exparts}
    
\end{ans}
\begin{ans}{三.III.1.27}
将该映射记为 \( \map{t}{\Re^2}{\Re^2} \)。
      \begin{exparts}
\partsitem 要相对于标准基表示这个映射，必须求出定义域基向量 $\vec{e}_1$ 和 $\vec{e}_2$ 的像，再表示这些像。$\vec{e}_1$ 的像已给出。

求 $\vec{e}_2$ 的像的一种方法是直接观察\Dash 可以看出
          \begin{equation*}
            \colvec[r]{1 \\ 1}-\colvec[r]{1 \\ 0}=\colvec[r]{0 \\ 1}
            \;\mapsunder{t}\;
            \colvec[r]{2 \\ 0}-\colvec[r]{-1 \\ 0}=\colvec[r]{3 \\ 0}
          \end{equation*}

更系统的方法是利用已知信息表示这个变换，再用该表示求像。取以下基，
          \begin{equation*}
            C=\sequence{\colvec[r]{1 \\ 1},\colvec[r]{1 \\ 0}}
          \end{equation*}
已知信息给出
          \begin{equation*}
            \rep{t}{C,\stdbasis_2}
            \begin{mat}[r]
              2  &-1  \\
              0  &0
            \end{mat}
          \end{equation*}
由于
          \begin{equation*}
            \rep{\vec{e}_2}{C}=\colvec[r]{1 \\ -1}_C
          \end{equation*}
可得
          \begin{equation*}
            \rep{t(\vec{e}_2)}{\stdbasis_2}
            =\begin{mat}[r]
               2  &-1  \\
               0  &0
             \end{mat}_{C,\stdbasis_2}
            \colvec[r]{1 \\ -1}_C
            =\colvec[r]{3 \\ 0}_{\stdbasis_2}
          \end{equation*}
因此 $t(\vec{e}_2)=3\cdot\vec{e}_1$（因为相对于标准基，这个向量的表示就是它自身）。所以，$t$ 相对于 $\stdbasis_2,\stdbasis_2$ 的表示如下。
          \begin{equation*}
            \rep{t}{\stdbasis_2,\stdbasis_2}
            =\begin{mat}[r]
               -1 &3   \\
               0  &0
             \end{mat}_{\stdbasis_2,\stdbasis_2}
          \end{equation*}
\partsitem 为了使用上一问的矩阵，注意到
          \begin{equation*}
            \rep{\colvec[r]{0 \\ 5}}{\stdbasis_2}=\colvec[r]{0 \\ 5}_{\stdbasis_2}
          \end{equation*}
所以，给定向量的像相对于陪域的基的表示如下。
          \begin{equation*}
            \rep{t(\colvec[r]{0 \\ 5})}{\stdbasis_2}
            =\begin{mat}[r]
              -1  &3   \\
               0  &0
             \end{mat}_{\stdbasis_2,\stdbasis_2}
            \colvec[r]{0 \\ 5}_{\stdbasis_2}
            =\colvec[r]{15 \\ 0}_{\stdbasis_2}
          \end{equation*}
由于陪域的基是标准基，陪域中向量的表示就是它们自身，故有
          \begin{equation*}
            t(\colvec[r]{0 \\ 5})
            =\colvec[r]{15 \\ 0}
          \end{equation*}
\partsitem 先求 \( B \) 中每个元素的像，再将这些像相对于 \( D \) 表示出来。第一步可以使用前面得到的矩阵。
          \begin{equation*}
            \rep{\colvec[r]{1 \\-1}}{\stdbasis_2}
            =\begin{mat}[r]
              -1  &3   \\
               0  &0
             \end{mat}_{\stdbasis_2,\stdbasis_2}
            \colvec[r]{1 \\ -1}_{\stdbasis_2}
            =\colvec[r]{-4 \\ 0}_{\stdbasis_2}
            \quad\text{so}\quad
            t(\colvec[r]{1 \\ -1})=\colvec[r]{-4 \\ 0}
          \end{equation*}
实际上，$B$ 的第二个元素的像已在题目中给出，无须使用矩阵计算。
          \begin{equation*}
            t(\colvec[r]{1 \\ 1})=\colvec[r]{2 \\ 0}
          \end{equation*}
现在，按常规方法求这些像相对于 $D$ 的表示。
          \begin{equation*}
            \rep{\colvec[r]{-4 \\ 0}}{D}=\colvec[r]{-1 \\ 2}_D
            \quad\text{and}\quad
            \rep{\colvec[r]{2 \\ 0}}{D}=\colvec[r]{1/2 \\ -1}_D
          \end{equation*}
因此矩阵如下。
          \begin{equation*}
            \rep{t}{B,D}=
            \begin{mat}[r]
              -1  &1/2  \\
               2  &-1
            \end{mat}_{B,D}
          \end{equation*}
\partsitem 从上一问已知定义域基向量的像。
          \begin{equation*}
             t(\colvec[r]{1 \\ -1})=\colvec[r]{-4 \\ 0}
             \qquad
             t(\colvec[r]{1 \\ 1})=\colvec[r]{2 \\ 0}
          \end{equation*}
可以计算这些像相对于陪域的基的表示。
          \begin{equation*}
             \rep{\colvec[r]{-4 \\ 0}}{B}=\colvec[r]{-2 \\ -2}_B
             \quad\text{and}\quad
             \rep{\colvec[r]{2 \\ 0}}{B}=\colvec[r]{1 \\ 1}_B
          \end{equation*}
因此矩阵如下。
          \begin{equation*}
            \rep{t}{B,B}=
            \begin{mat}[r]
              -2  &1  \\
              -2  &1
            \end{mat}_{B,B}
          \end{equation*}
      \end{exparts}
    
\end{ans}
\begin{ans}{三.III.1.28}
      \begin{exparts}
\partsitem 定义域基向量的像为
          \begin{equation*}
            \vec{\beta}_1\mapsto h(\vec{\beta}_1)
            \quad
            \vec{\beta}_2\mapsto h(\vec{\beta}_2)
            \quad\ldots\quad
            \vec{\beta}_n\mapsto h(\vec{\beta}_n)
          \end{equation*}
这些像相对于陪域的基的表示为
          \begin{multline*}
            \rep{\,h(\vec{\beta}_1)\,}{h(B)}=\colvec[r]{1 \\ 0 \\ \vdotswithin{0} \\ 0}
            \quad
            \rep{\,h(\vec{\beta}_2)\,}{h(B)}=\colvec[r]{0 \\ 1 \\ \vdotswithin{0} \\ 0}  \\
            \ldots\quad
            \rep{\,h(\vec{\beta}_n)\,}{h(B)}=\colvec[r]{0 \\ 0 \\ \vdotswithin{0} \\ 1}
          \end{multline*}
所以该矩阵是单位矩阵。
          \begin{equation*}
            \rep{h}{B,h(B)}
            =\begin{mat}
               1  &0  &\ldots  &0  \\
               0  &1  &        &0  \\
                  &   &\ddots      \\
               0  &0  &        &1
            \end{mat}
          \end{equation*}
\partsitem 使用上一问的矩阵，得到以下表示。
          \begin{equation*}
            \rep{\,h(\vec{v})\,}{h(B)}
             =\colvec{c_1 \\ \vdots \\ c_n}_{h(B)}
          \end{equation*}
        \end{exparts}
     
\end{ans}
\begin{ans}{三.III.1.29}
乘积
      \begin{equation*}
        \begin{mat}
          h_{1,1} &\ldots  &h_{1,i} &\ldots &h_{1,n}  \\
          h_{2,1} &\ldots  &h_{2,i} &\ldots &h_{2,n}  \\
                  &\vdots                             \\
          h_{m,1} &\ldots  &h_{m,i} &\ldots &h_{1,n}
        \end{mat}
        \colvec{0 \\ \vdots \\ 1 \\ \vdots \\ 0}
        =
        \colvec{h_{1,i} \\ h_{2,i} \\ \vdots \\ h_{m,i}}
      \end{equation*}
给出矩阵的第 \( i \) 列。
    
\end{ans}
\begin{ans}{三.III.1.30}
      \begin{exparts}
\partsitem 定义域基向量的像为 \( \cos x\mapsunder{d/dx}-\sin x \) 和 \( \sin x\mapsunder{d/dx}\cos x \)。
将它们相对于陪域的基（仍为 $B$）表示出来，再并排拼接，得到以下矩阵。
          \begin{equation*}
            \rep{\frac{d}{dx}}{B,B}=
            \begin{mat}[r]
                0  &1  \\
               -1  &0
            \end{mat}_{B,B}
          \end{equation*}
\partsitem 定义域基向量的像为 \( e^x\mapsunder{d/dx}e^x \) 和 \( e^{2x}\mapsunder{d/dx}2e^{2x} \)。
相对于陪域的基表示并拼接，得到以下矩阵。
          \begin{equation*}
            \rep{\frac{d}{dx}}{B,B}=
            \begin{mat}[r]
                1  &0  \\
                0  &2
            \end{mat}_{B,B}
          \end{equation*}
\partsitem 定义域基向量的像为 \( 1\mapsunder{d/dx}0 \)、\( x\mapsunder{d/dx}1 \)、\( e^{x}\mapsunder{d/dx}e^{x} \) 和 \( xe^{x}\mapsunder{d/dx}e^x+xe^x \)。
相对于 $B$ 表示并拼接，得到以下矩阵。
          \begin{equation*}
            \rep{\frac{d}{dx}}{B,B}=
            \begin{mat}[r]
                0  &1  &0  &0 \\
                0  &0  &0  &0 \\
                0  &0  &1  &1 \\
                0  &0  &0  &1
            \end{mat}_{B,B}
          \end{equation*}
      \end{exparts}
    
\end{ans}
\begin{ans}{三.III.1.31}
      \begin{exparts}
\partsitem 值域由陪域中相对于陪域的基具有以下表示的向量组成。
          \begin{equation*}
            \set{
              \begin{mat}[r]
                1  &0  \\
                0  &0
              \end{mat}
              \colvec{x  \\ y}
              \suchthat x,y\in\Re}
            =\set{\colvec{x  \\ 0}
                  \suchthat x,y\in\Re}
          \end{equation*}
由于陪域的基是 $\stdbasis_2$，每个向量的表示就是它自身，所以该变换的值域是 $x$ 轴。
\partsitem 值域由陪域中具有以下表示的向量组成。
          \begin{equation*}
            \set{
              \begin{mat}[r]
                0  &0  \\
                3  &2
              \end{mat}
              \colvec{x  \\ y}
              \suchthat x,y\in\Re}
            =\set{\colvec{0  \\ 3x+2y}
                  \suchthat x,y\in\Re}
          \end{equation*}
相对于 $\stdbasis_2$，向量的表示就是它自身，所以这个值域是 $y$ 轴。
\partsitem 相对于 $\stdbasis_2$ 具有以下表示的向量集合
          \begin{align*}
            \set{
              \begin{mat}
                a   &b  \\
                2a  &2b
              \end{mat}
              \colvec{x  \\ y}
              \suchthat x,y\in\Re}
            &=\set{\colvec{ax+by  \\ 2ax+2by}
                  \suchthat x,y\in\Re}                    \\
            &=\set{(ax+by)\cdot\colvec[r]{1  \\ 2}
                  \suchthat x,y\in\Re}
          \end{align*}
在 $a$ 或 $b$ 至少一个不为零时是直线 $y=2x$，在两者均为零时是仅含原点的集合。
      \end{exparts}
    
\end{ans}
\begin{ans}{三.III.1.32}
能，有两方面的原因。

首先，两个映射 $h$ 和 $\hat{h}$ 不必具有相同的定义域和陪域。例如，
      \begin{equation*}
        \begin{mat}[r]
          1  &2  \\
          3  &4
        \end{mat}
      \end{equation*}
相对于标准基表示映射 \( \map{h}{\Re^2}{\Re^2} \)，其作用为
      \begin{equation*}
        \colvec[r]{1 \\ 0}\mapsto\colvec[r]{1 \\ 3}
        \quad\text{and}\quad
        \colvec[r]{0 \\ 1}\mapsto\colvec[r]{2 \\ 4}
      \end{equation*}
同时，它还相对于 \( \sequence{1,x} \) 和 \( \stdbasis_2 \) 表示映射
\( \map{\hat{h}}{\polyspace_1}{\Re^2} \)，其作用如下。
      \begin{equation*}
        1\mapsto\colvec[r]{1 \\ 3}
        \quad\text{and}\quad
        x\mapsto\colvec[r]{2 \\ 4}
      \end{equation*}

其次，即使 \( h \) 和 \( \hat{h} \) 的定义域、陪域相同，不同的基也会产生不同的映射。例如，$\nbyn{2}$ 单位矩阵
      \begin{equation*}
        I=\begin{mat}[r]
          1  &0  \\
          0  &1
        \end{mat}
      \end{equation*}
相对于 $\stdbasis_2,\stdbasis_2$ 表示 $\Re^2$ 上的恒等映射。
但若定义域的基为 $\stdbasis_2$，陪域的基为 $D=\sequence{\vec{e}_2,\vec{e}_1}$，同一矩阵 $I$ 就表示交换第一、第二个分量的映射，
      \begin{equation*}
        \colvec{x \\ y}\mapsto\colvec{y \\ x}
      \end{equation*}
即关于直线 $y=x$ 的反射。
    
\end{ans}
\begin{ans}{三.III.1.33}
仿照 \nearbyexample{ex:TypLinMapRepByMat}，只需将数换成字母。

将 \( B \) 写成 \( \sequence{\vec{\beta}_1,\ldots,\vec{\beta}_n} \)，将 \( D \) 写成 \( \sequence{\vec{\delta}_1,\dots,\vec{\delta}_m} \)。
由映射相对于基的表示的定义，假设
      \begin{equation*}
        \rep{h}{B,D}
        =\begin{mat}
           h_{1,1} &\ldots  &h_{1,n}  \\
           \vdots  &        &\vdots   \\
           h_{m,1} &\ldots  &h_{m,n}
         \end{mat}
      \end{equation*}
意味着 $h(\vec{\beta}_i)=h_{i,1}\vec{\delta}_1+\dots+h_{i,n}\vec{\delta}_n$。
又由向量相对于基的表示的定义，假设
      \begin{equation*}
        \rep{\vec{v}}{B}=\colvec{c_1 \\ \vdots \\ c_n}
      \end{equation*}
意味着 \( \vec{v}=c_1\vec{\beta}_1+\cdots+c_n\vec{\beta}_n \)。代入得
      \begin{align*}
        h(\vec{v})
        &=h(c_1\cdot\vec{\beta}_1+\dots+c_n\cdot\vec{\beta}_n)      \\
        &=c_1\cdot h(\vec{\beta}_1)+\dots+c_n\cdot \vec{\beta}_n    \\
        &=c_1\cdot (h_{1,1}\vec{\delta}_1+\dots+h_{m,1}\vec{\delta}_m)
        +\dots
        +c_n\cdot (h_{1,n}\vec{\delta}_1+\dots+h_{m,n}\vec{\delta}_m) \\
        &=(h_{1,1}c_1+\dots+h_{1,n}c_n)\cdot\vec{\delta}_1
        +\cdots
        +(h_{m,1}c_1+\dots+h_{m,n}c_n)\cdot\vec{\delta}_m
      \end{align*}
因此 $h(\vec{v})$ 的表示正如所求。
    
\end{ans}
\begin{ans}{三.III.1.34}
      \begin{exparts}
        \partsitem The picture is this.
          \begin{center}  \small
            \includegraphics{map/mp/ch3.77}
         \end{center}
定义域基向量的像
         \begin{equation*}
           \colvec[r]{1 \\ 0 \\ 0}\mapsto\colvec[r]{1 \\ 0 \\ 0}
           \qquad
           \colvec[r]{0 \\ 1 \\ 0}\mapsto\colvec{0 \\ \cos\theta \\ -\sin\theta}
           \qquad
           \colvec[r]{0 \\ 0 \\ 1}\mapsto\colvec{0 \\ \sin\theta \\ \cos\theta}
         \end{equation*}
相对于陪域的基（仍为 $\stdbasis_3$）的表示就是它们自身，因此并排拼接这些表示得到以下矩阵。
         \begin{equation*}
            \rep{r_\theta}{\stdbasis_3,\stdbasis_3}=
            \begin{mat}
              1  &0       &0                \\
              0  &\cos\theta   &\sin\theta   \\
              0  &-\sin\theta  &\cos\theta
            \end{mat}
         \end{equation*}
\partsitem 图形与上一问类似。定义域基向量的像
         \begin{equation*}
           \colvec[r]{1 \\ 0 \\ 0}\mapsto\colvec{\cos\theta \\ 0 \\ \sin\theta}
           \qquad
           \colvec[r]{0 \\ 1 \\ 0}\mapsto\colvec[r]{0 \\ 1 \\ 0}
           \qquad
           \colvec[r]{0 \\ 0 \\ 1}\mapsto\colvec{-\sin\theta \\ 0 \\ \cos\theta}
         \end{equation*}
相对于陪域的基 $\stdbasis_3$ 的表示就是它们自身，因此矩阵如下。
         \begin{equation*}
            \begin{mat}
              \cos\theta  &0       &-\sin\theta   \\
              0           &1       &0             \\
              \sin\theta  &0       &\cos\theta
            \end{mat}
         \end{equation*}
\partsitem 对沿竖直 $z$ 轴站立的观察者，$xy$ 平面上的顺时针旋转从 $y$ 轴正半轴转向 $x$ 轴正半轴。
也就是说，其方向与 \nearbyexample{exam:RepsOfRigidPlaneMaps} 相反。定义域基向量的像
         \begin{equation*}
           \colvec[r]{1 \\ 0 \\ 0}\mapsto\colvec{\cos\theta \\ -\sin\theta \\ 0}
           \qquad
           \colvec[r]{0 \\ 1 \\ 0}\mapsto\colvec{\sin\theta \\ \cos\theta \\ 0}
           \qquad
           \colvec[r]{0 \\ 0 \\ 1}\mapsto\colvec[r]{0 \\ 0 \\ 1}
         \end{equation*}
相对于 $\stdbasis_3$ 的表示就是它们自身，所以矩阵如下。
         \begin{equation*}
            \begin{mat}
              \cos\theta    &\sin\theta  &0   \\
              -\sin\theta   &\cos\theta  &0    \\
              0             &0           &1
            \end{mat}
         \end{equation*}
        \partsitem
            $\begin{mat}
              \cos\theta  &\sin\theta &0 &0 \\
              -\sin\theta  &\cos\theta  &0 &0 \\
              0           &0           &1 &0 \\
              0           &0           &0 &1
            \end{mat}$
      \end{exparts}
    
\end{ans}
\begin{ans}{三.III.1.35}
      \begin{exparts}
        \partsitem
将基 \( B_U \) 写为 \( \sequence{\vec{\beta}_1,\dots,\vec{\beta}_k} \)，再将 $B_V$ 写为其扩充
\( \sequence{\vec{\beta}_1,\dots,\vec{\beta}_k,\vec{\beta}_{k+1},\dots,\vec{\beta}_n} \)。若
          \begin{equation*}
            \rep{\vec{v}}{B_U}=\colvec{c_1 \\ \vdots \\ c_k}
          \end{equation*}
即 $\vec{v}=c_1\cdot\vec{\beta}_1+\cdots+c_k\cdot\vec{\beta}_k$，则
          \begin{equation*}
            \rep{\vec{v}}{B_V}=\colvec{c_1 \\ \vdots\\ c_k \\ 0 \\ \vdots \\ 0}
          \end{equation*}
因为 $\vec{v}=c_1\cdot\vec{\beta}_1+\dots+c_k\cdot\vec{\beta}_k+0\cdot\vec{\beta}_{k+1}+\dots+0\cdot\vec{\beta}_n$。
        \partsitem
先明确问题的含义。比较 \( \map{h}{V}{W} \) 及其在子空间上的限制
\( \map{\restrictionmap{h}{U}}{U}{W} \)。限制映射的值域是 \( W \) 的子空间，故固定这个值域的基 \( D_{h(U)} \)，再将其扩充为 \( W \) 的基 \( D_V \)。我们要找出以下两者的关系。
          \begin{equation*}
            \rep{h}{B_V,D_V}
            \quad\text{and}\quad
            \rep{\restrictionmap{h}{U}}{B_U,D_{h(U)}}
          \end{equation*}
答案可直接由上一问得到：若
          \begin{equation*}
            \rep{\restrictionmap{h}{U}}{B_U,D_{h(U)}}
             =\begin{mat}
                h_{1,1}  &\ldots  &h_{1,k}  \\
                \vdots   &        &\vdots   \\
                h_{p,1}  &\ldots  &h_{p,k}
              \end{mat}
          \end{equation*}
则扩展映射的表示如下。
          \begin{equation*}
            \rep{h}{B_V,D_V}
             =\begin{mat}
                h_{1,1}  &\ldots  &h_{1,k}  &h_{1,k+1}  &\ldots  &h_{1,n}  \\
                \vdots   &        &         &           &        &\vdots   \\
                h_{p,1}  &\ldots  &h_{p,k}  &h_{p,k+1}  &\ldots  &h_{p,n}  \\
                0        &\ldots  &0        &h_{p+1,k+1}&\ldots  &h_{p+1,n}  \\
                \vdots   &        &         &           &        &\vdots   \\
                0        &\ldots  &0        &h_{m,k+1}  &\ldots  &h_{m,n}
              \end{mat}
          \end{equation*}
\partsitem 取 \( W_i \) 为 \( \set{h(\vec{\beta}_1),\dots,h(\vec{\beta}_i)} \) 的张成空间。
\partsitem 将第二问的答案应用于第三问。
\partsitem 不唯一。例如，向 \( x \) 轴的投影 \( \map{\pi_x}{\Re^2}{\Re^2} \) 可用以下两个上三角矩阵表示，
          \begin{equation*}
             \rep{\pi_x}{\stdbasis_2,\stdbasis_2}=
             \begin{mat}[r]
               1  &0  \\
               0  &0
             \end{mat}
             \quad\text{and}\quad
             \rep{\pi_x}{C,\stdbasis_2}=
             \begin{mat}[r]
               0  &1  \\
               0  &0
             \end{mat}
          \end{equation*}
其中 \( C=\sequence{\vec{e}_2,\vec{e}_1} \)。
      \end{exparts}
    
\end{ans}
\protect \subsection {三.III.2: 任意矩阵都表示线性映射}
\begin{ans}{三.III.2.12}
      \begin{exparts}
\partsitem 定义域：$2$，陪域：$2$
\partsitem 定义域：$3$，陪域：$2$
\partsitem 定义域：$2$，陪域：$3$
\partsitem 定义域：$3$，陪域：$2$
\partsitem 定义域：$4$，陪域：$2$
      \end{exparts}
    
\end{ans}
\begin{ans}{三.III.2.13}
每一问只需判断矩阵是否非奇异，必要时使用高斯消元。（实际上，这些矩阵都可以直接观察。）
      \begin{exparts}
\partsitem 该矩阵非奇异，因为第二行不是第一行的倍数，所以映射非奇异。
\partsitem 矩阵奇异，所以映射奇异。
\partsitem 非奇异。
\partsitem 非奇异。
      \end{exparts}
    
\end{ans}
\begin{ans}{三.III.2.14}
该向量相对于 $B$ 的表示如下。
     \begin{equation*}
       \rep{2x-1}{B}=\colvec{-1 \\ 3}
     \end{equation*}
利用矩阵与向量的乘积，可计算 $\rep{h(\vec{v})}{D}$。
     \begin{equation*}
       \rep{h(2x-1)}{D}
       =
       \begin{mat}
         2  &1  \\
         4  &2
       \end{mat}
       \colvec{-1 \\ 3}_B
       =
       \colvec{1 \\ 2}_D
     \end{equation*}
由这个表示可以求出 $h(\vec{v})$。
     \begin{equation*}
       h(2x-1)=1\cdot\colvec{1 \\ 1}+2\cdot\colvec{1 \\ 0}
       =\colvec{3 \\ 1}
     \end{equation*}
   
\end{ans}
\begin{ans}{三.III.2.15}
如本小节所述，相对于标准基，表示是直接的。例如，第一个矩阵描述以下映射。
     \begin{equation*}
       \colvec[r]{1 \\ 0 \\ 0}=\colvec[r]{1 \\ 0 \\ 0}_{\stdbasis_3}
          \!\mapsto\colvec[r]{1 \\ 0}_{\stdbasis_2}=\colvec[r]{1 \\ 0}
       \qquad
       %\colvec[r]{0 \\ 1 \\ 0}=\colvec[r]{0 \\ 1 \\ 0}_{\stdbasis_3}
       %   \!\mapsto\colvec[r]{1 \\ 1}_{\stdbasis_2}=\colvec[r]{1 \\ 1}
       \colvec[r]{0 \\ 1 \\ 0}
          \!\mapsto\colvec[r]{1 \\ 1}
       \qquad
       %\colvec[r]{0 \\ 0 \\ 1}=\colvec[r]{0 \\ 0 \\ 1}_{\stdbasis_3}
       %   \!\mapsto\colvec[r]{3 \\ 4}_{\stdbasis_2}=\colvec[r]{3 \\ 4}
       \colvec[r]{0 \\ 0 \\ 1}
          \!\mapsto\colvec[r]{3 \\ 4}
     \end{equation*}
因此，第一问就是问是否存在标量使得
     \begin{equation*}
       c_1\colvec[r]{1 \\ 0}+c_2\colvec[r]{1 \\ 1}
           +c_3\colvec[r]{3 \\ 4}=\colvec[r]{1 \\ 3}
     \end{equation*}
即该向量是否在矩阵的列空间中。
     \begin{exparts}
\partsitem 是。照常列出相应的线性方程组并使用高斯消元，即可得到结论。
也可观察列向量等式，发现取 $c_3=3/4$、$c_1=-5/4$、$c_2=0$ 即可。
第三种方法是注意到矩阵的秩为二，等于陪域维数，所以映射是满射\Dash 值域为整个 $\Re^2$，当然包含给定向量。
\partsitem 否。矩阵每一列的第二个分量都是第一个的两倍，而给定向量不满足这一点。也可计算矩阵的列空间，得到
         \begin{equation*}
           \set{c_1\colvec[r]{2 \\ 4}
                +c_2\colvec[r]{0 \\ 0}
                +c_3\colvec[r]{3 \\ 6} \suchthat c_1,c_2,c_3\in\Re}
           =\set{c\colvec[r]{1 \\ 2}\suchthat c\in\Re}
         \end{equation*}
（这正是刚才指出的事实，不过这里通过计算而非观察得到），给定向量不在此集合中。
     \end{exparts}
    
\end{ans}
\begin{ans}{三.III.2.16}
      \begin{exparts}
\partsitem 基的第一个元素
          \begin{equation*}
            \colvec[r]{0 \\ 1}=\colvec[r]{1 \\ 0}_B
          \end{equation*}
映到
          \begin{equation*}
            \colvec[r]{1/2 \\ -1/2}_D
          \end{equation*}
即陪域中的以下元素。
          \begin{equation*}
            \frac{1}{2}\cdot\colvec[r]{1 \\ 1}
              -\frac{1}{2}\cdot\colvec[r]{1 \\ -1}
              =\colvec[r]{0 \\ 1}
          \end{equation*}
\partsitem 基的第二个元素按如下方式映射，
          \begin{equation*}
            \colvec[r]{1 \\ 0}=\colvec[r]{0 \\ 1}_B
            \mapsto
            \colvec[r]{(1/2 \\ 1/2}_D
          \end{equation*}
得到陪域中的以下元素。
          \begin{equation*}
            \frac{1}{2}\cdot\colvec[r]{1 \\ 1}
              +\frac{1}{2}\cdot\colvec[r]{1 \\ -1}
              =\colvec[r]{1 \\ 0}
          \end{equation*}
\partsitem 由于矩阵表示的映射在基上是恒等映射，它在整个定义域上必定都是恒等映射。
也可用另一种方法得到同一结论：考虑
          \begin{equation*}
            \colvec{x \\ y}=\colvec{y \\ x}_B
          \end{equation*}
它映到
          \begin{equation*}
            \colvec{(x+y)/2 \\ (x-y)/2}_D
          \end{equation*}
即表示 $\Re^2$ 中的以下元素。
          \begin{equation*}
            \frac{x+y}{2}\cdot\colvec[r]{1 \\ 1}
              +\frac{x-y}{2}\cdot\colvec[r]{1 \\ -1}
            =\colvec{x \\ y}
          \end{equation*}
      \end{exparts}
    
\end{ans}
\begin{ans}{三.III.2.17}
由于相对于标准基~$\stdbasis_2$，向量的表示就是它自身，只需计算矩阵与向量的乘积。
      \begin{exparts}
        \partsitem
          \begin{align*}
            \rep{h(\colvec{2 \\ 3})}{\stdbasis_2}
             &=\rep{h}{\stdbasis_2,\stdbasis_2}\;
             \rep{\colvec{2 \\ 3}}{\stdbasis_2}        \\
             &=\begin{mat}
              1  &3  \\
              2  &4
             \end{mat}_{\stdbasis_2,\stdbasis_2}
             \colvec{2 \\ 3}_{\stdbasis_2}                 \\
             &=\colvec{11 \\ 16}_{\stdbasis_2}=
             \colvec{11 \\ 16}
          \end{align*}
        \partsitem
           $\begin{mat}
              1  &3  \\
              2  &4
             \end{mat}
             \colvec{0 \\ 1}=
             \colvec{3 \\ 4}$
        \partsitem
           $\begin{mat}
              1  &3  \\
              2  &4
             \end{mat}
             \colvec{-1 \\ 1}=
             \colvec{2 \\ 2}$
      \end{exparts}
    
\end{ans}
\begin{ans}{三.III.2.18}
定义域中的一般元素相对于定义域的基表示为
      \begin{equation*}
        a\cos\theta+b\sin\theta=\colvec{a \\ a+b}_B
      \end{equation*}
它映到
      \begin{equation*}
        \colvec{0 \\ a}_D
          \quad\text{representing}\quad
        0\cdot(\cos\theta+\sin\theta)+a\cdot(\cos\theta)
      \end{equation*}
所以，该矩阵相对于这些基表示的线性映射
      \begin{equation*}
        a\cos\theta+b\sin\theta
            \mapsto
        a\cos\theta
      \end{equation*}
是向第一个分量的投影。
    
\end{ans}
\begin{ans}{三.III.2.19}
将给定的 $\polyspace_2$ 的基记为 $B$。线性映射的作用由矩阵与向量的乘法表示。
因此，$\stdbasis_3$ 中第一个向量映到 $\polyspace_2$ 中相对于~$B$ 具有以下表示的元素，
       \begin{equation*}
       \begin{mat}[r]
         1  &3  &0  \\
         0  &1  &0  \\
         1  &0  &1
       \end{mat}
       \colvec[r]{1 \\ 0 \\ 0}
       =
       \colvec[r]{1 \\ 0 \\ 1}
       \end{equation*}
该元素就是 $1+x$。同样计算另两个基向量的像。
       \begin{equation*}
         \begin{mat}[r]
           1  &3  &0  \\
           0  &1  &0  \\
           1  &0  &1
         \end{mat}
         \colvec[r]{0 \\ 1 \\ 0}
         =
         \colvec[r]{3 \\ 1 \\ 0}=\rep{4+x^2}{B}
         \quad
         \begin{mat}[r]
           1  &3  &0  \\
           0  &1  &0  \\
           1  &0  &1
         \end{mat}
         \colvec[r]{0 \\ 0 \\ 1}
         =
         \colvec[r]{0 \\ 0 \\ 1}=\rep{x}{B}
       \end{equation*}
所以，$h$ 的值域由三个多项式 $1+x$、$4+x^2$ 和 $x$ 张成。
于是，要判断 $1+2x$ 是否在值域中，只需寻找标量 $c_1$、$c_2$ 和 $c_3$，使得
       \begin{equation*}
         c_1\cdot(1+x)+c_2\cdot(4+x^2)+c_3\cdot(x)=1+2x
       \end{equation*}
显然，取 $c_1=1$、$c_2=0$、$c_3=1$ 即可。因此 $1+2x$ 在值域中，因为它是以下向量的像。
       \begin{equation*}
         1\cdot\colvec[r]{1 \\ 0 \\ 0}+0\cdot\colvec[r]{0 \\ 1 \\ 0}
             +1\cdot\colvec[r]{0 \\ 0 \\ 1}
       \end{equation*}

\textit{评注。}更简洁的论证是：矩阵非奇异，故秩为~$3$，值域维数也为~$3$；而陪域维数为~$3$，所以映射是满射。
因此，每个多项式都是某个向量的像，特别地，$1+2x$ 也是定义域中某个向量的像。
     
\end{ans}
\begin{ans}{三.III.2.20}
设 $B=\sequence{\vec{\beta}_1,\vec{\beta}_2}$。对于分量为 $x$ 和~$y$ 的一般向量 $\vec{v}$，可直接观察求出 $\rep{\vec{v}}{B}$。
          \begin{equation*}
            \colvec{x \\ y}=a\cdot\colvec{1 \\ 0}+b\cdot\colvec{1 \\ 1}
            \quad\text{gives}\quad
            b=y, a=x-y
          \end{equation*}
因此，一般向量相对于~$B$ 的表示如下。
          \begin{equation*}
            \rep{\colvec{x \\ y}}{B}=\colvec{x-y \\ y}
          \end{equation*}
利用矩阵与向量的乘法计算映射的作用。
          \begin{equation*}
            \begin{mat}
              2  &1  \\
             -1 &0
            \end{mat}_B
            \colvec{x-y \\ y}_B
            =
            \colvec{2(x-y)+y \\ -(x-y)}_B
            =
            \colvec{2x-y \\ -x+y}_B
          \end{equation*}
最后转回标准的向量表示。
          \begin{equation*}
            \colvec{2x-y \\ -x+y}_B
            =
            (2x-y)\cdot\colvec{1 \\ 0}+(-x+y)\cdot\colvec{1 \\ 1}
            =\colvec{x \\ -x+y}
          \end{equation*}
    
\end{ans}
\begin{ans}{三.III.2.21}
将矩阵记为 $G$，设它相对于基 $B$ 和 $D$ 表示 $\map{g}{V}{W}$。
由于 $G$ 有两列，定义域 $V$ 为二维；由于 $G$ 有两行，陪域 $W$ 为二维。
$g$ 对定义域一般元素的表示 $\rep{\vec{v}}{B,D}$ 的作用如下。
      \begin{equation*}
        \colvec{x \\ y}_B
         \;\mapsto\;
        \colvec{x+2y \\ 3x+6y}_D
      \end{equation*}
      \begin{exparts}
\partsitem 无论陪域的基~$D$ 如何，零向量 $\zero_{W}$ 的唯一表示都是
           \begin{equation*}
             \rep{\zero}{D}=\colvec[r]{0 \\ 0}_D
           \end{equation*}
因此，零空间元素的表示构成以下集合。
           \begin{equation*}
             \set{\colvec{x \\ y}_B\suchthat \text{$x+2y=0$ and $3x+6y=0$}}
             =\set{y\cdot\colvec[r]{-1/2 \\ 1}_D\suchthat y\in\Re}
           \end{equation*}
\partsitem 零度为~$1$。表示映射 $\map{\mbox{Rep}_D}{W}{\Re^2}$ 及其逆映射都是同构，因而保持子空间的维数。
上一问中的 $\Re^2$ 的子空间是一维的，所以它在表示映射的逆映射下的像\Dash $G$ 的零空间\Dash 也是一维的。
\partsitem 值域空间元素的表示构成以下集合。
           \begin{align*}
             \set{\colvec{x+2y \\ 3x+6y}_D\suchthat x,y\in\Re}
             &=\set{x\cdot\colvec{1 \\ 3}_D+y\cdot\colvec{2 \\ 6}_D\suchthat x,y\in\Re}  \\
             &=\set{k\cdot\colvec[r]{1 \\ 3}_D\suchthat k\in\Re}
           \end{align*}
\partsitem 当然，由 \nearbytheorem{th:RankMatEqRankMap}，映射的秩等于矩阵的秩，即一。
也可用前面处理零空间的同一论证，说明值域空间的维数为一。
\partsitem 一加一等于二。
      \end{exparts}
    
\end{ans}
\begin{ans}{三.III.2.22}
       \begin{exparts}
         \partsitem
(i)~定义域维数为列数~$m=2$，陪域维数为行数~$n=2$。

其余各问考虑以下矩阵与向量的方程。
           \begin{equation*}
             \begin{mat}
               2  &1  \\
              -1  &3
             \end{mat}
             \colvec{x \\ y}
             =
             \colvec{a \\ b}
             \tag{$*$}
           \end{equation*}
解出 $x$ 和~$y$。
           \begin{equation*}
             \begin{amat}{2}
             2  &1 &a  \\
             -1 &3 &b
             \end{amat}
             \grstep{(1/2)\rho_1+\rho_2}
             \grstep[(2/7)\rho_2]{(1/2)\rho_1}
             \grstep{-(1/2)\rho_2+\rho_1}
             \begin{amat}{2}
               1  &0 &(3/7)a-(1/7)b  \\
               0  &1 &(1/7)a+(2/7)b
             \end{amat}
           \end{equation*}
(ii)~对于式~($*$) 中的所有
           \begin{equation*}
             \colvec{a \\ b}\in\Re^2
           \end{equation*}
由计算可知，方程组均有解。所以值域空间就是整个陪域，$\rangespace{h}=\Re^2$。
映射的秩是值域维数，即~$2$。由于值域空间等于整个陪域，映射是满射。

(iii)~同样由计算可知，为求零空间，在式~($*$) 中令 $a=b=0$，得到 $x=y=0$。
零空间是定义域的平凡子空间。
           \begin{equation*}
             \nullspace{h}=\set{\colvec{0 \\ 0}}
           \end{equation*}
零度为零空间的维数，即~$0$。由于零空间平凡，映射是单射。
         \partsitem
(i)~定义域维数为矩阵列数~$m=3$，陪域维数为行数~$n=3$。

计算如下。
            \begin{multline*}
              \begin{amat}{3}
                0  &1  &3  &a \\
                2  &3  &4  &b \\
               -2  &-1 &2  &c
              \end{amat}
              \grstep{\rho_1\leftrightarrow\rho_2}
              \grstep{\rho_1+\rho_3}
              \grstep{-2\rho_2+\rho_3}                 \\
              \grstep{(1/2)\rho_1}
              \grstep{-(3/2)\rho_2+\rho_1}
              \begin{amat}{3}
                1  &0 &-5/2 &-(3/2)a+(1/2)b  \\
                0  &1 &3    &a  \\
                0  &0 &0    &-2a+b+c
              \end{amat}
            \end{multline*}
(ii)~陪域中存在三元组
            \begin{equation*}
              \colvec{a \\ b \\ c}\in\Re^3
            \end{equation*}
使方程组无解。具体地，只有当 $-2a+b+c=0$ 时方程组才有解。
            \begin{equation*}
              \rangespace{h}=\set{\colvec{a \\ b \\ c}\suchthat a=(b+c)/2}
                            =\set{\colvec{1/2 \\ 1 \\ 0}b
                                  +\colvec{1/2 \\ 0 \\ 1}c\suchthat b,c\in\Re}
            \end{equation*}
映射的秩是值域维数，即~$2$。由于值域空间不等于整个陪域，映射不是满射。

(iii)~在计算中令 $a=b=c=0$，得到无穷多个解。以自由变量~$z$ 为参数，零空间可描述如下。
            \begin{equation*}
              \nullspace{h}=\set{\colvec{x \\ y \\ z}\suchthat
                               \text{$y=-3z$ and $x=(5/2)z$}}
                           =\set{\colvec{5/2 \\ -3 \\ 1}z\suchthat z\in\Re}
            \end{equation*}
零度是零空间的维数，即~$1$。由于零空间非平凡，映射不是单射。
          \partsitem
(i)~定义域维数为 $m=2$，陪域维数为~$n=3$。计算如下。
              \begin{equation*}
                \begin{amat}{2}
                  1  &1  &a \\
                  2  &1  &b \\
                  3  &1  &c
                \end{amat}
                \grstep[-3\rho_1+\rho_3]{-2\rho_1+\rho_2}
                \grstep{-2\rho_2+\rho_3}
                \grstep{-\rho_2}
                \grstep{-\rho_2+\rho_1}
                \begin{amat}{2}
                  1  &0 &-a+b  \\
                  0  &1 &2a-b  \\
                  0  &0 &a-2b+c
                \end{amat}
              \end{equation*}

(ii)~值域是陪域的以下子空间。
            \begin{equation*}
              \rangespace{h}=\set{\colvec{2b-c \\ b \\ c}\suchthat b,c\in\Re}
                  =\set{\colvec{2 \\ 1 \\ 0}b+\colvec{-1 \\ 0 \\ 1}c\suchthat b,c\in\Re}
            \end{equation*}
秩为~$2$，映射不是满射。

(iii)~零空间是定义域的平凡子空间。
            \begin{equation*}
              \nullspace{h}=\set{\colvec{x \\ y}=\colvec{0 \\ 0}}
            \end{equation*}
零度为~$0$，映射是单射。
       \end{exparts}
     
\end{ans}
\begin{ans}{三.III.2.23}
高斯–若尔当消元如下。
          \begin{align*}
            \begin{amat}{3}
                1 &0 &-1 &a \\
                2 &1 &0  &b \\
                2 &2 &2  &c
            \end{amat}
            &\grstep[-2\rho_1+\rho_3]{-2\rho_1+\rho_2}
            \begin{amat}{3}
                1 &0 &-1 &a \\
                0 &1 &2  &-2a+b \\
                0 &2 &4  &-2a+c
            \end{amat}                                    \\
            &\grstep{-2\rho_2+\rho_3}
            \begin{amat}{3}
                1 &0 &-1 &a \\
                0 &1 &2  &-2a+b \\
                0 &0 &0  &2a-2b+c
            \end{amat}
          \end{align*}
(i)~维数为 $m=n=3$。
(ii)~值域空间由使这个方程组有解的所有陪域元素组成。
          \begin{equation*}
            \rangespace{h}=\set{\colvec{b-(1/2)c \\ b \\ c}\suchthat b,c\in\Re}
          \end{equation*}
秩为 2。由于秩小于陪域维数~$n=3$，映射不是满射。

(iii)~零空间由定义域中映到 $a=0$、$b=0$ 和~$c=0$ 的元素组成。
          \begin{equation*}
            \nullspace{h}=\set{\colvec{z \\ -2z \\ z}\suchthat z\in\Re}
          \end{equation*}
零度为~$1$。由于零度不为~$0$，映射不是单射。
    
\end{ans}
\begin{ans}{三.III.2.24}
该矩阵表示的任意映射，其定义域与陪域维数均为~$3$。
要证明映射是同构，必须证明它既是满射又是单射。
为此，无须用 $a$、$b$ 和~$c$ 扩充矩阵；以下计算
             \begin{multline*}
               \begin{mat}
                 2  &1  &0 \\
                 3  &1  &1 \\
                 7  &2  &1
               \end{mat}
               \grstep[-(7/2)\rho_1+\rho_3]{-(3/2)\rho_1+\rho_2}
               \grstep{-3\rho_2+\rho_3}
               \grstep[-2\rho_2 \\ -(1/2)\rho_3]{(1/2)\rho_1}
               \grstep{2\rho_3+\rho_2}                          \\
               \grstep{-(1/2)\rho_2+\rho_1}
               \begin{mat}
                 1  &0 &0  \\
                 0  &1 &0  \\
                 0  &0 &1
               \end{mat}
             \end{multline*}
说明每个陪域向量都恰好对应一个定义域向量。
    
\end{ans}
\begin{ans}{三.III.2.25}
      \begin{exparts}
        \partsitem
所定义的映射 $h$ 是满射，当且仅当对每个 $\vec{w}\in W$，都有 $\vec{v}\in V$ 使 $h(\vec{v})=\vec{w}$。
由于每个向量都恰好有一个表示，转为表示后可知，$h$ 是满射，当且仅当对每个表示 $\rep{\vec{w}}{D}$，都有表示 $\rep{\vec{v}}{B}$ 使
$H\cdot\rep{\vec{v}}{B}=\rep{\vec{w}}{D}$。
       \partsitem
与上一问类似。
       \partsitem
如本小节开头所述，由定义，矩阵 $H$ 所定义的映射 $h$ 将以下定义域向量 $\vec{v}$ 对应到陪域向量 $\vec{w}$。
          \begin{equation*}
            \rep{\vec{v}}{B}=\colvec{v_1 \\ \vdots \\ v_n}
            \qquad
            \rep{\vec{w}}{D}
               =
               H\cdot\rep{\vec{v}}{B}
               =\colvec{h_{1,1}v_1+\dots+h_{1,n}v_n  \\
                          \vdots \\
                       h_{m,1}v_1+\dots+h_{m,n}v_n}
          \end{equation*}
固定 $\vec{w}\in W$，考虑上述等式定义的线性方程组。
          \begin{equation*}
            % H\cdot\rep{\vec{v}}{B}=\rep{\vec{w}}{D}
            % \qquad
            \begin{linsys}{3}
              h_{1,1}v_1 &+ &\cdots &+ &h_{1,n}v_n &= &w_1  \\
              h_{2,1}v_1 &+ &\cdots &+ &h_{2,n}v_n &= &w_2  \\
                       &  &        & &          &\vdotswithin{=} & \\
              h_{n,1}v_1 &+ &\cdots &+ &h_{n,n}v_n &= &w_n
            \end{linsys}
          \end{equation*}
（同样，这里的 $w_i$ 是给定值，$v_j$ 是未知数。）
$H$ 非奇异，当且仅当对所有 $w_1$、\ldots、$w_n$，此方程组都有唯一解。
由本题前两问，这等价于映射 $h$ 既是满射又是单射，也就等价于 $h$ 是同构。
       \end{exparts}
      
\end{ans}
\begin{ans}{三.III.2.26}
不一定，值域空间可以不同。\nearbyexample{ex:CngBasesChgMap} 就说明了这一点。
    
\end{ans}
\begin{ans}{三.III.2.27}
回顾，表示映射
      \begin{equation*}
        V\mapsunder{\text{Rep}_{B}}\Re^n
      \end{equation*}
是同构。因此其逆映射 $\map{\mbox{Rep}_B^{-1}}{\Re^n}{V}$ 也是同构。
所求的 $\Re^n$ 上的变换就是以下复合。
      \begin{equation*}
        \Re^n\mapsunder{\text{Rep}_{B}^{-1}}
        V\mapsunder{\text{Rep}_{D}}\Re^n
      \end{equation*}
由于同构的复合仍为同构，映射 $\composed{\mbox{Rep}_{D}}{\mbox{Rep}_{B}^{-1}}$ 是同构。
    
\end{ans}
\begin{ans}{三.III.2.28}
存在。考虑
      \begin{equation*}
        H=\begin{mat}[r]
            1  &0  \\
            0  &1
          \end{mat}
      \end{equation*}
表示从 \( \Re^2 \) 到 \( \Re^2 \) 的映射。相对于标准基
\( B_1=\stdbasis_2, D_1=\stdbasis_2 \)，该矩阵表示恒等映射。相对于
      \begin{equation*}
        B_2=D_2=\sequence{\colvec[r]{1 \\ 1},\colvec[r]{1 \\ -1}}
      \end{equation*}
它仍表示恒等映射。事实上，只要起始基与终止基相同\Dash 即 \( B_i=D_i \)\Dash $H$ 表示的映射就是恒等映射。
    
\end{ans}
\begin{ans}{三.III.2.29}
由 \nearbylemma{le:NonsingMatIffNonsingMap} 立即得到。
    
\end{ans}
\begin{ans}{三.III.2.30}
第一个映射
      \begin{equation*}
        \colvec{x \\ y}=\colvec{x \\ y}_{\stdbasis_2}
        \mapsto
        \colvec{3x \\ 2y}_{\stdbasis_2}=\colvec{3x \\ 2y}
      \end{equation*}
将向量在 \( x \) 方向伸长为原来的三倍，在 \( y \) 方向伸长为原来的两倍。
第二个映射
      \begin{equation*}
        \colvec{x \\ y}=\colvec{x \\ y}_{\stdbasis_2}
        \mapsto
        \colvec{x \\ 0}_{\stdbasis_2}=\colvec{x \\ 0}
      \end{equation*}
将向量投影到 \( x \) 轴。第三个映射
      \begin{equation*}
        \colvec{x \\ y}=\colvec{x \\ y}_{\stdbasis_2}
        \mapsto
        \colvec{y \\ x}_{\stdbasis_2}=\colvec{y \\ x}
      \end{equation*}
交换第一、第二个分量（即关于直线 \( y=x \) 的反射）。最后一个映射
      \begin{equation*}
        \colvec{x \\ y}=\colvec{x \\ y}_{\stdbasis_2}
        \mapsto
        \colvec{x+3y \\ y}_{\stdbasis_2}=\colvec{x+3y \\ y}
      \end{equation*}
沿平行于 \( y \) 轴的方向移动向量端点，移动量为端点到该轴距离的三倍（这是一个\definend{错切}）。
     
\end{ans}
\begin{ans}{三.III.2.31}
       \begin{exparts}
\partsitem 由 \nearbytheorem{th:RankMatEqRankMap} 立即得到。
\partsitem 存在。这由上一问立即得到。

给出一个具体例子：先取 $\stdbasis_3$ 为定义域的基，再寻找陪域 $\Re^3$ 的基 $D$。
矩阵 $H$ 给出映射的作用如下，
            \begin{multline*}
              \colvec[r]{1 \\ 0 \\ 0}=\colvec[r]{1 \\ 0 \\ 0}_{\stdbasis_3}
                 \mapsto\colvec[r]{1 \\ 2 \\ 0}_D
              \quad
              \colvec[r]{0 \\ 1 \\ 0}=\colvec[r]{0 \\ 1 \\ 0}_{\stdbasis_3}
                 \mapsto\colvec[r]{0 \\ 0 \\ 1}_D
              \\  %\quad
              \colvec[r]{0 \\ 0 \\ 1}=\colvec[r]{0 \\ 0 \\ 1}_{\stdbasis_3}
                 \mapsto\colvec[r]{0 \\ 0 \\ 0}_D
            \end{multline*}
可以寻找一组基 $D$ 使得
            \begin{equation*}
              \rep{\colvec[r]{1 \\ 0 \\ 0}}{D}=\colvec[r]{1 \\ 2 \\ 0}_D
              \quad\mbox{and}\quad
              \rep{\colvec[r]{0 \\ 1 \\ 0}}{D}=\colvec[r]{0 \\ 0 \\ 1}_D
            \end{equation*}
即使 $H$ 相对于 $\stdbasis_3,D$ 表示向 $xy$ 平面的投影。
第二个条件说明 $D$ 的第三个元素为 $\vec{e}_2$。
第一个条件说明 $D$ 的第一个元素加上第二个元素的两倍等于 $\vec{e}_1$，所以以下基符合要求。
            \begin{equation*}
              D=\sequence{\colvec{2 \\ 0 \\ 1},
                          \colvec{-1/2 \\ 0 \\ -1/2},
                          \colvec{0 \\ 1 \\ 0}}
            \end{equation*}
% sage: load('gauss_method.sage')
% sage: M = matrix(QQ, [[2,-1/2,0], [0,0,1], [1,-1/2,0]])
% sage: gauss_method(M)
% [   2 -1/2    0]
% [   0    0    1]
% [   1 -1/2    0]
% (' take', -1/2, 'times row', 1, 'plus row', 3)
% [   2 -1/2    0]
% [   0    0    1]
% [   0 -1/4    0]
% (' swap row', 2, 'with row', 3)
% [   2 -1/2    0]
% [   0 -1/4    0]
% [   0    0    1]
       \end{exparts}
     
\end{ans}
\begin{ans}{三.III.2.32}
      \begin{exparts}
\partsitem 回顾，表示映射 $\map{\mbox{Rep}_{B}}{V}{\Re^n}$ 是线性的（实际上是同构，但这里不需要单射性或满射性）。
将列向量 $x$ 看成 $\nbym{n}{1}$ 矩阵，可知从 $\Re^n$ 到 $\Re$、将列向量映到它与 $\vec{x}$ 的点积的映射是线性的（这是矩阵与向量的乘积，故可用 \nearbytheorem{th:MatIsLinMap}）。
因此，所考察的 $h_{\vec{x}}$ 是两个线性映射的复合，也是线性的。
          \begin{equation*}
            \vec{v}\mapsto \rep{\vec{v}}{B}
                   \mapsto \vec{x}\cdot\rep{\vec{v}}{B}
          \end{equation*}
\partsitem 任意线性映射 $\map{g}{V}{\Re}$ 都由某个矩阵表示，
          \begin{equation*}
            \begin{mat}
              g_1  &g_2 &\cdots &g_n
            \end{mat}
          \end{equation*}
（由于 $V$ 为 $n$ 维，该矩阵有 $n$ 列；由于 $\Re$ 为一维，它只有一行）。
取 $\vec{x}$ 为这个矩阵转置所得的列向量，
          \begin{equation*}
            \vec{x}=\colvec{g_1 \\ \vdots \\ g_n}
          \end{equation*}
就得到所需的作用。
          \begin{equation*}
            \vec{v}=\colvec{v_1 \\ \vdots \\ v_n}
             \mapsto
            \colvec{g_1 \\ \vdots \\ g_n}\dotprod\colvec{v_1 \\ \vdots \\ v_n}
            =g_1v_1+\dots+g_nv_n
          \end{equation*}
\partsitem 不能。如果 \( \vec{x} \) 有非零分量，则 \( h_{\vec{x}} \) 不可能是零映射；如果 \( \vec{x} \) 是零向量，则 \( h_{\vec{x}} \) 只能是零映射。
      \end{exparts}
    
\end{ans}
\begin{ans}{三.III.2.33}
参见下一节。
    
\end{ans}
\protect \section {第 IV 节：矩阵运算}
\protect \subsection {三.IV.1: 和与数乘}
\begin{ans}{三.IV.1.8}
      \begin{exparts*}
        \partsitem \( \begin{mat}[r]
                   7  &0  &6  \\
                   9  &1  &6
                 \end{mat}  \)
        \partsitem \( \begin{mat}[r]
                  12  &-6 &-6 \\
                   6  &12 &18
                 \end{mat}  \)
        \partsitem \( \begin{mat}[r]
                   4  &2  \\
                   0  &6
                 \end{mat}  \)
        \partsitem \( \begin{mat}[r]
                  -1  &28 \\
                   2  &1
                 \end{mat}  \)
        \partsitem Not defined.
      \end{exparts*}
   
\end{ans}
\begin{ans}{三.IV.1.9}
基的选择无关紧要，只需使大小正确。
     \begin{equation*}
       \begin{mat}
         0 &0 &0 &0 \\
         0 &0 &0 &0
       \end{mat}
     \end{equation*}
   
\end{ans}
\begin{ans}{三.IV.1.10}
照常相对于基 $B,D$ 表示定义域向量 $\vec{v}\in V$ 及映射 $\map{g,h}{V}{W}$。
      \begin{exparts}
\partsitem $(g+h)\,(\vec{v})=g(\vec{v})+h(\vec{v})$ 的表示
        \begin{multline*}
          \bigl( (g_{1,1}v_1+\dots+g_{1,n}v_n)\vec{\delta}_1
                  +\cdots
                  +(g_{m,1}v_1+\dots+g_{m,n}v_n)\vec{\delta}_m \bigr) \\
           +\bigl( (h_{1,1}v_1+\cdots+h_{1,n}v_n)\vec{\delta}_1
                   +\cdots
                   +(h_{m,1}v_1+\dots+h_{m,n}v_n)\vec{\delta}_m \bigr)
         \end{multline*}
重新组合为
         \begin{multline*}
          =((g_{1,1}+h_{1,1})v_1+\dots+
                 (g_{1,1}+h_{1,n})v_n)\cdot\vec{\delta}_1              \\
            +\cdots
           +((g_{m,1}+h_{m,1})v_1+\dots+
                 (g_{m,n}+h_{m,n})v_n)\cdot\vec{\delta}_m
        \end{multline*}
即 \( g(\vec{v}) \) 与 \( h(\vec{v}) \) 的表示逐项相加。
\partsitem $(r\cdot h)\,(\vec{v})=r\cdot \bigl(h(\vec{v})\bigr)$ 的表示
        \begin{align*}
          r\cdot
             &\bigl(
               (h_{1,1}v_1+h_{1,2}v_2+\dots+h_{1,n}v_n)\vec{\delta}_1   \\
             &\quad +\dots
             +(h_{m,1}v_1+h_{m,2}v_2+\dots+h_{m,n}v_n)\vec{\delta}_m\bigr)   \\
             &=(rh_{1,1}v_1+\dots+rh_{1,n}v_n)\cdot\vec{\delta}_1         \\
             &\quad+\dots
              +(rh_{m,1}v_1+\dots+rh_{m,n}v_n)\cdot\vec{\delta}_m
        \end{align*}
由 \( r \) 与 \( h \) 的表示逐项相乘得到。
     \end{exparts}
   
\end{ans}
\begin{ans}{三.IV.1.11}
首先，各性质都容易逐项验证。例如，写成
      \begin{equation*}
        G=
        \begin{mat}
          g_{1,1}  &\ldots  &g_{1,n}  \\
          \vdots   &        &\vdots   \\
          g_{m,1}  &\ldots  &g_{m,n}
        \end{mat}
        \qquad
        H=
        \begin{mat}
          h_{1,1}  &\ldots  &h_{1,n}  \\
          \vdots   &        &\vdots   \\
          h_{m,1}  &\ldots  &h_{m,n}
        \end{mat}
      \end{equation*}
由定义有
      \begin{equation*}
        G+H=
        \begin{mat}
          g_{1,1}+h_{1,1}  &\ldots  &g_{1,n}+h_{1,n}  \\
          \vdots           &        &\vdots           \\
          g_{m,1}+h_{m,1}  &\ldots  &g_{m,n}+h_{m,n}
        \end{mat}
      \end{equation*}
以及
      \begin{equation*}
        H+G=
        \begin{mat}
          h_{1,1}+g_{1,1}  &\ldots  &h_{1,n}+g_{1,n}  \\
          \vdots           &        &\vdots           \\
          h_{m,1}+g_{m,1}  &\ldots  &h_{m,n}+g_{m,n}
        \end{mat}
      \end{equation*}
两者相等，因为对应元素相等：$g_{i,j}+h_{i,j}=h_{i,j}+g_{i,j}$。
也就是说，仅使用 \nearbydefinition{def:SumScalarProdMats} 就容易验证各式。

不过，结合定义与 \nearbytheorem{th:MatOpsRepMapOps}，从所表示的映射出发也容易理解每个性质。
      \begin{exparts}
\partsitem 映射 $g+h$ 与 $h+g$ 相等，因为任意向量空间中的加法都满足交换律，故
$g(\vec{v})+h(\vec{v})=h(\vec{v})+g(\vec{v})$。映射相同，表示必定相同。
\partsitem 与上一问相同，只是改用向量空间加法的结合律。
\partsitem 与前面相同，注意 $g(\vec{v})+z(\vec{v})=g(\vec{v})+\zero=g(\vec{v})$。
\partsitem 利用 $0\cdot g(\vec{v})=\zero=z(\vec{v})$。
\partsitem 利用 $(r+s)\cdot g(\vec{v})=r\cdot g(\vec{v})+s\cdot g(\vec{v})$。
\partsitem 在前两问中取 $r=1$、$s=-1$。
\partsitem 利用 $r\cdot (g(\vec{v})+h(\vec{v}))=r\cdot g(\vec{v})+r\cdot h(\vec{v})$。
\partsitem 利用 $(rs)\cdot g(\vec{v})=r\cdot(s\cdot g(\vec{v}))$。
      \end{exparts}
    
\end{ans}
\begin{ans}{三.IV.1.12}
对任意具有基 \( B,D \) 的 \( V,W \)，大小适当的零矩阵表示以下映射。
      \begin{equation*}
        \vec{\beta}_1\mapsto 0\cdot\vec{\delta}_1+\dots+0\cdot\vec{\delta}_m
        \quad\cdots\quad
        \vec{\beta}_n\mapsto 0\cdot\vec{\delta}_1+\dots+0\cdot\vec{\delta}_m
      \end{equation*}
这就是零映射。

不存在其他只表示一个映射的矩阵。设 \( H \) 不是零矩阵，则它有非零元素，不妨设 \( h_{i,j}\neq 0 \)。
相对于基 $B,D$，它表示 \( \map{h_1}{V}{W} \)，其作用为
      \begin{equation*}
        \vec{\beta}_j\mapsto
        h_{1,j}\vec{\delta}_1+\dots+h_{i,j}\vec{\delta}_i
          +\dots+h_{m,j}\vec{\delta}_m
      \end{equation*}
相对于 $B,2\cdot D$，它又表示 \( \map{h_2}{V}{W} \)，其作用为
      \begin{equation*}
        \vec{\beta}_j\mapsto
        h_{1,j}\cdot(2\vec{\delta}_1)+\dots+h_{i,j}\cdot(2\vec{\delta}_i)
          +\dots+h_{m,j}\cdot(2\vec{\delta}_m)
      \end{equation*}
（记号 $2\cdot D$ 表示将 $D$ 中所有元素加倍）。容易看出这两个映射不同。
     
\end{ans}
\begin{ans}{三.IV.1.13}
固定 \( V \) 和 \( W \) 的基 \( B \) 和 \( D \)，考虑映射
\( \map{\mbox{Rep}_{B,D}}{\linmaps{V}{W}}{\matspace_{\nbym{m}{n}}} \)，将每个线性映射对应到其表示矩阵：$h\mapsto \rep{h}{B,D}$。
由上一节可知，固定基之后，矩阵与线性映射一一对应，所以表示映射既是单射又是满射。
它保持线性运算由 \nearbytheorem{th:MatOpsRepMapOps} 保证。
    
\end{ans}
\begin{ans}{三.IV.1.14}
固定基，将这些变换用 \( \nbyn{2} \) 矩阵表示。矩阵空间 \( \matspace_{\nbyn{2}} \) 的维数为四，所以任意六个元素必线性相关。
由上一题，这给出映射之间的线性相关关系。
（误导性只在于给了六个变换而非五个，比保证线性相关所需的更多。）
    
\end{ans}
\begin{ans}{三.IV.1.15}
和的迹等于迹的和，因为 \( \text{trace}(H+G) \) 与 \( \text{trace}(H)+\text{trace}(G) \)
都是 \( h_{1,1}+g_{1,1} \)、\( h_{2,2}+g_{2,2} \) 等的和。
对于数乘，有 \( \mbox{trace}(r\cdot H)=r\cdot\mbox{trace}(H) \)，证明很简单。
因此，迹映射是从 $\matspace_{\nbyn{n}}$ 到 $\Re$ 的同态。
    
\end{ans}
\begin{ans}{三.IV.1.16}
      \begin{exparts}
\partsitem \( \trans{(G+H)} \) 的 \( i,j \) 元素为 \( g_{j,i}+h_{j,i} \)，也就是 \( \trans{G}+\trans{H} \) 的 \( i,j \) 元素。
\partsitem \( \trans{(r\cdot H)} \) 的 \( i,j \) 元素为 \( rh_{j,i} \)，也就是 \( r\cdot\trans{H} \) 的 \( i,j \) 元素。
      \end{exparts}
   
\end{ans}
\begin{ans}{三.IV.1.17}
      \begin{exparts}
\partsitem \( H+\trans{H} \) 的 \( i,j \) 元素为 \( h_{i,j}+h_{j,i} \)，\( j,i \) 元素为 \( h_{j,i}+h_{i,j} \)。两者相等，因此 \( H+\trans{H} \) 对称。

每个对称矩阵确实都有这种形式，因为可写成 \( H=(1/2)\cdot(H+\trans{H}) \)。
\partsitem 对称矩阵的集合包含零矩阵，所以非空。
显然，对称矩阵的标量倍数仍对称。两个对称矩阵的和 \( H+G \) 也对称，因为
\( h_{i,j}+g_{i,j}=h_{j,i}+g_{j,i} \)（由 \( h_{i,j}=h_{j,i} \) 和 \( g_{i,j}=g_{j,i} \)）。
因此，该子集非空且对继承的运算封闭，故为子空间。
      \end{exparts}
    
\end{ans}
\begin{ans}{三.IV.1.18}
      \begin{exparts}
\partsitem 除乘以零会使秩变为零以外，数乘不改变矩阵的秩。
（这由本书第一个定理推出：一行乘以非零标量，不改变相应线性方程组的解集。）
\partsitem 秩为 \( n \) 的矩阵的和，秩可以小于 \( n \)。例如，对任意矩阵 \( H \)，和 \( H+(-1)\cdot H \) 的秩为零。

秩为 \( n \) 的矩阵的和，秩也可以大于 \( n \)。以下两个秩为一的矩阵相加，得到秩为二的矩阵。
          \begin{equation*}
            \begin{mat}[r]
              1  &0  \\
              0  &0
            \end{mat}
            +\begin{mat}[r]
              0  &0  \\
              0  &1
            \end{mat}
            =\begin{mat}[r]
              1  &0  \\
              0  &1
            \end{mat}
          \end{equation*}
      \end{exparts}
    
\end{ans}
\protect \subsection {三.IV.2: 矩阵乘法}
\begin{ans}{三.IV.2.14}
      \begin{exparts*}
        \partsitem
          $\begin{mat}[r]
            0  &15.5  \\
            0  &-19
          \end{mat}$
        \partsitem \(
          \begin{mat}[r]
            2  &-1  &-1  \\
           17  &-1  &-1
          \end{mat} \)
        \partsitem Not defined.
        \partsitem \( \begin{mat}[r]
                   1  &0  \\
                   0  &1
                 \end{mat}  \)
      \end{exparts*}
    
\end{ans}
\begin{ans}{三.IV.2.15}
      \begin{exparts*}
        \partsitem \(
          \begin{mat}[r]
            1  &-2   \\
            10 &4
          \end{mat}  \)
        \partsitem \(
          \begin{mat}[r]
            1  &-2   \\
            10 &4
          \end{mat}
          \begin{mat}[r]
            -2 &3    \\
            -4 &1
          \end{mat}=
          \begin{mat}[r]
            6  &1    \\
           -36 &34
          \end{mat}  \)
        \partsitem \(
          \begin{mat}[r]
           -18 &17   \\
           -24 &16
          \end{mat}  \)
        \partsitem \(
          \begin{mat}[r]
            1  &-1  \\
            2  &0
          \end{mat}
          \begin{mat}[r]
           -18 &17   \\
           -24 &16
          \end{mat}=
          \begin{mat}[r]
            6  &1    \\
           -36 &34
          \end{mat}  \)
      \end{exparts*}
    
\end{ans}
\begin{ans}{三.IV.2.16}
      \begin{exparts*}
\partsitem 有。
\partsitem 有。
\partsitem 无。
\partsitem 无。
      \end{exparts*}
    
\end{ans}
\begin{ans}{三.IV.2.17}
      \begin{exparts*}
\partsitem \( \nbym{2}{1} \)
\partsitem \( \nbym{1}{1} \)
\partsitem 没有定义。
\partsitem \( \nbym{2}{2} \)
      \end{exparts*}
     
\end{ans}
\begin{ans}{三.IV.2.18}
有
      \begin{equation*}
        \begin{linsys}{3}
            h_{1,1}\cdot(g_{1,1}y_1+g_{1,2}y_2)
               &+  &h_{1,2}\cdot (g_{2,1}y_1+g_{2,2}y_2)
               &+  &h_{1,3}\cdot (g_{3,1}y_1+g_{3,2}y_2)
               &=  &d_1  \\
            h_{2,1}\cdot(g_{1,1}y_1+g_{1,2}y_2)
               &+  &h_{2,2}\cdot (g_{2,1}y_1+g_{2,2}y_2)
               &+  &h_{2,3}\cdot (g_{3,1}y_1+g_{3,2}y_2)
               &=  &d_2
         \end{linsys}
      \end{equation*}
展开并按各个 $y$ 重新组合，得到
      \begin{equation*}
        \begin{linsys}{2}
            (h_{1,1}g_{1,1}+h_{1,2}g_{2,1}+h_{1,3}g_{3,1})y_1
              &+  &(h_{1,1}g_{1,2}+h_{1,2}g_{2,2}+h_{1,3}g_{3,2})y_2
               &=  &d_1  \\
            (h_{2,1}g_{1,1}+h_{2,2}g_{2,1}+h_{2,3}g_{3,1})y_1
              &+  &(h_{2,1}g_{1,2}+h_{2,2}g_{2,2}+h_{2,3}g_{3,2})y_2
               &=  &d_2
         \end{linsys}
      \end{equation*}
可用矩阵语言表示原方程组和用于代换的方程组，分别为
      \begin{equation*}
        \begin{mat}
          h_{1,1}  &h_{1,2}  &h_{1,3}  \\
          h_{2,1}  &h_{2,2}  &h_{2,3}
        \end{mat}
        \colvec{x_1 \\ x_2 \\ x_3}
        =H\colvec{x_1 \\ x_2 \\ x_3}=\colvec{d_1 \\ d_2}
      \end{equation*}
以及
      \begin{equation*}
        \begin{mat}
          g_{1,1}  &g_{1,2}  \\
          g_{2,1}  &g_{2,2}  \\
          g_{3,1}  &g_{3,2}
        \end{mat}
        \colvec{y_1 \\ y_2}
        =G\colvec{y_1 \\ y_2}=\colvec{x_1 \\ x_2 \\ x_3}
      \end{equation*}
于是代换可写成 $\vec{d}=H\vec{x}=H(G\vec{y})=(HG)\vec{y}$。
    
\end{ans}
\begin{ans}{三.IV.2.19}
      \begin{exparts}
        \partsitem
根据定义得到
          \begin{align*}
            \colvec{a \\ b \\ c}
            &\mapsto
            (a+b)x^2+(2a+2b)x+c                   \\
            &\mapsto
            \begin{mat}
              a+b &(a+b)-2(2a+2b) \\
              2a+2b &0
            \end{mat}
            =
            \begin{mat}
              a+b   &-3a-3b  \\
              2a+2b &0
            \end{mat}
          \end{align*}
        \partsitem
因为
          \begin{equation*}
            \colvec{1 \\ 1 \\ 1}\mapsto 2x^2+4x+1
            \quad
            \colvec{0 \\ 1 \\ 1}\mapsto x^2+2x+1
            \quad
            \colvec{0 \\ 0 \\ 1}\mapsto 0x^2+0x+1
          \end{equation*}
得到~$h$ 的以下表示。
          \begin{equation*}
            \rep{h}{B,C}
            =
            \begin{mat}
              5/2 &3/2  &1/2  \\
             -3/2 &-1/2 &1/2  \\
              2   &1    &0
            \end{mat}
          \end{equation*}
类似地，因为
          \begin{equation*}
            1+x\mapsto
            \begin{mat}
              0 &-2 \\
              1 &0
            \end{mat}
            \quad
            1-x\mapsto
            \begin{mat}
              0 &2 \\
              -1 &0
            \end{mat}
            \quad
            x^2\mapsto
            \begin{mat}
              1 &1 \\
              0 &0
            \end{mat}
          \end{equation*}
得到~$g$ 的以下表示。
          \begin{equation*}
            \rep{g}{C,D}=
            \begin{mat}
              0  &0    &1  \\
             -1  &1    &1/2 \\
             1/3 &-1/3 &0   \\
              0  &0    &0
            \end{mat}
          \end{equation*}
        \partsitem
$\composed{g}{h}$ 在定义域基上的作用如下。
          \begin{equation*}
            \colvec{1 \\ 1 \\ 1}\mapsto
              \begin{mat}
                2 &-6 \\
                4 &0
              \end{mat}
            \quad
            \colvec{0 \\ 1 \\ 1}\mapsto
              \begin{mat}
                1 &-3 \\
                2 &0
              \end{mat}
            \quad
            \colvec{0 \\ 0 \\ 1}\mapsto
              \begin{mat}
                0 &0 \\
                0 &0
              \end{mat}
          \end{equation*}
于是得到
          \begin{equation*}
            \rep{\composed{g}{h}}{B,D}=
            \begin{mat}
               2   &1    &0  \\
              -3   &-3/2 &0  \\
             4/3   &2/3  &0  \\
               0   &0    &0
            \end{mat}
          \end{equation*}
        \partsitem
矩阵乘法按常规计算即可，只需注意顺序。
          \begin{equation*}
            \begin{mat}
              0  &0    &1  \\
             -1  &1    &1/2 \\
             1/3 &-1/3 &0   \\
              0  &0    &0
            \end{mat}
            \begin{mat}
              5/2 &3/2  &1/2  \\
             -3/2 &-1/2 &1/2  \\
              2   &1    &0
            \end{mat}
            =
            \begin{mat}
               2   &1    &0  \\
              -3   &-3/2 &0  \\
             4/3   &2/3  &0  \\
               0   &0    &0
           \end{mat}
          \end{equation*}
      \end{exparts}
    
\end{ans}
\begin{ans}{三.IV.2.20}
严格说，不同。点积得到标量，矩阵乘积得到 $\nbyn{1}$ 矩阵。不过通常会忽略这个区别。
    
\end{ans}
\begin{ans}{三.IV.2.21}
$d/dx$ 在 $B$ 上的作用为 $1\mapsto0$、$x\mapsto 1$、$x^2\mapsto 2x$、\ldots，因此其 $\nbyn{(n+1)}$ 矩阵表示如下。
      \begin{equation*}
        \rep{\frac{d}{dx}}{B,B}=
        \begin{mat}
          0  &1 &0  &  &0  \\
          0  &0 &2  &  &0  \\
             &  &   &\ddots\\
          0  &0 &0  &  &n  \\
          0  &0 &0  &  &0
        \end{mat}
      \end{equation*}
这个矩阵是方阵，所以与自身的乘积有定义。
      \begin{equation*}
        \begin{mat}
          0  &1 &0  &  &0  \\
          0  &0 &2  &  &0  \\
             &  &   &\ddots\\
          0  &0 &0  &  &n  \\
          0  &0 &0  &  &0
        \end{mat}^2
        =
        \begin{mat}
          0  &0 &2  &0  &        &0       \\
          0  &0 &0  &6  &        &0       \\
             &  &   &   &\ddots  &        \\
          0  &0 &0  &  &         &n(n-1)  \\
          0  &0 &0  &  &         &0       \\
          0  &0 &0  &  &         &0
        \end{mat}
      \end{equation*}
所得矩阵表示复合
      \begin{equation*}
        p\;\mapsunder{\frac{d}{dx}}\;\frac{d\,p}{dx}
         \;\mapsunder{\frac{d}{dx}}\;\frac{d^2\,p}{dx^2}
      \end{equation*}
即二阶求导运算。
     
\end{ans}
\begin{ans}{三.IV.2.22}
      \begin{exparts}
\item ii
\item iii
\item 无
\item 无；若也允许从左边乘，则为 (i)。
      \end{exparts}
    
\end{ans}
\begin{ans}{三.IV.2.23}
对所有一维空间都成立。设 $f$、$g$ 是一维空间上的变换。需证明对所有向量都有
$\composed{g}{f}\,(\vec{v})=\composed{f}{g}\,(\vec{v})$。
固定空间的一组基 $B$，则这些变换由 $\nbyn{1}$ 矩阵表示。
      \begin{equation*}
        F=\rep{f}{B,B}=\begin{mat}
                        f_{1,1}
                     \end{mat}
        \qquad
        G=\rep{g}{B,B}=\begin{mat}
                        g_{1,1}
                     \end{mat}
      \end{equation*}
所以复合分别由 $GF$ 和 $FG$ 表示。
      \begin{equation*}
        GF=\rep{\composed{g}{f}}{B,B}=\begin{mat}
                        g_{1,1}f_{1,1}
                     \end{mat}
        \qquad
        FG=\rep{\composed{f}{g}}{B,B}=\begin{mat}
                        f_{1,1}g_{1,1}
                     \end{mat}
      \end{equation*}
两个矩阵相等，因此两个复合在空间中每个向量上的作用相同。
     
\end{ans}
\begin{ans}{三.IV.2.24}
那样就无法表示线性映射的复合，\nearbytheorem{th:MatMultRepComp} 将不成立。
    
\end{ans}
\begin{ans}{三.IV.2.25}
各式都容易由对应的映射性质推出。例如，先作用 $q$ 次变换 $h$，再作用 $p$ 次，就是总共作用 $p+q$ 次。
    
\end{ans}
\begin{ans}{三.IV.2.26}
虽然可以逐项计算，但从所表示的映射理解更好。固定空间和基，使这些矩阵表示线性映射 $f,g,h$。
       \begin{exparts}
\partsitem 均成立，因为
$r\cdot (\composed{g}{h})\,(\vec{v})=r\cdot g(\,h(\vec{v})\,)=\composed{(r\cdot g)}{h}\,(\vec{v})$，且
$\composed{g}{(r\cdot h)}\,(\vec{v})=g(\,r\cdot h(\vec{v})\,)=r\cdot g(h(\vec{v}))=r\cdot (\composed{g}{h})\,(\vec{v})$
（第二个等号由 $g$ 的线性性得到）。
\partsitem 均成立。首先，$\composed{f}{(rg+sh)}$ 与
$r\cdot(\composed{f}{g})+s\cdot(\composed{f}{h})$ 都将 $\vec{v}$ 映到
$r\cdot f(g(\vec{v}))+s\cdot f(h(\vec{v}))$；计算与上一问相同（前者使用 $f$ 的线性性）。
其次，$\composed{(rf+sg)}{h}$ 与 $r\cdot(\composed{f}{h})+s\cdot(\composed{g}{h})$ 都将
$\vec{v}$ 映到 $r\cdot f(h(\vec{v}))+s\cdot g(h(\vec{v}))$。
      \end{exparts}
    
\end{ans}
\begin{ans}{三.IV.2.27}
尚未从映射角度解释转置运算，因此通过考察元素来验证。
      \begin{exparts}
\partsitem \( \trans{GH} \) 的 \( i,j \) 元素是 $GH$ 的 $j,i$ 元素，即 \( G \) 的第 \( j \) 行与 \( H \) 的第 \( i \) 列的点积。
\( \trans{H}\trans{G} \) 的 \( i,j \) 元素是 \( \trans{H} \) 的第 \( i \) 行与 \( \trans{G} \) 的第 \( j \) 列的点积，
即 \( H \) 的第 \( i \) 列与 \( G \) 的第 \( j \) 行的点积。点积满足交换律，所以两者相等。
\partsitem 由上一问，两者都等于自身的转置，例如
$\trans{(H\trans{H})}=\trans{\trans{H}}\trans{H}=H\trans{H}$。
      \end{exparts}
    
\end{ans}
\begin{ans}{三.IV.2.28}
      Consider \( \map{r_x,r_y}{\Re^3}{\Re^3} \) rotating all vectors
      \( \pi/2 \)~radians
      绕 \( x \) 轴和 \( y \) 轴逆时针旋转
      (counterclockwise in the sense that a person whose head is at
      \( \vec{e}_1 \) or \( \vec{e}_2 \) and whose feet are at the origin
      从该方向朝向原点观察时，旋转看起来是
      逆时针）。
      \begin{center}  \small
        \includegraphics{map/mp/ch3.78}
        \qquad
        \includegraphics{map/mp/ch3.79}
      \end{center}
先作 $r_x$ 再作 $r_y$，与先作 $r_y$ 再作 $r_x$ 不同。
具体地，$r_x(\vec{e}_3)=-\vec{e}_2$，所以 $\composed{r_y}{r_x}(\vec{e}_3)=-\vec{e}_2$；
而 $r_y(\vec{e}_3)=\vec{e}_1$，所以 $\composed{r_x}{r_y}(\vec{e}_3)=\vec{e}_1$。
因此，这些映射不可交换。
    
\end{ans}
\begin{ans}{三.IV.2.29}
只要空间维数适当，具体选择无关紧要。

对于结合律，设 $F$ 为 $\nbym{m}{r}$，$G$ 为 $\nbym{r}{n}$，$H$ 为 $\nbym{n}{k}$。
可取任意 $r$ 维、$m$ 维、$n$ 维、$k$ 维空间\Dash 例如 $\Re^r$、$\Re^m$、$\Re^n$、$\Re^k$，并为它们任取基 $A$、$B$、$C$、$D$。
于是，矩阵 $H$ 相对于 $C,D$ 表示线性映射 $h$，$G$ 相对于 $B,C$ 表示 $g$，$F$ 相对于 $A,B$ 表示 $f$。
证明中可以使用这些映射。

第二部分类似，只是由于要将 $G$ 与 $H$ 相加，必须让它们表示定义域、陪域均相同的映射。
    
\end{ans}
\begin{ans}{三.IV.2.30}
     \begin{exparts}
\partsitem 秩为 \( n \) 的矩阵的乘积，秩可以小于或等于 \( n \)，但不能大于 \( n \)。

为说明秩可以降低，考虑向坐标轴的投影 \( \map{\pi_x,\pi_y}{\Re^2}{\Re^2} \)。
两者的秩均为一，但复合 $\composed{\pi_x}{\pi_y}$ 是零映射，秩为零。
转成表示这些映射的矩阵，就得到
        \begin{equation*}
          \rep{\pi_x}{\stdbasis_2,\stdbasis_2}\cdot\rep{\pi_y}{\stdbasis_2,\stdbasis_2}=
          \begin{mat}[r]
            1  &0  \\
            0  &0
          \end{mat}
          \begin{mat}[r]
            0  &0  \\
            0  &1
          \end{mat}=
          \begin{mat}[r]
            0  &0  \\
            0  &0
          \end{mat}
        \end{equation*}

为证明乘积的秩不能大于 \( n \)，可以使用线性相关向量组的像仍线性相关这一性质。
若 \( \map{h}{V}{W} \) 和 \( \map{g}{W}{X} \) 的秩均为 \( n \)，则值域
\( \rangespace{\composed{g}{h}} \) 中多于 \( n \) 个元素的集合，是 \( W \) 中多于 \( n \) 个元素的集合在 \( g \) 下的像；原像集合线性相关，因为 \( h \) 的秩为 \( n \)。
线性相关向量组的像仍线性相关，所以复合值域中任何多于 \( n \) 个元素的集合都线性相关。
（顺便注意，论证没有用到 \( g \) 的秩，参见下一问。）
\partsitem 固定空间和基，考虑相应线性映射 \( f \) 和 \( g \)。回顾，映射像的维数（映射的秩）不超过定义域维数，并考虑箭头图。
          \begin{equation*}
            \begin{matrix}
              V &\mapsunder{f} &\rangespace{f} &\mapsunder{g}
                &\rangespace{\composed{g}{f}}
            \end{matrix}
          \end{equation*}
首先，由刚才的结论，\( \rangespace{f} \) 的像的维数不超过 \( \rangespace{f} \) 的维数。
另一方面，\( \rangespace{f} \) 是 \( g \) 的定义域的子集，因此其像的维数不超过 \( g \) 的定义域的维数。
结合两者，复合的秩不超过两个秩的最小值。

矩阵的结论立即得到。
     \end{exparts}
    
\end{ans}
\begin{ans}{三.IV.2.31}
“可交换”关系具有自反性和对称性，但不具有传递性。例如，取
      \begin{equation*}
        G=\begin{mat}[r]
          1  &2  \\
          3  &4
        \end{mat}
        \quad
        H=\begin{mat}[r]
          1  &0  \\
          0  &1
        \end{mat}
        \quad
        J=\begin{mat}[r]
          5  &6  \\
          7  &8
        \end{mat}
      \end{equation*}
则 $G$ 与 $H$ 可交换，$H$ 与 $J$ 可交换，但 $G$ 与 $J$ 不可交换。
    
\end{ans}
\begin{ans}{三.IV.2.32}
      \begin{exparts}
\partsitem 以下两种顺序均可。
          \begin{equation*}
            \colvec{x \\ y \\ z}
              \mapsunder{\pi_x}
            \colvec{x \\ 0 \\ 0}
              \mapsunder{\pi_y}
            \colvec[r]{0 \\ 0 \\ 0}
            \qquad
            \colvec{x \\ y \\ z}
              \mapsunder{\pi_y}
            \colvec{0 \\ y \\ 0}
              \mapsunder{\pi_x}
            \colvec[r]{0 \\ 0 \\ 0}
          \end{equation*}
\partsitem 复合是在次数不超过四的多项式空间上的五阶求导映射 $d^5/dx^5$。
\partsitem 相对于自然基，
          \begin{equation*}
            \rep{\pi_x}{\stdbasis_3,\stdbasis_3}=\begin{mat}[r]
              1  &0  &0  \\
              0  &0  &0  \\
              0  &0  &0
            \end{mat}
            \qquad
            \rep{\pi_y}{\stdbasis_3,\stdbasis_3}=\begin{mat}[r]
              0  &0  &0  \\
              0  &1  &0  \\
              0  &0  &0
            \end{mat}
          \end{equation*}
它们无论以何种顺序相乘，都得到零矩阵。
\partsitem 取 \( B=\sequence{1,x,x^2,x^3,x^4} \)，有
          \begin{equation*}
            \rep{\frac{d^2}{dx^2}}{B,B}=\begin{mat}[r]
              0  &0  &2  &0  &0  \\
              0  &0  &0  &6  &0  \\
              0  &0  &0  &0  &12 \\
              0  &0  &0  &0  &0  \\
              0  &0  &0  &0  &0
            \end{mat}
            \qquad
            \rep{\frac{d^3}{dx^3}}{B,B}=\begin{mat}[r]
              0  &0  &0  &6  &0  \\
              0  &0  &0  &0  &24 \\
              0  &0  &0  &0  &0  \\
              0  &0  &0  &0  &0  \\
              0  &0  &0  &0  &0
            \end{mat}
          \end{equation*}
它们无论以何种顺序相乘，都得到零矩阵。
      \end{exparts}
    
\end{ans}
\begin{ans}{三.IV.2.33}
注意 \( (S+T)(S-T)=S^2-ST+TS-T^2 \)，因此可尝试寻找不可交换的矩阵，使 $-ST$ 与 $TS$ 不能抵消。取
      \begin{equation*}
        S=\begin{mat}[r]
            1  &2  \\
            3  &4
          \end{mat}
        \quad
        T=\begin{mat}[r]
            5  &6  \\
            7  &8
          \end{mat}
      \end{equation*}
便得到所需的不等式。
      \begin{equation*}
        (S+T)(S-T)=\begin{mat}[r]
            -56  &-56  \\
            -88  &-88
          \end{mat}
        \qquad
        S^2-T^2=\begin{mat}[r]
            -60  &-68  \\
            -76  &-84
          \end{mat}
      \end{equation*}
    
\end{ans}
\begin{ans}{三.IV.2.34}
由于恒等映射在基 $B$ 上的作用为 $\vec{\beta}_1\mapsto\vec{\beta}_1$、\ldots、$\vec{\beta}_n\mapsto\vec{\beta}_n$，其表示如下。
      \begin{equation*}
        \begin{mat}
          1  &0  &0  &  &0 \\
          0  &1  &0  &  &0 \\
          0  &0  &1  &  &0 \\
             &   &   &\ddots  \\
          0  &0  &0  &  &1
        \end{mat}
      \end{equation*}
第二部分由 \nearbytheorem{th:MatMultRepComp} 显然得到。
    
\end{ans}
\begin{ans}{三.IV.2.35}
      \begin{exparts}
\partsitem 向量空间 \( \matspace_{\nbyn{2}} \) 的维数为四。\( \set{T^4,\dots,T,I} \) 给出五个元素，所以它们线性相关。
\partsitem 若 \( T \) 为 \( \nbyn{n} \) 矩阵，推广上一问的论证可知，存在这样的多项式，其次数不超过 \( n^2 \)，因为
\( \set{T^{n^2},\dots,T,I} \) 给出了 $n^2$ 维空间 $\matspace_{\nbyn{n}}$ 中的 \( n^2+1 \) 个元素。
\partsitem 先计算各次幂，
           \begin{equation*}
              T^2=
              \begin{mat}[r]
                 1/2         &-\sqrt{3}/2  \\
                 \sqrt{3}/2  &1/2
              \end{mat}
              \qquad
              T^3=
              \begin{mat}[r]
                 0           &-1           \\
                 1           &0
              \end{mat}
             \qquad
              T^4=
              \begin{mat}[r]
                 -1/2        &-\sqrt{3}/2  \\
                 \sqrt{3}/2  &-1/2
              \end{mat}
           \end{equation*}
（注意，旋转 $\pi/6$ 三次相当于旋转 $\pi/2$，确实是 $T^3$ 所表示的变换）。
再令 \( c_4T^4+c_3T^3+c_2T^2+c_1T+c_0I \) 等于零矩阵，
           \begin{multline*}
             \begin{mat}[r]
                 -1/2        &-\sqrt{3}/2  \\
                 \sqrt{3}/2  &-1/2
              \end{mat}c_4
              +\begin{mat}[r]
                 0           &-1           \\
                 1           &0
              \end{mat}c_3
              +\begin{mat}[r]
                 1/2         &-\sqrt{3}/2  \\
                 \sqrt{3}/2  &1/2
              \end{mat}c_2                                     \\
              +\begin{mat}[r]
                 \sqrt{3}/2  &-1/2       \\
                 1/2         &\sqrt{3}/2
               \end{mat}c_1
              +\begin{mat}[r]
                 1  &0  \\
                 0  &1
               \end{mat}c_0
               =\begin{mat}[r]
                  0  &0  \\
                  0  &0
               \end{mat}
           \end{multline*}
得到以下线性方程组。
           \begin{equation*}
             \begin{linsys}{5}
                -(1/2)c_4         &  &  &+ &(1/2)c_2
                    &+ &(\sqrt{3}/2)c_1  &+  &c_0  &=  &0      \\
                -(\sqrt{3}/2)c_4  &- &c_3 &- &(\sqrt{3}/2)c_2
                    &- &(1/2)c_1         &   &          &=  &0        \\
                 (\sqrt{3}/2)c_4  &+ &c_3 &+ &(\sqrt{3}/2)c_2
                    &+ &(1/2)c_1  &   &            &=  &0      \\
                -(1/2)c_4              &  &  &+ &(1/2)c_2
                    &+ &(\sqrt{3}/2)c_1  &+  &c_0  &=  &0
              \end{linsys}
           \end{equation*}
作高斯消元。
           \begin{align*}
             &\grstep{-\rho_1+\rho_4}
             \repeatedgrstep{\rho_2+\rho_3}
             \begin{linsys}{5}
                -(1/2)c_4         &  &  &+ &(1/2)c_2
                    &+ &(\sqrt{3}/2)c_1  &+  &c_0  &=  &0      \\
                -(\sqrt{3}/2)c_4  &- &c_3 &- &(\sqrt{3}/2)c_2
                    &- &(1/2)c_1         &   &          &=  &0   \\
                                  &  &  &  &
                    &  &                 &   &0    &=  &0      \\
                                  &  &  &  &
                    &  &                 &   &0    &=  &0
             \end{linsys}                                               \\
             &\grstep{-\sqrt{3}\rho_1+\rho_2}
             \begin{linsys}{5}
                -(1/2)c_4         &  &  &+ &(1/2)c_2
                    &+ &(\sqrt{3}/2)c_1  &+  &c_0  &=  &0      \\
                                  &- &c_3 &- &\sqrt{3}c_2
                    &- &2c_1             &-  &\sqrt{3}c_0  &=  &0   \\
                                  &  &  &  &
                    &  &                 &   &0    &=  &0      \\
                                  &  &  &  &
                    &  &                 &   &0    &=  &0
              \end{linsys}
           \end{align*}
令 \( c_4 \)、\( c_3 \)、\( c_2 \) 为零，则 \( c_1 \)、\( c_0 \) 也为零，所以一次或零次多项式均不符合要求。
令 \( c_4 \)、\( c_3 \) 为零，$c_2$ 为一，得到线性方程组
           \begin{equation*}
             \begin{linsys}{3}
               (1/2)     &+  &(\sqrt{3}/2)c_1  &+  &c_0          &=  &0 \\
               -\sqrt{3} &-  &2c_1             &-  &\sqrt{3}c_0  &=  &0
             \end{linsys}
           \end{equation*}
其解为 $c_1=-\sqrt{3}$、$c_0=1$。
所以，$m(x)=x^2-\sqrt{3}x+1$ 是矩阵 $T$ 的最小多项式。
      \end{exparts}
    
\end{ans}
\begin{ans}{三.IV.2.36}
按常规方法验证：
      \begin{equation*}
        a_0+a_1x+\dots+a_nx^n\mapsunder{s}a_0x+a_1x^2+\dots+a_nx^{n+1}
                  \mapsunder{d/dx} a_0+2a_1x+\dots+(n+1)a_nx^n
      \end{equation*}
而
      \begin{equation*}
        a_0+a_1x+\dots+a_nx^n
                  \mapsunder{d/dx} a_1+\dots+na_nx^{n-1}
                  \mapsunder{s}a_1x+\dots+a_nx^n
      \end{equation*}
所以，映射 $(\composed{d/dx}{s})-(\composed{s}{d/dx})$ 的作用为
$a_0+a_1x+\dots+a_nx^n\mapsto a_0+a_1x+\dots+a_nx^n$。
    
\end{ans}
\begin{ans}{三.IV.2.37}
       \begin{exparts}
\partsitem 按照本小节末尾的评注，\( (FG)H \) 的 \( i,j \) 元素如下，
           \begin{multline*}
             \sum_{t=1}^s\bigl(\sum_{k=1}^{r} f_{i,k}g_{k,t}\bigr)h_{t,j}
             =\sum_{t=1}^s\sum_{k=1}^{r} (f_{i,k}g_{k,t})h_{t,j}
             =\sum_{t=1}^s\sum_{k=1}^{r} f_{i,k}(g_{k,t}h_{t,j})  \\
             =\sum_{k=1}^{r}\sum_{t=1}^s f_{i,k}(g_{k,t}h_{t,j})
             =\sum_{k=1}^{r}f_{i,k}\bigl(\sum_{t=1}^s g_{k,t}h_{t,j}\bigr)
           \end{multline*}
（第一个等号利用分配律将 $h$ 乘进去；第二个用实数结合律；第三个用实数交换律；第四个用分配律提出 $f$），这就是 \( F(GH) \) 的 \( i,j \) 元素。
\partsitem \( h(\vec{v})\) 的第 \( k \) 个分量为
           \begin{equation*}
             \sum_{j=1}^n h_{k,j}v_j
           \end{equation*}
所以，\( \composed{g}{h}\,(\vec{v}) \) 的第 \( i \) 个分量如下，
           \begin{multline*}
              \sum_{k=1}^r g_{i,k}\bigl(\sum_{j=1}^n h_{k,j}v_j\bigr)
              =\sum_{k=1}^r \sum_{j=1}^n g_{i,k}h_{k,j}v_j
              =\sum_{k=1}^r \sum_{j=1}^n (g_{i,k}h_{k,j})v_j
                                                              \\
              =\sum_{j=1}^n \sum_{k=1}^r (g_{i,k}h_{k,j})v_j
              =\sum_{j=1}^n (\sum_{k=1}^r g_{i,k}h_{k,j})\,v_j
           \end{multline*}
（第一个等号利用分配律将 $g$ 乘进去；第二个用实数结合律；第三个用实数交换律；第四个用分配律提出 $v$）。
       \end{exparts}
    
\end{ans}
\protect \subsection {三.IV.3: 矩阵乘法的运算机制}
\begin{ans}{三.IV.3.24}
      \begin{exparts}
        \partsitem  Acting from the left, the $P_{1,2}$ matrix swaps the
          第一、二行。
          $\begin{mat}
             3 &4 \\
             1 &2
           \end{mat}$
        \partsitem
          从右侧作用时，$P_{1,2}$ 交换第一、二
          columns.
          $\begin{mat}
             2 &1 \\
             4 &3
           \end{mat}$
        \partsitem From the left, $P_{2,3}$ swaps the second and third
          rows.
          $\begin{mat}
             1 &2 &3 \\
             7 &8 &9 \\
             4 &5 &6
           \end{mat}$
      \end{exparts}
    
\end{ans}
\begin{ans}{三.IV.3.25}
      \begin{exparts}
\partsitem 第二个矩阵的第一行乘以 \( 3 \)。
          \begin{equation*}
            \begin{mat}[r]
              3  &6  \\
              3  &4
            \end{mat}
          \end{equation*}
\partsitem 第二个矩阵的第二行乘以 \( 2 \)。
          \begin{equation*}
            \begin{mat}[r]
              1  &2  \\
              6  &8
            \end{mat}
          \end{equation*}
\partsitem 第二个矩阵执行行组合：将第一行的 \( -2 \) 倍加到第二行。
          \begin{equation*}
            \begin{mat}[r]
              1  &2  \\
              1  &0
            \end{mat}
          \end{equation*}
\partsitem 第一个矩阵执行列运算：将第一列的 $-1$ 倍加到第二列。
          \begin{equation*}
            \begin{mat}[r]
              1  &1  \\
              3  &1
            \end{mat}
          \end{equation*}
\partsitem 第一个矩阵的两列交换。
          \begin{equation*}
            \begin{mat}[r]
              2  &1  \\
              4  &3
            \end{mat}
          \end{equation*}
      \end{exparts}
    
\end{ans}
\begin{ans}{三.IV.3.26}
      \begin{exparts}
\partsitem 第二个矩阵的第一行乘以 \( -3 \)，第二行乘以 \( 0 \)。
          \begin{equation*}
            \begin{mat}[r]
              -3  &-6  \\
               0  &0
            \end{mat}
          \end{equation*}
\partsitem 第二个矩阵的第一行乘以 \( 4 \)，第二行乘以 \( 2 \)。
          \begin{equation*}
            \begin{mat}[r]
              4  &8  \\
              6  &8
            \end{mat}
          \end{equation*}
      \end{exparts}
    
\end{ans}
\begin{ans}{三.IV.3.27}
      \begin{exparts}
\partsitem 这个矩阵交换第一行与第三行。
          \begin{equation*}
            \begin{mat}
              0  &1 &0 \\
              1  &0 &0 \\
              0  &0 &1
            \end{mat}
            \begin{mat}
              a &b &c \\
              d &e &f \\
              g &h &i
            \end{mat}
            =
            \begin{mat}
              d &e &f \\
              a &b &c \\
              g &h &i
            \end{mat}
          \end{equation*}
\partsitem 这个矩阵交换第一列与第二列。
          \begin{equation*}
            \begin{mat}
              a  &b  \\
              c  &d
            \end{mat}
            \begin{mat}
              0  &1  \\
              1  &0
            \end{mat}
            =
            \begin{mat}
              b  &a  \\
              d  &c
            \end{mat}
          \end{equation*}
      \end{exparts}
    
\end{ans}
\begin{ans}{三.IV.3.28}
依次乘 $C_{1,2}(-2)$、$C_{1,3}(-7)$、$C_{2,3}(-3)$，注意从右向左的顺序。
      \begin{multline*}
        \begin{mat}
          1  &0 &0 \\
          0  &1 &0 \\
          0  &-3 &1
        \end{mat}
        \begin{mat}
          1  &0 &0 \\
          0  &1 &0 \\
          -7 &0 &1
        \end{mat}
        \begin{mat}
          1 &0 &0 \\
         -2 &1 &0 \\
          0 &0 &1
        \end{mat}
          \begin{mat}
            1 &2 &1  &0   \\
            2 &3 &1  &-1  \\
            7 &11 &4 &-3
          \end{mat}                 \\
        =
          \begin{mat}
            1 &2 &1  &0   \\
            0 &-1 &-1  &-1  \\
            0 &0  &0 &0
          \end{mat}
      \end{multline*}
    
\end{ans}
\begin{ans}{三.IV.3.29}
乘积为单位矩阵（回顾 $\cos^2\theta+\sin^2\theta =1$）。
其原因是：给定矩阵相对于标准基表示 \( \Re^2 \) 中旋转 \( \theta \) 弧度，而转置表示旋转 \( -\theta \) 弧度，两者相互抵消。
     
\end{ans}
\begin{ans}{三.IV.3.30}
      \begin{exparts}
\partsitem 邻接矩阵如下（例如，第一行表明伯灵顿只有一条连接道路，即通往威努斯基的道路）。
          \begin{equation*}
            \begin{mat}[r]
              0  &0  &0  &0  &1  \\
              0  &0  &1  &1  &1  \\
              0  &1  &0  &1  &0  \\
              0  &1  &1  &0  &0  \\
              1  &1  &0  &0  &0
            \end{mat}
          \end{equation*}
\partsitem 因为道路是双向的，连接城市~\( i \) 与城市~\( j \) 的道路，也连接城市~\( j \) 与城市~\( i \)。
\partsitem 邻接矩阵的平方给出各城市之间经过两条道路的行程数。
      \end{exparts}
   
\end{ans}
\begin{ans}{三.IV.3.31}
两数组的矩阵乘积给出每个人的应付工资。
    
\end{ans}
\begin{ans}{三.IV.3.32}
按常规方法作高斯–若尔当消元。
    \begin{equation*}
        \begin{mat}
          1 &2  &0 \\
          2 &-1 &0 \\
          3 &1 &2
        \end{mat}
      \grstep[-3\rho_1+\rho_3]{-2\rho_1+\rho_2}
      \grstep{-\rho_2+\rho_3}
      \grstep[(1/2)\rho_3]{-(1/5)\rho_2}
      \grstep{-2\rho_2+\rho_1}
      \begin{mat}
        1 &0 &0 \\
        0 &1 &0 \\
        0 &0 &1
      \end{mat}
    \end{equation*}
于是得到初等消元矩阵 $R_1,\ldots,R_6$，使 $R_6\cdot R_5\cdots R_1\cdot T=I$。
将这些矩阵移到另一边，注意顺序。例如，先在两边左乘 $R_6^{-1}$，得
$(R_6^{-1}R_6)\cdot R_5\cdots R_1\cdot T=R_6^{-1}I$，化简为
$R_5\cdots R_1\cdot T=R_6^{-1}$，依此类推。
    \begin{align*}
      T=
      &\begin{mat}
        1 &0 &0 \\
        -2 &1 &0 \\
        0 &0 &1
      \end{mat}^{-1}
      \begin{mat}
        1 &0 &0 \\
        0 &1 &0 \\
        -3 &0 &1
      \end{mat}^{-1}
      \begin{mat}
        1 &0 &0 \\
        0 &1 &0 \\
        0 &-1 &1
      \end{mat}^{-1}          \\
      &\qquad\cdot
      \begin{mat}
        1 &0 &0 \\
        0 &-1/5 &0 \\
        0 &0 &1
      \end{mat}^{-1}
      \begin{mat}
        1 &0 &0 \\
        0 &1 &0 \\
        0 &0 &1/2
      \end{mat}^{-1}
      \begin{mat}
        1 &-2 &0 \\
        0 &1 &0 \\
        0 &0 &1
      \end{mat}^{-1}
    \end{align*}
初等消元矩阵的逆很容易求。例如，要撤销将第~$i$ 行的 $k$ 倍加到第~$j$ 行的操作，只需将第~$i$ 行的 $-k$ 倍加到第~$j$ 行。
由引理~\ref{GrByMatMult}，$C_{i,j}(k)^{-1}=C_{i,j}(-k)$。
    \begin{equation*}
      =
      \begin{mat}
        1 &0 &0 \\
        2 &1 &0 \\
        0 &0 &1
      \end{mat}
      \begin{mat}
        1 &0 &0 \\
        0 &1 &0 \\
        3 &0 &1
      \end{mat}
      \begin{mat}
        1 &0 &0 \\
        0 &1 &0 \\
        0 &1 &1
      \end{mat}
      \begin{mat}
        1 &0 &0 \\
        0 &-5 &0 \\
        0 &0 &1
      \end{mat}
      \begin{mat}
        1 &0 &0 \\
        0 &1 &0 \\
        0 &0 &2
      \end{mat}
      \begin{mat}
        1 &2 &0 \\
        0 &1 &0 \\
        0 &0 &1
      \end{mat}
    \end{equation*}
    
\end{ans}
\begin{ans}{三.IV.3.33}
由单位矩阵得到此矩阵的一种方法是作列运算：先将第二列乘以三，再将所得第二列的相反数加到第一列。
      \begin{equation*}
        \begin{mat}[r]
          1  &0  \\
          0  &1
        \end{mat}
        \grstep{}
        \begin{mat}[r]
          1  &0  \\
          0  &3
        \end{mat}
        \grstep{}
        \begin{mat}[r]
          1  &0  \\
          -3  &3
        \end{mat}
      \end{equation*}
与行运算不同，列运算按从左到右的顺序书写，因此以下矩阵乘积表示上述两个操作。
      \begin{equation*}
        \begin{mat}[r]
          1  &0  \\
          0  &3
        \end{mat}
        \begin{mat}[r]
          1  &0  \\
         -1  &1
        \end{mat}
      \end{equation*}
\textit{评注。}也可以用行运算得到所需矩阵：从单位矩阵出发，先将第一行的相反数加到第二行，再将第二行乘以三。
由于连续行运算的矩阵乘积按从右到左书写，这两个行运算由同一个矩阵乘积表示。
    
\end{ans}
\begin{ans}{三.IV.3.34}
零矩阵是对角矩阵，所以该集合非空。显然它对数乘与加法封闭，因此是子空间。
维数为 \( n \)，以下是一组基。
      \begin{equation*}
        \set{\begin{mat}
               1  &0  &\ldots    \\
               0  &0             \\
                  &   &\ddots    \\
               0  &0  &      &0
             \end{mat},\ldots,
             \begin{mat}
               0  &0  &\ldots    \\
               0  &0             \\
                  &   &\ddots    \\
               0  &0  &      &1
             \end{mat}  }
      \end{equation*}
    
\end{ans}
\begin{ans}{三.IV.3.35}
不。在 \( \polyspace_1 \) 中，相对于不同的基 \( B=\sequence{1,x} \) 和 \( D=\sequence{1+x,1-x} \)，恒等变换由以下矩阵表示。
      \begin{equation*}
         \rep{\text{id}}{B,D}=
         \begin{mat}[r]
           1/2  &1/2  \\
           1/2  &-1/2
         \end{mat}_{B,D}
      \end{equation*}
    
\end{ans}
\begin{ans}{三.IV.3.36}
对任意标量 \( r \) 和方阵 \( H \)，有
\( (rI)H=r(IH)=rH=r(HI)=(Hr)I=H(rI) \)。

没有其他这样的矩阵。下面对 $\nbyn{2}$ 矩阵的论证很容易推广到 $\nbyn{n}$。
若一个矩阵与所有矩阵可交换，它就与以下矩阵单位可交换。
      \begin{equation*}
        \begin{mat}
          0  &a  \\
          0  &c
        \end{mat}
        =\begin{mat}
          a  &b  \\
          c  &d
        \end{mat}
        \begin{mat}[r]
          0  &1  \\
          0  &0
        \end{mat}
        =\begin{mat}[r]
          0  &1  \\
          0  &0
        \end{mat}
        \begin{mat}
          a  &b  \\
          c  &d
        \end{mat}
        =\begin{mat}
          c  &d  \\
          0  &0
        \end{mat}
      \end{equation*}
由此首先得到左上元素~$a$ 必须等于右下元素~$d$，还得到左下元素~$c$ 为零。
对右上元素~$b$ 的论证类似。
    
\end{ans}
\begin{ans}{三.IV.3.37}
错误。以下两个矩阵不可交换。
      \begin{equation*}
         \begin{mat}[r]
           1  &2  \\
           3  &4
         \end{mat}
         \qquad
         \begin{mat}[r]
           5  &6  \\
           7  &8
         \end{mat}
      \end{equation*}
    
\end{ans}
\begin{ans}{三.IV.3.38}
置换矩阵每行、每列恰有一个一，其余元素全为零。
固定这样的矩阵。设第 \( i \) 行的一位于第 \( j \) 列，则其他行在第 \( j \) 列都为零。
因此，第 \( i \) 行与其他任意一行的点积为零。

乘积的第 \( i \) 行由原矩阵第 \( i \) 行与转置各列的点积组成。
由上一段，这些点积除了第 \( i \) 个为一外，其余全为零。
    
\end{ans}
\begin{ans}{三.IV.3.39}
推广为第 $i_1$ 行与第 $i_2$ 行即可。\( GH \) 的第~$i$ 行由 \( G \) 的第~$i$ 行与 \( H \) 各列的点积组成。
因此，若 \( G \) 的第 \( i_1 \)、\( i_2 \) 行相同，则 \( GH \) 的对应两行也相同。
    
\end{ans}
\begin{ans}{三.IV.3.40}
若两个对角矩阵的乘积有定义\Dash 即两者均为 $\nbyn{n}$\Dash 则乘积的对角线由对应对角元素的乘积组成：
对相同大小的对角矩阵 \( G,H \)，\( GH \) 除 \( i,i \) 元素为 \( g_{i,i}h_{i,i} \) 外，其余均为零。
    
\end{ans}
\begin{ans}{三.IV.3.41}
\( GH \) 的第 \( i \) 行由 \( G \) 的第 \( i \) 行与 \( H \) 各列的点积组成。
零行与任意列的点积都为零。

对列需正确表述：若 \( H \) 有一列全为零，则 \( GH \)（若有定义）也有一列全为零。证明很简单。
    
\end{ans}
\begin{ans}{三.IV.3.42}
最简单的方法可能是证明每个 \( \nbym{n}{m} \) 矩阵都能唯一地写为矩阵单位的线性组合：
      \begin{equation*}
        c_1\begin{mat}
             1  &0  &\ldots  \\
             0  &0           \\
             \vdots
           \end{mat}
        +\dots+
        c_{n,m}\begin{mat}
             0  &0      &\ldots  \\
             \vdots              \\
             0  &\ldots &       &1
           \end{mat}
        =\begin{mat}
           a_{1,1} &a_{1,2} &\ldots          \\
           \vdots                            \\
           a_{n,1} &\ldots  &      &a_{n,m}
         \end{mat}
      \end{equation*}
具有唯一解 \( c_1=a_{1,1} \)、\( c_2=a_{1,2} \)，依此类推。
    
\end{ans}
\begin{ans}{三.IV.3.43}
将该矩阵记为 \( F \)。有
      \begin{equation*}
        F^2=\begin{mat}[r]
          2  &1  \\
          1  &1
        \end{mat}
        \quad
        F^3=\begin{mat}[r]
          3  &2  \\
          2  &1
        \end{mat}
        \quad
        F^4=\begin{mat}[r]
          5  &3  \\
          3  &2
        \end{mat}
      \end{equation*}
一般地，
      \begin{equation*}
        F^n=\begin{mat}
          f_{n+1} &f_n  \\
          f_n     &f_{n-1}
        \end{mat}
      \end{equation*}
其中 \( f_i \) 是第 \( i \) 个斐波那契数，满足 \( f_i=f_{i-1}+f_{i-2} \)、\( f_0=0 \)、\( f_1=1 \)。
根据以下等式可用归纳法验证。
      \begin{equation*}
        \begin{mat}
          f_{i-1}  &f_{i-2} \\
          f_{i-2}  &f_{i-3}
        \end{mat}
        \begin{mat}[r]
          1  &1  \\
          1  &0
        \end{mat}
        =\begin{mat}
           f_i     &f_{i-1}  \\
           f_{i-1} &f_{i-2}
        \end{mat}
      \end{equation*}
    
\end{ans}
\begin{ans}{三.IV.3.44}
\textit{第五章将给出一个不那么依赖计算的原因\Dash 矩阵的迹是特征多项式的第二个系数\Dash 目前先使用下标。}有
     \begin{align*}
       \trace (GH)
       &=(g_{1,1}h_{1,1}+g_{1,2}h_{2,1}+\dots+g_{1,n}h_{n,1})  \\
       &\text{}\quad+(g_{2,1}h_{1,2}+g_{2,2}h_{2,2}+\dots+g_{2,n}h_{n,2}) \\
       &\text{}\quad
        +\cdots+(g_{n,1}h_{1,n}+g_{n,2}h_{2,n}+\dots+g_{n,n}h_{n,n})
     \end{align*}
而
     \begin{align*}
       \trace (HG)
       &=(h_{1,1}g_{1,1}+h_{1,2}g_{2,1}+\dots+h_{1,n}g_{n,1})  \\
       &\text{}\quad+(h_{2,1}g_{1,2}+h_{2,2}g_{2,2}+\dots+h_{2,n}g_{n,2}) \\
       &\text{}\quad
         +\cdots+(h_{n,1}g_{1,n}+h_{n,2}g_{2,n}+\dots+h_{n,n}g_{n,n})
     \end{align*}
两者相等。
    
\end{ans}
\begin{ans}{三.IV.3.45}
矩阵为上三角矩阵，当且仅当 \( i>j \) 时其 $i,j$ 元素为零。
因此，若 \( G,H \) 为上三角矩阵，则 \( i>j \) 时 \( h_{i,j} \) 与 \( g_{i,j} \) 都为零。
乘积元素 \( p_{i,j}=g_{i,1}h_{1,j}+\dots+g_{i,n}h_{n,j} \) 要非零，至少有一项非零，
即至少对某个加项 \( g_{i,r}h_{r,j} \)，同时有 \( i\leq r \) 和 \( r\leq j \)。
当 \( i>j \) 时不可能如此，所以两个上三角矩阵的乘积仍为上三角矩阵。（下三角矩阵的论证类似。）
   
\end{ans}
\begin{ans}{三.IV.3.46}
乘积第 \( i \) 行的元素之和如下。
     \begin{align*}
       p_{i,1}+\cdots+p_{i,n}
       &=(h_{i,1}g_{1,1}+h_{i,2}g_{2,1}+\dots+h_{i,n}g_{n,1})  \\
       &\text{}\quad+(h_{i,1}g_{1,2}+h_{i,2}g_{2,2}+\dots+h_{i,n}g_{n,2}) \\
       &\text{}\quad
        +\dots+(h_{i,1}g_{1,n}+h_{i,2}g_{2,n}+\dots+h_{i,n}g_{n,n})  \\
       &=h_{i,1}(g_{1,1}+g_{1,2}+\dots+g_{1,n})  \\
       &\text{}\quad+h_{i,2}(g_{2,1}+g_{2,2}+\dots+g_{2,n}) \\
       &\text{}\quad
        +\dots+h_{i,n}(g_{n,1}+g_{n,2}+\dots+g_{n,n})  \\
       &=h_{i,1}\cdot 1+\dots+h_{i,n}\cdot 1  \\
       &=1
     \end{align*}
    
\end{ans}
\begin{ans}{三.IV.3.47}
固定基 $B=\sequence{\vec{\beta}_1, \vec{\beta}_2}$。表示以下映射的矩阵，
      \begin{equation*}
        \vec{\beta}_1 \mapsunder{h} \vec{\beta}_1
        \quad
        \vec{\beta}_2 \mapsunder{h} \zero
      \end{equation*}
以及
      \begin{equation*}
        \vec{\beta}_1 \mapsunder{g} \vec{\beta}_2
        \quad
        \vec{\beta}_2 \mapsunder{g} \zero
      \end{equation*}
就符合要求。例如，使用标准基得到以下矩阵。
      \begin{equation*}
        H=
        \begin{mat}
          1 &0 \\
          0 &0
        \end{mat}
        \qquad
        G=
        \begin{mat}
          0 &1 \\
          0 &0
        \end{mat}
      \end{equation*}
注意 $H^2$ 的秩为~$1$，而 $G^2$ 的秩为~$0$。
    
\end{ans}
\begin{ans}{三.IV.3.48}
      \begin{exparts}
\partsitem 每个元素 \(p_{i,j}=g_{i,1}h_{1,j}+\dots+g_{1,r}h_{r,1} \) 需要 \( r \) 次乘法，共有 \( m\cdot n \) 个元素，
所以共需 \( m\cdot n\cdot r \) 次乘法。
\partsitem 设 \( H_1 \) 为 \( \nbym{5}{10} \)，\( H_2 \) 为 \( \nbym{10}{20} \)，\( H_3 \) 为 \( \nbym{20}{5} \)，\( H_4 \) 为 \( \nbym{5}{1} \)。则
           \begin{center}
             \begin{tabular}{l|l}
               \multicolumn{1}{c}{\textit{this association}}
                &\multicolumn{1}{c}{\textit{uses this many multiplications}}\\
               \hline
               \( ((H_1H_2)H_3)H_4 \)     &\( 1000+500+25=1525 \)  \\
               \( (H_1(H_2H_3))H_4 \)     &\( 1000+250+25=1275 \)  \\
               \( (H_1H_2)(H_3H_4) \)     &\( 1000+100+100=1200 \)  \\
               \( H_1(H_2(H_3H_4)) \)     &\( 100+200+50=350    \)  \\
               \( H_1((H_2H_3)H_4) \)     &\( 1000+50+50=1100    \)
             \end{tabular}
           \end{center}
由表可知哪种最省。
\partsitem 下面是 S.~Winograd 对 V.~Strassen 公式的改进：令 \( w=aA-(a-c-d)(A-C+D) \)，则
           \begin{equation*}
               \begin{mat}
                 a  &b  \\
                 c  &d
               \end{mat}
               \begin{mat}
                 A  &B  \\
                 C  &D
                 \end{mat}    \\
               =
               \begin{mat}
                 \alpha  &\beta  \\
                 \gamma  &\delta
               \end{mat}
               \end{equation*}
其中 $\alpha=aA+bB$，$\beta=w+(c+d)(C-A)+(a+b-c-d)D$，
$\gamma=w+(a-c)(D-C)-d(A-B-C+D)$，$\delta=w+(a-c)(D-C)+(c+d)(C-A)$。
               % \begin{mat}
               %    aA+bB
               %    &w+(c+d)(C-A)+(a+b-c-d)D \\
               %    w+(a-c)(D-C)-d(A-B-C+D)
               %    &w+(a-c)(D-C)+(c+d)(C-A)
               % \end{mat}
这需要七次乘法和十五次加法（保存中间结果）。
      \end{exparts}
    
\end{ans}
\begin{ans}{三.IV.3.49}
\answerasgiven
不能。设 \( A \)、\( B \) 相对于标准基表示 \( \Re^3 \) 上的以下变换。
      \begin{equation*}
        \colvec{x \\ y \\ z}\mapsunder{a}\colvec{x \\ y \\ 0}
        \qquad
        \colvec{x \\ y \\ z}\mapsunder{b}\colvec{0 \\ x \\ y}
      \end{equation*}
注意
      \begin{equation*}
        \colvec{x \\ y \\ z}\mapsunder{abab}\colvec[r]{0 \\ 0 \\ 0}
        \quad\text{but}\quad
        \colvec{x \\ y \\ z}\mapsunder{baba}\colvec{0 \\ 0 \\ x}.
      \end{equation*}
    
\end{ans}
\begin{ans}{三.IV.3.50}
      \answerasgiven
      \begin{exparts}
\partsitem 显然。
\partsitem 若 \( \trans{A}A\vec{x}=\zero \)，则 \( \vec{y}\cdot\vec{y}=0 \)，其中 \( \vec{y}=A\vec{x} \)。由 (a)，\( \vec{y}=\zero \)。

逆命题显然成立。
\partsitem 由 (b)，\( A\vec{x}_1 \)、\ldots、\( A\vec{x}_n \) 线性无关，当且仅当
\( \trans{A}A\vec{x}_1 \)、\ldots、\( \trans{A}A\vec{x}_n \) 线性无关。
\partsitem 有
          \begin{multline*}
            \text{col rank}(A)=\text{col rank}(\trans{A}A)
            =\dim\set{\trans{A}(A\vec{x})\suchthat \text{all \(\vec{x}\)} }  \\
            \leq \dim\set{\trans{A}\vec{y}\suchthat \text{all \(\vec{y}\)} }
            =\text{col rank}(\trans{A}).
          \end{multline*}
因此也有 \( \text{列秩}(\trans{A})\leq\text{列秩}(\trans{\trans{A}}) \)，所以
\( \text{列秩}(A)=\text{列秩}(\trans{A})=\text{行秩}(A) \)。
      \end{exparts}
    
\end{ans}
\begin{ans}{三.IV.3.51}
\answerasgiven
设 \( \sequence{\vec{z}_1,\dots,\vec{z}_k} \) 是 \( \rangespace{A}\intersection\nullspace{A} \) 的一组基（\( k \) 可以为 \( 0 \)）。
取 \( \vec{x}_1,\dots,\vec{x}_k\in V \) 使 \( A\vec{x}_i=\vec{z}_i \)。
注意 \( \set{A\vec{x}_1,\dots,A\vec{x}_k} \) 线性无关，将其扩充为 \( \rangespace{A} \) 的基
\( A\vec{x}_1,\ldots,A\vec{x}_k,A\vec{x}_{k+1},\dots,A\vec{x}_{r_1} \)，其中 \( r_1=\dim(\rangespace{A}) \)。

任取 \( \vec{x}\in V \)，写成
      \begin{equation*}
        A\vec{x}=a_1(A\vec{x}_1)+\dots+a_{r_1}(A\vec{x}_{r_1})
      \end{equation*}
所以
      \begin{equation*}
        A^2\vec{x}=a_1(A^2\vec{x}_1)+\dots+a_{r_1}(A^2\vec{x}_{r_1}).
      \end{equation*}
但 \( A\vec{x}_1,\dots,A\vec{x}_k\in\nullspace{A} \)，所以 \( A^2\vec{x}_1=\zero,\dots,A^2\vec{x}_k=\zero \)。于是知道
      \begin{equation*}
        A^2\vec{x}_{k+1},\dots,A^2\vec{x}_{r_1}
      \end{equation*}
张成 \( \rangespace{A^2} \)。

为说明 \( \set{A^2\vec{x}_{k+1},\dots,A^2\vec{x}_{r_1}} \) 线性无关，写出
      \begin{align*}
        b_{k+1}A^2\vec{x}_{k+1}+\dots+b_{r_1}A^2\vec{x}_{r_1}
        &=\zero                                                 \\
        A[b_{k+1}A\vec{x}_{k+1}+\dots+b_{r_1}A\vec{x}_{r_1}]
        &=\zero
      \end{align*}
由于 \( b_{k+1}A\vec{x}_{k+1}+\dots+b_{r_1}A\vec{x}_{r_1}\in\nullspace{A} \)，除非它为 \( \zero \)，否则产生矛盾
（显然它也在 \( \rangespace{A} \) 中，而 \( A\vec{x}_1,\ldots,A\vec{x}_k \) 是 \( \rangespace{A}\intersection\nullspace{A} \) 的基）。

因此 \( \dim(\rangespace{A^2})=r_1-k=\dim(\rangespace{A})-\dim(\rangespace{A}\intersection\nullspace{A}) \)。
    
\end{ans}
\protect \subsection {三.IV.4: 逆矩阵}
\begin{ans}{三.IV.4.12}
以下是一种做法。先作
      \begin{equation*}
       \grstep{\rho_1\leftrightarrow\rho_2}
       \begin{pmat}{rrr|rrr}
              1  &0  &1   &0  &1  &0  \\
              0  &3  &-1  &1  &0  &0  \\
              1  &-1 &0   &0  &0  &1
           \end{pmat}
       \grstep{-\rho_1+\rho_3}
       \begin{pmat}{rrr|rrr}
              1  &0  &1   &0  &1  &0  \\
              0  &3  &-1  &1  &0  &0  \\
              0  &-1 &-1  &0  &-1 &1
           \end{pmat}
       \end{equation*}
再作
       \begin{align*}
         &\grstep{(1/3)\rho_2+\rho_3}
         \begin{pmat}{rrr|rrr}
                1  &0  &1     &0   &1  &0  \\
                0  &3  &-1    &1   &0  &0  \\
                0  &0  &-4/3  &1/3 &-1 &1
             \end{pmat}                                           \\
         &\grstep[-(3/4)\rho_3]{(1/3)\rho_2}
         \begin{pmat}{rrr|rrr}
                1  &0  &1     &0    &1   &0    \\
                0  &1  &-1/3  &1/3  &0   &0    \\
                0  &0  &1     &-1/4 &3/4 &-3/4
             \end{pmat}                                  \\
         &\grstep[-\rho_3+\rho_1]{(1/3)\rho_3+\rho_2}
         \begin{pmat}{rrr|rrr}
                1  &0  &0     &1/4  &1/4 &3/4  \\
                0  &1  &0     &1/4  &1/4 &-1/4 \\
                0  &0  &1     &-1/4 &3/4 &-3/4
             \end{pmat}
    \end{align*}
从右半部分读出答案。
   
\end{ans}
\begin{ans}{三.IV.4.13}
      \begin{exparts*}
\partsitem 有逆，因为 $ad-bc=2\cdot 1-1\cdot(-1)\neq 0$。
\partsitem 有。
\partsitem 无。
      \end{exparts*}
    
\end{ans}
\begin{ans}{三.IV.4.14}
       \begin{exparts}
         \partsitem
           $\displaystyle \frac{1}{2\cdot 1-1\cdot (-1)}
            \cdot\begin{mat}[r]
             1  &-1  \\
             1  &2
           \end{mat}
          =\frac{\displaystyle 1}{\displaystyle 3}\cdot\begin{mat}[r]
             1  &-1  \\
             1  &2
           \end{mat}
          =\begin{mat}[r]
             1/3  &-1/3  \\
             1/3  &2/3
           \end{mat}$
         \partsitem
           $\displaystyle \frac{1}{0\cdot (-3)-4\cdot 1}
           \cdot\begin{mat}[r]
             -3  &-4  \\
             -1  &0
           \end{mat}
           =\begin{mat}[r]
             3/4  &1  \\
             1/4  &0
           \end{mat}$
         \partsitem The prior question shows that no inverse exists.
       \end{exparts}
     
\end{ans}
\begin{ans}{三.IV.4.15}
      \begin{exparts}
\partsitem 按常规消元。
          \begin{align*}
            \begin{pmat}{rr|rr}
              3  &1  &1  &0 \\
              0  &2  &0  &1
            \end{pmat}
            &\grstep[(1/2)\rho_2]{(1/3)\rho_1}
            \begin{pmat}{rr|rr}
              1  &1/3  &1/3  &0   \\
              0  &1    &0    &1/2
            \end{pmat}                                 \\
            &\grstep{-(1/3)\rho_2+\rho_1}
            \begin{pmat}{rr|rr}
              1  &0    &1/3  &-1/6 \\
              0  &1    &0    &1/2
            \end{pmat}
          \end{align*}
与核对所得答案一致。
          \begin{equation*}
            \begin{mat}[r]
              3  &1  \\
              0  &2
            \end{mat}^{-1}
            =\frac{1}{3\cdot 2-0\cdot 1}\cdot
            \begin{mat}[r]
              2  &-1  \\
              0  &3
            \end{mat}
            =\frac{1}{6}\cdot
            \begin{mat}[r]
              2  &-1  \\
              0  &3
            \end{mat}
          \end{equation*}
\partsitem 消元很简单。
          \begin{multline*}
            \begin{pmat}{rr|rr}
              2  &1/2  &1  &0  \\
              3  &1    &0  &1
            \end{pmat}
            \grstep{-(3/2)\rho_1+\rho_2}
            \begin{pmat}{rr|rr}
              2  &1/2  &1     &0  \\
              0  &1/4  &-3/2  &1
            \end{pmat}                       \\
            \grstep[4\rho_2]{(1/2)\rho_1}
            \begin{pmat}{rr|rr}
              1  &1/4   &1/2   &0  \\
              0  &1     &-6    &4
            \end{pmat}
            \grstep{-(1/4)\rho_2+\rho_1}
            \begin{pmat}{rr|rr}
              1  &0     &2     &-1 \\
              0  &1     &-6    &4
            \end{pmat}
          \end{multline*}
核对结果一致。
          \begin{equation*}
            \frac{1}{2\cdot 1-3\cdot (1/2)}\cdot
            \begin{mat}[r]
              1  &-1/2  \\
              -3 &2
            \end{mat}
            =2\cdot
            \begin{mat}[r]
              1  &-1/2  \\
              -3 &2
            \end{mat}
          \end{equation*}
\partsitem 尝试高斯–若尔当消元，
          \begin{equation*}
            \begin{pmat}{rr|rr}
              2  &-4  &1  &0  \\
             -1  &2   &0  &1
            \end{pmat}
            \grstep{(1/2)\rho_1+\rho_2}
            \begin{pmat}{rr|rr}
              2  &-4  &1   &0  \\
              0  &0   &1/2 &1
            \end{pmat}
          \end{equation*}
发现左半部分不能化为单位矩阵，所以逆不存在。
核对 $ad-bc=2\cdot 2-(-4)\cdot (-1)=0$，结论一致。
\partsitem 以下计算得到逆矩阵。
          \begin{multline*}
            \begin{pmat}{rrr|rrr}
              1  &1  &3  &1  &0  &0  \\
              0  &2  &4  &0  &1  &0  \\
             -1  &1  &0  &0  &0  &1
            \end{pmat}
            \grstep{\rho_1+\rho_3}
            \begin{pmat}{rrr|rrr}
              1  &1  &3  &1  &0  &0  \\
              0  &2  &4  &0  &1  &0  \\
              0  &2  &3  &1  &0  &1
            \end{pmat}                                             \\
           \begin{aligned}
            &\grstep{-\rho_2+\rho_3}
            \begin{pmat}{rrr|rrr}
              1  &1  &3  &1  &0  &0  \\
              0  &2  &4  &0  &1  &0  \\
              0  &0  &-1 &1  &-1 &1
            \end{pmat}
           \grstep[-\rho_3]{(1/2)\rho_2}
            \begin{pmat}{rrr|rrr}
              1  &1  &3  &1  &0   &0  \\
              0  &1  &2  &0  &1/2 &0  \\
              0  &0  &1  &-1 &1   &-1
            \end{pmat}                                        \\
            &\grstep[-3\rho_3+\rho_1]{-2\rho_3+\rho_2}
            \begin{pmat}{rrr|rrr}
              1  &1  &0  &4  &-3   &3  \\
              0  &1  &0  &2  &-3/2 &2  \\
              0  &0  &1  &-1 &1    &-1
            \end{pmat}                                         \\
            &\grstep{-\rho_2+\rho_1}
            \begin{pmat}{rrr|rrr}
              1  &0  &0  &2  &-3/2 &1  \\
              0  &1  &0  &2  &-3/2 &2  \\
              0  &0  &1  &-1 &1    &-1
            \end{pmat}
           \end{aligned}
          \end{multline*}
\partsitem 以下是一种消元方法。
          \begin{multline*}
            \begin{pmat}{rrr|rrr}
              0  &1  &5  &1  &0  &0  \\
              0  &-2 &4  &0  &1  &0  \\
              2  &3  &-2 &0  &0  &1
            \end{pmat}
            \grstep{\rho_3\leftrightarrow\rho_1}\;
            \begin{pmat}{rrr|rrr}
              2  &3  &-2 &0  &0  &1  \\
              0  &-2 &4  &0  &1  &0  \\
              0  &1  &5  &1  &0  &0
            \end{pmat}                                            \\
            \begin{aligned}
              &\grstep{(1/2)\rho_2+\rho_3}
              \begin{pmat}{rrr|rrr}
                2  &3  &-2 &0  &0   &1  \\
                0  &-2 &4  &0  &1   &0  \\
                0  &0  &7  &1  &1/2 &0
              \end{pmat}                                             \\
              &\grstep[-(1/2)\rho_2 \\ (1/7)\rho_3]{(1/2)\rho_1}
              \begin{pmat}{rrr|rrr}
                1  &3/2  &-1 &0    &0     &1/2  \\
                0  &1    &-2 &0    &-1/2  &0    \\
                0  &0    &1  &1/7  &1/14  &0
              \end{pmat}                                   \\
              &\grstep[\rho_3+\rho_1]{2\rho_3+\rho_2}
              \begin{pmat}{rrr|rrr}
                1  &3/2  &0  &1/7  &1/14  &1/2  \\
                0  &1    &0  &2/7  &-5/14 &0    \\
                0  &0    &1  &1/7  &1/14  &0
              \end{pmat}                                   \\
              &\grstep{-(3/2)\rho_2+\rho_1}
              \begin{pmat}{rrr|rrr}
                1  &0    &0  &-2/7 &17/28 &1/2  \\
                0  &1    &0  &2/7  &-5/14 &0    \\
                0  &0    &1  &1/7  &1/14  &0
              \end{pmat}
            \end{aligned}
          \end{multline*}
\partsitem 逆不存在。
          \begin{align*}
            \begin{pmat}{rrr|rrr}
              2  &2  &3  &1  &0  &0  \\
              1  &-2 &-3 &0  &1  &0  \\
              4  &-2 &-3 &0  &0  &1
            \end{pmat}
            &\grstep[-2\rho_1+\rho_3]{-(1/2)\rho_1+\rho_2}
            \begin{pmat}{rrr|rrr}
              2  &2  &3    &1     &0  &0  \\
              0  &-3 &-9/2 &-1/2  &1  &0  \\
              0  &-6 &-9   &-2    &0  &1
            \end{pmat}                                          \\
            &\grstep{-2\rho_2+\rho_3}
            \begin{pmat}{rrr|rrr}
              2  &2  &3    &1     &0   &0  \\
              0  &-3 &-9/2 &-1/2  &1   &0  \\
              0  &0  &0    &-1    &-2  &1
            \end{pmat}
          \end{align*}
核对可知，原矩阵的第三列是第二列的 $3/2$ 倍，确实奇异，因此无逆。
      \end{exparts}
    
\end{ans}
\begin{ans}{三.IV.4.16}
可使用 \nearbycorollary{cor:TwoByTwoInv}。
       \begin{equation*}
         \frac{1}{1\cdot 5-2\cdot 3}\cdot
         \begin{mat}[r]
           5  &-3  \\
          -2  &1
         \end{mat}
         =\begin{mat}[r]
           -5  &3  \\
            2  &-1
         \end{mat}
       \end{equation*}
      
\end{ans}
\begin{ans}{三.IV.4.17}
      \begin{exparts}
\partsitem 容易证明逆为 \( r^{-1}H^{-1}=(1/r)\cdot H^{-1} \)（当然，须假定矩阵可逆）。
\partsitem 不成立。首先，$H+G$ 有逆并不意味着 $H$ 或 $G$ 有逆。以下两个矩阵均不可逆，但其和可逆。
          \begin{equation*}
            \begin{mat}[r]
              1  &0  \\
              0  &0
            \end{mat}
            \qquad
            \begin{mat}[r]
              0  &0  \\
              0  &1
            \end{mat}
          \end{equation*}
其次，$H$ 和 $G$ 各自有逆，也不意味着 $H+G$ 有逆，例如
          \begin{equation*}
            \begin{mat}[r]
              1  &0  \\
              0  &1
            \end{mat}
            \qquad
            \begin{mat}[r]
              -1  &0  \\
               0  &-1
            \end{mat}
          \end{equation*}
第三，即使两个矩阵及其和都有逆，也不意味着等式成立：
          \begin{equation*}
            \begin{mat}[r]
              2  &0  \\
              0  &2
            \end{mat}^{-1}
            =\begin{mat}[r]
              1/2  &0  \\
              0    &1/2
            \end{mat}^{-1}
            \qquad
            \begin{mat}[r]
              3  &0  \\
              0  &3
            \end{mat}^{-1}
            =\begin{mat}[r]
              1/3  &0  \\
              0    &1/3
            \end{mat}^{-1}
          \end{equation*}
但
          \begin{equation*}
            \begin{mat}[r]
              5  &0  \\
              0  &5
            \end{mat}^{-1}
            =\begin{mat}[r]
              1/5  &0  \\
              0    &1/5
            \end{mat}^{-1}
          \end{equation*}
而 $(1/2)_+(1/3)$ 不等于 $1/5$。
      \end{exparts}
    
\end{ans}
\begin{ans}{三.IV.4.18}
成立：\( T^k(T^{-1})^k=(TT\cdots T)\cdot (T^{-1}T^{-1}\cdots T^{-1})=T^{k-1}(TT^{-1})(T^{-1})^{k-1}=\dots=I \)。
    
\end{ans}
\begin{ans}{三.IV.4.19}
可逆，\( H^{-1} \) 的逆就是 \( H \)。
    
\end{ans}
\begin{ans}{三.IV.4.20}
第一个等式可用三角函数的和角公式验证。
      \begin{multline*}
        \begin{mat}
          \cos(\theta_1+\theta_2) &-\sin(\theta_1+\theta_2)  \\
          \sin(\theta_1+\theta_2) &\cos(\theta_1+\theta_2)
        \end{mat}                                                 \\
        \begin{aligned}
        &=\begin{mat}
          \cos\theta_1\cos\theta_2-\sin\theta_1\sin\theta_2
            &-\sin\theta_1\cos\theta_2-\cos\theta_1\sin\theta_2  \\
          \sin\theta_1\cos\theta_2+\cos\theta_1\sin\theta_2
            &\cos\theta_1\cos\theta_2-\sin\theta_1\sin\theta_2
        \end{mat}                                                    \\
        &=\begin{mat}
          \cos\theta_1 &-\sin\theta_1  \\
          \sin\theta_1 &\cos\theta_1
        \end{mat}
        \begin{mat}
          \cos\theta_2 &-\sin\theta_2  \\
          \sin\theta_2 &\cos\theta_2
        \end{mat}
        \end{aligned}
      \end{multline*}
第二个等式也可类似地验证。

当然，回顾 $t_\theta$ 是将向量绕原点旋转 \( \theta \) 弧度的映射，就不仅能验证这些等式，还能理解它们。
    
\end{ans}
\begin{ans}{三.IV.4.21}
分两种情形。第一种假设 \( a \) 非零，则
      \begin{equation*}
        \grstep{-(c/a)\rho_1+\rho_2}
        \begin{pmat}{cc|cc}
           a   &b           &1     &0   \\
           0   &-(bc/a)+d   &-c/a  &1
         \end{pmat}
        =\begin{pmat}{cc|cc}
           a   &b           &1     &0   \\
           0   &(ad-bc)/a   &-c/a  &1
         \end{pmat}
      \end{equation*}
说明在 \( a\neq 0 \) 时，矩阵可逆当且仅当 \( ad-bc\neq 0 \)。为求逆，继续完成若尔当消元部分。
      \begin{align*}
        &\grstep[(a/ad-bc)\rho_2]{(1/a)\rho_1}
        \begin{pmat}{cc|cc}
           1   &b/a     &1/a         &0   \\
           0   &1       &-c/(ad-bc)  &a/(ad-bc)
         \end{pmat}                                       \\
        &\grstep{-(b/a)\rho_2+\rho_1}
        \begin{pmat}{cc|cc}
           1   &0       &d/(ad-bc)   &-b/(ad-bc)   \\
           0   &1       &-c/(ad-bc)  &a/(ad-bc)
         \end{pmat}
      \end{align*}

另一种情形为 \( a=0 \)。交换两行，将 $c$ 移到 $1,1$ 位置。
      \begin{equation*}
        \grstep{\rho_1\leftrightarrow\rho_2}
        \begin{pmat}{cc|cc}
          c  &d  &0  &1  \\
          0  &b  &1  &0
        \end{pmat}
      \end{equation*}
此矩阵非奇异当且仅当 \( b \)、\( c \) 均非零（在 \( a=0 \) 的假设下，等价于 \( ad-bc\neq 0 \)）。
为求逆，完成若尔当消元部分。
      \begin{equation*}
        \grstep[(1/b)\rho_2]{(1/c)\rho_1}
        \begin{pmat}{cc|cc}
          1  &d/c  &0       &1/c  \\
          0  &1    &1/b     &0
        \end{pmat}
        \grstep{-(d/c)\rho_2+\rho_1}
        \begin{pmat}{cc|cc}
          1  &0  &-d/bc  &1/c  \\
          0  &1  &1/b    &0
        \end{pmat}
      \end{equation*}
（注意，这正是所需结果，因为 \( a=0 \) 时 \( ad-bc=-bc \)。）
    
\end{ans}
\begin{ans}{三.IV.4.22}
$H$ 为 $\nbym{2}{3}$ 矩阵。要找到 $G$ 使 $HG$ 为 $\nbyn{2}$ 单位矩阵，$G$ 必须为 $\nbym{3}{2}$。
列出方程
      \begin{equation*}
          \begin{mat}[r]
             1  &0   &1  \\
             0  &1   &0
           \end{mat}
          \begin{mat}
             m  &n  \\
             p  &q  \\
             r  &s
          \end{mat}
        =
          \begin{mat}[r]
             1  &0  \\
             0  &1
          \end{mat}
      \end{equation*}
并求解相应线性方程组，
      \begin{equation*}
        \begin{linsys}{6}
           m  &   &   &   &+r &   &=    &1  \\
              &n  &   &   &   &+s &=    &0  \\
              &   &p  &   &   &   &=    &0  \\
              &   &   &q  &   &   &=    &1
         \end{linsys}
      \end{equation*}
得到无穷多个解。
      \begin{equation*}
        \set{\colvec{m \\ n \\ p \\ q \\ r \\ s}
              =\colvec[r]{1 \\ 0 \\ 0 \\ 1 \\ 0 \\ 0}
              +r\cdot \colvec[r]{-1 \\ 0 \\ 0 \\ 0 \\ 1 \\ 0}
              +s\cdot \colvec[r]{0 \\ -1 \\ 0 \\ 0 \\ 0 \\ 1}
             \suchthat r,s\in\Re}
      \end{equation*}
所以 $H$ 有无穷多个右逆。

对于左逆，方程
      \begin{equation*}
         \begin{mat}
            a  &b  \\
            c  &d
         \end{mat}
         \begin{mat}[r]
             1  &0   &1  \\
             0  &1   &0
          \end{mat}
          =\begin{mat}[r]
             1  &0  &0  \\
             0  &1  &0  \\
             0  &0  &1
           \end{mat}
      \end{equation*}
给出一个含九个方程、四个未知数的线性方程组。
      \begin{equation*}
        \begin{linsys}{6}
          a & & & & & & & & & & &= &1 \\
            & &b& & & & & & & & &= &0 \\
          a & & & & & & & & & & &= &0 \\
            & & & &c& & & & & & &= &0 \\
            & & & & & &d& & & & &= &1 \\
            & & & &c& & & & & & &= &0 \\
            & & & & & & & &e& & &= &0 \\
            & & & & & & & & & &f&= &0 \\
            & & & & & & & &e& & &= &1
        \end{linsys}
      \end{equation*}
该方程组不相容（第一、第三个方程冲突，第七、第九个也冲突），所以没有左逆。
    
\end{ans}
\begin{ans}{三.IV.4.23}
相对于标准基，有
      \begin{equation*}
        \rep{\iota}{\stdbasis_2,\stdbasis_3}
        =\begin{mat}[r]
          1  &0  \\
          0  &1  \\
          0  &0
        \end{mat}
      \end{equation*}
列出求左逆矩阵的方程，
      \begin{equation*}
        \begin{mat}
          a &b &c \\
          d &e &f
        \end{mat}
        \begin{mat}[r]
          1 &0  \\
          0 &1  \\
          0 &0
        \end{mat}
        =\begin{mat}[r]
          1 &0  \\
          0 &1
        \end{mat}
        =\rep{\identity}{\stdbasis_2,\stdbasis_2}
      \end{equation*}
得到线性方程组。
      \begin{equation*}
        \begin{linsys}{6}
          a & & & & & & & & & & &= &1 \\
            & &b& & & & & & & & &= &0 \\
            & & & & & &d& & & & &= &0 \\
            & & & & & & & &e& & &= &1
        \end{linsys}
      \end{equation*}
关于 $a,\ldots,f$，方程组有无穷多个解，因为其中两个变量完全不受限制，
      \begin{equation*}
        \set{\colvec{a \\ b \\ c \\ d \\ e \\ f}
             =\colvec[r]{1 \\ 0 \\ 0 \\ 0 \\ 1 \\ 0}
              +c\cdot \colvec[r]{0 \\ 0 \\ 1 \\ 0 \\ 0 \\ 0}
              +f\cdot \colvec[r]{0 \\ 0 \\ 0 \\ 0 \\ 0 \\ 1}
             \suchthat c,f\in\Re}
      \end{equation*}
所以矩阵方程有无穷多个解。
      \begin{equation*}
        \set{\begin{mat}
               1  &0  &c  \\
               0  &1  &f
             \end{mat}
             \suchthat c,f\in\Re}
      \end{equation*}
基仍固定为 $\stdbasis_2,\stdbasis_2$，例如取 $c=2$、$f=3$，所得矩阵表示以下映射。
      \begin{equation*}
        \colvec{x \\ y \\ z}
        \;\mapsunder{f_{2,3}}\;
        \colvec{x+2z \\ y+3z}
      \end{equation*}
容易验证 $\composed{f_{2,3}}{\iota}$ 是 $\Re^2$ 上的恒等映射。
    
\end{ans}
\begin{ans}{三.IV.4.24}
由 \nearbylemma{le:LeftAndRightInvEqual}，不可能有无穷多个左逆，因为同时具有左右逆的矩阵，两种逆各自唯一，而且相等\Dash 即同时是左逆和右逆。
    
\end{ans}
\begin{ans}{三.IV.4.25}
      \begin{exparts}
\partsitem 真。它必须线性，如定理~II.\ref{th:OOHomoEquivalence} 的证明所示。
\partsitem 假。它可能线性，但不一定。考虑本小节开头的投影 $\map{\pi}{\Re^3}{\Re^2}$，定义 $\map{\eta}{\Re^2}{\Re^3}$ 如下。
          \begin{equation*}
            \colvec{x  \\  y}\mapsto\colvec{x  \\ y \\ 1}
          \end{equation*}
它是 $\pi$ 的右逆，因为 $\composed{\pi}{\eta}$ 的作用为
          \begin{equation*}
            \colvec{x  \\  y}\mapsto\colvec{x  \\ y \\ 1}\mapsto\colvec{x \\ y}
          \end{equation*}
但它不线性，因为没有将零向量映到零向量。
      \end{exparts}
    
\end{ans}
\begin{ans}{三.IV.4.26}
由矩阵乘法结合律，\( H^{-1}(HG)=H^{-1}Z=Z \)，同时
\( H^{-1}(HG)=(H^{-1}H)G=IG=G \)。
    
\end{ans}
\begin{ans}{三.IV.4.27}
在第一个等式两边分别乘 $H$。
    
\end{ans}
\begin{ans}{三.IV.4.28}
在 $T^4$ 为零矩阵的假设下，容易验证该表达式从左边或右边乘 $I-T$，结果均为单位矩阵。
显然可推广为：若 \( T^n \) 为零矩阵，则 \( (I-T)^{-1}=I+T+T^2+\cdots+T^{n-1} \)，同样容易验证。
    
\end{ans}
\begin{ans}{三.IV.4.29}
矩阵的幂由对角元素的相应幂构成。即 \( D^2 \) 除对角元素 \( {d_{1,1}}^2 \)、\( {d_{2,2}}^2 \) 等以外，其余均为零。
这提示将 \( D^0 \) 定义为单位矩阵。
    
\end{ans}
\begin{ans}{三.IV.4.30}
设 $B$ 与 $A$ 行等价，且 $A$ 可逆。由行等价，存在一系列行运算将一者化为另一者。
这些行运算可用矩阵实现，例如 $B=R_n\cdots R_1A$。
此式将 $B$ 表为可逆矩阵的乘积，由 \nearbylemma{lem:ProdInvIsInv}，$B$ 也可逆。
    
\end{ans}
\begin{ans}{三.IV.4.31}
      \begin{exparts}
\partsitem 参见 \nearbyexercise{exer:RankProdLeqRankFacts} 的答案。
\partsitem 我们将说明两个条件都等价于两矩阵非奇异。

由于 \( T \)、\( S \) 为方阵且乘积有定义，它们同阶，设均为 \( \nbyn{n} \)。
考虑 \( TS=I \)。由上一问，\( I \) 的秩不超过 \( T \)、\( S \) 秩的最小值。
但 \( I \) 的秩为 \( n \)，所以 \( T \)、\( S \) 的秩均为 \( n \)，因而均非奇异。

同样的论证说明 \( ST=I \) 蕴含两者均非奇异。
      \end{exparts}
    
\end{ans}
\begin{ans}{三.IV.4.32}
逆唯一，所以只需验证转置满足要求。该验证见前面的 \nearbyexercise{exer:PermTimesTransEqId}。
    
\end{ans}
\begin{ans}{三.IV.4.33}
      \begin{exparts}
\partsitem 参见 \nearbyexercise{exer:TranspAndMult} 的答案。
\partsitem 参见 \nearbyexercise{exer:TranspAndMult} 的答案。
\partsitem 对 \( I=AA^{-1} \) 应用第一问，得 $I=\trans{I}=\trans{(AA^{-1})}=\trans{(A^{-1})}\trans{A}$。
\partsitem 由于 \( A \) 对称，在上一问中使用 \( \trans{A}=A \)。
      \end{exparts}
    
\end{ans}
\begin{ans}{三.IV.4.34}
前半部分各问的答案见 \nearbyexercise{exer:ZeroDivisor}。
对于后半部分，设 $A$ 是零因子，则存在非零矩阵 $B$ 使 $AB=Z$（或 $BA=Z$，论证类似）。
若 $A$ 可逆，则 $A^{-1}(AB)=(A^{-1}A)B=IB=B$，同时 $A^{-1}(AB)=A^{-1}Z=Z$，与 $B$ 非零矛盾。
    
\end{ans}
\begin{ans}{三.IV.4.35}
平方为单位矩阵的 $\nbyn{2}$ 矩阵 $T$ 有无穷多个。以下是四个。
      \begin{equation*}
        \begin{mat}[r]
          1  &0  \\
          0      &1
        \end{mat}
        \quad
        \begin{mat}[r]
          -1  &0  \\
          0   &-1
        \end{mat}
        \quad
        \begin{mat}[r]
          1  &0  \\
          0  &-1
        \end{mat}
        \quad
        \begin{mat}[r]
          -1  &0  \\
          0    &1
        \end{mat}
      \end{equation*}
另外两个为
      \begin{equation*}
        \begin{mat}[r]
          0  &1  \\
          1  &0
        \end{mat}
        \quad
        \begin{mat}[r]
          0  &-1  \\
          -1 &0
        \end{mat}
      \end{equation*}
还有一个为
      \begin{equation*}
        \frac{1}{5}
        \begin{mat}[r]
          3  &4  \\
          4  &-3
        \end{mat}
      \end{equation*}
（最后一个例子的规律涉及勾股三元组，即满足 $r^2+s^2=t^2$ 的数 $r,s,t\in\R$）。
\textit{评注：}另见 \url{https://en.wikipedia.org/wiki/Square_root_of_a_2_by_2_matrix}。
    
\end{ans}
\begin{ans}{三.IV.4.36}
它不具有自反性，例如
     \begin{equation*}
       H=\begin{mat}[r]
         1  &0  \\
         0  &2
       \end{mat}
     \end{equation*}
就不是自身的双侧逆。同一例子还说明不具有传递性。该矩阵的双侧逆为
     \begin{equation*}
       G=\begin{mat}[r]
         1  &0  \\
         0  &1/2
       \end{mat}
     \end{equation*}
虽然 \( H \) 是 \( G \) 的双侧逆，\( G \) 是 \( H \) 的双侧逆，但 \( H \) 不是自身的双侧逆。
不过该关系具有对称性：若 \( G \) 是 \( H \) 的双侧逆，则 \( GH=I=HG \)，所以 \( H \) 也是 \( G \) 的双侧逆。
   
\end{ans}
\begin{ans}{三.IV.4.37}
\answerasgiven %
设 \( A \) 为满足所述性质的 \( \nbyn{m} \) 非奇异矩阵，\( B \) 为其逆。对 \( n\leq m \)，有
      \begin{equation*}
        1
        =\sum_{r=1}^{m}\delta_{nr}
        =\sum_{r=1}^{m}\sum_{s=1}^{m}b_{ns}a_{sr}
        =\sum_{s=1}^{m}\sum_{r=1}^{m}b_{ns}a_{sr}
        =k\sum_{s=1}^{m}b_{ns}
      \end{equation*}
（若 \( k=0 \)，则 \( A \) 奇异）。
   
\end{ans}
\protect \section {第 V 节：换基}
\protect \subsection {三.V.1: 改变向量的表示}
\begin{ans}{三.V.1.7}
为求从 $D$ 到 $\stdbasis_2$ 的换基矩阵，需要
$\rep{\identity(\vec{\delta}_1)}{\stdbasis_2}=\rep{\vec{\delta}_1}{\stdbasis_2}$ 和
$\rep{\identity(\vec{\delta}_2)}{\stdbasis_2}=\rep{\vec{\delta}_2}{\stdbasis_2}$。
当然，$\Re^2$ 中向量相对于标准基的表示很容易求。
      \begin{equation*}
        \rep{\vec{\delta}_1}{\stdbasis_2}=\colvec[r]{2 \\ 1}
        \qquad
        \rep{\vec{\delta}_2}{\stdbasis_2}=\colvec[r]{-2 \\ 4}
      \end{equation*}
将它们并排拼接为换基矩阵的列，得到
      \begin{equation*}
        \rep{\identity}{D,\stdbasis_2}
        =\begin{mat}[r]
          2     &-2    \\
          1     &4
        \end{mat}
      \end{equation*}
反方向的换基矩阵，可按常规计算
$\rep{\identity(\vec{e}_1)}{D}=\rep{\vec{e}_1}{D}$ 和
$\rep{\identity(\vec{e}_2)}{D}=\rep{\vec{e}_2}{D}$，也可求上面矩阵的逆。
利用 $\nbyn{2}$ 矩阵的求逆公式很容易得到。
      \begin{equation*}
        \rep{\identity}{\stdbasis_2,D}=
        \frac{1}{10}\cdot\begin{mat}[r]
          4  &2  \\
         -1  &2
        \end{mat}
        =\begin{mat}[r]
          4/10  &2/10  \\
         -1/10  &2/10
        \end{mat}
      \end{equation*}
    
\end{ans}
\begin{ans}{三.V.1.8}
非奇异矩阵都可以作为换基矩阵。以下各例均可直接观察判断。
     \begin{exparts}
\partsitem 非奇异，因为第二行不是第一行的倍数。
\partsitem 非奇异。
\partsitem 奇异，因为非奇异矩阵必须是方阵。
\partsitem 奇异，因为第一行的两倍加第二行等于第三行。
\partsitem 非奇异。
     \end{exparts}
   
\end{ans}
\begin{ans}{三.V.1.9}
将 $\rep{\identity(\vec{\beta}_1)}{D}=\rep{\vec{\beta}_1}{D}$ 与
$\rep{\identity(\vec{\beta}_2)}{D}=\rep{\vec{\beta}_2}{D}$ 并排拼接，得到换基矩阵 $\rep{\identity}{B,D}$。
      \begin{exparts*}
        \partsitem \( \begin{mat}[r]
                   0  &1  \\
                   1  &0
                 \end{mat} \)
        \partsitem \( \begin{mat}[r]
                   2  &-1/2  \\
                   -1 &1/2
                 \end{mat} \)
        \partsitem \( \begin{mat}[r]
                   1  &1     \\
                   2  &4
                 \end{mat} \)
        \partsitem \( \begin{mat}[r]
                   1  &-1    \\
                  -1  &2
                 \end{mat} \)
      \end{exparts*}
    
\end{ans}
\begin{ans}{三.V.1.10}
向量 $\rep{\identity(\vec{\beta}_1)}{D}=\rep{\vec{\beta}_1}{D}$、
$\rep{\identity(\vec{\beta}_2)}{D}=\rep{\vec{\beta}_2}{D}$ 和
$\rep{\identity(\vec{\beta}_3)}{D}=\rep{\vec{\beta}_3}{D}$ 构成换基矩阵 $\rep{\identity}{B,D}$ 的各列。
      \begin{exparts*}
        \partsitem \( \begin{mat}[r]
                   0  &0  &1  \\
                   1  &0  &0  \\
                   0  &1  &0
                 \end{mat} \)
        \partsitem \( \begin{mat}[r]
                   1  &-1 &0  \\
                   0  &1  &-1 \\
                   0  &0  &1
                 \end{mat} \)
        \partsitem \( \begin{mat}[r]
                   1  &-1 &1/2  \\
                   1  &1  &-1/2 \\
                   0  &2  &0
                 \end{mat} \)
      \end{exparts*}
例如，第一个矩阵的第一列来自 $1=0\cdot x^2+1\cdot 1+0\cdot x$。
     
\end{ans}
\begin{ans}{三.V.1.11}
一种方法是求 $\rep{\vec{\delta}_1}{B}$ 和 $\rep{\vec{\delta}_2}{B}$，再将它们作为列拼成所需换基矩阵。
另一种方法是求 \nearbyexercise{exer:ChngeBasMatsRTwo} 答案中矩阵的逆。
       \begin{exparts*}
        \partsitem
          $\begin{mat}[r]
            0  &1  \\
            1  &0
          \end{mat}$
        \partsitem \( \begin{mat}[r]
            1  &1  \\
            2  &4
          \end{mat} \)
        \partsitem \( \begin{mat}[r]
            2  &-1/2  \\
            -1 &1/2
          \end{mat} \)
        \partsitem \( \begin{mat}[r]
            2  &1  \\
            1  &1
          \end{mat} \)
      \end{exparts*}
    
\end{ans}
\begin{ans}{三.V.1.12}
矩阵实现换基当且仅当它非奇异。
       \begin{exparts}
\partsitem 该矩阵非奇异，所以能换基。求 \( \stdbasis_2 \) 换成的基，就是寻找 $D$ 使
          \begin{equation*}
            \rep{\identity}{\stdbasis_2,D}=
            \begin{mat}[r]
                 5  &0  \\
                 0  &4
            \end{mat}
          \end{equation*}
由矩阵表示线性映射的定义，有
          \begin{equation*}
            \rep{\identity(\vec{e}_1)}{D}=\rep{\vec{e}_1}{D}=\colvec[r]{5 \\ 0}
            \qquad
            \rep{\identity(\vec{e}_2)}{D}=\rep{\vec{e}_2}{D}=\colvec[r]{0 \\ 4}
          \end{equation*}
设
          \begin{equation*}
            D=\sequence{\colvec{x_1 \\ y_1},\colvec{x_2 \\ y_2}}
          \end{equation*}
可解方程组
          \begin{equation*}
            \colvec[r]{1 \\ 0}
            =5\colvec{x_1 \\ y_1}+0\colvec{x_2 \\ y_1}
            \qquad
            \colvec[r]{0 \\ 1}
            =0\colvec{x_1 \\ y_1}+4\colvec{x_2 \\ y_1}
          \end{equation*}
也可直接观察得到答案（参照 \nearbylemma{le:NonSingIsChBasis} 的证明）。
          \begin{equation*}
            D=\sequence{\colvec[r]{1/5 \\ 0},
                        \colvec[r]{0 \\ 1/4}}
          \end{equation*}
\partsitem 能，该矩阵非奇异。为求 $D$，与上面一样，从
          \begin{equation*}
            D=\sequence{\colvec{x_1 \\ y_1},
                        \colvec{x_2 \\ y_2}}
          \end{equation*}
出发求解
          \begin{equation*}
            \colvec[r]{1 \\ 0}
            =2\colvec{x_1 \\ y_1}+3\colvec{x_2 \\ y_1}
            \quad\text{and}\quad
            \colvec[r]{0 \\ 1}
            =1\colvec{x_1 \\ y_1}+1\colvec{x_2 \\ y_1}
          \end{equation*}
得到
          \begin{equation*}
            D=\sequence{\colvec[r]{-1 \\ 3},
                        \colvec[r]{1 \\ -2}}
          \end{equation*}
\partsitem 不能，因为矩阵奇异。
\partsitem 能，因为矩阵非奇异。终止基的计算与上面相同。
          \begin{equation*}
            D=\sequence{\colvec[r]{1/2 \\ -1/2},
                        \colvec[r]{1/2 \\ 1/2}}
          \end{equation*}
      \end{exparts}
    
\end{ans}
\begin{ans}{三.V.1.13}
      \begin{exparts}
        \partsitem
先求恒等函数在起始基~$B$ 各元素上的作用，显然如下。
          \begin{equation*}
            \colvec{1 \\ 0 \\ 0}\mapsunder{\identity}\colvec{1 \\ 0 \\ 0}
            \quad
            \colvec{0 \\ 1 \\ 0}\mapsunder{\identity}\colvec{0 \\ 1 \\ 0}
            \quad
            \colvec{0 \\ 0 \\ 1}\mapsunder{\identity}\colvec{0 \\ 0 \\ 1}
          \end{equation*}
再将三个输出相对于终止基表示出来。
          \begin{equation*}
            \rep{\colvec{1 \\ 0 \\ 0}}{D}=\colvec{-2/3 \\ 5/3 \\ -1/3}
            \quad
            \rep{\colvec{0 \\ 1 \\ 0}}{D}=\colvec{1/3 \\ -1/3 \\ 2/3}
            \quad
            \rep{\colvec{0 \\ 0 \\ 1}}{D}=\colvec{1/3 \\ -1/3 \\ -1/3}
          \end{equation*}
将它们并排拼接成矩阵。
          \begin{equation*}
            \rep{\identity}{B,D}=
            \begin{mat}
              -2/3 &1/3  &1/3  \\
               5/3 &-1/3 &-1/3 \\
              -1/3 &2/3  &-1/3
            \end{mat}
          \end{equation*}
        \partsitem
一种方法是求上一问矩阵的逆，因为它沿反方向换基。
也可计算以下三个表示，
          \begin{equation*}
            \rep{\colvec{1 \\ 2 \\ 3}}{\stdbasis_3}
              =\colvec{1 \\ 2 \\ 3}
            \quad
            \rep{\colvec{1 \\ 1 \\ 1}}{\stdbasis_3}
              =\colvec{1 \\ 1 \\ 1}
            \quad
            \rep{\colvec{0 \\ 1 \\ -1}}{\stdbasis_3}
               =\colvec{0 \\ 1 \\ -1}
          \end{equation*}
再将它们拼成矩阵。
          \begin{equation*}
            \rep{\identity}{B,D}=
            \begin{mat}
              1 &1 &0 \\
              2 &1 &1 \\
              3 &1 &-1
            \end{mat}
          \end{equation*}
        \partsitem
相对于终止基表示 $\identity(x^2)$、$\identity(x^2+x)$ 和~$\identity(x^2+x+1)$，得到
          \begin{equation*}
            \rep{x^2}{D}=\colvec{0 \\ 0 \\ 1}
            \quad
            \rep{x^2+x}{D}=\colvec{0 \\ -1 \\ 1}
            \quad
            \rep{x^2+x+1}{D}=\colvec{1/2 \\ -1 \\ 1}
          \end{equation*}
将它们拼接。
          \begin{equation*}
            \rep{\identity}{B,D}=
            \begin{mat}
              0 &0 &1/2 \\
              0 &-1 &-1 \\
              1 &1  &1
            \end{mat}
          \end{equation*}
      \end{exparts}
    
\end{ans}
\begin{ans}{三.V.1.14}
本题有许多不同答案。一种方法是在任意空间中任取基 $B$，再计算适当的 $D$（当然属于同一空间）。
更简单的方法是取陪域为 $\Re^3$，陪域的基为 $\stdbasis_3$。
由于任意向量相对于标准基的表示就是它自身，得到
      \begin{equation*}
        B=\sequence{\colvec[r]{3 \\ 2 \\ 0},
                    \colvec[r]{1 \\ -1 \\ 0},
                    \colvec[r]{4 \\ 1 \\ 4}  }
        \qquad
        D=\stdbasis_3
      \end{equation*}
    
\end{ans}
\begin{ans}{三.V.1.15}
按常规方法可验证 \( B=\sequence{2\sin(x)+\cos(x),3\cos(x)} \) 为基。将自然基记为 $D$。
为求换基矩阵 $\rep{\identity}{B,D}$，必须求 $\rep{2\sin(x)+\cos(x)}{D}$ 和 $\rep{3\cos(x)}{D}$，
即寻找 $x_1,y_1,x_2,y_2$ 使以下等式成立。
      \begin{align*}
        x_1\cdot \sin(x)+y_1\cdot\cos(x) &= 2\sin(x)+\cos(x) \\
        x_2\cdot \sin(x)+y_2\cdot\cos(x) &= 3\cos(x)
      \end{align*}
显然，答案如下。
      \begin{equation*}
        \rep{\identity}{B,D}
        =\begin{mat}[r]
          2  &0  \\
          1  &3
        \end{mat}
      \end{equation*}
反方向的换基矩阵可通过求解以下方程，得到 $\rep{\sin(x)}{B}$ 和 $\rep{\cos(x)}{B}$。
      \begin{align*}
        w_1\cdot (2\sin(x)+\cos(x))+z_1\cdot(3\cos(x)) &= \sin(x) \\
        w_2\cdot (2\sin(x)+\cos(x))+z_2\cdot(3\cos(x)) &= \cos(x)
      \end{align*}
更简单的方法是求上面矩阵的逆。
      \begin{equation*}
        \rep{\identity}{D,B}
        =\begin{mat}[r]
          2  &0  \\
          1  &3
        \end{mat}^{-1}
        =\frac{1}{6}\cdot\begin{mat}[r]
          3  &0  \\
          -1  &2
        \end{mat}
        =\begin{mat}[r]
           1/2  &0  \\
          -1/6  &1/3
         \end{mat}
      \end{equation*}
    
\end{ans}
\begin{ans}{三.V.1.16}
先求矩阵的逆，即确定所考察映射的逆。
      \begin{align*}
        \rep{\identity}{D,\stdbasis_2}
        \rep{\identity}{D,\stdbasis_2}^{-1}
        &=\frac{1}{-\cos^2(2\theta)-\sin^2(2\theta)}\cdot\begin{mat}
           -\cos(2\theta)   &-\sin(2\theta)  \\
           -\sin(2\theta)   &\cos(2\theta)
        \end{mat}                                         \\
        &=\begin{mat}
           \cos(2\theta)   &\sin(2\theta)  \\
           \sin(2\theta)   &-\cos(2\theta)
        \end{mat}
      \end{align*}
这比另一方向的表示更容易处理，因为该矩阵由以下两个列向量并排拼成，
      \begin{equation*}
        \rep{\vec{\delta}_1}{\stdbasis_2}
           =\colvec{\cos(2\theta) \\ \sin(2\theta)}
        \qquad
        \rep{\vec{\delta}_2}{\stdbasis_2}
           =\colvec{\sin(2\theta) \\ -\cos(2\theta)}
      \end{equation*}
而相对于 $\stdbasis_2$ 的表示是直接的。
      \begin{equation*}
        \vec{\delta}_1=\colvec{\cos(2\theta) \\ \sin(2\theta)}
        \qquad
        \vec{\delta}_2=\colvec{\sin(2\theta) \\ -\cos(2\theta)}
      \end{equation*}
      该图展示将 $D$ 变换为 $\stdbasis_2$ 的映射作用
      (it is, again, the inverse of the map that is the answer to
      此问题）。
      该直线与 $x$ 轴夹角为 $\theta$。
      \begin{center}  \small
        \includegraphics{map/mp/ch3.81}
      \end{center}
这个映射将向量关于该直线反射。反射的逆是自身，所以答案为：原映射是关于过原点、倾角为 $\theta$ 的直线的反射。
（当然，对任意基都是如此。）
    
\end{ans}
\begin{ans}{三.V.1.17}
大小适当的单位矩阵。
     
\end{ans}
\begin{ans}{三.V.1.18}
两者都等价于矩阵非奇异。
    
\end{ans}
\begin{ans}{三.V.1.19}
剩下需要证明：左乘消元矩阵表示从另一组基换到 \( B=\basis{\beta}{n} \)。

行倍乘矩阵 \( M_i(k) \) 将相对于基
\( \sequence{\vec{\beta}_1,\dots,k\vec{\beta}_i,\dots,\vec{\beta}_n} \) 的表示转换为相对于 \( B \) 的表示，如下。
      \begin{equation*}
         \vec{v}=c_1\cdot\vec{\beta}_1+\dots+c_i\cdot(k\vec{\beta}_i)
                   +\dots+c_n\cdot\vec{\beta}_n
         \;\mapsto\;
         c_1\cdot\vec{\beta}_1+\dots
             +(kc_i)\cdot\vec{\beta}_i+\dots+c_n\cdot\vec{\beta}_n=\vec{v}
      \end{equation*}
行交换矩阵 \( P_{i,j} \) 将相对于基
\( \sequence{\vec{\beta}_1,\dots,\vec{\beta}_j,\dots,\vec{\beta}_i,\dots,\vec{\beta}_n} \)
的表示转换为相对于
\( \sequence{\vec{\beta}_1,\dots,\vec{\beta}_i,\dots,\vec{\beta}_j,\dots,\vec{\beta}_n} \) 的表示。
最后，行组合矩阵 \( C_{i,j}(k) \) 将相对于
\( \sequence{\vec{\beta}_1,\dots,\vec{\beta}_i+k\vec{\beta}_j,\dots,\vec{\beta}_j,\dots,\vec{\beta}_n} \)
的表示转换为相对于 \( B \) 的表示。
      \begin{multline*}
         \vec{v}= c_1\cdot\vec{\beta}_1+\dots
                     +c_i\cdot(\vec{\beta}_i+k\vec{\beta}_j)
                       +\dots+c_j\vec{\beta}_j+\dots+c_n\cdot\vec{\beta}_n
         \\  \mapsto\;
         c_1\cdot\vec{\beta}_1+\dots+c_i\cdot\vec{\beta}_i
               +\dots+(kc_i+c_j)\cdot\vec{\beta}_j
               +\dots+c_n\cdot\vec{\beta}_n=\vec{v}
      \end{multline*}
（与正文中的证明相同，行运算的条件，例如标量 $k$ 非零，保证这些序列都是基。）
    
\end{ans}
\begin{ans}{三.V.1.20}
将 $H$ 看作换基矩阵 $H=\rep{\identity}{B,\stdbasis_n}$，其各列为
       \begin{equation*}
         \colvec{h_{1,i} \\ \vdots \\ h_{n,i}}
         =\rep{\identity(\vec{\beta}_i)}{\stdbasis_n}
         =\rep{\vec{\beta}_i}{\stdbasis_n}
       \end{equation*}
由于相对于标准基的表示是直接的，有
       \begin{equation*}
         \colvec{h_{1,i} \\ \vdots \\ h_{n,i}}
         =\vec{\beta}_i
       \end{equation*}
即所求基由 \( H \) 的各列组成。
    
\end{ans}
\begin{ans}{三.V.1.21}
      \begin{exparts}
\partsitem 可通过一系列行运算将起始向量表示转为终止表示，证明说明了基如何随之改变。
先交换相对于 $B$ 的表示的第一行与第二行，得到相对于新基 $B_1$ 的表示。
          \begin{equation*}
            \rep{1-x+3x^2-x^3}{B_1}=
              \colvec[r]{1 \\ 0 \\ 1 \\ 2}_{B_1}
            \qquad
            B_1=\sequence{1-x,1+x,x^2+x^3,x^2-x^3}
          \end{equation*}
再将表示向量第三行的 \( -2 \) 倍加到第四行。
          \begin{equation*}
            \rep{1-x+3x^2-x^3}{B_3}=
              \colvec[r]{1 \\ 0 \\ 1 \\ 0}_{B_2}
            \qquad
            B_2=\sequence{1-x,1+x,3x^2-x^3,x^2-x^3}
          \end{equation*}
（\( B_2 \) 的第三个元素是 \( B_1 \) 的第三个元素减去第四个元素的 \( -2 \) 倍。）
最后，将第三行加倍即可。
          \begin{equation*}
            \rep{1-x+3x^2-x^3}{D}=
              \colvec[r]{1 \\ 0 \\ 2 \\ 0}_{D}
            \qquad
            D=\sequence{1-x,1+x,(3x^2-x^3)/2,x^2-x^3}
          \end{equation*}
        \partsitem
这个结果有三种表述。第一种：设向量空间 $V$ 有基 $B$，且 $\vec{v}\in V$ 非零，则对任意非零列向量 $\vec{z}$（分量数等于 $V$ 的维数），
都存在换基矩阵 $M$ 使 $M\cdot \rep{\vec{v}}{B}=\vec{z}$。
第二种：对任意 $n$ 维向量空间 $V$、任意非零向量 \( \vec{v}\in V \) 以及非零 \( \vec{z}_1,\vec{z}_2\in\Re^n \)，
存在基 \( B,D\subset V \)，使 \( \rep{\vec{v}}{B}=\vec{z}_1 \)、\( \rep{\vec{v}}{D}=\vec{z}_2 \)。
第三种：对任意 $n$ 维向量空间中的非零向量 $\vec{v}$，以及任意具有 $n$ 个分量的非零列向量，存在一组基使 $\vec{v}$ 相对于它的表示就是该列向量。

前两种容易由第三种推出。第一种中，第三种给出基 $D$ 使 $\rep{\vec{v}}{D}=\vec{z}$，于是 $\rep{\identity}{B,D}$ 就是所需的~$M$。
第二种只需将第三种应用两次。

第三种可仿照本题第一问证明。概要如下：相对于任意基 $B$，将 $\vec{v}$ 表示为列向量 $\vec{z}_1$。
由于 $\vec{v}$ 非零，该列向量必有非零分量。利用这一分量，通过一系列行运算将 $\vec{z}_1$ 转为 $\vec{z}$。
（可对 $V$ 的维数使用归纳法，将此概要补全。）
       \end{exparts}
     
\end{ans}
\begin{ans}{三.V.1.22}
这就是下一小节的主题。
    
\end{ans}
\begin{ans}{三.V.1.23}
换基矩阵非奇异，所以秩等于列数。其各列在 $\Re^n$ 中构成大小为 $n$ 的线性无关子集，因此是一组基。
第二问的答案也为“能”，因为上述推理每一步均可逆，即都可写成“当且仅当”。
    
\end{ans}
\begin{ans}{三.V.1.24}
对于第一部分，这样的矩阵有无穷多个。其中一个相对于 \( \stdbasis_2 \) 表示 \( \Re^2 \) 上具有以下作用的变换。
      \begin{equation*}
        \colvec[r]{1 \\ 0}\mapsto\colvec[r]{4 \\ 0}
        \qquad
        \colvec[r]{0 \\ 1}\mapsto\colvec[r]{0 \\ -1/3}
      \end{equation*}
同时指定两对不同的输入、输出稍复杂一些。矩阵作用的线性性排除了一些可能性。
     \begin{exparts}
\partsitem 存在这样的矩阵。以下条件
         \begin{equation*}
           \begin{mat}
             a  &b  \\
             c  &d
           \end{mat}
           \colvec[r]{1 \\ 3}
           =
           \colvec[r]{1 \\ 1}
           \qquad
           \begin{mat}
             a  &b  \\
             c  &d
           \end{mat}
           \colvec[r]{2 \\ -1}
           =
           \colvec[r]{-1 \\ -1}
         \end{equation*}
可解为
         \begin{equation*}
           \begin{linsys}{4}
             a  &+  &3b  &   &   &   &   &=  &1  \\
                &   &    &   &c  &+  &3d &=  &1  \\
            2a  &-  &b   &   &   &   &   &=  &-1 \\
                &   &    &   &2c &-  &d  &=  &-1
           \end{linsys}
         \end{equation*}
得到这个矩阵。
         \begin{equation*}
           \begin{mat}[r]
             -2/7  &3/7 \\
             -2/7  &3/7
           \end{mat}
         \end{equation*}
\partsitem 不存在，因为
         \begin{equation*}
           2\cdot\colvec[r]{1 \\ 3}=\colvec[r]{2 \\ 6}
           \quad\text{but}\quad
           2\cdot\colvec[r]{1 \\ 1}\neq\colvec[r]{-1 \\ -1}
         \end{equation*}
线性作用不可能产生这种效果。
\partsitem \( \set{\vec{v}_1,\vec{v}_2} \) 线性无关是充分条件，但不是必要条件。
充要条件是：起始向量之间的任意线性相关关系，在终止向量之间也成立。即
         \begin{equation*}
           c_1\vec{v}_1+c_2\vec{v}_2=\zero
           \quad\text{implies}\quad
           c_1\vec{w}_1+c_2\vec{w}_2=\zero.
         \end{equation*}
该条件的证明可按常规方法完成。
     \end{exparts}
   
\end{ans}
\protect \subsection {三.V.2: 改变映射的表示}
\begin{ans}{三.V.2.11}
      \begin{exparts}
\partsitem 是，两者秩均为二。
\partsitem 是，秩相同。
\partsitem 否，秩不同。
      \end{exparts}
    
\end{ans}
\begin{ans}{三.V.2.12}
将矩阵按大小和秩分类，同类中的矩阵大小、秩均相同。
      \begin{exparts}
\partsitem $\nbyn{3}$，秩~$2$
\partsitem $\nbyn{2}$，秩~$1$
\partsitem $\nbym{2}{3}$，秩~$2$
\partsitem $\nbyn{3}$，秩~$2$
\partsitem $\nbyn{3}$，秩~$1$
      \end{exparts}
    
\end{ans}
\begin{ans}{三.V.2.13}
只需计算各自的秩。
      \begin{exparts*}
        \partsitem \( \begin{mat}[r]
                   1  &0  &0  \\
                   0  &0  &0
                 \end{mat} \)
        \partsitem \( \begin{mat}[r]
                   1  &0  &0  &0  \\
                   0  &1  &0  &0  \\
                   0  &0  &1  &0
                 \end{mat} \)
      \end{exparts*}
    
\end{ans}
\begin{ans}{三.V.2.14}
回顾箭头图和公式。
      \begin{equation*}
        \begin{CD}
          \Re^2_{\wrt{B}}                   @>t>T>        \Re^2_{\wrt{D}}       \\
          @V\text{\scriptsize$\identity$} VV   @V\text{\scriptsize$\identity$} VV \\
          \Re^2_{\wrt{\hat{B}}}             @>t>\hat{T}>  \Re^2_{\wrt{\hat{D}}}
        \end{CD}
        \qquad \hat{T}=
         \rep{\identity}{D,\hat{D}}\cdot T\cdot \rep{\identity}{\hat{B},B}
      \end{equation*}
      \begin{exparts}
\partsitem 以下两式
          \begin{equation*}
            \colvec[r]{1 \\ 1}=1\cdot\colvec[r]{-1 \\ 0}
                            +1\cdot\colvec[r]{2 \\ 1}
            \qquad
            \colvec[r]{1 \\ -1}=(-3)\cdot\colvec[r]{-1 \\ 0}
                            +(-1)\cdot\colvec[r]{2 \\ 1}
          \end{equation*}
说明
          \begin{equation*}
            \rep{\identity}{D,\hat{D}}=\begin{mat}[r]
              1  &-3  \\
              1  &-1
            \end{mat}
          \end{equation*}
类似地，以下两式
          \begin{equation*}
            \colvec[r]{0 \\ 1}=0\cdot\colvec[r]{1 \\ 0}
                            +1\cdot\colvec[r]{0 \\ 1}
            \qquad
            \colvec[r]{1 \\  1}=1\cdot\colvec[r]{1 \\ 0}
                            +1\cdot\colvec[r]{0 \\ 1}
          \end{equation*}
给出另一个非奇异矩阵。
          \begin{equation*}
            \rep{\identity}{\hat{B},B}=\begin{mat}[r]
              0  &1  \\
              1  &1
            \end{mat}
          \end{equation*}
因此答案如下。
          \begin{equation*}
            \hat{T}=
            \begin{mat}[r]
              1  &-3  \\
              1  &-1
            \end{mat}
            \begin{mat}[r]
              1  &2  \\
              3  &4
            \end{mat}
            \begin{mat}[r]
              0  &1  \\
              1  &1
            \end{mat}
            =\begin{mat}[r]
              -10  &-18  \\
              -2   &-4
            \end{mat}
          \end{equation*}
虽然并非严格必需，核对一下更令人放心。任取
          \begin{equation*}
            \vec{v}=\colvec[r]{3 \\ 2}
          \end{equation*}
有
          \begin{equation*}
            \rep{\vec{v}}{B}=\colvec[r]{3 \\ 2}_B
            \qquad
            \begin{mat}[r]
              1  &2  \\
              3  &4
            \end{mat}_{B,D}
            \colvec[r]{3 \\ 2}_B
            =\colvec[r]{7 \\ 17}_D
          \end{equation*}
所以 $t(\vec{v})$ 如下。
          \begin{equation*}
            7\cdot\colvec[r]{1 \\ 1}+17\cdot\colvec[r]{1 \\ -1}=\colvec[r]{24 \\ -10}
          \end{equation*}
相对于 $\hat{B},\hat{D}$ 计算，先求
          \begin{equation*}
            \rep{\vec{v}}{\hat{B}}=\colvec[r]{-1 \\ 3}_{\hat{B}}
            \qquad
            \begin{mat}[r]
              -10  &-18  \\
               -2  &-4
            \end{mat}_{\hat{B},\hat{D}}
            \colvec[r]{-1 \\ 3}_{\hat{B}}
            =\colvec[r]{-44 \\ -10}_{\hat{D}}
          \end{equation*}
再验证结果相同。
          \begin{equation*}
            -44\cdot\colvec[r]{-1 \\ 0}-10\cdot\colvec[r]{2 \\ 1}=\colvec[r]{24 \\ -10}
          \end{equation*}
\partsitem 以下两式
      \begin{equation*}
        \colvec[r]{1 \\ 1}=\frac{1}{3}\cdot\colvec[r]{1 \\ 2}
                  +\frac{1}{3}\cdot\colvec[r]{2 \\ 1}
        \qquad
        \colvec[r]{1 \\ -1}=-1\cdot\colvec[r]{1 \\ 2}
                  +1\cdot\colvec[r]{2 \\ 1}
      \end{equation*}
说明
      \begin{equation*}
        \rep{\identity}{D,\hat{D}}=\begin{mat}[r]
          1/3  &-1  \\
          1/3  &1
        \end{mat}
      \end{equation*}
而以下两式
      \begin{equation*}
        \colvec[r]{1 \\ 2}=1\cdot\colvec[r]{1 \\ 0}
                  +2\cdot\colvec[r]{0 \\ 1}
        \qquad
        \colvec[r]{1 \\ 0}=-1\cdot\colvec[r]{1 \\ 0}
                  +0\cdot\colvec[r]{0 \\ 1}
      \end{equation*}
说明
      \begin{equation*}
        \rep{\identity}{\hat{B},B}=\begin{mat}[r]
          1  &1  \\
          2  &0
        \end{mat}
      \end{equation*}
由此，转换如下。
      \begin{equation*}
        \hat{T}=\begin{mat}[r]
          1/3  &-1  \\
          1/3  &1
        \end{mat}
        \begin{mat}[r]
          1  &2  \\
          3  &4
        \end{mat}
        \begin{mat}[r]
          1  &1  \\
          2  &0
        \end{mat}
        =\begin{mat}[r]
          -28/3  &-8/3  \\
          38/3   &10/3
        \end{mat}
      \end{equation*}
与上一问一样，核对可帮助确认计算没有出错。例如，固定向量
      \begin{equation*}
        \vec{v}=\colvec[r]{-1 \\ 2}
      \end{equation*}
（这是任意选取的）。有
      \begin{equation*}
        \rep{\vec{v}}{B}=\colvec[r]{-1 \\ 2}
        \qquad
        \begin{mat}[r]
          1  &2  \\
          3  &4
        \end{mat}_{B,D}
        \colvec[r]{-1  \\ 2}_{B}
        =\colvec[r]{3  \\ 5}_D
      \end{equation*}
所以 $t(\vec{v})$ 为以下向量。
      \begin{equation*}
        3\cdot\colvec[r]{1 \\ 1}+5\cdot\colvec[r]{1 \\ -1}=\colvec[r]{8 \\ -2}
      \end{equation*}
相对于 $\hat{B},\hat{D}$，先计算
      \begin{equation*}
        \rep{\vec{v}}{\hat{B}}=\colvec[r]{1 \\ -2}
        \qquad
        \begin{mat}[r]
          -28/3  &-8/3  \\
          38/3   &10/3
        \end{mat}_{\hat{B},\hat{D}}
        \colvec[r]{1 \\ -2}_{\hat{B}}
        =\colvec[r]{-4 \\ 6}_{\hat{D}}
      \end{equation*}
果然，得到相同的 $t(\vec{v})$。
      \begin{equation*}
        -4\cdot\colvec[r]{1 \\ 2}+6\cdot\colvec[r]{2 \\ 1}=\colvec[r]{8 \\ -2}
      \end{equation*}
      \end{exparts}
    
\end{ans}
\begin{ans}{三.V.2.15}
若 \( H \)、\( \hat{H} \) 为 \( \nbym{m}{n} \)，则 \( P \) 为 \( \nbyn{m} \)，\( Q \) 为 \( \nbyn{n} \)。
    
\end{ans}
\begin{ans}{三.V.2.16}
      \begin{exparts}
        \partsitem
箭头图如下。
          \begin{equation*}
            \begin{CD}
              V_{\wrt{B}}                   @>h>H>        W_{\wrt{D}}       \\
              @V{\text{\scriptsize$\identity$}} VV      @V{\text{\scriptsize$\identity$}} VV \\
              V_{\wrt{\hat{B}}}             @>h>\hat{H}>   W_{\wrt{\hat{D}}}
            \end{CD}
          \end{equation*}
由图得 $\hat{H}=PHQ$，其中 $Q=\rep{\identity}{\hat{B},B}$、$P=\rep{\identity}{D,\hat{D}}$
（记住先执行的运算写在右边）。
        \partsitem
要求 $Q=\rep{\identity}{\hat{B},B}$ 和 $P=\rep{\identity}{D,\hat{D}}$。
对 $Q$，作以下计算（这里直接观察）。
          \begin{equation*}
            \rep{\identity(1)}{B}=\colvec{1 \\ 0 \\ 0}
            \quad
            \rep{\identity(x)}{B}=\colvec{-1 \\ 1 \\ 0}
            \quad
            \rep{\identity(x^2)}{B}=\colvec{-1 \\ 0 \\ 1}
          \end{equation*}
以下计算给出 $P$。
          \begin{multline*}
            \rep{\identity(
              \begin{mat}
                0 &0 \\
                0 &1
              \end{mat})}{\hat{D}}=\colvec{0 \\ 0 \\ 0 \\ 1}
            \quad
            \rep{\identity(
              \begin{mat}
                0 &0 \\
                1 &1
              \end{mat})}{\hat{D}}=\colvec{0 \\ 0 \\ 1 \\ 1}
            \\
            \rep{\identity(
              \begin{mat}
                0 &1 \\
                1 &1
              \end{mat})}{\hat{D}}=\colvec{0 \\ -1 \\ 1 \\ 1}
            \rep{\identity(
              \begin{mat}
                1 &1 \\
                1 &1
              \end{mat})}{\hat{D}}=\colvec{-1 \\ -1 \\ 1 \\ 1}
          \end{multline*}
答案如下。
          \begin{equation*}
            P=
            \begin{mat}
              0 &0 &0  &-1 \\
              0 &0 &-1 &-1  \\
              0 &1 &1  &1  \\
              1 &1 &1  &1
            \end{mat}
            \qquad
            Q=\begin{mat}
              1 &-1 &-1 \\
              0 &1  &0  \\
              0 &0  &1
            \end{mat}
          \end{equation*}
      \end{exparts}
    
\end{ans}
\begin{ans}{三.V.2.17}
箭头图如下。
          \begin{equation*}
            \begin{CD}
              V_{\wrt{B}}                   @>t>T>        W_{\wrt{D}}       \\
              @V{\text{\scriptsize$\identity$}} VV      @V{\text{\scriptsize$\identity$}} VV \\
              V_{\wrt{\hat{B}}}             @>t>\hat{T}>   W_{\wrt{\hat{D}}}
            \end{CD}
          \end{equation*}
取 $Q=\rep{\identity}{\hat{B},B}$、$P=\rep{\identity}{D,\hat{D}}$，则 $\hat{T}=PTQ$。

以下计算给出 $Q=\rep{\identity}{\hat{B},B}$（直接观察）。
          \begin{equation*}
            \rep{\identity(\colvec{1 \\ 0})}{B}=\colvec{1 \\ 0}
            \quad
            \rep{\identity(\colvec{0 \\ 1})}{B}=\colvec{-1 \\ 1}
          \end{equation*}
所以得到
          \begin{equation*}
            Q=\begin{mat}
              1 &-1 \\
              0 &1
             \end{mat}
          \end{equation*}
以下计算给出 $P=\rep{\identity}{D,\hat{D}}$。
          \begin{equation*}
            \rep{\identity(\colvec{0 \\ 1})}{\hat{D}}=\colvec{0 \\ 1}
            \quad
            \rep{\identity(\colvec{-1 \\ 0})}{\hat{D}}=\colvec{-1 \\ -1}
          \end{equation*}
将它们并排拼接成另一个矩阵。
          \begin{equation*}
            P=\begin{mat}
              0 &-1 \\
              1 &-1
             \end{mat}
          \end{equation*}
    
\end{ans}
\begin{ans}{三.V.2.18}
高斯消元得到
      \begin{equation*}
          \begin{mat}
            2 &1  &1  \\
            3 &-1 &0  \\
            1 &3  &2
          \end{mat}
          \grstep[-(1/2)\rho_1+\rho_3]{-(3/2)\rho_1+\rho_2}
          \grstep{\rho_2+\rho_3}
          \grstep[-(2/5)\rho_2]{(1/2)\rho_1}
          \begin{mat}
            1 &1/2  &1/2  \\
            0 &1    &3/5  \\
            0 &0    &0
          \end{mat}
      \end{equation*}
再用列运算化为矩阵等价的标准形。
      \begin{equation*}
        \grstep{-(3/5)\text{col}_2+\text{col}_3}
        \grstep[-(1/5)\text{col}_1+\text{col}_3]{-(1/2)\text{col}_1+\text{col}_2}
          \begin{mat}
            1 &0    &0  \\
            0 &1 &0  \\
            0 &0  &0
          \end{mat}
      \end{equation*}
于是两个矩阵如下。
      \begin{align*}
        P&=
        \begin{mat}
          1 &0 &0 \\
          0 &-2/5 &0 \\
          0 &0 &1
         \end{mat}
        \begin{mat}
          1/2 &0 &0 \\
          0 &1 &0 \\
          0 &0 &1
         \end{mat}
        \begin{mat}
          1 &0 &0 \\
          0 &1 &0 \\
          0 &1 &1
         \end{mat}
        \begin{mat}
          1    &0 &0 \\
          0    &1 &0 \\
          -1/2 &0 &1
         \end{mat}
        \begin{mat}
          1    &0 &0 \\
          -3/2 &1 &0 \\
          0    &0 &1
         \end{mat}                                 \\
         &=
         \begin{mat}
           1/2   &0     &0 \\
           3/5   &-2/5  &0 \\
           -2    &1     &1
         \end{mat}                              \\
        Q&=
        \begin{mat}
          1 &0 &0 \\
          0 &1 &-3/5 \\
          0 &0 &1
        \end{mat}
        \begin{mat}
          1 &-1/2 &0 \\
          0 &1    &0 \\
          0 &0    &1
        \end{mat}
        \begin{mat}
          1 &0 &-1/5 \\
          0 &1 &0 \\
          0 &0 &1
        \end{mat}
        =
        \begin{mat}
          1 &-1/2 &-1/5 \\
          0 &1    &-3/5 \\
          0 &0    &1
        \end{mat}
      \end{align*}
    
\end{ans}
\begin{ans}{三.V.2.19}
任意 \( \nbyn{n} \) 矩阵非奇异，当且仅当秩为 \( n \)；由 \nearbytheorem{th:CanonFormForMatEquiv}，这等价于它矩阵等价于对角线全为一的 $\nbyn{n}$ 矩阵。
    
\end{ans}
\begin{ans}{三.V.2.20}
若 \( PAQ=I \)，则 \( QPAQ=Q \)，从而 \( QPA=I \)，所以 \( QP=A^{-1} \)。
    
\end{ans}
\begin{ans}{三.V.2.21}
由 \nearbyexample{ex:DiagizedMat} 后的定义，矩阵 $M$ 可对角化，是指它表示变换 $M=\rep{t}{B,D}$，且存在基 $\hat{B}$ 使
$\rep{t}{\hat{B},\hat{B}}$ 为对角矩阵\Dash 起始基和终止基必须相同。
但 \nearbytheorem{th:CanonFormForMatEquiv} 只说明存在 $\hat{B}$、$\hat{D}$，使转换后的表示 $\rep{t}{\hat{B},\hat{D}}$ 为对角矩阵。
没有理由认为能选取相同的 $\hat{B}$ 和 $\hat{D}$。
    
\end{ans}
\begin{ans}{三.V.2.22}
是。行秩等于列秩，所以转置的秩等于原矩阵的秩。相同大小、相同秩的矩阵矩阵等价。
    
\end{ans}
\begin{ans}{三.V.2.23}
只有零矩阵的秩为零。
    
\end{ans}
\begin{ans}{三.V.2.24}
自反性：取 \( P \)、\( Q \) 为单位矩阵，则任意矩阵都与自身矩阵等价。
对称性：若 \( H_1=PH_2Q \)，则 \( H_2=P^{-1}H_1Q^{-1} \)（$P$、$Q$ 非奇异，故逆存在）。
传递性：设 \( H_1=P_2H_2Q_2 \)、\( H_2=P_3H_3Q_3 \)，代入得
\( H_1=P_2(P_3H_3Q_3)Q_2=(P_2P_3)H_3(Q_3Q_2) \)。
非奇异矩阵的乘积非奇异（已证明可逆矩阵的乘积可逆，并给出求逆方法），所以 \( H_1 \) 与 \( H_3 \) 矩阵等价。
     
\end{ans}
\begin{ans}{三.V.2.25}
由 \nearbytheorem{th:CanonFormForMatEquiv}，零矩阵是其类中唯一的元素，因为它是唯一秩为零的 \( \nbym{m}{n} \) 矩阵。
其他矩阵都不是类中唯一的元素，因为非零标量倍数与原矩阵的秩相同。
    
\end{ans}
\begin{ans}{三.V.2.26}
\( \nbyn{1} \) 矩阵有两个矩阵等价类\Dash 秩为零和秩为一。
\( \nbyn{3} \) 矩阵分为四个矩阵等价类。
    
\end{ans}
\begin{ans}{三.V.2.27}
对 \( \nbym{m}{n} \) 矩阵，每个可能的秩对应一类。设 \( k \) 为 \( m \)、\( n \) 的最小值，则秩可为 \( 0 \)、\( 1 \)、\ldots、\( k \)，共 \( k+1 \) 类。
（当然，将所有大小的矩阵合计，就有无穷多个类。）
    
\end{ans}
\begin{ans}{三.V.2.28}
它们对非零数乘封闭，因为非零标量倍数的秩不变。
对加法不封闭，例如 \( H+(-H) \) 的秩为零。
    
\end{ans}
\begin{ans}{三.V.2.29}
图示如下。
        \begin{equation*}
          \begin{CD}
            \Re^2_{\wrt{\stdbasis_2}}         @>t>T>     \Re^2_{\wrt{\stdbasis_2}}    \\
            @V\text{\scriptsize$\identity$} VV   @V\text{\scriptsize$\identity$} VV \\
           \Re^2_{\wrt{B}}        @>t>\hat{T}>  \Re^2_{\wrt{B}}
          \end{CD}
        \end{equation*}
从左下到右下有两条路径。第一条直接使用 $\hat{T}=\rep{t}{B,B}$；第二条先向上、再向右、再向下，使用
$\rep{\identity}{\stdbasis_2,B}\cdot T\cdot \rep{\identity}{B,\stdbasis_2}$。
记住从右向左书写，因此“向上”的矩阵在右边，这样作用于 $\rep{\vec{v}}{B}$ 时，先作用的是 $\rep{\identity}{B,\stdbasis_2}$。
将 $\rep{\identity}{B,\stdbasis_2}$ 记为 $S$；两个换基矩阵中选择这一个，是因为它更容易计算。
于是 $\hat{T}=S^{-1}TS$。
      \begin{exparts}
\partsitem 有
          \begin{equation*}
            \rep{\identity}{B,\stdbasis_2}=S=
            \begin{mat}[r]
              1  &-1  \\
              2  &-1
            \end{mat}
          \end{equation*}
以及
          \begin{equation*}
            \rep{\identity}{\stdbasis_2,B}
            =\rep{\identity}{B,\stdbasis_2}^{-1}=S^{-1}=
            \begin{mat}[r]
              1  &-1  \\
              2  &-1
            \end{mat}^{-1}
            =\begin{mat}[r]
              -1  &1  \\
              -2  &1
            \end{mat}
          \end{equation*}
因此答案如下。
          \begin{equation*}
            \rep{t}{B,B}
            =
            \begin{mat}[r]
              -1  &1  \\
              -2  &1
            \end{mat}
            \begin{mat}[r]
              1  &1  \\
              3  &-1
            \end{mat}
            \begin{mat}[r]
              1  &-1  \\
              2  &-1
            \end{mat}
            =\begin{mat}[r]
              -2  &0   \\
              -5  &2
            \end{mat}
          \end{equation*}
为快速核对，可任取向量
          \begin{equation*}
            \vec{v}=\colvec[r]{4  \\  5}
          \end{equation*}
得到
          \begin{equation*}
            \rep{\vec{v}}{\stdbasis_2}=\colvec[r]{4 \\ 5}
            \qquad
            \begin{mat}[r]
              1  &1  \\
              3  &-1
            \end{mat}
            \colvec[r]{4 \\ 5}
            =\colvec[r]{9 \\ 7}=t(\vec{v})
          \end{equation*}
而相对于 $B,B$ 计算，
          \begin{equation*}
            \rep{\vec{v}}{B}=\colvec[r]{1 \\ -3}
            \qquad
            \begin{mat}[r]
              -2  &0   \\
              -5  &2
            \end{mat}_{B,B}
            \colvec[r]{1 \\ -3}_B
            =\colvec[r]{-2 \\ -11}_B
          \end{equation*}
得到相同结果。
          \begin{equation*}
            -2\cdot\colvec[r]{1 \\ 2}-11\cdot\colvec[r]{-1 \\ -1}
               =\colvec[r]{9 \\ 7}
          \end{equation*}
\partsitem 与本题第一问一样，
          \begin{equation*}
            S=\rep{\identity}{B,\stdbasis_2}
            =\begin{pmat}{c|c|c}
              \vec{\beta}_1 &\;\cdots\; &\vec{\beta}_n
            \end{pmat}
            \qquad
            \rep{\identity}{\stdbasis_2,B}=S^{-1}=
             \rep{\identity}{B,\stdbasis_2}^{-1}
          \end{equation*}
因此，将各列为基向量的矩阵记为 $S$，有 $\rep{t}{B,B}=\hat{T}=S^{-1}TS$。
      \end{exparts}
    
\end{ans}
\begin{ans}{三.V.2.30}
       \begin{exparts}
\partsitem 相应的箭头图如下。
           \begin{equation*}
             \begin{CD}
               V_{\wrt{B_1}}                   @>h>H>        W_{\wrt{D}}       \\
               @V\text{\scriptsize$\identity$} VQV      @V\text{\scriptsize$\identity$} VPV \\
               V_{\wrt{B_2}}             @>h>\hat{H}>  W_{\wrt{D}}
             \end{CD}
           \end{equation*}
由于 \( W \) 中无须换基（即换基矩阵 $P$ 为单位矩阵），有
\( \rep{h}{B_2,D}=\rep{h}{B_1,D}\cdot Q \)，其中 \( Q=\rep{\identity}{B_2,B_1} \)。
\partsitem 此时箭头图如下。
           \begin{equation*}
             \begin{CD}
               V_{\wrt{B}}                   @>h>H>   W_{\wrt{D_1}}       \\
               @V\text{\scriptsize$\identity$} VQV  @V\text{\scriptsize$\identity$} VPV \\
               V_{\wrt{B}}             @>h>\hat{H}>  W_{\wrt{D_2}}
             \end{CD}
           \end{equation*}
有 \( \rep{h}{B,D_2}=P\cdot \rep{h}{B,D_1} \)，其中 \( P=\rep{\identity}{D_1,D_2} \)。
       \end{exparts}
      
\end{ans}
\begin{ans}{三.V.2.31}
      \begin{exparts}
\partsitem 以下是箭头图及对应的逆函数版本。
          \begin{equation*}
           \begin{CD}
             V_{\wrt{B}}                   @>h>H>      W_{\wrt{D}}       \\
             @V\text{\scriptsize$\identity$} VQV  @V\text{\scriptsize$\identity$} VPV \\
             V_{\wrt{\hat{B}}}             @>h>\hat{H}> W_{\wrt{\hat{D}}}
            \end{CD}
            \hspace*{4em}\qquad
           \begin{CD}
             V_{\wrt{B}}              @<h^{-1}<H^{-1}<    W_{\wrt{D}}       \\
             @V\text{\scriptsize$\identity$} VQV  @V\text{\scriptsize$\identity$} VPV \\
             V_{\wrt{\hat{B}}}        @<h^{-1}<\hat{H}^{-1}< W_{\wrt{\hat{D}}}
            \end{CD}
           \end{equation*}
是。逆矩阵表示逆映射。从右下到左下，可以先向上、再向左、再向下。
即若 \( \hat{H}=PHQ \)，其中 \( P,Q \) 可逆，且 \( H,\hat{H} \) 可逆，则
\( \hat{H}^{-1}=Q^{-1}H^{-1}P^{-1} \)。
\partsitem 是，这是上一问的另一种表述。
\partsitem 不一定。需要额外假设：若 \( H \) 相对于相同的起始、终止基 \( B,B \) 表示 \( h \)，则 \( H^2 \) 表示 \( \composed{h}{h} \)。
具体反例是以下两个矩阵，它们的秩均为一，所以矩阵等价，
          \begin{equation*}
             \begin{mat}[r]
               1  &0  \\
               0  &0
             \end{mat}
             \qquad
             \begin{mat}[r]
               0  &0  \\
               1  &0
             \end{mat}
          \end{equation*}
但平方不矩阵等价\Dash 第一个平方的秩为一，第二个平方的秩为零。
\partsitem 不一定。以下两个矩阵不矩阵等价，但其平方矩阵等价。
          \begin{equation*}
             \begin{mat}[r]
               0  &0  \\
               0  &0
             \end{mat}
             \qquad
             \begin{mat}[r]
               0  &0  \\
               1  &0
             \end{mat}
          \end{equation*}
      \end{exparts}
    
\end{ans}
\begin{ans}{三.V.2.32}
      \begin{exparts}
\partsitem 箭头图提示以下定义。
          \begin{equation*}
            \begin{CD}
              V_{\wrt{B_1}}                   @>t>T>        V_{\wrt{B_1}}       \\
              @V\text{\scriptsize$\identity$} VV  @V\text{\scriptsize$\identity$} VV \\
              V_{\wrt{B_2}}             @>t>\hat{T}>  V_{\wrt{B_2}}
            \end{CD}
          \end{equation*}
若存在非奇异矩阵 \( P \) 使 \( \hat{T}=P^{-1}TP \)，则称矩阵 \( T,\hat{T} \) \definend{相似}。
\item 在矩阵等价的定义中，用 \( P^{-1} \) 作 \( P \)，用 \( P \) 作 \( Q \)。
\item \textit{与 \nearbyexercise{exer:MatEqIsEqRel} 类似。}
自反性显然：\( T=I^{-1}TI \)。对称性也容易：\( \hat{T}=P^{-1}TP \) 蕴含 \( T=P\hat{T}P^{-1} \)
（第一式右乘 $P^{-1}$、左乘 $P$）。
对于传递性，设 \( T_1={P_2}^{-1}T_2P_2 \)、\( T_2={P_3}^{-1}T_3P_3 \)，则
\( T_1={P_2}^{-1}({P_3}^{-1}T_3P_3)P_2=({P_2}^{-1}{P_3}^{-1})T_3(P_3P_2) \)。
注意 \( P_3P_2 \) 可逆，逆为 \( {P_2}^{-1}{P_3}^{-1} \)，证明完成。
\partsitem 设 \( \hat{T}=P^{-1}TP \)。对于平方，
\( \hat{T}^2=(P^{-1}TP)(P^{-1}TP)=P^{-1}T(PP^{-1})TP=P^{-1}T^2P \)。更高次幂用归纳法得到。
\partsitem 以下两个矩阵矩阵等价，但其平方不矩阵等价。
          \begin{equation*}
             \begin{mat}[r]
               1  &0  \\
               0  &0
             \end{mat}
             \qquad
             \begin{mat}[r]
               0  &0  \\
               1  &0
             \end{mat}
          \end{equation*}
由上一问可知，矩阵相似与矩阵等价不同。
      \end{exparts}
   
\end{ans}
\protect \section {第 VI 节：投影}
\protect \subsection {三.VI.1: 正交投影到直线}
\begin{ans}{三.VI.1.6}
       % Each is a straightforward application of the formula from
       % \nearbydefinition{def:ProjIntoLine}.
       \begin{exparts}
        \partsitem
          $\displaystyle \frac{\colvec[r]{2 \\ 1}\dotprod\colvec[r]{3 \\ -2}}{%
                               \colvec[r]{3 \\ -2}\dotprod\colvec[r]{3 \\ -2}}
             \cdot\colvec[r]{3 \\ -2}
           =\frac{4}{13}
             \cdot\colvec[r]{3 \\ -2}
           =\colvec[r]{12/13 \\ -8/13}$
        \partsitem
          $\displaystyle \frac{\colvec[r]{2 \\ 1}\dotprod\colvec[r]{3 \\ 0}}{%
                               \colvec[r]{3 \\ 0}\dotprod\colvec[r]{3 \\ 0}}
             \cdot\colvec[r]{3 \\ 0}
           =\frac{2}{3}
             \cdot\colvec[r]{3 \\ 0}
           =\colvec[r]{2 \\ 0}$
        \partsitem
          $\displaystyle
             \frac{\colvec[r]{1 \\ 1 \\ 4}\dotprod\colvec[r]{1 \\ 2 \\ -1}}{%
                   \colvec[r]{1 \\ 2 \\ -1}\dotprod\colvec[r]{1 \\ 2 \\ -1}}
             \cdot\colvec[r]{1 \\ 2 \\ -1}
           =\frac{-1}{6}
             \cdot\colvec[r]{1 \\ 2 \\ -1}
           =\colvec[r]{-1/6 \\ -1/3 \\ 1/6}$
        \partsitem
          $\displaystyle
             \frac{\colvec[r]{1 \\ 1 \\ 4}\dotprod\colvec[r]{3 \\ 3 \\ 12}}{%
                   \colvec[r]{3 \\ 3 \\ 12}\dotprod\colvec[r]{3 \\ 3 \\ 12}}
             \cdot\colvec[r]{3 \\ 3 \\ 12}
           =\frac{1}{3}
             \cdot\colvec[r]{3 \\ 3 \\ 12}
           =\colvec[r]{1 \\ 1 \\ 4}$
      \end{exparts}
    
\end{ans}
\begin{ans}{三.VI.1.7}
      \begin{exparts}
        \partsitem
          $\displaystyle
             \frac{\colvec[r]{2 \\ -1 \\ 4}\dotprod\colvec[r]{-3 \\ 1 \\ -3}}{%
                   \colvec[r]{-3 \\ 1 \\ -3}\dotprod\colvec[r]{-3 \\ 1 \\ -3}}
             \cdot\colvec[r]{-3 \\ 1 \\ -3}
           =\frac{-19}{19}\cdot\colvec[r]{-3 \\ 1 \\ -3}
           =\colvec[r]{3 \\ -1 \\ 3}$
         \partsitem Writing the line as
           \begin{equation*}
             \set{c\cdot \colvec[r]{1 \\ 3}\suchthat c\in\Re}
           \end{equation*}
将直线写成以下形式即可得到这个投影。
           \begin{equation*}
             \frac{\colvec[r]{-1 \\ -1}\dotprod\colvec[r]{1 \\ 3}}{%
                   \colvec[r]{1 \\ 3}\dotprod\colvec[r]{1 \\ 3}}
              \cdot\colvec[r]{1 \\ 3}
             =\frac{-4}{10}\cdot\colvec[r]{1 \\ 3}
             =\colvec[r]{-2/5 \\ -6/5}
           \end{equation*}
      \end{exparts}
    
\end{ans}
\begin{ans}{三.VI.1.8}
      $\displaystyle
         \frac{\colvec[r]{1 \\ 2 \\ 1 \\ 3}\dotprod\colvec[r]{-1 \\ 1 \\ -1 \\ 1}}{%
               \colvec[r]{-1 \\ 1 \\ -1 \\ 1}\dotprod\colvec[r]{-1 \\ 1 \\ -1 \\ 1}}
         \cdot\colvec[r]{-1 \\ 1 \\ -1 \\ 1}
       =\frac{3}{4}
         \cdot\colvec[r]{-1 \\ 1 \\ -1 \\ 1}
       =\colvec[r]{-3/4 \\ 3/4 \\ -3/4 \\ 3/4}$
    
\end{ans}
\begin{ans}{三.VI.1.9}
      \begin{exparts}
        \partsitem $\displaystyle
          \frac{\colvec[r]{1 \\ 2}\dotprod\colvec[r]{3 \\ 1}}{%
                 \colvec[r]{3 \\ 1}\dotprod\colvec[r]{3 \\ 1}}
          \cdot \colvec[r]{3 \\ 1}
          =\frac{1}{2}\cdot\colvec[r]{3 \\ 1}
          =\colvec[r]{3/2 \\ 1/2}$
        \partsitem $\displaystyle
          \frac{\colvec[r]{0 \\ 4}\dotprod\colvec[r]{3 \\ 1}}{%
                 \colvec[r]{3 \\ 1}\dotprod\colvec[r]{3 \\ 1}}
          \cdot \colvec[r]{3 \\ 1}
          =\frac{2}{5}\cdot\colvec[r]{3 \\ 1}
          =\colvec[r]{6/5 \\ 2/5}$
      \end{exparts}
一般地，投影为
      \begin{equation*}
          \frac{\colvec{x_1 \\ x_2}\dotprod\colvec[r]{3 \\ 1}}{%
                 \colvec[r]{3 \\ 1}\dotprod\colvec[r]{3 \\ 1}}
          \cdot \colvec[r]{3 \\ 1}
          =\frac{3x_1+x_2}{10}\cdot\colvec[r]{3 \\ 1}
          =\colvec{(9x_1+3x_2)/10 \\ (3x_1+x_2)/10}
      \end{equation*}
相应矩阵为
      \begin{equation*}
        \begin{mat}[r]
          9/10  &3/10  \\
          3/10  &1/10
        \end{mat}
      \end{equation*}
     
\end{ans}
\begin{ans}{三.VI.1.10}
设两个向量非零且正交，考虑 $c_1\vec{v}_1+c_2\vec{v}_2=\zero$。两边与 $\vec{v}_1$ 作点积。
      \begin{multline*}
        \vec{v}_1\dotprod(c_1\vec{v}_1+c_2\vec{v}_2)
        =c_1\cdot(\vec{v}_1\dotprod\vec{v}_1)
           +c_2\cdot(\vec{v}_1\dotprod\vec{v}_2)            \\
        =c_1\cdot (\vec{v}_1\dotprod\vec{v}_1)+c_2\cdot 0
        =c_1\cdot (\vec{v}_1\dotprod\vec{v}_1)
      \end{multline*}
左边等于右边的 $\vec{v}_1\dotprod\zero=\zero$，由 $\vec{v}_1\ne\zero$ 得 $c_1=0$；同理 $c_2=0$。
    
\end{ans}
\begin{ans}{三.VI.1.11}
      \begin{exparts}
\partsitem 若 $\vec{v}$ 在线上，正交投影就是 \(\vec{v}\)。因为 $\vec{v}=c_{\vec{v}}\vec{s}$。
         \begin{equation*}
           \frac{\vec{v}\dotprod\vec{s}}{\vec{s}\dotprod\vec{s}}\cdot\vec{s}
           =\frac{c_{\vec{v}}\cdot \vec{s}\dotprod\vec{s}}{%
                         \vec{s}\dotprod\vec{s}}\cdot\vec{s}
           =c_{\vec{v}}\cdot\frac{\vec{s}\dotprod\vec{s}}{%
                                  \vec{s}\dotprod\vec{s}}\cdot\vec{s}
           =c_{\vec{v}}\cdot 1 \cdot\vec{s}
           =\vec{v}
         \end{equation*}
\partsitem 设投影为 $c_{\vec{p}}\vec{s}$。当 $\vec{v}$ 不在线上时，两向量均非零；若 $c_{\vec{p}}=0$，集合只有非零向量 $\vec{v}$，显然线性无关。因此只需考虑 $c_{\vec{p}}\ne0$。

设线性关系
         \begin{equation*}
           a_1(\vec{v})+a_2(\vec{v}-c_{\vec{p}}\vec{s})=\zero
         \end{equation*}
得到 $(a_1+a_2)\vec{v}=a_2c_{\vec{p}}\vec{s}$。由于 $\vec{v}$ 不在线上，两边系数均为零，故 $a_2=0$ 且 $a_1=0$，集合线性无关。
     \end{exparts}
    
\end{ans}
\begin{ans}{三.VI.1.12}
若 $\vec{s}$ 是零向量，以下表达式
          \begin{equation*}
            \proj{\vec{v}}{\spanof{\vec{s}\,}}=
            \frac{ \vec{v}\dotprod\vec{s} }{%
                    \vec{s}\dotprod\vec{s} }\cdot\vec{s}
          \end{equation*}
含有除零，未定义。投影必须属于零向量的张成空间，所以合理定义为 \(\zero\)。
    
\end{ans}
\begin{ans}{三.VI.1.13}
只要 $n>1$，$\Re^n$ 中任意向量都是某个其他向量到某条直线的投影。

设 $\vec{v}\in\Re^n$ 且 $n>1$。若 $\vec{v}\ne\zero$，取其张成的直线；若为零向量，取任意非退化直线。若某个分量 $v_i=0$，取 $\vec{w}=\vec{e}_i$；若所有分量非零，取分量为 $(v_2,-v_1,0,\dots,0)$ 的 $\vec{w}$。两种情况下 $\vec{w}\perp\vec{v}$，且 $\vec{v}$ 是 $\vec{v}+\vec{w}$ 在该直线上的投影。
      \begin{equation*}
        \frac{ (\vec{v}+\vec{w})\dotprod\vec{v} }{
               \vec{v}\dotprod\vec{v} }\cdot\vec{v}
        =\bigl(\frac{ \vec{v}\dotprod\vec{v} }{
                     \vec{v}\dotprod\vec{v} }+
               \frac{ \vec{w}\dotprod\vec{v} }{ \
                      \vec{v}\dotprod\vec{v} }\bigr)\cdot\vec{v}
        =\frac{ \vec{v}\dotprod\vec{v} }{ \vec{v}\dotprod\vec{v} }\cdot\vec{v}
        =\vec{v}
      \end{equation*}

$n=0$ 时空间只有一个向量，不能是不同向量的投影；$n=1$ 时只有一条非退化直线，每个向量只投影到自身。
    
\end{ans}
\begin{ans}{三.VI.1.14}
直接计算即可。回顾 $|\vec{u}|^2=\vec{u}\dotprod\vec{u}$。
      \begin{equation*}
        \absval{\frac{\vec{v}\dotprod\vec{s}}{\vec{s}\dotprod\vec{s}}
              \cdot\vec{s}\,}^2
        =
        \bigl(\frac{\vec{v}\dotprod\vec{s}}{\vec{s}\dotprod\vec{s}}
           \cdot \vec{s}\bigr)
        \dotprod
        \bigl(\frac{\vec{v}\dotprod\vec{s}}{\vec{s}\dotprod\vec{s}}
           \cdot \vec{s}\bigr)
        =
        \bigl(\frac{\vec{v}\dotprod\vec{s}}{\vec{s}\dotprod\vec{s}}\bigr)^2
        \cdot(\vec{s}\dotprod\vec{s})
        =
        \frac{(\vec{v}\dotprod\vec{s})^2}{(\vec{s}\dotprod\vec{s})^2}
        \cdot(\vec{s}\dotprod\vec{s})
        =
        \frac{(\vec{v}\dotprod\vec{s})^2}{\absval{\vec{s}}^2}
      \end{equation*}
    
\end{ans}
\begin{ans}{三.VI.1.15}
由于投影为
      \begin{equation*}
        \frac{\vec{v}\dotprod\vec{s}}{\vec{s}\dotprod\vec{s}}\cdot\vec{s}
      \end{equation*}
点到直线的距离平方为
      \begin{align*}
        \absval{\vec{v}-
                \frac{\vec{v}\dotprod\vec{s}}{\vec{s}\dotprod\vec{s}}
                \cdot\vec{s}\, }^2
        &=\big( \vec{v}-
                \frac{\vec{v}\dotprod\vec{s}}{\vec{s}\dotprod\vec{s}}
                \cdot\vec{s} \big)
          \dotprod
          \big( \vec{v}-
                \frac{\vec{v}\dotprod\vec{s}}{\vec{s}\dotprod\vec{s}}
                \cdot\vec{s} \big)                                  \\
        &=\vec{v}\dotprod\vec{v}
        -\vec{v}\dotprod(
              \frac{\vec{v}\dotprod\vec{s}}{\vec{s}\dotprod\vec{s}}\cdot\vec{s}
           )
        -(
              \frac{\vec{v}\dotprod\vec{s}}{\vec{s}\dotprod\vec{s}}
               \cdot\vec{s}\,
         )\dotprod\vec{v}
        +\big(\frac{\vec{v}\dotprod\vec{s}}{
                \vec{s}\dotprod\vec{s}}\cdot\vec{s}\,\big)
          \dotprod
          \big(\frac{\vec{v}\dotprod\vec{s}}{
                \vec{s}\dotprod\vec{s}}\cdot\vec{s}\,\big)     \\
        &=\vec{v}\dotprod\vec{v}
        -2\cdot\big(\frac{\vec{v}\dotprod\vec{s}}{\vec{s}\dotprod\vec{s}}\big)
           \cdot\vec{v}\dotprod\vec{s}
        +\big(\frac{\vec{v}\dotprod\vec{s}}{
                \vec{s}\dotprod\vec{s}}\big)^2\cdot(\vec{s}\dotprod\vec{s})  \\
        &=\frac{(\vec{v}\dotprod\vec{v}\,)\cdot(\vec{s}\dotprod\vec{s}\,)
                 -2\cdot(\vec{v}\dotprod\vec{s}\,)^2
                 +(\vec{v}\dotprod\vec{s}\,)^2}{%
               \vec{s}\dotprod\vec{s}}                               \\
        &=\frac{(\vec{v}\dotprod\vec{v}\,)(\vec{s}\dotprod\vec{s}\,)
                -(\vec{v}\dotprod\vec{s}\,)^2}{\vec{s}\dotprod\vec{s}}
      \end{align*}
    
\end{ans}
\begin{ans}{三.VI.1.16}
平方根严格递增，所以最小化 $d(c)=(cs_1-v_1)^2+(cs_2-v_2)^2$ 即可。求导并令导数为零，得到唯一临界点。
      \begin{equation*}
        c=\frac{v_1s_1+v_2s_2}{{s_1}^2+{s_2}^2}
         =\frac{\vec{v}\dotprod\vec{s}}{\vec{s}\dotprod\vec{s}}
      \end{equation*}
现在对 $c$ 求二阶导数
      \begin{equation*}
        \frac{d^2\,d}{dc^2}=2{s_1}^2+2{s_2}^2
      \end{equation*}
只要 $s_1,s_2$ 不同时为零，它严格为正，因此临界点是最小值。

推广到 $\Re^n$ 只需对 $d_n(c)=(cs_1-v_1)^2+\dots+(cs_n-v_n)^2$ 求导。
    
\end{ans}
\begin{ans}{三.VI.1.17}
任取 $\vec{w}\in\ell$。由于正交投影，$\vec{v}-\vec{p}$ 与 $\vec{p}-\vec{w}$ 垂直。三角不等式说明斜边 $\vec{v}-\vec{w}$ 至少与 $\vec{v}-\vec{p}$ 一样长。
    
\end{ans}
\begin{ans}{三.VI.1.18}
任意向量满足 $|\vec{u}|^2=\vec{u}\dotprod\vec{u}$，所以投影长度平方为
      \begin{equation*}
        \absval{\frac{\vec{v}\dotprod\vec{s}}{\vec{s}\dotprod\vec{s}}
              \cdot\vec{s}\,}^2
        =
        \bigl(\frac{\vec{v}\dotprod\vec{s}}{\vec{s}\dotprod\vec{s}}
           \cdot \vec{s}\,\bigr)
        \dotprod
        \bigl(\frac{\vec{v}\dotprod\vec{s}}{\vec{s}\dotprod\vec{s}}
           \cdot \vec{s}\,\bigr)
        =
        \bigl(\frac{\vec{v}\dotprod\vec{s}}{\vec{s}\dotprod\vec{s}}\bigr)^2
        \cdot(\vec{s}\dotprod\vec{s})
        =
        \frac{(\vec{v}\dotprod\vec{s})^2}{(\vec{s}\dotprod\vec{s})^2}
        \cdot(\vec{s}\dotprod\vec{s})
        =
        \frac{(\vec{v}\dotprod\vec{s})^2}{\absval{\vec{s}\,}^2}
      \end{equation*}
于是长度为 $|\vec{v}\dotprod\vec{s}|/|\vec{s}|$。由柯西–施瓦茨不等式可知它不超过 $|\vec{v}|$。
    
\end{ans}
\begin{ans}{三.VI.1.19}
写 $\vec{q}=c\vec{s}$，直接计算即可。
    
\end{ans}
\begin{ans}{三.VI.1.20}
      \begin{exparts}
固定
          \begin{equation*}
            \vec{s}=\colvec[r]{1 \\ 1}
          \end{equation*}
作为直线方向向量，公式给出以下作用，
          \begin{equation*}
            \colvec{x \\ y}
             \mapsto
            \frac{\colvec{x \\ y}\dotprod\colvec[r]{1 \\ 1}}{%
                  \colvec[r]{1 \\ 1}\dotprod\colvec[r]{1 \\ 1}}\cdot\colvec[r]{1 \\ 1}
            =\frac{x+y}{2}\cdot\colvec[r]{1 \\ 1}
            =\colvec{(x+y)/2 \\ (x+y)/2}
          \end{equation*}
这正是该矩阵的作用。
          \begin{equation*}
            \begin{mat}[r]
              1/2  &1/2  \\
              1/2  &1/2
            \end{mat}
          \end{equation*}
顺时针旋转 $\pi/4$ 将 $y=x$ 变到 $x$ 轴；投影后再旋转回来即可。
      \end{exparts}
    
\end{ans}
\begin{ans}{三.VI.1.21}
序列不一定稳定。取
      \begin{equation*}
        \vec{a}=\colvec[r]{1 \\ 0}
        \qquad
        \vec{b}=\colvec[r]{1 \\ 1}
      \end{equation*}
投影为
      \begin{equation*}
        \vec{v}_1=\colvec[r]{1/2 \\ 1/2},
        \quad
        \vec{v}_2=\colvec[r]{1/2 \\ 0},
        \quad
        \vec{v}_3=\colvec[r]{1/4 \\ 1/4},
        \quad
        \ldots
      \end{equation*}
      这个序列不会重复。
      \begin{center}  \small
        \includegraphics{map/mp/ch3.82}
      \end{center}
      
\end{ans}
\protect \subsection {三.VI.2: 格拉姆–施密特正交化}
\begin{ans}{三.VI.2.10}
      \begin{exparts}
        \partsitem $\colvec{1/\sqrt{5} \\ 2/\sqrt{5}}
             =\frac{1}{\sqrt{5}}\cdot\colvec{1 \\ 2}$
        \partsitem $\frac{1}{\sqrt{10}}\cdot\colvec{-1 \\ 3 \\ 0}$
        \partsitem $\frac{1}{\sqrt{2}}\cdot\colvec{1 \\ -1}$
      \end{exparts}
    
\end{ans}
\begin{ans}{三.VI.2.11}
将给定基记为 $B=\sequence{\vec{\beta}_1,\vec{\beta}_2}$。先取 $\vec{\kappa}_1=\vec{\beta}_1$，再取 $\vec{\kappa}_2=\vec{\beta}_2-\proj{\vec{\beta}_2}{\spanof{\vec{\kappa}_1}}$。
      \begin{equation*}
        \vec{\kappa}_2=
        \colvec{-1 \\ 2}-\frac{\colvec{-1 \\ 2}\dotprod\colvec{1 \\ 1}}{\colvec{1 \\ 1}\dotprod\colvec{1 \\ 1}}\cdot\colvec{1 \\ 1}
        =\colvec{-3/2 \\ 3/2}
      \end{equation*}
只需计算两者点积为零即可验证正交。
    
\end{ans}
\begin{ans}{三.VI.2.12}
记给定基为 $B=\sequence{\vec{\beta}_1,\vec{\beta}_2,\vec{\beta}_3}$。先取 $\vec{\kappa}_1=\vec{\beta}_1$。

接着取 $\vec{\kappa}_2=\vec{\beta}_2-\proj{\vec{\beta}_2}{\spanof{\vec{\kappa}_1}}$。
      \begin{equation*}
        \vec{\kappa}_2=
        \colvec{2 \\ 1 \\ -3}-\frac{\colvec{2 \\ 1 \\ -3}\dotprod\colvec{1 \\ 2 \\ 3}}{\colvec{1 \\ 2 \\ 3}\dotprod\colvec{1 \\ 2 \\ 3}}\cdot\colvec{1 \\ 2 \\ 3}
        =\colvec{33/14 \\ 24/14 \\ -27/14}
      \end{equation*}

第三个向量的公式为
      \begin{equation*}
         \vec{\kappa}_3
          =\vec{\beta}_3
           -\proj{\vec{\beta}_3}{\spanof{\vec{\kappa}_1}}
           -\proj{\vec{\beta}_3}{\spanof{\vec{\kappa}_2}}
      \end{equation*}
计算如下。
      \begin{equation*}
        \colvec{3 \\ 3 \\ 3}
        -\frac{\colvec{3 \\ 3 \\ 3}\dotprod\colvec{1 \\ 2 \\ 3}}{\colvec{1 \\ 2 \\ 3}\dotprod\colvec{1 \\ 2 \\ 3}}\cdot\colvec{1 \\ 2 \\ 3}
        -\frac{\colvec{3 \\ 3 \\ 3}\dotprod\colvec{33/14 \\ 24/14 \\ -27/14}}{\colvec{33/14 \\ 24/14 \\ -27/14}\dotprod\colvec{33/14 \\ 24/14 \\ -27/14}}\cdot\colvec{33/14 \\ 24/14 \\ -27/14}
       =\colvec{9/19 \\ -9/19 \\ 3/19}
      \end{equation*}
      % 1 - 18/14-(0/14)/(1274/196) times \colvec{3 \\ 3 \\ 3}
    
\end{ans}
\begin{ans}{三.VI.2.13}
      \begin{exparts}
       \partsitem
        \begin{align*}
          \vec{\kappa}_1 &= \colvec[r]{1 \\ 1}           \\
          \vec{\kappa}_2
            &=
            \colvec[r]{2 \\ 1}
            -\proj{\colvec[r]{2 \\ 1}}{\spanof{\vec{\kappa}_1}}
            =
            \colvec[r]{2 \\ 1}
            -\frac{\colvec[r]{2 \\ 1}\dotprod\colvec[r]{1 \\ 1}}{%
                    \colvec[r]{1 \\ 1}\dotprod\colvec[r]{1 \\ 1}}
            \cdot\colvec[r]{1 \\ 1}                                \\
            &\quad=
            \colvec[r]{2 \\ 1}
            -\frac{3}{2}
            \cdot\colvec[r]{1 \\ 1}
            =\colvec[r]{1/2 \\ -1/2}
        \end{align*}
对应的标准正交基为
        \begin{equation*}
          \sequence{
              \colvec[r]{1/\sqrt{2} \\ 1/\sqrt{2}},
              \colvec[r]{\sqrt{2}/2 \\ -\sqrt{2}/2}
                }
        \end{equation*}
       \partsitem
        \begin{align*}
          \vec{\kappa}_1 &= \colvec[r]{0 \\ 1}           \\
          \vec{\kappa}_2
            &=
            \colvec[r]{-1 \\ 3}
            -\proj{\colvec[r]{-1 \\ 3}}{\spanof{\vec{\kappa}_1}}
            =
            \colvec[r]{-1 \\ 3}
            -\frac{\colvec[r]{-1 \\ 3}\dotprod\colvec[r]{0 \\ 1}}{%
                    \colvec[r]{0 \\ 1}\dotprod\colvec[r]{0 \\ 1}}
            \cdot\colvec[r]{0 \\ 1}                                \\
            &\quad=
            \colvec[r]{-1 \\ 3}
            -\frac{3}{1}
            \cdot\colvec[r]{0 \\ 1}
            =\colvec[r]{-1 \\ 0}
        \end{align*}
这里是标准正交基。
        \begin{equation*}
          \sequence{
                \colvec[r]{0 \\ 1},
                \colvec[r]{-1 \\ 0}
                }
        \end{equation*}
       \partsitem
        \begin{align*}
          \vec{\kappa}_1 &= \colvec[r]{0 \\ 1}           \\
          \vec{\kappa}_2
            &=
            \colvec[r]{-1 \\ 0}
            -\proj{\colvec[r]{-1 \\ 0}}{\spanof{\vec{\kappa}_1}}
            =
            \colvec[r]{-1 \\ 0}
            -\frac{\colvec[r]{-1 \\ 0}\dotprod\colvec[r]{0 \\ 1}}{%
                    \colvec[r]{0 \\ 1}\dotprod\colvec[r]{0 \\ 1}}
            \cdot\colvec[r]{0 \\ 1}                               \\
            &\quad=
            \colvec[r]{-1 \\ 0}
            -\frac{0}{1}
            \cdot\colvec[r]{0 \\ 1}
            =\colvec[r]{-1 \\ 0}
        \end{align*}
      \end{exparts}
相应的标准正交基为
      \begin{equation*}
        \sequence{
              \colvec[r]{0 \\ 1},
              \colvec[r]{-1 \\ 0}
              }
      \end{equation*}
    
\end{ans}
\begin{ans}{三.VI.2.14}
      \begin{exparts}
       \partsitem The first basis vector is unchanged.
        \begin{equation*}
          \vec{\kappa}_1 = \colvec[r]{2 \\ 2 \\ 2}
        \end{equation*}
        第二个向量由以下计算得到。
        \begin{align*}
          \vec{\kappa}_2
            &=
            \colvec[r]{1 \\ 0 \\ -1}
            -\proj{\colvec[r]{1 \\ 0 \\ -1}}{\spanof{\vec{\kappa}_1}}
            =
            \colvec[r]{1 \\ 0 \\ -1}
            -\frac{\colvec[r]{1 \\ 0 \\ -1}\dotprod\colvec[r]{2 \\ 2 \\ 2}}{%
                    \colvec[r]{2 \\ 2 \\ 2}\dotprod\colvec[r]{2 \\ 2 \\ 2}}
            \cdot\colvec[r]{2 \\ 2 \\ 2}                                   \\
            &=
            \colvec[r]{1 \\ 0 \\ -1}
            -\frac{0}{12}
            \cdot\colvec[r]{2 \\ 2 \\ 2}
            =\colvec[r]{1 \\ 0 \\ -1}
        \end{align*}
        第三个向量的算术稍繁琐，但计算过程直接。
        \begin{align*}
          \vec{\kappa}_3
            &=
            \colvec[r]{0 \\ 3 \\ 1}
            -\proj{\colvec[r]{0 \\ 3 \\ 1}}{\spanof{\vec{\kappa}_1}}
            -\proj{\colvec[r]{0 \\ 3 \\ 1}}{\spanof{\vec{\kappa}_2}}   \\
            &=
            \colvec[r]{0 \\ 3 \\ 1}
            -\frac{\colvec[r]{0 \\ 3 \\ 1}\dotprod\colvec[r]{2 \\ 2 \\ 2}}{%
                    \colvec[r]{2 \\ 2 \\ 2}\dotprod\colvec[r]{2 \\ 2 \\ 2}}
            \cdot\colvec[r]{2 \\ 2 \\ 2}
            -\frac{\colvec[r]{0 \\ 3 \\ 1}\dotprod\colvec[r]{1 \\ 0 \\ -1}}{%
                    \colvec[r]{1 \\ 0 \\ -1}\dotprod\colvec[r]{1 \\ 0 \\ -1}}
            \cdot\colvec[r]{1 \\ 0 \\ -1}                                   \\
            &\hbox{}\quad=
            \colvec[r]{0 \\ 3 \\ 1}
            -\frac{8}{12}
            \cdot\colvec[r]{2 \\ 2 \\ 2}
            -\frac{-1}{2}
            \cdot\colvec[r]{1 \\ 0 \\ -1}
            =\colvec[r]{-5/6 \\ 5/3 \\ -5/6}
        \end{align*}
        这就是标准正交基。
        \begin{equation*}
          \sequence{
              \colvec[r]{1/\sqrt{3} \\ 1/\sqrt{3} \\ 1/\sqrt{3}},
              \colvec[r]{1/\sqrt{2} \\ 0 \\ -1/\sqrt{2}},
              \colvec[r]{-1/\sqrt{6} \\ 2/\sqrt{6} \\ -1/\sqrt{6}}
                }
        \end{equation*}
       \partsitem 第一个基向量就是题目给出的向量。
        \begin{equation*}
          \vec{\kappa}_1 = \colvec[r]{1 \\ -1 \\ 0}
        \end{equation*}
        第二个向量为
        \begin{align*}
          \vec{\kappa}_2
            &=
            \colvec[r]{0 \\ 1 \\ 0}
            -\proj{\colvec[r]{0 \\ 1 \\ 0}}{\spanof{\vec{\kappa}_1}}
            =
            \colvec[r]{0 \\ 1 \\ 0}
            -\frac{\colvec[r]{0 \\ 1 \\ 0}\dotprod\colvec[r]{1 \\ -1 \\ 0}}{%
                    \colvec[r]{1 \\ -1 \\ 0}\dotprod\colvec[r]{1 \\ -1 \\ 0}}
            \cdot\colvec[r]{1 \\ -1 \\ 0}                                   \\
            &=
            \colvec[r]{0 \\ 1 \\ 0}
            -\frac{-1}{2}
            \cdot\colvec[r]{1 \\ -1 \\ 0}
            =\colvec[r]{1/2 \\ 1/2 \\ 0}
         \end{align*}
         第三个向量为
        \begin{align*}
          \vec{\kappa}_3
            &=
            \colvec[r]{2 \\ 3 \\ 1}
            -\proj{\colvec[r]{2 \\ 3 \\ 1}}{\spanof{\vec{\kappa}_1}}
            -\proj{\colvec[r]{2 \\ 3 \\ 1}}{\spanof{\vec{\kappa}_2}}       \\
            &=
            \colvec[r]{2 \\ 3 \\ 1}
            -\frac{\colvec[r]{2 \\ 3 \\ 1}\dotprod\colvec[r]{1 \\ -1 \\ 0}}{%
                    \colvec[r]{1 \\ -1 \\ 0}\dotprod\colvec[r]{1 \\ -1 \\ 0}}
            \cdot\colvec[r]{1 \\ -1 \\ 0}
            -\frac{\colvec[r]{2 \\ 3 \\ 1}\dotprod\colvec[r]{1/2 \\ 1/2 \\ 0}}{%
                    \colvec[r]{1/2 \\ 1/2 \\ 0}\dotprod\colvec[r]{1/2 \\ 1/2 \\ 0}}
            \cdot\colvec[r]{1/2 \\ 1/2 \\ 0}                   \\
            &=
            \colvec[r]{2 \\ 3 \\ 1}
            -\frac{-1}{2}
            \cdot\colvec[r]{1 \\ -1 \\ 0}
            -\frac{5/2}{1/2}
            \cdot\colvec[r]{1/2 \\ 1/2 \\ 0}
            =\colvec[r]{0 \\ 0 \\ 1}
        \end{align*}
        相应的标准正交基为
        \begin{equation*}
          \sequence{
              \colvec[r]{1/\sqrt{2} \\ -1/\sqrt{2} \\ 0},
              \colvec[r]{1/\sqrt{2} \\ 1/\sqrt{2} \\ 0}
              \colvec[r]{0 \\ 0 \\ 1}
                }
        \end{equation*}
      \end{exparts}
    
\end{ans}
\begin{ans}{三.VI.2.15}
         将给定空间参数化为
         \begin{equation*}
           \set{\colvec{x \\ y \\ z} \suchthat x=y-z}
           =\set{\colvec[r]{1 \\ 1 \\ 0}\cdot y+\colvec[r]{-1 \\ 0 \\ 1}\cdot z
                     \suchthat y,z\in\Re}
         \end{equation*}
         因而取基
         \begin{equation*}
           \sequence{\colvec[r]{1 \\ 1 \\ 0},
                     \colvec[r]{-1 \\ 0 \\ 1}}
         \end{equation*}
         应用 Gram–Schmidt 过程，得到第一个基向量
         \begin{equation*}
           \vec{\kappa}_1 = \colvec[r]{1 \\ 1 \\ 0}
         \end{equation*}
         以及第二个基向量。
         \begin{align*}
           \vec{\kappa}_2
           &=
           \colvec[r]{-1 \\ 0 \\ 1}-\proj{\colvec[r]{-1 \\ 0 \\ 1}}{%
                                        \spanof{\vec{\kappa}_1}}
           =
           \colvec[r]{-1 \\ 0 \\ 1}
            -\frac{\colvec[r]{-1 \\ 0 \\ 1}\dotprod\colvec[r]{1 \\ 1 \\ 0}}{%
                    \colvec[r]{1 \\ 1 \\ 0}\dotprod\colvec[r]{1 \\ 1 \\ 0}}
              \cdot\colvec[r]{1 \\ 1 \\ 0}                                \\
           &=
           \colvec[r]{-1 \\ 0 \\ 1}
            -\frac{-1}{2}\cdot\colvec[r]{1 \\ 1 \\ 0}
           =\colvec[r]{-1/2 \\ 1/2 \\ 1}
         \end{align*}
         最后将它们单位化。
         \begin{equation*}
           \sequence{
                    \colvec[r]{1/\sqrt{2} \\ 1/\sqrt{2} \\ 0},
                    \colvec[r]{-1/\sqrt{6} \\ 1/\sqrt{6} \\ 2/\sqrt{6}}
                    }
         \end{equation*}
       
\end{ans}
\begin{ans}{三.VI.2.16}
       化简线性方程组
       \begin{equation*}
         \begin{linsys}{4}
           x  &-  &y  &-  &z  &+  &w  &=  &0  \\
           x  &   &   &+  &z  &   &   &=  &0
         \end{linsys}
         \grstep{-\rho_1+\rho_2}
         \begin{linsys}{4}
           x  &-  &y  &-  &z  &+  &w  &=  &0  \\
              &   &y  &+  &2z &-  &w  &=  &0
         \end{linsys}
       \end{equation*}
       参数化后得到该子空间的如下描述。
       \begin{equation*}
         \set{\colvec[r]{-1 \\ -2 \\ 1 \\ 0}\cdot z
              +\colvec[r]{0 \\ 1 \\ 0 \\ 1}\cdot w
              \suchthat z,w\in\Re}
       \end{equation*}
       因而取基
       \begin{equation*}
         \sequence{
                  \colvec[r]{-1 \\ -2 \\ 1 \\ 0},
                  \colvec[r]{0 \\ 1 \\ 0 \\ 1}
                  }
       \end{equation*}
       对其执行 Gram–Schmidt 过程，第一个
       \begin{equation*}
         \vec{\kappa}_1 = \colvec[r]{-1 \\ -2 \\ 1 \\ 0}
       \end{equation*}
       和第二个基向量分别为
       \begin{align*}
         \vec{\kappa}_2
          &=
          \colvec[r]{0 \\ 1 \\ 0 \\ 1}
          -\proj{\colvec[r]{0 \\ 1 \\ 0 \\ 1}}{\spanof{\vec{\kappa}_1}}
          =
          \colvec[r]{0 \\ 1 \\ 0 \\ 1}
         -\frac{\colvec[r]{0 \\ 1 \\ 0 \\ 1}\dotprod\colvec[r]{-1 \\ -2 \\ 1 \\ 0}}{%
             \colvec[r]{-1 \\ -2 \\ 1 \\ 0}\dotprod\colvec[r]{-1 \\ -2 \\ 1 \\ 0}}
          \cdot\colvec[r]{-1 \\ -2 \\ 1 \\ 0}                            \\
          &=
          \colvec[r]{0 \\ 1 \\ 0 \\ 1}
         -\frac{-2}{6}
          \cdot\colvec[r]{-1 \\ -2 \\ 1 \\ 0}
          =\colvec[r]{-1/3 \\ 1/3 \\ 1/3 \\ 1}
       \end{align*}
       最后将结果单位化。
       \begin{equation*}
         \sequence{
               \colvec[r]{-1/\sqrt{6} \\ -2/\sqrt{6} \\ 1/\sqrt{6} \\ 0},
               \colvec[r]{-\sqrt{3}/6 \\ \sqrt{3}/6 \\ \sqrt{3}/6 \\ \sqrt{3}/2}
                  }
       \end{equation*}
     
\end{ans}
\begin{ans}{三.VI.2.17}
       $\Re^n$ 的线性无关子集是其张成空间的一组基，因此应用定理 \nearbytheorem{th:GramSchmidt}。

       \textit{Remark}.
       题目要求“线性无关”是有原因的。删去这一条件后，需要担心两件事。第一，对线性相关集合执行 Gram–Schmidt 过程会产生零向量。
       例如, with
       \begin{equation*}
         S=\set{\colvec[r]{1 \\ 2},\colvec[r]{3 \\ 6}}
       \end{equation*}
       得到：
       \begin{equation*}
         \vec{\kappa}_1   =   \colvec[r]{1 \\ 2}
         \qquad
         \vec{\kappa}_2
           =
           \colvec[r]{3 \\ 6}-\proj{\colvec[r]{3 \\ 6}}{\spanof{\vec{\kappa}_1}}
           =
           \colvec[r]{0 \\ 0}
       \end{equation*}
       这并不严重，因为零向量按定义与所有向量正交；我们可以接受所得的正交集合（尽管零向量无法单位化），也可以修改 Gram–Schmidt 过程，删去所有零向量。第二个问题是集合可能无限。
       有限维 $\Re^n$ 的任意子空间也有限维，因此其中只有有限多个向量能线性无关；不过，逐个检查无限集合中向量的“过程”仍需进一步说明。$\Re^n$ 的线性无关子集自动是有限的\Dash 事实上大小不超过 $n$\Dash 所以这些担忧都不存在。
     
\end{ans}
\begin{ans}{三.VI.2.18}
      如果集合线性相关，就会得到零向量。否则（集合线性无关但不张成整个空间），我们是在一个子空间的基上执行 Gram–Schmidt，因此得到该子空间的一组正交基。
    
\end{ans}
\begin{ans}{三.VI.2.19}
       该过程保持基不变。
     
\end{ans}
\begin{ans}{三.VI.2.20}
       \begin{exparts}
         \partsitem 论证与定理 \nearbytheorem{th:GramSchmidt} 证明中的 $i=3$ 情形相同。
           点积
           \begin{equation*}
             \vec{\kappa}_i\dotprod
             \left(\vec{v}-\proj{\vec{v}\,}{\spanof{\vec{\kappa}_1}}
               -\dots-\proj{\vec{v}\,}{\spanof{\vec{v}_k}}\right)
           \end{equation*}
           可以写成以下两类项之和：
           $-\vec{\kappa}_i\dotprod\proj{\vec{v}\,}{\spanof{\vec{\kappa}_j}}$
           其中 $j\neq i$，
           以及项
           $\vec{\kappa}_i\dotprod(\vec{v}-
                                 \proj{\vec{v}\,}{\spanof{\vec{\kappa}_i}})$.
           第一类项因各 $\vec\kappa$ 两两正交而为零。另一类项因投影正交而为零
           （即投影定义使其为零：
           $\vec{\kappa}_i\dotprod(\vec{v}
               -\proj{\vec{v}\,}{\spanof{\vec{\kappa}_i}})
            =\vec{\kappa}_i\dotprod\vec{v}-
              \vec{\kappa}_i\dotprod
               \left((\vec{v}\dotprod\vec{\kappa}_i)/(%
                      \vec{\kappa}_i\dotprod\vec{\kappa}_i)\right)
                   \cdot\vec{\kappa}_i$
           经过约分后等于零）。
         \partsitem 向量 $\vec v$ 用黑色表示，而
           vector $\proj{\vec{v}\,}{\spanof{\vec{\kappa}_1}}
                    +\proj{\vec{v}\,}{\spanof{\vec{v}_2}}
                   =1\cdot\vec{e}_1+2\cdot\vec{e}_2$ is in gray.
           \begin{center}  \small
             \includegraphics{map/mp/ch3.83}
              \end{center}
              The vector
              $\vec{v}-(\proj{\vec{v}\,}{\spanof{\vec{\kappa}_1}}
                    +\proj{\vec{v}\,}{\spanof{\vec{v}_2}})$
              位于连接黑色向量与灰色向量的点划线上，也就是与 $xy$ 平面正交。
          \partsitem 按提示即可得到该图。
           \begin{center}  \small
             \includegraphics{map/mp/ch3.84}
            \end{center}
            虚线三角形的直角位于
            灰色向量 $1\cdot\vec{e}_1+2\cdot\vec{e}_2$
            与竖直虚线相交处；
            $\vec{v}-(1\cdot\vec{e}_1+2\cdot\vec{e}_2)$; this is what
            这正是本题第一问证明的结论。由勾股定理，斜边\Dash 从 $\vec v$ 到任意其他向量的线段\Dash 长于竖直虚线。

            更形式化地，将 $\proj{\vec{v}\,}{\spanof{\vec{\kappa}_1}}
                 +\dots+\proj{\vec{v}\,}{\spanof{\vec{v}_k}}$ as
            $c_1\cdot\vec{\kappa}_1+\dots+c_k\cdot\vec{\kappa}_k$,
            再任取张成空间中的另一个向量
            $d_1\cdot\vec{\kappa}_1+\dots+d_k\cdot\vec{\kappa}_k$.
            Note that
            \begin{multline*}
              \vec{v}-(d_1\cdot\vec{\kappa}_1+\dots+d_k\cdot\vec{\kappa}_k)  \\
              \begin{aligned}
              &=
              \bigl(\vec{v}
                    -(c_1\cdot\vec{\kappa}_1+\dots+c_k\cdot\vec{\kappa}_k)
              \bigr)                                                     \\
              &\quad+\bigl((c_1\cdot\vec{\kappa}_1+\dots+c_k\cdot\vec{\kappa}_k)
                     -(d_1\cdot\vec{\kappa}_1+\dots+d_k\cdot\vec{\kappa}_k)
               \bigr)
              \end{aligned}
            \end{multline*}
            and that
            $
              \bigl(\vec{v}
                 -(c_1\cdot\vec{\kappa}_1+\dots+c_k\cdot\vec{\kappa}_k)
              \bigr)
              \dotprod
              \bigl((c_1\cdot\vec{\kappa}_1+\dots+c_k\cdot\vec{\kappa}_k)
                     -(d_1\cdot\vec{\kappa}_1+\dots+d_k\cdot\vec{\kappa}_k)
              \bigr)
              =0
            $
            (because the first item shows the $\vec{v}
                 -(c_1\cdot\vec{\kappa}_1+\dots+c_k\cdot\vec{\kappa}_k)$
            与每个 $\vec\kappa$ 正交，因此也与
            to th是线性的 combination of the $\vec{\kappa}$'s).
            最后应用勾股定理（即三角不等式）。
       \end{exparts}
     
\end{ans}
\begin{ans}{三.VI.2.21}
       一种方法是找第三个向量，使三个向量一起
       构成 $\Re^3$ 的一组基，例如，
       \begin{equation*}
         \vec{\beta}_3=\colvec[r]{1 \\ 0 \\ 0}
       \end{equation*}
       (the second vector is not dependent on the third because it has a
       构成一组基（第二个向量的第二个分量非零，因而不依赖于第三个向量；第一个向量的第三个分量非零，因而不依赖于后两者），然后应用 Gram–Schmidt 过程。新基的第一个向量为
       \begin{equation*}
         \vec{\kappa}_1 = \colvec[r]{1 \\ 5 \\ -1}
       \end{equation*}
       第二个向量为
       \begin{align*}
         \vec{\kappa}_2
           &=
           \colvec[r]{2 \\ 2 \\ 0}
           -\proj{\colvec[r]{2 \\ 2 \\ 0}}{\spanof{\vec{\kappa}_1}}
           =\colvec[r]{2 \\ 2 \\ 0}
           -\frac{\colvec[r]{2 \\ 2 \\ 0}\dotprod\colvec[r]{1 \\ 5 \\ -1}}{%
                  \colvec[r]{1 \\ 5 \\ -1}\dotprod\colvec[r]{1 \\ 5 \\ -1}}
            \cdot\colvec[r]{1 \\ 5 \\ -1}                               \\
           &\hbox{}\quad=
           \colvec[r]{2 \\ 2 \\ 0}
           -\frac{12}{27}\cdot\colvec[r]{1 \\ 5 \\ -1}
           =\colvec[r]{14/9 \\ -2/9 \\ 4/9}
       \end{align*}
       最后一个向量为
       \begin{align*}
        \vec{\kappa}_3
           &=
           \colvec[r]{1 \\ 0 \\ 0}
             -\proj{\colvec[r]{1 \\ 0 \\ 0}}{\spanof{\vec{\kappa}_1}}
             -\proj{\colvec[r]{1 \\ 0 \\ 0}}{\spanof{\vec{\kappa}_2}}       \\
           &\hbox{}\quad=
           \colvec[r]{1 \\ 0 \\ 0}
           -\frac{\colvec[r]{1 \\ 0 \\ 0}\dotprod\colvec[r]{1 \\ 5 \\ -1}}{%
                    \colvec[r]{1 \\ 5 \\ -1}\dotprod\colvec[r]{1 \\ 5 \\ -1}}
             \cdot\colvec[r]{1 \\ 5 \\ -1}
           -\frac{\colvec[r]{1 \\ 0 \\ 0}\dotprod\colvec[r]{14/9 \\ -2/9 \\  4/9}}{%
            \colvec[r]{14/9 \\ -2/9 \\ 4/9}\dotprod\colvec[r]{14/9 \\ -2/9 \\ 4/9}}
            \cdot\colvec[r]{14/9 \\ -2/9 \\ 4/9}              \\
         &=
         \colvec[r]{1 \\ 0 \\ 0}
           -\frac{1}{27}\cdot\colvec[r]{1 \\ 5 \\ -1}
           -\frac{7}{12}\cdot\colvec[r]{14/9 \\ -2/9 \\ 4/9}
         =\colvec[r]{1/18 \\ -1/18  \\ -4/18}
       \end{align*}
       结果 $\vec{\kappa}_3$ is orthogonal to both $\vec{\kappa}_1$
       and $\vec{\kappa}_2$.
       因此它与张成空间中的每个向量正交，
       属于集合 $\set{\vec{\kappa}_1,\vec{\kappa}_2}$ 的张成空间，
       包括题目给出的两个向量。
     
\end{ans}
\begin{ans}{三.VI.2.22}
      \begin{exparts}
        \partsitem We can do the representation by eye.
          \begin{equation*}
            \colvec[r]{2 \\ 3}=3\cdot\colvec[r]{1 \\ 1}+(-1)\cdot\colvec[r]{1 \\ 0}
            \qquad
            \rep{\vec{v}\,}{B}=\colvec[r]{3 \\ -1}_B
          \end{equation*}
          两个投影也很容易求出。
          \begin{align*}
            \proj{\colvec[r]{2 \\ 3}}{\spanof{\vec{\beta}_1}}
            &=\frac{\colvec[r]{2 \\ 3}\dotprod\colvec[r]{1 \\ 1}}{%
                   \colvec[r]{1 \\ 1}\dotprod\colvec[r]{1 \\ 1}}
              \cdot\colvec[r]{1 \\ 1}
            =\frac{5}{2}\cdot\colvec[r]{1 \\ 1}            \\
            \proj{\colvec[r]{2 \\ 3}}{\spanof{\vec{\beta}_2}}
            &=\frac{\colvec[r]{2 \\ 3}\dotprod\colvec[r]{1 \\ 0}}{%
                   \colvec[r]{1 \\ 0}\dotprod\colvec[r]{1 \\ 0}}
              \cdot\colvec[r]{1 \\ 0}
            =\frac{2}{1}\cdot\colvec[r]{1 \\ 0}
          \end{align*}
        \partsitem As above, we can do the representation by eye
          \begin{equation*}
            \colvec[r]{2 \\ 3}=(5/2)\cdot\colvec[r]{1 \\ 1}
                             +(-1/2)\cdot\colvec[r]{1 \\ -1}
          \end{equation*}
          两个投影也很容易求出。
          \begin{align*}
            \proj{\colvec[r]{2 \\ 3}}{\spanof{\vec{\beta}_1}}
            &=\frac{\colvec[r]{2 \\ 3}\dotprod\colvec[r]{1 \\ 1}}{%
                   \colvec[r]{1 \\ 1}\dotprod\colvec[r]{1 \\ 1}}
              \cdot\colvec[r]{1 \\ 1}
            =\frac{5}{2}\cdot\colvec[r]{1 \\ 1}                    \\
            \proj{\colvec[r]{2 \\ 3}}{\spanof{\vec{\beta}_2}}
            &=\frac{\colvec[r]{2 \\ 3}\dotprod\colvec[r]{1 \\ -1}}{%
                   \colvec[r]{1 \\ -1}\dotprod\colvec[r]{1 \\ -1}}
              \cdot\colvec[r]{1 \\ -1}
            =\frac{-1}{2}\cdot\colvec[r]{1 \\ -1}
          \end{align*}
          注意 $5/2$ 和 $-1/2$ 再次出现。
        \partsitem Represent $\vec{v}$ with respect to the basis
          \begin{equation*}
            \rep{\vec{v}\,}{K}=\colvec{r_1 \\ \vdots \\ r_k}
          \end{equation*}
          so that
          $\vec{v}=r_1\vec{\kappa}_1+\dots+r_k\vec{\kappa}_k$.
          为确定 $r_i$，
          将等式两边与 $\vec{\kappa}_i$ 作点积。
          \begin{equation*}
            \vec{v}\dotprod\vec{\kappa}_i
               =\left(r_1\vec{\kappa}_1+\dots+r_k\vec{\kappa}_k\right)
                 \dotprod\vec{\kappa}_i
               =r_1\cdot 0+\dots+r_i\cdot(\vec{\kappa}_i\dotprod\vec{\kappa}_i)
                    +\dots+r_k\cdot 0
          \end{equation*}
          解出 $r_i$ 即得所需系数。
        \partsitem 这是 a restatement of the prior item.
      \end{exparts}
    
\end{ans}
\begin{ans}{三.VI.2.23}
      首先，$\absval{\vec v}^2=4^2+3^2+2^2+1^2=50$。
      \begin{exparts*}
        \partsitem $c_1=4$
        \partsitem $c_1=4$, $c_2=3$
        \partsitem $c_1=4$, $c_2=3$, $c_3=2$, $c_4=1$
      \end{exparts*}
      证明只讨论 $k=2$ 的情形，因为一般情形更繁琐但没有增加新的思想。按照提示，并回忆任意向量 $\vec w$ 都满足 $\absval{\vec w}^2=\vec w\dotprod\vec w$。
      \begin{align*}
        0
         &\leq
         \left(\vec{v}-\bigl(\frac{\vec{v}\dotprod\vec{\kappa}_1}{%
                              \vec{\kappa}_1\dotprod\vec{\kappa}_1}
                          \cdot\vec{\kappa}_1
                         +\frac{\vec{v}\dotprod\vec{\kappa}_2}{%
                              \vec{\kappa}_2\dotprod\vec{\kappa}_2}
                          \cdot\vec{\kappa}_2
                       \bigr)
         \right)
         \dotprod
         \left(\vec{v}-\bigl(\frac{\vec{v}\dotprod\vec{\kappa}_1}{%
                              \vec{\kappa}_1\dotprod\vec{\kappa}_1}
                          \cdot\vec{\kappa}_1
                         +\frac{\vec{v}\dotprod\vec{\kappa}_2}{%
                              \vec{\kappa}_2\dotprod\vec{\kappa}_2}
                          \cdot\vec{\kappa}_2
                       \bigr)
         \right)                                                       \\
         &=\begin{aligned}[t]
            &\vec{v}\dotprod\vec{v}
             -2\cdot\vec{v}
              \dotprod\left(\frac{\vec{v}\dotprod\vec{\kappa}_1}{%
                              \vec{\kappa}_1\dotprod\vec{\kappa}_1}
                          \cdot\vec{\kappa}_1
                       +\frac{\vec{v}\dotprod\vec{\kappa}_2}{%
                              \vec{\kappa}_2\dotprod\vec{\kappa}_2}
                          \cdot\vec{\kappa}_2\right)               \\
            &+\left(\frac{\vec{v}\dotprod\vec{\kappa}_1}{%
                              \vec{\kappa}_1\dotprod\vec{\kappa}_1}
                          \cdot\vec{\kappa}_1
                       +\frac{\vec{v}\dotprod\vec{\kappa}_2}{%
                              \vec{\kappa}_2\dotprod\vec{\kappa}_2}
                          \cdot\vec{\kappa}_2\right)
           \dotprod
           \left(\frac{\vec{v}\dotprod\vec{\kappa}_1}{%
                              \vec{\kappa}_1\dotprod\vec{\kappa}_1}
                          \cdot\vec{\kappa}_1
                       +\frac{\vec{v}\dotprod\vec{\kappa}_2}{%
                              \vec{\kappa}_2\dotprod\vec{\kappa}_2}
                          \cdot\vec{\kappa}_2\right)
          \end{aligned}                                              \\
         &=
          \vec{v}\dotprod\vec{v}
          -2\cdot
             \left(\frac{\vec{v}\dotprod\vec{\kappa}_1}{%
                              \vec{\kappa}_1\dotprod\vec{\kappa}_1}
                          \cdot(\vec{v}\dotprod\vec{\kappa}_1)
                       +\frac{\vec{v}\dotprod\vec{\kappa}_2}{%
                              \vec{\kappa}_2\dotprod\vec{\kappa}_2}
                          \cdot(\vec{v}\dotprod\vec{\kappa}_2)\right)  \\
          &\quad+\left((\frac{\vec{v}\dotprod\vec{\kappa}_1}{%
                              \vec{\kappa}_1\dotprod\vec{\kappa}_1})^2
                          \cdot(\vec{\kappa}_1\dotprod\vec{\kappa}_1)
                       +(\frac{\vec{v}\dotprod\vec{\kappa}_2}{%
                              \vec{\kappa}_2\dotprod\vec{\kappa}_2})^2
                          \cdot(\vec{\kappa}_2\dotprod\vec{\kappa}_2)\right)
      \end{align*}
      （第三行第三部分中的两个混合项为零，因为 $\vec\kappa_1,\vec\kappa_2$ 正交。）合并同类项，并注意它们属于标准正交集，故 $\vec\kappa_1\dotprod\vec\kappa_1=\vec\kappa_2\dotprod\vec\kappa_2=1$，即可得到结论。
    
\end{ans}
\begin{ans}{三.VI.2.24}
      除零向量外命题成立。$\Re^n$ 中每个非零向量都属于某组基，而该基可以正交化。
     
\end{ans}
\begin{ans}{三.VI.2.25}
      $\nbyn 3$ 的情形说明了思路。
      The set
      \begin{equation*}
        \set{\colvec{a \\ d \\ g},\colvec{b \\ e \\ h},\colvec{c \\ f \\ i}}
      \end{equation*}
      标准正交，当且仅当以下九个条件全部成立
      \begin{equation*}
        \begin{array}{ccc}
          \rowvec{a &d &g}\dotprod\colvec{a \\ d \\ g}=1
            &\rowvec{a &d &g}\dotprod\colvec{b \\ e \\ h}=0
            &\rowvec{a &d &g}\dotprod\colvec{c \\ f \\ i}=0    \\
          \rowvec{b &e &h}\dotprod\colvec{a \\ d \\ g}=0
            &\rowvec{b &e &h}\dotprod\colvec{b \\ e \\ h}=1
            &\rowvec{b &e &h}\dotprod\colvec{c \\ f \\ i}=0    \\
          \rowvec{c &f &i}\dotprod\colvec{a \\ d \\ g}=0
            &\rowvec{c &f &i}\dotprod\colvec{b \\ e \\ h}=0
            &\rowvec{c &f &i}\dotprod\colvec{c \\ f \\ i}=1
        \end{array}
      \end{equation*}
      （左下角的三个条件是冗余的，但同样正确。）这些条件又等价于
      \begin{equation*}
        \begin{mat}
          a &d &g \\
          b &e &h \\
          c &f &i
        \end{mat}
        \begin{mat}
          a &b &c  \\
          d &e &f  \\
          g &h &i
        \end{mat}
        =
        \begin{mat}
          1 &0 &0  \\
          0 &1 &0  \\
          0 &0 &1
        \end{mat}
      \end{equation*}
      如所需。

      这是一个例子：该矩阵的逆就是它的转置。
      \begin{equation*}
        \begin{mat}[r]
          1/\sqrt{2}  &1/\sqrt{2}  &0  \\
         -1/\sqrt{2}  &1/\sqrt{2}  &0  \\
          0           &0           &1
        \end{mat}
      \end{equation*}
    
\end{ans}
\begin{ans}{三.VI.2.26}
      若集合为空，左侧求和按定义是空集向量的线性组合，即零向量。第二句话中不存在这样的 $i$，所以“若……则……”蕴含命题真空成立。
    
\end{ans}
\begin{ans}{三.VI.2.27}
      \begin{exparts}
        \partsitem 归纳证明
        定理 \nearbytheorem{th:GramSchmidt} 的归纳证明说明 $\vec\kappa_i$ 属于
          $\sequence{\vec{\beta}_1,\dots,\vec{\beta}_i}$.
          （证明中的 $i=3$ 情形展示了这一点。）因此，在换基矩阵 $\rep{\identity}{K,B}$ 中，第 $i$ 列 $\rep{\vec{\kappa}_i}{B}$ 的第 $i+1$ 至第 $k$ 个分量均为零。
        \partsitem 可以回忆求逆时的计算过程来理解这一点：
          可以回忆求逆时的计算过程：将矩阵与其右侧的单位矩阵并排写出，再进行 Gauss–Jordan 消元。如果初始矩阵是上三角矩阵，消元只涉及左半部分；将这些步骤作用于单位矩阵，就会得到上三角的逆矩阵。
      \end{exparts}
    
\end{ans}
\begin{ans}{三.VI.2.28}
      对归纳步骤，假设对所有 $j\in[1..i]$，每个 $\vec\kappa_j$ 都满足三个条件：
      (i)~each $\vec{\kappa}_j$ is nonzero,
      (ii)~each $\vec{\kappa}_j$ is a linear combination of
         向量 $\vec{\beta}_1,\dots,\vec{\beta}_j$，
      and (iii)~each $\vec{\kappa}_j$ is orthogonal to all of the
         $\vec{\kappa}_m$'s prior to it (that is, with $m<j$).
      在这些归纳假设下，考察 $\vec\kappa_{i+1}$。
      \begin{align*}
        \vec{\kappa}_{i+1}
         &=
        \vec{\beta}_{i+1}
          -\proj{\beta_{i+1}}{\spanof{\vec{\kappa}_1}}
          -\proj{\beta_{i+1}}{\spanof{\vec{\kappa}_2}}
          -\dots
          -\proj{\beta_{i+1}}{\spanof{\vec{\kappa}_i}}  \\
         &=
        \vec{\beta}_{i+1}
          -\frac{\beta_{i+1}\dotprod\vec{\kappa}_1}{%
                  \vec{\kappa}_1\dotprod\vec{\kappa}_1}
              \cdot\vec{\kappa}_1
          -\frac{\beta_{i+1}\dotprod\vec{\kappa}_2}{%
                  \vec{\kappa}_2\dotprod\vec{\kappa}_2}
              \cdot\vec{\kappa}_2
          -\dots
          -\frac{\beta_{i+1}\dotprod\vec{\kappa}_i}{%
                  \vec{\kappa}_i\dotprod\vec{\kappa}_i}
              \cdot\vec{\kappa}_i
       \end{align*}
       由归纳假设（ii），可将每个 $\vec\kappa_j$ 展开为 $\vec\beta_1,\ldots,\vec\beta_j$ 的线性组合。
       \begin{align*}
         &=
        \vec{\beta}_{i+1}
          -\frac{\vec{\beta}_{i+1}\dotprod\vec{\kappa}_1}{%
                  \vec{\kappa}_1\dotprod\vec{\kappa}_1}
              \cdot\vec{\beta}_1                          \\
          &\quad-\frac{\vec{\beta}_{i+1}\dotprod\vec{\kappa}_2}{%
                  \vec{\kappa}_2\dotprod\vec{\kappa}_2}
              \cdot
                 \left(\text{
                   $\vec{\beta}_1,\vec{\beta}_2$ 的线性组合}
                 \right)                                                 \\
          &\quad-\;\cdots\;
          -\frac{\vec{\beta}_{i+1}\dotprod\vec{\kappa}_i}{%
                  \vec{\kappa}_i\dotprod\vec{\kappa}_i}
              \cdot
                 \left(\text{
                   $\vec{\beta}_1,\dots,\vec{\beta}_i$ 的线性组合}
                 \right)
       \end{align*}
       这些分式是标量，因此上式是 $\vec\beta_1,\dots,\vec\beta_{i+1}$ 的线性组合的线性组合，也就是它们的线性组合。
       Now, (i)~it cannot sum to the zero vector because the equation would
       否则该等式会描述基向量之间的非平凡线性关系（之所以非平凡，是因为 $\vec\beta_{i+1}$ 的系数为 $1$）。
       同时，（ii）中的等式把 $\vec\kappa_{i+1}$ 表示为 $\vec\beta_1,\dots,\vec\beta_{i+1}$ 的线性组合。最后，对（iii）考察 $\vec\kappa_j\dotprod\vec\kappa_{i+1}$；与 $i=3$ 时一样，$\vec\kappa_j$ 与
       $\vec{\kappa}_{i+1}=\vec{\beta}_{i+1}
          -\proj{\vec{\beta}_{i+1}}{\spanof{\vec{\kappa}_1}}
          -\dots
          -\proj{\vec{\beta}_{i+1}}{\spanof{\vec{\kappa}_i}}$
        可以改写为两类项之和：
       $\vec{\kappa}_j\dotprod
           \left(\vec{\beta}_{i+1}
                  -\proj{\vec{\beta}_{i+1}}{\spanof{\vec{\kappa}_j}}
           \right)$
       （第一类因投影的正交性为零），以及
       $\vec{\kappa}_j\dotprod
               \proj{\vec{\beta}_{i+1}}{\spanof{\vec{\kappa}_m}}$
       其中 $m\neq j$ 且 $m<i+1$（第二类因归纳假设（iii）中 $\vec\kappa_j$ 与 $\vec\kappa_m$ 正交而为零）。
     
\end{ans}
\protect \subsection {三.VI.3: 投影到子空间}
\begin{ans}{三.VI.3.10}
       \begin{exparts}
         \partsitem 将两个子空间的基
           \begin{equation*}
             B_M=\sequence{\colvec[r]{1 \\ -1}}
             \qquad
             B_N=\sequence{\colvec[r]{2 \\ -1}}
           \end{equation*}
           连接起来
           \begin{equation*}
             B=\cat{B_M}{B_N}=\sequence{\colvec[r]{1 \\ -1},
                                        \colvec[r]{2 \\ -1}}
           \end{equation*}
           并表示给定向量
           \begin{equation*}
             \colvec[r]{3 \\ -2}=1\cdot\colvec[r]{1 \\ -1}+1\cdot\colvec[r]{2 \\ -1}
           \end{equation*}
           答案就是保留 $M$ 部分并舍去 $N$ 部分。
           \begin{equation*}
             \proj{\colvec[r]{3 \\ -2}}{M,N}
             =\colvec[r]{1 \\ -1}
           \end{equation*}
         \partsitem 将这两组基连接起来并表示向量后，
           \begin{equation*}
             B_M=\sequence{\colvec[r]{1 \\ 1}}
             \qquad
             B_N\sequence{\colvec[r]{1 \\ -2}}
           \end{equation*}
           连接两组基并表示向量后，
           \begin{equation*}
             \colvec[r]{1 \\ 2}
               =(4/3)\cdot\colvec[r]{1 \\ 1}
                -(1/3)\cdot\colvec[r]{1 \\ -2}
           \end{equation*}
           保留 $M$ 部分即得答案。
           \begin{equation*}
             \proj{\colvec[r]{1 \\ 2}}{M,N}=\colvec[r]{4/3 \\ 4/3}
           \end{equation*}
         \partsitem 对于这些基，
           \begin{equation*}
             B_M=\sequence{\colvec[r]{1 \\ -1 \\ 0},
                           \colvec[r]{0 \\ 0 \\ 1}}
             \qquad
             B_N=\sequence{\colvec[r]{1 \\ 0 \\ 1}}
           \end{equation*}
           关于连接基的表示为
           \begin{equation*}
             \colvec[r]{3 \\ 0 \\ 1}
               =0\cdot\colvec[r]{1 \\ -1 \\ 0}
                -2\cdot\colvec[r]{0 \\ 0 \\ 1}
                +3\cdot\colvec[r]{1 \\ 0 \\ 1}
           \end{equation*}
           因而投影为
           \begin{equation*}
             \proj{\colvec[r]{3 \\ 0 \\ 1}}{M,N}=\colvec[r]{0 \\ 0 \\ -2}
           \end{equation*}
       \end{exparts}
     
\end{ans}
\begin{ans}{三.VI.3.11}
       与例 \nearbyexample{ex:OrthoCompOne} 一样，只需找出与 $M$ 的基中所有向量垂直的向量空间，即可简化计算。
       \begin{exparts}
         \partsitem 参数化得到
           \begin{equation*}
             M=\set{c\cdot\colvec[r]{-1 \\ 1}\suchthat c\in\Re}
           \end{equation*}
           从而
           \begin{equation*}
             M^\perp
              \set{\colvec{u \\ v}
                   \suchthat
                      0=\colvec{u \\ v}\dotprod\colvec[r]{-1 \\ 1}}
              =\set{\colvec{u \\ v}\suchthat 0=-u+v}
           \end{equation*}
           将这个单方程线性系统参数化，得到
           \begin{equation*}
             M^\perp=\set{k\cdot\colvec[r]{1 \\ 1}\suchthat k\in\Re}
           \end{equation*}
         \partsitem 如上一问答案，可将 $M$ 描述为张成空间
           \begin{equation*}
             M=\set{c\cdot\colvec[r]{3/2 \\ 1}\suchthat c\in\Re}
             \qquad
             B_M=\sequence{\colvec[r]{3/2 \\ 1}}
           \end{equation*}
           因而 $M^\perp$ 是与该基中唯一向量垂直的向量集合。
           \begin{equation*}
             M^\perp
             =\set{\colvec{u \\ v}\suchthat (3/2)\cdot u+1\cdot v=0}
             =\set{k\cdot\colvec[r]{-2/3 \\ 1}\suchthat k\in\Re}
           \end{equation*}
         \partsitem 将 $M$ 的描述中的线性条件参数化，得到以下基。
            \begin{equation*}
              M=\set{c\cdot\colvec[r]{1 \\ 1}\suchthat c\in\Re}
              \qquad
              B_M=\sequence{\colvec[r]{1 \\ 1}}
            \end{equation*}
            现在，$M^\perp$ 是与 $B_M$ 中唯一向量垂直的向量集合。
            \begin{equation*}
              M^\perp
              =\set{\colvec{u \\ v}\suchthat u+v=0}
              =\set{k\cdot\colvec[r]{-1 \\ 1}\suchthat k\in\Re}
            \end{equation*}
            （这个答案也与本题第一问相符。）
          \partsitem 空间中的每个向量都与零向量垂直，因此 $M^\perp=\Re^n$。
          \partsitem $M$ 的适当描述和基很容易得到。
            \begin{equation*}
              M=\set{y\cdot\colvec[r]{0 \\ 1}\suchthat y\in\Re}
              \qquad
              B_M=\sequence{\colvec[r]{0 \\ 1}}
            \end{equation*}
            于是
            \begin{equation*}
              M^\perp
              =\set{\colvec{u \\ v}\suchthat 0\cdot u+1\cdot v=0}
              =\set{k\cdot\colvec[r]{1 \\ 0}\suchthat k\in\Re}
            \end{equation*}
            因而 $(\text{$y$-axis})^\perp=\text{$x$-axis}$。
          \partsitem 将 $M$ 参数化即可容易地得到其描述。
            \begin{equation*}
              M=\set{c\cdot\colvec[r]{3 \\ 1 \\ 0}
                     +d\cdot\colvec[r]{1 \\ 0 \\ 1}\suchthat c,d\in\Re}
              \qquad
              B_M=\sequence{\colvec[r]{3 \\ 1 \\ 0},\colvec[r]{1 \\ 0 \\ 1}}
            \end{equation*}
            此处只需解一个含两个方程的线性方程组即可求出 $M^\perp$，
            \begin{equation*}
              \begin{linsys}{3}
                3u  &+  &v  &   &   &=  &0  \\
                 u  &   &   &+  &w  &=  &0
              \end{linsys}
              \grstep{-(1/3)\rho_1+\rho_2}
              \begin{linsys}{3}
                3u  &+  &v         &   &   &=  &0  \\
                    &   &-(1/3)v   &+  &w  &=  &0
              \end{linsys}
            \end{equation*}
            再进行参数化。
            \begin{equation*}
              M^\perp
              =\set{k\cdot\colvec[r]{-1 \\ 3 \\ 1}\suchthat k\in\Re}
            \end{equation*}
          \partsitem 这里 $M$ 是一维的
            \begin{equation*}
              M=\set{c\cdot\colvec[r]{0 \\ -1 \\ 1}\suchthat c\in\Re}
              \qquad
              B_M=\sequence{\colvec[r]{0 \\ -1 \\ 1}}
            \end{equation*}
            因此 $M^\perp$ 是二维的。
            \begin{equation*}
              M^\perp
              =\set{\colvec{u \\ v \\ w}
                     \suchthat 0\cdot u-1\cdot v+1\cdot w=0}
              =\set{j\cdot\colvec[r]{1 \\ 0 \\ 0}
                    +k\cdot\colvec[r]{0 \\ 1 \\ 1}\suchthat j,k\in\Re}
            \end{equation*}
       \end{exparts}
     
\end{ans}
\begin{ans}{三.VI.3.12}
      其中
      \begin{equation*}
        A=
        \begin{mat}
          0 &1 \\
          2 &-1 \\
          0 &1
        \end{mat}
      \end{equation*}
      直接但稍繁琐的计算给出
      % sage: A = matrix(QQ, [[0,1], [2,-1], [0,1]])
      % sage: A*(A.transpose()*A).inverse()*A.transpose()
      % [1/2    0  1/2]
      % [  0    1    0]
      % [1/2    0  1/2]
      \begin{equation*}
        A\bigl(\trans{A}A\bigr)^{-1}\trans{A}=
        \begin{mat}
          1/2 &0 &1/2 \\
           0  &1 &0    \\
          1/2 &0 &1/2
        \end{mat}
      \end{equation*}
      将其作用于该向量即可得到投影。
      \begin{equation*}
        \begin{mat}
          1/2 &0 &1/2 \\
           0  &1 &0    \\
          1/2 &0 &1/2
        \end{mat}
        \colvec{1 \\ 2 \\ 0}
        =
        \colvec{1/2 \\ 2 \\  1/2}
      \end{equation*}
      % sage: v = vector(QQ, [1, 2, 0])
      % sage: v
      % (1, 2, 0)
      % sage: A*(A.transpose()*A).inverse()*A.transpose()*v
      % (1/2, 2, 1/2)
    
\end{ans}
\begin{ans}{三.VI.3.13}
      使用上一答案中的矩阵计算，得到
      \begin{equation*}
        \begin{mat}
          1/2 &0 &-1/2 \\
           0  &1 &0    \\
         -1/2 &0 &1/2
        \end{mat}
        \colvec{1 \\ 2 \\ -1}
        =
        \colvec{1 \\ 2 \\ -1}
      \end{equation*}
      这说明该向量属于子空间 $S$。
    
\end{ans}
\begin{ans}{三.VI.3.14}
      设 $\vec w\in S$。由于这是正交投影，向量 $\vec v-\vec p$ 与 $\vec p-\vec w$ 垂直。由三角不等式，斜边 $\vec v-\vec w$ 至少与 $\vec v-\vec p$ 一样长。
    
\end{ans}
\begin{ans}{三.VI.3.15}
      \begin{exparts}
        \partsitem 将方程参数化，得到 $P$ 的如下基。
          \begin{equation*}
            B_P=\sequence{\colvec[r]{1 \\ 0 \\ 3},
                          \colvec[r]{0 \\ 1 \\ 2} }
          \end{equation*}
        \partsitem 由于 $\Re^3$ 为三维而 $P$ 为二维，补空间 $P^\perp$ 必为一条直线。按照例 \nearbyexample{ex:OrthoCompOne} 的计算，
          \begin{equation*}
            P^\perp=\set{\colvec{x \\ y \\ z}
                         \suchthat
                         \begin{mat}[r]
                           1  &0 &3  \\
                           0  &1 &2
                         \end{mat}
                         \colvec{x \\ y \\ z}
                         =\colvec[r]{0 \\ 0}\;}
          \end{equation*}
          得到 $P^\perp$ 的如下基。
          \begin{equation*}
            B_{P^\perp}=\sequence{\colvec[r]{3 \\ 2 \\ -1} }
          \end{equation*}
        \partsitem $ \displaystyle
             \colvec[r]{1 \\ 1 \\ 2}=
             (5/14)\cdot\colvec[r]{1 \\ 0 \\ 3}+
             (8/14)\cdot\colvec[r]{0 \\ 1 \\ 2}+
             (3/14)\cdot\colvec[r]{3 \\ 2 \\ -1}$
        \partsitem $\proj{\colvec[r]{1 \\ 1 \\ 2}}{P}
          =\colvec[r]{5/14 \\ 8/14 \\ 31/14}$
        \partsitem 投影矩阵
          \begin{multline*}
            \begin{mat}[r]
              1  &0 \\
              0  &1 \\
              3  &2
            \end{mat}
            \bigl(
              \begin{mat}[r]
                1  &0  &3 \\
                0  &1  &2
              \end{mat}
              \begin{mat}[r]
                1  &0 \\
                0  &1 \\
                3  &2
              \end{mat}
            \bigr)^{-1}
            \begin{mat}[r]
              1  &0  &3 \\
              0  &1  &2
            \end{mat}                               \\
            \begin{aligned}
            &=
            \begin{mat}[r]
              1  &0 \\
              0  &1 \\
              3  &2
            \end{mat}
              \begin{mat}[r]
                10  &6  \\
                6   &5
              \end{mat}^{-1}
            \begin{mat}[r]
              1  &0  &3 \\
              0  &1  &2
            \end{mat}                                     \\
            &=
            \frac{1}{14}
            \begin{mat}[r]
              5  &-6 &3  \\
             -6  &10 &2  \\
              3  &2  &13
            \end{mat}
            \end{aligned}
          \end{multline*}
          将它作用于该向量，得到预期结果。
          \begin{equation*}
            \frac{1}{14}
            \begin{mat}[r]
              5  &-6 &3  \\
             -6  &10 &2  \\
              3  &2  &13
            \end{mat}
            \colvec[r]{1 \\ 1 \\ 2}
            =\colvec[r]{5/14 \\ 8/14 \\ 31/14}
          \end{equation*}
      \end{exparts}
    
\end{ans}
\begin{ans}{三.VI.3.16}
       \begin{exparts}
         \partsitem 参数化得到
           \begin{equation*}
             M=\set{c\cdot\colvec[r]{-1 \\ 1}\suchthat c\in\Re}
           \end{equation*}

           第一种方法取张成直线 $M$ 的向量为
           \begin{equation*}
             \vec{s}=\colvec[r]{-1 \\ 1}
           \end{equation*}
           由定义 \nearbydefinition{def:ProjIntoLine} 的公式得到
           \begin{equation*}
             \proj{\colvec[r]{1 \\ -3}}{\spanof{\vec{s}\,}}
              =\frac{\colvec[r]{1 \\ -3}\dotprod\colvec[r]{-1 \\ 1}}{
                     \colvec[r]{-1 \\ 1}\dotprod\colvec[r]{-1 \\ 1}}
                \cdot\colvec[r]{-1 \\ 1}
              =\frac{-4}{2}
                \cdot\colvec[r]{-1 \\ 1}
              =\colvec[r]{2 \\ -2}
           \end{equation*}

           第二种方法中，取
           \begin{equation*}
             B_M=\sequence{\colvec[r]{-1 \\ 1}}
           \end{equation*}
           因而（如例 \nearbyexample{ex:OrthoCompOne} 和~\ref{ex:OrthoCompTwo}，只需找出与基中所有向量垂直的向量）
           \begin{equation*}
             M^\perp
             =\set{\colvec{u \\ v}\suchthat -1\cdot u+1\cdot v=0}
             =\set{k\cdot\colvec[r]{1 \\ 1}\suchthat k\in\Re}
             \qquad
             B_{M^\perp}=\sequence{\colvec[r]{1 \\ 1}}
           \end{equation*}
           关于连接基表示该向量，得到
           \begin{equation*}
             \colvec[r]{1 \\ -3}=-2\cdot\colvec[r]{-1 \\ 1}-1\cdot\colvec[r]{1 \\ 1}
           \end{equation*}
           保留 $M$ 部分即得答案。
           \begin{equation*}
             \proj{\colvec[r]{1 \\ -3}}{M,M^\perp}=\colvec[r]{2 \\ -2}
           \end{equation*}

           第三种方法同样只是简单计算（中间出现一个 $\nbyn 1$ 矩阵，其逆也是 $\nbyn 1$ 矩阵）。
           \begin{multline*}
             A\left(\trans{A}A\right)^{-1}\trans{A}           \\
             \begin{aligned}
             &=
             \begin{mat}[r]
               -1 \\ 1
             \end{mat}
             \left(
               \begin{mat}[r]
                 -1  &1
               \end{mat}
               \begin{mat}[r]
                 -1  \\
                  1
               \end{mat}
             \right)^{-1}
             \begin{mat}[r]
               -1  &1
             \end{mat}
             =
             \begin{mat}[r]
               -1 \\ 1
             \end{mat}
               \begin{mat}[r]
                 2
               \end{mat}^{-1}
             \begin{mat}[r]
               -1  &1
             \end{mat}               \\
             &=
             \begin{mat}[r]
               -1 \\ 1
             \end{mat}
               \begin{mat}[r]
                 1/2
               \end{mat}
             \begin{mat}[r]
               -1  &1
             \end{mat}
             =
             \begin{mat}[r]
               -1 \\ 1
             \end{mat}
             \begin{mat}[r]
               -1/2  &1/2
             \end{mat}
             =
             \begin{mat}[r]
               1/2  &-1/2  \\
              -1/2  &1/2
             \end{mat}
             \end{aligned}
           \end{multline*}
           这当然给出相同答案。
           \begin{equation*}
             \proj{\colvec[r]{1 \\ -3}}{M}
             =
             \begin{mat}[r]
               1/2  &-1/2  \\
              -1/2  &1/2
             \end{mat}
             \colvec[r]{1 \\ -3}
             =\colvec[r]{2 \\ -2}
           \end{equation*}
         \partsitem 参数化得到
           \begin{equation*}
             M=\set{c\cdot\colvec[r]{-1 \\ 0 \\ 1}\suchthat c\in\Re}
           \end{equation*}
           由此，第一种方法的公式给出
           \begin{equation*}
             \frac{\colvec[r]{0 \\ 1 \\ 2}\dotprod\colvec[r]{-1 \\ 0 \\ 1}}{
                   \colvec[r]{-1 \\ 0 \\ 1}\dotprod\colvec[r]{-1 \\ 0 \\ 1}}
               \cdot\colvec[r]{-1 \\ 0 \\ 1}
             =\frac{2}{2}
               \cdot\colvec[r]{-1 \\ 0 \\ 1}
             =\colvec[r]{-1 \\ 0 \\ 1}
           \end{equation*}
           第二种方法先求 $M^\perp$，
           \begin{equation*}
             M^\perp
             =\set{\colvec{u \\ v \\ w}\suchthat -u+w=0}
             =\set{j\cdot\colvec[r]{1 \\ 0 \\ 1}
                   +k\cdot\colvec[r]{0 \\ 1 \\ 0}\suchthat j,k\in\Re}
           \end{equation*}
           再关于 $B_M$ 与 $B_{M^\perp}$ 的连接基表示给定向量，
           \begin{equation*}
             \colvec[r]{0 \\ 1 \\ 2}
              =1\cdot\colvec[r]{-1 \\ 0 \\ 1}
               +1\cdot\colvec[r]{1 \\ 0 \\ 1}
               +1\cdot\colvec[r]{0 \\ 1 \\ 0}
           \end{equation*}
           只保留 $M$ 部分。
           \begin{equation*}
             \proj{\colvec[r]{0 \\ 1 \\ 2}}{M}=1\cdot\colvec[r]{-1 \\ 0 \\ 1}
                       =\colvec[r]{-1 \\ 0 \\ 1}
           \end{equation*}
           最后，第三种方法进行矩阵计算，
           \begin{multline*}
             A\left(\trans{A}A\right)^{-1}\trans{A}                 \\
             \begin{aligned}
             &=
             \begin{mat}[r]
               -1 \\ 0  \\  1
             \end{mat}
             \bigl(
               \begin{mat}[r]
                 -1  &0  &1
               \end{mat}
               \begin{mat}[r]
                 -1  \\
                  0  \\
                  1
               \end{mat}
             \bigr)^{-1}
             \begin{mat}[r]
               -1  &0  &1
             \end{mat}
             =
             \begin{mat}[r]
               -1 \\ 0  \\  1
             \end{mat}
               \begin{mat}[r]
                 2
               \end{mat}^{-1}
             \begin{mat}[r]
               -1  &0  &1
             \end{mat}              \\
             &=
             \begin{mat}[r]
               -1 \\ 0  \\  1
             \end{mat}
               \begin{mat}[r]
                 1/2
               \end{mat}
             \begin{mat}[r]
               -1  &0  &1
             \end{mat}
             =
             \begin{mat}[r]
               -1 \\ 0  \\  1
             \end{mat}
             \begin{mat}[r]
               -1/2  &0  &1/2
             \end{mat}                                 \\
             &=
             \begin{mat}[r]
               1/2  &0  &-1/2  \\
               0    &0  &0     \\
              -1/2  &0  &1/2
             \end{mat}
             \end{aligned}
           \end{multline*}
           再进行矩阵与向量的乘法
           \begin{equation*}
             \proj{\colvec[r]{0 \\ 1 \\ 2}}{M}
             \begin{mat}[r]
               1/2  &0  &-1/2  \\
               0    &0  &0     \\
              -1/2  &0  &1/2
             \end{mat}
             \colvec[r]{0 \\ 1 \\ 2}
             =\colvec[r]{-1 \\ 0 \\ 1}
           \end{equation*}
           即得答案。
       \end{exparts}
     
\end{ans}
\begin{ans}{三.VI.3.17}
      不依赖。向量分解 $\vec v=\vec m+\vec n$（其中 $\vec m\in M,\vec n\in N$）不依赖子空间所选的基，正如“直和”小节所证明的那样。
    
\end{ans}
\begin{ans}{三.VI.3.18}
      向量到子空间的正交投影属于该子空间。平凡子空间只有一个元素 $\zero$，因此任意向量的投影都必须等于 $\zero$。
    
\end{ans}
\begin{ans}{三.VI.3.19}
      对 $\vec v\in M$ 沿 $N$ 投影到 $M$，结果就是 $\vec v$。分解 $\vec v=\vec m+\vec n$ 得到 $\vec m=\vec v,\vec n=\zero$；舍去 $N$ 部分并保留 $M$ 部分，投影即为 $\vec m=\vec v$。
    
\end{ans}
\begin{ans}{三.VI.3.20}
      引理 \nearbylemma{le:OrthoProjWellDefd} 的证明表明，每个 $\vec v\in\Re^n$ 都等于其到各基向量所张成直线的正交投影之和。
      \begin{equation*}
        \vec{v}=\proj{\vec{v}\,}{\spanof{\vec{\kappa}_1}}
                 +\dots+\proj{\vec{v}\,}{\spanof{\vec{\kappa}_n}}
               =\frac{\vec{v}\dotprod\vec{\kappa}_1}{
                      \vec{\kappa}_1\dotprod\vec{\kappa}_1}
                   \cdot\vec{\kappa}_1
               +\dots
               +\frac{\vec{v}\dotprod\vec{\kappa}_n}{
                      \vec{\kappa}_n\dotprod\vec{\kappa}_n}
                   \cdot\vec{\kappa}_n
      \end{equation*}
      由于该基是标准正交基，每个分式的分母 $\vec\kappa_i\dotprod\vec\kappa_i=1$。
    
\end{ans}
\begin{ans}{三.VI.3.21}
       若 $V=M\directsum N$，每个向量都唯一分解为 $\vec v=\vec m+\vec n$。对所有 $\vec v$，$p(\vec v)=\vec m$ 当且仅当 $\vec v-p(\vec v)=\vec n$，这正是所需结论。
     
\end{ans}
\begin{ans}{三.VI.3.22}
      设 $\vec v$ 与每个 $\vec w\in S$ 垂直。则
      $\vec{v}\dotprod(c_1\vec{w}_1+\dots+c_n\vec{w}_n)
       =\vec{v}\dotprod(c_1\vec{w}_1)
          +\dots+\vec{v}\dotprod (c_n\dotprod\vec{w}_n)
       =c_1(\vec{v}\dotprod\vec{w}_1)
          +\dots+c_n(\vec{v}\dotprod\vec{w}_n)
       =c_1\cdot 0+\dots+c_n\cdot 0=0$.
    
\end{ans}
\begin{ans}{三.VI.3.23}
      正确；唯一与自身正交的向量是零向量。
    
\end{ans}
\begin{ans}{三.VI.3.24}
      这直接由引理 \nearbylemma{le:OrthoProjWellDefd} 中“空间是两者直和”的结论得到。
    
\end{ans}
\begin{ans}{三.VI.3.25}
      两者必相等，即使只假设
      它们在所有正交投影下结果相同也如此。
      考虑由集合
      $\set{\vec{v}_1,\vec{v}_2}$.
      由于两者都在 $M$ 中，$\vec{v}_1$ 到 $M$ 的正交投影
      是 $\vec{v}_1$，而 $\vec{v}_2$ 到 $M$ 的正交投影
      $M$ is $\vec{v}_2$.
      因此投影相等只能说明两者相等。
    
\end{ans}
\begin{ans}{三.VI.3.26}
      \begin{exparts}
        \partsitem
          证明两个集合互相包含，即
          $M\subseteq (M^\perp)^\perp$ 且 $(M^\perp)^\perp \subseteq M$。
          首先，
          若 $\vec m\in M$，由正交补运算的定义，
          $\vec m$ 与每个 $\vec v\in M^\perp$ 垂直，
          因而（同样由正交补运算的定义）
          $\vec m\in(M^\perp)^\perp$。
          反过来，取 $\vec v\in(M^\perp)^\perp$。
          引理 \nearbylemma{le:OrthoProjWellDefd} 的证明表明
          $\Re^n=M\directsum M^\perp$，且
          空间有一组正交基
          $\sequence{\vec{\kappa}_1,\dots,\vec{\kappa}_k,
                        \vec{\kappa}_{k+1},\dots,\vec{\kappa}_n}$
          其中前一半
          $\sequence{\vec{\kappa}_1,\dots,\vec{\kappa}_k}$ 是
          $M$ 的一组基，后一半是 $M^\perp$ 的一组基。
          证明还表明，空间中每个向量都等于其到这些基向量所张成直线的正交投影之和。
          \begin{equation*}
             \vec{v}=\proj{\vec{v}\,}{\spanof{\vec{\kappa}_1}}
                          +\dots+\proj{\vec{v}\,}{\spanof{\vec{\kappa}_n}}
          \end{equation*}
          由于 $\vec v\in(M^\perp)^\perp$，它与 $M^\perp$ 中每个向量都垂直，所以
          后一半投影全为零。
          因此 $\vec v=\proj{\vec{v}\,}{\spanof{\vec{\kappa}_1}}
                          +\dots+\proj{\vec{v}\,}{\spanof{\vec{\kappa}_k}}$,
          这是 $M$ 中向量的线性组合，所以 $\vec v\in M$。
          (\textit{Remark}.
           还有一种更简洁的证明：将空间写成
           既可写成 $M\directsum M^\perp$，也可写成 $M^\perp\directsum (M^\perp)^\perp$。
           由于前半部分已证明 $M\subseteq(M^\perp)^\perp$
           且前一句说明两个
           子空间 $M$ 与 $(M^\perp)^\perp$ 维数相等，因此可以得出
           $M=(M^\perp)^\perp$。）
        \partsitem 由于 $M\subseteq N$，与 $N$ 中每个向量垂直的任意 $\vec v$ 也与 $M$ 中每个向量垂直。这句话正是 $N^\perp\subseteq M^\perp$。
        \partsitem 再次通过相互包含证明两个集合相等。首先，若 $\vec v$ 与
           $M+N=\set{\vec{m}+\vec{n}\suchthat \vec{m}\in M,\, \vec{n}\in N}$
           与形如 $\vec m+\zero$ 的每个向量
           （即 $M$ 中每个向量）以及形如 $\zero+\vec n$ 的每个向量（即 $N$ 中每个向量）垂直，因此
           $(M+N)^\perp\subseteq M^\perp\intersection N^\perp$.
           反向包含同样直接：任意 $\vec v\in M^\perp\intersection N^\perp$ 都与形如 $c\vec m+d\vec n$ 的向量垂直，因为
           $\vec{v}\dotprod (c\vec{m}+d\vec{n})
             =c\cdot(\vec{v}\dotprod\vec{m})+d\cdot(\vec{v}\dotprod\vec{n})
             =c\cdot 0+d\cdot 0
             =0$.
      \end{exparts}
    
\end{ans}
\begin{ans}{三.VI.3.27}
      \begin{exparts}
        \partsitem 该映射的表示为
          \begin{equation*}
            \colvec{v_1 \\ v_2 \\ v_3}
            \mapsunder{f}
            1v_1+2v_2+ 3v_3
          \end{equation*}

          \begin{equation*}
            \rep{f}{\stdbasis_3,\stdbasis_1}
            =\begin{mat}[r]
               1  &2  &3
             \end{mat}
          \end{equation*}
          由 $f$ 的定义，
          \begin{equation*}
            \nullspace{f}
            =\set{\colvec{v_1 \\ v_2 \\ v_3}
                  \suchthat 1v_1+2v_2+3v_3=0}
            =\set{\colvec{v_1 \\ v_2 \\ v_3}
                  \suchthat
                   \colvec[r]{1 \\ 2 \\ 3}\dotprod\colvec{v_1 \\ v_2 \\ v_3}=0}
          \end{equation*}
          而第二种描述恰好说明了这一点。
          \begin{equation*}
            \nullspace{f}^\perp
            =\spanof{\set{\colvec[r]{1 \\ 2 \\ 3}}}
          \end{equation*}
        \partsitem 一般地，对任意 $\map f{\Re^n}{\Re}$，存在向量 $\vec h$ 使
           \begin{equation*}
             \colvec{v_1 \\ \vdots \\ v_n}
             \mapsunder{f}
             h_1v_1+\dots+h_nv_n
           \end{equation*}
           且 $\vec h\in\nullspace{f}^\perp$。
           证明方法与上一问相同：关于标准基表示 $f$，将其矩阵表示的唯一一行转置为列向量，所得即为 $\vec h$。
         \partsitem 显然，
           \begin{equation*}
             \rep{f}{\stdbasis_3,\stdbasis_2}
             =
             \begin{mat}[r]
               1  &2  &3  \\
               4  &5  &6
             \end{mat}
           \end{equation*}
           因而零空间为
           \begin{equation*}
             \nullspace{f}
             \set{\colvec{v_1 \\ v_2 \\ v_3}
                  \suchthat
                   \;\begin{mat}[r]
                     1  &2  &3  \\
                     4  &5  &6
                   \end{mat}
                   \colvec{v_1 \\ v_2 \\ v_3}
                   =\colvec[r]{0  \\ 0}\;}
           \end{equation*}
           由此可见
           \begin{equation*}
             \colvec[r]{1 \\ 2 \\ 3},
              \,\colvec[r]{4 \\ 5 \\ 6}\in\nullspace{f}^\perp
           \end{equation*}
           由于 $\nullspace{f}^\perp$ 是 $\Re^n$ 的子空间，这两个向量的张成空间包含于零空间的正交补。为证明这是等式，取
           \begin{equation*}
             M=\spanof{\set{\colvec[r]{1 \\ 2 \\ 3}}}
             \qquad
             N=\spanof{\set{\colvec[r]{4 \\ 5 \\ 6}}}
           \end{equation*}
           并按提示应用习题 \nearbyexercise{exer:AlgOfPerps} 的第三问。
        \partsitem 同上，从特殊情形推广很容易：对任意 $\map f{\Re^n}{\Re^m}$，
           关于标准基表示该映射的矩阵 $H$ 描述如下作用：
           \begin{equation*}
             \colvec{v_1 \\ \vdots \\ v_n}
             \mapsunder{f}
             \colvec{h_{1,1}v_1+h_{1,2}v_2+\dots+h_{1,n}v_n \\
                     \vdots                                 \\
                     h_{m,1}v_1+h_{m,2}v_2+\dots+h_{m,n}v_n}
           \end{equation*}
           零空间的描述表明，对 $H$ 的 $m$ 行分别转置为
           将 $H$ 的 $m$ 行分别转置为
           \begin{equation*}
             \vec{h}_1=\colvec{h_{1,1} \\ h_{1,2} \\ \vdots \\ h_{1,n}},
              \dots
             \vec{h}_m=\colvec{h_{m,1} \\ h_{m,2} \\ \vdots \\ h_{m,n}}
           \end{equation*}
           有 $\nullspace{f}^\perp=\spanof{\set{\vec h_1,\dots,\vec h_m}}$。（\cite{Strang93} 将该空间称为 $H$ 的行空间的转置。）
      \end{exparts}
    
\end{ans}
\begin{ans}{三.VI.3.28}
      \begin{exparts}
        \partsitem 首先，若向量 $\vec v$ 已经在线
          上，则正交投影就是 $\vec v$ 本身。
          验证方法是应用
          向量 $\vec s$ 张成直线的投影公式，
          即 $(\vec v\dotprod\vec s/\vec s\dotprod\vec s)\cdot\vec s$。
          取直线为
          $\set{k\cdot\vec{v}\suchthat k\in\Re}$
          （$\vec v=\zero$ 情形另行处理且很简单），得到
          $(\vec{v}\dotprod\vec{v}/\vec{v}\dotprod\vec{v})\cdot\vec{v}$,
          化简为 $\vec v$，即得结论。

          这就回答了问题，因为投影一次后
          结果 $\proj{\vec v}{\ell}$ 已在线上。
          上一段说明再次投影到同一直线
          不会产生任何变化。
        \partsitem 论证与上一问类似。
          当 $V=M\directsum N$ 时，$\vec v=\vec m+\vec n$ 的投影
          为 $\proj{\vec v}{M,N}=\vec m$。
          再次投影得到 $\proj{\vec m}{M,N}=\vec m$，因为 $M$ 中元素分解为 $M$ 中元素与 $N$ 中元素之和时就是 $\vec m=\vec m+\zero$。因此沿 $N$ 投影到 $M$ 两次与一次效果相同。
        \partsitem 如前面各问所示，该条件说明 $t$ 保持值域空间中的向量不变，因此可取
          $\vec{\beta}_1$, \ldots, $\vec{\beta}_r$
          作为值域空间的基向量，也就是取 $t$ 的值域空间为 $M$（所以 $\dim(M)=r$）。对于补空间，记 $t$ 的零空间为 $N$，并证明 $V=M\directsum N$。

          为此，只需证明它们的交为平凡空间 $M\intersection N=\set{\zero}$，且它们的和为整个空间 $M+N=V$。
          首先，若向量 $\vec m$ 属于值域空间，则存在 $\vec v\in V$ 使 $t(\vec v)=\vec m$；由 $t$ 的条件，
          $t(\vec{m})=(\composed{t}{t})\,(\vec{v})=t(\vec{v})=\vec{m}$, while
          若该向量同时属于零空间，则 $t(\vec m)=\zero$，因此值域空间与零空间的交是平凡空间。
          其次，要将任意 $\vec v$ 写成值域空间向量与零空间向量之和，
          条件 $t(\vec v)=t(t(\vec v))$ 可改写为 $t(\vec v-t(\vec v))=\zero$，因此取 $\vec v=t(\vec v)+(\vec v-t(\vec v))$。

          最后取 $V$ 的一组基 $B=\sequence{\vec\beta_1,\ldots,\vec\beta_n}$，其中前 $r$ 个向量构成值域空间 $M$ 的一组基，剩余向量构成零空间 $N$ 的一组基。
        \partsitem 每个投影（按本题定义）都是沿其零空间投影到其值域空间。
        \partsitem 这也可由第三问立即得到。
      \end{exparts}
    
\end{ans}
\begin{ans}{三.VI.3.29}
      对任意矩阵 $M$，有 $\trans{(M^{-1})}=(\trans M)^{-1}$；对任意两个矩阵 $M,N$，有 $\trans{MN}=\trans N\trans M$（当然，前提是逆和乘积有定义）。应用这两条性质即可证明该矩阵等于其转置。
      \begin{multline*}
         \trans{\bigl(A(\trans{A}A)^{-1}\trans{A}\bigr)}
         =(\trans{\trans{A}})
           (\trans{\bigl((\trans{A}A)^{-1}\bigr)})
           (\trans{A})                               \\
         =(\trans{\trans{A}})
           ( \bigl(\trans{(\trans{A}A)}\bigr)^{-1} )
           (\trans{A})
         =A
           (\trans{A}\trans{\trans{A}})^{-1}
           \trans{A}
         =A
           (\trans{A}A)^{-1}
           \trans{A}
      \end{multline*}
    
\end{ans}
\protect \topic {最佳拟合直线}
\begin{ans}{1}
      与前面的第一个例子一样，我们要寻找
      最佳 $m$ 来“求解”该方程组。
      \begin{equation*}
        \begin{linsys}{5}
          8m  &=  &4  \\
          16m &=  &9  \\
          24m &=  &13 \\
          32m &=  &17 \\
          40m &=  &20
        \end{linsys}
      \end{equation*}
      投影到线性子空间得到
      \begin{equation*}
        \frac{\colvec[r]{4 \\ 9 \\ 13 \\ 17 \\20}
          \dotprod
          \colvec[r]{8  \\ 16 \\ 24 \\ 32 \\ 40}}{%
          \colvec[r]{8  \\ 16 \\ 24 \\ 32 \\ 40}
          \dotprod
          \colvec[r]{8  \\ 16 \\ 24 \\ 32 \\ 40}}
        \cdot
        \colvec[r]{8  \\ 16 \\ 24 \\ 32 \\ 40}
        =
        \frac{1832}{3520}
        \cdot
        \colvec[r]{8  \\ 16 \\ 24 \\ 32 \\ 40}
      \end{equation*}
      因此最佳拟合直线斜率约为 $0.52$。
      \begin{center}  \small
        \includegraphics{map/mp/ch3.85}
      \end{center}
    
\end{ans}
\begin{ans}{2}
     使用以下输入
     \begin{equation*}
       A =
       \begin{mat}
           1 & 1852.71 \\
           1 & 1858.88 \\
          \vdots  &\vdots      \\
           1 & 1985.54 \\
           1 & 1993.71
        \end{mat}
        \qquad
        b = \colvec{ 292.0 \\
                     285.0 \\
                    \vdots  \\
                     226.32 \\
                     224.39}
     \end{equation*}
     (the dates have been rounded to months, e.g., for a September record,
     小数 $.71\approx (8.5/12)$），Maple 给出
     截距为 $b=994.8276974$，斜率为
     $m=-0.3871993827$.
     \begin{center}  \small
        \includegraphics{map/mp/ch3.86}
     \end{center}
     
\end{ans}
\begin{ans}{3}
   使用以下输入（年份以 $1900$ 年为零点）
   \begin{equation*}
     A :=
     \begin{mat}
        1 &   .38   \\
        1 &   .54   \\
       \vdots \vdots \\
        1 & 92.71 \\
        1 & 95.54
     \end{mat}
     \qquad
     b = \colvec{ 249.0 \\
                  246.2 \\
                 \vdots \\
                  208.86 \\
                  207.37}
   \end{equation*}
     (the dates have been rounded to months, e.g., for a September record,
     使用小数 $.71\approx (8.5/12)$），
     Maple 给出截距 $b=243.1590327$，斜率
   $m=-0.401647703$.
   本专题正文中男子一英里赛的斜率与此
   与此非常接近。
   \begin{center}  \small
     \includegraphics{map/mp/ch3.87}
   \end{center}
   
\end{ans}
\begin{ans}{4}
    使用以下输入（年份以 $1900$ 年为零点）
    \begin{equation*}
      A =
      \begin{mat}
         1 & 21.46 \\
         1 & 32.63 \\
        \vdots  &\vdots     \\
         1 & 89.54  \\
         1 & 96.63
      \end{mat}
      \qquad
      b =
      \colvec{ 373.2 \\
               327.5 \\
              \vdots  \\
               255.61 \\
               252.56}
    \end{equation*}
     (the dates have been rounded to months, e.g., for a September record,
     使用小数 $.71\approx (8.5/12)$），
     MAPLE 给出截距 $b=378.7114894$，斜率
    $m=-1.445753225$.
    \begin{center}  \small
      \includegraphics{map/mp/ch3.88}
    \end{center}
   
\end{ans}
\begin{ans}{5}
    男子和女子一英里赛直线方程为
    （女子一英里赛方程的纵截距项已根据上面的答案调整，使其在年份 $0$ 处为零，因为男子一英里赛方程采用了同样处理。）
    \begin{align*}
        y &=994.8276974-0.3871993827x  \\
        y &=3125.6426-1.445753225x
    \end{align*}
    显然两直线相交。
    使用计算机程序最容易完成
    运算：MuPAD 给出 $x=2012.949004$ 和 $y=215.4150856$
    ($215$~seconds is $3$~minutes and $35$~seconds).
    \textit{Remark.}
    当然，这种投影存在很大疑问 \Dash  例如，
    女子方程受早期较慢纪录的影响
    \Dash  但过程仍然很有趣。
     \begin{center}  \small
       \includegraphics{map/mp/ch3.89}
     \end{center}
  
\end{ans}
\begin{ans}{6}
      \textit{Sage} gives the line of best fit as
      $\text{toll}=-0.05\cdot\text{dist}+5.63$.
\begin{lstlisting}
sage: dist = [2, 7, 8, 16, 27, 47, 67, 82, 102, 120]
sage: toll = [6, 6, 6, 6.5, 2.5, 1, 1, 1, 1, 1]
sage: var('a,b,t')
(a, b, t)
sage: model(t) = a*t+b
sage: data = zip(dist,toll)
sage: fit = find_fit(data, model, solution_dict=True)
sage: model.subs(fit)
t |--> -0.0508568169130319*t + 5.630955848442933
sage: p = plot(model.subs(fit), (t,0,120))+points(data,size=25,color='red')
sage: p.save('bridges.pdf')
\end{lstlisting}
      但图表显示该方程几乎没有预测价值。
      \begin{center}
        \includegraphics[width=.5\textwidth]{map/pix/bridges.pdf}
      \end{center}
      显然更好的模型是：
      （只有一个中间例外）市区的 crossing
      费用大致相同，
      州北部 crossing 的费用也彼此相同。
    
\end{ans}
\begin{ans}{7}
      \begin{exparts}
        \partsitem MAPLE 或 MuPAD 等计算机代数系统给出
          截距为 $b=4259/1398\approx 3.239628$
          斜率为 $m=-71/2796\approx -0.025393419$
          将 $x=31$ 代入方程，预测
          O 形密封圈失效次数为 $y=2.45$（四舍五入到两位）。
          代入 $y=4$ 并求解，得到温度
          $x=-29.94^\circ$F.
        \partsitem On the basis of this information
          \begin{equation*}
            A =
            \begin{mat}
               1 & 53 \\
               1 & 75 \\
%              1 & 57 \\
%              1 & 58 \\
%              1 & 63 \\
%              1 & 70 \\
%              1 & 70 \\
%              1 & 66 \\
%              1 & 67 \\
%              1 & 67 \\
%              1 & 67 \\
%              1 & 68 \\
%              1 & 69 \\
%              1 & 70 \\
%              1 & 70 \\
%              1 & 72 \\
%              1 & 73 \\
%              1 & 75 \\
%              1 & 76 \\
%              1 & 76 \\
%              1 & 78 \\
%              1 & 79 \\
               \vdots \\
               1 & 80 \\
               1 & 81
            \end{mat}
            \qquad
            b =
            \colvec{ 3 \\
                     2 \\
%                    1 \\
%                    1 \\
%                    1 \\
%                    1 \\
%                    1 \\
%                    0 \\
%                    0 \\
%                    0 \\
%                    0 \\
%                    0 \\
%                    0 \\
%                    0 \\
%                    0 \\
%                    0 \\
%                    0 \\
%                    0 \\
%                    0 \\
%                    0 \\
%                    0 \\
%                    0 \\
                     \vdots \\
                     0 \\
                     0}
          \end{equation*}
          MAPLE 给出截距 $b=187/40=4.675$，
          slope $m=-73/1200\approx -0.060833$.
          此时将 $x=31$ 代入方程预测
          $y=2.79$ O-ring failures (rounded to two places).
          代入 $y=4$ 次失效得到温度
          $x=11^\circ$F。
     \begin{center}  \small
       \includegraphics{map/mp/ch3.90}
      \end{center}
      \end{exparts}
    
\end{ans}
\begin{ans}{8}
        \begin{exparts}
          \partsitem The plot is nonlinear.
            \begin{center}  \small
              \includegraphics{map/mp/ch3.91}
            \end{center}
          \partsitem Here is the plot.
           \begin{center}  \small
             \includegraphics{map/mp/ch3.92}
           \end{center}
           行星 $4$ 与 $5$ 之间可能有向上的突变。
          \partsitem This plot seems even more linear.
           \begin{center}  \small
             \includegraphics{map/mp/ch3.93}
           \end{center}
          \partsitem
            使用以下输入
            \begin{equation*}
              A =
              \begin{mat}[r]
                1 & 1 \\
                1 & 2 \\
                1 & 3 \\
                1 & 4 \\
                1 & 6 \\
                1 & 7 \\
                1 & 8
              \end{mat}
              \qquad
              b = \colvec[r]{-0.40893539 \\
                          -0.1426675 \\
                           0 \\
                           0.18184359 \\
                           0.71600334 \\
                           0.97954837 \\
                           1.2833012}
            \end{equation*}
            MuPAD 给出截距 $b= -0.6780677466$，
            斜率为 $m=0.2372763818$。
          \partsitem Plugging $x=9$ into the equation
            $y= -0.6780677466+0.2372763818x$ from the prior item gives
            距离对数为 $1.4574197$，因此预测
            距离为 $28.669472$。
            实际距离约为 $30.003$。
          \partsitem Plugging $x=10$ into the same equation
            给出距离对数为 $1.6946961$，因此预测
            距离为 $49.510362$。
            实际距离约为 $39.503$。
        \end{exparts}
      
\end{ans}
\begin{ans}{9}
    For any $\vec{w}\in W$, the vectors $\vec{v}-\vec{p}$ and
    $\vec{p}-\vec{w}$ are orthogonal.
    因此三角不等式可应用于以这些向量为边的三角形，
    两边，且以 $\vec{v}-\vec{w}$ 为斜边。
    Therefore $\vec{v}-\vec{w}$ is at least as long as $\vec{v}-\vec{p}$.
  
\end{ans}
\protect \topic {线性映射的几何}
\begin{ans}{1}
      下面的高斯消元
      \begin{multline*}
        \grstep[-\rho_1+\rho_3]{-3\rho_1+\rho_2}
        \begin{mat}[r]
          1  &2  &1  \\
          0  &0  &-3 \\
          0  &0  &1
        \end{mat}
        \grstep{(1/3)\rho_2+\rho_3}
        \begin{mat}[r]
          1  &2  &1  \\
          0  &0  &-3 \\
          0  &0  &0
        \end{mat}                                     \\
        \grstep{(-1/3)\rho_2}
        \begin{mat}[r]
          1  &2  &1  \\
          0  &0  &1 \\
          0  &0  &0
        \end{mat}
        \grstep{-\rho_2+\rho_1}
        \begin{mat}[r]
          1  &2  &0  \\
          0  &0  &1 \\
          0  &0  &0
        \end{mat}
      \end{multline*}
      得到该矩阵的简化阶梯形。
      接着，对列作两次操作：先将第一列乘以 $-2$
      加到第二列，再交换第二、三列
      得到如下部分恒等矩阵。
      \begin{equation*}
        B=\begin{mat}[r]
          1  &0  &0  \\
          0  &1  &0  \\
          0  &0  &0
        \end{mat}
      \end{equation*}
      以上内容用矩阵表示为：
      \begin{equation*}
        P=
        \begin{mat}[r]
          1  &-1    &0  \\
          0  &1     &0  \\
          0  &0     &1
        \end{mat}
        \begin{mat}[r]
          1  &0    &0  \\
          0  &-1/3 &0  \\
          0  &0    &1
        \end{mat}
        \begin{mat}[r]
          1  &0    &0  \\
          0  &1    &0  \\
          0  &1/3  &1
        \end{mat}
        \begin{mat}[r]
          1  &0  &0  \\
          0  &1  &0  \\
         -1  &0  &1
        \end{mat}
        \begin{mat}[r]
          1  &0  &0  \\
         -3  &1  &0  \\
          0  &0  &1
        \end{mat}
      \end{equation*}
      和
      \begin{equation*}
        Q=
        \begin{mat}[r]
          1  &-2    &0  \\
          0  &1     &0  \\
          0  &0     &1
        \end{mat}
        \begin{mat}[r]
          0  &1     &0  \\
          1  &0     &0  \\
          0  &0     &1
        \end{mat}
      \end{equation*}
      因此给定矩阵可分解为 $PBQ$。
    
\end{ans}
\begin{ans}{2}
      首先用矩阵 $H$ 表示该映射，
      进行行变换，并在需要时进行列变换
      将其化为部分恒等矩阵。
      然后将其写成分解 $H=PBQ$。
      代入旋转矩阵的一般形式
          \begin{equation*}
            \rep{r_\theta}{\stdbasis_2,\stdbasis_2}
            \begin{mat}
              \cos\theta  &-\sin\theta  \\
              \sin\theta  &\cos\theta
            \end{mat}
          \end{equation*}
          得到如下表示。
          \begin{equation*}
            \rep{r_{2\pi/3}}{\stdbasis_2,\stdbasis_2}
            \begin{mat}[r]
              -1/2        &-\sqrt{3}/2  \\
              \sqrt{3}/2  &-1/2
            \end{mat}
          \end{equation*}
          Gauss's Method is routine.
          \begin{equation*}
            \grstep{\sqrt{3}\rho_1+\rho_2}
            \begin{mat}[r]
              -1/2        &-\sqrt{3}/2  \\
               0          &-2
            \end{mat}
            \grstep[(-1/2)\rho_2]{-2\rho_1}
            \begin{mat}[r]
               1          &\sqrt{3}    \\
               0          &1
            \end{mat}
            \grstep{-\sqrt{3}\rho_2+\rho_1}
            \begin{mat}[r]
               1          &0   \\
               0          &1
            \end{mat}
          \end{equation*}
          这对应于如下矩阵方程。
          \begin{equation*}
            \begin{mat}[r]
              1  &-\sqrt{3}  \\
              0  &1
            \end{mat}
            \begin{mat}[r]
              -2  &0    \\
               0  &-1/2
            \end{mat}
            \begin{mat}[r]
               1         &0  \\
               \sqrt{3}  &1
            \end{mat}
            \begin{mat}[r]
              -1/2        &-\sqrt{3}/2  \\
              \sqrt{3}/2  &-1/2
            \end{mat}
            =I
          \end{equation*}
          对方程取逆并解出 $H$，得到以下分解。
          \begin{equation*}
            \begin{mat}[r]
              -1/2        &-\sqrt{3}/2  \\
              \sqrt{3}/2  &-1/2
            \end{mat}
            =
            \begin{mat}[r]
                1         &0  \\
               -\sqrt{3}  &1
            \end{mat}
            \begin{mat}[r]
              -1/2  &0    \\
               0    &-2
            \end{mat}
            \begin{mat}[r]
              1  &\sqrt{3}  \\
              0  &1
            \end{mat}
            I
          \end{equation*}
    
\end{ans}
\begin{ans}{3}
      \begin{exparts}
        \partsitem 回忆，逆时针旋转
          $\theta$ 弧度相对于标准基表示为
          可表示为：
          \begin{equation*}
            \rep{t_{\pi/4}}{\stdbasis_2,\stdbasis_2}
            =\begin{mat}
              \cos\theta  &-\sin\theta  \\
              \sin\theta  &\cos\theta
             \end{mat}
          \end{equation*}
          顺时针角是逆时针角的相反数，
          one.
          \begin{equation*}
            \rep{t_{-pi/4}}{\stdbasis_2,\stdbasis_2}
            =\begin{mat}
              \cos(-\pi/4)  &-\sin(-\pi/4)  \\
              \sin(-\pi/4)  &\cos(-\pi/4)
            \end{mat}
            =\begin{mat}[r]
              \sqrt{2}/2  &\sqrt{2}/2  \\
              -\sqrt{2}/2 &\sqrt{2}/2
            \end{mat}
          \end{equation*}
        \partsitem
          下面的高斯-约旦消元
          \begin{equation*}
            \grstep{\rho_1+\rho_2}
            \begin{mat}[r]
              \sqrt{2}/2  &\sqrt{2}/2  \\
              0           &\sqrt{2}
            \end{mat}
            \grstep[(1/\sqrt{2})\rho_2]{(2/\sqrt{2})\rho_1}
            \begin{mat}[r]
              1  &1  \\
              0  &1
            \end{mat}
            \grstep{-\rho_2+\rho_1}
            \begin{mat}[r]
              1  &0  \\
              0  &1
            \end{mat}
          \end{equation*}
          得到单位矩阵。
          因此无需交换列
          即可得到部分恒等矩阵。
        \partsitem In matrix multiplication the reduction is
          \begin{equation*}
            \begin{mat}[r]
              1  &-1 \\
              0  &1
            \end{mat}
            \begin{mat}[r]
              2/\sqrt{2}  &0         \\
              0           &1/\sqrt{2}
            \end{mat}
            \begin{mat}[r]
              1  &0 \\
              1  &1
            \end{mat}
            H
            =I
          \end{equation*}
          （注意，高斯操作的复合从右向左进行）。
        \partsitem  Taking inverses
          \begin{equation*}
            H
            =
            \underbrace{
              \begin{mat}[r]
                1  &0 \\
                -1  &1
              \end{mat}
              \begin{mat}[r]
                \sqrt{2}/2  &0         \\
                0           &\sqrt{2}
              \end{mat}
              \begin{mat}[r]
                1  &1 \\
                0  &1
              \end{mat}
             }_P
            I
          \end{equation*}
          得到 $H$ 所需的分解。
          部分恒等矩阵就是 $I$。
          % , 和 $Q$ is trivial, that is, it is also an identity
          % matrix.
        \partsitem 从右向左读取复合（忽略
          单位矩阵视为平凡矩阵），可知 $H$ 的作用等同于
          先执行以下剪切
          \begin{center}
            \includegraphics{map/mp/ch3.95}
         \end{center}
         再进行伸缩，将所有第一分量乘以
         $\sqrt{2}/2$（这是收缩，因为 $\sqrt{2}/2\approx0.707$
         小于 $1$）。
         并将所有第二分量乘以 $\sqrt{2}$，
         最后再进行一次剪切。
          \begin{center}
            \includegraphics{map/mp/ch3.96}
         \end{center}
         例如，取与
         $x$ 轴夹角为 $\pi/6$ 的单位向量，依次应用 $H$ 的各个分量。
          \begin{center}
            \includegraphics{map/mp/ch3.97}
         \end{center}
         容易验证所得向量长度为 1，
         与 $x$ 轴夹角为 $-\pi/12$，确实形成
         即顺时针旋转 $\pi/4$ 弧度，
         因为 $(\pi/6)-(\pi/4)=-\pi/12$。
      \end{exparts}
    
\end{ans}
\begin{ans}{4}
      相对于
      标准基 $\stdbasis_1,\stdbasis_1$。
      这得到一个 $\nbyn{1}$ 矩阵。
      唯一元素就是标量~$k$。
    
\end{ans}
\begin{ans}{5}
      \begin{exparts}
        \partsitem 直线是 $\Re^n$ 中形如
          $\set{\vec{v}=\vec{u}+t\cdot\vec{w}\suchthat t\in\Re}$.
          该直线上点的像为
          $h(\vec{v})=h(\vec{u}+t\cdot\vec{w})=h(\vec{u})+t\cdot h(\vec{w})$,
          当 $t$ 遍历实数时，这些向量的集合是
          一条直线（若 $h(\vec{w})=\zero$ 则为退化直线）。
        \partsitem 这是前面论证的直接推广。
        \partsitem 若点~$B$ 位于点~$A$ 与~$C$ 之间，则
          从 $A$ 到 $C$ 的线段包含 $B$。
          也就是说，存在 $t\in (0\,..\,1)$ 使得
          $\vec{b}=\vec{a}+t\cdot (\vec{c}-\vec{a})$（其中 $B$ 是
          （其中 $B$ 是 $\vec b$ 的端点，其他类似）。
          如第一小题的论证，线性性说明
          $h(\vec{b})=h(\vec{a})+t\cdot h(\vec{c}-\vec{a})$.
      \end{exparts}
    
\end{ans}
\begin{ans}{6}
      二者互为倒数。
      例如，对固定的 $x\in\Re$，
      若 $f^\prime (x)=k$（其中 $k\neq 0$），则
      $(f^{-1})^\prime (x)=1/k$.
      \begin{center}
        \includegraphics{map/mp/ch3.98}
     \end{center}
    
\end{ans}
\begin{ans}{7}
      对向量分量数作归纳即可证明。
      vector.
      In the $n=1$ base case the only permutation is the trivial one,
      且映射
      \begin{equation*}
        \colvec{x_1}
        \mapsto
        \colvec{x_1}
      \end{equation*}
      可表示为交换的复合，即零次交换。
      归纳步骤中，假设由
      少于 $n$ 个数的任意排列诱导的映射
      $n$ 个数的排列可仅用交换表示；现在考虑映射
      由
      $n$ 个数的排列 $p$ 诱导。
      \begin{equation*}
        \colvec{x_1 \\ x_2 \\ \vdots \\ x_n}
        \mapsto
        \colvec{x_{p(1)} \\ x_{p(2)} \\ \vdots \\ x_{p(n)}}
      \end{equation*}
      考虑满足 $p(i)=n$ 的编号 $i$。
      映射
      \begin{equation*}
        \colvec{x_1      \\ x_2      \\ \vdots \\ x_i      \\ \vdots \\ x_n}
        \mapsunder{\hat{p}}
        \colvec{x_{p(1)} \\ x_{p(2)} \\ \vdots \\ x_{p(n)} \\ \vdots  \\ x_{n}}
      \end{equation*}
      再交换第 $i$ 个与第 $n$ 个分量后，
      得到映射~$p$。
      根据归纳假设，$\hat{p}$ 可表示为交换的复合。
      的交换复合。
    
\end{ans}
\protect \topic {幻方}
\begin{ans}{1}
      \begin{exparts}
        \partsitem $M$ 的元素之和等于三行元素和之和。

        \partsitem 涉及中心元素的约束给出方程组：
          \begin{equation*}
            \begin{linsys}{3}
              m_{2,1}  &+  &m_{2,2}  &+  &m_{2,3}  &=  &s  \\
              m_{1,2}  &+  &m_{2,2}  &+  &m_{3,2}  &=  &s  \\
              m_{1,1}  &+  &m_{2,2}  &+  &m_{3,3}  &=  &s  \\
              m_{1,3}  &+  &m_{2,2}  &+  &m_{3,1}  &=  &s
            \end{linsys}
          \end{equation*}
          将这四个方程相加时，每个元素只计算一次，
          但中心元素计算了四次。
          因此左端为 $3s+3m_{2,2}$，右端为 $4s$。
          所以 $3m_{2,2}=s$。
        \partsitem
          第二行之和为 $s$，所以 $m_{2,1}+m_{2,2}+m_{2,3}=3m_{2,2}$,
          从而 $(1/2)\cdot(m_{2,1}+m_{2,3})=m_{2,2}$.
          对列和对角线同理。
        \partsitem
          根据上一小题，$m_{2,1}$ 和 $m_{2,3}$ 要么都等于 $m_{2,2}$，要么一个大于而另一个小于 $m_{2,2}$。
          Thus $m_{2,2}$ is the median of the set
          $\set{m_{2,1},m_{2,2},m_{2,3}}$.
          对第二列作同样推理可得，
          Thus $m_{2,2}$ is the median of the set
          $\set{m_{1,2},m_{2,1},m_{2,2},m_{2,3},m_{3,2}}$.
          对两条对角线作推广可知它是所有元素的中位数。
          所有元素的中位数。
      \end{exparts}
    
\end{ans}
\begin{ans}{2}
      对任意 $k$，有如下消元：
      \begin{multline*}
        \begin{amat}{4}
          1  &1  &0  &0  &s  \\
          0  &0  &1  &1  &s  \\
          1  &0  &1  &0  &s  \\
          0  &1  &0  &1  &s  \\
          1  &0  &0  &1  &s  \\
          0  &1  &1  &0  &s
        \end{amat}
        \grstep[-\rho_1+\rho_5]{-\rho_1+\rho_3}
        \begin{amat}{4}
          1  &1  &0  &0  &s  \\
          0  &0  &1  &1  &s  \\
          0  &-1 &1  &0  &0  \\
          0  &1  &0  &1  &s  \\
          0  &-1 &0  &1  &0  \\
          0  &1  &1  &0  &s
        \end{amat}                                            \\
        \grstep{-\rho_2\leftrightarrow\rho_6}
        \begin{amat}{4}
          1  &1  &0  &0  &s  \\
          0  &1  &1  &0  &s  \\
          0  &-1 &1  &0  &0  \\
          0  &1  &0  &1  &s  \\
          0  &-1 &0  &1  &0  \\
          0  &0  &1  &1  &s
        \end{amat}
        \grstep[-\rho_2+\rho_4 \\ \rho_2+\rho_5]{-\rho_2+\rho_3}
        \begin{amat}{4}
          1  &1  &0  &0  &s  \\
          0  &1  &1  &0  &s  \\
          0  &0  &2  &0  &s  \\
          0  &1  &-1 &1  &0  \\
          0  &0  &1  &1  &s  \\
          0  &0  &1  &1  &s
        \end{amat}
      \end{multline*}
      唯一解为 $a=b=c=d=s/2$。
    
\end{ans}
\begin{ans}{3}
       根据上一小题，唯一元素是 $Z_{\nbyn{2}}$。
    
\end{ans}
\begin{ans}{4}
      \begin{exparts}
        \partsitem
          对 $M,N\in\matspace_{\nbyn{n}}$，有
          $\trace(cM+dN)=(cm_{1,1}+dn_{1,1})+\cdots+(cm_{n,n}+dn_{n,n})
          =(cm_{1,1}+\cdots+cm_{n,n})+(dn_{1,1}+\cdots+dn_{n,n})
          =c\cdot\trace(M)+d\cdot\trace(N)$其中所有数均为实数，因此迹保持线性组合。
         对 $\trace^*$ 的论证类似。
       \partsitem
         它保持线性组合；所有数均为实数，
          $\theta(cM+dN)=(\trace(cM+dN),\trace^*(cM+dN))
          =(c\cdot\trace(M)+d\cdot\trace(N), c\cdot\trace^*(M)+d\cdot\trace^*(N))
          =c\cdot\theta(M)+d\cdot\theta(N)$.
       \partsitem
         若 $\map{h_1,\ldots,h_n}{V}{W}$ 都是线性映射，则
         $\map{g}{V}{W^n}$ given by
         $g(\vec{v})=(h_1(\vec{v}), \ldots, h_n(\vec{v}))$.
         证明同上一小题。
      \end{exparts}
    
\end{ans}
\begin{ans}{5}
       \begin{exparts*}
         \partsitem
           两个半幻方之和仍是半幻方，半幻方的标量倍数也仍是半幻方。
         \partsitem
           与上一小题一样，两个
           幻和为零的半幻方的线性组合仍是此类矩阵。
       \end{exparts*}
     
\end{ans}
\protect \topic {马尔可夫链}
\begin{ans}{1}
      \begin{exparts}
        \partsitem With this file \texttt{coin.m}
\begin{computercode}
# Octave function for Markov coin game.  p is chance of going down.
function w = coin(p,v)
  q = 1-p;
  A=[1,p,0,0,0,0;
     0,0,p,0,0,0;
     0,q,0,p,0,0;
     0,0,q,0,p,0;
     0,0,0,q,0,0;
     0,0,0,0,q,1];
  w = A * v;
endfunction
\end{computercode}
         本次 Octave 会话产生如下输出。
\begin{computercode}
octave:1> v0=[0;0;0;1;0;0]
v0 =
  0
  0
  0
  1
  0
  0
octave:2> p=.5
p = 0.50000
octave:3> v1=coin(p,v0)
v1 =
  0.00000
  0.00000
  0.50000
  0.00000
  0.50000
  0.00000
octave:4> v2=coin(p,v1)
v2 =
  0.00000
  0.25000
  0.00000
  0.50000
  0.00000
  0.25000
\end{computercode}
后续步骤过多，此处不一一列出。
\begin{computercode}
octave:26> v24=coin(p,v23)
v24 =
  0.39600
  0.00276
  0.00000
  0.00447
  0.00000
  0.59676
\end{computercode}
        \partsitem 利用以下公式
          \begin{alignat*}{2}
             p_{1}(n+1) &= 0.5\cdot p_{2}(n)
             &\qquad
             p_{2}(n+1) &= 0.5\cdot p_{1}(n)+0.5\cdot p_{3}(n)
                                                              \\
             p_{3}(n+1) &= 0.5\cdot p_{2}(n)+0.5\cdot p_{4}(n)
             &\qquad
             p_{5}(n+1) &= 0.5\cdot p_{4}(n)
          \end{alignat*}
          以及以下初始条件
          \begin{equation*}
            \colvec{p_{0}(0) \\ p_{1}(0) \\ p_{2}(0) \\ p_{3}(0) \\
                        p_{4}(0) \\ p_{5}(0)}
            =
            \colvec{0 \\ 0 \\ 0 \\ 1 \\
                        0 \\ 0}
          \end{equation*}
          我们将用归纳法证明：当 $n$ 为奇数时，
          $p_{1}(n)=p_{3}(n)=0$；当 $n$ 为偶数时，$p_{2}(n)=p_{4}(n)=0$。
          首先，由初始条件可知 $n=0$ 基础情形成立。
          的情形成立。
          归纳步骤中，假设结论在
          $n=0,1,\ldots,k$ 情形成立，考察 $n=k+1$。
          若 $k+1$ 为奇数，则以下两个
          \begin{align*}
             p_{1}(k+1) &= 0.5\cdot p_{2}(k)=0.5\cdot 0=0
                                                          \\
             p_{3}(k+1) &=0.5\cdot p_{2}(k)+0.5\cdot p_{4}(k)
                        =0.5\cdot 0+0.5\cdot 0
                        =0
          \end{align*}
          根据归纳假设可得
          $p_{2}(k)=p_{4}(k)=0$，因为 $k$ 为偶数。
          当 $k+1$ 为偶数时同理。
        \partsitem 例如取 $n=100$。
          以下 Octave 会话
\begin{computercode}
octave:1> B=[1,.5,0,0,0,0;
>             0,0,.5,0,0,0;
>             0,.5,0,.5,0,0;
>             0,0,.5,0,.5,0;
>             0,0,0,.5,0,0;
>             0,0,0,0,.5,1];
octave:2> B100=B**100
B100 =
  1.00000  0.80000  0.60000  0.40000  0.20000  0.00000
  0.00000  0.00000  0.00000  0.00000  0.00000  0.00000
  0.00000  0.00000  0.00000  0.00000  0.00000  0.00000
  0.00000  0.00000  0.00000  0.00000  0.00000  0.00000
  0.00000  0.00000  0.00000  0.00000  0.00000  0.00000
  0.00000  0.20000  0.40000  0.60000  0.80000  1.00000
octave:3> B100*[0;1;0;0;0;0]
octave:4> B100*[0;1;0;0;0;0]
octave:5> B100*[0;0;0;1;0;0]
octave:6> B100*[0;1;0;0;0;0]
\end{computercode}
        得到以下输出。
        \begin{equation*}
          \begin{array}{r|*{4}{c}}
            \multicolumn{1}{r}{\textit{starting with:}}
            &\text{\$1}  &\text{\$2}  &\text{\$3}  &\text{\$4} \\
             \hline
               \begin{array}{c}
                 s_{0}(100) \\
                 s_{1}(100) \\
                 s_{2}(100) \\
                 s_{3}(100) \\
                 s_{4}(100) \\
                 s_{5}(100)
               \end{array}
               &\begin{aligncolondecimal}{5}
                 0.80000 \\
                 0.00000 \\
                 0.00000 \\
                 0.00000 \\
                 0.00000 \\
                 0.20000
               \end{aligncolondecimal}
               &\begin{aligncolondecimal}{5}
                    0.60000 \\
                    0.00000 \\
                    0.00000 \\
                    0.00000 \\
                    0.00000 \\
                    0.40000
               \end{aligncolondecimal}
               &\begin{aligncolondecimal}{5}
                      0.40000 \\
                      0.00000 \\
                      0.00000 \\
                      0.00000 \\
                      0.00000 \\
                      0.60000
               \end{aligncolondecimal}
               &\begin{aligncolondecimal}{5}
                   0.20000 \\
                   0.00000 \\
                   0.00000 \\
                   0.00000 \\
                   0.00000 \\
                   0.80000
               \end{aligncolondecimal}
          \end{array}
        \end{equation*}
      \end{exparts}
    
\end{ans}
\begin{ans}{2}
      \begin{exparts}
        \partsitem From these equations
          \begin{equation*}
            \begin{linsys}{6}
              s_{1}(n)/6 &+ &0s_{2}(n) &+ &0s_{3}(n)
                  &+ &0s_{4}(n) &+ &0s_{5}(n)  &+ &0s_{6}(n)  &= &s_{1}(n+1) \\
              s_{1}(n)/6 &+ &2s_{2}(n)/6 &+ &0s_{3}(n)
                  &+ &0s_{4}(n) &+ &0s_{5}(n)  &+ &0s_{6}(n)  &= &s_{2}(n+1) \\
              s_{1}(n)/6 &+ &s_{2}(n)/6 &+ &3s_{3}(n)/6
                  &+ &0s_{4}(n) &+ &0s_{5}(n)  &+ &0s_{6}(n)  &= &s_{3}(n+1) \\
              s_{1}(n)/6 &+ &s_{2}(n)/6 &+ &s_{3}(n)/6
                  &+ &4s_{4}(n)/6 &+ &0s_{5}(n)  &+ &0s_{6}(n)
                  &=  &s_{4}(n+1) \\
              s_{1}(n)/6 &+ &s_{2}(n)/6 &+ &s_{3}(n)/6
                  &+ &s_{4}(n)/6 &+ &5s_{5}(n)/6  &+ &0s_{6}(n)
                  &=  &s_{5}(n+1) \\
              s_{1}(n)/6 &+ &s_{2}(n)/6 &+ &s_{3}(n)/6
                  &+ &s_{4}(n)/6 &+ &s_{5}(n)/6  &+ &6s_{6}(n)/6
                  &=  &s_{6}(n+1)
            \end{linsys}
          \end{equation*}
          得到如下转移矩阵。
          \begin{equation*}
            \begin{mat}
              1/6  &0   &0   &0   &0   &0  \\
              1/6  &2/6 &0   &0   &0   &0  \\
              1/6  &1/6 &3/6 &0   &0   &0  \\
              1/6  &1/6 &1/6 &4/6 &0   &0  \\
              1/6  &1/6 &1/6 &1/6 &5/6 &0  \\
              1/6  &1/6 &1/6 &1/6 &1/6 &6/6
            \end{mat}
          \end{equation*}
       \partsitem 下面是 Octave 会话，
         省略输出，并将结果汇总在末尾表格中。
\begin{computercode}
octave:1>   F=[1/6,  0,   0,   0,   0,   0;
>      1/6,  2/6, 0,   0,   0,   0;
>      1/6,  1/6, 3/6, 0,   0,   0;
>      1/6,  1/6, 1/6, 4/6, 0,   0;
>      1/6,  1/6, 1/6, 1/6, 5/6, 0;
>      1/6,  1/6, 1/6, 1/6, 1/6, 6/6];
octave:2> v0=[1;0;0;0;0;0]
octave:3> v1=F*v0
octave:4> v2=F*v1
octave:5> v3=F*v2
octave:6> v4=F*v3
octave:7> v5=F*v4
\end{computercode}
        结果如下。
        \begin{equation*}
          \begin{array}{c|*{5}{c}}
            \multicolumn{1}{c}{\ }
            &1  &2  &3  &4 &5  \\
            \hline
               \begin{aligncolondecimal}{0}
                   1 \\
                   0 \\
                   0 \\
                   0 \\
                   0 \\
                   0
               \end{aligncolondecimal}
               &\begin{aligncolondecimal}{5}
                    0.16667 \\
                    0.16667 \\
                    0.16667 \\
                    0.16667 \\
                    0.16667 \\
                    0.16667
               \end{aligncolondecimal}
               &\begin{aligncolondecimal}{6}
                   0.027778 \\
                   0.083333 \\
                   0.138889 \\
                   0.194444 \\
                   0.250000 \\
                   0.305556
               \end{aligncolondecimal}
               &\begin{aligncolondecimal}{7}
                    0.0046296 \\
                    0.0324074 \\
                    0.0879630 \\
                    0.1712963 \\
                    0.2824074 \\
                    0.4212963
               \end{aligncolondecimal}
               &\begin{aligncolondecimal}{8}
                    0.00077160 \\
                    0.01157407 \\
                    0.05015432 \\
                    0.13503086 \\
                    0.28472222 \\
                    0.51774691
               \end{aligncolondecimal}
               &\begin{aligncolondecimal}{8}
                   0.00012860 \\
                   0.00398663 \\
                   0.02713477 \\
                   0.10043724 \\
                   0.27019033 \\
                   0.59812243
               \end{aligncolondecimal}
          \end{array}
        \end{equation*}
      \end{exparts}
    
\end{ans}
\begin{ans}{3}
     \begin{exparts}
      \partsitem 直觉上合理的是：公司的当前位置应当强烈影响它下一时刻的位置（例如是否留下），而更早阶段的位置影响应当很小。也就是说，公司可能因为当前所在位置而移动或留下，却不太可能因为过去所在的位置而移动或留下。
      \partsitem 下面是略作编辑的 Octave 会话，输出汇总在末尾的表格中。
\begin{lstlisting}
octave:1> M=[.787,0,0,.111,.102;
>            0,.966,.034,0,0;
>            0,.063,.937,0,0;
>            0,0,.074,.612,.314;
>            .021,.009,.005,.010,.954]
M =
  0.78700  0.00000  0.00000  0.11100  0.10200
  0.00000  0.96600  0.03400  0.00000  0.00000
  0.00000  0.06300  0.93700  0.00000  0.00000
  0.00000  0.00000  0.07400  0.61200  0.31400
  0.02100  0.00900  0.00500  0.01000  0.95400
octave:2> v0=[.025;.025;.025;.025;.900]
octave:3> v1=M*v0
octave:4> v2=M*v1
octave:5> v3=M*v2
octave:6> v4=M*v3
\end{lstlisting}
        下表总结了结果。
        \begin{center}
           \begin{tabular}{c|cccc}
             $\vec{p}_0$   &$\vec{p}_1$    &$\vec{p}_2$
                    &$\vec{p}_3$   &$\vec{p}_4$    \\ \hline
             $\left(\begin{aligncolondecimal}{6}
                    0.025000 \\
                    0.025000 \\
                    0.025000 \\
                    0.025000 \\
                    0.900000
              \end{aligncolondecimal}\right)$
             &$\left(\begin{aligncolondecimal}{6}
                  0.114250 \\
                  0.025000 \\
                  0.025000 \\
                  0.299750 \\
                  0.859725
              \end{aligncolondecimal}\right)$
             &$\left(\begin{aligncolondecimal}{6}
                  0.210879 \\
                  0.025000 \\
                  0.025000 \\
                  0.455251 \\
                  0.825924
              \end{aligncolondecimal}\right)$
             &$\left(\begin{aligncolondecimal}{6}
                  0.300739 \\
                  0.025000 \\
                  0.025000 \\
                  0.539804 \\
                  0.797263
              \end{aligncolondecimal}\right)$
             &$\left(\begin{aligncolondecimal}{6}
                  0.377920 \\
                  0.025000 \\
                  0.025000 \\
                  0.582550 \\
                  0.772652
              \end{aligncolondecimal}\right)$
           \end{tabular}
       \end{center}
       \partsitem 以下是上一小题 Octave 会话的继续。
\begin{lstlisting}
octave:7> p0=[.0000;.6522;.3478;.0000;.0000]
octave:8> p1=M*p0
octave:9> p2=M*p1
octave:10> p3=M*p2
octave:11> p4=M*p3
\end{lstlisting}
        以上概括了输出。
        \begin{center}
           \begin{tabular}{c|cccc}
             $\vec{p}_0$   &$\vec{p}_1$    &$\vec{p}_2$
                    &$\vec{p}_3$   &$\vec{p}_4$    \\ \hline
             $\left(\begin{aligncolondecimal}{5}
                  0.00000 \\
                  0.65220 \\
                  0.34780 \\
                  0.00000 \\
                  0.00000
              \end{aligncolondecimal}\right)$
             &$\left(\begin{aligncolondecimal}{5}
                  0.00000 \\
                  0.64185 \\
                  0.36698 \\
                  0.02574 \\
                  0.00761
              \end{aligncolondecimal}\right)$
             &$\left(\begin{aligncolondecimal}{7}
                  0.0036329 \\
                  0.6325047 \\
                  0.3842942 \\
                  0.0452966 \\
                  0.0151277
              \end{aligncolondecimal}\right)$
             &$\left(\begin{aligncolondecimal}{7}
                  0.0094301 \\
                  0.6240656 \\
                  0.3999315 \\
                  0.0609094 \\
                  0.0225751
              \end{aligncolondecimal}\right)$
             &$\left(\begin{aligncolondecimal}{6}
                  0.016485 \\
                  0.616445 \\
                  0.414052 \\
                  0.073960 \\
                  0.029960
              \end{aligncolondecimal}\right)$
           \end{tabular}
       \end{center}
        \partsitem 以下是同一 Octave 会话的后续。
\begin{lstlisting}
octave:12> M50=M**50
M50 =
  0.03992  0.33666  0.20318  0.02198  0.37332
  0.00000  0.65162  0.34838  0.00000  0.00000
  0.00000  0.64553  0.35447  0.00000  0.00000
  0.03384  0.38235  0.22511  0.01864  0.31652
  0.04003  0.33316  0.20029  0.02204  0.37437
octave:13> p50=M50*p0
p50 =
  0.29024
  0.54615
  0.54430
  0.32766
  0.28695
octave:14> p51=M*p50
p51 =
  0.29406
  0.54609
  0.54442
  0.33091
  0.29076
\end{lstlisting}
        这已经接近稳态。
      \end{exparts}
    
\end{ans}
\begin{ans}{4}
      \begin{exparts}
        \partsitem 相关方程组如下。
          \begin{equation*}
            \begin{linsys}{5}
              (1-2p)\cdot s_{U}(n)  &+ &p\cdot t_{A}(n)  &+  &p\cdot t_{B}(n)
                &  &  &  &
                &=  &s_{U}(n+1)            \\
              p\cdot s_{U}(n)  &+  &(1-2p)\cdot t_{A}(n)  &  &  &  &  &  &
                &=  &t_{A}(n+1)            \\
              p\cdot s_{U}(n)  &  &  &+  &(1-2p)\cdot t_{B}(n)  &  &  &  &
                &=  &t_{B}(n+1)            \\
                &  &p\cdot t_{A}(n)  &  &  &+  &s_{A}(n)  &  &
                &=  &s_{A}(n+1)            \\
                &  &  &   &p\cdot t_{B}(n)  &  &  &+  &s_{B}(n)
                &=  &s_{B}(n+1)
            \end{linsys}
          \end{equation*}
          因此得到：
          \begin{equation*}
            \begin{pmatrix}
              1-2p &p     &p    &0  &0  \\
              p    &1-2p  &0    &0  &0  \\
              p    &0     &1-2p &0  &0  \\
              0    &p     &0    &1  &0  \\
              0    &0     &p    &0  &1
            \end{pmatrix}
            \colvec{s_{U}(n) \\ t_{A}(n) \\ t_{B}(n) \\ s_{A}(n) \\ s_{B}(n)}
            =
            \colvec{s_{U}(n+1) \\ t_{A}(n+1) \\ t_{B}(n+1)
                                     \\ s_{A}(n+1) \\ s_{B}(n+1)}
          \end{equation*}
        \partsitem 以下是 Octave 代码（省略输出）。
\begin{lstlisting}
octave:1> T=[.5,.25,.25,0,0;
>            .25,.5,0,0,0;
>            .25,0,.5,0,0;
>            0,.25,0,1,0;
>            0,0,.25,0,1]
T =
  0.50000  0.25000  0.25000  0.00000  0.00000
  0.25000  0.50000  0.00000  0.00000  0.00000
  0.25000  0.00000  0.50000  0.00000  0.00000
  0.00000  0.25000  0.00000  1.00000  0.00000
  0.00000  0.00000  0.25000  0.00000  1.00000
octave:2> p0=[1;0;0;0;0]
octave:3> p1=T*p0
octave:4> p2=T*p1
octave:5> p3=T*p2
octave:6> p4=T*p3
octave:7> p5=T*p4
\end{lstlisting}
        输出如下。
        最终到达 $s_A$ 的概率约为 $0.23$。
        \begin{equation*}
          \begin{array}{cc|ccccc}
            &\vec{p}_0 &\vec{p}_1 &\vec{p}_2 &\vec{p}_3
                &\vec{p}_4 &\vec{p}_5 \\
            \hline
            \begin{array}{c}
              s_U \\
              t_A \\
              t_B \\
              s_A \\
              s_B
            \end{array}
            &\begin{array}{c}
               1 \\
               0 \\
               0 \\
               0 \\
               0
            \end{array}
            &\begin{array}{c}
              0.50000 \\
              0.25000 \\
              0.25000 \\
              0.00000 \\
              0.00000
            \end{array}
            &\begin{array}{c}
              0.375000 \\
              0.250000 \\
              0.250000 \\
              0.062500 \\
              0.062500
            \end{array}
            &\begin{array}{c}
              0.31250 \\
              0.21875 \\
              0.21875 \\
              0.12500 \\
              0.12500
            \end{array}
            &\begin{array}{c}
              0.26562 \\
              0.18750 \\
              0.18750 \\
              0.17969 \\
              0.17969
            \end{array}
            &\begin{array}{c}
              0.22656 \\
              0.16016 \\
              0.16016 \\
              0.22656 \\
              0.22656
            \end{array}
          \end{array}
        \end{equation*}
       \partsitem 将以下文件保存为 \texttt{learn.m} 后
\begin{lstlisting}
# Octave script file for learning model.
function w = learn(p)
   T = [1-2*p,p,    p,    0, 0;
        p,    1-2*p,0,    0, 0;
        p,    0,    1-2*p,0, 0;
        0,    p,    0,    1, 0;
        0,    0,    p,    0, 1];
  T5 = T**5;
  p5 = T5*[1;0;0;0;0];
  w = p5(4);
endfunction
\end{lstlisting}
        执行命令 \texttt{octave:1> learn(.20)} 得到
        \texttt{ans = 0.17664}.
       \partsitem This Octave session
\begin{lstlisting}
octave:1> x=(.01:.01:.50)';
octave:2> y=(.01:.01:.50)';
octave:3> for i=.01:.01:.50
>           y(100*i)=learn(i);
>         endfor
octave:4> z=[x, y];
octave:5> gplot z
\end{lstlisting}
        得到下图。
        不存在阈值 \Dash 没有某个概率使
        曲线突然上升。
        \begin{center}
          \includegraphics[width=.55\textwidth]{map/pix/learn5.eps}
        \end{center}
      \end{exparts}
    
\end{ans}
\begin{ans}{5}
       \begin{exparts}
         \partsitem 由以下方程组
           \begin{equation*}
             \begin{linsys}{2}
               0.90\cdot p_{T}(n)  &+  &0.01\cdot p_{C}(n) &= &p_{T}(n+1)   \\
               0.10\cdot p_{T}(n)  &+  &0.99\cdot p_{C}(n) &= &p_{C}(n+1)
             \end{linsys}
           \end{equation*}
           得到该矩阵。
           \begin{equation*}
             \begin{pmatrix}
               0.90  &0.01  \\
               0.10  &0.99
             \end{pmatrix}
             \colvec{p_{T}(n) \\ p_{C}(n)}
             =
             \colvec{p_{T}(n+1) \\ p_{C}(n+1)}
           \end{equation*}
         \partsitem Octave 的结果如下。
           \begin{center}
             \begin{tabular}{c|ccccc}
               $n=0$ &1  &2  &3  &4  &5  \\
               \hline
               $\begin{array}{c}  0.30000 \\ 0.70000 \end{array}$
               &$\begin{array}{c} 0.27700 \\ 0.72300 \end{array}$
               &$\begin{array}{c}  0.25653 \\ 0.74347 \end{array}$
               &$\begin{array}{c}  0.23831 \\ 0.76169 \end{array}$
               &$\begin{array}{c}  0.22210 \\ 0.77790 \end{array}$
               &$\begin{array}{c}  0.20767 \\ 0.79233 \end{array}$
             \end{tabular}                                     \\[1ex]
             \begin{tabular}{|ccccc}
               6  &7  &8  &9 &10 \\
               \hline
               $\begin{array}{c}  0.19482 \\ 0.80518 \end{array}$
               &$\begin{array}{c}  0.18339 \\ 0.81661 \end{array}$
               &$\begin{array}{c}  0.17322 \\ 0.82678 \end{array}$
               &$\begin{array}{c}  0.16417 \\ 0.83583 \end{array}$
               &$\begin{array}{c}  0.15611 \\ 0.84389 \end{array}$
             \end{tabular}
           \end{center}
         \partsitem 这是 $s_T=0.2$ 时的结果。
           \begin{center}
             \begin{tabular}{c|ccccc}
               $n=0$ &1  &2  &3  &4  &5  \\
               \hline
               $\begin{array}{c}  0.20000 \\ 0.80000 \end{array}$
               &$\begin{array}{c}   0.18800 \\ 0.81200 \end{array}$
               &$\begin{array}{c}   0.17732 \\ 0.82268 \end{array}$
               &$\begin{array}{c}   0.16781 \\ 0.83219 \end{array}$
               &$\begin{array}{c}   0.15936 \\ 0.84064 \end{array}$
               &$\begin{array}{c}   0.15183 \\ 0.84817 \end{array}$
             \end{tabular}                                         \\[1ex]
             \begin{tabular}{|ccccc}
               6  &7  &8  &9 &10 \\
               \hline
               $\begin{array}{c}  0.14513 \\ 0.85487  \end{array}$
               &$\begin{array}{c}  0.13916 \\ 0.86084  \end{array}$
               &$\begin{array}{c}  0.13385 \\ 0.86615  \end{array}$
               &$\begin{array}{c}  0.12913 \\ 0.87087  \end{array}$
               &$\begin{array}{c}  0.12493 \\ 0.87507  \end{array}$
             \end{tabular}
           \end{center}
          \partsitem Although the probability vectors start $0.1$ apart,
            二者相差仅 $0.032$。
            因此二者相似。
       \end{exparts}
     
\end{ans}
\begin{ans}{6}
     这些分别是 $p=.55$ 和 $p=0.60$ 时的向量。
     \begin{center}\footnotesize
       \begin{tabular}{@{}r@{}cccccccc@{}}
         &$n=0$  &$n=1$  &$n=2$  &$n=3$  &$n=4$  &$n=5$  &$n=6$  &$n=7$  \\
        \hline
        \begin{tabular}{@{}c@{}}
           0-0 \\
           1-0 \\
           0-1 \\
           2-0 \\
           1-1 \\
           0-2 \\
           3-0 \\
           2-1 \\
           1-2 \\
           0-3 \\
           4-0 \\
           3-1 \\
           2-2 \\
           1-3 \\
           0-4 \\
           4-1 \\
           3-2 \\
           2-3 \\
           1-4 \\
           4-2 \\
           3-3 \\
           2-4 \\
           4-3 \\
           3-4
         \end{tabular}
          &$\begin{aligncolondecimal}{0}
             1 \\
             0 \\
             0 \\
             0 \\
             0 \\
             0 \\
             0 \\
             0 \\
             0 \\
             0 \\
             0 \\
             0 \\
             0 \\
             0 \\
             0 \\
             0 \\
             0 \\
             0 \\
             0 \\
             0 \\
             0 \\
             0 \\
             0 \\
             0
           \end{aligncolondecimal}$
         &$\begin{aligncolondecimal}{5}
           0 \\
           0.55000 \\
           0.45000 \\
           0 \\
           0 \\
           0 \\
           0 \\
           0 \\
           0 \\
           0 \\
           0 \\
           0 \\
           0 \\
           0 \\
           0 \\
           0 \\
           0 \\
           0 \\
           0 \\
           0 \\
           0 \\
           0 \\
           0 \\
           0
         \end{aligncolondecimal}$
         &$\begin{aligncolondecimal}{5}
           0 \\
           0 \\
           0 \\
           0.30250 \\
           0.49500 \\
           0.20250 \\
           0 \\
           0 \\
           0 \\
           0 \\
           0 \\
           0 \\
           0 \\
           0 \\
           0 \\
           0 \\
           0 \\
           0 \\
           0 \\
           0 \\
           0 \\
           0 \\
           0 \\
           0
         \end{aligncolondecimal}$
         &$\begin{aligncolondecimal}{5}
           0 \\
           0 \\
           0 \\
           0 \\
           0 \\
           0 \\
           0.16638 \\
           0.40837 \\
           0.33412 \\
           0.09112 \\
           0 \\
           0 \\
           0 \\
           0 \\
           0 \\
           0 \\
           0 \\
           0 \\
           0 \\
           0 \\
           0 \\
           0 \\
           0 \\
           0
         \end{aligncolondecimal}$
         &$\begin{aligncolondecimal}{5}
           0 \\
           0 \\
           0 \\
           0 \\
           0 \\
           0 \\
           0 \\
           0 \\
           0 \\
           0 \\
           0.09151 \\
           0.29948 \\
           0.36754 \\
           0.20047 \\
           0.04101 \\
           0 \\
           0 \\
           0 \\
           0 \\
           0 \\
           0 \\
           0 \\
           0 \\
           0
         \end{aligncolondecimal}$
         &$\begin{aligncolondecimal}{5}
          0 \\
          0 \\
          0 \\
          0 \\
          0 \\
          0 \\
          0 \\
          0 \\
          0 \\
          0 \\
          0.09151 \\
          0 \\
          0 \\
          0 \\
          0.04101 \\
          0.16471 \\
          0.33691 \\
          0.27565 \\
          0.09021 \\
          0 \\
          0 \\
          0 \\
          0 \\
          0
         \end{aligncolondecimal}$
         &$\begin{aligncolondecimal}{5}
          0 \\
          0 \\
          0 \\
          0 \\
          0 \\
          0 \\
          0 \\
          0 \\
          0 \\
          0 \\
          0.09151 \\
          0 \\
          0 \\
          0 \\
          0.04101 \\
          0.16471 \\
          0 \\
          0 \\
          0.09021 \\
          0.18530 \\
          0.30322 \\
          0.12404 \\
          0 \\
          0
         \end{aligncolondecimal}$
         &$\begin{aligncolondecimal}{5}
          0 \\
          0 \\
          0 \\
          0 \\
          0 \\
          0 \\
          0 \\
          0 \\
          0 \\
          0 \\
          0.09151 \\
          0 \\
          0 \\
          0 \\
          0.04101 \\
          0.16471 \\
          0 \\
          0 \\
          0.09021 \\
          0.18530 \\
          0 \\
          0.12404 \\
          0.16677 \\
          0.13645
         \end{aligncolondecimal}$
       \end{tabular}                      \\
       % \end{center}
       % and these are the $p=.60$ vectors.
       % \begin{center}\small
       \begin{tabular}{@{}r@{}cccccccc@{}}
         &$n=0$  &$n=1$  &$n=2$  &$n=3$  &$n=4$  &$n=5$  &$n=6$  &$n=7$  \\
        \hline
        \begin{tabular}{@{}c@{}}
           0-0 \\
           1-0 \\
           0-1 \\
           2-0 \\
           1-1 \\
           0-2 \\
           3-0 \\
           2-1 \\
           1-2 \\
           0-3 \\
           4-0 \\
           3-1 \\
           2-2 \\
           1-3 \\
           0-4 \\
           4-1 \\
           3-2 \\
           2-3 \\
           1-4 \\
           4-2 \\
           3-3 \\
           2-4 \\
           4-3 \\
           3-4
         \end{tabular}
          &$\begin{aligncolondecimal}{0}
             1 \\
             0 \\
             0 \\
             0 \\
             0 \\
             0 \\
             0 \\
             0 \\
             0 \\
             0 \\
             0 \\
             0 \\
             0 \\
             0 \\
             0 \\
             0 \\
             0 \\
             0 \\
             0 \\
             0 \\
             0 \\
             0 \\
             0 \\
             0
           \end{aligncolondecimal}$
         &$\begin{aligncolondecimal}{5}
          0 \\
          0.60000 \\
          0.40000 \\
          0 \\
          0 \\
          0 \\
          0 \\
          0 \\
          0 \\
          0 \\
          0 \\
          0 \\
          0 \\
          0 \\
          0 \\
          0 \\
          0 \\
          0 \\
          0 \\
          0 \\
          0 \\
          0 \\
          0 \\
          0
         \end{aligncolondecimal}$
         &$\begin{aligncolondecimal}{5}
            0 \\
            0 \\
            0 \\
            0.36000 \\
            0.48000 \\
            0.16000 \\
            0 \\
            0 \\
            0 \\
            0 \\
            0 \\
            0 \\
            0 \\
            0 \\
            0 \\
            0 \\
            0 \\
            0 \\
            0 \\
            0 \\
            0 \\
            0 \\
            0 \\
            0
         \end{aligncolondecimal}$
         &$\begin{aligncolondecimal}{5}
            0 \\
            0 \\
            0 \\
            0 \\
            0 \\
            0 \\
            0.21600 \\
            0.43200 \\
            0.28800 \\
            0.06400 \\
            0 \\
            0 \\
            0 \\
            0 \\
            0 \\
            0 \\
            0 \\
            0 \\
            0 \\
            0 \\
            0 \\
            0 \\
            0 \\
            0
         \end{aligncolondecimal}$
         &$\begin{aligncolondecimal}{5}
            0 \\
            0 \\
            0 \\
            0 \\
            0 \\
            0 \\
            0 \\
            0 \\
            0 \\
            0 \\
            0.12960 \\
            0.34560 \\
            0.34560 \\
            0.15360 \\
            0.02560 \\
            0 \\
            0 \\
            0 \\
            0 \\
            0 \\
            0 \\
            0 \\
            0 \\
            0
         \end{aligncolondecimal}$
         &$\begin{aligncolondecimal}{5}
             0 \\
             0 \\
             0 \\
             0 \\
             0 \\
             0 \\
             0 \\
             0 \\
             0 \\
             0 \\
             0.12960 \\
             0 \\
             0 \\
             0 \\
             0.02560 \\
             0.20736 \\
             0.34560 \\
             0.23040 \\
             0.06144 \\
             0 \\
             0 \\
             0 \\
             0 \\
             0
         \end{aligncolondecimal}$
         &$\begin{aligncolondecimal}{5}
            0 \\
            0 \\
            0 \\
            0 \\
            0 \\
            0 \\
            0 \\
            0 \\
            0 \\
            0 \\
            0.12960 \\
            0 \\
            0 \\
            0 \\
            0.02560 \\
            0.20736 \\
            0 \\
            0 \\
            0.06144 \\
            0.20736 \\
            0.27648 \\
            0.09216 \\
            0 \\
            0
         \end{aligncolondecimal}$
         &$\begin{aligncolondecimal}{5}
            0 \\
            0 \\
            0 \\
            0 \\
            0 \\
            0 \\
            0 \\
            0 \\
            0 \\
            0 \\
            0.12960 \\
            0 \\
            0 \\
            0 \\
            0.02560 \\
            0.20736 \\
            0 \\
            0 \\
            0.06144 \\
            0.20736 \\
            0 \\
            0.09216 \\
            0.16589 \\
            0.11059
         \end{aligncolondecimal}$
       \end{tabular}
     \end{center}
     \begin{exparts}
      \partsitem We can adapt the script from the end of this Topic.
\begin{lstlisting}
# Octave script file to compute chance of World Series outcomes.
function w = markov(p,v)
  q = 1-p;
  A=[0,0,0,0,0,0, 0,0,0,0,0,0, 0,0,0,0,0,0, 0,0,0,0,0,0;  # 0-0
     p,0,0,0,0,0, 0,0,0,0,0,0, 0,0,0,0,0,0, 0,0,0,0,0,0;  # 1-0
     q,0,0,0,0,0, 0,0,0,0,0,0, 0,0,0,0,0,0, 0,0,0,0,0,0;  # 0-1_
     0,p,0,0,0,0, 0,0,0,0,0,0, 0,0,0,0,0,0, 0,0,0,0,0,0;  # 2-0
     0,q,p,0,0,0, 0,0,0,0,0,0, 0,0,0,0,0,0, 0,0,0,0,0,0;  # 1-1
     0,0,q,0,0,0, 0,0,0,0,0,0, 0,0,0,0,0,0, 0,0,0,0,0,0;  # 0-2__
     0,0,0,p,0,0, 0,0,0,0,0,0, 0,0,0,0,0,0, 0,0,0,0,0,0;  # 3-0
     0,0,0,q,p,0, 0,0,0,0,0,0, 0,0,0,0,0,0, 0,0,0,0,0,0;  # 2-1
     0,0,0,0,q,p, 0,0,0,0,0,0, 0,0,0,0,0,0, 0,0,0,0,0,0;  # 1-2_
     0,0,0,0,0,q, 0,0,0,0,0,0, 0,0,0,0,0,0, 0,0,0,0,0,0;  # 0-3
     0,0,0,0,0,0, p,0,0,0,1,0, 0,0,0,0,0,0, 0,0,0,0,0,0;  # 4-0
     0,0,0,0,0,0, q,p,0,0,0,0, 0,0,0,0,0,0, 0,0,0,0,0,0;  # 3-1__
     0,0,0,0,0,0, 0,q,p,0,0,0, 0,0,0,0,0,0, 0,0,0,0,0,0;  # 2-2
     0,0,0,0,0,0, 0,0,q,p,0,0, 0,0,0,0,0,0, 0,0,0,0,0,0;  # 1-3
     0,0,0,0,0,0, 0,0,0,q,0,0, 0,0,1,0,0,0, 0,0,0,0,0,0;  # 0-4_
     0,0,0,0,0,0, 0,0,0,0,0,p, 0,0,0,1,0,0, 0,0,0,0,0,0;  # 4-1
     0,0,0,0,0,0, 0,0,0,0,0,q, p,0,0,0,0,0, 0,0,0,0,0,0;  # 3-2
     0,0,0,0,0,0, 0,0,0,0,0,0, q,p,0,0,0,0, 0,0,0,0,0,0;  # 2-3__
     0,0,0,0,0,0, 0,0,0,0,0,0, 0,q,0,0,0,0, 1,0,0,0,0,0;  # 1-4
     0,0,0,0,0,0, 0,0,0,0,0,0, 0,0,0,0,p,0, 0,1,0,0,0,0;  # 4-2
     0,0,0,0,0,0, 0,0,0,0,0,0, 0,0,0,0,q,p, 0,0,0,0,0,0;  # 3-3_
     0,0,0,0,0,0, 0,0,0,0,0,0, 0,0,0,0,0,q, 0,0,0,1,0,0;  # 2-4
     0,0,0,0,0,0, 0,0,0,0,0,0, 0,0,0,0,0,0, 0,0,p,0,1,0;  # 4-3
     0,0,0,0,0,0, 0,0,0,0,0,0, 0,0,0,0,0,0, 0,0,q,0,0,1]; # 3-4
  v7 = (A**7) * v;
  w = v7(11)+v7(16)+v7(20)+v7(23)
endfunction
\end{lstlisting}
       当美国联盟球队
       $p=0.55$ probability of winning each game then their probability
       赢得系列赛的概率为 $0.60829$。
       当其单场获胜概率为 $p=0.6$ 时，
       其赢得系列赛的概率为
       $0.71021$.
      \partsitem From this Octave session and its graph
\begin{lstlisting}
octave:1> v0=[1;0;0;0;0;0;0;0;0;0;0;0;0;0;0;0;0;0;0;0;0;0;0;0];
octave:2> x=(.01:.01:.99)';
octave:3> y=(.01:.01:.99)';
octave:4> for i=.01:.01:.99
>          y(100*i)=markov(i,v0);
>         endfor
octave:5> z=[x, y];
octave:6> gplot z
\end{lstlisting}
       目测判断，当 $p>0.7$ 时球队几乎可以确定
       的系列赛。
       \begin{center}
         \includegraphics[width=0.4\textwidth]{map/pix/ws.eps}
       \end{center}
     \end{exparts}
    
\end{ans}
\begin{ans}{7}
      \begin{exparts}
        \partsitem 它们必须满足该条件，因为状态转移（包括回到
          原状态）的总概率
          为 $100\%$。
        \partsitem 见第三小题答案。
        \partsitem 我们只证明
          $\nbyn{2}$ 情形；更高阶情形只是记号
          更复杂。
          该乘积
          \begin{equation*}
            \begin{pmatrix}
              a_{1,1}  &a_{1,2}  \\
              a_{2,1}  &a_{2,2}
            \end{pmatrix}
            \begin{pmatrix}
              b_{1,1}  &b_{1,2}  \\
              b_{2,1}  &b_{2,2}
            \end{pmatrix}
            =\begin{pmatrix}
              a_{1,1}b_{1,1}+a_{1,2}b_{2,1}  &a_{1,1}b_{1,2}+a_{1,2}b_{2,2} \\
              a_{2,1}b_{1,1}+a_{2,2}b_{2,1}  &a_{2,1}b_{1,2}+a_{2,2}b_{2,2} \\
            \end{pmatrix}
           \end{equation*}
            其两列之和分别为
            \begin{multline*}
              (a_{1,1}b_{1,1}+a_{1,2}b_{2,1})
              +(a_{2,1}b_{1,1}+a_{2,2}b_{2,1})
              =(a_{1,1}+a_{2,1})\cdot b_{1,1}
                +(a_{1,2}+a_{2,2})\cdot b_{2,1}            \\
              =1\cdot b_{1,1}+1\cdot b_{2,1}
              =1
            \end{multline*}
            以及
            \begin{multline*}
              (a_{1,1}b_{1,2}+a_{1,2}b_{2,2})
              +(a_{2,1}b_{1,2}+a_{2,2}b_{2,2})
              =(a_{1,1}+a_{2,1})\cdot b_{1,2}
                +(a_{1,2}+a_{2,2})\cdot b_{2,2}              \\
              =1\cdot b_{1,2}+1\cdot b_{2,2}
              =1
            \end{multline*}
            如所需。
      \end{exparts}
    
\end{ans}
\protect \topic {标准正交矩阵}
\begin{ans}{1}
      \begin{exparts}
        \partsitem Yes.
        \partsitem No, the columns do not have length one.
        \partsitem Yes.
      \end{exparts}
    
\end{ans}
\begin{ans}{2}
      其中一些是非线性的，
      因为包含非平凡平移。
      \begin{exparts}
        \partsitem
          $\colvec{x \\ y}
           \mapsto
           \begin{mat}
             x\cdot\cos(\pi/6)-y\cdot\sin(\pi/6) \\
             x\cdot\sin(\pi/6)+y\cdot\cos(\pi/6)
           \end{mat}
           +\colvec[r]{0 \\ 1}
           =\begin{mat}
             x\cdot(\sqrt{3}/2)-y\cdot(1/2)+0 \\
             x\cdot(1/2)+y\cdot\cos(\sqrt{3}/2)+1
           \end{mat}$
        \partsitem The line $y=2x$ makes an angle of $\arctan(2/1)$
          with the $x$-axis.
          因此 $\sin\theta=2/\sqrt{5}$，$\cos\theta=1/\sqrt{5}$。
          \begin{equation*}
          \colvec{x \\ y}
           \mapsto
           \begin{mat}
             x\cdot(1/\sqrt{5})-y\cdot(2/\sqrt{5}) \\
             x\cdot(2/\sqrt{5})+y\cdot(1/\sqrt{5})
           \end{mat}
           \end{equation*}
        \partsitem
           $\colvec{x \\ y}
           \mapsto
           \begin{mat}
             x\cdot(1/\sqrt{5})-y\cdot(-2/\sqrt{5}) \\
             x\cdot(-2/\sqrt{5})+y\cdot(1/\sqrt{5})
           \end{mat}
           +\colvec[r]{1 \\ 1}
           =\begin{mat}
             x/\sqrt{5}+2y/\sqrt{5}+1 \\
             -2x/\sqrt{5}+y/\sqrt{5}+1
           \end{mat}$
      \end{exparts}
    
\end{ans}
\begin{ans}{3}
       \begin{exparts}
         \partsitem Let $f$ be distance-preserving ，而consider $f^{-1}$.
           陪域中任意两点可写成 $f(P_1)$ 和
           $f(P_2)$.
           因为 $f$ 保持距离，从 $f(P_1)$
           到 $f(P_2)$ 的距离等于从 $P_1$ 到 $P_2$ 的距离。
           这正是 $f^{-1}$
           保持距离所需的条件。
         \partsitem 任意平面图形 $F$ 通过显然保持距离的单位映射 $\map{\identity}{\Re^2}{\Re^2}$ 与自身全等。若 $F_1$ 通过某个 $f$ 与 $F_2$ 全等，则 $F_2$ 通过保持距离的 $f^{-1}$ 与 $F_1$ 全等。最后，若 $F_1$ 通过 $f$ 与 $F_2$ 全等且 $F_2$ 通过 $g$ 与 $F_3$ 全等，则 $F_1$ 通过 $\composed{g}{f}$ 与 $F_3$ 全等，且该复合映射保持距离。
       \end{exparts}
     
\end{ans}
\begin{ans}{4}
       两种计算的前两个分量分别为 $ax+cy+e$ 和 $bx+dy+f$。
     
\end{ans}
\begin{ans}{5}
      \begin{exparts}
        \partsitem The Pythagorean Theorem gives that
          三点
          collinear 当且仅当
          （按某种顺序记为 $P_1$、$P_2$、$P_3$）共线，当且仅当
          $\dist(P_1,P_2)+\dist(P_2,P_3)=\dist(P_1,P_3)$.
          当然，当 $f$ 保持距离时，这等价于
          $\dist(f(P_1),f(P_2))+\dist(f(P_2),f(P_3))=\dist(f(P_1),f(P_3))$,
          根据勾股定理，这又等价于
          $f(P_1)$、$f(P_2)$ 和 $f(P_3)$ 共线。

          对“介于其间”关系的论证类似（上式中 $P_2$ 位于
          $P_1$ 与 $P_3$ 之间）。

          若图形 $F$ 是三角形，则它是三条
          线段 $P_1P_2$、$P_2P_3$ 和 $P_1P_3$ 的并。
          前两段论证共同表明，线段这一性质保持不变。
          因此 $f(F)$ 是三条线段的并，也是
          三角形。

          以 $P$ 为圆心、半径为 $r$ 的圆 $C$ 是集合
          $\dist(P,Q)=r$ 的所有点 $Q$。
          应用保持距离映射 $f$，其像
          $f(C)$ 是满足条件
          $\dist(P,Q)=r$.
          由于 $\dist(P,Q)=\dist(f(P),f(Q))$，集合 $f(C)$ 也是
          圆，圆心为 $f(P)$、半径为 $r$。
        \partsitem 以下两个性质容易验证：(i)
          是直角三角形；(ii)
          两条直线平行。
        \partsitem 本节提到的一个性质是图形的“方向”。

          顶点按顺时针顺序为 $P_1$、$P_2$、$P_3$ 的三角形
          在保持距离映射下可能映到一个
          顶点按逆时针顺序为 $P_1$、$P_2$、$P_3$ 的三角形。
      \end{exparts}
    
\end{ans}
\protect \chapter {行列式}
\protect \section {第 I 节：定义}
\protect \subsection {四.I.1: 探索}
\begin{ans}{四.I.1.1}
       \begin{exparts*}
         \partsitem \( 4 \)
         \partsitem \( 3 \)
         \partsitem \( -12 \)
       \end{exparts*}
     
\end{ans}
\begin{ans}{四.I.1.2}
      \begin{exparts*}
        \partsitem \( 6 \)
        \partsitem \( 21 \)
        \partsitem \( 27 \)
      \end{exparts*}
    
\end{ans}
\begin{ans}{四.I.1.3}
      对于第一个，应用本节中的公式，注意任何
      含有 \( d \)、\( g \) 或 \( h \) 的项都是零，再化简即可。
      下三角矩阵也是如此。
    
\end{ans}
\begin{ans}{四.I.1.4}
       \begin{exparts}
         \partsitem 非奇异的，行列式为 \( -1 \)。
         \partsitem 非奇异的，行列式为 \( -1 \)。
         \partsitem 奇异的，行列式为 \( 0 \)。
       \end{exparts}
     
\end{ans}
\begin{ans}{四.I.1.5}
      \begin{exparts}
         \partsitem 非奇异的，行列式为 \( 3 \)。
         \partsitem 奇异的，行列式为 \( 0 \)。
         \partsitem 奇异的，行列式为 \( 0 \)。
       \end{exparts}
     
\end{ans}
\begin{ans}{四.I.1.6}
       \begin{exparts}
         \partsitem \( \det(B)=\det(A) \)，通过 \( -2\rho_1+\rho_2 \)
         \partsitem \( \det(B)=-\det(A) \)，通过
             \( \rho_2\leftrightarrow\rho_3 \)
         \partsitem \( \det(B)=(1/2)\cdot \det(A) \)，通过 \( (1/2)\rho_2 \)
       \end{exparts}
     
\end{ans}
\begin{ans}{四.I.1.7}
      高斯消元法正是这样做的。
      % [ 0 -7  1 -2]
      %  swap row 2 with row 4
      % [ 1  2  0  2]
      % [ 0 -7  1 -2]
      % [ 0  0 -1  3]
      % [ 0  0  1 -4]
      %  take 1 times row 3 plus row 4
      % [ 1  2  0  2]
      % [ 0 -7  1 -2]
      % [ 0  0 -1  3]
      % [ 0  0  0 -1]
      \begin{equation*}
        \begin{mat}
          1 &2 &0  &2 \\
          2 &4 &1  &0 \\
          0 &0 &-1 &3 \\
          3 &-1 &1 &4
        \end{mat}
        \grstep[-3\rho_1+\rho_4]{-2\rho_1+\rho_2}
        \grstep{\rho_2\leftrightarrow\rho_4}
        \grstep{\rho_3+\rho_4}
        \begin{mat}
          1 &2 &0  &2 \\
          0 &-7 &1  &-2 \\
          0 &0 &-1 &3 \\
          0 &0 &0 &-1
        \end{mat}
      \end{equation*}
      阶梯形矩阵的对角元之积为
       $1\cdot(-7)\cdot(-1)\cdot(-1)=-7$。
      在高斯消元的过程中没有行被倍乘，但有
      一次行对换，所以求行列式时须变号，
      得 $+7$。
    
\end{ans}
\begin{ans}{四.I.1.8}
       利用 $\nbyn{3}$ 矩阵的行列式公式，
       我们将左边展开
       \begin{equation*}
         1\cdot b\cdot c^2+1\cdot c\cdot a^2+1\cdot a\cdot b^2
          -b^2\cdot c\cdot 1 -c^2\cdot a\cdot 1-a^2\cdot b\cdot 1
       \end{equation*}
       而将右边展开则得
       \begin{equation*}
         (bc-ba-ac+a^2)\cdot(c-b)
         =c^2b-b^2c-bac+b^2a-ac^2+acb+a^2c-a^2b
       \end{equation*}
       于是只需验证两者相等即可。
       （\textit{注}。
       这是
       \definend{Vandermonde 行列式}\index{Vandermonde!determinant}%
       \index{determinant!Vandermonde}的 \( \nbyn{3} \) 情形，
       它出现在许多应用中）。
     
\end{ans}
\begin{ans}{四.I.1.9}
         这个方程
         \begin{equation*}
           0=
           \det(
             \begin{mat}
                12-x  &4  \\
                8    &8-x
             \end{mat}
           )
           =64-20x+x^2
           =(x-16)(x-4)
       \end{equation*}
       有根 \( x=16 \) 和 \( x=4 \)。
     
\end{ans}
\begin{ans}{四.I.1.10}
      我们先把矩阵化成阶梯形。
      首先，假设 \( a\neq 0 \) 且 \( ae-bd\neq 0 \)。
      \begin{multline*}
        \grstep{(1/a)\rho_1}
        \begin{mat}
           1   &b/a   &c/a   \\
           d   &e     &f     \\
           g   &h     &i
         \end{mat}
        \grstep[-g\rho_1+\rho_3]{-d\rho_1+\rho_2}
        \begin{mat}
           1   &b/a           &c/a           \\
           0   &(ae-bd)/a     &(af-cd)/a     \\
           0   &(ah-bg)/a     &(ai-cg)/a
         \end{mat}                                            \\
        \grstep{(a/(ae-bd))\rho_2}
        \begin{mat}
           1   &b/a           &c/a             \\
           0   &1             &(af-cd)/(ae-bd) \\
           0   &(ah-bg)/a     &(ai-cg)/a
         \end{mat}
      \end{multline*}
      这一步完成了计算。
      \begin{equation*}
        \grstep{((ah-bg)/a)\rho_2+\rho_3}
        \begin{mat}
           1   &b/a    &c/a             \\
           0   &1      &(af-cd)/(ae-bd)      \\
           0   &0      &(aei+bgf+cdh-hfa-idb-gec)/(ae-bd)
         \end{mat}
      \end{equation*}
      现在假设 $a\neq 0$ 且 \( ae-bd\neq 0 \)，
      原矩阵是非奇异的
      当且仅当上面第 \( 3,3 \) 处的元素非零。
      也就是说，在这些假设下，原矩阵
      是非奇异的当且仅当 $aei+bgf+cdh-hfa-idb-gec\neq 0$，
      即为所证。

      最后我们逐一看一下，若为方便起见而在前一段中所作的假设不成立，
      会发生什么。
      首先，若 \( a\neq 0 \) 但 \( ae-bd=0 \)，那么我们可以对换
      \begin{equation*}
        \begin{mat}
           1   &b/a           &c/a           \\
           0   &0             &(af-cd)/a     \\
           0   &(ah-bg)/a     &(ai-cg)/a
         \end{mat}
        \grstep{\rho_2\leftrightarrow\rho_3}
        \begin{mat}
           1   &b/a           &c/a           \\
           0   &(ah-bg)/a     &(ai-cg)/a     \\
           0   &0             &(af-cd)/a
         \end{mat}
      \end{equation*}
      从而得出结论：矩阵是非奇异的当且仅当
      \( ah-bg=0 \) 或 \( af-cd=0 \)。
      条件‘\( ah-bg=0 \) 或 \( af-cd=0 \)’等价于
      条件‘\( (ah-bg)(af-cd)=0 \)’。
      乘开并利用本情形的假设 $ae-bd=0$，以 $ae$ 代 $bd$，便得
      \begin{multline*}
         0=ahaf-ahcd-bgaf+bgcd
          =ahaf-ahcd-bgaf+aegc            \\
          =a(haf-hcd-bgf+egc)
      \end{multline*}
      由于 \( a\neq 0 \)，故矩阵
      是非奇异的当且仅当 \( haf-hcd-bgf+egc=0 \)。
      因此，在 \( a\neq 0 \) 且 \( ae-bd=0 \) 的情形下，
      当
      \( haf-hcd-bgf+egc-i(ae-bd)=0 \) 时矩阵
      是非奇异的。

      余下的情形是常规的。
      对于 \( a=0 \) 但 \( d\neq 0 \) 的情形，以及 \( a=0 \)、\( d=0 \)
      但 \( g\neq 0 \) 的情形，都先对换两行，然后
      照上面那样继续即可。
      \( a=0 \)、\( d=0 \) 且 \( g=0 \) 的情形则很简单\Dash 该矩阵是
      奇异的，因为它的列构成一个线性相关的组，而且
      行列式算出来为零。
    
\end{ans}
\begin{ans}{四.I.1.11}
       算出行列式并做一点代数运算，便得到
       \begin{align*}
          0
          &=y_1x+x_2y+x_1y_2-y_2x-x_1y-x_2y_1     \\
          (x_2-x_1)\cdot y
          &=(y_2-y_1)\cdot x+x_2y_1-x_1y_2              \\
          y
          &=\frac{y_2-y_1}{x_2-x_1}\cdot x+\frac{x_2y_1-x_1y_2}{x_2-x_1}
       \end{align*}
       注意这是一个直线的方程（特别地，
       其中含有熟悉的斜率表达式），
       并注意 \( (x_1,y_1) \) 与 \( (x_2,y_2) \) 都满足它。
     
\end{ans}
\begin{ans}{四.I.1.12}
      \begin{exparts}
        \partsitem 与本节引言中给出的公式比较，很容易。
        \partsitem 它对 \( \nbyn{2} \) 矩阵成立
          \begin{align*}
            \begin{pmat}{cc|c}
              h_{1,1} &h_{1,2} &h_{1,1} \\
              h_{2,1} &h_{2,2} &h_{2,1}
            \end{pmat}
            &=\begin{array}[t]{@{}l@{}}
              h_{1,1}h_{2,2}+h_{1,2}h_{2,1}  \\
              \>-h_{2,1}h_{1,2}-h_{2,2}h_{1,1}
            \end{array}                                     \\
            &=h_{1,1}h_{2,2}-h_{1,2}h_{2,1}
          \end{align*}
          但对 \( \nbyn{4} \) 矩阵不成立。
          例如，这个矩阵是
          奇异的，因为第二行与第三行相等
          \begin{equation*}
            \begin{mat}[r]
              1  &0  &0  &1   \\
              0  &1  &1  &0   \\
              0  &1  &1  &0   \\
             -1  &0  &0  &1
            \end{mat}
          \end{equation*}
          但按口诀的方案算并不得到零。
          \begin{equation*}
            \begin{pmat}{rrrr|rrr}
                       1  &0  &0  &1  &1  &0  &0  \\
                       0  &1  &1  &0  &0  &1  &1  \\
                       0  &1  &1  &0  &0  &1  &1  \\
                      -1  &0  &0  &1  &-1 &0  &0
            \end{pmat}
            =\begin{array}[t]{@{}l@{}}
               1+0+0+0 \\
               \>-(-1)-0-0-0
              \end{array}
           \end{equation*}
      \end{exparts}
     
\end{ans}
\begin{ans}{四.I.1.13}
      该行列式为
      $
        (x_2y_3-x_3y_2)\vec{e}_1
          +(x_3y_1-x_1y_3)\vec{e}_2
          +(x_1y_2-x_2y_1)\vec{e}_3
      $。
      要检验垂直性，我们检验它与第一个向量的点积
      为零
      \begin{equation*}
        \colvec{x_1 \\ x_2 \\ x_3}
        \dotprod
        \colvec{x_2y_3-x_3y_2 \\ x_3y_1-x_1y_3 \\ x_1y_2-x_2y_1}
        =x_1x_2y_3-x_1x_3y_2+x_2x_3y_1-x_1x_2y_3+x_1x_3y_2-x_2x_3y_1=0
      \end{equation*}
      而它与第二个向量的点积也为零。
      \begin{equation*}
        \colvec{y_1 \\ y_2 \\ y_3}
        \dotprod
        \colvec{x_2y_3-x_3y_2 \\ x_3y_1-x_1y_3 \\ x_1y_2-x_2y_1}
        =x_2y_1y_3-x_3y_1y_2+x_3y_1y_2-x_1y_2y_3+x_1y_2y_3-x_2y_1y_3=0
      \end{equation*}
    
\end{ans}
\begin{ans}{四.I.1.14}
       \begin{exparts}
        \partsitem 直接代入计算：
          乘积的行列式为
          \begin{align*}
             \det(\begin{mat}
                       a  &b  \\
                       c  &d
                    \end{mat}
                    \begin{mat}
                       w  &x  \\
                       y  &z
                    \end{mat}  )
             &=
             \det(\begin{mat}
                 aw+by  &ax+bz  \\
                 cw+dy  &cx+dz
              \end{mat} )                 \\
             &=
             \begin{array}[t]{@{}l@{}}
                 acwx+adwz+bcxy+bdyz  \\
                 \> -acwx-bcwz-adxy-bdyz
             \end{array}
          \end{align*}
          而行列式的乘积为
          \begin{equation*}
             \det(\begin{mat}
                a  &b  \\
                c  &d
             \end{mat})
             \cdot\det(\begin{mat}
                w  &x  \\
                y  &z
             \end{mat})
             =
             (ad-bc)\cdot (wz-xy)
          \end{equation*}
          验证两者相等很容易。
        \partsitem 利用前一小题。
       \end{exparts}
       \noindent 由上面两条立即得到：相似矩阵有相同的行列式
       $
          \det(PTP^{-1})=\det(P)\cdot\det(T)\cdot\det(P^{-1})
       $。
     
\end{ans}
\begin{ans}{四.I.1.15}
      一种方法是数出这些面积
      \begin{center}
        \includegraphics{det/mp/ch4.31}
      \end{center}
      即取整个矩形的面积，再减去
      左上矩形 $A$、上中三角形 $B$、
      右上三角形 $D$、左下三角形 $C$、
      下中三角形 $E$ 与右下矩形 $F$ 的面积
      \( (x_1+x_2)(y_1+y_2)-x_2y_1-(1/2)x_1y_1-(1/2)x_2y_2
              -(1/2)x_2y_2-(1/2)x_1y_1-x_2y_1 \)。
      化简即得行列式公式。

      这个行列式是上面那个的相反数；该公式
      区分了第二列是否在第一列的
      逆时针方向。
     
\end{ans}
\begin{ans}{四.I.1.16}
      对 \( \nbyn{2} \) 矩阵，利用
      引言中引用的公式，计算很容易。
      它对 \( \nbyn{3} \) 矩阵也成立；其
      计算是常规的。
    
\end{ans}
\begin{ans}{四.I.1.17}
      不成立。
      我们用 $\nbyn{2}$ 行列式来说明。
      回想常数可以每次提出一行。
      \begin{equation*}
         \det(
         \begin{mat}[r]
            2  &4  \\
            2  &6  \\
         \end{mat})
         =
         2\cdot\det(\begin{mat}[r]
            1  &2  \\
            2  &6  \\
         \end{mat})
         =
         2\cdot 2\cdot \det(\begin{mat}[r]
            1  &2  \\
            1  &3  \\
         \end{mat})
      \end{equation*}
      这与线性性矛盾（这里我们不需要 \( S \)，也就是说，我们可以
      取 $S$ 为零矩阵）。
    
\end{ans}
\begin{ans}{四.I.1.18}
       把 \( c \) 一次一行地提出来。
    
\end{ans}
\begin{ans}{四.I.1.19}
      不存在使该矩阵奇异的实数 \( \theta \)，
      因为矩阵的行列式
      \( \cos^2\theta+\sin^2\theta \) 永不为 $0$，它对所有 $\theta$
      都等于 $1$。
      从几何上看，关于标准基，
      这个矩阵表示
      平面旋转角度 \( \theta \) 的转动。
      每个这样的映射都是单射 \Dash  例如，它是可逆的。
    
\end{ans}
\begin{ans}{四.I.1.20}
      \answerasgiven
      设 \( P \) 为行列式中三个正项之和，
      \( -N \) 为三个负项之和。
      \( P \) 的最大值为
      \begin{equation*}
        9\cdot 8\cdot 7 +6\cdot 5\cdot 4 +3\cdot 2\cdot 1=630.
      \end{equation*}
      与 \( P \) 相容的 \( N \) 的最小值为
      \begin{equation*}
        9\cdot 6\cdot 1 +8\cdot 5\cdot 2 +7\cdot 4\cdot 3=218.
      \end{equation*}
      \( P \) 的任何改变都会使该和的降低量超过
      \( 4 \)。
      因此 \( 412 \) 是该行列式的最大值，而该行列式的一种形式是
      \begin{equation*}
         \begin{vmat}[r]
            9  &4  &2  \\
            3  &8  &6  \\
            5  &1  &7
         \end{vmat}.
      \end{equation*}
    
\end{ans}
\protect \subsection {四.I.2: 行列式的性质}
\begin{ans}{四.I.2.8}
      \begin{exparts}
       \partsitem 作 $2\rho_1+\rho_2$ 化为阶梯形，然后沿对角线相乘。
          行列式为~$0$。
          % sage: M = matrix(QQ, [[1,-2,1,2], [2,-4,1,0], [0,0,-1,0], [0,0,0,5]])
          % sage: M.determinant()
          % 0
        \partsitem 对换第二行与第三行即可将其化为阶梯形（同时改变行列式的符号）。
          沿对角线相乘得~$12$，故所给矩阵的行列式为~$-12$。
          % sage: M = matrix(QQ, [[1,1,-2], [0,0,4], [0,3,-6]])
          % sage: M.determinant()
          % -12
      \end{exparts}
    
\end{ans}
\begin{ans}{四.I.2.9}
      \begin{exparts}
      \partsitem \( \begin{vmat}[r]
                 3  &1  &2  \\
                 3  &1  &0  \\
                 0  &1  &4
               \end{vmat}
               =
               \begin{vmat}[r]
                 3  &1  &2  \\
                 0  &0  &-2 \\
                 0  &1  &4
               \end{vmat}
               =
               -\begin{vmat}[r]
                 3  &1  &2  \\
                 0  &1  &4  \\
                 0  &0  &-2
               \end{vmat}
               =6  \)
      \partsitem \( \begin{vmat}[r]
                 1  &0  &0  &1 \\
                 2  &1  &1  &0 \\
                -1  &0  &1  &0 \\
                 1  &1  &1  &0
               \end{vmat}
               =
               \begin{vmat}[r]
                 1  &0  &0  &1 \\
                 0  &1  &1  &-2\\
                 0  &0  &1  &1 \\
                 0  &1  &1  &-1
               \end{vmat}
               =
               \begin{vmat}[r]
                 1  &0  &0  &1 \\
                 0  &1  &1  &-2\\
                 0  &0  &1  &1 \\
                 0  &0  &0  &1
               \end{vmat}
               =1    \)
      \end{exparts}
    
\end{ans}
\begin{ans}{四.I.2.10}
     \begin{exparts}
      \partsitem \(
        \begin{vmat}[r]
                 2  &-1  \\
                 -1 &-1
               \end{vmat}
             =\begin{vmat}[r]
                 2  &-1  \\
                 0  &-3/2
               \end{vmat}=-3  \);
      \partsitem \( \begin{vmat}[r]
                 1  &1  &0  \\
                 3  &0  &2  \\
                 5  &2  &2
               \end{vmat}
              =\begin{vmat}[r]
                 1  &1  &0  \\
                 0  &-3 &2  \\
                 0  &-3 &2
               \end{vmat}
              =\begin{vmat}[r]
                 1  &1  &0  \\
                 0  &-3 &2  \\
                 0  &0  &0
               \end{vmat}
              =0 \)
     \end{exparts}
    
\end{ans}
\begin{ans}{四.I.2.11}
     什么时候行列式不为零？
      \begin{equation*}
         \begin{vmat}
           1  &0  &1  &-1  \\
           0  &1  &-2 &0   \\
           1  &0  &k  &0   \\
           0  &0  &1  &-1
         \end{vmat}
         =
         \begin{vmat}
           1  &0  &1  &-1  \\
           0  &1  &-2 &0   \\
           0  &0  &k-1&1   \\
           0  &0  &1  &-1
         \end{vmat}
      \end{equation*}
      显然，
      $k=1$ 给出非奇异性，从而行列式非零。
      若 $k\neq 1$，则
      用倍加 $(-1/k-1)\rho_3+\rho_4$ 可得到阶梯形。
      \begin{equation*}
         =
         \begin{vmat}
           1  &0  &1  &-1  \\
           0  &1  &-2 &0   \\
           0  &0  &k-1&1   \\
           0  &0  &0  &-1-(1/k-1)
         \end{vmat}
      \end{equation*}
      沿对角线相乘得 $(k-1)(-1-(1/k-1))=-(k-1)-1=-k$。
      因此，该矩阵的行列式非零，从而方程组
      有唯一解，当且仅当 \( k\neq 0 \)。
    
\end{ans}
\begin{ans}{四.I.2.12}
      \begin{exparts}
        \partsitem 行列式定义中的条件~(2)
          经对换 $\rho_1\leftrightarrow\rho_3$ 适用。
          \begin{equation*}
            \begin{vmat}
              h_{3,1}  &h_{3,2} &h_{3,3} \\
              h_{2,1}  &h_{2,2} &h_{2,3} \\
              h_{1,1}  &h_{1,2} &h_{1,3}
            \end{vmat}
            =
            -\begin{vmat}
              h_{1,1}  &h_{1,2} &h_{1,3} \\
              h_{2,1}  &h_{2,2} &h_{2,3} \\
              h_{3,1}  &h_{3,2} &h_{3,3}
            \end{vmat}
          \end{equation*}
        \partsitem 条件~(3) 适用。
          \begin{multline*}
            \begin{vmat}
             -h_{1,1}   &-h_{1,2}  &-h_{1,3} \\
             -2h_{2,1}  &-2h_{2,2} &-2h_{2,3} \\
             -3h_{3,1}  &-3h_{3,2} &-3h_{3,3}
            \end{vmat}
            =
            (-1)\cdot(-2)\cdot(-3)\cdot
            \begin{vmat}
             h_{1,1}   &h_{1,2}  &h_{1,3} \\
             h_{2,1}   &h_{2,2}  &h_{2,3} \\
             h_{3,1}   &h_{3,2}  &h_{3,3}
            \end{vmat}                                  \\
            =
            (-6)\cdot
            \begin{vmat}
             h_{1,1}   &h_{1,2}  &h_{1,3} \\
             h_{2,1}   &h_{2,2}  &h_{2,3} \\
             h_{3,1}   &h_{3,2}  &h_{3,3}
            \end{vmat}
          \end{multline*}
        \partsitem
          \begin{multline*}
          \begin{vmat}
            h_{1,1}+h_{3,1}  &h_{1,2}+h_{3,2} &h_{1,3}+h_{3,3} \\
            h_{2,1}          &h_{2,2}         &h_{2,3} \\
            5h_{3,1}         &5h_{3,2}        &5h_{3,3}
          \end{vmat}                                                 \\
          \begin{aligned}
          &=
          5\cdot
          \begin{vmat}
            h_{1,1}+h_{3,1}  &h_{1,2}+h_{3,2} &h_{1,3}+h_{3,3} \\
            h_{2,1}          &h_{2,2}         &h_{2,3} \\
            h_{3,1}          &h_{3,2}         &h_{3,3}
          \end{vmat}                                     \\
          &=5\cdot
          \begin{vmat}
            h_{1,1}          &h_{1,2}         &h_{1,3} \\
            h_{2,1}          &h_{2,2}         &h_{2,3} \\
            h_{3,1}          &h_{3,2}         &h_{3,3}
          \end{vmat}
          \end{aligned}
          \end{multline*}
      \end{exparts}
     
\end{ans}
\begin{ans}{四.I.2.13}
       对角矩阵已是阶梯形，所以
       行列式就是对角线上元素的乘积。
    
\end{ans}
\begin{ans}{四.I.2.14}
      它就是平凡子空间。
    
\end{ans}
\begin{ans}{四.I.2.15}
      把第二行加到第一行，得到的矩阵其第一行
      是第三行的 $x+y+z$ 倍。
    
\end{ans}
\begin{ans}{四.I.2.16}
      \begin{exparts}
        \partsitem
          $\begin{mat}[r]
            1
          \end{mat}$,
          $\begin{mat}[r]
            1  &-1  \\
            -1  &1
          \end{mat}$,
          $\begin{mat}[r]
            1   &-1  &1   \\
            -1  &1   &-1  \\
            1   &-1  &1
          \end{mat}$
        \partsitem $\nbyn{1}$ 情形的行列式为 $1$。
          在其余各种情形中，第二行都是第一行的相反数，
          因此矩阵是奇异的，行列式为零。
      \end{exparts}
    
\end{ans}
\begin{ans}{四.I.2.17}
      \begin{exparts}
        \partsitem
          $\begin{mat}[r]
            2
          \end{mat}$,
          $\begin{mat}[r]
            2  &3  \\
            3  &4
          \end{mat}$,
          $\begin{mat}[r]
            2  &3  &4  \\
            3  &4  &5  \\
            4  &5  &6
          \end{mat}$
        \partsitem $\nbyn{1}$ 与 $\nbyn{2}$ 情形给出如下结果。
          \begin{equation*}
            \begin{vmat}[r]
              2
            \end{vmat}
            =2
            \qquad
            \begin{vmat}[r]
              2  &3  \\
              3  &4
            \end{vmat}=-1
          \end{equation*}
          而 $n\geq 3$ 的 $\nbyn{n}$ 矩阵都是奇异的，例如
          \begin{equation*}
            \begin{vmat}[r]
              2  &3  &4  \\
              3  &4  &5  \\
              4  &5  &6
            \end{vmat}=0
          \end{equation*}
          因为第二行的两倍减去第一行
          等于第三行。
          验证这一点是例行的。
      \end{exparts}
    
\end{ans}
\begin{ans}{四.I.2.18}
      取这一个
      \begin{equation*}
         A=B=
         \begin{mat}[r]
           1  &2  \\
           3  &4
         \end{mat}
      \end{equation*}
      容易验证
      \begin{equation*}
         \deter{A+B}
         =
         \begin{vmat}[r]
           2  &4  \\
           6  &8
         \end{vmat}
         =-8
         \qquad
         \deter{A}+\deter{B}
         =
         -2-2
         =-4
      \end{equation*}
      顺便说一下，这也给出了一个标量乘法不被保持的例子：
      $\deter{2\cdot A}\neq 2\cdot\deter{A}$。
     
\end{ans}
\begin{ans}{四.I.2.19}
      不行，不能替换。
      \nearbyremark{rem:SwapRowsRedun}
      表明替换之后的四个条件会
      相互矛盾 \Dash  没有函数同时满足这四个条件。
    
\end{ans}
\begin{ans}{四.I.2.20}
      上三角矩阵已是阶梯形。

      下三角矩阵要么奇异要么非奇异。
      若它奇异，则其对角线上有零，从而其
      行列式（即零）确实是对角线上元素的乘积。
      若它非奇异，则其对角线上没有零，并且
      我们可以用高斯消元法将其化简为阶梯形，
      而不改变对角线。
    
\end{ans}
\begin{ans}{四.I.2.21}
      \begin{exparts}
        \partsitem 行列式定义中的诸性质表明
          \( \deter{M_i(k)}=k \)、
          \( \deter{P_{i,j}}=-1 \)
          以及
          \( \deter{C_{i,j}(k)}=1 \)。
        \partsitem 回想左乘每类矩阵的作用，三种情形都容易验证。
        \partsitem 若 \( TS \) 可逆，即 \( (TS)M=I \)，则
          由矩阵乘法的结合律，
          \( T(SM)=I \) 表明 \( T \) 可逆。
          所以若 \( T \) 不可逆，则 \( TS \) 也不可逆。
        \partsitem 若 \( T \) 奇异，则应用前面的答案：
          \( \deter{TS}=0 \) 且
          \( \deter{T}\cdot\deter{S}=0\cdot\deter{S}=0 \)。
          若 \( T \) 不奇异，则可将它写成初等矩阵的乘积
          $
            \deter{TS}=\deter{E_r\cdots E_1S}=
            \deter{E_r}\cdots\deter{E_1}\cdot\deter{S}=
            \deter{E_r\cdots E_1}\deter{S}=\deter{T}\deter{S}
          $。
        \partsitem \( 1=\deter{I}=\deter{T\cdot T^{-1}}
                     =\deter{T}\deter{T^{-1}} \)
      \end{exparts}
     
\end{ans}
\begin{ans}{四.I.2.22}
      \begin{exparts}
        \partsitem 我们必须证明：若
          \begin{equation*}
            T\grstep{k\rho_i+\rho_j}\hat{T}
          \end{equation*}
          则 $d(T)=\deter{TS}/\deter{S}
          =\deter{\hat{T}S}/\deter{S}=d(\hat{T})$。
          如果我们能证明先组合行、再相乘得到 \( \hat{T}S \)，
          与先相乘得到 \( TS \)、再组合行，结果相同
          （因为行列式 \( \deter{TS} \) 不受该倍加的影响，所以我们就有
          \( \deter{\hat{T}S}=\deter{TS} \)，从而 \( d(\hat{T})=d(T) \)），
          便完成证明。
          该论证如下：~把 \( TS \) 的第~\( i \) 行的 \( k \) 倍加到
          \( TS \) 的第~$j$ 行之后，第 \( j,p \) 元为
          \( (kt_{i,1}+t_{j,1})s_{1,p}+\dots+(kt_{i,r}+t_{j,r})s_{r,p} \)，
          这正是 \( \hat{T}S \) 的第 \( j,p \) 元。
        \partsitem 我们只需证明：对换
          $
            T\smash[b]{\grstep{\rho_i\swap\rho_j}}\hat{T}
          $
          之后相乘得到 \( \hat{T}S \)，与先把 \( T \) 乘以 \( S \)、再对换行，
          结果相同（因为行列式 \( \deter{TS} \) 在行对换下变号，
         所以我们就有 \( \deter{\hat{T}S}=-\deter{TS} \)，
          于是 \( d(\hat{T})=-d(T) \)）。
          该论证与前一个论证的进行方式相同。
        \partsitem 正如现在所期待的，我们只需证明：
          把某行乘以一个标量
          $
            T\smash[b]{\grstep{k\rho_i}}\hat{T}
          $
          之后计算 \( \hat{T}S \)，与先计算 \( TS \)、再把该行乘以 \( k \)，
          结果相同（由于行列式 \( \deter{TS} \) 在该倍乘下被乘以 \( k \)，
          我们就有 \( \deter{\hat{T}S}=k\deter{TS} \)，
          所以 \( d(\hat{T})=k\,d(T) \)）。
          论证的进行方式与上面完全一样。
        \partsitem 显然。
        \partsitem 由于我们已经证明 \( d(T) \) 是一个行列式，
          并且行列式函数（若存在）是唯一的，
          我们便有
          所以 \( \deter{T}=d(T)=\deter{TS}/\deter{S} \)。
      \end{exparts}
    
\end{ans}
\begin{ans}{四.I.2.23}
      我们首先论证：秩为 \( r \) 的矩阵有一个行列式非零的
      \( \nbyn{r} \) 子矩阵。
      秩为 \( r \) 的矩阵有一个由 \( r \) 行组成的线性无关集合。
      由这些行构成的矩阵行秩为 \( r \)，因而其
      列秩也为 \( r \)。
      结论：从这 \( r \) 行中可以取出一组由 \( r \) 列组成的线性无关集合，
      于是原矩阵有一个秩为 \( r \) 的
      \( \nbyn{r} \) 子矩阵。

      最后我们证明：若 \( r \) 是具有该性质的最大整数，
      则矩阵的秩为 \( r \)。
      由 \( r \) 的最大性，我们只需证明：
      若一个矩阵有行列式非零的
      \( \nbyn{k} \) 子矩阵，则该矩阵的秩至少为 \( k \)。
      考虑这样一个 \( \nbyn{k} \) 子矩阵。
      它的行是原矩阵的行的一部分，显然
      完整的行的集合是线性无关的。
      于是原矩阵的行秩至少为 \( k \)，而矩阵的行秩
      等于其秩。
   
\end{ans}
\begin{ans}{四.I.2.24}
      元素全为有理数的矩阵用高斯
      消元法化简时，只用到有理数运算即可化为阶梯形矩阵。
      因此对角线上的元素必定是有理数，从而
      对角线上元素的乘积是有理数。
    
\end{ans}
\begin{ans}{四.I.2.25}
       \answerasgiven
       行列式的值 \( (1-a^4)^3 \) 与
       \( B \)、\( C \)、\( D \) 的取值无关。
       因此操作~(e) 不改变行列式的值，
       只是改变了它的外观。
       于是在 (a)、(b)、(c)、(d)、(e) 中，
       共同点仅仅在于主要对象的外观被改变了。
       同样的共同点还出现在 (f)~改换玫瑰的名称标签、
       (g)~把一个十进制整数写成 \( 12 \) 进制、(h)~给百合花镀金、
       (i)~给政客粉饰门面，以及 (j)~授予荣誉学位。
    
\end{ans}
\protect \subsection {四.I.3: 置换展开}
\begin{ans}{四.I.3.17}
      记该矩阵为~$M$。
      \begin{center}
        \begin{tabular}{r|cccccc}
          置换
            &$\phi_1$ &$\phi_2$ &$\phi_3$ &$\phi_4$ &$\phi_5$ &$\phi_6$  \\
          \hline
          项
           &$1\cdot 5\cdot 9$
           &$1\cdot 6\cdot 8$
           &$2\cdot 4\cdot 9$
           &$2\cdot 6\cdot 7$
           &$3\cdot 4\cdot 8$
           &$3\cdot 5\cdot 7$
        \end{tabular}
      \end{center}
    
\end{ans}
\begin{ans}{四.I.3.18}
      我们可以把每个 $P_{\phi}$ 经过行对换变为单位矩阵。
      如果对换的次数是偶数，那么它的行列式为~$+1$；
      而如果对换的次数是奇数，那么行列式为~$-1$。
      \begin{center}
        \begin{tabular}{r|cccccc}
          置换
            &$\phi_1$ &$\phi_2$ &$\phi_3$ &$\phi_4$ &$\phi_5$ &$\phi_6$  \\
          \hline
          对换次数 &$0$
                  &$1$
                  &$1$
                  &$2$
                  &$2$
                  &$1$  \\
          $|P_{\phi_i}|$
           &$+1$
           &$-1$
           &$-1$
           &$+1$
           &$+1$
           &$-1$
        \end{tabular}
      \end{center}
      \textit{注。}
      在上表中，
      我们只是不断地对换，直到找到一个把我们带回到`$1,2,3$'的数为止。
      但不同的对换次数是否可能？
      如果一个人找到 $2$ 次对换而另一个人找到
      $4$ 次，这没有问题，因为两者给出的行列式都是~$+1$。
      但如果我们能用两种不同的方式对换，其中一种次数为偶数而另一种为奇数，
      那就有问题了。
      下一节将表明偶数与奇数混合的情形不可能出现。
    
\end{ans}
\begin{ans}{四.I.3.19}
      \begin{align*}
        \begin{vmat}
          2  &3 \\
          1  &5
        \end{vmat}
        &=
        2\cdot 5\cdot\deter{P_{\phi_1}}
        +
        1\cdot 3\cdot\deter{P_{\phi_2}}      \\
        &=
        2\cdot 5\cdot\begin{vmat}[r]
                           1  &0 \\
                           0  &1
                         \end{vmat}
        +
        1\cdot 3\cdot\begin{vmat}[r]
                           0  &1 \\
                           1  &0
                         \end{vmat}               \\
        &=2\cdot 5\cdot 1+1\cdot 3\cdot (-1)=7
        \end{align*}
    
\end{ans}
\begin{ans}{四.I.3.20}
          \begin{align*}
            \begin{vmat}[r]
              -1  &0  &1  \\
               3  &1  &4  \\
               2  &1  &5
            \end{vmat}
            &=
            \begin{aligned}[t]
              &(-1)(1)(5)\deter{P_{\phi_1}}
              +(-1)(4)(1)\deter{P_{\phi_2}}
              +(0)(3)(5)\deter{P_{\phi_3}}  \\
              &\hbox{}\quad\hbox{}
              +(0)(4)(2)\deter{P_{\phi_4}}
              +(1)(3)(1)\deter{P_{\phi_5}}
              +(1)(1)(2)\deter{P_{\phi_6}}
            \end{aligned}                        \\
            &=
            \begin{aligned}[t]
              &(-1)(1)(5)\cdot 1
              +(-1)(4)(1)\cdot (-1)
              +(0)(3)(5)\cdot (-1)  \\
              &\hbox{}\quad\hbox{}
              +(0)(4)(2)\cdot 1
              +(1)(3)(1)\cdot 1
              +(1)(1)(2)\cdot (-1)
            \end{aligned}                        \\
            &=-5+4+0+0+3-2=0
          \end{align*}
    
\end{ans}
\begin{ans}{四.I.3.21}
      \begin{exparts}
        \partsitem 该矩阵是奇异的。
          \begin{align*}
            \begin{vmat}[r]
              1  &2  &3  \\
              4  &5  &6  \\
              7  &8  &9
            \end{vmat}
            &=
            \begin{aligned}[t]
              &(1)(5)(9)\deter{P_{\phi_1}}
              +(1)(6)(8)\deter{P_{\phi_2}}
              +(2)(4)(9)\deter{P_{\phi_3}}  \\
              &\hbox{}\quad\hbox{}
              +(2)(6)(7)\deter{P_{\phi_4}}
              +(3)(4)(8)\deter{P_{\phi_5}}
              +(7)(5)(3)\deter{P_{\phi_6}}
            \end{aligned}                        \\
            &=0
          \end{align*}
        \partsitem 该矩阵是非奇异的。
          \begin{align*}
            \begin{vmat}[r]
              2  &2  &1  \\
              3  &-1 &0  \\
              -2 &0  &5
            \end{vmat}
            &=
            \begin{aligned}[t]
              &(2)(-1)(5)\deter{P_{\phi_1}}
              +(2)(0)(0)\deter{P_{\phi_2}}
              +(2)(3)(5)\deter{P_{\phi_3}} \\
              &\hbox{}\quad\hbox{}
              +(2)(0)(-2)\deter{P_{\phi_4}}
              +(1)(3)(0)\deter{P_{\phi_5}}
              +(-2)(-1)(1)\deter{P_{\phi_6}}
            \end{aligned}                      \\
            &=-42
        \end{align*}
      \end{exparts}
    
\end{ans}
\begin{ans}{四.I.3.22}
      \begin{exparts}
        \partsitem 高斯消元法给出
          \begin{equation*}
             \begin{vmat}[r]
                 2  &1  \\
                 3  &1
             \end{vmat}
             =
             \begin{vmat}[r]
                 2  &1  \\
                 0  &-1/2
             \end{vmat}
             =-1
          \end{equation*}
          而置换展开给出
          \begin{equation*}
             \begin{vmat}[r]
                 2  &1  \\
                 3  &1
             \end{vmat}
             =
             \begin{vmat}[r]
                 2  &0  \\
                 0  &1
             \end{vmat} +
             \begin{vmat}[r]
                 0  &1  \\
                 3  &0
             \end{vmat}
             =
             (2)(1)\begin{vmat}[r]
                 1  &0  \\
                 0  &1
             \end{vmat} +
             (1)(3)\begin{vmat}[r]
                 0  &1  \\
                 1  &0
             \end{vmat}
             =
             -1
          \end{equation*}
        \partsitem 高斯消元法给出
          \begin{equation*}
            \begin{vmat}[r]
              0  &1  &4  \\
              0  &2  &3  \\
              1  &5  &1
            \end{vmat}
            =
            -\begin{vmat}[r]
              1  &5  &1  \\
              0  &2  &3  \\
              0  &1  &4
            \end{vmat}
            =
            -\begin{vmat}[r]
              1  &5  &1  \\
              0  &2  &3  \\
              0  &0  &5/2
            \end{vmat}
            =-5
          \end{equation*}
          而置换展开给出
          \begin{align*}
            \begin{vmat}[r]
              0  &1  &4  \\
              0  &2  &3  \\
              1  &5  &1
            \end{vmat}
            &=
            \begin{aligned}[t]
            &(0)(2)(1)\deter{P_{\phi_1}}
            +(0)(3)(5)\deter{P_{\phi_2}}
            +(1)(0)(1)\deter{P_{\phi_3}}  \\
            &\hbox{}\quad\hbox{}
            +(1)(3)(1)\deter{P_{\phi_4}}
            +(4)(0)(5)\deter{P_{\phi_5}}
            +(4)(2)(1)\deter{P_{\phi_6}}
            \end{aligned}                         \\
            &=-5
          \end{align*}
      \end{exparts}
    
\end{ans}
\begin{ans}{四.I.3.23}
        仿照 \nearbyexample{ex:SamplePermExp} 可得
          \begin{align*}
             \begin{vmat}
                t_{1,1}  &t_{1,2}  &t_{1,3}  \\
                t_{2,1}  &t_{2,2}  &t_{2,3}  \\
                t_{3,1}  &t_{3,2}  &t_{3,3}
             \end{vmat}
             &=\begin{aligned}[t]
                 &t_{1,1}t_{2,2}t_{3,3}\deter{P_{\phi_1}}
                   +t_{1,1}t_{2,3}t_{3,2}\deter{P_{\phi_2}} \\
                 &\hbox{}\quad\hbox{}
                  +t_{1,2}t_{2,1}t_{3,3}\deter{P_{\phi_3}}
                             +t_{1,2}t_{2,3}t_{3,1}\deter{P_{\phi_4}} \\
                 &\hbox{}\quad\hbox{}
                  +t_{1,3}t_{2,1}t_{3,2}\deter{P_{\phi_5}}
                             +t_{1,3}t_{2,2}t_{3,1}\deter{P_{\phi_6}}
               \end{aligned}                                               \\
             &=\begin{aligned}[t]
                 & t_{1,1}t_{2,2}t_{3,3}(+1)
                   +t_{1,1}t_{2,3}t_{3,2}(-1)  \\
                 &\hbox{}\quad\hbox{}
                   +t_{1,2}t_{2,1}t_{3,3}(-1)
                    +t_{1,2}t_{2,3}t_{3,1}(+1) \\
                 &\hbox{}\quad\hbox{}
                   +t_{1,3}t_{2,1}t_{3,2}(+1)
                         +t_{1,3}t_{2,2}t_{3,1}(-1)
               \end{aligned}
          \end{align*}
     
\end{ans}
\begin{ans}{四.I.3.24}
      下面是所有满足 $\phi(1)=1$ 的置换
      \begin{align*}
        &\phi_1=\sequence{1,2,3,4}
        \quad
        \phi_2=\sequence{1,2,4,3}
        \quad
        \phi_3=\sequence{1,3,2,4}  \\
        &\quad
        \phi_4=\sequence{1,3,4,2}
        \quad
        \phi_5=\sequence{1,4,2,3}
        \quad
        \phi_6=\sequence{1,4,3,2}
      \end{align*}
      下面是满足 $\phi(1)=1$ 的那些
      \begin{align*}
        &\phi_7=\sequence{2,1,3,4}
        \quad
        \phi_8=\sequence{2,1,4,3}
        \quad
        \phi_9=\sequence{2,3,1,4}    \\
        &\quad
        \phi_{10}=\sequence{2,3,4,1}
        \quad
        \phi_{11}=\sequence{2,4,1,3}
        \quad
        \phi_{12}=\sequence{2,4,3,1}
      \end{align*}
      下面是满足 $\phi(1)=3$ 的那些
      \begin{align*}
        &\phi_{13}=\sequence{3,1,2,4}
        \quad
        \phi_{14}=\sequence{3,1,4,2}
        \quad
        \phi_{15}=\sequence{3,2,1,4}  \\
        &\quad
        \phi_{16}=\sequence{3,2,4,1}
        \quad
        \phi_{17}=\sequence{3,4,1,2}
        \quad
        \phi_{18}=\sequence{3,4,2,1}
      \end{align*}
      以及满足 $\phi(1)=4$ 的那些。
      \begin{align*}
        &\phi_{19}=\sequence{4,1,2,3}
        \quad
        \phi_{20}=\sequence{4,1,3,2}
        \quad
        \phi_{21}=\sequence{4,2,1,3}  \\
        &\quad
        \phi_{22}=\sequence{4,2,3,1}
        \quad
        \phi_{23}=\sequence{4,3,1,2}
        \quad
        \phi_{24}=\sequence{4,3,2,1}
      \end{align*}
    
\end{ans}
\begin{ans}{四.I.3.25}
      下列每一个都容易验证。
      \begin{exparts}
        \partsitem
          \begin{tabular}[t]{r|cc}
            置换 &$\phi_1$  &$\phi_2$ \\
             \hline
            逆     &$\phi_1$  &$\phi_2$
          \end{tabular}
        \partsitem
          \begin{tabular}[t]{r|cccccc}
            置换
              &$\phi_1$ &$\phi_2$ &$\phi_3$ &$\phi_4$ &$\phi_5$ &$\phi_6$ \\
            \hline
            逆
              &$\phi_1$ &$\phi_2$ &$\phi_3$ &$\phi_5$ &$\phi_4$ &$\phi_6$
          \end{tabular}
      \end{exparts}
    
\end{ans}
\begin{ans}{四.I.3.26}
         对于「如果」部分，取 $k_1=k_2=1$ 即得
         \nearbydefinition{def:multilinear} 的第一个条件，
         取 $k_2=0$ 即得第二个条件。

         「仅当」部分也是常规推导。
         由
         $
           f(\vec{\rho}_1,\dots,k_1\vec{v}_1+k_2\vec{v}_2,
             \dots,\vec{\rho}_n)
         $
         出发，\nearbydefinition{def:multilinear} 的第一个条件给出
         $
           =
           f(\vec{\rho}_1,\dots,k_1\vec{v}_1,\dots,\vec{\rho}_n)+
           f(\vec{\rho}_1,\dots,k_2\vec{v}_2,\dots,\vec{\rho}_n)
         $
         再对第二个条件应用两次，即得结论。
       
\end{ans}
\begin{ans}{四.I.3.27}
       它们都会变成原来的两倍。
    
\end{ans}
\begin{ans}{四.I.3.28}
      对于第二个论断，
      给定一个矩阵，先转置，再对换行，然后转置回来。
      所得结果是对换了列，而行列式改变了因子
      \( -1 \)。
      第三个论断类似：~给定
      一个矩阵，先转置，对如今已是行的东西应用多重线性，
      然后把所得的各个矩阵转置回来。
    
\end{ans}
\begin{ans}{四.I.3.29}
      行列式非零的 \( \nbyn{n} \) 矩阵的秩为
      \( n \)，所以它的各列构成 \( \Re^n \) 的一组基。
    
\end{ans}
\begin{ans}{四.I.3.30}
      假。
      \begin{equation*}
        \begin{vmat}[r]
          0 &1  &1  \\
          1 &0  &1 \\
          1  &1 &0
        \end{vmat}
        =2
      \end{equation*}
     
\end{ans}
\begin{ans}{四.I.3.31}
      \begin{exparts}
        \partsitem 第一行元素的列下标有
          五种选择。
          于是第二行元素的列下标有
          四种选择。
          依此类推，我们得到 $5\cdot 4\cdot 3\cdot 2\cdot 1=120$。
          （另见下一题。）
        \partsitem 一旦选定了第一行的第二列，
          其余元素可以有 \( 4\cdot 3\cdot 2\cdot 1=24 \)
          种选法。
      \end{exparts}
    
\end{ans}
\begin{ans}{四.I.3.32}
       \( n\cdot(n-1)\cdots 2\cdot 1=n! \)
    
\end{ans}
\begin{ans}{四.I.3.33}
      \cite{SchmidtSO}
      我们将证明 $P\trans{P}=I$；$\trans{P}P=I$ 的论证类似。
      $P\trans{P}$ 的 $i,j$ 元是形如
      $p_{i,k}q_{k,j}$ 的项之和，其中 $\trans{P}$ 的元素用 $q$ 记，
      即 $q_{k,j}=p_{j,k}$。
      于是 $P\trans{P}$ 的 $i,j$ 元是和
      \(
        \sum_{k=1}^n p_{i,k}p_{j,k}
      \)。
      但 $p_{i,k}$ 通常为 $0$，所以 $P_{i,k}P_{j,k}$ 通常为 $0$。
      $P_{i,k}$ 非零仅当它为 $1$，
      而此时再无其他 $i^\prime\neq i$ 使 $P_{i^\prime,k}$
      非零（$i$ 是第~$k$ 列中含 $1$ 的唯一一行）。
      换句话说，
      \begin{equation*}
        \sum_{k=1}^n p_{i,k}p_{j,k}=
        \begin{cases}
          1  &i=j  \\
          0  &\text{其他}
        \end{cases}
      \end{equation*}
      而这正是单位矩阵各元素的公式。
    
\end{ans}
\begin{ans}{四.I.3.34}
      在
      \( \deter{A}=\deter{\trans{A}}=\deter{-A}=(-1)^n\deter{A} \)
      中指数 $n$ 必须为偶数。
    
\end{ans}
\begin{ans}{四.I.3.35}
      证明三个零的任何放置都不够是常规性的。
      四个零确实够了；把它们全放在同一
      行或同一列。
    
\end{ans}
\begin{ans}{四.I.3.36}
      $n=3$ 的情形表明了该怎么做。
      行倍加变换
      $-x_1\rho_2+\rho_3$ 与 $-x_1\rho_1+\rho_2$
      给出
      \begin{multline*}
        \begin{vmat}
          1     &1     &1     \\
          x_1   &x_2   &x_3   \\
          x_1^2 &x_2^2 &x_3^2
        \end{vmat}
        =
        \begin{vmat}
          1     &1             &1             \\
          x_1   &x_2           &x_3           \\
          0     &(-x_1+x_2)x_2 &(-x_1+x_3)x_3
        \end{vmat}                                           \\
        =
        \begin{vmat}
          1     &1             &1             \\
          0     &-x_1+x_2      &-x_1+x_3      \\
          0     &(-x_1+x_2)x_2 &(-x_1+x_3)x_3
        \end{vmat}
      \end{multline*}
      再作行倍加变换 $x_2\rho_2+\rho_3$ 便得
      所要的结果。
      \begin{equation*}
        =
        \begin{vmat}
          1     &1             &1                     \\
          0     &-x_1+x_2      &-x_1+x_3              \\
          0     &0             &(-x_1+x_3)(-x_2+x_3)
        \end{vmat}
        =(x_2-x_1)(x_3-x_1)(x_3-x_2)
      \end{equation*}
    
\end{ans}
\begin{ans}{四.I.3.37}
      设 \( T \) 为 \( \nbyn{n} \) 的，
      \( J \) 为 \( \nbyn{p} \) 的，
      \( K \) 为 \( \nbyn{q} \) 的。
      应用置换展开公式
      \begin{equation*}
        \deter{T}=
        \sum_{\text{\scriptsize 置换 }\phi}
                t_{1,\phi(1)}t_{2,\phi(2)}\dots t_{n,\phi(n)}
                \deter{P_{\phi}}
      \end{equation*}
      由于 \( T \) 的右上部分全为零，如果一个
      \( \phi \) 的前
      \( p \) 个列编号 \( \phi(1),\dots,\phi(p) \) 中至少有一个取自 \( p+1,\dots,n \)，
      那么由 \( \phi \) 产生的项为 \( 0 \)
      （例如，若 \( \phi(1)=n \) 则
      \( t_{1,\phi(1)}t_{2,\phi(2)}\dots t_{n,\phi(n)} \)
      为 \( 0 \)）。
      所以上式化简为对具有两段的全部置换求和：
      先重排 \( 1,\dots,p \)，随后是
      \( p+1,\dots,p+q \)。
      为看出这给出 \( \deter{J}\cdot\deter{K}  \)，将其展开。
      \begin{multline*}
         \bigg[\sum_{\substack{\text{\scriptsize 置换 }\phi_1 \\
                              \text{\scriptsize 作用于 } 1,\dots,p}}
               t_{1,\phi_1(1)}\cdots t_{p,\phi_1(p)}
               \deter{P_{\phi_1}} \bigg]                                  \\
         \cdot
         \bigg[\sum_{\substack{\text{\scriptsize 置换 }\phi_2 \\
                              \text{\scriptsize 作用于 } p+1,\dots,p+q}}
               t_{p+1,\phi_2(p+1)}\cdots t_{p+q,\phi_2(p+q)}
               \deter{P_{\phi_2}}                          \bigg]
      \end{multline*}
    
\end{ans}
\begin{ans}{四.I.3.38}
      $n=3$ 的情形表明了会发生什么。
      \begin{equation*}
        \deter{T-rI}
        =
        \begin{vmat}
          t_{1,1}-x  &t_{1,2}   &t_{1,3}  \\
          t_{2,1}    &t_{2,2}-x &t_{2,3}  \\
          t_{3,1}    &t_{3,2}   &t_{3,3}-x
        \end{vmat}
      \end{equation*}
      置换展开中的每一项都有三个取自
      矩阵元素的因子（例如，$(t_{1,1}-x)(t_{2,2}-x)(t_{3,3}-x)$
      与 $(t_{1,1}-x)(t_{2,3})(t_{3,2})$），所以该行列式
      可以表示为 $x$ 的 $3$ 次多项式。
      这样的多项式至多有 $3$ 个根。

      一般地，置换展开表明
      行列式是一些项的和，每一项
      有 \( n \) 个因子，给出一个 $n$ 次多项式。
      \( n \) 次多项式至多有 \( n \) 个根。
    
\end{ans}
\begin{ans}{四.I.3.39}
      \answerasgiven
      当行列式的两行对换时，行列式的
      符号改变。
      当一个三阶行列式的行作置换时，得到 \( 3 \)
      个正的与 \( 3 \) 个绝对值相等的
      负行列式。
      于是 \( 9! \) 个行列式分成 \( 9!/6 \) 组，每一组
      之和为零。
    
\end{ans}
\begin{ans}{四.I.3.40}
      \answerasgiven
      当任一列的元素分别从另外两列的元素中减去后，所得
      行列式的两列元素成比例，所以行列式为零。
      也就是，
      \begin{equation*}
        \begin{vmat}
          2  &1    &x-4  \\
          4  &2    &x-3  \\
          6  &3    &x-10
        \end{vmat}=
        \begin{vmat}
          1    &x-3  &-1   \\
          2    &x-1  &-2   \\
          3    &x-7  &-3
        \end{vmat}=
        \begin{vmat}
          x-2  &-1   &-2   \\
          x+1  &-2   &-4   \\
          x-4  &-3   &-6
        \end{vmat}=0.
      \end{equation*}
    
\end{ans}
\begin{ans}{四.I.3.41}
      \answerasgiven
      设
      \begin{equation*}
        \begin{array}{ccc}
          a  &b  &c  \\
          d  &e  &f  \\
          g  &h  &i
        \end{array}
      \end{equation*}
      其幻和为 \( N=S/3 \)。
      那么
      \begin{align*}
        N
        &=(a+e+i)+(d+e+f)+(g+e+c)   \\
        &\quad\text{}-(a+d+g)-(c+f+i)=3e
      \end{align*}
      且 \( S=9e \)。
      因此，把行与列相加，
      \begin{equation*}
        D=
        \begin{vmat}
          a  &b  &c  \\
          d  &e  &f  \\
          g  &h  &i
        \end{vmat}
        =\begin{vmat}
          a  &b  &c  \\
          d  &e  &f  \\
         3e  &3e &3e
        \end{vmat}
        =\begin{vmat}
          a  &b  &3e \\
          d  &e  &3e \\
         3e  &3e &9e
        \end{vmat}
        =\begin{vmat}
          a  &b  &e  \\
          d  &e  &e  \\
          1  &1  &1
        \end{vmat}S.
      \end{equation*}
    
\end{ans}
\begin{ans}{四.I.3.42}
      \answerasgiven
      记所论行列式为 \( D_n \)，记
      \( i \) 行 \( j \) 列的元素为 \( a_{i,j} \)。
      那么由元素的构造法则我们有
      \begin{equation*}
        a_{i,j}=a_{i,j-1}+a_{i-1,j},
        \qquad a_{1,j}=a_{i,1}=1.
      \end{equation*}
      从最后一对行开始，将 \( D_n \) 的每一行从其后一行中减去。
      经过 \( n-1 \) 次这样的减法，上面的等式表明元素
      \( a_{i,j} \) 被元素 \( a_{i,j-1} \) 所取代，而第一列中除
      \( a_{1,1}=1 \) 外的所有元素都变成零。
      现在从最后一对列开始，将每一列从其后一列中减去。
      经过这一过程，如上面的关系式所示，元素 \( a_{i,j-1} \) 被
      \( a_{i-1,j-1} \) 所取代。
      这两个运算的结果是把 \( a_{i,j} \) 换成
      \( a_{i-1,j-1} \)，并使第一行与
      第一列中的每个元素都化为零。
      于是 \( D_n=D_{n+i} \)，从而
      \begin{equation*}
        D_n=D_{n-1}=D_{n-2}=\dots=D_2=1.
      \end{equation*}
    
\end{ans}
\protect \subsection {四.I.4: 行列式的存在性}
\begin{ans}{四.I.4.11}
      这是 $\nbyn{2}$ 矩阵的行列式的置换展开
      \begin{equation*}
        \begin{vmat}
          a  &b  \\
          c  &d
        \end{vmat}
         =ad\cdot \begin{vmat}[r]
                    1  &0  \\
                    0  &1
                  \end{vmat}
         +bc\cdot \begin{vmat}[r]
                    0  &1 \\
                    1  &0
                  \end{vmat}
      \end{equation*}
      以及它的转置的行列式的置换展开。
      \begin{equation*}
         \begin{vmat}
            a  &c  \\
            b  &d
         \end{vmat}
         =ad\cdot\begin{vmat}[r]
                    1  &0  \\
                    0  &1
                 \end{vmat}
         +cb\cdot\begin{vmat}[r]
                   0  &1 \\
                   1  &0
                 \end{vmat}
       \end{equation*}
       与小节中描述的 $\nbyn{3}$ 展开一样，对应项的置换矩阵互为转置
       （尽管每个矩阵都是自转置的这一事实把这一点掩盖了）。
    
\end{ans}
\begin{ans}{四.I.4.12}
      每一个都容易验证。
      \begin{exparts}
        \partsitem
          \begin{tabular}[t]{r|cc}
            \textit{置换} &$\phi_1$  &$\phi_2$ \\
             \hline
            \textit{逆}     &$\phi_1$  &$\phi_2$
          \end{tabular}
        \partsitem
          \begin{tabular}[t]{r|cccccc}
            \textit{置换}
              &$\phi_1$ &$\phi_2$ &$\phi_3$ &$\phi_4$ &$\phi_5$ &$\phi_6$ \\
            \hline
            \textit{逆}
              &$\phi_1$ &$\phi_2$ &$\phi_3$ &$\phi_5$ &$\phi_4$ &$\phi_6$
          \end{tabular}
      \end{exparts}
    
\end{ans}
\begin{ans}{四.I.4.13}
      \begin{exparts}
        \partsitem \( \sgn(\phi_1)=+1 \), \( \sgn(\phi_2)=-1 \)
        \partsitem \( \sgn(\phi_1)=+1 \), \( \sgn(\phi_2)=-1 \),
              \( \sgn(\phi_3)=-1 \), \( \sgn(\phi_4)=+1 \),
              \( \sgn(\phi_5)=+1 \), \( \sgn(\phi_6)=-1 \)
      \end{exparts}
     
\end{ans}
\begin{ans}{四.I.4.14}
      要在置换展开中得到非零项，我们必须取 \( 1,2 \) 处的元素与
      \( 4,3 \) 处的元素。
      选定这两处之后，我们还必须取 \( 2,1 \) 处的元素与 \( 3,4 \) 处的元素。
      \( \sequence{2,1,4,3} \) 的符号是 \( +1 \)，因为由
      \begin{equation*}
        \begin{mat}[r]
          0  &1  &0  &0  \\
          1  &0  &0  &0  \\
          0  &0  &0  &1  \\
          0  &0  &1  &0
        \end{mat}
      \end{equation*}
      经过两次行对换 $\rho_1\leftrightarrow\rho_2$ 与
      $\rho_3\leftrightarrow\rho_4$ 即可得到单位矩阵。
    
\end{ans}
\begin{ans}{四.I.4.15}
      规律如下。
      \begin{center}
         \begin{tabular}[b]{c|ccccccc}
           \( i \) &\( 1 \) &\( 2 \) &\( 3 \) &\( 4 \) &\( 5 \)
              &\( 6 \) &\ldots \\
           \hline
           \( \sgn(\phi_i) \) &\( +1 \) &\( -1 \) &\( -1 \) &\( +1 \)
              &\( +1 \) &\( -1 \) &\ldots
         \end{tabular}
      \end{center}
      于是为求 $\phi_{n!}$ 的符号，我们用 $n!-1$ 减一，再看它除以四的余数。
      若余数是 $1$ 或 $2$，则符号为 $-1$；否则为 $+1$。
      对 $n>4$，数 $n!$ 能被四整除，所以 $n!-1$ 除以四的余数是
      $-1$（更准确地说，余数是 $3$），因此符号是 $+1$。
      $n=1$ 的情形符号为 $+1$，$n=2$ 的情形符号为 $-1$，
      $n=3$ 的情形符号为 $-1$。
    
\end{ans}
\begin{ans}{四.I.4.16}
      \begin{exparts}
       \partsitem 我们可以把置换看成
         \( \map{\phi}{\set{1,\ldots,n}}{\set{1,\ldots,n}} \)
         这样单射且满射的映射。
         任何单射且满射的映射都有逆。
       \partsitem 若从 \( P_\phi \)
         变到单位矩阵总需要奇数次对换，那么从单位矩阵变到
         \( P_\phi \) 也总需要奇数次对换（任何对换都是可逆的）。
       \partsitem 这就是第一问。
      \end{exparts}
    
\end{ans}
\begin{ans}{四.I.4.17}
        若 $\phi(i)=j$ 则 $\phi^{-1}(j)=i$。
        于是注意到 $P_{\phi}$ 在 $i,j$ 处有一个 $1$ 当且仅当 $\phi(i)=j$，
        而 $P_{\phi^{-1}}$ 在 $j,i$ 处有一个 $1$ 当且仅当 $\phi^{-1}(j)=i$，
        结论即得。
      
\end{ans}
\begin{ans}{四.I.4.18}
      这并不是说 \( m \) 是变成单位矩阵所需的最少对换次数，
      也不是说 \( m \) 是最多的。
      它说的是存在一种恰好经 \( m \) 步对换变成单位矩阵的方法。

      设 \( \iota_j \) 是相对于前面某行来说第一个出现逆序的行，
      并设 \( \iota_k \) 是给出该逆序的第一行。
      我们有这样一个行区间。
      \begin{equation*}
         \begin{mat}
           \vdots      \\
           \iota_k     \\
           \iota_{r_1} \\
           \vdots      \\
           \iota_{r_s} \\
           \iota_j     \\
           \vdots
         \end{mat}
         \qquad j<k<r_1<\dots <r_s
      \end{equation*}
      作对换。
      \begin{equation*}
         \begin{mat}
           \vdots      \\
           \iota_j     \\
           \iota_{r_1} \\
           \vdots      \\
           \iota_{r_s} \\
           \iota_k     \\
           \vdots
         \end{mat}
      \end{equation*}
      第二个矩阵的逆序少了一个，因为该区间内的逆序少了一个
      （\( s \) 对 \( s+1 \)），而涉及区间之外各行的逆序不受影响。

      如此进行下去，每作一次行对换就把逆序的个数减少一个。
      当逆序不复存在时，结果就是单位矩阵。

      与\nearbycorollary{cor:ParityInversEqParitySwaps}的对比在于，
      本练习的陈述是一个“存在”命题：
      存在一种恰好经~$m$ 步对换变成单位矩阵的方法。
      而该推论是一个“对所有”命题：对变成单位矩阵的所有对换方式，
      奇偶性（是偶还是奇）都相同。
    
\end{ans}
\begin{ans}{四.I.4.19}
      \begin{exparts}
        \partsitem 首先，$g(\phi_1)$ 是单独一个因子 $2-1$ 的乘积，
          所以 $g(\phi_1)=1$。
          其次，$g(\phi_2)$ 是单独一个因子 $1-2$ 的乘积，
          所以 $g(\phi_2)=-1$。
        \partsitem
            \begin{tabular}[t]{r|cccccc}
               \textit{置换 $\phi$}
                 &$\phi_{1}$ &$\phi_{2}$ &$\phi_{3}$
                    &$\phi_{4}$ &$\phi_{5}$ &$\phi_{6}$ \\
               \hline
               $g(\phi)$ &$2$ &$-2$ &$-2$ &$2$ &$2$ &$-2$
            \end{tabular}
        \partsitem 它是一些非零项的乘积。
        \partsitem 注意 \( \phi(j)-\phi(i) \) 为负当且仅当
           \( \iota_{\phi(j)} \) 与 \( \iota_{\phi(i)} \) 处于它们通常顺序的一个逆序之中。
      \end{exparts}
     
\end{ans}
\protect \section {第 II 节：行列式的几何}
\protect \subsection {四.II.1: 作为大小函数的行列式}
\begin{ans}{四.II.1.8}
      求解
      \begin{equation*}
        c_1\colvec[r]{3 \\ 3 \\ 1}
        +c_2\colvec[r]{2 \\ 6 \\ 1}
        +c_3\colvec[r]{1 \\ 0 \\ 5}
        =\colvec[r]{4 \\ 1 \\ 2}
      \end{equation*}
      得到唯一解
      \( c_3=11/57 \)、\( c_2=-40/57 \) 以及 \( c_1=99/57 \)。
      由于 \( c_1>1 \)，该向量不在箱体内。
    
\end{ans}
\begin{ans}{四.II.1.9}
      对每一小题，求出行列式并取绝对值。
      \begin{exparts*}
        \partsitem $7$
        \partsitem $0$
         % sage: A = matrix(QQ, [[2, 3, 8], [1, -2, -3], [0, 4, 8]])
         % sage: A
         % [ 2  3  8]
         % [ 1 -2 -3]
         % [ 0  4  8]
         % sage: det(A)
         % 0
        \partsitem $-58$
         % sage: A = matrix(QQ, [[1,2,-1,0], [2,2,3,1], [0,2,0,0], [1,2,5,7]])
         % sage: A
         % [ 1  2 -1  0]
         % [ 2  2  3  1]
         % [ 0  2  0  0]
         % [ 1  2  5  7]
         % sage: det(A)
         % -58
      \end{exparts*}
    
\end{ans}
\begin{ans}{四.II.1.10}
      我们画这幅图就是为了误导读者。
      左图并不是由两个向量形成的箱体。
      若把它平移到原点，它就变成由下面这个序列形成的箱体。
      \begin{equation*}
        \sequence{
          \colvec[r]{0 \\ 1},
          \colvec[r]{2 \\ 0}
       }
      \end{equation*}
      而矩阵作用后的像是由下面这个序列形成的箱体。
      \begin{equation*}
        \sequence{
          \colvec[r]{1 \\ 1},
          \colvec[r]{4 \\ 0}
         }
      \end{equation*}
      其面积为 $4$。
     
\end{ans}
\begin{ans}{四.II.1.11}
      将该平行四边形平移至从原点出发，
      它就变成由下列向量形成的箱体
      \begin{equation*}
        \sequence{
          \colvec[r]{3 \\ 0},
          \colvec[r]{2 \\ 1}
        }
      \end{equation*}
      于是这个行列式的绝对值
      容易算出为 $3$。
      \begin{equation*}
        \begin{vmat}[r]
          3  &2  \\
          0  &1
        \end{vmat}=3
      \end{equation*}
     
\end{ans}
\begin{ans}{四.II.1.12}
     \begin{exparts*}
        \partsitem \( 3 \)
        \partsitem \( 9 \)
        \partsitem $1/9$
      \end{exparts*}
    
\end{ans}
\begin{ans}{四.II.1.13}
      \begin{exparts}
        \partsitem
          用高斯消元法
          \begin{equation*}
          \grstep[\rho_1+\rho_3]{-3\rho_1+\rho_2}
          \begin{mat}
            1 &0   &-1    \\
            0 &1   &4     \\
            0 &0   &2
          \end{mat}
          \end{equation*}
          得到行列式为~$+2$。
          符号为正，所以该变换保持定向。
        \partsitem
          箱体的大小就是下述行列式的值。
          \begin{equation*}
            \begin{vmat}
              1 &2  &1 \\
             -1 &0  &1 \\
              2 &-1 &0
            \end{vmat}
            =+6
          \end{equation*}
          定向为正。
        \partsitem
          由于该变换关于标准基由所给矩阵表示，
          而关于标准基，向量就表示它们自身，
          所以要求这些向量在该变换下的像，
          只需从左边用矩阵去乘它们。
          \begin{equation*}
            \colvec{1 \\ -1 \\ 2}\mapsto\colvec{-1 \\ 4 \\ 5}
            \qquad
            \colvec{2 \\ 0 \\ -1}\mapsto\colvec{3 \\ 5 \\ -5}
            \qquad
            \colvec{1 \\ 1 \\ 0}\mapsto\colvec{1 \\ 4 \\ -1}
          \end{equation*}
          然后计算所得箱体的大小。
          \begin{equation*}
            \begin{vmat}
              -1 &3  &1 \\
               4 &5  &4 \\
               5 &-5 &-1
            \end{vmat}
            =+12
          \end{equation*}
          起始箱体是正定向的，该变换
          保持定向（因为矩阵的行列式为
          正），而终止箱体也是正定向的。
      \end{exparts}
    
\end{ans}
\begin{ans}{四.II.1.14}
      将每个变换关于标准基表出，并求出行列式。
      \begin{exparts*}
        \partsitem $6$
        \partsitem $-1$
        \partsitem $-5$
      \end{exparts*}
    
\end{ans}
\begin{ans}{四.II.1.15}
      起始面积为 \( 6 \)，该矩阵把大小改变
      \( -14 \) 倍。
      因此像的面积为 \( 84 \)。
    
\end{ans}
\begin{ans}{四.II.1.16}
        改变 \( 21/2 \) 倍。
     
\end{ans}
\begin{ans}{四.II.1.17}
      对于箱体，我们取的是一个向量的序列（如注中所说，
      向量的顺序是重要的），
      而对于张成，我们取的是一个向量的集合。
      此外，$\Re^n$ 的箱体子集必须有 $n$ 个向量；
      当然，对于张成，向量的个数可以是任意的。
      最后，对于箱体，系数 $t_1$、~\ldots、$t_n$
      取自区间 $[0..1]$，而对于
      张成，系数可以取遍整个 $\Re$。
    
\end{ans}
\begin{ans}{四.II.1.18}
      两问的答案都是肯定的。
      例如，第一问是 \( \deter{TS}=\deter{T}\cdot\deter{S}=
                     \deter{S}\cdot\deter{T}=\deter{ST} \)。
    
\end{ans}
\begin{ans}{四.II.1.19}
  % due to math.stackexchange.com user dgrasines517
     因为 $\deter{AB}=\deter{A}\cdot\deter{B}=\deter{BA}$，而这两个
     矩阵的行列式不同。
   
\end{ans}
\begin{ans}{四.II.1.20}
      \begin{exparts}
        \partsitem 若它有定义，则它为
           \( (3^2)\cdot (2)\cdot (2^{-2})\cdot (3) \)。
        \partsitem \( \deter{6A^3+5A^2+2A}=\deter{A}\cdot\deter{6A^2+5A+2I} \)。
      \end{exparts}
    
\end{ans}
\begin{ans}{四.II.1.21}
       \(\begin{vmat}
                \cos\theta  &-\sin\theta  \\
                \sin\theta  &\cos\theta
              \end{vmat}=1 \)
    
\end{ans}
\begin{ans}{四.II.1.22}
      不是，例如
      \begin{equation*}
        T=\begin{mat}[r]
          2  &0  \\
          0  &1/2
        \end{mat}
      \end{equation*}
      的行列式为 \( 1 \)，所以它保持面积，但向量 \( T\vec{e}_1 \)
      的长度为 \( 2 \)。
    
\end{ans}
\begin{ans}{四.II.1.23}
       为零。
    
\end{ans}
\begin{ans}{四.II.1.24}
      三角形三条边中有两条由这些向量构成。
      \begin{equation*}
        \colvec[r]{2 \\ 2 \\ 2}-\colvec[r]{1 \\ 2 \\ 1}=\colvec[r]{1 \\ 0 \\ 1}
        \qquad
        \colvec[r]{3 \\ -1 \\ 4}-\colvec[r]{1 \\ 2 \\ 1}=\colvec[r]{2 \\ -3 \\ 3}
      \end{equation*}
      求这个三角形面积的一种方法是：先求由这两个向量与一个
      与它们都正交的单位向量所形成的平行六面体的体积，
      再取其一半。
      为求与那两个向量都正交的向量族，可以从两个关系式
      \begin{equation*}
        \colvec[r]{1 \\ 0 \\ 1}
        \dotprod\colvec{x \\ y \\ z}
        =\colvec[r]{0 \\ 0 \\ 0}
        \qquad
        \colvec[r]{2 \\ -3 \\ 3}
        \dotprod\colvec{x \\ y \\ z}
        =\colvec[r]{0 \\ 0 \\ 0}
      \end{equation*}
      出发，解方程组
      \begin{equation*}
        \begin{linsys}{3}
          x  &   &   &+  &z  &=  &0  \\
         2x  &-  &3y &+  &3z &=  &0
        \end{linsys}
        \grstep{-2\rho_1+\rho_2}
        \begin{linsys}{3}
          x  &   &   &+  &z  &=  &0  \\
             &   &-3y&+  &z  &=  &0
        \end{linsys}
      \end{equation*}
      得到这个解集。
      \begin{equation*}
        \set{\colvec[r]{-1 \\ 1/3 \\ 1}z\suchthat z\in\Re}
      \end{equation*}
      下面是一个长度为一的解。
      \begin{equation*}
        \frac{1}{\sqrt{19/9}}\cdot\colvec[r]{-1 \\ 1/3 \\ 1}
        =\colvec[r]{-3/\sqrt{19} \\ 1/\sqrt{19} \\ 3/\sqrt{19}}
      \end{equation*}
      于是三角形的面积是这个行列式
      绝对值的一半。
      \begin{equation*}
        \begin{vmat}[r]
             1  &2   &-3/\sqrt{19}   \\
             0  &-3  &1/\sqrt{19}   \\
             1  &3   &3/\sqrt{19}
        \end{vmat}
        =-19/\sqrt{19}
      \end{equation*}
      绝对值的一半是 $\sqrt{19}/2$。
    
\end{ans}
\begin{ans}{四.II.1.25}
      \begin{exparts}
        \partsitem 因为线性相关集的像是
          线性相关的，
          所以若构成 $S$ 的向量构成线性相关集，
          即 $\deter{S}=0$，
          则构成 $t(S)$ 的向量构成线性相关集，
          从而 $\deter{TS}=0$，此时等式成立。
        \partsitem 我们必须验证：若
          $T\smash[b]{\grstep{k\rho_i+\rho_j}}\hat{T}$，则
          $d(T)=\deter{TS}/\deter{S}=\deter{\hat{T}S}/\deter{S}=d(\hat{T})$。
          为此可以验证：先把行做倍加再相乘得到 \( \hat{T}S \)，与
          先相乘得到 \( TS \) 再做倍加，
          结果相同（因为行列式 \( \deter{TS} \) 不受行倍加的
          影响，
          这样便有 \( \deter{\hat{T}S}=\deter{TS} \)，从而
          \( d(\hat{T})=d(T) \)）。
          具体验证如下：~把 \( TS \) 的第~\( i \) 行的
          \( k \) 倍加到
          \( TS \) 的第~$j$ 行之后，其 \( j,p \) 元为
          \( (kt_{i,1}+t_{j,1})s_{1,p}+\dots+(kt_{i,r}+t_{j,r})s_{r,p} \)，
          这正是 \( \hat{T}S \) 的 \( j,p \) 元。
        \partsitem 对于第二条性质，只需验证：对换
          $T\smash[b]{\grstep{\rho_i\swap\rho_j}}\hat{T}$
          后再相乘得到 \( \hat{T}S \)，与先把 \( T \) 乘以 \( S \) 再做对换，
          结果相同（因为行列式 \( \deter{TS} \) 在
          行对换下变号，
          这样便有 \( \deter{\hat{T}S}=-\deter{TS} \)，
          从而 \( d(\hat{T})=-d(T) \)）。
          该验证与第一条性质的验证完全类似。

          对于第三条性质，只需证明：先做
          $T\smash[b]{\grstep{k\rho_i}}\hat{T}$
          再计算 \( \hat{T}S \)，与先计算 \( TS \) 再做那次
          倍乘，
          结果相同（由于行列式 \( \deter{TS} \) 被缩放 \( k \) 倍，
          我们将有 \( \deter{\hat{T}S}=k\deter{TS} \)，
          从而 \( d(\hat{T})=k\,d(T) \)）。
          这里的论证同样与上面完全类似。

          第四条性质，即若 $T$ 为 $I$ 则结果为 $1$，
          是显然的。
        \partsitem 行列式函数是唯一的，所以
          \( \deter{TS}/\deter{S}=d(T)=\deter{T} \)，
          从而 $\deter{TS}=\deter{T}\deter{S}$。
      \end{exparts}
    
\end{ans}
\begin{ans}{四.II.1.26}
      任何置换矩阵都具有矩阵的转置等于其逆的性质。

      对于这个蕴含关系，我们知道 \( \deter{\trans{A}}=\deter{A} \)。
      于是 \( 1=\deter{A\cdot A^{-1}}=\deter{A\cdot\trans{A}}
               =\deter{A}\cdot\deter{\trans{A}}=\deter{A}^2 \)。

      逆命题不成立；下面是一个反例。
      \begin{equation*}
        \begin{mat}[r]
          3  &1  \\
          2  &1
        \end{mat}
      \end{equation*}
    
\end{ans}
\begin{ans}{四.II.1.27}
      当箱体的边长变为原来的 \( c \) 倍时，箱体的
      体积就有原来的 \( c^3 \) 倍立方单位。
    
\end{ans}
\begin{ans}{四.II.1.28}
      若 \( H=P^{-1}GP \)，
      则 \( \deter{H}=\deter{P^{-1}}\deter{G}\deter{P}
        =\deter{P^{-1}}\deter{P}\deter{G}=\deter{P^{-1}P}\deter{G}
        =\deter{G} \)。
    
\end{ans}
\begin{ans}{四.II.1.29}
      \begin{exparts}
        \partsitem 新基是旧基旋转 \( \pi/4 \) 得到的。
        \partsitem
          $
             \sequence{\colvec[r]{-1 \\ 0},
                       \colvec[r]{0 \\ -1}}
          $、$
             \sequence{\colvec[r]{0 \\ -1},
                       \colvec[r]{1 \\ 0}}
          $
        \partsitem 每种情形下行列式都是 \( +1 \)
          （我们说这些基
          具有\definend{正定向}）。
        \partsitem 因为一次只能改变一个符号，仅有的另一种
          可能的循环是
          \begin{equation*}
             \cdots
             \;\longrightarrow\colvec{+ \\ +}
             \;\longrightarrow\colvec{+ \\ -}
             \;\longrightarrow\colvec{- \\ -}
             \;\longrightarrow\colvec{- \\ +}
             \;\longrightarrow\cdots\,.
          \end{equation*}
          这里每个相应的行列式都是 \( -1 \)
          （我们说这样的基具有\definend{负定向}）。
        \partsitem 有一个正定向的基 \( \sequence{(1)} \)
          和一个负定向的基 \( \sequence{(-1)} \)。
        \partsitem 共有 \( 48 \) 个基（第一个单位向量有
          \( 6 \) 种半轴可选，
          第二个有 \( 4 \) 种，最后一个有 \( 2 \) 种）。
          其中一半像下面左图中的标准基那样是正定向的，
          另一半像右图中的那样是负定向的
         \begin{center}
           \includegraphics{det/mp/ch4.47}
           \hspace*{4em}
           \includegraphics{det/mp/ch4.48}
          \end{center}
          在 \( \Re^3 \) 中，正定向有时被称为
          “右手定向”，因为如果一个人伸出
          右手，
          手指从 \( \vec{e}_1 \) 卷向 \( \vec{e}_2 \)，
          那么拇指就会指向 \( \vec{e}_3 \)。
      \end{exparts}
    
\end{ans}
\begin{ans}{四.II.1.30}
      我们将比较 \( \det(\vec{s}_1,\dots,\vec{s}_n) \) 与
      \( \det(t(\vec{s}_1),\dots,t(\vec{s}_n)) \)，以证明后者
      与前者相差一个因子 $\deter{T}$。
      我们关于标准基表示那些 \( \vec{s}\, \)：
      \begin{equation*}
        \rep{\vec{s}_i}{\stdbasis_n}=
           \colvec{s_{1,i} \\ s_{2,i} \\ \vdots \\ s_{n,i}}
      \end{equation*}
      然后用矩阵—向量乘法
      表示映射的作用
      \begin{align*}
        \rep{\,t(\vec{s}_i)\,}{\stdbasis_n}
         &=\generalmatrix{t}{n}{n}
           \colvec{s_{1,j} \\ s_{2,j} \\ \vdots \\ s_{n,j}}     \\
         &=s_{1,j}\colvec{t_{1,1} \\ t_{2,1} \\ \vdots \\ t_{n,1}}
          +s_{2,j}\colvec{t_{1,2} \\ t_{2,2} \\ \vdots \\ t_{n,2}}
          +\dots
          +s_{n,j}\colvec{t_{1,n} \\ t_{2,n} \\ \vdots \\ t_{n,n}} \\
         &=s_{1,j}\vec{t}_1+s_{2,j}\vec{t}_2+\dots+s_{n,j}\vec{t}_n
      \end{align*}
      其中 \( \vec{t}_i \) 是 \( T \) 的第~$i$ 列。
      于是 $\det(t(\vec{s}_1),\,\dots,\,t(\vec{s}_n))$ 等于
      $
        \det(s_{1,1}\vec{t}_1\!+\!s_{2,1}\vec{t}_2\!
                +\!\dots\!+\!s_{n,1}\vec{t}_n,\,
             \dots,\,
             s_{1,n}\vec{t}_1\!+\!s_{2,n}\vec{t}_2
                \!+\!\dots\!+\!s_{n,n}\vec{t}_n)
     $。

     如同置换展开公式的推导那样，我们
     利用多重线性，
     先沿第一个变元中的和式拆分
     \begin{multline*}
         \det(s_{1,1}\vec{t}_1,\,
            \dots,\,
            s_{1,n}\vec{t}_1+s_{2,n}\vec{t}_2+\dots+s_{n,n}\vec{t}_n)  \\
        +\cdots{}
        +\det(s_{n,1}\vec{t}_n,\,
           \ldots,\,
            s_{1,n}\vec{t}_1+s_{2,n}\vec{t}_2+\dots+s_{n,n}\vec{t}_n)
     \end{multline*}
     再沿第二个变元中的和式拆分这 $n$ 个和式中的每一个，如此等等。
     与置换展开的推导一样，我们最终得到
     \( n^n \) 个和式行列式，每个形如
     $\det(s_{i_1,1}\vec{t}_{i_1},s_{i_2,2}\vec{t}_{i_2},
            \,\dots,\,
            s_{i_n,n}\vec{t}_{i_n})$。
     把每个 $s_{i,j}$ 提出来，得
     $=
       s_{i_1,1}s_{i_2,2}\dots s_{i_n,n}
        \cdot\det(\vec{t}_{i_1},\vec{t}_{i_2},
       \,\dots,\,
       \vec{t}_{i_n})
      $.

      如同置换展开的推导那样，
      只要 $i_1$、\ldots、$i_n$ 中有两个指标相等，
      该行列式
      就有两个相同的变元，其值为 $0$。
      因此只需考虑 $i_1$、\ldots、$i_n$ 构成数
      $1$、\ldots、$n$ 的一个置换的情形。
      于是我们有
      \begin{equation*}
        \det(t(\vec{s}_1),\dots,t(\vec{s}_n))=
          \sum_{\text{permutations\ } \phi}
            s_{\phi(1),1}\dots s_{\phi(n),n}
            \det(\vec{t}_{\phi(1)},\dots,\vec{t}_{\phi(n)}).
      \end{equation*}
      把 $\det(\vec{t}_{\phi(1)},\ldots,\vec{t}_{\phi(n)})$
      中的列对换，便还原出矩阵 \( T \)，这会按因子
      $\sgn{\phi}$ 改变符号，
      然后再提出 $T$ 的行列式。
      \begin{equation*}
        =\sum_\phi
          s_{\phi(1),1}\dots s_{\phi(n),n}
            \det(\vec{t}_1,\dots,\vec{t}_n)\cdot\sgn{\phi}
        =\det(T)\sum_\phi
          s_{\phi(1),1}\dots s_{\phi(n),n}\cdot\sgn{\phi}.
      \end{equation*}
      如同在证明矩阵的行列式
      等于其转置的
      行列式时那样，我们交换各个 $s$ 的位置，改为按行号递增
      （而不是按列号递增）的顺序列出它们
      （并把 $\sgn(\phi^{-1})$ 代入 $\sgn(\phi)$ 的位置）。
      \begin{equation*}
        =\det(T)\sum_\phi
          s_{1,\phi^{-1}(1)}\dots s_{n,\phi^{-1}(n)}\cdot\sgn{\phi^{-1}}
        =\det(T)\det(\vec{s}_1,\vec{s}_2,\dots,\vec{s}_n)
      \end{equation*}
    
\end{ans}
\begin{ans}{四.II.1.31}
      \begin{exparts}
        \partsitem 代数上的验证是容易的。
        \begin{equation*}
          0
          =xy_2+x_2y_3+x_3y-x_3y_2-xy_3-x_2y
          =x\cdot (y_2-y_3)+y\cdot (x_3-x_2)+x_2y_3-x_3y_2
        \end{equation*}
        化简后得到熟悉的形式
        \begin{equation*}
          y=x\cdot (x_3-x_2)/(y_3-y_2)+(x_2y_3-x_3y_2)/(y_3-y_2)
        \end{equation*}
        （$y_3-y_2=0$ 的情形也容易处理）。

        从几何上看，这幅
        图显示了由这三个向量所形成的箱体。
        注意，所有
        三个向量的终点都在 $z=1$ 平面内。
        右边两个向量的下方是过点
        $(x_2,y_2)$ 与 $(x_3,y_3)$ 的直线。
        \begin{center}
          \includegraphics{det/mp/ch4.49}
        \end{center}
        该箱体
        将有非零的体积，除非这三个向量的终点所构成的三角形
        是退化的。
        这只会在
        （假定 $(x_2,y_3)\neq (x_3,y_3)$ 时）
        $(x,y)$ 位于过另两点的直线上时发生。
       \partsitem \answerasgiven %
        我们按通常的办法，从上式的法向形式出发，
        求以 $(x_1,y_1)$、$(x_2,y_2)$ 与 $(x_3,y_3)$ 为顶点的
        三角形过 $(x_1,y_1)$ 的高：
        \begin{equation*}
          \frac{1}{\sqrt{(x_2-x_3)^2+(y_2-y_3)^2}}
          \begin{vmat}
            x_1  &x_2  &x_3  \\
            y_1  &y_2  &y_3  \\
            1    &1    &1
          \end{vmat}.
        \end{equation*}
        再走一步就得到三角形的面积为
        \begin{equation*}
          \frac{1}{2}
          \begin{vmat}
            x_1  &x_2  &x_3  \\
            y_1  &y_2  &y_3  \\
            1    &1    &1
          \end{vmat}.
        \end{equation*}
        与通常那种证明一列项与行列式恒等的办法相比，
        这里的阐述更清楚地揭示了\textit{modus operandi}（做法）。
       \partsitem  \answerasgiven %
        设
        \begin{equation*}
          D=
          \begin{vmat}
            x_1  &x_2  &x_3  \\
            y_1  &y_2  &y_3  \\
            1    &1    &1
          \end{vmat}
        \end{equation*}
        则三角形的面积为 $(1/2)\deter{D}$。
        于是，若坐标全为整数，则 $D$ 是一个整数。
      \end{exparts}
    
\end{ans}
\protect \section {第 III 节：Laplace 公式}
\protect \subsection {四.III.1: Laplace 展开}
\begin{ans}{四.III.1.10}
      \begin{exparts*}
       \partsitem \( (-1)^{2+3}\begin{vmat}[r]
                            1  &0  \\
                            0  &2
                          \end{vmat}=-2  \)
       \partsitem \( (-1)^{3+2}\begin{vmat}[r]
                            1  &2  \\
                           -1  &3
                          \end{vmat}=-5  \)
       \partsitem \( (-1)^{4}\begin{vmat}[r]
                           -1  &1  \\
                            0  &2
                          \end{vmat}=-2  \)
      \end{exparts*}
    
\end{ans}
\begin{ans}{四.III.1.11}
      % sage: M = matrix(QQ, [[1,0,2], [-1,1,3], [0,2,-1]])
      % sage: M.adjoint()
      % [-7  4 -2]
      % [-1 -1 -5]
      % [-2 -2  1]
      {\renewcommand{\arraystretch}{1.2}
      \begin{equation*}
        \adj(T)=
        \begin{mat}
          \begin{vmatrix}
            1 &3 \\
            2 &-1
          \end{vmatrix}
          &-\begin{vmatrix}
            0 &2 \\
            2 &-1
          \end{vmatrix}
          &\begin{vmatrix}
             0 &2 \\
              1 &3
           \end{vmatrix}                  \\
          -\begin{vmatrix}
             -1 &3 \\
             0 &-1
           \end{vmatrix}
          &\begin{vmatrix}
            1 &2 \\
            0 &-1
           \end{vmatrix}
          &-\begin{vmatrix}
            1 &2 \\
            -1 &3
           \end{vmatrix}                  \\
          \begin{vmatrix}
            -1 &1 \\
            0 &2
          \end{vmatrix}
          &-\begin{vmatrix}
            1  &0 \\
            0 &2
           \end{vmatrix}
          &\begin{vmatrix}
            1 &0 \\
            -1 &1
           \end{vmatrix}
        \end{mat}
        =
        \begin{mat}
          -7 &4 &-2 \\
          -1  &-1 &-5 \\
          -2 &-2 &1
        \end{mat}
      \end{equation*}}
    
\end{ans}
\begin{ans}{四.III.1.12}
      \begin{equation*}
        \begin{vmat}[r]
          1  &2  &3  \\
          4  &5  &6  \\
          7  &8  &9
        \end{vmat}=
        1\cdot(+1)\cdot
          \begin{vmatrix}
            5 &6 \\
            8 &9
          \end{vmatrix}
        +2\cdot(-1)\cdot
          \begin{vmatrix}
            4 &6 \\
            7 &9
          \end{vmatrix}
        +3\cdot(+1)\cdot
          \begin{vmatrix}
            4 &5 \\
            7 &8
          \end{vmatrix}
      \end{equation*}
      $\nbyn{2}$ 矩阵的公式给出
      $1\cdot (-3)-2\cdot (-6)+3\cdot(-3)=0$。
    
\end{ans}
\begin{ans}{四.III.1.13}
      \begin{exparts}
         \partsitem \( 3\cdot(+1)\begin{vmat}[r]
                          2  &2  \\
                          3  &0
                       \end{vmat}
                 +0\cdot (-1)\begin{vmat}[r]
                          1  &2  \\
                         -1  &0
                       \end{vmat}
                 +1\cdot (+1)\begin{vmat}[r]
                          1  &2  \\
                         -1  &3
                       \end{vmat} =-13 \)
         \partsitem \( 1\cdot (-1)\begin{vmat}[r]
                          0  &1  \\
                          3  &0
                       \end{vmat}
                 +2\cdot (+1)\begin{vmat}[r]
                          3  &1  \\
                         -1  &0
                       \end{vmat}
                 +2\cdot (-1)\begin{vmat}[r]
                          3  &0  \\
                         -1  &3
                       \end{vmat} =-13 \)
         \partsitem \( 1\cdot (+1)\begin{vmat}[r]
                          1  &2  \\
                         -1  &3
                       \end{vmat}
                 +2\cdot (-1)\begin{vmat}[r]
                          3  &0  \\
                         -1  &3
                       \end{vmat}
                 +0\cdot (+1)\begin{vmat}[r]
                          3  &0  \\
                          1  &2
                       \end{vmat} =-13 \)
       \end{exparts}
     
\end{ans}
\begin{ans}{四.III.1.14}
      这就是 $\adj(T)$。
      \begin{align*}
         \begin{mat}
           T_{1,1}  &T_{2,1}  &T_{3,1} \\
           T_{1,2}  &T_{2,2}  &T_{3,2} \\
           T_{1,3}  &T_{2,3}  &T_{3,3}
         \end{mat}
         &=\begin{mat}
            \begin{vmat}
               5  &6  \\ 8  &9
            \end{vmat}
            &-\begin{vmat}
               2  &3  \\  8  &9
            \end{vmat}
            &+\begin{vmat}
               2  &3  \\  5  &6
            \end{vmat}        \\[2.5ex]
            -\begin{vmat}
               4  &6  \\  7  &9
            \end{vmat}
            &+\begin{vmat}
               1  &3  \\  7  &9
            \end{vmat}
            &-\begin{vmat}
               1  &3  \\  4  &6
            \end{vmat}        \\[2.5ex]
            +\begin{vmat}
               4  &5  \\  7  &8
            \end{vmat}
            &-\begin{vmat}
               1  &2  \\  7  &8
            \end{vmat}
            &+\begin{vmat}
               1  &2  \\  4  &5
            \end{vmat}
         \end{mat}                                   \\
         &=\begin{mat}[r]
          -3  &6   &-3 \\
           6  &-12 &6  \\
          -3  &6   &-3
         \end{mat}
      \end{align*}
    
\end{ans}
\begin{ans}{四.III.1.15}
      \begin{exparts}
        \partsitem 这就是伴随矩阵。
          \begin{align*}
          \begin{mat}
            T_{1,1}  &T_{2,1}  &T_{3,1} \\
            T_{1,2}  &T_{2,2}  &T_{3,2} \\
            T_{1,3}  &T_{2,3}  &T_{3,3}
          \end{mat}
          &=\begin{mat}
            \begin{vmat}[r]
              0  &2  \\
              0  &1
            \end{vmat}
            &-\begin{vmat}[r]
               1  &4  \\
               0  &1
            \end{vmat}
            &\begin{vmat}[r]
               1  &4  \\
               0  &2
            \end{vmat}        \\[2.5ex]
            -\begin{vmat}[r]
              -1  &2  \\
               1  &1
            \end{vmat}
            &\begin{vmat}[r]
               2  &4  \\
               1  &1
            \end{vmat}
            &-\begin{vmat}[r]
               2  &4  \\
              -1  &2
            \end{vmat}         \\[2.5ex]
            \begin{vmat}[r]
              -1  &0  \\
               1  &0
            \end{vmat}
            &-\begin{vmat}[r]
               2  &1  \\
               1  &0
            \end{vmat}
            &\begin{vmat}[r]
               2  &1  \\
              -1  &0
            \end{vmat}
          \end{mat}                         \\
          &=\begin{mat}[r]
             0  &-1  &2  \\
             3  &-2  &-8 \\
             0  &1   &1
          \end{mat}
          \end{align*}
        \partsitem 子式是 $\nbyn{1}$ 的。
          $\begin{mat}
             T_{1,1}  &T_{2,1} \\
             T_{1,2}  &T_{2,2}
          \end{mat}
          =
          \begin{mat}
            \begin{vmat}
              4
            \end{vmat}
            &-\begin{vmat}
              -1
            \end{vmat}        \\[1.5ex]
            -\begin{vmat}
               2
            \end{vmat}
            &\begin{vmat}
               3
            \end{vmat}
          \end{mat}
          =
          \begin{mat}[r]
            4  &1  \\
           -2  &3
          \end{mat}$
        \partsitem
          $\begin{mat}[r]
               0  &-1 \\
              -5  &1
          \end{mat}$
        \partsitem 子式是 $\nbyn{2}$ 的。
          \begin{align*}
          \begin{mat}
            T_{1,1}  &T_{2,1}  &T_{3,1} \\
            T_{1,2}  &T_{2,2}  &T_{3,2} \\
            T_{1,3}  &T_{2,3}  &T_{3,3}
          \end{mat}
          &=\begin{mat}
            \begin{vmat}[r]
              0  &3  \\
              8  &9
            \end{vmat}
            &-\begin{vmat}[r]
              4  &3  \\
              8  &9
            \end{vmat}
            &\begin{vmat}[r]
              4  &3  \\
              0  &3
            \end{vmat}       \\[2.5ex]
            -\begin{vmat}[r]
              -1  &3  \\
               1  &9
            \end{vmat}
            &\begin{vmat}[r]
               1  &3  \\
               1  &9
            \end{vmat}
            &-\begin{vmat}[r]
               1  &3 \\
              -1  &3
            \end{vmat}        \\[2.5ex]
            \begin{vmat}[r]
               -1  &0  \\
                1  &8
            \end{vmat}
            &-\begin{vmat}[r]
                1  &4 \\
                1  &8
            \end{vmat}
            &\begin{vmat}[r]
                1  &4  \\
               -1  &0
            \end{vmat}
          \end{mat}                       \\
          &=
          \begin{mat}[r]
            -24  &-12  &12  \\
             12  &6    &-6   \\
             -8  &-4   &4
          \end{mat}
        \end{align*}
      \end{exparts}
    
\end{ans}
\begin{ans}{四.III.1.16}
      \begin{exparts}
        \partsitem $(1/3)\cdot
           \begin{mat}[r]
             0  &-1  &2  \\
             3  &-2  &-8 \\
             0  &1   &1
          \end{mat}
          =
          \begin{mat}[r]
            0  &-1/3  &2/3  \\
            1  &-2/3  &-8/3 \\
            0  &1/3   &1/3
          \end{mat}$
        \partsitem $(1/14)\cdot
          \begin{mat}[r]
            4  &1  \\
           -2  &3
          \end{mat}
          =
          \begin{mat}[r]
            2/7  &1/14  \\
           -1/7  &3/14
          \end{mat}$
        \partsitem $(1/-5)\cdot
          \begin{mat}[r]
               0  &-1 \\
              -5  &1
          \end{mat}
          =
          \begin{mat}[r]
            0  &1/5 \\
            1  &-1/5
          \end{mat}$
        \partsitem 该矩阵的行列式为零，因此没有逆矩阵。
      \end{exparts}
    
\end{ans}
\begin{ans}{四.III.1.17}
      $\begin{mat}
        T_{1,1}  &T_{2,1}  &T_{3,1}  &T_{4,1}  \\
        T_{1,2}  &T_{2,2}  &T_{3,2}  &T_{4,2}  \\
        T_{1,3}  &T_{2,3}  &T_{3,3}  &T_{4,3}  \\
        T_{1,4}  &T_{2,4}  &T_{3,4}  &T_{4,4}
      \end{mat}
      =
      \begin{mat}[r]
        4  &-3  &2  &-1  \\
       -3  &6   &-4 &2   \\
        2  &-4  &6  &-3  \\
        -1 &2  &-3  &4
      \end{mat}$
    
\end{ans}
\begin{ans}{四.III.1.18}
       行列式
       \begin{equation*}
         \begin{vmat}
           a  &b  \\
           c  &d
         \end{vmat}
       \end{equation*}
       按第一行展开给出
       \( a\cdot (+1)\deter{d}+b\cdot (-1)\deter{c}=ad-bc \)
       （注意那两个 $\nbyn{1}$ 子式）。
     
\end{ans}
\begin{ans}{四.III.1.19}
       矩阵
       \begin{equation*}
         \begin{mat}
           a  &b  &c  \\
           d  &e  &f  \\
           g  &h  &i
         \end{mat}
       \end{equation*}
       的行列式如下。
       \begin{equation*}
         a\cdot \begin{vmat}
             e  &f  \\
             h  &i
          \end{vmat}
         -b\cdot \begin{vmat}
             d  &f  \\
             g  &i
          \end{vmat}
         +c\cdot \begin{vmat}
             d  &e  \\
             g  &h
          \end{vmat}
         =a(ei-fh)-b(di-fg)+c(dh-eg)
       \end{equation*}
    
\end{ans}
\begin{ans}{四.III.1.20}
      \begin{exparts}
        \partsitem
          $\begin{mat}
             T_{1,1}  &T_{2,1}  \\
             T_{1,2}  &T_{2,2}
            \end{mat}
            =\begin{mat}
            \begin{vmat}
              t_{2,2}
            \end{vmat}
           &-\begin{vmat}
               t_{1,2}
             \end{vmat}       \\
           -\begin{vmat}
               t_{2,1}
            \end{vmat}
           &\begin{vmat}
              t_{1,1}
            \end{vmat}
          \end{mat}
          =\begin{mat}
             t_{2,2}  &-t_{1,2}  \\
            -t_{2,1} &t_{1,1}
           \end{mat}$
        \partsitem
          $(1/t_{1,1}t_{2,2}-t_{1,2}t_{2,1})\cdot
          \begin{mat}
             t_{2,2}  &-t_{1,2}  \\
            -t_{2,1} &t_{1,1}
           \end{mat}$
      \end{exparts}
    
\end{ans}
\begin{ans}{四.III.1.21}
      不能。
      下面这个行列式的值
      \begin{equation*}
        \begin{vmat}[r]
          1  &0  &0  \\
          0  &1  &0  \\
          0  &0  &1
        \end{vmat}=1
      \end{equation*}
      并不等于沿对角线展开的结果。
      \begin{equation*}
        1\cdot (+1)\begin{vmat}[r]
               1  &0  \\
               0  &1
             \end{vmat}
       +1\cdot (+1)\begin{vmat}[r]
               1  &0  \\
               0  &1
             \end{vmat}
       +1\cdot (+1)\begin{vmat}[r]
               1  &0  \\
               0  &1
             \end{vmat}=3
      \end{equation*}
    
\end{ans}
\begin{ans}{四.III.1.22}
      考虑这个对角矩阵。
      \begin{equation*}
        D=
        \begin{mat}
          d_1  &0   &0   &\ldots    \\
          0    &d_2 &0   &          \\
          0    &0   &d_3            \\
               &    &    &\ddots    \\
               &    &    &      &d_n
        \end{mat}
      \end{equation*}
      若 $i\neq j$，那么 $i,j$~子式是一个 $\nbyn{(n-1)}$ 矩阵，
      其中只有 $n-2$ 个非零元，因为我们同时删去了 $d_i$ 与 $d_j$。
      于是，该子式至少有一行或一列全为零，所以代数余子式 $D_{i,j}$ 为零。
      若 $i=j$，那么该子式是对角元为
      $d_1$, \ldots, $d_{i-1}$, $d_{i+1}$, \ldots, $d_n$ 的对角矩阵。
      它的行列式显然等于 $(-1)^{i+j}=(-1)^{2i}=1$
      乘以这些元的乘积。
      \begin{equation*}
        \adj(D)
        =
        \begin{mat}
          d_2\cdots d_n    &0                 &      &0    \\
          0                &d_1d_3\cdots d_n  &      &0    \\
                           &                  &\ddots      \\
                           &                  &       &d_1\cdots d_{n-1}
        \end{mat}
      \end{equation*}

      顺便说一下，\nearbytheorem{th:MatTimesAdjEqDiagDets} 提供了
      一种更简洁的方法来推导这个结论。
    
\end{ans}
\begin{ans}{四.III.1.23}
      只需注意，若 $S=\trans{T}$，则代数余子式
      $S_{j,i}$ 等于代数余子式 $T_{i,j}$，这是因为 $(-1)^{j+i}=(-1)^{i+j}$，
      并且两个子式互为转置（而转置的行列式等于原矩阵的行列式）。
    
\end{ans}
\begin{ans}{四.III.1.24}
      该命题不成立；这里有一个例子。
      \begin{equation*}
         T=\begin{mat}[r]
           1  &2   &3  \\
           4  &5   &6  \\
           7  &8   &9
         \end{mat}
         \qquad
         \adj(T)=\begin{mat}[r]
          -3  &6   &-3 \\
           6  &-12 &6  \\
          -3  &6   &-3
         \end{mat}
         \qquad
         \adj(\adj(T))=\begin{mat}[r]
           0  &0   &0  \\
           0  &0   &0  \\
           0  &0   &0
         \end{mat}
      \end{equation*}
    
\end{ans}
\begin{ans}{四.III.1.25}
      \begin{exparts}
        \partsitem 例子
          \begin{equation*}
             M=
            \begin{mat}[r]
              1  &2  &3  \\
              0  &4  &5  \\
              0  &0  &6
            \end{mat}
          \end{equation*}
          提示了正确的答案。
          \begin{align*}
            \adj(M)=
            \begin{mat}
             M_{1,1}  &M_{2,1}  &M_{3,1}  \\
             M_{1,2}  &M_{2,2}  &M_{3,2}  \\
             M_{1,3}  &M_{2,3}  &M_{3,3}
            \end{mat}
            &=\begin{mat}
              \begin{vmat}
                4  &5 \\ 0 &6
              \end{vmat}
              &-\begin{vmat}
                2  &3  \\  0  &6
              \end{vmat}
              &\begin{vmat}
                2  &3  \\  4  &5
              \end{vmat}          \\[1.5ex]
              -\begin{vmat}
                0  &5  \\  0  &6
              \end{vmat}
              &\begin{vmat}
                1  &3  \\  0  &6
              \end{vmat}
              &-\begin{vmat}
                1  &3  \\  0  &5
              \end{vmat}           \\[1.5ex]
              \begin{vmat}
                0  &4  \\  0  &0
              \end{vmat}
              &-\begin{vmat}
                1  &2  \\  0  &0
              \end{vmat}
              &\begin{vmat}
                1  &2  \\  0  &4
              \end{vmat}
            \end{mat}                                \\
            &=
            \begin{mat}[r]
              24  &-12 &-2 \\
               0  &6   &-5  \\
               0  &0   &4
            \end{mat}
          \end{align*}
          结果确实是上三角的。

          这项检验是细致的，但并不难。
          伴随矩阵上三角部分的元是
          $M_{a,b}$，其中 $a>b$。
          我们需要验证：若 $a>b$，则代数余子式 $M_{a,b}$ 为零。
          当 $a>b$ 时，把 $M$ 的第~$a$ 行与第~$b$ 列
          \begin{equation*}
            \begin{mat}
              m_{1,1} &\ldots   &m_{1,b}  &\ldots       &        \\
              m_{2,1} &\ldots   &m_{2,b}  &       &        \\
              \vdots  &         &\vdots   &       &        \\
              m_{a,1} &\ldots   &m_{a,b}  &\ldots &m_{a,n} \\
              \vdots        &         &\vdots   &       &        \\
                      &         &m_{n,b}  &       &
            \end{mat}
          \end{equation*}
          删去后，得到一个上三角的子式，
          这是因为该子式的第~$i,j$ 元，要么是 $M$ 的第~$i,j$ 元
          （当 $a>i$ 且~$b>j$ 时发生这种情形；
          此时 $i<j$ 蕴含该元为零），
          要么是 $M$ 的第~$i,j+1$ 元（当 $i<a$ 且
          $j>b$ 时发生这种情形；此时 $i<j$ 蕴含 $i<j+1$，从而
          该元为零），要么是 $M$ 的第~$i+1,j+1$ 元
          （最后这种情形在 $i>a$ 且~$j>b$ 时发生；显然这里
          $i<j$ 蕴含 $i+1<j+1$，所以该元为零）。
          于是该子式的行列式就是沿对角线的乘积。
          注意，$M$ 的 $a-1,a$ 元就是该子式的
          $a-1,a-1$~元（由于关系 $a>b$ 是严格的，所以它
          不会被删去）。
          但这个元为零，因为 $M$ 是上三角的且
          $a-1<a$。
          因此该代数余子式为零，伴随矩阵是上三角的。
          （下三角的情形类似。）
        \partsitem 利用\nearbytheorem{th:MatTimesAdjEqDiagDets}，
          这可由前一小问立即得出。
      \end{exparts}
    
\end{ans}
\begin{ans}{四.III.1.26}
      我们将证明每个行列式都可以沿第~\( i \) 行展开。
      对第~$j$ 列的论证是类似的。

      置换展开中的每一项恰好含有来自每一行的一个元。
      如\nearbyexample{ex:ExpThreeFirstRow}那样，
      把每个第~$i$ 行的元提出来，得到
      \( \deter{T}
           =t_{i,1}\cdot\hat{T}_{i,1}+\dots+t_{i,n}\cdot\hat{T}_{i,n} \)，
      其中每个 \( \hat{T}_{i,j} \) 是一些不含第 \( i \) 行任何元素的项之和。
      我们将证明 \( \hat{T}_{i,j} \) 就是 \( i,j \) 代数余子式。

      先考虑 \( i,j=n,n \) 的情形：
      \begin{equation*}
        t_{n,n}\cdot\hat{T}_{n,n}
        =t_{n,n}\cdot
            \sum_{\phi}t_{1,\phi(1)}t_{2,\phi(2)}\dots\,t_{n-1,\phi(n-1)}
                      \sgn(\phi)
      \end{equation*}
      其中求和取遍所有满足
      \( \phi(n)=n \) 的 \( n \)-置换 \( \phi \)。
      为了证明
      \( \hat{T}_{i,j} \) 是子式 \( T_{i,j} \)，我们只需证明：
      若 \( \phi \) 是满足
      \( \phi(n)=n \) 的 \( n \)-置换，
      且 \( \sigma \) 是满足
      \( \sigma(1)=\phi(1) \), \ldots, \( \sigma(n-1)=\phi(n-1) \)
      的 \( n-1 \)-置换，
      则 \( \sgn(\sigma)=\sgn(\phi) \)。
      但这是成立的，因为 $\phi$ 与 $\sigma$
      有相同数目的逆序。

      回到一般的 \( i,j \) 情形。
      对换相邻的行，直到第 \( i \) 行变为最后一行；再对换相邻的列，
      直到第 \( j \) 列变为最后一列。
      注意，第 \( i,j \) 个子式的行列式不受这些相邻
      对换的影响，因为逆序被保持（由于该子式
      略去了第 \( i \) 行与第 \( j \) 列）。
      另一方面，\( \deter{T} \) 与 \( \hat{T}_{i,j} \) 的符号
      改变了 \( n-i \) 加 \( n-j \) 次。
      于是 \( \hat{T}_{i,j}=(-1)^{n-i+n-j}\deter{T_{i,j}}
                =(-1)^{i+j}\deter{T_{i,j}} \)。
    
\end{ans}
\begin{ans}{四.III.1.27}
      对 \( \nbyn{1} \) 的基准情形，这是显然的。

      归纳的情形：
      假设对所有 \( \nbyn{1} \), \ldots, \( \nbyn{(n-1)} \)
      矩阵，矩阵的行列式都等于其转置的行列式。
      沿第~\( i \) 行展开给出
      \( \deter{T}=t_{i,1}T_{i,1}+\dots\,+t_{i,n}T_{i,n} \)，
      而沿第~\( i \) 列展开给出
      \( \deter{\trans{T}}=t_{1,i}(\trans{T})_{1,i}
           +\dots+t_{n,i}(\trans{T})_{n,i} \)。
      由于 \( (-1)^{i+j}=(-1)^{j+i} \)，两个和式中的符号是相同的。
      由于 \( \trans{T} \) 的 \( j,i \) 子式是 \( T \) 的 \( i,j \) 子式的转置，
      归纳假设给出 \( \deter{(\trans{T})_{i,j}}=\deter{T_{i,j}} \)。
    
\end{ans}
\begin{ans}{四.III.1.28}
      \answerasgiven %
      把上面的行列式记作 \( D_n \)，可以看到
      \( D_2=1 \)，\( D_3=2 \)。
      剩下要证明的是 \( D_n=D_{n-1}+D_{n-2},\; n\geq 4 \)。
      在 \( D_n \) 中，从第 \( (n-1) \) 列减去第 \( (n-3) \) 列，
      从第 \( (n-2) \) 列减去第 \( (n-4) \) 列，\ldots，从第三列减去
      第一列，得到
      \begin{equation*}
        F_n=
        \begin{vmat}
          1  &-1  &0  &0   &0  &0   &\ldots  \\
          1  &1   &-1 &0   &0  &0   &\ldots  \\
          0  &1   &1  &-1  &0  &0   &\ldots  \\
          0  &0   &1  &1   &-1 &0   &\ldots  \\
          .  &.   &.  &.   &.  &.   &\ldots
        \end{vmat}.
      \end{equation*}
      按第一行展开这个行列式，即得所要的关系式。
    
\end{ans}
\protect \topic {Cramer 法则}
\begin{ans}{1}
      \begin{exparts*}
        \partsitem 分别解出各个变量。
          \begin{equation*}
            x=
             \frac{ \begin{vmat}[r]
                       4  &-1  \\
                      -7  &2
                    \end{vmat}  }{
                    \begin{vmat}[r]
                       1  &-1  \\
                      -1  &2
                    \end{vmat}  }
            =\frac{1}{1}=1
            \qquad
            y=
             \frac{ \begin{vmat}[r]
                       1  &4  \\
                      -1  &-7
                    \end{vmat}  }{
                    \begin{vmat}[r]
                       1  &-1  \\
                      -1  &2
                    \end{vmat}  }
            =\frac{-3}{1}=-3
          \end{equation*}
        \partsitem $x=2$，$y=2$
      \end{exparts*}
    
\end{ans}
\begin{ans}{2}
      \( z=1 \)
    
\end{ans}
\begin{ans}{3}
      行列式在倍加下不变，
      包括列的倍加，于是
      \( \det(B_i)=\det(\vec{a}_1,\dots,
      x_1\vec{a}_1+\dots+x_i\vec{a}_i+\dots+x_n\vec{a}_n,\dots,\vec{a}_n) \)。
      用“取第一列的 $-x_1$ 倍加到第~$i$ 列”等操作，可以看出它
      等于
      $\det(\vec{a}_1,\dots,x_i\vec{a}_i,\dots,\vec{a}_n)$。
      进而，它又等于
      $x_i\cdot\det(\vec{a}_1,\dots,\vec{a}_i,\dots,\vec{a}_n)
         =x_i\cdot\det(A)$，
      这正是所要证的。
    
\end{ans}
\begin{ans}{4}
     \begin{exparts*}
       \partsitem
         下面是 $\nbyn{2}$ 方程组取 $i=2$ 的情形。
         \begin{equation*}
           \begin{linsys}{2}
             a_{1,1}x_1 &+ &a_{1,2}x_2 &= &b_1 \\
             a_{2,1}x_1 &+ &a_{2,2}x_2 &= &b_2
           \end{linsys}
           \quad\Longleftrightarrow\quad
           \begin{mat}
             a_{1,1}  &a_{1,2}  \\
             a_{2,1}  &a_{2,2}
           \end{mat}
           \begin{mat}
             1  &x_1 \\
             0  &x_2
           \end{mat}
           =
           \begin{mat}
             a_{1,1}  &b_1 \\
             a_{2,1}  &b_2
           \end{mat}
         \end{equation*}
      \partsitem 行列式函数是乘性的：
        $\det(B_i)=\det(AX_i)=\det(A)\cdot\det(X_i)$。
        Laplace 展开表明 $\det(X_i)=x_i$，
        解出 $x_i$ 便得到 Cramer 法则。
     \end{exparts*}
   
\end{ans}
\begin{ans}{5}
      由于 $A$ 的行列式非零，Cramer 法则适用，并给出 $x_i=\deter{B_i}/1$。
      由于 $B_i$ 是整数矩阵，故其行列式是整数。
    
\end{ans}
\begin{ans}{6}
      方程组
      \begin{equation*}
        \begin{linsys}{2}
           ax  &+  by  &=  &e  \\
           cx  &+  dy  &=  &f
        \end{linsys}
      \end{equation*}
      的解为
      \begin{equation*}
        x=\frac{ed-fb}{ad-bc}
        \qquad
        y=\frac{af-ec}{ad-bc}
      \end{equation*}
      当然前提是分母不为零。
    
\end{ans}
\begin{ans}{7}
      当然，奇异方程组的 \( \deter{A} \) 等于零；
      而无穷多解的情形可以用
      所有的 \( \deter{B_i} \) 也都为零这一事实来刻画。
    
\end{ans}
\begin{ans}{8}
      我们可以把两种非奇异的情形与下面这个
      方程组一起考虑
      \begin{equation*}
        \begin{linsys}{2}
           x_1  &+  &2x_2  &=  &6  \\
           x_1  &+  &2x_2  &=  &c
        \end{linsys}
      \end{equation*}
      其中 $c=6$ 当然给出无穷多解，而 $c$ 取任何其他值
      都给出无解。
      相应的向量方程
      \begin{equation*}
        x_1\cdot\colvec[r]{1 \\ 1}+x_2\cdot\colvec[r]{2 \\ 2}=\colvec[r]{6 \\ c}
      \end{equation*}
      给出两条重叠向量的图景。
      它们都位于直线 $y=x$ 上。
      在 $c=6$ 的情形，右边的向量也位于
      直线 $y=x$ 上，而在其他任何情形下都不是。
    
\end{ans}
\protect \topic {行列式计算的速度}
\begin{ans}{1}
      你的计时结果一部分取决于所用的计算机代数系统，
      一部分取决于你在其上进行计算的计算机的能力。
      但你应得到一条与所示曲线相似的曲线。
%       \begin{exparts}
%         \partsitem Under Octave, \texttt{rank(rand(5))} finds the
%           rank of a $\nbyn{5}$ matrix whose entries are (uniformly
%           distributed) in the interval $[0..1)$.
%           This loop which runs the test $5000$ times
% \begin{lstlisting}
% octave:1> for i=1:5000
% > if rank(rand(5))<5 printf("That's one."); endif
% > endfor
% \end{lstlisting}
%           produces (after a few seconds) returns the prompt, with no output.

%           The Octave script
% \begin{lstlisting}
% function elapsed_time = detspeed (size)
%   a=rand(size);
%   tic();
%   for i=1:10
%      det(a);
%   endfor
%   elapsed_time=toc();
% endfunction
% \end{lstlisting}
%           lead to this session (obviously, your times will vary).
% \begin{lstlisting}
% octave:1> detspeed(5)
% ans = 0.019505
% octave:2> detspeed(15)
% ans = 0.0054691
% octave:3> detspeed(25)
% ans = 0.0097431
% octave:4> detspeed(35)
% ans = 0.017398
% \end{lstlisting}
%           \partsitem Here is the data (rounded a bit), and the graph.
%             \begin{center}
%               \begin{tabular}{l}
%               \begin{tabular}{r|cccccc}
%                  \textit{matrix rows}
%                     &$15$ &$25$ &$35$ &$45$ &$55$ &$65$  \\
%                  \hline
%                  \textit{time per ten}
%                     &$0.003\,4$
%                     &$0.009\,8$
%                     &$0.067\,5$
%                     &$0.028\,5$
%                     &$0.044\,3$
%                     &$0.066\,3$
%               \end{tabular}                                       \\
%               \qquad\begin{tabular}{r|ccc}
%                  \textit{matrix rows}
%                     &$75$ &$85$ &$95$ \\
%                  \hline
%                  \textit{time per ten}
%                     &$0.142\,8$
%                     &$0.228\,2$
%                     &$0.168\,6$
%               \end{tabular}
%               \end{tabular}
%             \end{center}
%           This data is from an average of twenty runs of the above script,
%           because of the possibility that the randomly chosen matrix
%           happens to take an unusually long or short time.
%           Even so, the timing cannot be relied on too heavily; this is
%           just an experiment.
%           \begin{center}
%             \includegraphics[width=0.4\textwidth]{det/mp/ch4.28}
%           \end{center}
%       \end{exparts}
    
\end{ans}
\begin{ans}{2}
       操作次数取决于我们具体如何执行这些操作。
       \begin{exparts}
         \partsitem 行列式为 $-11$。
           行化简只需一次倍加，
           含两次乘法
           （$-5/2$ 乘 $2$ 加 $5$，以及 $-5/2$ 乘 $1$ 加 $-3$），
           沿对角线取乘积还需一次乘法。
           置换展开需要两次乘法（$2$ 乘
           $-3$ 与 $5$ 乘 $1$）。
         \partsitem 行列式为 $-39$。
           清点操作次数是例行的。
         \partsitem 行列式为 $4$。
       \end{exparts}
     
\end{ans}
\begin{ans}{3}
      你应当只构造矩阵一次，然后对它被多次化简
      进行计时。
      注意，如上所写的 \lstinline[style=inline]!gauss_method!
      会改变输入矩阵，因此要小心，不要用
      该例程的输出做循环，否则除第一次循环外，
      其余各次循环都将在一个阶梯形矩阵上进行。
    
\end{ans}
\begin{ans}{4}
      单位矩阵通常化简起来用不了多久。
      对于化简起来很慢的矩阵，在 Python 中你可以试试元素
      相当大的矩阵。
      （一些计算机代数系统还具有生成
      特殊矩阵的能力；可以试试那些。）
      % One way to get started is to compare these under Octave:
      % \texttt{det(rand(10));}, versus
      % \texttt{det(hilb(10));}, versus
      % \texttt{det(eye(10));}, versus
      % \texttt{det(zeroes(10));}.
      % You can time them as in \texttt{tic(); det(rand(10)); toc()}.
    
\end{ans}
\begin{ans}{5}
      没有，上面的代码是按行处理这些数的。
    
\end{ans}
\protect \topic {Chiò 方法}
\begin{ans}{1}
      \begin{exparts*}
        \partsitem Chiò 矩阵是
        \begin{equation*}
          C=
          \begin{mat}
            -3 &-6 \\
            -6 &-12
          \end{mat}
        \end{equation*}
        其行列式为 $0$
        \partsitem 从
        \begin{equation*}
          C_3=
          \begin{mat}
            2 &8  &0 \\
            1 &-2 &2 \\
            4 &2  &2
          \end{mat}
        \end{equation*}
        开始，再下一步得
        \begin{equation*}
          C_2=
          \begin{mat}
            -12 &4 \\
            -28 &4
          \end{mat}
        \end{equation*}
        其行列式 $\det(C_2)=64$。
        故原矩阵的行列式为 $64/(2^2\cdot 2^1)=8$
      \end{exparts*}
    
\end{ans}
\begin{ans}{2}
      与上面 $\nbyn{3}$~情形相同的构造表明，
      我们可以选取任一非零元素~$a_{i,j}$ 来代替 $a_{1,1}$。
      Chiò 矩阵的元素 $c_{p,q}$ 是行列式
      \begin{equation*}
        \begin{vmat}
          a_{1,1}    &a_{1,q+1} \\
          a_{p+1,1}  &a_{p+1,q+1}
        \end{vmat}
      \end{equation*}
      的值，其中 $p+1\neq i$ 且 $q+1\neq j$。
    
\end{ans}
\begin{ans}{3}
     Sarrus 公式用 $12$~次乘法和 $5$ 次加法（减法计入加法之中）。
     Chiò 公式对四个 $\nbyn{2}$~行列式中的每一个要用两次乘法和一次加法
     （实际上是减法），
     对 $\nbyn{2}$~Chiò 行列式还要再用两次乘法和一次加法，
     最后还要除以 $a_{1,1}$。
     总共是十一次乘除和五次加减。
     所以 Chiò 方法胜出。
   
\end{ans}
\begin{ans}{4}
      考虑 $\nbyn{n}$ 矩阵
      \begin{equation*}
        A=
        \begin{mat}
           a_{1,1}  &a_{1,2}   &\cdots &a_{1,n-1}  &a_{1,n} \\
           a_{2,1}  &a_{2,2}   &\cdots &a_{2,n-1}  &a_{2,n} \\
                  &\vdots                         \\
           a_{n-1,1} &a_{n-1,2} &\cdots &a_{n-1,n-1} &a_{n-1,n}  \\
           a_{n,1}  &a_{n,2}   &\cdots &a_{n,n-1}  &a_{n,n}
        \end{mat}
      \end{equation*}
      把除第一行外的每一行都乘以~$a_{1,1}$。
      \begin{equation*}
        \grstep[a_{1,1}\rho_3 \\ \vdots \\ a_{1,1}\rho_{n}]{a_{1,1}\rho_2}
        \begin{mat}
          a_{1,1}         &a_{1,2}        &\cdots &a_{1,n-1}        &a_{1,n} \\
          a_{2,1}a_{1,1}   &a_{2,2}a_{1,1}  &\cdots &a_{2,n-1}a_{1,1}  &a_{2,n}a_{1,1} \\
                        &\vdots                         \\
          a_{n-1,1}a_{1,1} &a_{n-1,2}a_{1,1} &\cdots &a_{n-1,n-1}a_{1,1} &a_{n-1,n}a_{1,1}\\
          a_{n,1}a_{1,1}  &a_{n,2}a_{1,1}   &\cdots &a_{n,n-1}a_{1,1}  &a_{n,n}a_{1,1}
        \end{mat}
      \end{equation*}
      这将行列式缩放了~$a_{1,1}^{n-1}$ 倍。

      接下来对每个行~$i>1$ 执行行变换
      $-a_{i,1}\rho_1+\rho_i$。
      这些行变换不改变行列式。
      \begin{equation*}
        \grstep[-a_{3,1}\rho_1+\rho_3 \\ \vdotswithin{-a_{3,1}\rho_1+\rho_3} \\ -a_{n,1}\rho_1+\rho_n]{-a_{2,1}\rho_1+\rho_2}
      \end{equation*}
      结果是一个矩阵~$B$，其第一行不变，第一列全为零（除了
      $1,1$ 元素为 $a_{1,1}$），其余元素如下。
      \begin{equation*}
        \begin{mat}
              a_{1,2}
              &\cdots
              &a_{1,n-1}
              &a_{1,n}
              \\
              a_{2,2}a_{1,1}-a_{2,1}a_{1,2}
              &\cdots
              &a_{2,n-1}a_{n,n}-a_{2,n-1}a_{1,n-1}
              &a_{2,n}a_{n,n}-a_{2,n}a_{1,n}
              \\
              \vdots                         \\
             a_{n,2}a_{1,1}-a_{n,1}a_{1,2}
             &\cdots
             &a_{n,n-1}a_{1,1}-a_{n,1}a_{1,n-1}
             &a_{n,n}a_{1,1}-a_{n,1}a_{1,n}
        \end{mat}
      \end{equation*}
       这个 $\nbyn{n}$ 矩阵 $B$ 的行列式是
       $A$ 的行列式的 $a_{1,1}^{n-1}$~倍。

     用~$C$ 表示矩阵 $B$ 的 $1,1$ 子式，
     即由最后 $n-1$ 行和最后 $n-1$ 列构成的子矩阵。
     沿 $B$ 的第一列作 Laplace 展开
     得到其行列式为 $\det(B) = (-1)^{1+1}a_{1,1}\det(C)$。

      若 $a_{1,1}\neq 0$，则令 $\set(B)$ 的两个表达式相等并约去公因子，
      得 $\det(A)=\det(C)/a_{n,n}^{n-2}$。
    
\end{ans}
\protect \topic {射影几何}
\begin{ans}{1}
      由点积
      \begin{equation*}
        0=\colvec[r]{1 \\ 0 \\ 0}\dotprod\rowvec{L_1 &L_2 &L_3}
         =L_1
      \end{equation*}
      我们得到方程为 $L_1=0$。
    
\end{ans}
\begin{ans}{2}
      \begin{exparts}
        \partsitem 这个行列式
          \begin{equation*}
            0=\begin{vmat}
              1  &4  &x \\
              2  &5  &y \\
              3  &6  &z
            \end{vmat}
            =-3x+6y-3z
          \end{equation*}
          表明这条直线是 $L=\rowvec{-3 &6 &-3}$。
        \partsitem $\colvec[r]{-3 \\ 6 \\ -3}$
      \end{exparts}
    
\end{ans}
\begin{ans}{3}
      与
      \begin{equation*}
        u=\colvec{u_1 \\ u_2 \\ u_3}
        \qquad
        v=\colvec{v_1 \\ v_2 \\ v_3}
      \end{equation*}
      关联的直线来自这个行列式方程。
      \begin{equation*}
        0=\begin{vmat}
          u_1  &v_1  &x  \\
          u_2  &v_2  &y  \\
          u_3  &v_3  &z
        \end{vmat}
        =(u_2v_3-u_3v_2)\cdot x
          + (u_3v_1-u_1v_3)\cdot y
          + (u_1v_2-u_2v_1)\cdot z
      \end{equation*}
      与两条直线关联的点的方程与此相同。
    
\end{ans}
\begin{ans}{4}
      若 $p_1$、$p_2$、$p_3$ 与 $q_1$、$q_2$、$q_3$ 是 $p$ 的两个三元组
      齐次坐标，则这两个列向量
      成比例，也就是说，落在同一条过
      原点的直线上。
      类似地，两个行向量也成比例。
      \begin{equation*}
        k\cdot\colvec{p_1 \\ p_2 \\ p_3}
          =\colvec{q_1 \\ q_2 \\ q_3}
        \qquad
        m\cdot\rowvec{L_1 &L_2 &L_3}
          =\rowvec{M_1 &M_2 &M_3}
      \end{equation*}
      于是相乘即得答案
      $(km)\cdot (p_1L_1+p_2L_2+p_3L_3)=q_1M_1+q_2M_2+q_3M_3=0$。
    
\end{ans}
\begin{ans}{5}
      日食的那幅图\Dash 除非
      图像平面恰好垂直于
      从太阳穿过小孔的连线\Dash 显示太阳的圆
      投影成一个椭圆的像。
      （另一个例子是，在本
      专题的许多图中，我们把作为球面赤道的圆画成了椭圆，
      也就是说，观看这幅图的人把圆看成了椭圆。）

      日食的那幅图也展示了其逆命题。
      如果我们把投影设想成从左到右
      穿过小孔，
      那么椭圆 $I$ 穿过 $P$ 投影成一个圆~$S$。
    
\end{ans}
\begin{ans}{6}
      单位球面上的一个点
      \begin{equation*}
        \colvec{p_1 \\ p_2 \\ p_3}
      \end{equation*}
      是非赤道的当且仅当 $p_3\neq 0$。
      在那种情形，它对应于 $z=1$ 平面上的这个点
      \begin{equation*}
        \colvec{p_1/p_3 \\ p_2/p_3  \\ 1}
      \end{equation*}
      因为那就是包含该向量的直线与该
      平面的交点。
    
\end{ans}
\begin{ans}{7}
      \begin{exparts}
        \partsitem 其他的图也是可能的，但这是其中之一。
          \begin{center}
            \includegraphics{det/mp/ch4.54}
          \end{center}
          各交点
          $
              T_0U_1\,\intersection T_1U_0=V_2
          $、$
              T_0V_1\,\intersection T_1V_0=U_2
          $ 与 $
              U_0V_1\,\intersection U_1V_0=T_2
          $
          这样标记，使得每条直线上都有一个 $T$、一个 $U$ 和一个 $V$。
        \partsitem Desargues 定理中所用的引理给出了一个
          基 $B$，关于这个基各点具有以下
          齐次坐标向量。
          \begin{equation*}
            \rep{\vec{t}_0}{B}=\colvec[r]{1 \\ 0 \\ 0}
            \quad
            \rep{\vec{t}_1}{B}=\colvec[r]{0 \\ 1 \\ 0}
            \quad
            \rep{\vec{t}_2}{B}=\colvec[r]{0 \\ 0 \\ 1}
            \quad
            \rep{\vec{v}_0}{B}=\colvec[r]{1 \\ 1 \\ 1}
          \end{equation*}
        \partsitem
          首先，$T_0V_0$ 上的任何 $U_0$
          \begin{equation*}
            \rep{\vec{u}_0}{B}=a\colvec[r]{1 \\ 0 \\ 0}
                               +b\colvec[r]{1 \\ 1 \\ 1}
                              =\colvec{a+b \\ b \\ b}
          \end{equation*}
          都具有这种形式的齐次坐标向量
          \begin{equation*}
            \colvec{u_0 \\ 1 \\ 1}
          \end{equation*}
          （$u_0$ 是一个参数；它取决于点 $U_0$ 在 $T_0V_0$ 直线上的位置，
          但那条直线上的任何点都具有某个 $u_0\in\Re$ 的
          这种形式的齐次坐标向量）。
          类似地，$U_2$ 在 $T_1V_0$ 上
          \begin{equation*}
            \rep{\vec{u}_2}{B}=c\colvec[r]{0 \\ 1 \\ 0}
                                +d\colvec[r]{1 \\ 1 \\ 1}
                              =\colvec{d \\ c+d \\ d}
          \end{equation*}
          因而具有这个齐次坐标向量。
          \begin{equation*}
            \colvec{1 \\ u_2 \\ 1}
          \end{equation*}
          同样类似地，$U_1$ 与 $T_2V_0$ 关联
          \begin{equation*}
            \rep{\vec{u}_1}{B}=e\colvec[r]{0 \\ 0 \\ 1}
                                +f\colvec[r]{1 \\ 1 \\ 1}
                              =\colvec{f \\ f \\ e+f}
          \end{equation*}
          并具有这个齐次坐标向量。
          \begin{equation*}
            \colvec{1 \\ 1 \\ u_1}
          \end{equation*}
        \partsitem
          因为 $V_1$ 是 $T_0U_2\,\intersection\,U_0T_2$，我们有下式。
          \begin{equation*}
            g\colvec[r]{1 \\ 0 \\ 0}+h\colvec{1 \\ u_2 \\ 1}
            =i\colvec{u_0 \\ 1 \\ 1}+j\colvec[r]{0 \\ 0 \\ 1}
            \qquad\Longrightarrow\qquad
            \begin{aligned}
              g+h  &= iu_0 \\
              hu_2 &= i    \\
              h    &= i+j
            \end{aligned}
          \end{equation*}
          把 $hu_2$ 代入第一个方程中的 $i$
          \begin{equation*}
            \colvec{hu_0u_2 \\ hu_2 \\ h}
          \end{equation*}
          可知 $V_1$ 具有这个
          双参数的齐次坐标向量。
          \begin{equation*}
            \colvec{u_0u_2 \\ u_2 \\ 1}
          \end{equation*}
        \partsitem
           由于 $V_2$ 是交点
           $T_0U_1\,\intersection\,T_1U_0$，
           \begin{equation*}
             k\colvec[r]{1 \\ 0 \\ 0}+l\colvec{1 \\ 1 \\ u_1}
              =m\colvec[r]{0 \\ 1 \\ 0}+n\colvec{u_0 \\ 1 \\ 1}
             \qquad\Longrightarrow\qquad
             \begin{aligned}
               k+l  &= nu_0 \\
               l    &= m+n    \\
               lu_1 &= n
             \end{aligned}
            \end{equation*}
           把 $lu_1$ 代入第一个方程中的 $n$
           \begin{equation*}
             \colvec{lu_0u_1 \\ l \\ lu_1}
           \end{equation*}
           可得
           $V_2$ 具有这个双参数的齐次坐标向量。
           \begin{equation*}
             \colvec{u_0u_1 \\ 1 \\ u_1}
           \end{equation*}
        \partsitem
           因为 $V_1$ 在 $T_1U_1$ 直线上，它的
           齐次坐标向量具有形式
           \begin{equation*}
             p\colvec[r]{0 \\ 1 \\ 0}+q\colvec{1 \\ 1 \\ u_1}
             =\colvec{q \\ p+q \\ qu_1}
           \tag*{($*$)}\end{equation*}
           而本题前一小问已经确立了 $V_1$ 的
           齐次坐标向量具有形式
           \begin{equation*}
            \colvec{u_0u_2 \\ u_2 \\ 1}
           \end{equation*}
           于是下面这是 $V_1$ 的一个齐次坐标向量。
           \begin{equation*}
             \colvec{u_0u_1u_2 \\ u_1u_2 \\ u_1}
           \tag*{($**$)}\end{equation*}
           由 ($*$) 与 ($**$)，三个参数之间存在一个
           关系：$u_0u_1u_2=1$。
         \partsitem
           $V_2$ 的齐次坐标向量可以这样写出。
           \begin{equation*}
             \colvec{u_0u_1u_2 \\ u_2 \\ u_1u_2}
             =\colvec{1 \\ u_2 \\ u_1u_2}
           \end{equation*}
           现在，$T_2U_2$ 直线由齐次
           坐标具有这种形式的点组成。
           \begin{equation*}
             r\colvec[r]{0 \\ 0 \\ 1}+s\colvec{1 \\ u_2 \\ 1}
             =\colvec{s \\ su_2 \\ r+s}
           \end{equation*}
           取 $s=1$ 与 $r=u_1u_2-1$ 即表明
           $V_2$ 的齐次坐标向量具有这种形式。
      \end{exparts}
      %\textit{Author's note.}
      % Good thing that there are no more parts because we've used all of the
      % lower-case letters.
    
\end{ans}
\protect \topic {计算机图形学}
\begin{ans}{1}
      直接计算得乘积如下。
      \begin{equation*}
        \begin{mat}
          1   &0  &0  \\
          0   &1  &2  \\
          0   &0  &1
        \end{mat}
        \begin{mat}
          0   &-1  &0  \\
          -1  &0   &0  \\
          0   &0   &1
        \end{mat}
        \begin{mat}
          1   &0  &0  \\
          0   &1  &-2  \\
          0   &0  &1
        \end{mat}
        =
        \begin{mat}
         0 &-1 &2 \\
        -1 &0  &2 \\
         0 &0  &1
        \end{mat}
      \end{equation*}
    
\end{ans}
\begin{ans}{2}
      在 $\R^2$ 中工作，设该矩阵为 $M$。
      我们得到这两个。
      \begin{equation*}
        \colvec{1 \\ 2}\mapsto \colvec{1 \\ 2}
        \qquad
        \colvec{-2 \\ 1}\mapsto \colvec{2 \\ -1}
      \end{equation*}
      （对于第二个，起始向量位于过原点且
      垂直于 $y=2x$ 的直线上。）
      我们有
      \begin{equation*}
        M
        \begin{mat}
          1  &-2  \\
          2  &1
        \end{mat}
        =
        \begin{mat}
          1  &2  \\
          2  &-1
        \end{mat}
      \end{equation*}
      求解得
      \begin{equation*}
        M =
            \begin{mat}
              1  &2  \\
              2  &-1
            \end{mat}
            \begin{mat}
              1  &-2  \\
              2  &1
            \end{mat}^{-1}
          =
            \begin{mat}
              -3/5  &4/5 \\
              4/5   &3/5
            \end{mat}
      \end{equation*}
      于是对齐次坐标而言，矩阵是这个。
      \begin{equation*}
      \begin{mat}
        -3/5  &4/5  &0  \\
         4/5  &3/5  &0  \\
           0  &0    &1
      \end{mat}
      \end{equation*}
    
\end{ans}
\begin{ans}{3}
      先把所有点移动 $4$，用前一题的结果关于直线 $y=2x$ 作反射，然后把它们移回去。
      \begin{equation*}
      \begin{mat}
        1  &0  &0  \\
        0  &1  &-4 \\
        0  &0  &1
      \end{mat}
      \begin{mat}
        -3/5  &4/5  &0  \\
         4/5  &3/5  &0  \\
           0  &0    &1
      \end{mat}
      \begin{mat}
        1  &0  &0  \\
        0  &1  &4 \\
        0  &0  &1
      \end{mat}
      =
      \begin{mat}
        -3/5  &4/5  &16/5  \\
         4/5  &3/5  &-8/5 \\
           0  &0    &1
      \end{mat}
      \end{equation*}
    
\end{ans}
\begin{ans}{4}
    \begin{equation*}
      \begin{mat}
        1  &0  &t_x  \\
        0  &1  &t_y  \\
        0  &0  &1
      \end{mat}
      \begin{mat}
        \cos\theta &-\sin\theta  &0 \\
        \sin\theta &\cos\theta   &0 \\
        0          &0            &1
      \end{mat}
      =
      \begin{mat}
        \cos\theta  &-\sin\theta  &t_x  \\
        \sin\theta  &\cos\theta   &t_y  \\
        0           &0            &1
      \end{mat}
    \end{equation*}
    
\end{ans}
\begin{ans}{5}
    下面是两个 $\nbyn{4}$ 矩阵。
    \begin{equation*}
      \begin{mat}
        \cos\theta  &-\sin\theta  &0  &0  \\
        \sin\theta  &\cos\theta   &0  &0  \\
        0           &0            &0  &0  \\
        0           &0            &0  &1
      \end{mat}
      \qquad
      \begin{mat}
        \cos\theta  &0  &-\sin\theta  &0   \\
        0           &0  &0            &0  \\
        \sin\theta  &0  &\cos\theta   &0   \\
        0           &0  &0            &1
      \end{mat}
    \end{equation*}
    
\end{ans}
\protect \chapter {相似}
\protect \section {第 I 节：复向量空间}
\protect \section {第 II 节：相似}
\protect \subsection {五.II.1: 定义与例题}
\begin{ans}{五.II.1.5}
      一种做法是从左到右计算。
      \begin{multline*}
        PTP^{-1}=
        \begin{mat}[r]
          4  &2  \\
         -3  &2
        \end{mat}
        \begin{mat}[r]
          1  &3  \\
         -2  &-6
        \end{mat}
        \begin{mat}[r]
          2/14  &-2/14  \\
          3/14  &4/14
        \end{mat}                                 \\
        =
        \begin{mat}[r]
          0  &0  \\
         -7  &-21
        \end{mat}
        \begin{mat}[r]
          2/14  &-2/14  \\
          3/14  &4/14
        \end{mat}
        =
        \begin{mat}[r]
          0    &0  \\
         -11/2 &-5
        \end{mat}
      \end{multline*}
    
\end{ans}
\begin{ans}{五.II.1.6}
     \begin{exparts}
      \partsitem 由于矩阵 $(2)$ 是 $\nbyn{1}$ 的，矩阵
         $P$ 与 $P^{-1}$ 也是 $\nbyn{1}$ 的，因此当
         $P=(p)$ 时，其逆为 $P^{-1}=(1/p)$。
         于是 $P(2)P^{-1}=(p)(2)(1/p)=(2)$。
       \partsitem 是的：~回忆一下，我们可以把标量倍数提到矩阵外面
        \( P(cI)P^{-1}=cPIP^{-1}=cI \)。
        顺便指出，零矩阵与单位矩阵分别是特殊情形
        $c=0$ 与 $c=1$。
      \partsitem 不是，如下面这个例子所示。
        \begin{equation*}
           \begin{mat}[r]
              1  &-2  \\
             -1  &1
            \end{mat}
           \begin{mat}[r]
             -1  &0   \\
              0  &-3
           \end{mat}
           \begin{mat}[r]
              -1  &-2   \\
              -1  &-1
           \end{mat}
           =
           \begin{mat}[r]
              -5  &-4   \\
              2   &1
           \end{mat}
        \end{equation*}
    \end{exparts}
   
\end{ans}
\begin{ans}{五.II.1.7}
      \begin{exparts}
        \partsitem
          \begin{equation*}
            \begin{CD}
              \C^3_{\wrt{B}}                   @>t>T>        \C^3_{\wrt{B}}       \\
              @V{\scriptstyle\identity} VV              @V{\scriptstyle\identity} VV \\
              \C^3_{\wrt{D}}                   @>t>\hat{T}>        \C^3_{\wrt{D}}
            \end{CD}
          \end{equation*}
        \partsitem
          对于起始基~$B$ 中的每个元素，求
          变换作用于其后的结果
          \begin{equation*}
             \colvec{1 \\ 2 \\ 3}\mapsunder{t}\colvec{-2 \\ 3 \\ 4}
             \qquad
             \colvec{0 \\ 1 \\ 0}\mapsunder{t}\colvec{0 \\ 0 \\ 2}
             \qquad
             \colvec{0 \\ 0 \\ 1}\mapsunder{t}\colvec{-1 \\ 1 \\ 0}
          \end{equation*}
          并将这些输出关于终止基~$B$ 表示
          \begin{equation*}
            \rep{\colvec{-2 \\ 3 \\ 4}}{B}=\colvec{-2 \\ 7 \\ 10}
            \qquad
            \rep{\colvec{0 \\ 0 \\ 2}}{B}=\colvec{0 \\ 0 \\ 2}
            \qquad
            \rep{\colvec{-1 \\ 1 \\ 0}}{B}=\colvec{-1 \\ 3 \\ 3}
          \end{equation*}
          由此得到矩阵。
          \begin{equation*}
            T=\rep{t}{B,B}=
            \begin{mat}
              -2 &0 &-1 \\
               7 &0 &3  \\
              10 &2 &3
            \end{mat}
          \end{equation*}
        \partsitem
          求变换作用于~$D$ 中元素的结果
          \begin{equation*}
             \colvec{1 \\ 0 \\ 0}\mapsunder{t}\colvec{1 \\ 0 \\ 0}
             \qquad
             \colvec{1 \\ 1 \\ 0}\mapsunder{t}\colvec{1 \\ 0 \\ 2}
             \qquad
             \colvec{1 \\ 0 \\ 1}\mapsunder{t}\colvec{0 \\ 1 \\ 0}
          \end{equation*}
          并将它们关于终止基~$D$ 表示
          \begin{equation*}
            \rep{\colvec{1 \\ 0 \\ 0}}{D}=\colvec{1 \\ 0 \\ 0}
            \qquad
            \rep{\colvec{1 \\ 0 \\ 2}}{D}=\colvec{-1 \\ 0 \\ 2}
            \qquad
            \rep{\colvec{0 \\ 1 \\ 0}}{D}=\colvec{-1 \\ 1 \\ 0}
          \end{equation*}
          由此得到矩阵。
          \begin{equation*}
            \hat{T}=\rep{t}{D,D}=
            \begin{mat}
               1 &-1 &-1 \\
               0 &0  &1  \\
               0 &2  &0
            \end{mat}
          \end{equation*}
        \partsitem
          要沿右侧向下，我们需要
          $\rep{\identity}{B,D}$
          因此我们先计算恒等映射对
          ~$B$ 的每个元素的作用——它没有任何作用，
          然后将结果关于
          ~$D$ 表示。
          \begin{equation*}
            \rep{\colvec{1 \\ 2 \\ 3}}{D}=\colvec{-4 \\ 2 \\ 3}
            \qquad
            \rep{\colvec{0 \\ 1 \\ 0}}{D}=\colvec{-1 \\ 1 \\ 0}
            \qquad
            \rep{\colvec{0 \\ 0 \\ 1}}{D}=\colvec{-1 \\ 0 \\ 1}
          \end{equation*}
          所以这就是~$P$。
          \begin{equation*}
            P=
            \begin{mat}
              -4 &-1 &-1 \\
               2 &1  &0  \\
               3 &0  &1
            \end{mat}
          \end{equation*}
          对于另一个矩阵~$\rep{\identity}{D,B}$，我们既可以
          像刚才对~$P$ 那样直接求出它，也可以进行
          常规的矩阵求逆计算。
          \begin{equation*}
            P^{-1}=
            \begin{mat}
               1 &1 &1 \\
               -2 &-1  &-2  \\
               -3 &-3  &-2
            \end{mat}
          \end{equation*}
      \end{exparts}
    
\end{ans}
\begin{ans}{五.II.1.8}
      \begin{exparts}
        \partsitem 由于我们用 $B$ 的成员来描述 $t$，
          求矩阵表示很容易：
          \begin{equation*}
            \rep{t(x^2)}{B}=\colvec[r]{0 \\ 1 \\ 1}_B
            \quad
            \rep{t(x)}{B}=\colvec[r]{1 \\ 0 \\ -1}_B
            \quad
            \rep{t(1)}{B}=\colvec[r]{0 \\ 0 \\ 3}_B
          \end{equation*}
          由此得到下式。
          \begin{equation*}
            \rep{t}{B,B}
            \begin{mat}[r]
              0  &1  &0  \\
              1  &0  &0  \\
              1  &-1 &3
            \end{mat}
          \end{equation*}
        \partsitem 我们将求 $t(1)$、$t(1+x)$ 和 $t(1+x+x^2$，
          以求出每一个关于 $D$ 的表示。
          已知 $t(1)=3$，另外两个容易看出：
          $t(1+x)=x^2+2$ 以及 $t(1+x+x^2)=x^2+x+3$。
          通过观察，我们得到每个向量的表示
          \begin{equation*}
            \rep{t(1)}{D}=\colvec[r]{3 \\ 0 \\ 0}_D
            \quad
            \rep{t(1+x)}{D}=\colvec[r]{2  \\ -1 \\  1}_D
            \quad
            \rep{t(1+x+x^2)}{D}=\colvec[r]{2 \\ 0 \\ 1}_D
          \end{equation*}
          从而得到该映射的表示。
          \begin{equation*}
            \rep{t}{D,D}
            =
            \begin{mat}[r]
              3  &2  &2  \\
              0  &-1 &0  \\
              0  &1  &1
            \end{mat}
          \end{equation*}
         \partsitem 下面的图表
           \begin{equation*}
             \begin{CD}
               V_{\wrt{B}}                  @>t>T>  V_{\wrt{B}}       \\
               @V\scriptstyle\identity VPV      @V\scriptstyle\identity VPV \\
               V_{\wrt{D}}                  @>t>\hat{T}>  V_{\wrt{D}}
             \end{CD}
           \end{equation*}
           表明它们是 $P=\rep{\identity}{B,D}$ 与
           $P^{-1}=\rep{\identity}{D,B}$。
           \begin{equation*}
             P=
             \begin{mat}[r]
               0  &-1  &1  \\
               -1  &1  &0  \\
               1  &0  &0
             \end{mat}
             \qquad
             P^{-1}=
             \begin{mat}[r]
               0  &0  &1  \\
               0  &1  &1  \\
               1  &1  &1
             \end{mat}
           \end{equation*}
      \end{exparts}
    
\end{ans}
\begin{ans}{五.II.1.9}
      \begin{exparts}
        \partsitem
          \begin{equation*}
            \begin{CD}
              \C^2_{\wrt{B}}                   @>t>T>        \C^2_{\wrt{B}}       \\
              @V{\scriptstyle\identity} VV              @V{\scriptstyle\identity} VV \\
              \C^2_{\wrt{D}}                   @>t>\hat{T}>        \C^2_{\wrt{D}}
            \end{CD}
          \end{equation*}
        \partsitem
          对于右边，我们求恒等映射的作用，
          它没有任何作用，
          \begin{equation*}
            \colvec{1 \\ 0}\mapsunder{\identity}\colvec{1 \\ 0}
            \qquad
            \colvec{1 \\ 1}\mapsunder{\identity}\colvec{1 \\ 1}
          \end{equation*}
          并将它们关于~$D$ 表示
          \begin{equation*}
            \rep{\colvec{1 \\ 0}}{D}=\colvec{1/2 \\ 0}
            \qquad
            \rep{\colvec{1 \\ 1}}{D}=\colvec{1/2 \\ -1/2}
          \end{equation*}
          于是我们得到下式。
          \begin{equation*}
            P=\rep{\identity}{B,D}=
            \begin{mat}
              1/2 &1/2 \\
              0   &-1/2
            \end{mat}
          \end{equation*}

          对于左边的矩阵，我们既可以像
          上一段那样直接计算，
          也可以取其逆。
          \begin{equation*}
            P^{-1}=\rep{\identity}{D,B}=
            \frac{1}{(-1/4)}\cdot
            \begin{mat}
              -1/2 &-1/2 \\
              0   &1/2
            \end{mat}
            =
            \begin{mat}
              2 &2 \\
              0 &-2
            \end{mat}
          \end{equation*}
        \partsitem
          与前一问一样，我们既可以
          由定义直接计算，
          也可以利用矩阵运算来计算。
          \begin{equation*}
            PTP^{-1}=
            \begin{mat}
              2 &2 \\
              0 &-2
            \end{mat}
            \begin{mat}
              1 &-1 \\
              2 &1
            \end{mat}
            \begin{mat}
              2 &2 \\
              0 &-2
            \end{mat}
            =
            \begin{mat}
              3 &3  \\
             -2 &-1
            \end{mat}
          \end{equation*}
      \end{exparts}
    
\end{ans}
\begin{ans}{五.II.1.10}
        基的一种可行选取是
        \begin{equation*}
          B=\sequence{\colvec[r]{1 \\ 2},\colvec[r]{-1 \\ 1}}
          \qquad
          D=\stdbasis_2=\sequence{\colvec[r]{1 \\ 0},\colvec[r]{0 \\ 1}}
        \end{equation*}
        （这个 $B$ 来自对该映射的描述）。
        为求矩阵 $\hat{T}=\rep{t}{B,B}$，解关系式
        \begin{equation*}
           c_1\colvec[r]{1 \\ 2}+c_2\colvec[r]{-1 \\ 1}=\colvec[r]{3 \\ 0}
           \qquad
          \hat{c}_1\colvec[r]{1 \\ 2}+\hat{c}_2\colvec[r]{-1 \\ 1}=\colvec[r]{-1 \\ 2}
        \end{equation*}
        得 \( c_1=1 \)、\( c_2=-2 \)、\( \hat{c}_1=1/3 \) 及
        \( \hat{c}_2=4/3 \)。
        \begin{equation*}
           \rep{t}{B,B}=
           \begin{mat}[r]
              1  &1/3 \\
             -2  &4/3
           \end{mat}
        \end{equation*}

        求 \( \rep{t}{D,D} \) 需要稍多一些计算。
        我们先求 \( t(\vec{e}_1) \)。
        关系式
        \begin{equation*}
          c_1\colvec[r]{1 \\ 2}+c_2\colvec[r]{-1 \\ 1}=\colvec[r]{1 \\ 0}
        \end{equation*}
        给出 \( c_1=1/3 \) 与 \( c_2=-2/3 \)，于是
        \begin{equation*}
          \rep{\vec{e}_1}{B}=\colvec[r]{1/3 \\ -2/3}_B
        \end{equation*}
        从而
        \begin{equation*}
           \rep{t(\vec{e}_1)}{B}=
                  \begin{mat}[r]
                     1  &1/3  \\
                    -2  &4/3
                  \end{mat}_{B,B}
                  \colvec[r]{1/3 \\ -2/3}_B
                  =
                  \colvec[r]{1/9 \\ -14/9}_B
        \end{equation*}
        因此 $t$ 对第一个基向量 $\vec{e}_1$ 的作用如下。
        \begin{equation*}
          t(\vec{e}_1)
          =(1/9)\cdot\colvec[r]{1 \\ 2}-(14/9)\cdot\colvec[r]{-1 \\ 1}
          =\colvec[r]{5/3 \\ -4/3}
        \end{equation*}
        \( t(\vec{e}_2) \) 的计算类似。
        关系式
        \begin{equation*}
          c_1\colvec[r]{1 \\ 2}+c_2\colvec[r]{-1 \\ 1}=\colvec[r]{0 \\ 1}
        \end{equation*}
        给出 \( c_1=1/3 \) 与 \( c_2=1/3 \)，所以
        \begin{equation*}
          \rep{\vec{e}_1}{B}=\colvec[r]{1/3 \\ 1/3}_B
        \end{equation*}
        从而
        \begin{equation*}
          \rep{t(\vec{e}_1)}{B}=
                  \begin{mat}[r]
                     1  &1/3  \\
                    -2  &4/3
                  \end{mat}_{B,B}
                  \colvec[r]{1/3 \\ 1/3}_B
                  =
                  \colvec[r]{4/9 \\ -2/9}_B
        \end{equation*}
        因此 $t$ 对第二个基向量 $\vec{e}_2$ 的作用如下。
        \begin{equation*}
          t(\vec{e}_2)
          =(4/9)\cdot\colvec[r]{1 \\ 2}-(2/9)\cdot\colvec[r]{-1 \\ 1}
          =\colvec[r]{2/3 \\ 2/3}
        \end{equation*}
        因此
        \begin{equation*}
           \rep{t}{D,D}=
           \begin{mat}[r]
              5/3  &2/3  \\
             -4/3  &2/3
           \end{mat}
        \end{equation*}
        于是得到这个矩阵。
        \begin{equation*}
          P=\rep{\identity}{B,D}
          =\begin{mat}[r]
            1  &-1  \\
            2  &1
          \end{mat}
        \end{equation*}
       而这一个用于换基。
       \begin{equation*}
          P^{-1}=\bigl(\rep{\identity}{B,D}\bigr)^{-1}
          =\begin{mat}[r]
             1  &-1 \\
             2  &1
          \end{mat}^{-1}
          =
          \begin{mat}[r]
            1/3  &1/3  \\
            -2/3 &1/3
          \end{mat}
       \end{equation*}
       这些计算的检验是例行的。
       \begin{equation*}
          \begin{mat}[r]
             1  &-1 \\
             2  &1
          \end{mat}
          \begin{mat}[r]
             1  &1/3 \\
            -2  &4/3
          \end{mat}
          \begin{mat}[r]
            1/3 &1/3 \\
           -2/3 &1/3
          \end{mat}
          =
          \begin{mat}[r]
            5/3 &2/3 \\
           -4/3 &2/3
          \end{mat}
       \end{equation*}
    
\end{ans}
\begin{ans}{五.II.1.11}
      高斯消元法表明，
      第一个矩阵表示的映射秩为二，而第二个
      矩阵表示的映射秩为三。
    
\end{ans}
\begin{ans}{五.II.1.12}
      零映射的表示只有零矩阵，
      无论基对是什么都有 $\rep{z}{B,D}=Z$，
      因此特别地，对任何单独一个基 $B$ 我们有 $\rep{z}{B,B}=Z$。
      恒等映射的情形稍有不同：~恒等映射关于任何 $B,B$ 的
      唯一表示
      是单位矩阵 $\rep{\identity}{B,B}=I$。
      （\textit{注：}~当然，我们已经见过 $B\neq D$ 且
      $\rep{\identity}{B,D}\neq I$ 的例子 \Dash 事实上，我们已经看到任何
      非奇异矩阵都是恒等映射关于某个
      $B,D$ 的表示。）
    
\end{ans}
\begin{ans}{五.II.1.13}
      不存在。
      若 \( A=PBP^{-1} \)，则 \( A^2=(PBP^{-1})(PBP^{-1})=PB^2P^{-1} \)。
    
\end{ans}
\begin{ans}{五.II.1.14}
       矩阵相似是矩阵等价的特殊情形
       （若矩阵相似则它们矩阵等价），
       而矩阵等价保持非奇异性。
    
\end{ans}
\begin{ans}{五.II.1.15}
       矩阵与它自身相似；取 \( P \) 为单位
       矩阵：~$P=IPI^{-1}=IPI$。

       若 \( \hat{T} \) 相似于 \( T \)，则 \( \hat{T}=PTP^{-1} \)
       于是 \( P^{-1}\hat{T}P=T \)。
       将其改写为 \( T=(P^{-1})\hat{T}(P^{-1})^{-1} \)，即可得出
       $T$ 相似于 \( \hat{T} \)。

        对于传递性，
        若 \( T \) 相似于 \( S \) 且 \( S \) 相似于 \( U \)，
        则 \( T=PSP^{-1} \) 且 \( S=QUQ^{-1} \)。
        于是 \( T=PQUQ^{-1}P^{-1}=(PQ)U(PQ)^{-1} \)，这表明 \( T \)
        相似于 \( U \)。
      
\end{ans}
\begin{ans}{五.II.1.16}
         设 $f_x$ 与 $f_y$ 为这两个反射映射（有时称为「翻转」）。
         对任何基
         \( B \) 与 \( D \)，矩阵 \( \rep{f_x}{B,B}  \) 与
         \( \rep{f_y}{D,D} \) 相似。
         首先注意
         \begin{equation*}
            S=\rep{f_x}{\stdbasis_2,\stdbasis_2}=
            \begin{mat}[r]
               1  &0  \\
               0  &-1
            \end{mat}
            \qquad
            T=\rep{f_y}{\stdbasis_2,\stdbasis_2}=
            \begin{mat}[r]
              -1  &0  \\
               0  &1
            \end{mat}
         \end{equation*}
         相似，因为第二个矩阵是 $f_x$
         关于基 \( A=\sequence{\vec{e}_2,\vec{e}_1} \) 的表示：
         \begin{equation*}
            \begin{mat}[r]
               1  &0  \\
               0  &-1
            \end{mat}
            =
            P
            \begin{mat}[r]
              -1  &0  \\
               0  &1
            \end{mat}
            P^{-1}
         \end{equation*}
         其中 $P=\rep{\identity}{A,\stdbasis_2}$。
         \begin{equation*}
           \begin{CD}
             \C^2_{\wrt{A}}
                @>f_x>T>
                V\C^2_{\wrt{A}}       \\
             @V\scriptstyle\identity VPV
                @V\scriptstyle\identity VPV \\
             \C^2_{\wrt{\stdbasis_2}}
                @>f_x>S>
                \C^2_{\wrt{\stdbasis_2}}
           \end{CD}
         \end{equation*}
         现在结论由
         \nearbyexercise{exer:SimIsEquivRel} 的传递性部分得出。

        我们也可以不依赖那道练习来完成。
        记
        $\rep{f_x}{B,B}=QTQ^{-1}=Q\rep{f_x}{\stdbasis_2,\stdbasis_2}Q^{-1}$
        以及
        $\rep{f_y}{D,D}=RSR^{-1}=R\rep{f_y}{\stdbasis_2,\stdbasis_2}R^{-1}$。
        由第一段中的等式，
        这两个中的第一个是
        $\rep{f_x}{B,B}=QP\rep{f_y}{\stdbasis_2,\stdbasis_2}P^{-1}Q^{-1}$。
        将这两个中的第二个改写为
        $R^{-1}\cdot\rep{f_y}{D,D}\cdot R=\rep{f_y}{\stdbasis_2,\stdbasis_2}$
        并代入，就得到所需的关系
        \begin{multline*}
          \rep{f_x}{B,B}
          =QP\rep{f_y}{\stdbasis_2,\stdbasis_2}P^{-1}Q^{-1}  \\
          =QPR^{-1}\cdot\rep{f_y}{D,D}\cdot RP^{-1}Q^{-1}
          =(QPR^{-1})\cdot\rep{f_y}{D,D}\cdot (QPR^{-1})^{-1}
        \end{multline*}
        因此矩阵 \( \rep{f_x}{B,B}  \) 与 \( \rep{f_y}{D,D} \)
        相似。
      
\end{ans}
\begin{ans}{五.II.1.17}
        我们必须证明：若两个矩阵相似，则它们有相同的
        秩和相同的行列式。

        首先，我们已经证明
        映射的秩等于表示该映射的任何矩阵的秩，
        所以秩对所有表示都相同。
        （换句话说，秩是映射的性质，
        因为它是值域空间的维数。）

        至于行列式，
        \( \deter{PSP^{-1}}=\deter{P}\cdot\deter{S}\cdot\deter{P^{-1}}
           =\deter{P}\cdot\deter{S}\cdot\deter{P}^{-1}=\deter{S} \)。

       该命题的逆命题不成立。
       一方面，
       存在行列式相同却不相似的矩阵。
       举个例子，考虑一个行列式为
       零的非零矩阵。
       它不与零矩阵相似，因为零矩阵只与
       它自身相似，但二者行列式相同。
       同样的思路给出秩的反例。
     
\end{ans}
\begin{ans}{五.II.1.18}
       包含所有秩为
       零的 \( \nbyn{n} \) 矩阵的矩阵等价类只含一个矩阵，即零矩阵。
       因此它作为子集只含有一个相似类。

       相比之下，秩为一的
       \( \nbyn{1} \) 矩阵的矩阵等价类由那些
       \( k\neq 0 \) 的 $\nbyn{1}$ 矩阵 \( (k) \) 组成。
       对任何基 \( B \)，
       用标量 \( k \) 相乘的映射的表示
       为 \( \rep{t_k}{B,B}=(k) \)，
       所以每个这样的矩阵独自构成其相似类。
       所以这是矩阵等价类分裂成
       无穷多个相似类的一个例子。
      
\end{ans}
\begin{ans}{五.II.1.19}
       能，下面两个矩阵相似
       \begin{equation*}
          \begin{mat}[r]
            1  &0  \\
            0  &3
          \end{mat}
          \qquad
          \begin{mat}[r]
            3  &0  \\
            0  &1
          \end{mat}
       \end{equation*}
       因为，当第一个矩阵是关于
       $B=\sequence{\vec{\beta}_1,\vec{\beta}_2}$ 的 $\rep{t}{B,B}$ 时，
       第二个矩阵就是关于
       $D=\sequence{\vec{\beta}_2,\vec{\beta}_1}$ 的 $\rep{t}{D,D}$。
      
\end{ans}
\begin{ans}{五.II.1.20}
       \( k \) 次幂相似，因为当每个矩阵表示
       映射 $t$ 时，\( k \) 次幂表示
       \( t^k \)，即 $k$ 个 $t$ 的复合。
       （例如，若 $T=rep{t}{B,B}$，则 $T^2=\rep{\composed{t}{t}}{B,B}$。）

       用更计算化的说法重述：若 \( T=PSP^{-1} \)，则
       \( T^2=(PSP^{-1})(PSP^{-1})=PS^2P^{-1} \)。
       归纳法将此推广到所有幂。

       对于 $k\leq 0$ 的情形，
       设 \( S \) 可逆且 \( T=PSP^{-1} \)。
       注意 \( T \) 可逆：
       \( T^{-1}=(PSP^{-1})^{-1}=PS^{-1}P^{-1} \)，
       而同一个等式表明
       \( T^{-1} \) 相似于 \( S^{-1} \)。
       其余负次幂则由第一段给出。
     
\end{ans}
\begin{ans}{五.II.1.21}
        从概念上说，二者都表示某个
        变换 \( t \) 的 \( p(t) \)。
        从计算上说，我们有下式。
        \begin{align*}
           p(T)
           &=c_n(PSP^{-1})^n+\dots+c_1(PSP^{-1})+c_0I   \\
           &=c_nPS^nP^{-1}+\dots+c_1PSP^{-1}+c_0I   \\
           &=Pc_nS^nP^{-1}+\dots+Pc_1SP^{-1}+Pc_0P^{-1}   \\
           &=P(c_nS^n+\dots+c_1S+c_0)P^{-1}
        \end{align*}
      
\end{ans}
\begin{ans}{五.II.1.22}
       有两个等价类：(i)~秩为零的矩阵的类，
       其中只有一个矩阵：
       $\mathscr{C}_1=\set{(0)}$，
       以及 (2)~秩为一的矩阵的类，
       其中有无穷多个：
       \( \mathscr{C}_2=\set{(k)\suchthat k\neq0} \)。

      每个 \( \nbyn{1} \) 矩阵都独自构成其相似类。
      这是因为一维空间上的任何变换
      都是乘以一个标量的映射 $\map{t_k}{V}{V}$，由
      $\vec{v}\mapsto k\cdot\vec{v}$ 给出。
      于是，对任何基 \( B=\sequence{\vec{\beta}} \)，
      表示变换 \( t_k \) 的矩阵
      关于 \( B,B \) 为
      \( (\rep{t_k(\vec{\beta})}{B})=(k) \)。

      所以，矩阵等价类
      $\mathscr{C}_1$ 中（显然）只含有由矩阵 $(0)$ 构成的
      唯一相似类。
      而矩阵等价类 $\mathscr{C}_2$ 中含有
      无穷多个各含一个成员的相似类，由
      $k\neq0$ 的 $(k)$ 构成。
    
\end{ans}
\begin{ans}{五.II.1.23}
       不保持。
       下面是一个例子，其中有两对矩阵，每对中的两个矩阵相似：
       \begin{equation*}
          \begin{mat}[r]
             1  &-1  \\
             1  &2
          \end{mat}
          \begin{mat}[r]
             1  &0   \\
             0  &3
          \end{mat}
          \begin{mat}[r]
             2/3  &1/3   \\
            -1/3  &1/3
          \end{mat}
          =
          \begin{mat}[r]
             5/3  &-2/3  \\
            -4/3  &7/3
          \end{mat}
       \end{equation*}
       以及
       \begin{equation*}
          \begin{mat}[r]
             1  &-2  \\
            -1  &1
          \end{mat}
          \begin{mat}[r]
            -1  &0   \\
             0  &-3
          \end{mat}
          \begin{mat}[r]
             -1  &-2   \\
             -1  &-1
          \end{mat}
          =
          \begin{mat}[r]
             -5  &-4   \\
              2  &1
          \end{mat}
       \end{equation*}
       （这个例子并非完全任意，因为
       两个等号左边的中间矩阵相加为零矩阵）。
       注意这些相似矩阵的和并不相似
       \begin{equation*}
          \begin{mat}[r]
             1  &0   \\
             0  &3
          \end{mat}
          +
          \begin{mat}[r]
            -1  &0   \\
             0  &-3
          \end{mat}
          =
          \begin{mat}[r]
            0  &0  \\
            0  &0
          \end{mat}
          \qquad
          \begin{mat}[r]
             5/3  &-2/3   \\
             -4/3 &7/3
          \end{mat}
          +
          \begin{mat}[r]
             -5  &-4   \\
              2  &1
          \end{mat}
          \neq
          \begin{mat}[r]
            0  &0  \\
            0  &0
          \end{mat}
       \end{equation*}
       因为零矩阵只与它自身相似。
     
\end{ans}
\begin{ans}{五.II.1.24}
    若 \( N=P(T-\lambda I)P^{-1} \)，则
    \( N=PTP^{-1}-P(\lambda I)P^{-1} \)。
    对角矩阵 \( \lambda I \) 与任何矩阵可交换，所以
    \( P(\lambda I)P^{-1}=PP^{-1}(\lambda I)=\lambda I \)。
    于是 \( N=PTP^{-1}-\lambda I \)，
    从而 \( N+\lambda I=PTP^{-1} \)。
    （所以二者不仅相似，事实上它们是经由
    同一个 \( P \) 相似的。）
    
\end{ans}
\protect \subsection {五.II.2: 可对角化性}
\begin{ans}{五.II.2.6}
      由于基向量是任意选取的，因此可能有许多不同的答案。
      不过，下面给出一种做法；
      要对角化
      \begin{equation*}
         T=\begin{mat}[r]
            4  &-2  \\
            1  &1
         \end{mat}
      \end{equation*}
      把它视为某个变换在标准基下的表示 $T=\rep{t}{\stdbasis_2,\stdbasis_2}$，并寻找
      \( B=\sequence{\vec{\beta}_1,\vec{\beta}_2} \) 使得
      \begin{equation*}
        \rep{t}{B,B}
        =
        \begin{mat}
          \lambda_1  &0          \\
          0          &\lambda_2
        \end{mat}
      \end{equation*}
      也就是说，使得
      $t(\vec{\beta}_1)=\lambda_1$ 且 $t(\vec{\beta}_2)=\lambda_2$。
      \begin{equation*}
        \begin{mat}[r]
           4  &-2  \\
           1  &1
        \end{mat}
        \vec{\beta}_1=\lambda_1\cdot\vec{\beta}_1
        \qquad
        \begin{mat}[r]
           4  &-2  \\
           1  &1
        \end{mat}
        \vec{\beta}_2=\lambda_2\cdot\vec{\beta}_2
      \end{equation*}
      我们要寻找标量 \( x \)，使得方程
      \begin{equation*}
       \begin{mat}[r]
          4  &-2  \\
          1  &1
       \end{mat}
       \colvec{b_1 \\ b_2}=x\cdot\colvec{b_1 \\ b_2}
      \end{equation*}
      有不全为零的解 $b_1$ 与 $b_2$。
      将其改写为线性方程组
      \begin{equation*}
        \begin{linsys}{2}
           (4-x)\cdot b_1  &+  &-2\cdot b_2       &=  &0  \\
           1\cdot b_1      &+   &(1-x)\cdot b_2   &=  &0
        \end{linsys}
      \end{equation*}
      若 $x=4$，则第一个方程给出 $b_2=0$，进而第二个方程给出 $b_1=0$。
      我们已排除两个 $b$ 都为零的情形，
      故可假设 $x\neq 4$。
      \begin{equation*}
        \grstep{(-1/(4-x))\rho_1+\rho_2}
        \begin{linsys}{2}
           (4-x)\cdot b_1  &+   &-2\cdot b_2                   &=  &0  \\
                           &    &((x^2-5x+6)/(4-x))\cdot b_2   &=  &0
        \end{linsys}
      \end{equation*}
      考虑最下面的方程。
      若 \( b_2=0 \)，则第一个方程给出 $b_1=0$ 或 $x=4$。
      $b_1=b_2=0$ 的情形是不允许的。
      最下面方程的另一种可能是分式 $x^2-5x+6=(x-2)(x-3)$ 的分子为零。
      $x=2$ 的情形给出第一个方程 $2b_1-2b_2=0$，因此与 $x=2$ 相关联的是
      第一、第二分量相等的向量：
      \begin{equation*}
         \vec{\beta}_1=\colvec[r]{1 \\ 1}
         \qquad\text{（于是 \(
           \begin{mat}[r]
              4  &-2  \\
              1  &1
           \end{mat}
           \colvec[r]{1 \\ 1}=2\cdot\colvec[r]{1 \\ 1} \)，且 $\lambda_1=2$）。}
      \end{equation*}
      若 \( x=3 \)，则第一个方程为
      $b_1-2b_2=0$，于是相关的向量
      是那些第一分量为第二分量两倍的向量：
      \begin{equation*}
         \vec{\beta}_2=\colvec[r]{2 \\ 1}
         \qquad\text{（于是 \(
           \begin{mat}[r]
              4  &-2  \\
              1  &1
           \end{mat}
           \colvec[r]{2 \\ 1}=3\cdot\colvec[r]{2 \\ 1} \)，从而 $\lambda_2=3$）。}
      \end{equation*}
      下图
        \begin{equation*}
          \begin{CD}
            \C^2_{\wrt{\stdbasis_2}}
               @>t>T>
               \C^2_{\wrt{\stdbasis_2}}       \\
            @V{\scriptstyle\identity} VV
               @V{\scriptstyle\identity} VV \\
            \C^2_{\wrt{B}}
               @>t>D>
               \C^2_{\wrt{B}}
          \end{CD}
        \end{equation*}
      展示了如何得到该对角化。
      \begin{equation*}
         \begin{mat}[r]
           2  &0  \\
           0  &3
         \end{mat}
         =
         \begin{mat}[r]
           1  &2  \\
           1  &1
         \end{mat}^{-1}
         \begin{mat}[r]
           4  &-2  \\
           1  &1
         \end{mat}
         \begin{mat}[r]
           1  &2  \\
           1  &1
         \end{mat}
      \end{equation*}
      \textit{评注}。
      在下述改名之下，这个等式与 $T=PSP^{-1}$ 的定义相符。
      \begin{equation*}
        T=
         \begin{mat}[r]
           2  &0  \\
           0  &3
         \end{mat}
         \quad
         P=
         \begin{mat}[r]
           1  &2  \\
           1  &1
         \end{mat}^{-1}
         \quad
         P^{-1}=
         \begin{mat}[r]
           1  &2  \\
           1  &1
         \end{mat}
         \quad
         S=
         \begin{mat}[r]
           4  &-2  \\
           1  &1
         \end{mat}
      \end{equation*}
    
\end{ans}
\begin{ans}{五.II.2.7}
     \begin{exparts}
        \partsitem 我们要一个基，使其中的向量满足
          \begin{equation*}
            \begin{mat}
              1 &1 \\
              0 &0
            \end{mat}
            \colvec{b_1 \\ b_2}
            =x\cdot \colvec{b_1 \\ b_2}
          \end{equation*}
          于是考虑由此得到的线性方程组。
          \begin{equation*}
            \begin{linsys}{2}
              b_1 &+ &b_2 &= &x\cdot b_1 \\
                  &  &0   &= &x\cdot b_2 \\
            \end{linsys}
          \end{equation*}
        \partsitem
          若 $x=0$，则 $b_1+b_2=0$，即 $b_1=-b_2$。
          也就是说，与这种情形相关联的基元素，
          其分量符号相反。
          （分量必须非零，因为零向量
          不能是基的成员。）

          若 $x\neq 0$，则 $b_2=0$，把它代入上面的
          方程得到 $b_1=x\cdot b_1$。
          所以在这种情形，要么 $b_1=0$，要么 $x=1$。
          但由于此情形中 $b_2=0$，$b_1=0$ 是不可能的，因为
          零向量不在任何基中。
          于是~$x=1$，方程组变为
          \begin{equation*}
            \begin{linsys}{2}
              b_1 &+ &b_2 &= &1\cdot b_1 \\
                  &  &0   &= &1\cdot b_2 \\
            \end{linsys}
          \end{equation*}
          所以这种情形相关联的基元素满足 $b_2=0$，
          而~$b_1$ 可取任意非零值。

          因此，合适的基是
          \begin{equation*}
            B=\sequence{\colvec{1 \\ -1},
                        \colvec{1 \\ 0}}
          \end{equation*}
        \partsitem
          有了这个基，
          为得到对角化，画出相似图。
          \begin{equation*}
            \begin{CD}
            \C^2_{\wrt{\stdbasis_2}}
               @>t>T>
               \C^2_{\wrt{\stdbasis_2}}       \\
            @V{\scriptstyle\identity} VV
               @V{\scriptstyle\identity} VV \\
            \C^2_{\wrt{B}}
               @>t>D>
               \C^2_{\wrt{B}}
           \end{CD}
          \end{equation*}
          在此图中，
          \begin{equation*}
             T=
             \begin{mat}
               1  &1  \\
               0  &0
             \end{mat}
          \end{equation*}
          我们将利用基~$B$，
          沿着左下到左上、再到右上、
          再到右下的路径来计算矩阵~$D$。

          我们先要从左下走到左上。
          注意到由于我们把图上方的基
          选为标准基 $\stdbasis_2$，
          表示该映射的矩阵很容易写出。
          \begin{equation*}
            \rep{\identity}{B,\stdbasis_2}=
            \begin{mat}
              1 &1 \\
             -1 &0
            \end{mat}
          \end{equation*}
          整条路径的计算确实给出了
          对角的结果
          （当然，左下到左上这一段出现在
          三个矩阵乘积的右边）。
          \begin{align*}
                 D
                 &=
                 \begin{mat}
                    1  &1  \\
                    -1  &0
                 \end{mat}^{-1}
                 \begin{mat}
                    1  &1  \\
                    0  &0
                 \end{mat}
                 \begin{mat}
                    1  &1  \\
                    -1  &0
                 \end{mat}                \\
                 &=
                 \begin{mat}
                    0  &-1  \\
                    1  &1
                 \end{mat}
                 \begin{mat}
                    1  &1  \\
                    0  &0
                 \end{mat}
                 \begin{mat}
                    1  &1  \\
                    -1  &0
                 \end{mat}
                =\begin{mat}
                    0  &0  \\
                    0  &1
                 \end{mat}
          \end{align*}
     \end{exparts}
    
\end{ans}
\begin{ans}{五.II.2.8}
     \begin{exparts}
        \partsitem
          基中的成员满足
          \begin{equation*}
            \begin{mat}
              0 &1 \\
              1 &0
            \end{mat}
            \colvec{b_1 \\ b_2}
            =x\cdot \colvec{b_1 \\ b_2}
          \end{equation*}
          因此由此得到线性方程组
          \begin{equation*}
            \begin{linsys}{2}
                   &  &b_2 &= &x\cdot b_1 \\
               b_1 &  &    &= &x\cdot b_2 \\
            \end{linsys}
            \qquad\Longrightarrow\qquad
            \begin{linsys}{2}
               -xb_1 &+ &b_2  &= &0 \\
                 b_1 &- &xb_2 &= &0 \\
            \end{linsys}
          \end{equation*}
        \partsitem
          若 $x=0$，则 $b_1=b_2=0$，但零向量不是任何
          基的成员，所以这种情形不可能发生。

          另一种情形是 $x\neq 0$。
          \begin{equation*}
            \begin{linsys}{2}
               -xb_1 &+ &b_2  &= &0 \\
                 b_1 &- &xb_2 &= &0 \\
            \end{linsys}
            \grstep{(1/x)\rho_1+\rho_2}
            \begin{linsys}{2}
               -xb_1 &+ &b_2           &= &0 \\
                     &  &(-x+(1/x))b_2 &= &0 \\
            \end{linsys}
          \end{equation*}
          关注最下面的方程。
          要么 $b_2=0$，要么 $-x+(1/x)=0$。
          若 $b_2=0$，则代入上面的方程给出
          $b_1=0$，因为 $x\neq 0$ 是这种情形的假设。
          但零向量不是任何基的成员，所以这不可能
          发生。

          于是只剩这种可能。
          \begin{equation*}
            \frac{1}{x}-x=0
            \qquad\Longrightarrow\qquad
            0=\frac{1-x^2}{x}=\frac{(1-x)(1+x)}{x}
          \end{equation*}
          若 $x=1$，则方程组为
          \begin{equation*}
            \begin{linsys}{2}
               -1\cdot b_1 &+ &b_2  &= &0 \\
                           &  &0    &= &0 \\
            \end{linsys}
          \end{equation*}
          于是基向量满足 $b_2=b_1$。
          若 $x=-1$，则方程组为
          \begin{equation*}
            \begin{linsys}{2}
               b_1 &+ &b_2  &= &0 \\
                   &  &0    &= &0 \\
            \end{linsys}
          \end{equation*}
          且基向量满足 $b_2=-b_1$。
          因此，下面是一个合适的基的例子。
          \begin{equation*}
            B=\sequence{\colvec{1 \\ 1},
                        \colvec{1 \\ -1}}
          \end{equation*}
        \partsitem
          为得到对角化，画出相似图
          \begin{equation*}
            \begin{CD}
            \C^2_{\wrt{\stdbasis_2}}
               @>t>T>
               \C^2_{\wrt{\stdbasis_2}}       \\
            @V{\scriptstyle\identity} VV
               @V{\scriptstyle\identity} VV \\
            \C^2_{\wrt{B}}
               @>t>D>
               \C^2_{\wrt{B}}
           \end{CD}
          \end{equation*}
          其中 $T$ 是题目中给出的矩阵。
          \begin{equation*}
             T=
             \begin{mat}
               0  &1  \\
               1  &0
             \end{mat}
          \end{equation*}
          利用基~$B$ 计算矩阵~$D$，
          从左下走到左上，再到右上，
          然后到右下。

          注意表示左下到左上
          映射的矩阵很容易写出。
          \begin{equation*}
            \rep{\identity}{B,\stdbasis_2}=
            \begin{mat}
              1 &1 \\
              1 &-1
            \end{mat}
          \end{equation*}
          有了它，整条路径的计算确实给出
          对角的结果。
          \begin{align*}
                 D
                 &=
                 \begin{mat}
                    1  &1  \\
                    1  &-1
                 \end{mat}^{-1}
                 \begin{mat}
                    0  &1  \\
                    1  &0
                 \end{mat}
                 \begin{mat}
                    1  &1  \\
                    1  &-1
                 \end{mat}                \\
                 &=
                 -\frac{1}{2}
                 \begin{mat}
                    -1  &-1  \\
                    -1  &1
                 \end{mat}
                 \begin{mat}
                    0  &1  \\
                    1  &0
                 \end{mat}
                 \begin{mat}
                    1  &1  \\
                    1  &-1
                 \end{mat}
                =\begin{mat}
                    1  &0  \\
                    0  &-1
                 \end{mat}
          \end{align*}
      \end{exparts}
    
\end{ans}
\begin{ans}{五.II.2.9}
      \begin{exparts}
        \partsitem
          写出
          \begin{equation*}
            \begin{mat}[r]
              -2  &1  \\
               0  &2
            \end{mat}
            \colvec{b_1 \\ b_2}
            =x\cdot\colvec{b_1 \\ b_2}
            \qquad\Longrightarrow\qquad
            \begin{linsys}{2}
               (-2-x)\cdot b_1  &+  &b_2            &=  &0  \\
                                &   &(2-x)\cdot b_2 &= &0
            \end{linsys}
          \end{equation*}
          给出两种可能：$b_2=0$ 与 $x=2$。
          沿 $b_2=0$ 这种可能，第一个方程为
          $(-2-x)b_1=0$，有两种情形：$b_1=0$ 或
          $x=-2$。
          于是，在第一种可能之下，我们得到 $x=-2$ 以及
          相关联的向量，其第二分量为零，而
          第一分量自由。
          \begin{equation*}
            \begin{mat}[r]
              -2  &1  \\
               0  &2
            \end{mat}
            \colvec{b_1 \\ 0}
            =-2\cdot\colvec{b_1 \\ 0}
            \qquad
            \vec{\beta}_1=\colvec[r]{1 \\ 0}
          \end{equation*}
          沿另一种可能，得到第一个方程
          $-4b_1+b_2=0$，因此与这个
          解相关联的向量的第二分量是其第一分量
          的四倍。
          \begin{equation*}
            \begin{mat}[r]
              -2  &1  \\
               0  &2
            \end{mat}
            \colvec{b_1 \\ 4b_1}
            =2\cdot\colvec{b_1 \\ 4b_1}
            \qquad
            \vec{\beta}_2=\colvec{1 \\ 4}
          \end{equation*}
          对角化如下。
          \begin{equation*}
            \begin{mat}[r]
              1  &1  \\
              0  &4
            \end{mat}
            \begin{mat}[r]
              -2  &1  \\
               0  &2
            \end{mat}
            \begin{mat}[r]
              1  &1  \\
              0  &4
            \end{mat}^{-1}
            =
            \begin{mat}[r]
              -2  &0  \\
               0  &2
            \end{mat}
          \end{equation*}
      \partsitem 计算与前一小题类似。
          \begin{equation*}
            \begin{mat}[r]
               5  &4  \\
               0  &1
            \end{mat}
            \colvec{b_1 \\ b_2}
            =x\cdot\colvec{b_1 \\ b_2}
            \qquad\Longrightarrow\qquad
            \begin{linsys}{2}
               (5-x)\cdot b_1  &+  &4\cdot b_2     &=  &0  \\
                               &   &(1-x)\cdot b_2 &=  &0
            \end{linsys}
          \end{equation*}
          最下面的方程
          给出两种可能：$b_2=0$ 与 $x=1$。
          沿 $b_2=0$ 这种可能，并舍去
          $b_2$ 与 $b_1$ 都为零的情形，得到
          $x=5$，与之相关联的向量第二分量
          为零，而第一分量自由。
          \begin{equation*}
            \vec{\beta}_1=\colvec[r]{1 \\ 0}
          \end{equation*}
          $x=1$ 这种可能给出第一个方程
          $4b_1+4b_2=0$，因此相关联的向量的
          第二分量是其第一分量的相反数。
          \begin{equation*}
            \vec{\beta}_1=\colvec[r]{1 \\ -1}
          \end{equation*}
          于是我们有如下对角化。
          \begin{equation*}
            \begin{mat}[r]
              1  &1  \\
              0  &-1
            \end{mat}
            \begin{mat}[r]
               5  &4  \\
               0  &1
            \end{mat}
            \begin{mat}[r]
              1  &1  \\
              0  &-1
            \end{mat}^{-1}
            =
            \begin{mat}[r]
               5  &0  \\
               0  &1
            \end{mat}
          \end{equation*}
      \end{exparts}
    
\end{ans}
\begin{ans}{五.II.2.10}
      \begin{exparts}
        \item  在这个相似图
          \begin{equation*}
            \begin{CD}
            \C^2_{\wrt{\stdbasis_2}}
               @>n>N>
               \C^2_{\wrt{\stdbasis_2}}       \\
            @V{\scriptstyle\identity} VV
               @V{\scriptstyle\identity} VV \\
            \C^2_{\wrt{B}}
               @>n>D>
               \C^2_{\wrt{B}}
           \end{CD}
          \end{equation*}
          中，我们将尝试计算矩阵~$D$ 并找出合适的
          基~$B$，
          即沿左下到左上、再到右上、
          再到右下的路径行进。
        \item 该图表明，当
          $B=\sequence{\vec{\beta}_1,\vec{\beta}_2}$ 时，我们要
          $n(\vec{\beta}_1)=\lambda_1\vec{\beta}_1$ 且
          $n(\vec{\beta}_2)=\lambda_2\vec{\beta}_2$。
          为确定 $\lambda_1$ 和 $\lambda_2$，
          由于矩阵~$N$ 表示的是关于
          标准基的映射，我们要求这个方程组的解。
          \begin{equation*}
            \begin{mat}
              0 &0 \\
              1 &0
            \end{mat}
            \colvec{b_1 \\ b_2}
            =
            x\cdot\colvec{b_1 \\ b_2}
          \end{equation*}
          同解的线性方程组为
          \begin{equation*}
            \begin{linsys}{2}
              0b_1 &+  &0b_2 &= &xb_1 \\
              b_1  &+  &0b_2 &= &xb_2 \\
            \end{linsys}
            \qquad
            \Longrightarrow
            \qquad
            \begin{linsys}{2}
              -xb_1 &   &     &= &0 \\
              b_1   &-  &xb_2 &= &0 \\
            \end{linsys}
          \end{equation*}
          由第一行，要么 $x=0$，要么~$b_1=0$。
          若 $x=0$，则第二个方程变为 $b_1-0=0$，故 $b_1=0$，
          于是得到这个解集。
          \begin{equation*}
            \set{\colvec{0 \\ 1}b_2 \suchthat b_2\in\C}
          \end{equation*}
          我们可以取基的第一个成员为
          $\vec{\beta}_1=\vec{e}_2$。

          若 $x\neq 0$，则线性方程组的
          第一个方程给出 $b_1=0$。
          方程组的第一行变为 $0=0$，第二行变为
          $-xb_2=0$。
          由于这种情形取 $x\neq 0$，我们得到 $b_2=0$，
          解集是平凡的。
          \begin{equation*}
            \set{\colvec{0 \\ 0}}
          \end{equation*}
          所以出问题的地方在于：我们得不到
          \nearbylemma{lm:DiagIffBasisOfEigens}所要求的基。
      \end{exparts}
    
\end{ans}
\begin{ans}{五.II.2.11}
      对于任意整数 \( p \)，我们有下式。
      \begin{equation*}
        \begin{mat}
          d_1  &0      &   \\
          0    &\ddots &   \\
               &       &d_n
        \end{mat}^p=
        \begin{mat}
          d_1^p  &0      &   \\
          0      &\ddots &   \\
                 &       &d_n^p
        \end{mat}
     \end{equation*}
    
\end{ans}
\begin{ans}{五.II.2.12}
       这两个矩阵不相似
       \begin{equation*}
          \begin{mat}[r]
             0  &0  \\
             0  &0
          \end{mat}
          \qquad
          \begin{mat}[r]
             1  &0  \\
             0  &1
          \end{mat}
       \end{equation*}
       因为每一个在各自的相似类中都是单独的。

       对于后半问，这两个
       \begin{equation*}
          \begin{mat}[r]
             2  &0  \\
             0  &3
          \end{mat}
          \qquad
          \begin{mat}[r]
             3  &0  \\
             0  &2
          \end{mat}
       \end{equation*}
       通过把基从
       \( \sequence{\vec{\beta}_1,\vec{\beta}_2} \) 变为
       \( \sequence{\vec{\beta}_2,\vec{\beta}_1} \) 的矩阵而相似。
       （\textit{问题。}
        两个对角矩阵相似当且仅当它们的对角元互为排列？）
    
\end{ans}
\begin{ans}{五.II.2.13}
      对比这两个矩阵。
      \begin{equation*}
         \begin{mat}[r]
           2  &0  \\
           0  &1
         \end{mat}
         \qquad
         \begin{mat}[r]
           2  &0  \\
           0  &0
         \end{mat}
      \end{equation*}
      第一个是非奇异的，第二个是奇异的。
     
\end{ans}
\begin{ans}{五.II.2.14}
       要验证对角矩阵的逆是由各元素的逆构成的对角矩阵，只需直接相乘。
       \begin{equation*}
          \begin{mat}
             a_{1,1}  &0                \\
             0        &a_{2,2}          \\
                      &       &\ddots    \\
                      &       &      &a_{n,n}
          \end{mat}
          \begin{mat}
            1/a_{1,1}  &0                \\
             0        &1/a_{2,2}          \\
                      &       &\ddots    \\
                      &       &      &1/a_{n,n}
          \end{mat}
      \end{equation*}
      （证明它是左逆同样容易。）

      若某个对角元为零，则该对角矩阵是奇异的；它的行列式为零。
    
\end{ans}
\begin{ans}{五.II.2.15}
     \begin{exparts}
       \partsitem 验证很容易。
         \begin{equation*}
           \begin{mat}[r]
             1  &1  \\
             0  &-1
           \end{mat}
           \begin{mat}[r]
             3  &2  \\
             0  &1
           \end{mat}
           =
           \begin{mat}[r]
             3  &3  \\
             0  &-1
           \end{mat}
           \qquad
           \begin{mat}[r]
             3  &3  \\
             0  &-1
           \end{mat}
           \begin{mat}[r]
             1  &1  \\
             0  &-1
           \end{mat}^{-1}
           =
           \begin{mat}[r]
             3  &0  \\
             0  &1
           \end{mat}
         \end{equation*}
        \partsitem 这是一个巧合，也就是说，若 $T=PSP^{-1}$，
          则 $T$ 不必等于 $P^{-1}SP$。
          甚至对于对角矩阵~$D$，条件 $D=PTP^{-1}$
          也不能推出 $D$ 等于 $P^{-1}TP$。
          来自 \nearbyexample{ex:DiagTwoByTwo} 的矩阵说明了这一点。
          \begin{equation*}
            \begin{mat}[r]
              1  &2  \\
              1  &1
            \end{mat}
            \begin{mat}[r]
              4  &-2  \\
              1  &1
            \end{mat}
            =
            \begin{mat}[r]
              6  &0  \\
              5  &-1
            \end{mat}
            \qquad
            \begin{mat}[r]
              6  &0  \\
              5  &-1
            \end{mat}
            \begin{mat}[r]
              1  &2  \\
              1  &1
            \end{mat}^{-1}
            =
            \begin{mat}[r]
              -6  &12  \\
              -6  &11
            \end{mat}
          \end{equation*}
     \end{exparts}
   
\end{ans}
\begin{ans}{五.II.2.16}
      该矩阵的列就是与各个 $x$ 相联系的向量。
      具体的选取，以及选取的顺序，都是任意的。
      例如，交换这两列就能得到另一个矩阵。
    
\end{ans}
\begin{ans}{五.II.2.17}
      先对角化，再取对角矩阵的幂，可知
      \begin{equation*}
        \begin{mat}[r]
          -3  &1  \\
          -4  &2
        \end{mat}^k
        =
        \frac{1}{3}
        \begin{mat}[r]
          -1  &1  \\
          -4  &4
        \end{mat}
        +(\frac{-2}{3})^k
        \begin{mat}[r]
           4  &-1 \\
           4  &-1
        \end{mat}.
      \end{equation*}
     
\end{ans}
\begin{ans}{五.II.2.18}
      是的，由本小节最后一个定理可知 \( ct \) 可对角化。

      不是，\( t+s \) 不一定可对角化。
      直观地说，当这两个映射关于不同的基对角化时
      （即当它们不是
      \definend{可同时对角化的}）问题就出现了。
      具体地，这两个矩阵都可对角化，但它们的和不可对角化：
      \begin{equation*}
         \begin{mat}[r]
            1  &1  \\
            0  &0
         \end{mat}
         \qquad
         \begin{mat}[r]
           -1  &0  \\
            0  &0
         \end{mat}
      \end{equation*}
      （第二个已经是对角矩阵；第一个
      参见 \nearbyexercise{exer:DiagTheseA}）。
      其和不可对角化，因为它的平方是零矩阵。

     同样的直觉提示 \( \composed{t}{s} \) 也不
     可对角化。
     这两个矩阵都可对角化，但它们的乘积不可对角化：
     \begin{equation*}
        \begin{mat}[r]
           1  &0  \\
           0  &0
        \end{mat}
        \qquad
        \begin{mat}[r]
           0  &1  \\
           1  &0
        \end{mat}
     \end{equation*}
     （第二个参见 \nearbyexercise{exer:DiagTheseB}）。
    
\end{ans}
\begin{ans}{五.II.2.19}
      若
      \begin{equation*}
         P
         \begin{mat}
            1  &c  \\
            0  &1
         \end{mat}
         P^{-1}
         =
         \begin{mat}
            a  &0  \\
            0  &b
         \end{mat}
      \end{equation*}
      则
      \begin{equation*}
         P
         \begin{mat}
            1  &c  \\
            0  &1
         \end{mat}
         =
         \begin{mat}
            a  &0  \\
            0  &b
         \end{mat}
         P
      \end{equation*}
      于是
      \begin{align*}
         \begin{mat}
            p  &q  \\
            r  &s
         \end{mat}
         \begin{mat}
            1  &c  \\
            0  &1
         \end{mat}
         &=
         \begin{mat}
            a  &0  \\
            0  &b
         \end{mat}
         \begin{mat}
            p  &q  \\
            r  &s
         \end{mat}        \\
         \begin{mat}
            p  &cp+q  \\
            r  &cr+s
         \end{mat}
         &=
         \begin{mat}
            ap  &aq  \\
            br  &bs
         \end{mat}
      \end{align*}
      \( 1,1 \) 元表明 \( a=1 \)，于是 \( 1,2 \) 元
      表明 \( pc=0 \)。
      由于 \( c\neq 0 \)，这意味着 \( p=0 \)。
      \( 2,1 \) 元表明
      \( b=1 \)，于是 \( 2,2 \) 元表明
      \( rc=0 \)。
      由于 \( c\neq 0 \)，这意味着 \( r=0 \)。
      但若 \( p \) 与 \( r \) 都是 \( 0 \)，则 \( P \) 不
      可逆。
     
\end{ans}
\begin{ans}{五.II.2.20}
      \begin{exparts}
      \partsitem 使用 $\nbyn{2}$
        矩阵的逆的公式可得下式。
        \begin{multline*}
           \begin{mat}
              a  &b  \\
              c  &d
           \end{mat}
           \begin{mat}[r]
              1  &2  \\
              2  &1
           \end{mat}
           \cdot\frac{1}{ad-bc}\cdot
           \begin{mat}
              d  &-b \\
             -c  &a
           \end{mat}                                  \\
           =\frac{1}{ad-bc}
           \begin{mat}
              ad+2bd-2ac-bc    &-ab-2b^2+2a^2+ab \\
              cd+2d^2-2c^2-cd  &-bc-2bd+2ac+ad
           \end{mat}
        \end{multline*}
        现在选取标量 \( a,\ldots,d \)，使得
        \( ad-bc\neq 0 \) 且 \( 2d^2-2c^2=0 \) 且 \( 2a^2-2b^2=0 \)。
        例如，下面这组就行。
        \begin{equation*}
          \begin{mat}[r]
            1  &1  \\
            1  &-1
          \end{mat}
          \begin{mat}[r]
            1  &2  \\
            2  &1
          \end{mat}
          \cdot\frac{1}{-2}\cdot
          \begin{mat}[r]
            -1  &-1  \\
            -1  &1
          \end{mat}
          =
          \frac{1}{-2}
          \begin{mat}[r]
            -6  &0   \\
             0  &2
          \end{mat}
        \end{equation*}
      \partsitem 同上，
        \begin{multline*}
           \begin{mat}
              a  &b  \\
              c  &d
           \end{mat}
           \begin{mat}
              x  &y  \\
              y  &z
           \end{mat}
           \cdot\frac{1}{ad-bc}\cdot
           \begin{mat}
              d  &-b \\
             -c  &a
           \end{mat}                               \\
           =\frac{1}{ad-bc}
           \begin{mat}
              adx+bdy-acy-bcz    &-abx-b^2y+a^2y+abz \\
              cdx+d^2y-c^2y-cdz  &-bcx-bdy+acy+adz
           \end{mat}
        \end{multline*}
        我们要找标量 \( a,\ldots,d \)，使得
        \( ad-bc\neq 0 \) 且
        \( -abx-b^2y+a^2y+abz=0 \)
        且 \( cdx+d^2y-c^2y-cdz=0 \)，无论
        \( x \)、\( y \)、\( z \)
        取什么值。

        首先，我们假定 \( y\neq 0 \)，否则所给矩阵已经
        是对角矩阵。
        我们将使用该假定，因为若我们（随意地）令
        \( a=1 \) 则得到
        \begin{align*}
           -bx-b^2y+y+bz
           &=0              \\
           (-y)b^2+(z-x)b+y
           &=0
        \end{align*}
        而求根公式给出
        \begin{equation*}
           b=\frac{-(z-x)\pm\sqrt{(z-x)^2-4(-y)(y)} }{-2y}
           \qquad
           y\neq 0
        \end{equation*}
        （注意若 \( x \)、\( y \)、\( z \) 为实数，则这两个
         \( b \) 都是实数，因为判别式为正）。
        同理，若我们（随意地）令 \( c=1 \) 则
        \begin{align*}
           dx+d^2y-y-dz
           &=0              \\
           (y)d^2+(x-z)d-y
           &=0
        \end{align*}
        在这里我们得到
        \begin{equation*}
           d=\frac{-(x-z)\pm\sqrt{(x-z)^2-4(y)(-y)} }{2y}
           \qquad
           y\neq 0
        \end{equation*}
        （同上，若 \( x,y,z\in\Re \)，则此判别式为正，
        故实对称的 \( \nbyn{2} \) 矩阵相似于一个实的
        对角矩阵）。

        作为检验，我们取 \( x=1 \)、\( y=2 \)、\( z=1 \)。
        \begin{equation*}
           b=\frac{0\pm\sqrt{0+16} }{-4}=\mp 1
           \qquad
           d=\frac{0\pm\sqrt{0+16} }{4}=\pm 1
        \end{equation*}
        注意并非所有四个选择 \( (b,d)=(+1,+1),\dots,(-1,-1) \)
        都满足 \( ad-bc\neq 0 \)。
      \end{exparts}
    
\end{ans}
\protect \subsection {五.II.3: 特征值与特征向量}
\begin{ans}{五.II.3.22}
      我们可以调换列的位置，这只需调换做表示时所用的基。
      \begin{equation*}
        \begin{mat}
          3 &0 \\
          0 &-4
        \end{mat}
        \qquad
        \begin{mat}
          -4 &0 \\
          0 &3
        \end{mat}
      \end{equation*}
    
\end{ans}
\begin{ans}{五.II.3.23}
      \begin{exparts}
         \partsitem 这个
           \begin{equation*}
             0=
             \begin{vmatrix}
               10-x  &-9  \\
               4     &-2-x
             \end{vmatrix}
             =(10-x)(-2-x)-(-36)
           \end{equation*}
           化简后得到特征方程 \( x^2-8x+16=0 \)。
           由于该方程可分解为 $(x-4)^2$，所以只有一个特征值
           \( \lambda_1=4 \)。
         \partsitem $0=(1-x)(3-x)-8=x^2-4x-5$；$\lambda_1=5$，$\lambda_2=-1$
         \partsitem \( x^2-21=0 \)；
           \( \lambda_1=\sqrt{21} \)，$\lambda_2=-\sqrt{21}$
         \partsitem \( x^2=0 \)；\( \lambda_1=0 \)
         \partsitem \( x^2-2x+1=0 \)；\( \lambda_1=1 \)
       \end{exparts}
     
\end{ans}
\begin{ans}{五.II.3.24}
       \begin{exparts}
         \partsitem 特征方程是 \( (3-x)(-1-x)=0 \)。
           它的根即特征值，为 \( \lambda_1=3 \) 与
           \( \lambda_2=-1 \)。
           为求特征向量，我们考虑这个方程。
           \begin{equation*}
             \begin{mat}
               3-x  &0    \\
               8    &-1-x
             \end{mat}
             \colvec{b_1  \\  b_2}
             =\colvec[r]{0  \\  0}
           \end{equation*}
           为求与 $\lambda_1=3$ 相对应的特征向量，
           我们考虑由此得到的线性方程组。
           \begin{equation*}
             \begin{linsys}{2}
               0\cdot b_1  &+  &0\cdot b_2  &=  &0  \\
               8\cdot b_1  &+  &-4\cdot b_2 &=  &0
             \end{linsys}
           \end{equation*}
           特征子空间是第二个分量等于第一个分量两倍的向量组成的集合。
           \begin{equation*}
             \set{\colvec{b_2/2 \\ b_2}\suchthat b_2\in\C}
             \qquad
             \begin{mat}[r]
               3  &0  \\
               8  &-1
             \end{mat}
             \colvec{b_2/2  \\ b_2}
             =3\cdot\colvec{b_2/2 \\ b_2}
           \end{equation*}
           （这里取 $b_2$ 作为参数，只是因为它是上面方程组中的自由变量。）
           因此，这是一个与特征值 $3$ 相对应的特征向量。
           \begin{equation*}
             \colvec[r]{1 \\ 2}
           \end{equation*}

           求与 $\lambda_2=-1$ 相对应的特征向量与此类似。方程组
           \begin{equation*}
             \begin{linsys}{2}
               4\cdot b_1  &+  &0\cdot b_2  &=  &0  \\
               8\cdot b_1  &+  &0\cdot b_2 &=  &0
             \end{linsys}
           \end{equation*}
           导出由第一个分量为零的向量组成的集合。
           \begin{equation*}
             \set{\colvec{0 \\ b_2}\suchthat b_2\in\C}
             \qquad
             \begin{mat}[r]
               3  &0  \\
               8  &-1
             \end{mat}
             \colvec{0  \\ b_2}
             =-1\cdot\colvec{0 \\ b_2}
           \end{equation*}
           于是这是与 $\lambda_2$ 相对应的一个特征向量。
           \begin{equation*}
             \colvec[r]{0 \\ 1}
           \end{equation*}
         \partsitem 特征方程是
           \begin{equation*}
             0=
             \begin{vmatrix}
               3-x  &2  \\
               -1   &-x
             \end{vmatrix}
             =x^2-3x+2=(x-2)(x-1)
           \end{equation*}
           因此特征值为 $\lambda_1=2$ 和 $\lambda_2=1$。
           为求特征向量，考虑这个方程组。
           \begin{equation*}
             \begin{linsys}{2}
               (3-x)\cdot b_1  &+  &2\cdot b_2  &=  &0  \\
               -1\cdot b_1     &-  &x\cdot b_2  &=  &0
             \end{linsys}
           \end{equation*}
           取 $\lambda_1=2$ 我们得到
           \begin{equation*}
             \begin{linsys}{2}
                1\cdot b_1  &+  &2\cdot b_2  &=  &0  \\
               -1\cdot b_1  &-  &2\cdot b_2  &=  &0
             \end{linsys}
           \end{equation*}
           由此得到这个特征子空间和特征向量。
           \begin{equation*}
             \set{\colvec{-2b_2 \\ b_2}
                   \suchthat b_2\in\C}
             \qquad
             \colvec[r]{-2 \\ 1}
           \end{equation*}
           取 $\lambda_2=1$，方程组为
           \begin{equation*}
             \begin{linsys}{2}
                2\cdot b_1  &+  &2\cdot b_2  &=  &0  \\
               -1\cdot b_1  &-  &1\cdot b_2  &=  &0
             \end{linsys}
           \end{equation*}
           由此得到这个。
           \begin{equation*}
             \set{\colvec{-b_2 \\ b_2}
                   \suchthat b_2\in\C}
             \qquad
             \colvec[r]{-1 \\ 1}
           \end{equation*}
       \end{exparts}
     
\end{ans}
\begin{ans}{五.II.3.25}
         特征方程
           \begin{equation*}
             0=
             \begin{vmatrix}
               -2-x  &-1  \\
               5     &2-x
             \end{vmatrix}
             =x^2+1
           \end{equation*}
           有复根 $\lambda_1=i$ 和 $\lambda_2=-i$。
           对于方程组
           \begin{equation*}
             \begin{linsys}{2}
               (-2-x)\cdot b_1  &-  &1\cdot b_2      &=  &0  \\
               5\cdot b_1       &   &(2-x)\cdot b_2  &=  &0
             \end{linsys}
           \end{equation*}
           取 $\lambda_1=i$，用高斯消元法得到如下化简。
           \begin{equation*}
             \begin{linsys}{2}
                (-2-i)\cdot b_1  &-  &1\cdot b_2      &=  &0  \\
                5\cdot b_1       &-  &(2-i)\cdot b_2  &=  &0
             \end{linsys}
             \grstep{(-5/(-2-i))\rho_1+\rho_2}
             \begin{linsys}{2}
                (-2-i)\cdot b_1  &-  &1\cdot b_2      &=  &0  \\
                                 &   &0               &=  &0
             \end{linsys}
           \end{equation*}
           （对于右下角的计算，通分后取公分母
           \begin{equation*}
             \frac{5}{-2-i}-(2-i)
             =
             \frac{5}{-2-i}-\frac{-2-i}{-2-i}\cdot (2-i)
             =
             \frac{5-(-5)}{-2-i}
           \end{equation*}
           即可看出它给出一个 $0=0$ 的方程。）
           由此得到特征子空间和特征向量。
           \begin{equation*}
             \set{\colvec{(1/(-2-i))b_2 \\ b_2}
                   \suchthat b_2\in\C}
             \qquad
             \colvec{1/(-2-i) \\ 1}
           \end{equation*}
           取 $\lambda_2=-i$，方程组
           \begin{equation*}
             \begin{linsys}{2}
                (-2+i)\cdot b_1  &-  &1\cdot b_2      &=  &0  \\
                5\cdot b_1       &-  &(2+i)\cdot b_2  &=  &0
             \end{linsys}
             \grstep{(-5/(-2+i))\rho_1+\rho_2}
             \begin{linsys}{2}
                (-2+i)\cdot b_1  &-  &1\cdot b_2      &=  &0  \\
                                 &   &0               &=  &0
             \end{linsys}
           \end{equation*}
           导出这个。
           \begin{equation*}
             \set{\colvec{(1/(-2+i))b_2 \\ b_2}
                   \suchthat b_2\in\C}
             \qquad
             \colvec{1/(-2+i) \\ 1}
           \end{equation*}
    
\end{ans}
\begin{ans}{五.II.3.26}
      特征方程是
      \begin{equation*}
        0=
        \begin{vmatrix}
          1-x  &1   &1   \\
          0    &-x  &1   \\
          0    &0   &1-x
        \end{vmatrix}
        =(1-x)^2(-x)
      \end{equation*}
      因此特征值为 $\lambda_1=1$（它是方程的一个重根）和 $\lambda_2=0$。
      对于其余部分，考虑这个方程组。
      \begin{equation*}
        \begin{linsys}{3}
          (1-x)\cdot b_1  &+  &b_2         &+  &b_3            &=  &0  \\
                          &   &-x\cdot b_2 &+  &b_3            &=  &0  \\
                          &   &            &   &(1-x)\cdot b_3 &= &0
        \end{linsys}
      \end{equation*}
      当 $x=\lambda_1=1$ 时，解集是这个特征子空间。
      \begin{equation*}
        \set{\colvec{b_1 \\ 0 \\ 0}\suchthat b_1\in\C}
      \end{equation*}
      当 $x=\lambda_2=0$ 时，解集是这个特征子空间。
      \begin{equation*}
        \set{\colvec{-b_2 \\ b_2 \\ 0}\suchthat b_2\in\C}
      \end{equation*}
      所以这些就是与 $\lambda_1=1$ 和
      $\lambda_2=0$ 相对应的特征向量。
      \begin{equation*}
        \colvec[r]{1 \\ 0 \\ 0}
        \qquad
        \colvec[r]{-1 \\ 1 \\ 0}
      \end{equation*}
    
\end{ans}
\begin{ans}{五.II.3.27}
      \begin{exparts}
         \partsitem 特征方程是
           \begin{equation*}
             0=
             \begin{vmatrix}
              3-x  &-2   &0  \\
             -2    &3-x  &0  \\
              0    &0    &5-x
             \end{vmatrix}
             =x^3-11x^2+35x-25=(x-1)(x-5)^2
           \end{equation*}
           因此特征值为 $\lambda_1=1$，还有重复的特征值
           $\lambda_2=5$。
           为求特征向量，考虑这个方程组。
           \begin{equation*}
             \begin{linsys}{3}
               (3-x)\cdot b_1  &-  &2\cdot b_2      &   &   &=  &0  \\
               -2\cdot b_1     &+  &(3-x)\cdot b_2  &   &   &=  &0  \\
                               &   &                &   &(5-x)\cdot b_3 &= &0
             \end{linsys}
           \end{equation*}
           取 $\lambda_1=1$ 我们得到
           \begin{equation*}
             \begin{linsys}{3}
                2\cdot b_1     &-  &2\cdot b_2   &   &   &=  &0  \\
               -2\cdot b_1     &+  &2\cdot b_2   &   &   &=  &0  \\
                               &   &             &   &4\cdot b_3 &= &0
             \end{linsys}
           \end{equation*}
           由此得到这个特征子空间和特征向量。
           \begin{equation*}
             \set{\colvec{b_2 \\ b_2 \\ 0}
                   \suchthat b_2\in\C}
             \qquad
             \colvec[r]{1 \\ 1 \\ 0}
           \end{equation*}
           取 $\lambda_2=5$，方程组为
           \begin{equation*}
             \begin{linsys}{3}
                -2\cdot b_1  &-  &2\cdot b_2   &   &   &=  &0  \\
               -2\cdot b_1   &-  &2\cdot b_2   &   &   &=  &0  \\
                             &   &             &   &0\cdot b_3 &= &0
             \end{linsys}
           \end{equation*}
           由此得到这个。
           \begin{equation*}
             \set{\colvec{-b_2 \\ b_2 \\ 0}+\colvec{0 \\ 0 \\ b_3}
                   \suchthat b_2,b_3\in\C}
             \qquad
             \colvec[r]{-1 \\ 1 \\ 0},\,\colvec[r]{0 \\ 0 \\ 1}
           \end{equation*}
         \partsitem 特征方程是
           \begin{equation*}
             0=
             \begin{vmatrix}
              -x   &1    &0  \\
              0    &-x   &1  \\
              4    &-17  &8-x
             \end{vmatrix}
             =-x^3+8x^2-17x+4=-1\cdot(x-4)(x^2-4x+1)
           \end{equation*}
           特征值为 $\lambda_1=4$ 以及（用求根公式）
           $\lambda_2=2+\sqrt{3}$ 和
           $\lambda_3=2-\sqrt{3}$。
           为求特征向量，考虑这个方程组。
           \begin{equation*}
             \begin{linsys}{3}
               -x\cdot b_1  &+  &b_2          &   &               &=  &0  \\
                            &   &-x\cdot b_2  &+  &b_3            &=  &0  \\
               4\cdot b_1   &-  &17\cdot b_2  &+  &(8-x)\cdot b_3 &= &0
             \end{linsys}
           \end{equation*}
           代入 $x=\lambda_1=4$ 得到方程组
           \begin{align*}
             \begin{linsys}{3}
               -4\cdot b_1  &+  &b_2          &   &           &= &0  \\
                            &   &-4\cdot b_2  &+  &b_3        &= &0  \\
               4\cdot b_1   &-  &17\cdot b_2  &+  &4\cdot b_3 &= &0
             \end{linsys}
             &\grstep{\rho_1+\rho_3}
             \begin{linsys}{3}
               -4\cdot b_1  &+  &b_2          &   &           &= &0  \\
                            &   &-4\cdot b_2  &+  &b_3        &= &0  \\
                            &   &-16\cdot b_2  &+ &4\cdot b_3 &= &0
             \end{linsys}                                                \\
             &\grstep{-4\rho_2+\rho_3}
             \begin{linsys}{3}
               -4\cdot b_1  &+  &b_2          &   &           &= &0  \\
                            &   &-4\cdot b_2  &+  &b_3        &= &0  \\
                            &   &              &  &0          &= &0
             \end{linsys}
           \end{align*}
           由此得到这个特征子空间和特征向量。
           \begin{equation*}
             V_4=\set{\colvec{(1/16)\cdot b_3 \\ (1/4)\cdot b_3 \\ b_3}
                   \suchthat b_2\in\C}
             \qquad
             \colvec[r]{1 \\ 4 \\ 16}
           \end{equation*}

           代入 $x=\lambda_2=2+\sqrt{3}$ 得到方程组
           \begin{multline*}
             \begin{linsys}{3}
               (-2-\sqrt{3})\cdot b_1
                      &+  &b_2
                          &   &           &= &0  \\
                      &   &(-2-\sqrt{3})\cdot b_2
                          &+  &b_3        &= &0  \\
               4\cdot b_1
                      &-  &17\cdot b_2
                          &+  &(6-\sqrt{3})\cdot b_3 &= &0
             \end{linsys}                                   \\
             \grstep{(-4/(-2-\sqrt{3}))\rho_1+\rho_3}
             \begin{linsys}{3}
               (-2-\sqrt{3})\cdot b_1
                      &+  &b_2
                          &   &           &= &0  \\
                      &   &(-2-\sqrt{3})\cdot b_2
                          &+  &b_3        &= &0  \\
                      &+  &(-9-4\sqrt{3})\cdot b_2
                          &+  &(6-\sqrt{3})\cdot b_3 &= &0
             \end{linsys}
           \end{multline*}
           （第三个方程的中间系数等于数 $(-4/(-2-\sqrt{3}))-17$；
           先取 $-2-\sqrt{3}$ 为公分母，
           再将分数的分子分母同乘 $-2+\sqrt{3}$ 以使分母有理化）
           \begin{equation*}
             \grstep{((9+4\sqrt{3})/(-2-\sqrt{3}))\rho_2+\rho_3}
             \begin{linsys}{3}
               (-2-\sqrt{3})\cdot b_1
                      &+  &b_2
                          &   &           &= &0  \\
                      &   &(-2-\sqrt{3})\cdot b_2
                          &+  &b_3        &= &0  \\
                      &   &
                          &   &0                     &= &0
             \end{linsys}
           \end{equation*}
           由此得到这个特征子空间和特征向量。
           \begin{equation*}
             V_{2+\sqrt{3}}
             =\set{\colvec{(1/(2+\sqrt{3})^2)\cdot b_3  \\
                           (1/(2+\sqrt{3}))\cdot b_3    \\
                           b_3}
                    \suchthat b_3\in\C}
             \qquad
             \colvec{(1/(2+\sqrt{3})^2)  \\
                           (1/(2+\sqrt{3}))  \\
                           1}
           \end{equation*}

           最后，代入 $x=\lambda_3=2-\sqrt{3}$ 得到方程组
           \begin{multline*}
             \begin{linsys}{3}
               (-2+\sqrt{3})\cdot b_1
                      &+  &b_2
                          &   &           &= &0  \\
                      &   &(-2+\sqrt{3})\cdot b_2
                          &+  &b_3        &= &0  \\
               4\cdot b_1
                      &-  &17\cdot b_2
                          &+  &(6+\sqrt{3})\cdot b_3 &= &0
             \end{linsys}                                       \\
             \begin{aligned}
               &\grstep{(-4/(-2+\sqrt{3}))\rho_1+\rho_3}
                \begin{linsys}{3}
                  (-2+\sqrt{3})\cdot b_1
                         &+  &b_2
                             &   &           &= &0  \\
                         &   &(-2+\sqrt{3})\cdot b_2
                             &+  &b_3        &= &0  \\
                         &   &(-9+4\sqrt{3})\cdot b_2
                             &+  &(6+\sqrt{3})\cdot b_3 &= &0
                \end{linsys}                                       \\
                &\grstep{((9-4\sqrt{3})/(-2+\sqrt{3}))\rho_2+\rho_3}
                \begin{linsys}{3}
                  (-2+\sqrt{3})\cdot b_1
                         &+  &b_2
                             &   &           &= &0  \\
                         &   &(-2+\sqrt{3})\cdot b_2
                             &+  &b_3        &= &0  \\
                         &   &
                             &   &0                     &= &0
                \end{linsys}
              \end{aligned}
           \end{multline*}
           由此得到这个特征子空间和特征向量。
           \begin{equation*}
             V_{2-\sqrt{3}}
             =\set{\colvec{(1/(-2+\sqrt{3})^2)\cdot b_3  \\
                           (1/(2-\sqrt{3}))\cdot b_3    \\
                           b_3}
                    \suchthat b_3\in\C}
             \qquad
             \colvec{(1/(-2+\sqrt{3})^2)  \\
                           (1/(2-\sqrt{3}))  \\
                           1}
           \end{equation*}
       \end{exparts}
     
\end{ans}
\begin{ans}{五.II.3.28}
      \begin{exparts}
        \item 特征多项式可分解为
          $x^2-20x+75=(x-5)(x-15)$，
          因此特征值为 $\lambda_1=5$ 和~$\lambda_2=15$。
          下面是相应的特征子空间。
          \begin{equation*}
            V_5=\set{k\colvec{1 \\ 2}\suchthat k\in\C}
            \quad
            V_{15}=\set{k\colvec{-2 \\ 1}\suchthat k\in\C}
          \end{equation*}
          对每个特征值，代数重数与几何重数都是~$1$。
       \item 特征多项式 $x^3 - 6x^2 + 32$ 可分解为
          $(x+2)(x-4)^2$。
          特征向量为 $\lambda_1=-2$ 和 $\lambda_2=4$。
          下面是相应的特征子空间。
          \begin{equation*}
            V_{-2}=\set{k\colvec{1 \\ 1 \\ 2}\suchthat k\in\C}
            \quad
            V_4=\set{k_1\colvec{1 \\ 1 \\ 0}
                      +k_2\colvec{1 \\ 0 \\ -1} \suchthat k_1,k_2\in\C}
          \end{equation*}
          对于 $\lambda_1=-2$，代数重数与几何重数
          都是~$1$。
          对于 $\lambda_2=4$，代数重数与几何重数
          都是~$2$。
       \item 特征多项式 $x^3 - 5x^2 + 8x - 4$ 可分解为
          $(x-1)(x-2)^2$。
          特征值为 $\lambda_1=1$ 和 $\lambda_2=2$。
          下面是相应的特征子空间。
          \begin{equation*}
            V_1=\set{k\colvec{-6 \\ 3 \\ 1}\suchthat k\in \C}
            \quad
            V_2=\set{k\colvec{1 \\ 0 \\ 0}\suchthat k\in\C}
          \end{equation*}
          对于 $\lambda_1=1$，代数重数与几何重数
          都是~$1$。
          对于 $\lambda_2=2$，代数重数是~$2$，但
          几何重数是~$1$。
      \end{exparts}
    
\end{ans}
\begin{ans}{五.II.3.29}
      关于自然基 $B=\sequence{1,x,x^2}$
      矩阵表示如下。
      \begin{equation*}
        \rep{t}{B,B}
        =
        \begin{mat}[r]
          5  &6  &2  \\
          0  &-1 &-8 \\
          1  &0  &-2
        \end{mat}
      \end{equation*}
      于是特征方程
      \begin{equation*}
        0
        =
        \begin{mat}
          5-x  &6    &2  \\
          0    &-1-x &-8 \\
          1    &0    &-2-x
        \end{mat}
        =(5-x)(-1-x)(-2-x)-48-2\cdot(-1-x)
      \end{equation*}
      为 $0=-x^3+2x^2+15x-36=-1\cdot (x+4)(x-3)^2$。
      为求相应的特征向量，考虑这个方程组。
      \begin{equation*}
        \begin{linsys}{3}
          (5-x)\cdot b_1 &+ &6b_2      &+ &2b_3      &= &0 \\
                         &  &(-1-x)\cdot b_2 &- &8b_3      &= &0 \\
          b_1            &  &                &+ &(-2-x)\cdot b_3 &= &0
        \end{linsys}
      \end{equation*}
      将 $\lambda_1=-4$ 代入 $x$ 得到
      \begin{equation*}
        \begin{linsys}{3}
                    9b_1 &+ &6b_2      &+ &2b_3      &= &0 \\
                         &  &3b_2      &- &8b_3      &= &0 \\
                     b_1 &  &          &+ &2b_3     &= &0
        \end{linsys}
        \grstep{-(1/9)\rho_1+\rho_3}\quad
        \grstep{(2/9)\rho_2+\rho_3}\quad
        \begin{linsys}{3}
                     9b_1 &+ &6b_2      &+ &2b_3      &= &0 \\
                         &  &3b_2       &- &8b_3      &= &0
        \end{linsys}
      \end{equation*}
      特征子空间与一个特征向量如下。
           \begin{equation*}
             V_{-4}
             =\set{\colvec{2\cdot b_3  \\
                           (8/3)\cdot b_3    \\
                           b_3}
                    \suchthat b_3\in\C}
             \qquad
             \colvec{2  \\
                     8/3  \\
                      1}
           \end{equation*}

      同样，代入 $x=\lambda_2=3$ 得
      \begin{equation*}
        \begin{linsys}{3}
                    2b_1 &+ &6\cdot b_2      &+ &2\cdot b_3      &= &0 \\
                         &  &-4\cdot b_2      &- &8\cdot b_3      &= &0 \\
                     b_1 &  &                &- & 5\cdot b_3     &= &0
        \end{linsys}
        \grstep{-(1/2)\rho_1+\rho_3}
        \repeatedgrstep{-(3/4)\rho_2+\rho_3}
        \begin{linsys}{3}
                     2b_1 &+ &6\cdot b_2      &+ &2\cdot b_3      &= &0 \\
                         &  &-4\cdot b_2       &- &8\cdot b_3      &= &0
        \end{linsys}
      \end{equation*}
      特征子空间与特征向量如下。
           \begin{equation*}
             V_{3}
             =\set{\colvec{5\cdot b_3  \\
                           -2\cdot b_3    \\
                           b_3}
                    \suchthat b_3\in\C}
             \qquad
             \colvec[r]{5  \\
                     -2  \\
                       1}
           \end{equation*}
    
\end{ans}
\begin{ans}{五.II.3.30}
        $\lambda=1,
          \begin{mat}
                   0  &0  \\
                   0  &1
          \end{mat} \text{ 和 }
          \begin{mat}[r]
                   2  &3  \\
                   1  &0
          \end{mat}$,
          $\lambda=-2,
          \begin{mat}[r]
                  -1  &0  \\
                   1  &0
          \end{mat}$,
          $\lambda=-1,
          \begin{mat}[r]
                  -2  &1  \\
                   1  &0
          \end{mat}$
       
\end{ans}
\begin{ans}{五.II.3.31}
       取自然基 $B=\sequence{1,x,x^2,x^3}$。
       该映射的作用为 $1\mapsto 0$，$x\mapsto 1$，$x^2\mapsto 2x$，
       $x^3\mapsto 3x^2$，其表示容易计算。
       \begin{equation*}
         T=\rep{d/dx}{B,B}=
         \begin{mat}[r]
           0  &1  &0  &0  \\
           0  &0  &2  &0  \\
           0  &0  &0  &3  \\
           0  &0  &0  &0
         \end{mat}_{B,B}
       \end{equation*}
       我们用如下计算求特征值。
       \begin{equation*}
         0=\deter{T-xI}=
         \begin{vmatrix}
           -x &1  &0  &0  \\
           0  &-x &2  &0  \\
           0  &0  &-x &3  \\
           0  &0  &0  &-x
         \end{vmatrix}
         =x^4
       \end{equation*}
       因此该映射只有唯一特征值 $\lambda=0$。
       为求相应的特征向量，我们求解
       \begin{equation*}
         \begin{mat}[r]
           0  &1  &0  &0  \\
           0  &0  &2  &0  \\
           0  &0  &0  &3  \\
           0  &0  &0  &0
         \end{mat}_{B,B}
         \colvec{b_1 \\ b_2 \\ b_3 \\ b_4}_B
         =0\cdot\colvec{b_1 \\ b_2 \\ b_3 \\ b_4}_B
         \qquad\Longrightarrow\qquad
         \text{$b_2=0$, $b_3=0$, $b_4=0$}
       \end{equation*}
       由此得到这个特征子空间。
       \begin{equation*}
         \set{\colvec{b_1 \\ 0 \\ 0 \\ 0}_B
               \suchthat b_1\in\C}
         =\set{b_1+0\cdot x+0\cdot x^2+0\cdot x^3
               \suchthat b_1\in\C}
         =\set{b_1
               \suchthat b_1\in\C}
       \end{equation*}
      
\end{ans}
\begin{ans}{五.II.3.32}
       三角矩阵 $T-xI$ 的行列式是沿对角线取的乘积，
       因此它可以分解为诸项 $t_{i,i}-x$ 的乘积。
     
\end{ans}
\begin{ans}{五.II.3.33}
      \begin{exparts}
       \partsitem
       计算特征值，得 $\lambda_1=3$，$\lambda_2=0$，
       $\lambda_3=-4$。
       于是该矩阵的对角化如下。
       \begin{equation*}
         \begin{mat}
           3 &0 &0 \\
           0 &0 &0 \\
           0 &0 &-4
         \end{mat}
       \end{equation*}
      \partsitem
       对每个特征值求一个与之相伴的特征向量，并用它们
       构成一个基。

       对 $\lambda_1=3$，我们求解这个方程组。
       \begin{equation*}
         \begin{mat}
           1-3  &2    &1 \\
           6    &-1-3 &0 \\
          -1    &-2   &-1-3
         \end{mat}
         \colvec{x \\ y \\ z}
         =\colvec{0 \\ 0 \\ 0}
       \end{equation*}
       用高斯消元法即可。
       % sage: M = matrix(QQ, [[-2,2,1,0], [6,-4,0,0], [-1,-2,-4,0]])
       % sage: gauss_method(M)
       % [-2  2  1  0]
       % [ 6 -4  0  0]
       % [-1 -2 -4  0]
       % take 3 times row 1 plus row 2
       % take -1/2 times row 1 plus row 3
       % [  -2    2    1    0]
       % [   0    2    3    0]
       % [   0   -3 -9/2    0]
      % take 3/2 times row 2 plus row 3
      % [-2  2  1  0]
      % [ 0  2  3  0]
      % [ 0  0  0  0]
      \begin{equation*}
        \begin{amat}{3}
          -2 &2  &1  &0  \\
           6 &-4 &0  &0  \\
          -1 &-2 &-4 &0
        \end{amat}
        \grstep[-(1/2)\rho_1+\rho_3]{3\rho_1+\rho_2}
        \grstep{(3/2)\rho_1+\rho_3}
        \begin{amat}{3}
          -2 &2 &1   &0  \\
           0 &2 &3   &0  \\
           0 &0 &0   &0
        \end{amat}
      \end{equation*}
      参数化即得解空间。
      \begin{equation*}
        S_1=\set{\colvec{-1 \\ -3/2 \\ 1}\cdot z \suchthat z\in\C}
      \end{equation*}

      这是 $\lambda_2=0$ 的方程组。
      \begin{equation*}
        \begin{mat}
          1    &2    &1 \\
          6    &-1 &0 \\
         -1    &-2   &-1
        \end{mat}
        \colvec{x \\ y \\ z}
        =\colvec{0 \\ 0 \\ 0}
      \end{equation*}
      高斯消元法
      % sage: M = matrix(QQ, [[1,2,1,0], [6,-1,0,0], [-1,-2,-1,0]])
      % sage: gauss_method(M)
      % [ 1  2  1  0]
      % [ 6 -1  0  0]
      % [-1 -2 -1  0]
      % take -6 times row 1 plus row 2
      % take 1 times row 1 plus row 3
      % [  1   2   1   0]
      % [  0 -13  -6   0]
      % [  0   0   0   0]
      \begin{equation*}
        \begin{amat}{3}
           1 &2  &1  &0  \\
           6 &-1 &0  &0  \\
          -1 &-2 &-1 &0
        \end{amat}
        \grstep[\rho_1+\rho_3]{-6\rho_1+\rho_2}
        \begin{amat}{3}
           1 &2 &1   &0  \\
           0 &-13 &-6   &0  \\
           0 &0 &0   &0
        \end{amat}
      \end{equation*}
      再加参数化即得解空间。
      \begin{equation*}
        S_2=\set{\colvec{-1/13 \\ -6/13 \\ 1}\cdot z \suchthat z\in\C}
      \end{equation*}

      $\lambda_3=-4$ 的方程组做法相同。
      \begin{equation*}
        \begin{mat}
          1-(-4)    &2    &1 \\
          6    &-1-(-4) &0 \\
         -1    &-2   &-1-(-4)
        \end{mat}
        \colvec{x \\ y \\ z}
        =\colvec{0 \\ 0 \\ 0}
      \end{equation*}
      高斯消元法
      % sage: M = matrix(QQ, [[5,2,1,0], [6,3,0,0], [-1,-2,3,0]])
      % sage: gauss_method(M)
      % [ 5  2  1  0]
      % [ 6  3  0  0]
      % [-1 -2  3  0]
      % take -6/5 times row 1 plus row 2
      % take 1/5 times row 1 plus row 3
      % [   5    2    1    0]
      % [   0  3/5 -6/5    0]
      % [   0 -8/5 16/5    0]
      % take 8/3 times row 2 plus row 3
      % [   5    2    1    0]
      % [   0  3/5 -6/5    0]
      % [   0    0    0    0]
      \begin{equation*}
        \begin{amat}{3}
           5 &2  &1  &0  \\
           6 &3 &0  &0  \\
          -1 &-2 &3 &0
        \end{amat}
        \grstep[(1/5)\rho_1+\rho_3]{-(6/5)\rho_1+\rho_2}
        \grstep{(8/3)\rho_2+\rho_3}
        \begin{amat}{3}
           5 &2 &1   &0  \\
           0 &3/5 &-6/5   &0  \\
           0 &0 &0   &0
        \end{amat}
      \end{equation*}
      参数化后得到
      \begin{equation*}
        S_3=\set{\colvec{-1 \\ 2 \\ 1}\cdot z \suchthat z\in\C}
      \end{equation*}

      于是，当该矩阵表示同一个变换、但关于下面这个 $D$
      取 $D,D$ 为基时，它就具有那个对角形状。
      \begin{equation*}
        D=\sequence{\colvec{-1 \\ -3/2 \\ 1},
                    \colvec{-1/13 \\ -6/13 \\ 1},
                    \colvec{-1 \\ 2 \\ 1}}
      \end{equation*}
     \partsitem
     % sage: M = matrix(QQ, [[1,2,1], [6,-1,0], [-1,-2,-1]])
     % sage: M.eigenvalues()
     % [3, 0, -4]
     % sage: P = matrix(QQ, [[-1,-1/13,-1], [-3/2,-6/13,2], [1,1,1]])
     % sage: P.inverse()*M*P
     % [ 3  0  0]
     % [ 0  0  0]
     % [ 0  0 -4]
     % sage: P.inverse()
     % [-16/21   -2/7  -4/21]
     % [ 13/12      0  13/12]
     % [ -9/28    2/7   3/28]
       将该矩阵记为~$T$。
       取定义域空间与陪域空间为 $\Re^3$，
       并取矩阵为 $T=\rep{t}{\stdbasis_3,\stdbasis_3}$，
       便得到下图。
       \begin{equation*}
       \begin{CD}
          \C^3_{\wrt{\stdbasis_3}}            @>t>T>      \C^3_{\wrt{\stdbasis_3}}       \\
          @V{\scriptstyle\identity} VV           @V{\scriptstyle\identity} VV \\
          \C^3_{\wrt{B}}             @>t>D>  \C^3_{\wrt{B}}
       \end{CD}
       \end{equation*}
       从图中可读出 $D=P^{-1}MP$，其中
       $P=\rep{\identity}{\stdbasis_3,B}$。
       用~$E_3$ 表示 $B$ 中的每个向量很容易，
       因为关于标准基，每个向量的表示就是它自身。
       \begin{multline*}
         \rep{\colvec{-1 \\ -3/2 \\ 1}}{\stdbasis_3}=\colvec{-1 \\ -3/2 \\ 1}
         \quad
         \rep{\colvec{-1/13 \\ -6/13 \\ 1}}{\stdbasis_3}=\colvec{-1/13 \\ -6/13 \\ 1}
         \\
         \rep{\colvec{-1 \\ 2 \\ 1}}{\stdbasis_3}=\colvec{-1 \\ 2 \\ 1}
       \end{multline*}
       于是该矩阵就是把这三个基向量依次并排放置。
       \begin{equation*}
         P=
         \begin{mat}
           -1   &-1/13 &-1 \\
           -3/2 &-6/13 &2  \\
           1    &1     &1
         \end{mat}
        \quad\text{于是有}\quad
        P^{-1}
        \begin{mat}
          -16/21 &-2/7 &-4/21 \\
          13/12  &0    &13/12 \\
          -9/28  &2/7  &3/28
        \end{mat}
       \end{equation*}





     \end{exparts}
    
\end{ans}
\begin{ans}{五.II.3.34}
      只需展开 $T-xI$ 的行列式。
      \begin{equation*}
        \begin{vmatrix}
          a-x  &c  \\
          b    &d-x
        \end{vmatrix}
        =(a-x)(d-x)-bc
        =x^2+(-a-d)\cdot x +(ad-bc)
      \end{equation*}
    
\end{ans}
\begin{ans}{五.II.3.35}
      该变换的任意两个表示都是相似的，
      而相似矩阵具有相同的特征多项式。
    
\end{ans}
\begin{ans}{五.II.3.36}
      这不成立。
      这个矩阵的所有特征值都是 $0$。
      \begin{equation*}
        \begin{mat}[r]
          0  &1  \\
          0  &0
        \end{mat}
      \end{equation*}
    
\end{ans}
\begin{ans}{五.II.3.37}
      \begin{exparts}
        \partsitem 取 \( \lambda=1 \) 与恒等映射即可。
        \partsitem 可以，使用把所有向量都乘以
          标量 \( \lambda \) 的那个变换。
      \end{exparts}
     
\end{ans}
\begin{ans}{五.II.3.38}
      若 $t(\vec{v})=\lambda\cdot\vec{v}$，则在映射 $t-\lambda\cdot\identity$
      下 $\vec{v}\mapsto\zero$。
    
\end{ans}
\begin{ans}{五.II.3.39}
      特征方程
      \begin{equation*}
        0=
        \begin{vmatrix}
          a-x  &b  \\
          c    &d-x
        \end{vmatrix}
        =(a-x)(d-x)-bc
      \end{equation*}
      化简为 $x^2+(-a-d)\cdot x + (ad-bc)$。
      验证 $x=a+b$ 与 $x=a-c$ 满足该方程
      （在 $a+b=c+d$ 的条件下）是容易的。
    
\end{ans}
\begin{ans}{五.II.3.40}
      考虑特征子空间 $V_{\lambda}$。
      任意 $\vec{w}\in V_{\lambda}$ 都是某个 $\vec{v}\in V_{\lambda}$
      的像 $\vec{w}=\lambda\cdot\vec{v}$
      （即 $\vec{v}=(1/\lambda)\cdot\vec{w}$）。
      于是，在 $V_{\lambda}$（它是一个非平凡子空间）上
      $t^{-1}$ 的作用为
      $t^{-1}(\vec{w})=\vec{v}=(1/\lambda)\cdot\vec{w}$，
      因此 $1/\lambda$ 是 $t^{-1}$ 的特征值。
    
\end{ans}
\begin{ans}{五.II.3.41}
      \begin{exparts}
        \partsitem 我们有
          $(cT+dI)\vec{v}=cT\vec{v}+dI\vec{v}=c\lambda\vec{v}+d\vec{v}
               =(c\lambda+d)\cdot \vec{v}$。
        \partsitem 设 $S=PTP^{-1}$ 是对角的。
          则 $P(cT+dI)P^{-1}=P(cT)P^{-1}+P(dI)P^{-1}
                 =cPTP^{-1}+dI=cS+dI$ 也是对角的。
      \end{exparts}
    
\end{ans}
\begin{ans}{五.II.3.42}
      标量 $\lambda$ 是特征值当且仅当变换
      $t-\lambda \identity$ 是奇异的。
      变换是奇异的当且仅当它不是同构
      （也就是说，变换是同构当且仅当它是
      非奇异的）。
    
\end{ans}
\begin{ans}{五.II.3.43}
      \begin{exparts}
        \partsitem 当特征值 $\lambda$ 与特征向量 $\vec{x}$ 相伴时，
          $A^k\vec{x}=A\cdots A\vec{x}=A^{k-1}\lambda\vec{x}
            =\lambda A^{k-1}\vec{x}=\cdots=\lambda^k\vec{x}$。
          （完整的细节需要对 $k$ 作归纳。）
        \partsitem 与 $\lambda$ 相伴的特征向量
          未必是与 $\mu$ 相伴的特征向量。
      \end{exparts}
    
\end{ans}
\begin{ans}{五.II.3.44}
      没有。
      这是两个同尺寸、同秩但特征值不同的矩阵。
      \begin{equation*}
        \begin{mat}[r]
          1  &0  \\
          0  &1
        \end{mat}
        \qquad
        \begin{mat}[r]
          1  &0  \\
          0  &2
        \end{mat}
      \end{equation*}
    
\end{ans}
\begin{ans}{五.II.3.45}
      特征多项式具有奇数次幂，因此
      至少有一个实根。
    
\end{ans}
\begin{ans}{五.II.3.46}
      特征多项式 $x^3+5x^2+6x$ 有互不相同的根
      \( \lambda_1=0 \)、\( \lambda_2=-2 \)、\( \lambda_3=-3 \)。
      于是该矩阵可以对角化成这个形状。
      \begin{equation*}
         \begin{mat}[r]
            0  &0  &0  \\
            0  &-2 &0  \\
            0  &0  &-3
         \end{mat}
      \end{equation*}
    
\end{ans}
\begin{ans}{五.II.3.47}
      我们必须证明它是单射且满射，并且它保持
      矩阵加法与标量乘法运算。

      为证它是单射，设 $t_P(T)=t_P(S)$，
      即设 $PTP^{-1}=PSP^{-1}$，注意两边
      左乘 $P^{-1}$、右乘 $P$ 便得
      $T=S$。
      为证它是满射，考虑 $S\in\matspace_{\nbyn{n}}$，并观察
      $S=t_P(P^{-1}SP)$。

      映射 $t_P$ 保持矩阵加法，因为
      $t_P(T+S)=P(T+S)P^{-1}=(PT+PS)P^{-1}=PTP^{-1}+PSP^{-1}=t_P(T+S)$
      可由我们已见过的矩阵乘法与加法的性质
      得到。
      标量乘法是
      类似的：~$t_P(cT)=P(c\cdot T)P^{-1}=c\cdot (PTP^{-1})=c\cdot t_P(T)$。
    
\end{ans}
\begin{ans}{五.II.3.48}
      \answerasgiven %
      若把 \( A \) 的特征函数中的变元取为
      \( c \)，把前 \( (n-1) \) 行（列）加到
      第 \( n \) 行（列）上，就得到一个第 \( n \) 行
      （列）全为零的行列式。
      因此 \( c \) 是 \( A \) 的一个特征根。
    
\end{ans}
\protect \section {第 III 节：幂零性}
\protect \subsection {五.III.1: 自复合}
\begin{ans}{五.III.1.9}
      对于零变换，无论空间是什么，值域空间链都是
      \( V\supset\set{\vec{0}}=\set{\vec{0}}=\cdots\, \)，
      而零空间链是 \( \set{\vec{0}}\subset V=V=\cdots\, \)。
      对于恒等变换，两条链分别是
      \( V=V=V=\cdots \) 与
      \( \set{\vec{0}}=\set{\vec{0}}=\cdots\, \)。
    
\end{ans}
\begin{ans}{五.III.1.10}
       \begin{exparts}
         \partsitem 将 $t_0$ 迭代两次
           $a+bx+cx^2\mapsto b+cx^2\mapsto cx^2$
           得到
           \begin{equation*}
             a+bx+cx^2\mapsunder{t_0^2}cx^2
           \end{equation*}
           而任何更高的迭代都是同一个映射。
           于是，虽然 $\rangespace{t_0}$ 是没有一次项的
           二次多项式构成的空间
           $\set{p+rx^2\suchthat p,r\in \C}$，
           且
           $\rangespace{t_0^2}$ 是纯二次多项式构成的空间
           $\set{rx^2\suchthat r\in \C}$，
           但链就在这里稳定下来：
           $\genrangespace{t_0}=\set{rx^2\suchthat r\in \C}$。

           至于零空间，
           $\nullspace{t_0}$ 是仅含一次项的
           多项式构成的空间 $\set{qx\suchthat q\in \C}$，
           而
           $\nullspace{t_0^2}$ 是一次多项式构成的空间
           $\set{p+qx\suchthat p,q\in \C}$。
           到此为止：$\gennullspace{t_0}=\nullspace{t_0^2}$。
         \partsitem 二次幂
           \begin{equation*}
             \colvec{a \\ b}
              \mapsunder{t_1}\colvec{0 \\ a}
              \mapsunder{t_1}\colvec{0 \\ 0}
           \end{equation*}
           是零映射。
           因此，值域空间链
           \begin{equation*}
             \Re^2
               \supset\set{\colvec{0 \\ p}\suchthat p\in\C}
               \supset\set{\zero}
               =\cdots
           \end{equation*}
           与零空间链
           \begin{equation*}
             \set{\zero}
               \subset\set{\colvec{q \\ 0}\suchthat q\in\C}
               \subset\Re^2
               =\cdots
           \end{equation*}
           的长度都是二。
           广义值域空间是平凡子空间，而
           广义零空间是整个空间。
         \partsitem 这个映射的迭代循环往复
           \begin{equation*}
             a+bx+cx^2
               \mapsunder{t_2} b+cx+ax^2
               \mapsunder{t_2} c+ax+bx^2
               \mapsunder{t_2} a+bx+cx^2
               \;\cdots
           \end{equation*}
           并且值域空间链与零空间链都是平凡的。
           \begin{equation*}
             \polyspace_2=\polyspace_2=\cdots
              \qquad
              \set{\zero}=\set{\zero}=\cdots
           \end{equation*}
           于是，显然，
           广义空间是 $\genrangespace{t_2}=\polyspace_2$
           与 $\gennullspace{t_2}=\set{\zero}$。
         \partsitem 我们有
           \begin{equation*}
             \colvec{a \\ b \\ c}
                \mapsto\colvec{a \\ a \\ b}
                \mapsto\colvec{a \\ a \\ a}
                \mapsto\colvec{a \\ a \\ a}
                \mapsto\cdots
           \end{equation*}
           于是值域空间链
           \begin{equation*}
             \Re^3
               \supset\set{\colvec{p \\ p \\ r}\suchthat p,r\in\C}
               \supset\set{\colvec{p \\ p \\ p}\suchthat p\in\C}
               =\cdots
           \end{equation*}
           以及零空间链
           \begin{equation*}
             \set{\zero}
                \subset\set{\colvec{0 \\ 0 \\ r}\suchthat r\in \C}
                \subset\set{\colvec{0 \\ q \\ r}\suchthat q,r\in \C}
                =\cdots
           \end{equation*}
           的长度都是二。
           广义空间就是上面每条链中最后所示的那个空间。
       \end{exparts}
      
\end{ans}
\begin{ans}{五.III.1.11}
      两者都把 \( x\mapsto t(t(t(x))) \)。
    
\end{ans}
\begin{ans}{五.III.1.12}
      回忆一下，若 $W$ 是 $V$ 的子空间，则我们可以把
      $W$ 的任意一组基 $B_W$
      扩充为 $V$ 的一组基 $B_V$。
      由此第一句话立即可得。
      第二句话也不难：~$W$ 是 $B_W$ 的张成空间，
      若 $W$ 是真子空间，则 $V$ 不是 $B_W$ 的张成空间，
      所以 $B_V$ 至少比 $B_W$ 多一个向量。
    
\end{ans}
\begin{ans}{五.III.1.13}
      “若”与“仅当”都成立。
      线性映射是非奇异的
      当且仅当它保持维数，也就是说，其值域的维数
      等于其定义域的维数。
      对于变换 $\map{t}{V}{V}$，这意味着
      该映射是非奇异的当且仅当它是满射：
      $\rangespace{t}=V$（从而 $\rangespace{t^2}=V$，等等）。
    
\end{ans}
\begin{ans}{五.III.1.14}
      零空间构成链，因为
      若 $\vec{v}\in\nullspace{t^j}$，则 $t^j(\vec{v})=\zero$，
      且 $t^{j+1}(\vec{v})=t(\,t^j(\vec{v})\,)=t(\zero)=\zero$，
      于是 $\vec{v}\in\nullspace{t^{j+1}}$。

      现在，
      零空间的“进一步”性质由它对值域空间成立这一事实
      结合前面的练习即可得出。
      由于 $\rangespace{t^j}$ 的维数加上
      $\nullspace{t^j}$ 的维数等于起始空间~$V$ 的维数~$n$，
      当值域空间的维数停止下降时，
      零空间的维数也不再变化。
      前面的练习表明，从这一点~$k$ 开始，
      链中的包含关系都不是真包含\Dash 零空间
      彼此相等。
    
\end{ans}
\begin{ans}{五.III.1.15}
      （许多例子都正确，这里给出一个。）
      一个例子是实三元组上的平移算子
      \( (x,y,z)\mapsto (0,x,y) \)。
      零空间是所有以两个零开头的三元组。
      该映射在三次迭代后稳定。
     
\end{ans}
\begin{ans}{五.III.1.16}
        求微分算子
        \( \map{d/dx}{\polyspace_1}{\polyspace_1} \) 的值域空间与零空间相同。
        要看一个二者不相交\Dash
        除零向量外\Dash 的例子，
        可考虑恒等映射，或任意非奇异映射。
      
\end{ans}
\protect \subsection {五.III.2: 链}
\begin{ans}{五.III.2.20}
         是三。
         说它至少是三，是因为 $\ell^2(\,(1,1,1)\,)=(0,0,1)\neq \zero$。
         说它至多是三，是因为
         $(x,y,z)\mapsto (0,x,y)\mapsto (0,0,x)\mapsto (0,0,0)$。
       
\end{ans}
\begin{ans}{五.III.2.21}
       \begin{exparts}
         \partsitem 定义域的维数是四。
           该映射的作用是：空间中任意向量
           $c_1\cdot \vec{\beta}_1+c_2\cdot \vec{\beta}_2
             +c_3\cdot \vec{\beta}_3+c_4\cdot \vec{\beta}_4$
           被送到
           $c_1\cdot \vec{\beta}_2+c_2\cdot \zero
             +c_3\cdot \vec{\beta}_4+c_4\cdot \zero
            =c_1\cdot \vec{\beta}_3+c_3\cdot\vec{\beta}_4$。
           映射的第一次作用把两个基向量 $\vec{\beta}_2$ 与
           $\vec{\beta}_4$ 送到零，
           因此零空间的维数为~二，值域空间的维数也为~二。
           第二次作用后，四个基向量全部被送到零，
           于是二次幂的零空间维数为~四，而二次幂的值域空间维数为零。
           因此幂零指数是二。
           这就是标准形。
           \begin{equation*}
             \begin{mat}[r]
               0  &0  &0  &0  \\
               1  &0  &0  &0  \\
               0  &0  &0  &0  \\
               0  &0  &1  &0
             \end{mat}
           \end{equation*}
         \partsitem 该映射定义域的维数是六。
           对一次幂而言，零空间的维数是四，值域空间的维数是二。
           对二次幂而言，零空间的维数是五，值域空间的维数是一。
           第三次迭代得到的零空间维数为六，值域空间维数为零。
           幂零指数是三，这就是标准形。
           \begin{equation*}
             \begin{mat}[r]
               0  &0  &0  &0  &0  &0  \\
               1  &0  &0  &0  &0  &0  \\
               0  &1  &0  &0  &0  &0  \\
               0  &0  &0  &0  &0  &0  \\
               0  &0  &0  &0  &0  &0  \\
               0  &0  &0  &0  &0  &0
             \end{mat}
           \end{equation*}
         \partsitem 定义域的维数是三，幂零指数是三。
           一次幂的零空间维数为一，值域空间维数为二。
           二次幂的零空间维数为二，值域空间维数为一。
           最后，三次幂的零空间维数为三，值域空间维数为零。
           这是标准形矩阵。
           \begin{equation*}
             \begin{mat}[r]
               0  &0  &0  \\
               1  &0  &0  \\
               0  &1  &0
             \end{mat}
          \end{equation*}
      \end{exparts}
    
\end{ans}
\begin{ans}{五.III.2.22}
       由 \nearbylemma{le:RangeAndNullChains}，零度会在第 $n$ 次迭代时增长到尽可能大，其中 $n$ 是定义域的维数。
       于是，对 $\nbyn{2}$ 矩阵，
       我们只需检查其平方是否为零矩阵。
       对 $\nbyn{3}$ 矩阵，只需检查其立方。
       \begin{exparts}
         \partsitem 是，该矩阵是幂零的，因为
           它的平方是零矩阵。
         \partsitem 否，它的平方不是零矩阵。
           \begin{equation*}
             \begin{mat}[r]
                3  &1  \\
                1  &3
             \end{mat}^2
             =\begin{mat}[r]
                10  &6  \\
                6   &10
             \end{mat}
           \end{equation*}
         \partsitem 是，它的立方是零矩阵。
           事实上，它的平方就是零。
         \partsitem 否，它的三次幂不是零矩阵。
           \begin{equation*}
             \begin{mat}[r]
               1  &1  &4  \\
               3  &0  &-1 \\
               5  &2  &7
             \end{mat}^3
             =\begin{mat}[r]
               206  &86  &304  \\
                26  &8   &26   \\
               438  &180 &634
             \end{mat}
           \end{equation*}
         \partsitem 是，该矩阵的立方是零矩阵。
       \end{exparts}
       要看出第二与第四个矩阵不是幂零的，另一种方法是注意它们是非奇异的。
     
\end{ans}
\begin{ans}{五.III.2.23}
 计算表格
       \begin{center}
           \begin{tabular}{c|cc}
              \multicolumn{1}{c}{\( p \)}  &\( N^p \)  &\( \nullspace{N^p}  \) \\
              \hline
              \( 1 \)
                &\matrixvenlarge{\begin{mat}[r]
                    0  &1  &1  &0  &1  \\
                    0  &0  &1  &1  &1  \\
                    0  &0  &0  &0  &0  \\
                    0  &0  &0  &0  &0  \\
                    0  &0  &0  &0  &0
                   \end{mat}}
                &\( \set{\matrixvenlarge{\colvec{r \\ u \\ -u-v \\ u \\ v}}
                               \suchthat r,u,v\in\C}  \) \\
              \( 2 \)
                &\matrixvenlarge{\begin{mat}[r]
                    0  &0  &1  &1  &1  \\
                    0  &0  &0  &0  &0  \\
                    0  &0  &0  &0  &0  \\
                    0  &0  &0  &0  &0  \\
                    0  &0  &0  &0  &0
                   \end{mat}}
                &\( \set{\matrixvenlarge{\colvec{r \\ s \\ -u-v \\ u \\ v}}
                               \suchthat r,s,u,v\in\C}  \) \\
             \( 2 \)
                &\textit{--零矩阵--}
                &\( \C^5 \)
           \end{tabular}
         \end{center}
         由此得到对链基的如下要求：~有三个基向量直接被映射到零，
         另有一个基向量经第二次作用被映射到零，
         最后一个基向量
         则经第三次作用被映射到零。
         因此，链基具有如下形式。
         \begin{equation*}
           \begin{strings}{ccccccc}
             \vec{\beta}_1 &\mapsto &\vec{\beta}_2 &\mapsto
               &\vec{\beta}_3 &\mapsto &\zero               \\
             \vec{\beta}_4 &\mapsto &\zero                  \\
             \vec{\beta}_5 &\mapsto &\zero
           \end{strings}
         \end{equation*}
         据此立得标准形。
         \begin{equation*}
           \begin{mat}[r]
             0  &0  &0  &0  &0  \\
             1  &0  &0  &0  &0  \\
             0  &1  &0  &0  &0  \\
             0  &0  &0  &0  &0  \\
             0  &0  &0  &0  &0
           \end{mat}
         \end{equation*}
     
\end{ans}
\begin{ans}{五.III.2.24}
       \begin{exparts}
         \partsitem 标准形含一个 $\nbyn{3}$ 块与一个
           $\nbyn{2}$ 块
           \begin{equation*}
             \begin{pmat}{rrr|rr}
               0  &0  &0  &0  &0  \\
               1  &0  &0  &0  &0  \\
               0  &1  &0  &0  &0  \\ \hline
               0  &0  &0  &0  &0  \\
               0  &0  &0  &1  &0  \\
             \end{pmat}
           \end{equation*}
           分别对应于基中长度为三的链与长度为二的
           链。
         \partsitem 设 $N$ 是底层映射关于标准基的表示。
           设 $B$ 是我们即将换到的基。
           由相似图
           \begin{equation*}
             \begin{CD}
               \C^2_{\wrt{\stdbasis_2}}
                   @>n>N>
                   \C^2_{\wrt{\stdbasis_2}}     \\
              @V\scriptstyle\identity V\scriptstyle PV
                  @V\scriptstyle\identity V\scriptstyle PV \\
              C^2_{\wrt{B}}
                  @>n>>
              \C^2_{\wrt{B}}
            \end{CD}
          \end{equation*}
          可知标准形矩阵为 $PNP^{-1}$，其中
          \begin{equation*}
            P^{-1}
            =\rep{\identity}{B,\stdbasis_5}
            =\begin{mat}[r]
               1 &0 &0 &0 &0 \\
               0 &1 &0 &1 &0 \\
               1 &0 &1 &0 &0 \\
               0 &0 &1 &1 &1 \\
               0 &0 &0 &0 &1
             \end{mat}
          \end{equation*}
          而 $P$ 是它的逆。
          \begin{equation*}
            P=\rep{\identity}{\stdbasis_5,B}
             =(P^{-1})^{-1}
            =\begin{mat}[r]
               1 &0 &0 &0 &0 \\
              -1 &1 &1 &-1&1 \\
              -1 &0 &1 &0 &0 \\
               1 &0 &-1&1 &-1\\
               0 &0 &0 &0 &1
             \end{mat}
          \end{equation*}
         \partsitem 检验它的计算是例行的。
       \end{exparts}
     
\end{ans}
\begin{ans}{五.III.2.25}
       \begin{exparts}
       \partsitem 计算
         \begin{center}
           \begin{tabular}{c|cc}
              \multicolumn{1}{c}{\( p \)}  &\( N^p \)  &\( \nullspace{N^p} \) \\
             \hline
              \( 1 \)
                &\matrixvenlarge{\begin{mat}[r]
                     1/2  &-1/2  \\
                     1/2  &-1/2
                   \end{mat}}
                &\( \set{\matrixvenlarge{\colvec{u \\ u}}
                               \suchthat u\in\C}  \) \\
             \( 2 \)
                &\textit{--零矩阵--}
                &\( \C^2 \)
           \end{tabular}
         \end{center}
         表明由该矩阵表示的任何映射
         必须这样作用在链基上
         \begin{equation*}
           \begin{strings}{ccccccc}
             \vec{\beta}_1 &\mapsto &\vec{\beta}_2  &\mapsto &\zero
           \end{strings}
         \end{equation*}
         因为作用一次后零空间的维数为~一，
         且恰好有一个基向量 $\vec{\beta}_2$ 被映射到零。
         因此，这个关于
         $\sequence{\vec{\beta}_1,\vec{\beta}_2}$ 的表示就是标准形。
         \begin{equation*}
           \begin{mat}[r]
             0    &0   \\
             1    &0
           \end{mat}
         \end{equation*}
       \partsitem 这里的计算与前一题类似。
         \begin{center}
           \begin{tabular}{c|cc}
              \multicolumn{1}{c}{\( p \)}  &\( N^p \)  &\( \nullspace{N^p} \) \\
              \hline
              \( 1 \)
                &\matrixvenlarge{\begin{mat}[r]
                     0  &0  &0  \\
                     0  &-1 &1  \\
                     0  &-1 &1
                   \end{mat}}
                &\( \set{\matrixvenlarge{\colvec{u \\ v \\ v}}
                               \suchthat u,v\in\C}  \) \\
            \( 2 \)
                &\textit{--零矩阵--}
                &\( \C^3 \)
           \end{tabular}
        \end{center}
        表格表明链基形如
         \begin{equation*}
           \begin{strings}{ccccccc}
             \vec{\beta}_1 &\mapsto &\vec{\beta}_2 &\mapsto &\zero  \\
             \vec{\beta}_3 &\mapsto &\zero
           \end{strings}
         \end{equation*}
         因为映射作用一次后零空间的
         维数为二\Dash
         --$\vec{\beta}_2$ 与 $\vec{\beta}_3$ 都被送到零\Dash 而
         再迭代一次又使另一个向量
         被送到零。
       \partsitem 计算
         \begin{center}
           \begin{tabular}{c|cc}
              \multicolumn{1}{c}{\( p \)}  &\( N^p \)  &\( \nullspace{N^p} \) \\
              \hline
              \( 1 \)
                &\matrixvenlarge{\begin{mat}[r]
                     -1  &1  &-1  \\
                      1  &0  &1   \\
                      1  &-1 &1
                   \end{mat}}
                &\( \set{\matrixvenlarge{\colvec{u \\ 0 \\ -u}}
                               \suchthat u\in\C}  \) \\
              \( 2 \)
                &\matrixvenlarge{\begin{mat}[r]
                      1  &0  &1   \\
                      0  &0  &0   \\
                     -1  &0  &-1
                   \end{mat}}
                &\( \set{\matrixvenlarge{\colvec{u \\ v \\ -u}}
                               \suchthat u,v\in\C}  \) \\
             \( 3 \)
                &\textit{--零矩阵--}
                &\( \C^3 \)
           \end{tabular}
         \end{center}
         表明由该基表示的任何映射必须这样作用在
         链基上。
         \begin{equation*}
           \begin{strings}{ccccccc}
             \vec{\beta}_1 &\mapsto &\vec{\beta}_2 &\mapsto
               &\vec{\beta}_3 &\mapsto &\zero
           \end{strings}
         \end{equation*}
         因此，这就是标准形。
         \begin{equation*}
           \begin{mat}[r]
             0    &0   &0   \\
             1    &0   &0   \\
             0    &1   &0
           \end{mat}
         \end{equation*}
       \end{exparts}
      
\end{ans}
\begin{ans}{五.III.2.26}
       下面几个例子
       \begin{equation*}
         \begin{mat}[r]
           0  &0  \\
           1  &0
         \end{mat}
         \begin{mat}
           a  &b  \\
           c  &d
         \end{mat}
         =
         \begin{mat}
           0  &0  \\
           a  &b
         \end{mat}
         \qquad
         \begin{mat}[r]
           0  &0  &0 \\
           1  &0  &0 \\
           0  &1  &0
         \end{mat}
         \begin{mat}
           a  &b  &c \\
           d  &e  &f \\
           g  &h  &i
         \end{mat}
         =
         \begin{mat}
           0  &0  &0 \\
           a  &b  &c \\
           d  &e  &f
         \end{mat}
       \end{equation*}
       表明用一块次对角线上的 $1$ 左乘一个矩阵，
       会把该矩阵的行向下移。
       而不同的块
       \begin{equation*}
         \begin{mat}[r]
           0  &0  &0  &0  \\
           1  &0  &0  &0  \\
           0  &0  &0  &0  \\
           0  &0  &1  &0
         \end{mat}
         \begin{mat}
           a  &b  &c  &d  \\
           e  &f  &g  &h  \\
           i  &j  &k  &l  \\
           m  &n  &o  &p
         \end{mat}
         =
         \begin{mat}
           0  &0  &0  &0  \\
           a  &b  &c  &d  \\
           0  &0  &0  &0  \\
           i  &j  &k  &l
         \end{mat}
       \end{equation*}
       则会把矩阵的不同部分向下移。

       右乘对列做的是类似的事情。
       参见 \nearbyexercise{exer:IndNilLftShift}。
     
\end{ans}
\begin{ans}{五.III.2.27}
       是。
       把 \nearbyexample{ex:NilMatNotCanon} 的最后一句推广即可。
       至于指数，同样的最后一句表明新矩阵的指数小于或等于 $\hat{N}$ 的指数，而把
       两个矩阵的角色互换便得到另一方向的不等式。

       这个问题的另一种回答方式是证明：一个矩阵是幂零的当且仅当与之相关的映射是幂零的，且
       指数相同。
       于是，由于相似矩阵表示同一个映射，结论
       随之而来。
       这就是下面的 \nearbyexercise{exer:MatNilIffMapNil}。
     
\end{ans}
\begin{ans}{五.III.2.28}
       注意标准形的幂零矩阵只有
       零特征值；例如，这个下三角矩阵
       \begin{equation*}
          \begin{mat}
            -x  &0  &0  \\
             1  &-x &0  \\
             0  &1  &-x
          \end{mat}
       \end{equation*}
       的行列式是 \( (-x)^3 \)，它的唯一的根是零。
       而相似矩阵有相同的特征值，且每个幂零
       矩阵都与一个标准形的矩阵相似。

       看出这一点的另一种方式是注意：幂零矩阵在若干次迭代后把所有
       向量都送到零，但这与在特征子空间上的作用 $\vec{v}\mapsto \lambda\vec{v}$
       相矛盾，除非
       $\lambda$ 为零。
     
\end{ans}
\begin{ans}{五.III.2.29}
       不存在，由 \nearbylemma{le:RangeAndNullChains}，对二维
       空间上的映射，零度
       在第二次迭代时就会增长到尽可能大。
     
\end{ans}
\begin{ans}{五.III.2.30}
       一个变换的幂零指数为零，仅当开启链的
       向量
       必须是 $\zero$ 时，也就是说，仅当 $V$ 是平凡空间时。
     
\end{ans}
\begin{ans}{五.III.2.31}
       \begin{exparts}
         \partsitem 张成空间中的任意成员 $\vec{w}$ 都是
           一个线性组合
           $\vec{w}=c_0\cdot \vec{v}+c_1\cdot t(\vec{v})+\dots
            +c_{k-1}\cdot t^{k-1}(\vec{v})$。
           于是，由映射的线性性，
           $t(\vec{w})=c_0\cdot t(\vec{v})+c_1\cdot t^2(\vec{v})+\dots
            +c_{k-2}\cdot t^{k-1}(\vec{v})+c_{k-1}\cdot \zero$
           也在这个张成空间中。
         \partsitem 前一小题中的操作迭代 $k$ 次后，
           会得到零的线性组合。
         \partsitem 若 \( \vec{v}=\zero \)，则该集合是空集，按定义
           线性无关。
           否则，写 \( c_1\vec{v}+\dots+c_{k-1}t^{k-1}(\vec{v})=\zero \)
           并对两边施加 \( t^{k-1} \)。
           右边得到 \( \zero \)，而左边得到
           \( c_1t^{k-1}(\vec{v}) \)；于是 \( c_1=0 \)。
           如此继续，对两边施加 \( t^{k-2} \)，
           依此类推。
         \partsitem 当然，$t$ 把这条长为 $k$ 的单一 $t$-链作为基作用在其张成空间上。
           \begin{equation*}
             \begin{mat}
               0 &0 &0 &0 &\ldots &0 &0 \\
               1 &0 &0 &0 &\ldots &0 &0 \\
               0 &1 &0 &0 &\ldots &0 &0 \\
               0 &0 &1 &0 &       &0 &0 \\
                 &  &  &\ddots          \\
               0 &0 &0 &0 &       &1 &0 \\
             \end{mat}
           \end{equation*}
       \end{exparts}
     
\end{ans}
\begin{ans}{五.III.2.32}
       我们必须证明
       \( B\union\hat{C}\union\set{\vec{v}_1,\dots,\vec{v}_j} \) 线性
       无关，其中 \( B \) 是
       \( \rangespace{t} \) 的 \( t \)-链基，\( \hat{C} \) 是
       \( \nullspace{t} \) 的基，且
       \( t(\vec{v}_1)=\vec{\beta}_1,\dots,t(\vec{v}_i)=\vec{\beta}_i \)。
       写出
       \begin{equation*}
         \zero=c_{1,-1}\vec{v}_1+c_{1,0}\vec{\beta}_1
                +c_{1,1}t(\vec{\beta}_1)+\dots+
                c_{1,{h_1}}t^{h_1}(\vec{\vec{\beta}}_1)
                +c_{2,-1}\vec{v}_2+\dots+c_{j,h_i}t^{h_i}(\vec{\beta_i})
       \end{equation*}
       并施加 \( t \)。
       \begin{multline*}
         \zero=c_{1,-1}\vec{\beta}_1+c_{1,0}t(\vec{\beta}_1)
                +\dots+
                c_{1,h_1-1}t^{h_1}(\vec{\vec{\beta}}_1)+c_{1,h_1}\zero      \\
             +c_{2,-1}\vec{\beta}_2+\cdots+c_{i,h_i-1}t^{h_i}(\vec{\beta_i})
             +c_{i,h_i}\zero
       \end{multline*}
       由于 \( B\union\hat{C} \)
       是一组基，可断定系数 \( c_{1,-1},\dots,c_{1,h_i-1},
       c_{2,-1},\dots,c_{i,h_i-1} \) 全为零。
       代回第一个行间公式，可知
       其余系数也为零。
     
\end{ans}
\begin{ans}{五.III.2.33}
       对任意基 $B$，
       变换~$n$ 是幂零的当且仅当
       $N=\rep{n}{B,B}$ 是幂零矩阵。
       这是因为只有零矩阵才表示零映射，
       于是 \( n^j \) 是零映射当且仅当 \( N^j \)
       是零矩阵。
     
\end{ans}
\begin{ans}{五.III.2.34}
       它可以是任何大于或等于一的维数。
       要构造一个指数为四的幂零变换，
       其立方有维数为~$k$ 的值域空间，取一个向量空间、
       该空间的一组基，以及一个按如下方式
       作用于这组基的变换。
       \begin{equation*}
          \begin{strings}{ccccccccc}
             \vec{\beta}_1 &\mapsto &\vec{\beta}_2 &\mapsto &\vec{\beta}_3
               &\mapsto &\vec{\beta}_4 &\mapsto &\zero  \\
             \vec{\beta}_5 &\mapsto &\vec{\beta}_6 &\mapsto &\vec{\beta}_7
               &\mapsto &\vec{\beta}_8 &\mapsto &\zero  \\
             &\vdotswithin{\vec{\beta}_{4k-3}}                       \\
             \vec{\beta}_{4k-3} &\mapsto &\vec{\beta}_{4k-2}
               &\mapsto &\vec{\beta}_{4k-1}
               &\mapsto &\vec{\beta}_{4k} &\mapsto &\zero  \\
             &\vdotswithin{\vec{\beta}_{4k-3}}              \\
             \multicolumn{8}{l}{%
               \text{\textit{--可能还有别的更短的链--}}}
          \end{strings}
       \end{equation*}
       所以 $T^3$ 的值域空间的维数可以任意大。
       它最小为一\Dash 这里
       必须至少有一条链，否则该映射的幂零指数
       就不会是四。
     
\end{ans}
\begin{ans}{五.III.2.35}
       下面两个矩阵的特征值都只有零
       \begin{equation*}
         \begin{mat}
           0  &0  \\
           0  &0
         \end{mat}
         \qquad
         \begin{mat}
           0  &0  \\
           1  &0
         \end{mat}
       \end{equation*}
       但它们并不相似（它们有不同的标准形
       代表元，即它们自身）。
     
\end{ans}
\begin{ans}{五.III.2.36}
       由 \nearbylemma{le:RangeAndNullChains}，它是满射。
       它不必是恒等映射：考虑这个映射 $\map{t}{\Re^2}{\Re^2}$。
       \begin{equation*}
         \colvec{x \\ y}\mapsunder{t}\colvec{y \\ x}
       \end{equation*}
       对该映射有 $\genrangespace{t}=\Re^2$，而 $t$ 不是恒等映射。
     
\end{ans}
\begin{ans}{五.III.2.37}
       对链基作一个简单的重排即可。
       例如，与这个链基相关的映射
       \begin{equation*}
          \begin{strings}{ccccccc}
             \vec{\beta}_1 &\mapsto &\vec{\beta}_2 &\mapsto &\zero
          \end{strings}
       \end{equation*}
       关于
       $B=\sequence{\vec{\beta}_1,\vec{\beta}_2}$ 由这个矩阵表示
       \begin{equation*}
         \begin{mat}
           0  &0  \\
           1  &0
         \end{mat}
       \end{equation*}
       但关于
       $B=\sequence{\vec{\beta}_2,\vec{\beta}_1}$ 却这样表示。
       \begin{equation*}
         \begin{mat}
           0  &1  \\
           0  &0
         \end{mat}
       \end{equation*}
     
\end{ans}
\begin{ans}{五.III.2.38}
       设 $\map{t}{V}{V}$ 是该变换。
       若 \( \rank (t)=\nullity (t) \)，则等式
       \( \rank(t)+\nullity(t)=\dim (V) \) 表明
       \( \dim (V) \) 是偶数。
     
\end{ans}
\begin{ans}{五.III.2.39}
       矩阵要是幂零的，它们必须是方阵。
       要可交换，它们必须大小相同。
       于是它们的乘积与和都有定义。

       把这两个矩阵称为 \( A \) 与 \( B \)。
       要看出 \( AB \) 是幂零的，计算
       $
          (AB)^2=ABAB=AABB=A^2B^2$，
       以及
          $(AB)^3=A^3B^3$，依此类推，
       又因 \( A \) 是幂零的，该乘积最终为零。

       和的情形类似；利用二项式定理。
      
\end{ans}
\protect \section {第 IV 节：Jordan 形式}
\protect \subsection {五.IV.1: 映射与矩阵的多项式}
\begin{ans}{五.IV.1.13}
      Cayley-Hamilton 定理\nearbytheorem{th:CayHam}指出，
      极小多项式必须含有与特征多项式相同的线性因子，
      尽管次数可以更低，但不能是零次。
      \begin{exparts}
        \partsitem 可能的是
          $m_1(x)=x-3$，$m_2(x)=(x-3)^2$，$m_3(x)=(x-3)^3$
          以及 $m_4(x)=(x-3)^4$。
          第一个是一次多项式，第二个是二次的，
          第三个是三次的，第四个是四次的。
        \partsitem 可能的是 $m_1(x)=(x+1)(x-4)$、
          $m_2(x)=(x+1)^2(x-4)$ 以及 $m_3(x)=(x+1)^3(x-4)$。
          第一个是二次多项式，即次数为二。
          第二个次数为三，第三个次数为四。
        \partsitem 我们有 $m_1(x)=(x-2)(x-5)$、$m_2(x)=(x-2)^2(x-5)$、
          $m_3(x)=(x-2)(x-5)^2$ 以及 $m_4(x)=(x-2)^2(x-5)^2$。
          它们分别是二次、三次、三次与四次多项式。
        \partsitem 可能的是 \( m_1(x)=(x+3)(x-1)(x-2) \)、
          \( m_2(x)=(x+3)^2(x-1)(x-2) \)、
          \( m_3(x)=(x+3)(x-1)(x-2)^2 \)
          以及 \( m_4(x)=(x+3)^2(x-1)(x-2)^2 \)。
          $m_1$ 的次数是三，$m_2$ 的次数是四，
          $m_3$ 的次数是四，$m_4$ 的次数是五。
      \end{exparts}
    
\end{ans}
\begin{ans}{五.IV.1.14}
    特征多项式
    \begin{equation*}
      \deter{T-xI}=
      \begin{vmat}
        -x &1  &0  &1 \\
         1 &-x &1  &0 \\
         0 &1  &-x &1 \\
         1 &0  &1  &-x \\
      \end{vmat}
    \end{equation*}
    是 $x^2(x^2-4)$。

    由 Cayley-Hamilton 定理，特征多项式的所有互不相同的根
    也都是极小多项式的根，所以极小多项式的根为 $0$、$2$ 以及~$-2$。
    因此极小多项式要么是 $x(x^2-4)$，要么是 $x^2(x^2-4)$。

    求极小多项式的一种办法是，把矩阵代入这两个表达式试算，
    直到得到零矩阵为止。
    \begin{equation*}
      T(T^2-4I)=
      \begin{mat}
        0 &1 &0 &1 \\
        1 &0 &1 &0 \\
        0 &1 &0 &1 \\
        1 &0 &1 &0 \\
      \end{mat}
      \begin{mat}
        -2 &0  &1  &0 \\
        0  &-2 &0  &2 \\
        2  &0  &-2 &0 \\
        0  &2  &0  &-2 \\
      \end{mat}
      =
      \begin{mat}
        0 &0 &0 &0 \\
        0 &0 &0 &0 \\
        0 &0 &0 &0 \\
        0 &0 &0 &0 \\
      \end{mat}
    \end{equation*}
    所以极小多项式是 $x(x^2-4)$。
    
\end{ans}
\begin{ans}{五.IV.1.15}
       它的特征多项式有复根。
       \begin{equation*}
          \begin{vmat}
                   -x  &1  &0  \\
                    0  &-x &1  \\
                    1  &0  &-x
          \end{vmat}
          =(1-x)\cdot (x-(-\frac{1}{2}+\frac{\sqrt{3}}{2}i))
                \cdot (x-(-\frac{1}{2}-\frac{\sqrt{3}}{2}i))
       \end{equation*}
       由于根互不相同，特征多项式就等于极小多项式。
     
\end{ans}
\begin{ans}{五.IV.1.16}
      每种情形我们都将使用\nearbyexample{ex:MinPolyUsingCH}的方法。
      \begin{exparts}
       \partsitem 因为 $T$ 是三角矩阵，所以 $T-xI$ 也是三角矩阵
         \begin{equation*}
           T-xI=
           \begin{mat}
             3-x  &0    &0   \\
             1    &3-x  &0   \\
             0    &0    &4-x
           \end{mat}
         \end{equation*}
         特征多项式容易算出 $c(x)=\deter{T-xI}=(3-x)^2(4-x)=-1\cdot (x-3)^2(x-4)$。
         极小多项式只有两种可能，
         $m_1(x)=(x-3)(x-4)$ 与 $m_2(x)=(x-3)^2(x-4)$。
         （注意特征多项式带有负号，而极小多项式没有，
         因为它必须首项系数为一。）
         由于 $m_1(T)$ 不是零矩阵
         \begin{equation*}
           (T-3I)(T-4I)
           =
           \begin{mat}[r]
             0  &0  &0  \\
             1  &0  &0  \\
             0  &0  &1
           \end{mat}
           \begin{mat}[r]
             -1  &0  &0  \\
              1  &-1 &0  \\
              0  &0  &0
           \end{mat}
           =
           \begin{mat}[r]
             0  &0  &0  \\
            -1  &0  &0  \\
             0  &0  &0
           \end{mat}
         \end{equation*}
         所以极小多项式是 $m(x)=m_2(x)$。
         \begin{multline*}
           (T-3I)^2(T-4I)
           =(T-3I)\cdot\bigl((T-3I)(T-4I)\bigr)              \\
           =
           \begin{mat}
             0  &0  &0  \\
             1  &0  &0  \\
             0  &0  &1
           \end{mat}
           \begin{mat}[r]
              0  &0  &0  \\
             -1  &0  &0  \\
              0  &0  &0
           \end{mat}
           =
           \begin{mat}[r]
             0  &0  &0  \\
             0  &0  &0  \\
             0  &0  &0
           \end{mat}
         \end{multline*}
       \partsitem 与前一（小）题一样，矩阵是三角矩阵这一事实
        使得特征多项式的计算很容易。
         \begin{equation*}
           c(x)=\deter{T-xI}
               =
               \begin{vmat}
                 3-x  &0   &0   \\
                 1    &3-x &0   \\
                 0    &0   &3-x
               \end{vmat}
               =(3-x)^3=-1\cdot (x-3)^3
         \end{equation*}
         极小多项式有三种可能：
         $m_1(x)=(x-3)$、$m_2(x)=(x-3)^2$ 以及 $m_3(x)=(x-3)^3$。
         我们通过计算 $m_1(T)$
         \begin{equation*}
           T-3I=
           \begin{mat}[r]
             0  &0  &0  \\
             1  &0  &0  \\
             0  &0  &0
           \end{mat}
         \end{equation*}
         与 $m_2(T)$ 来解决这个问题。
         \begin{equation*}
           (T-3I)^2=
           \begin{mat}[r]
             0  &0  &0  \\
             1  &0  &0  \\
             0  &0  &0
           \end{mat}
           \begin{mat}[r]
             0  &0  &0  \\
             1  &0  &0  \\
             0  &0  &0
           \end{mat}
           =
           \begin{mat}[r]
             0  &0  &0  \\
             0  &0  &0  \\
             0  &0  &0
           \end{mat}
         \end{equation*}
         因为 $m_2(T)$ 是零矩阵，所以 $m_2(x)$ 就是极小多项式。
       \partsitem 同样，该矩阵是三角矩阵。
         \begin{equation*}
           c(x)=\deter{T-xI}
               =
               \begin{vmat}
                 3-x  &0   &0   \\
                 1    &3-x &0   \\
                 0    &1   &3-x
               \end{vmat}
               =(3-x)^3=-1\cdot (x-3)^3
         \end{equation*}
         同样，极小多项式有三种可能：
         $m_1(x)=(x-3)$、$m_2(x)=(x-3)^2$ 以及 $m_3(x)=(x-3)^3$。
         我们计算 $m_1(T)$
         \begin{equation*}
           T-3I=
           \begin{mat}[r]
             0  &0  &0  \\
             1  &0  &0  \\
             0  &1  &0
           \end{mat}
         \end{equation*}
         与 $m_2(T)$
         \begin{equation*}
           (T-3I)^2=
           \begin{mat}[r]
             0  &0  &0  \\
             1  &0  &0  \\
             0  &1  &0
           \end{mat}
           \begin{mat}[r]
             0  &0  &0  \\
             1  &0  &0  \\
             0  &1  &0
           \end{mat}
           =
           \begin{mat}[r]
             0  &0  &0  \\
             0  &0  &0  \\
             1  &0  &0
           \end{mat}
         \end{equation*}
         与 $m_3(T)$。
         \begin{equation*}
           (T-3I)^3
           =(T-3I)^2(T-3I)
           =
           \begin{mat}[r]
             0  &0  &0  \\
             0  &0  &0  \\
             1  &0  &0
           \end{mat}
           \begin{mat}[r]
             0  &0  &0  \\
             1  &0  &0  \\
             0  &1  &0
           \end{mat}
           =
           \begin{mat}[r]
             0  &0  &0  \\
             0  &0  &0  \\
             0  &0  &0
           \end{mat}
         \end{equation*}
         因此，极小多项式是 $m(x)=m_3(x)=(x-3)^3$。
       \partsitem 这一情形也是三角矩阵，这里是上三角矩阵。
          \begin{equation*}
            c(x)=\deter{T-xI}=
            \begin{vmat}
              2-x  &0   &1     \\
              0    &6-x &2     \\
              0    &0   &2-x
            \end{vmat}
            =(2-x)^2(6-x)=-(x-2)^2(x-6)
          \end{equation*}
          极小多项式有两种可能，
          $m_1(x)=(x-2)(x-6)$ 与 $m_2(x)=(x-2)^2(x-6)$。
          计算表明极小多项式不是 $m_1(x)$。
          \begin{equation*}
            (T-2I)(T-6I)=
            \begin{mat}[r]
              0  &0  &1  \\
              0  &4  &2  \\
              0  &0  &0
            \end{mat}
            \begin{mat}[r]
              -4  &0  &1  \\
               0  &0  &2  \\
               0  &0  &-4
            \end{mat}
            =
            \begin{mat}[r]
              0  &0  &-4  \\
              0  &0  &0   \\
              0  &0  &0
            \end{mat}
          \end{equation*}
          因此必有 $m(x)=m_2(x)=(x-2)^2(x-6)$。
          下面给出验证。
          \begin{multline*}
            (T-2I)^2(T-6I)=(T-2I)\cdot\bigl((T-2I)(T-6I)\bigr)             \\
            =
            \begin{mat}[r]
              0  &0  &1  \\
              0  &4  &2  \\
              0  &0  &0
            \end{mat}
            \begin{mat}[r]
               0  &0  &-4   \\
               0  &0  &0   \\
               0  &0  &0
            \end{mat}
            =
            \begin{mat}[r]
              0  &0  &0  \\
              0  &0  &0   \\
              0  &0  &0
            \end{mat}
          \end{multline*}
       \partsitem 特征多项式是
          \begin{equation*}
            c(x)=\deter{T-xI}=
            \begin{vmat}
              2-x  &2   &1     \\
              0    &6-x &2     \\
              0    &0   &2-x
            \end{vmat}
            =(2-x)^2(6-x)=-(x-2)^2(x-6)
          \end{equation*}
          极小多项式有两种可能，
          $m_1(x)=(x-2)(x-6)$ 与 $m_2(x)=(x-2)^2(x-6)$。
          检验第一个
          \begin{equation*}
            (T-2I)(T-6I)=
            \begin{mat}[r]
              0  &2  &1  \\
              0  &4  &2  \\
              0  &0  &0
            \end{mat}
            \begin{mat}[r]
              -4  &2  &1  \\
               0  &0  &2  \\
               0  &0  &-4
            \end{mat}
            =
            \begin{mat}[r]
              0  &0  &0  \\
              0  &0  &0   \\
              0  &0  &0
            \end{mat}
          \end{equation*}
          这表明极小多项式是
          $m(x)=m_1(x)=(x-2)(x-6)$。
       \partsitem 特征多项式如下。
          \begin{equation*}
            c(x)=\deter{T-xI}=
            \begin{vmat}
               -1-x &4    &0    &0    &0    \\
                0   &3-x  &0    &0    &0    \\
                0   &-4   &-1-x &0    &0    \\
                3   &-9   &-4   &2-x  &-1   \\
                1   &5    &4    &1    &4-x
            \end{vmat}
            =(x-3)^3(x+1)^2
          \end{equation*}
          下面是极小多项式的各种可能，按次数升序列出：
          $m_1(x)=(x-3)(x+1)$，$m_1(x)=(x-3)^2(x+1)$，$m_1(x)=(x-3)(x+1)^2$，
          $m_1(x)=(x-3)^3(x+1)$，$m_1(x)=(x-3)^2(x+1)^2$，
          以及 $m_1(x)=(x-3)^3(x+1)^2$。
          第一个行不通
          \begin{align*}
            (T-3I)(T+1I)
            &=
            \begin{mat}[r]
               -4   &4    &0    &0    &0    \\
                0   &0    &0    &0    &0    \\
                0   &-4   &-4   &0    &0    \\
                3   &-9   &-4   &-1   &-1   \\
                1   &5    &4    &1    &1
            \end{mat}
            \begin{mat}[r]
                0   &4    &0    &0    &0    \\
                0   &4    &0    &0    &0    \\
                0   &-4   &0    &0    &0    \\
                3   &-9   &-4   &3    &-1   \\
                1   &5    &4    &1    &5
            \end{mat}                           \\
            &=
            \begin{mat}[r]
                0   &0    &0    &0    &0    \\
                0   &0    &0    &0    &0    \\
                0   &0    &0    &0    &0    \\
               -4   &-4   &0    &-4   &-4   \\
                4   &4    &0    &4    &4
            \end{mat}
          \end{align*}
          而第二个则行得通。
          \begin{multline*}
            (T-3I)^2(T+1I)=(T-3I)\bigl((T-3I)(T+1I)\bigr) \\
            \begin{aligned}
            &=
            \begin{mat}[r]
               -4   &4    &0    &0    &0    \\
                0   &0    &0    &0    &0    \\
                0   &-4   &-4   &0    &0    \\
                3   &-9   &-4   &-1   &-1   \\
                1   &5    &4    &1    &1
            \end{mat}
            \begin{mat}[r]
                0   &0    &0    &0    &0    \\
                0   &0    &0    &0    &0    \\
                0   &0    &0    &0    &0    \\
               -4   &-4   &0    &-4   &-4   \\
                4   &4    &0    &4    &4
            \end{mat}                          \\
            &=
            \begin{mat}[r]
                0   &0    &0    &0    &0    \\
                0   &0    &0    &0    &0    \\
                0   &0    &0    &0    &0    \\
                0   &0    &0    &0    &0    \\
                0   &0    &0    &0    &0
            \end{mat}
            \end{aligned}
          \end{multline*}
          极小多项式是 \( m(x)=(x-3)^2(x+1) \)。
       \end{exparts}
     
\end{ans}
\begin{ans}{五.IV.1.17}
       我们知道 $\polyspace_n$ 是一个维数为 $n+1$ 的空间，
       并且求导算子是幂零指数为~$n+1$ 的幂零算子
       （例如取 $n=3$，
       $\polyspace_3=\set{c_3x^3+c_2x^2+c_1x+c_0\suchthat c_3,\ldots,c_0\in\C}$，
       而三次多项式的四阶导数是零多项式）。
       用幂零变换的标准
       形式来表示这个算子。
       \begin{equation*}
         \begin{mat}
           0  &0  &0  &\ldots &  &0  \\
           1  &0  &0  &       &  &0  \\
           0  &1  &0  &       &  &   \\
              &   &\ddots            \\
           0  &0  &0  &       &1 &0
         \end{mat}
       \end{equation*}
       这是一个 $\nbyn{(n+1)}$ 矩阵，其特征多项式容易算出，
       为 $c(x)=x^{n+1}$。
       （\textit{注：}这个矩阵是 $\rep{d/dx}{B,B}$，其中
        $B=\sequence{x^n,nx^{n-1},n(n-1)x^{n-2},\ldots,n!}$。）
       为像\nearbyexample{ex:MinPolyUsingCH}那样求出极小多项式，
       我们考虑 $T-0I=T$ 的各次幂。
       但是，当然，$T$ 的第一次变成零矩阵的幂
       是 $n+1$ 次幂。
       所以极小多项式也是 \( x^{n+1} \)。
     
\end{ans}
\begin{ans}{五.IV.1.18}
      记该矩阵为 $T$，并设它是 \( \nbyn{n} \) 的。
      由于 $T$ 是三角矩阵，从而 $T-xI$ 也是三角矩阵，
      所以特征多项式是 $c(x)=(x-\lambda)^n$。
      为看出极小多项式与此相同，考虑
      $T-\lambda I$。
      \begin{equation*}
        \begin{mat}
          0        &0        &0          &\ldots  &0  \\
          1        &0        &0          &\ldots  &0  \\
          0        &1        &0                       \\
                   &         &\ddots                  \\
          0        &0        &\ldots     &1       &0
        \end{mat}
      \end{equation*}
      将它识别为一个幂零指数为~$n$ 的变换的标准形；
      $(T-\lambda I)^j$ 第一次变为零是在 $j$ 等于 $n$ 时。
    
\end{ans}
\begin{ans}{五.IV.1.19}
      $n=3$ 的情形提供了一个提示。
      $\polyspace_3$ 的一个自然基是
      $B=\sequence{1,x,x^2,x^3}$。
      该变换的作用是
      \begin{equation*}
        1\mapsto 1
        \quad
        x\mapsto x+1
        \quad
        x^2\mapsto x^2+2x+1
        \quad
        x^3\mapsto x^3+3x^2+3x+1
      \end{equation*}
      所以表示 $\rep{t}{B,B}$ 是这个上三角矩阵。
      \begin{equation*}
        \begin{mat}[r]
          1  &1  &1  &1  \\
          0  &1  &2  &3  \\
          0  &0  &1  &3  \\
          0  &0  &0  &1
        \end{mat}
      \end{equation*}
      由于它是三角矩阵，显然特征多项式是
      $c(x)=(x-1)^4$。
      至于极小多项式，候选者是 $m_1(x)=(x-1)$、
      \begin{equation*}
        T-1I=
        \begin{mat}[r]
          0  &1  &1  &1  \\
          0  &0  &2  &3  \\
          0  &0  &0  &3  \\
          0  &0  &0  &0
        \end{mat}
      \end{equation*}
      $m_2(x)=(x-1)^2$、
      \begin{equation*}
        (T-1I)^2=
        \begin{mat}[r]
          0  &0  &2  &6  \\
          0  &0  &0  &6  \\
          0  &0  &0  &0  \\
          0  &0  &0  &0
        \end{mat}
      \end{equation*}
      $m_3(x)=(x-1)^3$、
      \begin{equation*}
        (T-1I)^3=
        \begin{mat}[r]
          0  &0  &0  &6  \\
          0  &0  &0  &0  \\
          0  &0  &0  &0  \\
          0  &0  &0  &0
        \end{mat}
      \end{equation*}
      以及 $m_4(x)=(x-1)^4$。
      由于 $m_1$、$m_2$ 和 $m_3$ 都不对，所以 $m_4$ 必定是对的，
      这也容易验证。

      在一般的 $n$ 的情形，表示是一个对角线上全为 $1$ 的
      上三角矩阵。
      于是特征多项式是 $c(x)=(x-1)^{n+1}$。
      验证极小多项式等于特征多项式的一种办法是
      给出类似这样的论证：
      若上三角矩阵的对角线上有非零元素，就称它为 $0$-上三角的；
      若对角线上只有零而对角线上方紧邻处有非零元素，
      就称它为 $1$-上三角的，等等。
      正如上面的例子所说明的，用归纳论证可以证明：
      当 $T$ 只有非负元素时，
      $T^j$ 是 $j$-上三角的。
      % We leave that argument to the reader.
    
\end{ans}
\begin{ans}{五.IV.1.20}
        该映射复合两次与复合一次相同：~$\composed{\pi}{\pi}=\pi$，
        即 $\pi^2=\pi$，所以极小多项式的次数
        至多是二，因为 \( m(x)=x^2-x \) 就行得通。
        没有一次多项式行得通这一事实，来自把映射作用到
        $c_1\cdot \pi+c_0\cdot \identity=z$（其中 $z$ 是零映射）
        的左右两边于下面这两个向量。
        \begin{equation*}
          \colvec[r]{0 \\ 0 \\ 1}
          \qquad
          \colvec[r]{1 \\ 0 \\ 0}
        \end{equation*}
        于是极小多项式是 $m$。
      
\end{ans}
\begin{ans}{五.IV.1.21}
         这是一个答案。
         \begin{equation*}
             \begin{mat}[r]
               0  &0  &0  \\
               1  &0  &0  \\
               0  &0  &0
             \end{mat}
         \end{equation*}
       
\end{ans}
\begin{ans}{五.IV.1.22}
        这里的 \( x \) 必须是标量，而不是矩阵。
      
\end{ans}
\begin{ans}{五.IV.1.23}
      矩阵
      \begin{equation*}
         T=\begin{mat}
              a  &b  \\
              c  &d
           \end{mat}
      \end{equation*}
      的特征多项式是 \( (a-x)(d-x)-bc=x^2-(a+d)x+(ad-bc) \)。
      代入
      \begin{multline*}
         \begin{mat}
              a  &b  \\
              c  &d
         \end{mat}^2
         -
         (a+d)\begin{mat}
            a  &b  \\
            c  &d
         \end{mat}
         +
         (ad-bc)\begin{mat}[r]
            1  &0  \\
            0  &1
         \end{mat}                            \\
         =
         \begin{mat}
            a^2+bc  &ab+bd  \\
            ac+cd   &bc+d^2
         \end{mat}
         -
         \begin{mat}
            a^2+ad  &ab+bd   \\
            ac+cd   &ad+d^2
         \end{mat}
         +
         \begin{mat}
            ad-bc  &0      \\
            0      &ad-bc
         \end{mat}
      \end{multline*}
      然后只需逐项核对每个元素的和，看结果是不是零矩阵。
    
\end{ans}
\begin{ans}{五.IV.1.24}
      由 Cayley-Hamilton 定理，极小多项式的次数
      小于或等于特征多项式的次数，
      即 \( n \)。
      \nearbyexample{ex:MinPolyForRotMat} 表明 \( n \) 可以取到。
    
\end{ans}
\begin{ans}{五.IV.1.25}
       设该线性变换为 $\map{t}{V}{V}$。
       若 $t$ 是幂零的，则存在 $n$ 使得 $t^n$ 是零映射，
       于是 $t$ 满足多项式 $p(x)=x^n=(x-0)^n$。
       由 \nearbylemma{le:tSatisImpMinPolyDivides}，$t$ 的极小多项式
       整除 $p$，所以极小多项式的根只有零。
       由 Cayley-Hamilton 定理，即 \nearbytheorem{th:CayHam}，
       特征多项式的根也只有零。
       于是 $t$ 仅有的特征值是零。

       反之，若 \( n \) 维空间上的变换 \( t \) 只有零这一个特征值，
       则它的特征多项式是 \( x^n \)。
       Cayley-Hamilton 定理说映射满足其特征多项式，
       所以 \( t^n \) 是零映射。
       于是 $t$ 是幂零的。
     
\end{ans}
\begin{ans}{五.IV.1.26}
         极小多项式必须首项系数为 $1$，
         因此若一个映射或矩阵的极小多项式
         是零次多项式，则它将是 $m(x)=1$。
         但恒等映射或单位矩阵等于零映射或零矩阵，
         仅在平凡向量空间上才如此。

         所以在非平凡的情形，极小多项式的次数
         至少是一。
         零映射或零矩阵的极小多项式是 \( m(x)=x \)，而
         恒等映射或单位矩阵的极小多项式是 \( m(x)=x-1 \)。
       
\end{ans}
\begin{ans}{五.IV.1.27}
       对于对角矩阵
       \begin{equation*}
          T=
          \begin{mat}
             t_{1,1}   &0        \\
             0         &t_{2,2}  \\
                       &        &\ddots  \\
                       &        &      &t_{n,n}
          \end{mat}
       \end{equation*}
       特征多项式是
       $(t_{1,1}-x)(t_{2,2}-x)\cdots (t_{n,n}-x)$。
       当然，其中一些因子可能重复，例如矩阵可能
       有 $t_{1,1}=t_{2,2}$。
       例如，
       \begin{equation*}
          D=
          \begin{mat}[r]
             3 &0 &0  \\
             0 &3 &0  \\
             0 &0 &1
          \end{mat}
       \end{equation*}
       的特征多项式是 \( (3-x)^2(1-x)=-1\cdot (x-3)^2(x-1) \)。

       要构造极小多项式，
       取各项 \( x-t_{i,i} \)，去掉重复项，
       然后把它们乘起来。
       例如，$D$ 的极小多项式
       是 \( (x-3)(x-1) \)。
       为验证这一点，首先注意 \nearbytheorem{th:CayHam}，
       即 Cayley-Hamilton 定理，要求特征多项式中的每个线性因子
       在极小多项式中至少出现一次。
       检验另一个方向\Dash 即在对角矩阵的情形，
       每个线性因子至多出现一次\Dash 的一个办法是
       使用矩阵论证。
       对角矩阵从左边相乘时按
       对角元进行倍乘。
       而在乘积 $(T-t_{1,1}I)\cdots\hbox{}$ 中，即便没有任何重复
       因子，每一行也至少在其中一个因子中为零。

       例如，在乘积
       \begin{equation*}
         (D-3I)(D-1I)=(D-3I)(D-1I)I=
         \begin{mat}[r]
           0  &0  &0  \\
           0  &0  &0  \\
           0  &0  &-2
         \end{mat}
         \begin{mat}[r]
           2  &0  &0  \\
           0  &2  &0  \\
           0  &0  &0
         \end{mat}
         \begin{mat}[r]
           1  &0  &0  \\
           0  &1  &0  \\
           0  &0  &1
         \end{mat}
       \end{equation*}
       中，因为第一个矩阵 $D-3I$ 的第一行与第二行为零，
       所以整个乘积的第一行与第二行都是零。
       又因为中间那个矩阵 $D-1I$ 的第三行为零，
       所以整个乘积的第三行是零。
    
\end{ans}
\begin{ans}{五.IV.1.28}
      本小节开头指出，线性变换的幂不可能永远攀升而没有
      ``重复''，也就是说，对某个幂~$n$ 存在线性关系
      $c_n\cdot t^n+\dots+c_1\cdot t+c_0\cdot \identity=z$，其中 $z$ 是
      零变换。
      投影的定义是，对这样的映射，
      有一个线性关系是二次的：$t^2-t=z$。
      为完成讨论，我们只需考虑这个关系是否可能不是极小的，
      即是否存在极小多项式为常数或一次的投影？

      若极小多项式是常数，则该映射必须满足
      $c_0\cdot\identity=z$，其中 $c_0=1$，因为极小多项式的
      首项系数是 $1$。
      这只在平凡空间上被零变换满足。
      这是一个投影，但不是令人感兴趣的投影。

      若一个变换的极小多项式是一次
      的，则给出 $c_1\cdot t+c_0\cdot\identity=z$，其中 $c_1=1$。
      这个方程给出 $t=-c_0\cdot \identity$。
      再结合要求 $t^2=t$，得
      $t^2=(-c_0)^2\cdot\identity=-c_0\cdot\identity$，由此得
      $c_0=0$ 且 $t$ 是零变换，或者 $c_0=1$ 且
      $t$ 是恒等变换。

      于是，除投影是零映射或恒等映射的情形外，
      极小多项式是 $m(x)=x^2-x$。
    
\end{ans}
\begin{ans}{五.IV.1.29}
      \begin{exparts}
       \partsitem \textit{这是一般函数的性质，
          而不仅是线性函数的性质。}
          设 $f$ 和 $g$ 是单射函数，且
          $\composed{f}{g}$ 有定义。
          设 $\composed{f}{g}(x_1)=\composed{f}{g}(x_2)$，则
          $f(g(x_1))=f(g(x_2))$。
          由于 $f$ 是单射的，这蕴含 $g(x_1)=g(x_2)$。
          又因 $g$ 也是单射的，这进而蕴含
          $x_1=x_2$。
          于是，总之，$\composed{f}{g}(x_1)=\composed{f}{g}(x_2)$
          蕴含 $x_1=x_2$，所以 $\composed{f}{g}$ 是单射。
        \partsitem 若线性映射 $h$
          不是单射，则存在不相等的
          向量 $\vec{v}_1$、$\vec{v}_2$ 被映到相同的值
          $h(\vec{v}_1)=h(\vec{v}_2)$。
          由于 $h$ 是线性的，我们有
          $\zero=h(\vec{v}_1)-h(\vec{v}_2)=h(\vec{v}_1-\vec{v}_2)$，
          于是 $\vec{v}_1-\vec{v}_2$ 是定义域中一个非零向量，
          $h$ 把它映到陪域的零向量
          （$\vec{v}_1-\vec{v}_2$
          不等于定义域的零向量，因为 $\vec{v}_1$
          不等于 $\vec{v}_2$）。
        \partsitem 极小多项式
          $m(t)$ 把定义域中的每个向量都映到
          零，所以它不是单射（在平凡空间中除外，这里
          不予考虑）。
          由本题的第一（小）题，
          由于复合 $m(t)$ 不是单射，
          所以分量 $t-\lambda_i$ 中至少有一个不是单射。
          由第二（小）题，$t-\lambda_i$ 有一个非平凡的零空间。
          因为 $(t-\lambda_i)(\vec{v})=\zero$ 成立当且仅当
          $t(\vec{v})=\lambda_i\cdot\vec{v}$，前一句话给出：
          $\lambda_i$ 是特征值（回忆特征值的
          定义要求该关系至少对一个
          非零 $\vec{v}$ 成立）。
      \end{exparts}
    
\end{ans}
\begin{ans}{五.IV.1.30}
      这是错的。
      不可对角化变换的自然例子在这里适用。
      考虑 $\C^2$ 上关于标准基由这个矩阵
      表示的变换。
      \begin{equation*}
        N=
        \begin{mat}[r]
          0  &1  \\
          0  &0
        \end{mat}
      \end{equation*}
      特征多项式是 $c(x)=x^2$。
      于是极小多项式要么是 $m_1(x)=x$，要么是 $m_2(x)=x^2$。
      第一个不对，因为 $N-0\cdot I$ 不是零矩阵；
      于是在这个例子中极小多项式的次数等于
      底空间的维数，而且如前所述，
      我们知道这个矩阵不可对角化，因为它是幂零的。
    
\end{ans}
\begin{ans}{五.IV.1.31}
       设 \( A \) 与 \( B \) 相似，即 $A=PBP^{-1}$。
       由这样的事实：
       \begin{multline*}
           A^n=(PBP^{-1})^n=(PBP^{-1})(PBP^{-1})\cdots(PBP^{-1})   \\
                           =PB(P^{-1}P)B(P^{-1}P)\cdots (P^{-1}P)BP^{-1}
                           =PB^nP^{-1}
       \end{multline*}
       以及 $c\cdot A=c\cdot(PBP^{-1})=P(c\cdot B)P^{-1}$，
       即得所需的事实：对任意多项式函数 $f$，我们有
       \( f(A)=P\,f(B)\,P^{-1} \)。
       例如，若 $f(x)=x^2+2x+3$，则
       \begin{multline*}
         A^2+2A+3I=(PBP^{-1})^2+2\cdot PBP^{-1}+3\cdot I           \\
                  =(PBP^{-1})(PBP^{-1})+P(2B)P^{-1}+3\cdot PP^{-1}
                  =P(B^2+2B+3I)P^{-1}
       \end{multline*}
       这表明 $f(A)$ 与 $f(B)$ 相似。
       \begin{exparts}
         \partsitem 取 $f$ 为一次多项式，我们得到
           $A-xI$ 与 $B-xI$ 相似。
           相似矩阵有相等的行列式（因为
           $\deter{A}=\deter{PBP^{-1}}
               =\deter{P}\cdot\deter{B}\cdot\deter{P^{-1}}
               =1\cdot\deter{B}\cdot 1=\deter{B}$）。
           于是特征多项式相等。
         \partsitem
           由于 \( P \) 与 \( P^{-1} \) 可逆，\( f(A) \) 是
           零矩阵当且仅当 \( f(B) \) 是零矩阵。
         \partsitem 它们不可能相似，因为它们没有相同的
           特征多项式。
           第一个的特征多项式是
           $x^2-4x-3$，而第二个的
           特征多项式是 $x^2-5x+5$。
       \end{exparts}
     
\end{ans}
\begin{ans}{五.IV.1.32}
      设 \( m(x)=x^n+m_{n-1}x^{n-1}+\dots+m_1x+m_0 \)
      是 \( T \) 的极小多项式。
      \begin{exparts}
       \partsitem
         对于「若」的论证，
         由于 \( T^n+\dots+m_1T+m_0I \) 是零矩阵，我们
         得到 \( I=(T^n+\dots+m_1T)/(-m_0)=
         T\cdot (T^{n-1}+\dots+m_1I)/(-m_0) \)，所以
         矩阵 $(-1/m_0)\cdot (T^{n-1}+\dots+m_1I)$ 是 $T$ 的逆。
         对于「仅当」，设 \( m_0=0 \)
         （我们把 \( n=1 \) 的情形搁置一旁，但它是容易的），
         于是
         \( T^n+\dots+m_1T=(T^{n-1}+\dots+m_1I)T \) 是零矩阵。
         注意 \( T^{n-1}+\dots+m_1I \) 不是零矩阵，因为
         极小多项式的次数是 \( n \)。
         若 \( T^{-1} \) 存在，则把
         \( (T^{n-1}+\dots+m_1I)T \) 与零矩阵从右边
         乘以 $T^{-1}$ 就得出矛盾。
       \partsitem 若 \( T \) 不可逆，则其
         极小多项式中的常数项为零。
         于是，
         \begin{equation*}
            T^n+\dots+m_1T=(T^{n-1}+\dots+m_1I)T=T(T^{n-1}+\dots+m_1I)
         \end{equation*}
         是零矩阵。
       \end{exparts}
     
\end{ans}
\begin{ans}{五.IV.1.33}
        \begin{exparts}
          \partsitem
            对于归纳步骤，设 \nearbylemma{le:PolyMapsFactor}
            对次数为 \( i,\ldots,k-1 \) 的多项式
            成立，并考虑次数为 \( k \) 的多项式 \( f(x) \)。
            把 $f(x)$ 分解为 $f(x)=k(x-\lambda_1)^{q_1}\cdots(x-\lambda_z)^{q_z}$，
            并设
            \( k(x-\lambda_1)^{q_1-1}\cdots(x-\lambda_z)^{q_z} \)
            为 \( c_{n-1}x^{n-1}+\cdots+c_1x+c_0 \)。
            代入：
            \begin{align*}
               k\composed{\composed{(t-\lambda_1)^{q_1}}{\cdots}}{
                                            (t-\lambda_z)^{q_z}}(\vec{v})
               &=
               \composed{(t-\lambda_1)}{
                  \composed{\composed{(t-\lambda_1)^{q_1}}{\cdots}}{
                  (t-\lambda_z)^{q_z}} }
                          (\vec{v})    \\
               &=
               (t-\lambda_1)\,(c_{n-1}t^{n-1}(\vec{v})+\cdots+c_0\vec{v}) \\
               &=
               f(t)(\vec{v})
            \end{align*}
           （第二个等号来自归纳假设，
           第三个等号来自 \( t \) 的线性）。
         \partsitem 一个例子是考虑平方映射
            $\map{s}{\Re}{\Re}$，由 $s(x)=x^2$ 给出。
            它是非线性的。
            多项式 $f(t)=t^2-1$ 定义的作用
            把 $s$ 变成 $f(s)=s^2-1$，即这个映射。
            \begin{equation*}
              x\mapsunder{s^2-1} \composed{s}{s}(x)-1=x^4-1
            \end{equation*}
            注意这个映射不同于映射
            $\composed{(s-1)}{(s+1)}$；例如，第一个映射把
            $x=5$ 变为 $624$，而第二个把 $x=5$ 变为 $675$。
       \end{exparts}
     
\end{ans}
\begin{ans}{五.IV.1.34}
      成立。
      按最后一列展开可以检验
      \( x^n+m_{n-1}x^{n-1}+\dots+m_1x+m_0 \) 是下面这个行列式的
      正值或负值。
      \begin{equation*}
         \begin{mat}
            -x  &0  &0  &      &    &m_0    \\
             0  &1-x&0  &      &    &m_1    \\
             0  &0  &1-x&      &    &m_2    \\
                &   &   &\ddots             \\
                &   &   &      &1-x &m_{n-1}
         \end{mat}
      \end{equation*}
     
\end{ans}
\protect \subsection {五.IV.2: Jordan 标准形}
\begin{ans}{五.IV.2.15}
      我们必须检验
      \begin{equation*}
         \begin{mat}[r]
           3  &0  \\
           1  &3
         \end{mat}
         =
         N+3I=PTP^{-1}
         =
         \begin{mat}[r]
           1/2  &1/2  \\
          -1/4 &1/4
         \end{mat}
         \begin{mat}[r]
           2  &-1  \\
           1  &4
         \end{mat}
         \begin{mat}[r]
           1  &-2  \\
           1  &2
         \end{mat}
      \end{equation*}
      这个计算很容易。
    
\end{ans}
\begin{ans}{五.IV.2.16}
      \begin{exparts}
        \partsitem 特征多项式为 $c(x)=(x-3)^2$，
          极小多项式与之相同。
        \partsitem 特征多项式为 $c(x)=(x+1)^2$。
          极小多项式为 $m(x)=x+1$。
        \partsitem 特征多项式为
          $c(x)=(x+(1/2))(x-2)^2$，
          极小多项式与之相同。
        \partsitem 特征多项式为 $c(x)=(x-3)^3$
          极小多项式与之相同。
        \partsitem 特征多项式为 $c(x)=(x-3)^4$。
          极小多项式为 $m(x)=(x-3)^2$。
        \partsitem 特征多项式为 $c(x)=(x+4)^2(x-4)^2$，
          极小多项式与之相同。
        \partsitem 特征多项式为
          $c(x)=(x-2)^2(x-3)(x-5)$，
          极小多项式为 $m(x)=(x-2)(x-3)(x-5)$。
        \partsitem 特征多项式为
          $c(x)=(x-2)^2(x-3)(x-5)$，
          极小多项式与之相同。
      \end{exparts}
    
\end{ans}
\begin{ans}{五.IV.2.17}
      \begin{exparts}
         \partsitem 变换 $t-3$ 是幂零的
          （即 $\gennullspace{t-3}$ 是整个空间），
          它通过两条链作用在一个链基上：
          $\vec{\beta}_1\mapsto\vec{\beta}_2\mapsto\vec{\beta}_3
            \mapsto\vec{\beta}_4\mapsto\zero$
          以及 $\vec{\beta}_5\mapsto\zero$。
          因此，$t-3$ 可以用这个标准形来表示。
          \begin{equation*}
            N_3=
            \begin{mat}[r]
               0  &0  &0  &0  &0  \\
               1  &0  &0  &0  &0  \\
               0  &1  &0  &0  &0  \\
               0  &0  &1  &0  &0  \\
               0  &0  &0  &0  &0
             \end{mat}
          \end{equation*}
          从而 $T$ 相似于这个标准形矩阵。
          \begin{equation*}
           J_3=N_3+3I=
           \begin{mat}[r]
               3  &0  &0  &0  &0  \\
               1  &3  &0  &0  &0  \\
               0  &1  &3  &0  &0  \\
               0  &0  &1  &3  &0  \\
               0  &0  &0  &0  &3
             \end{mat}
          \end{equation*}
         \partsitem 变换 $s+1$ 在子空间 $\gennullspace{s+1}$ 上的限制
          是幂零的，它在链基上的作用为 $\vec{\beta}_1\mapsto\zero$。
          变换 $s-2$ 在子空间 $\gennullspace{s-2}$ 上的限制
          是幂零的，它在链基上的作用为 $\vec{\beta}_2\mapsto\vec{\beta}_3\mapsto\zero$
          与 $\vec{\beta}_4\mapsto\vec{\beta}_5\mapsto\zero$。
          因此 Jordan 形式如下。
          \begin{equation*}
            \begin{mat}[r]
              -1  &0  &0  &0  &0  \\
               0  &2  &0  &0  &0  \\
               0  &1  &2  &0  &0  \\
               0  &0  &0  &2  &0  \\
               0  &0  &0  &1  &2
             \end{mat}
          \end{equation*}
      \end{exparts}
    
\end{ans}
\begin{ans}{五.IV.2.18}
      对每一小题，由于基的取法有很多种，所以还可以有许多别的答案。
      当然，检验一个答案是否给出 Jordan 形式的 $PTP^{-1}$
      的计算才是正确与否的裁判。
      \begin{exparts}
        \partsitem 下面是箭头图。
          \begin{equation*}
            \begin{CD}
              \C^3_{\wrt{\stdbasis_3}}      @>t>T>    \C^3_{\wrt{\stdbasis_3}}   \\
                @V\scriptstyle\identity V\scriptstyle PV
                                    @V\scriptstyle\identity V\scriptstyle PV \\
              \C^3_{\wrt{B}}                 @>t>J>         \C^3_{\wrt{B}}
            \end{CD}
          \end{equation*}
          从左下角移动到左上角的矩阵是这个。
          \begin{equation*}
            P^{-1}=\bigl(\rep{\identity}{\stdbasis_3,B}\bigr)^{-1}
                  =\rep{\identity}{B,\stdbasis_3}
                  =\begin{mat}[r]
                     1  &-2   &0 \\
                     1  &0   &1 \\
                    -2  &0   &0
                   \end{mat}
          \end{equation*}
          从右上角移动到右下角的矩阵 $P$ 是 $P^{-1}$ 的逆。
        \partsitem 我们需要这个矩阵及其逆。
          \begin{equation*}
            P^{-1}=
            \begin{mat}[r]
              1  &0  &3  \\
              0  &1  &4  \\
              0  &-2 &0
            \end{mat}
          \end{equation*}
        \partsitem 将这些广义零空间的基拼接起来，
          就可以作为整个空间的基。
          \begin{equation*}
             B_{-1}=\sequence{\colvec[r]{-1\\ 0 \\  0 \\ 1 \\ 0},
                 \colvec[r]{-1\\ 0 \\ -1 \\ 0 \\ 1}}
            \qquad
             B_3=\sequence{\colvec[r]{1 \\ 1 \\ -1 \\ 0 \\ 0},
                 \colvec[r]{0 \\ 0 \\  0 \\-2 \\ 2},
                 \colvec[r]{-1\\-1 \\  1 \\ 2 \\ 0}}
          \end{equation*}
          换基矩阵是这一个矩阵及其逆。
          \begin{equation*}
            P^{-1}=
            \begin{mat}[r]
              -1  &-1  &1  &0  &-1  \\
              0   &0   &1  &0  &-1  \\
              0   &-1  &-1 &0  &1   \\
              1   &0   &0  &-2 &2   \\
              0   &1   &0  &2  &0   \\
            \end{mat}
          \end{equation*}
      \end{exparts}
    
\end{ans}
\begin{ans}{五.IV.2.19}
      一般的做法是将特征多项式
      $c(x)=(x-\lambda_1)^{p_1}(x-\lambda_2)^{p_2}\cdots $
      因式分解，得到特征值 $\lambda_1$、$\lambda_2$ 等。
      然后，对每个 $\lambda_i$ 我们求一个
      链基，它是变换 $t-\lambda_i$
      限制在 $\gennullspace{t-\lambda_i}$ 上时的作用的链基：
      计算矩阵 $T-\lambda_iI$ 的各次幂并求出
      相应的零空间，直到这些零空间稳定下来
      （不再变化）为止，此时我们就得到了广义零空间。
      这些零空间的维数（零度）告诉我们
      $t-\lambda_i$ 在广义零空间的某个链基上的作用，
      于是我们可以写出次对角线上 $1$ 的分布模式，
      得到 $N_{\lambda_i}$。
      由此矩阵，与 $\lambda_i$ 相关联的 Jordan 块
      $J_{\lambda_i}$ 立即可得：$J_{\lambda_i}=N_{\lambda_i}+\lambda_iI$。
      最后，在对每个特征值都这样做之后，
      我们把它们拼成标准形。
      \begin{exparts}
        \partsitem 这个矩阵的特征多项式
           是 $c(x)=(-10-x)(10-x)+100=x^2$，
           所以它只有唯一的特征值 $\lambda=0$。
           \begin{center}
             \renewcommand{\arraystretch}{1.25}
             \begin{tabular}{r|ccc}
                \multicolumn{1}{c}{\textit{幂}~$p$}
                    &$(T+0\cdot I)^p$ &$\nullspace{(t-0)^p}$
                    &\textit{零度}                   \\
                \hline
                $1$
                &\matrixvenlarge{\begin{mat}[r]
                  -10  &4  \\
                  -25  &10
                \end{mat}}
                &$\set{\matrixvenlarge{\colvec{2y/5 \\ y}}
                                     \suchthat
                                     y\in\C}$
                &$1$                                       \\
                $2$
                &\matrixvenlarge{\begin{mat}[r]
                    0  &0  \\
                    0  &0
                \end{mat}}
                &$\C^2$
                &$2$
             \end{tabular}
           \end{center}
           （因此，这个变换是幂零的：
           $\gennullspace{t-0}$ 是整个空间）。
           由这些零度我们知道 $t$ 在某个链基上的作用是
           $\vec{\beta}_1\mapsto\vec{\beta}_2\mapsto\zero$。
           这是 $t-0$ 在 $\gennullspace{t-0}=\C^2$ 上的作用的
           标准形矩阵
           \begin{equation*}
             N_0=
             \begin{mat}[r]
               0  &0  \\
               1  &0
            \end{mat}
           \end{equation*}
           而这就是该矩阵的 Jordan 形式。
           \begin{equation*}
             J_0=N_0+0\cdot I=
             \begin{mat}[r]
               0  &0  \\
               1  &0
            \end{mat}
           \end{equation*}
           注意，若一个矩阵是幂零的，则它的标准形
           等于它的 Jordan 形式。

           我们可以用前一小节的方法找到这样一个链基。
           \begin{equation*}
                 B=\sequence{\colvec[r]{1 \\ 0},
                             \colvec[r]{-10 \\ -25}}
           \end{equation*}
           我们取第一个基向量使得它在 $t^2$ 的零空间中，
           但不在 $t$ 的零空间中。
           第二个基向量是第一个基向量在 $t$ 下的像。
        \partsitem 这个矩阵的特征多项式
           是 \( c(x)=(x+1)^2 \)，所以它是单特征值矩阵。
           （即 $t+1$ 的广义零空间是整个空间。）
           我们有
           \begin{equation*}
             \nullspace{t+1}=\set{\colvec{2y/3 \\ y}\suchthat
                                     y\in\C}
             \qquad
             \nullspace{(t+1)^2}=\C^2
           \end{equation*}
           于是 $t+1$ 在某个相应的链基上的作用是
           $\vec{\beta}_1\mapsto\vec{\beta}_2\mapsto\zero$。
           因此，
           \begin{equation*}
             N_{-1}
             =
             \begin{mat}[r]
                0  &0  \\
                1  &0
             \end{mat}
           \end{equation*}
           $T$ 的 Jordan 形式是
           \begin{equation*}
             J_{-1}=N_{-1}+-1\cdot I
             =
             \begin{mat}[r]
               -1  &0  \\
                1  &-1
             \end{mat}
           \end{equation*}
           而从上面的零空间中取向量就得到
           这个链基（也可以有别的取法）。
           \begin{equation*}
             B=\sequence{\colvec[r]{1 \\ 0},
                         \colvec[r]{6 \\ 9}}
           \end{equation*}
        \partsitem 特征多项式
            \( c(x)=(1-x)(4-x)^2=-1\cdot (x-1)(x-4)^2 \) 有两个根，
            它们就是特征值 $\lambda_1=1$ 与 $\lambda_2=4$。

            我们分别处理这两个特征值。
            对 $\lambda_1$，计算 $T-1I$ 的各次幂
            得到
            \begin{equation*}
              \nullspace{t-1}=\set{\colvec{0 \\ y \\ 0}
                                      \suchthat y\in\C}
            \end{equation*}
            而 $(t-1)^2$ 的零空间与之相同。
            于是这个集合就是广义零空间
            $\gennullspace{t-1}$。
            这些零度表明 $t-1$ 限制在广义零空间上后
            在某个链基上的作用
            是 $\vec{\beta}_1\mapsto\zero$。

            对 $\lambda_2=4$ 作类似的计算，得到这些零空间。
            \begin{equation*}
              \nullspace{t-4}=\set{\colvec{0 \\ z \\ z}
                                      \suchthat z\in\C}
              \qquad
              \nullspace{(t-4)^2}=\set{\colvec{z-y \\ y \\ z}
                                          \suchthat y,z\in\C}
            \end{equation*}
            （$(t-4)^3$ 的零空间与之相同，这也必然如此，因为
            特征多项式中与 $\lambda_2=4$ 相关联的项的
            幂是二，所以
            $t-2$ 限制在广义零空间 $\gennullspace{t-2}$ 上
            是幂零的且指数至多为二\Dash 零空间至多经过
            两次 $t-2$ 的作用就稳定下来。）
            零度上升的模式告诉我们
             $t-4$ 在 $\gennullspace{t-4}$ 的某个相应链基上的作用
            是
            $\vec{\beta}_2\mapsto\vec{\beta}_3\mapsto\zero$。

            把这两个特征值的信息
            合在一起，就得到变换 $t$ 的 Jordan 形式。
            \begin{equation*}
              \begin{mat}[r]
                1  &0  &0  \\
                0  &4  &0  \\
                0  &1  &4
              \end{mat}
            \end{equation*}
            我们可以取零空间中的元素来得到一个合适的基。
            \begin{equation*}
              B=\cat{B_{1}}{B_4}=
               \sequence{\colvec[r]{0 \\ 1 \\ 0},
                          \colvec[r]{1 \\ 0 \\ 1},
                          \colvec[r]{0 \\ 5 \\ 5}}
            \end{equation*}
        \partsitem 特征多项式是
            \( c(x)=(-2-x)(4-x)^2=-1\cdot (x+2)(x-4)^2 \)。

            对特征值 $\lambda_{-2}$
            （这里的下标是 \( -2\)），计算
            $T+2I$ 的各次幂得到这个。
            \begin{equation*}
              \nullspace{t+2}=\set{\colvec{z \\ z \\ z}
                                      \suchthat z\in\C}
            \end{equation*}
            $(t+2)^2$ 的零空间与之相同，所以
            这就是广义零空间 $\gennullspace{t+2}$。
            于是 $t+2$ 限制在
            $\gennullspace{t+2}$ 上后，在某个相应的
            链基上的作用是 $\vec{\beta}_1\mapsto\zero$。

            对 $\lambda_2=4$，
            计算 $T-4I$ 的各次幂得到
            \begin{equation*}
              \nullspace{t-4}=\set{\colvec{z \\ -z \\ z}
                                      \suchthat z\in\C}
              \qquad
              \nullspace{(t-4)^2}=\set{\colvec{x \\ -z \\ z}
                                           \suchthat x,z\in\C}
            \end{equation*}
            于是 $t-4$ 在 $\gennullspace{t-4}$ 的某个
            链基上的作用是
            $\vec{\beta}_2\mapsto\vec{\beta}_3\mapsto\zero$。

            因此 Jordan 形式为
            \begin{equation*}
              \begin{mat}[r]
                -2  &0  &0  \\
                 0  &4  &0  \\
                 0  &1  &4
              \end{mat}
            \end{equation*}
            而一个合适的基如下。
            \begin{equation*}
              B=\cat{B_{-2}}{B_4}=
                \sequence{\colvec[r]{1 \\ 1 \\ 1},
                          \colvec[r]{0 \\ -1 \\ 1},
                          \colvec[r]{-1 \\ 1 \\ -1}}
            \end{equation*}
        \partsitem 该矩阵的特征多项式
            为 \( c(x)=(2-x)^3=-1\cdot (x-2)^3 \)。
            这个矩阵只有一个特征值 $\lambda=2$。
            通过求 $T-2I$ 的各次幂可得
            \begin{equation*}
              \nullspace{t-2}=\set{\colvec{-y \\ y \\ 0}
                                      \suchthat y\in\C}
              \qquad
              \nullspace{(t-2)^2}=\set{\colvec{-y-(1/2)z \\ y \\ z}
                                          \suchthat y,z\in\C}
            \end{equation*}
            以及
            \begin{equation*}
              \nullspace{(t-2)^3}=\C^3
            \end{equation*}
            所以
            $t-2$ 在一个相伴的
            链基上的作用是
            $\vec{\beta}_1\mapsto\vec{\beta}_2\mapsto
                   \vec{\beta}_3\mapsto\zero$。
            Jordan 形式为
            \begin{equation*}
                  \begin{mat}[r]
                    2  &0  &0  \\
                    1  &2  &0  \\
                    0  &1  &2
                  \end{mat}
            \end{equation*}
            而一个可选的基如下。
            \begin{equation*}
              B=\sequence{\colvec[r]{0 \\ 1 \\ 0},
                          \colvec[r]{7 \\ -9 \\ 4},
                          \colvec[r]{-2 \\ 2 \\ 0}}
            \end{equation*}
        \partsitem 特征多项式
            \( c(x)=(1-x)^3=-(x-1)^3 \) 只有一个单根，
            因此该矩阵只有一个特征值 $\lambda=1$。
            求 $T-1I$ 的各次幂
            并计算零空间
            \begin{equation*}
               \nullspace{t-1}=\set{\colvec{-2y+z \\ y \\ z}
                                      \suchthat y,z\in\C}
               \qquad
               \nullspace{(t-1)^2}=\C^3
            \end{equation*}
            可知幂零映射 $t-1$ 在一个链基上
            的作用是
            $\vec{\beta}_1\mapsto\vec{\beta}_2\mapsto\zero$ 与
            $\vec{\beta}_3\mapsto\zero$。
            因此 Jordan 形式为
            \begin{equation*}
                  J=
                  \begin{mat}[r]
                    1  &0  &0  \\
                    1  &1  &0  \\
                    0  &0  &1
                  \end{mat}
            \end{equation*}
            而一个合适的基（与 $t-1$ 相伴的
            链基）如下。
            \begin{equation*}
              B=\sequence{\colvec[r]{0 \\ 1 \\ 0},
                          \colvec[r]{2 \\ -2 \\ -2},
                          \colvec[r]{1 \\ 0 \\ 1}}
            \end{equation*}
        \partsitem 特征多项式用手算
            有点大，但尚可应付：
            \( c(x)=x^4-24x^3+216x^2-864x+1296=(x-6)^4 \)。
            这是单特征值映射，所以
            变换 $t-6$ 是幂零的。
            零空间
            \begin{equation*}
               \nullspace{t-6}=\set{\colvec{-z-w \\ -z-w \\ z \\ w}
                                      \suchthat z,w\in\C}
               \qquad
               \nullspace{(t-6)^2}=\set{\colvec{x \\ -z-w \\ z \\ w}
                                      \suchthat x,z,w\in\C}
            \end{equation*}
            以及
            \begin{equation*}
               \nullspace{(t-6)^3}=\C^4
            \end{equation*}
            连同各零度一起表明
            $t-6$ 在一个链基上的作用是
            $\vec{\beta}_1\mapsto\vec{\beta}_2\mapsto
                   \vec{\beta}_3\mapsto\zero$ 与
            $\vec{\beta}_4\mapsto\zero$。
            Jordan 形式为
            \begin{equation*}
              \begin{mat}[r]
                6  &0  &0  &0  \\
                1  &6  &0  &0  \\
                0  &1  &6  &0  \\
                0  &0  &0  &6  \\
              \end{mat}
            \end{equation*}
            而找一个合适的链基是例行的操作。
            \begin{equation*}
              B=\sequence{\colvec[r]{0 \\ 0 \\ 0 \\ 1},
                          \colvec[r]{2 \\ -1 \\ -1 \\ 2},
                          \colvec[r]{3 \\ 3 \\ -6 \\ 3},
                          \colvec[r]{-1 \\ -1 \\ 1 \\ 0}}
            \end{equation*}
      \end{exparts}
    
\end{ans}
\begin{ans}{五.IV.2.20}
      有两个特征值，$\lambda_1=-2$ 与 $\lambda_2=1$。
      $t+2$ 限制在
      $\gennullspace{t+2}$ 上，在一个相伴的
      链基上可能有下面两种作用。
      \begin{equation*}
        \vec{\beta}_1\mapsto\vec{\beta}_2\mapsto\zero
        \qquad
        \begin{array}[t]{l}
          \vec{\beta}_1\mapsto\zero  \\
          \vec{\beta}_2\mapsto\zero
        \end{array}
      \end{equation*}
      $t-1$ 限制在
      $\gennullspace{t-1}$ 上，在一个相伴的
      链基上可能有下面两种作用。
      \begin{equation*}
        \vec{\beta}_3\mapsto\vec{\beta}_4\mapsto\zero
        \qquad
        \begin{array}[t]{l}
          \vec{\beta}_3\mapsto\zero  \\
          \vec{\beta}_4\mapsto\zero
        \end{array}
      \end{equation*}
      合起来共有四种可能的 Jordan 形式：
      两个都取第一种作用，第二种与第一种，第一种与第二种，
      以及两个都取第二种作用。
      \begin{equation*}
        \begin{mat}[r]
          -2  &0  &0  &0  \\
           1  &-2 &0  &0  \\
           0  &0  &1  &0  \\
           0  &0  &1  &1
        \end{mat}
        \quad
        \begin{mat}[r]
          -2  &0  &0  &0  \\
           0  &-2 &0  &0  \\
           0  &0  &1  &0  \\
           0  &0  &1  &1
        \end{mat}
        \quad
        \begin{mat}[r]
          -2  &0  &0  &0  \\
           1  &-2 &0  &0  \\
           0  &0  &1  &0  \\
           0  &0  &0  &1
        \end{mat}
        \quad
        \begin{mat}[r]
          -2  &0  &0  &0  \\
           0  &-2 &0  &0  \\
           0  &0  &1  &0  \\
           0  &0  &0  &1
        \end{mat}
     \end{equation*}
    
\end{ans}
\begin{ans}{五.IV.2.21}
     $t+2$ 限制在
     $\gennullspace{t+2}$ 上只能有作用
     $\vec{\beta}_1\mapsto\zero$。
     $t-1$ 限制在 $\gennullspace{t-1}$ 上，在一个相伴的链基上可能有
     下面三种作用中的任意一种。
     \begin{equation*}
        \vec{\beta}_2\mapsto\vec{\beta}_3\mapsto\vec{\beta}_4\mapsto\zero
        \qquad
        \begin{array}[t]{l}
          \vec{\beta}_2\mapsto\vec{\beta}_3\mapsto\zero  \\
          \vec{\beta}_4\mapsto\zero
        \end{array}
        \qquad
        \begin{array}[t]{l}
          \vec{\beta}_2\mapsto\zero  \\
          \vec{\beta}_3\mapsto\zero  \\
          \vec{\beta}_4\mapsto\zero
        \end{array}
     \end{equation*}
     合起来共有三种可能的 Jordan 形式：
     由 $t-1$ 的第一种作用（连同 $t+2$ 唯一的
     作用）产生的形式、由第二种作用产生的形式，
     以及由第三种作用产生的形式。
     \begin{equation*}
       \begin{mat}[r]
         -2  &0  &0  &0  \\
          0  &1  &0  &0  \\
          0  &1  &1  &0  \\
          0  &0  &1  &1
       \end{mat}
       \quad
       \begin{mat}[r]
         -2  &0  &0  &0  \\
          0  &1  &0  &0  \\
          0  &1  &1  &0  \\
          0  &0  &0  &1
       \end{mat}
       \quad
       \begin{mat}[r]
         -2  &0  &0  &0  \\
          0  &1  &0  &0  \\
          0  &0  &1  &0  \\
          0  &0  &0  &1
       \end{mat}
     \end{equation*}
    
\end{ans}
\begin{ans}{五.IV.2.22}
      $t+1$ 在 $\gennullspace{t+1}$ 的一个链基上的作用
      必为 $\vec{\beta}_1\mapsto\zero$。
      由于极小多项式中 \( x-2 \) 的幂次，$t-2$ 的一个链基
      长度为二，因此
      \( t-2 \) 在 \( \gennullspace{t-2} \) 上的作用
      必须具有如下形式。
      \begin{equation*}
        \begin{array}[t]{l}
          \vec{\beta}_2\mapsto\vec{\beta}_3\mapsto\zero  \\
          \vec{\beta}_4\mapsto\zero
        \end{array}
      \end{equation*}
      因此只可能有一种 Jordan 形式。
      \begin{equation*}
          \begin{mat}[r]
            -1  &0  &0  &0  \\
             0  &2  &0  &0  \\
             0  &1  &2  &0  \\
             0  &0  &0  &2
          \end{mat}
       \end{equation*}
      
\end{ans}
\begin{ans}{五.IV.2.23}
      有两种可能的 Jordan 形式。
      $t+1$ 在 $\gennullspace{t+1}$ 的一个链基上的作用
      必为 $\vec{\beta}_1\mapsto\zero$。
      对于 $t-2$ 在
      $\gennullspace{t-2}$ 的一个链基上的作用，在这个特征
      多项式与极小多项式之下有两种可能的情形。
      \begin{equation*}
        \begin{array}[t]{l}
          \vec{\beta}_2\mapsto\vec{\beta}_3\mapsto\zero  \\
          \vec{\beta}_4\mapsto\vec{\beta}_5\mapsto\zero
        \end{array}
        \qquad
        \begin{array}[t]{l}
          \vec{\beta}_2\mapsto\vec{\beta}_3\mapsto\zero  \\
          \vec{\beta}_4\mapsto\zero                      \\
          \vec{\beta}_5\mapsto\zero
        \end{array}
      \end{equation*}
      由此得到的 Jordan 形式矩阵如下。
      \begin{equation*}
        \begin{mat}[r]
          -1  &0  &0  &0  &0  \\
           0  &2  &0  &0  &0  \\
           0  &1  &2  &0  &0  \\
           0  &0  &0  &2  &0  \\
           0  &0  &0  &1  &2
        \end{mat}
        \qquad
        \begin{mat}[r]
          -1  &0  &0  &0  &0  \\
           0  &2  &0  &0  &0  \\
           0  &1  &2  &0  &0  \\
           0  &0  &0  &2  &0  \\
           0  &0  &0  &0  &2
        \end{mat}
     \end{equation*}
    
\end{ans}
\begin{ans}{五.IV.2.24}
      \begin{exparts}
        \partsitem 特征多项式为 \( c(x)=x(x-1) \)。
          对于 $\lambda_1=0$，有
          \begin{equation*}
            \nullspace{t-0}=\set{\colvec{-y \\ y}
                                 \suchthat y\in\C }
          \end{equation*}
          （当然，$t^2$ 的零空间与之相同）。
          对于 $\lambda_2=1$，
          \begin{equation*}
            \nullspace{t-1}=\set{\colvec{x \\ 0}
                                     \suchthat x\in\C }
          \end{equation*}
          （而 $(t-1)^2$ 的零空间与之相同）。
          我们可以取这个基
          \begin{equation*}
            B=\sequence{\colvec[r]{1 \\ -1},\colvec[r]{1 \\ 0}}
          \end{equation*}
          来得到对角化。
          \begin{equation*}
            \begin{mat}[r]
              1  &1  \\
             -1  &0
            \end{mat}^{-1}
            \begin{mat}[r]
              1  &1  \\
              0  &0
            \end{mat}
            \begin{mat}[r]
              1  &1  \\
             -1  &0
            \end{mat}
            =
            \begin{mat}[r]
              0  &0  \\
              0  &1
            \end{mat}
          \end{equation*}
        \partsitem 特征多项式为
          \( c(x)=x^2-1=(x+1)(x-1) \)。
          对于 $\lambda_1=-1$，
          \begin{equation*}
            \nullspace{t+1}=\set{\colvec{-y \\ y}
                                    \suchthat y\in\C }
          \end{equation*}
          且 $(t+1)^2$ 的零空间与之相同。
          对于 $\lambda_2=1$，
          \begin{equation*}
            \nullspace{t-1}=\set{\colvec{y \\ y}
                                    \suchthat y\in\C }
          \end{equation*}
          且 $(t-1)^2$ 的零空间与之相同。
          我们可以取这个基
          \begin{equation*}
            B=\sequence{\colvec[r]{1 \\ -1},\colvec[r]{1 \\ 1}}
          \end{equation*}
          来得到一个对角化。
          \begin{equation*}
            \begin{mat}[r]
              1  &1  \\
              1  &-1
            \end{mat}^{-1}
            \begin{mat}[r]
              0  &1  \\
              1  &0
            \end{mat}
            \begin{mat}[r]
              1   &1  \\
              -1  &1
            \end{mat}
            =
            \begin{mat}[r]
              -1  &0  \\
              0   &1
            \end{mat}
          \end{equation*}
      \end{exparts}
     
\end{ans}
\begin{ans}{五.IV.2.25}
      变换 $\map{d/dx}{\polyspace_3}{\polyspace_3}$
      是幂零的。
      它在基 \( B=\sequence{x^3,3x^2,6x,6} \) 上的作用
      是 $x^3\mapsto 3x^2\mapsto 6x\mapsto 6\mapsto 0$。
      它的 Jordan 形式就是它作为幂零矩阵的标准形。
      \begin{equation*}
         J=
         \begin{mat}[r]
           0  &0  &0  &0  \\
           1  &0  &0  &0  \\
           0  &1  &0  &0  \\
           0  &0  &1  &0
         \end{mat}
      \end{equation*}
    
\end{ans}
\begin{ans}{五.IV.2.26}
        相似。
        两者都具有特征多项式 $(x+1)^2$。
        计算 $T_1+1\cdot I$ 与
        $T_2+1\cdot I$ 的各次幂得到下面两个。
        \begin{equation*}
          \nullspace{t_1+1}=\set{\colvec{y/2 \\ y} \suchthat y\in\C}
          \qquad
          \nullspace{t_2+1}=\set{\colvec{0 \\ y} \suchthat y\in\C}
        \end{equation*}
        （当然，对每一个来说，其平方的零空间
        都是整个空间。）
        零度的增长方式表明每一个都
        相似于这个 Jordan 形式矩阵
        \begin{equation*}
           \begin{mat}[r]
             -1  &0  \\
              1  &-1 \\
           \end{mat}
        \end{equation*}
        因而它们彼此相似。
      
\end{ans}
\begin{ans}{五.IV.2.27}
       它的特征多项式是
       \( c(x)=x^2+1 \)，它有复根
       \( x^2+1=(x+i)(x-i) \)。
       因为根互不相同，
       所以该矩阵可对角化，其 Jordan 形式就是那个
       对角矩阵。
       \begin{equation*}
         \begin{mat}[r]
           -i  &0  \\
            0  &i
         \end{mat}
       \end{equation*}
       为了找一个相伴的基，我们计算零空间。
       \begin{equation*}
         \nullspace{t+i}=\set{\colvec{-iy \\ y}
                                        \suchthat y\in\C}
         \qquad
         \nullspace{t-i}=\set{\colvec{iy \\ y}
                                        \suchthat y\in\C}
       \end{equation*}
       例如，
       \begin{equation*}
         T+i\cdot I=
         \begin{mat}[r]
           i  &-1  \\
           1  &i
         \end{mat}
       \end{equation*}
       于是通过求解这个线性方程组，
       我们得到 $t+i$ 的零空间的描述。
       \begin{equation*}
         \begin{linsys}{2}
           ix  &-  &y  &=  &0  \\
            x  &+  &iy &=  &0
         \end{linsys}
         \grstep{i\rho_1+\rho_2}
         \begin{linsys}{2}
           ix  &-  &y  &=  &0  \\
               &   &0  &=  &0
         \end{linsys}
       \end{equation*}
       （要把关系式 $ix=y$ 改写为用自由变量 $y$ 表示首变量 $x$
       的形式，可以将两
       边乘以 $-i$。）

       结果，其中一个这样的基如下。
       \begin{equation*}
         B=\sequence{\colvec[r]{-i \\ 1},
                     \colvec[r]{i \\ 1}}
       \end{equation*}
     
\end{ans}
\begin{ans}{五.IV.2.28}
     我们可以通过数可能的规范代表元，也就是数可能的
     Jordan 形式矩阵，来数出可能的类的个数。
     特征多项式必定是 $c_1(x)=(x+3)^2(x-4)$
     或者 $c_2(x)=(x+3)(x-4)^2$。
     在 $c_1$ 的情形，$t+3$ 在 $\gennullspace{t+3}$ 的一个链基上
     有两种可能的作用。
     \begin{equation*}
       \vec{\beta}_1\mapsto\vec{\beta}_2\mapsto \zero
       \qquad
       \begin{array}[t]{l}
         \vec{\beta}_1\mapsto\zero \\
         \vec{\beta}_2\mapsto\zero
       \end{array}
     \end{equation*}
     有两个相伴的 Jordan 形式矩阵。
     \begin{equation*}
        \begin{mat}[r]
          -3  &0  &0  \\
           1  &-3 &0  \\
           0  &0  &4
        \end{mat}
        \qquad
        \begin{mat}[r]
          -3  &0  &0  \\
           0  &-3 &0  \\
           0  &0  &4
        \end{mat}
      \end{equation*}
      类似地，也有两个可能由 $c_2$ 产生的
      Jordan 形式矩阵。
      \begin{equation*}
        \begin{mat}[r]
          -3  &0  &0  \\
           0  &4  &0  \\
           0  &1  &4
        \end{mat}
        \qquad
        \begin{mat}[r]
          -3  &0  &0  \\
           0  &4  &0  \\
           0  &0  &4
        \end{mat}
     \end{equation*}
     所以总共有四种可能的 Jordan 形式。
    
\end{ans}
\begin{ans}{五.IV.2.29}
       Jordan 形式是唯一的。
       对角矩阵本身就是 Jordan 形式。
       因此可对角化矩阵的 Jordan 形式就是它的对角化。
       若极小多项式有幂次高于一次的因式，
       则 Jordan 形式在次对角线上有 \( 1 \)，从而
       不是对角的。
     
\end{ans}
\begin{ans}{五.IV.2.30}
       一个例子是 \( \C \) 上把 \( x \) 映到 \( -x \) 的
        变换。
      
\end{ans}
\begin{ans}{五.IV.2.31}
       两次应用 \nearbylemma{le:tInvIfftMinLambdaInv}；
       该子空间是 $t-\lambda_1$~不变的当且仅当它是
       $t$~不变的，而这又当且仅当它是
       $t-\lambda_2$~不变的。
     
\end{ans}
\begin{ans}{五.IV.2.32}
       不成立；这两个 $\nbyn{4}$ 矩阵都有 $c(x)=(x-3)^4$
       与 $m(x)=(x-3)^2$。
       \begin{equation*}
          \begin{mat}[r]
             3  &0  &0  &0  \\
             1  &3  &0  &0  \\
             0  &0  &3  &0  \\
             0  &0  &1  &3
          \end{mat}
          \quad
          \begin{mat}[r]
             3  &0  &0  &0  \\
             1  &3  &0  &0  \\
             0  &0  &3  &0  \\
             0  &0  &0  &3
          \end{mat}
       \end{equation*}
     
\end{ans}
\begin{ans}{五.IV.2.33}
      \begin{exparts}
         \partsitem 特征多项式如下。
           \begin{equation*}
             \begin{vmat}
               a-x  &b  \\
               c  &d-x
             \end{vmat}
             =(a-x)(d-x)-bc=ad-(a+d)x+x^2-bc
             =x^2-(a+d)x+(ad-bc)
           \end{equation*}
           注意行列式作为常数项出现。
         \partsitem 回忆特征多项式
            \( \deter{T-xI} \) 在相似下不变。
            利用置换展开公式证明迹
            是 \( x^{n-1} \) 的系数的相反数。
         \partsitem 不，存在矩阵 $T$ 与 $S$，它们
            等价 $S=PTQ$（对某个非奇异的 $P$ 与 $Q$），
            但迹不同。
            一个简单的例子如下。
            \begin{equation*}
               PTQ=
               \begin{mat}[r]
                  2  &0  \\
                  0  &1
               \end{mat}
               \begin{mat}[r]
                  1  &0  \\
                  0  &1
               \end{mat}
               \begin{mat}[r]
                  1  &0  \\
                  0  &1
               \end{mat}
               =
               \begin{mat}[r]
                  2  &0  \\
                  0  &1
               \end{mat}
            \end{equation*}
            用 $\nbyn{1}$ 矩阵甚至能得到更简单的例子。
         \partsitem 把矩阵化成 Jordan 形式。
            由第一小题，迹不变。
         \partsitem 第一部分容易；用第三小题。
            逆命题不成立：~这个矩阵
            \begin{equation*}
               \begin{mat}[r]
                  1  &0  \\
                  0  &-1
               \end{mat}
            \end{equation*}
            迹为零但不是幂零的。
       \end{exparts}
     
\end{ans}
\begin{ans}{五.IV.2.34}
      设 \( B_M \) 是某个
      向量空间的子空间 \( M \) 的一个基。
      一个方向的蕴含是显然的；若 \( M \) 是 \( t \) 不变的，则
      特别地，若 \( \vec{m}\in B_M \)，那么 \( t(\vec{m})\in M \)。
      对于另一个方向，设
      \( B_M=\sequence{\vec{\beta}_1,\dots,\vec{\beta}_q} \)，注意
      \( t(\vec{m})=t(m_1\vec{\beta}_1+\dots+m_q\vec{\beta}_q)
                   =m_1t(\vec{\beta}_1)+\dots+m_qt(\vec{\beta}_q) \)
      在 \( M \) 中，因为任何子空间对线性
      组合封闭。
    
\end{ans}
\begin{ans}{五.IV.2.35}
      是的，$t$ 不变子空间的交
      是 $t$~不变的。
      设 \( M \) 与 \( N \) 都是 \( t \) 不变的。
      若 \( \vec{v}\in M\intersection N \)，则 \( t(\vec{v})\in M \)
      由 \( M \) 的不变性有 \( t(\vec{v})\in M \)，由
      \( N \) 的不变性有 \( t(\vec{v})\in N \)。

      当然，两个子空间的并未必是子空间
      （记住 $x$ 轴\hbox{} 与 $y$ 轴是平面
      $\Re^2$ 的子空间，但两轴的并对向量加法不封闭；
      例如它不包含
      $\vec{e}_1+\vec{e}_2$。）
      然而，不变子集的并是不变子集；若
      \( \vec{v}\in M\union N \)，则 \( \vec{v}\in M \) 或 \( \vec{v}\in N \)，
      于是 \( t(\vec{v})\in M \) 或 \( t(\vec{v})\in N \)。

      不，不变子空间的补未必是不变的。
      考虑零变换之下 \( \C^2 \) 的子空间
      \begin{equation*}
        \set{\colvec{x \\ 0}\suchthat x\in\C}
      \end{equation*}
      。

      是的，两个不变子空间的和是不变的。
      验证很容易。
    
\end{ans}
\begin{ans}{五.IV.2.36}
       一种这样的顺序是\definend{字典序}。
       先按实部排序，再按 \( i \) 的系数排序。
       例如，\( 3+2i<4+1i \) 但 \( 4+1i<4+2i \)。
     
\end{ans}
\begin{ans}{五.IV.2.37}
     前一半容易\Dash 任何实系数多项式的导数是次数更低的
      实系数多项式。
      后一半的答案是「没有」；\( \polyspace_j(\Re) \) 的任何补
      必须包含一个次数为 \( j+1 \) 的多项式，
      而该多项式的导数在 \( \polyspace_j(\Re)\) 中。
     
\end{ans}
\begin{ans}{五.IV.2.38}
      对于前一半，证明每一个都是子空间，然后观察
      到任何多项式都可以唯一地
      写成偶次项与奇次项之和（零
      多项式两者都是）。
      后一半的答案是「不」：\( x^2 \) 是偶的，而
      \( 2x \) 是奇的。
     
\end{ans}
\begin{ans}{五.IV.2.39}
         把矩阵化成 Jordan 形式。
         由非奇异性，对角线上没有零特征值。
         仿照这个例子：
          \begin{equation*}
             \begin{mat}[r]
                9  &0  &0 \\
                1  &9  &0 \\
                0  &0  &4
             \end{mat}
             =
             \begin{mat}[r]
                3  &0  &0 \\
               1/6 &3  &0 \\
                0  &0  &2
             \end{mat}^2
          \end{equation*}
         来构造一个平方根。
         证明它在相似下仍然成立：若 \( S^2=T \)，则
         \( (PSP^{-1})(PSP^{-1})=PTP^{-1} \)。
    
\end{ans}
\protect \topic {幂方法}
\begin{ans}{1}
     \begin{exparts}
       \partsitem 一眼即可看出，最大特征值是 $4$。
         \Sage{} 给出如下结果。
\begin{lstlisting}
sage: def eigen(M,v,num_loops=10):
....:     for p in range(num_loops):
....:         v_normalized = (1/v.norm())*v
....:         v = M*v
....:     return v
....:
sage: M = matrix(RDF, [[1,5], [0,4]])
sage: v = vector(RDF, [1, 2])
sage: v = eigen(M,v)
sage: (M*v).dot_product(v)/v.dot_product(v)
4.00000147259
\end{lstlisting}
       \partsitem 简单的计算表明，最大特征值
          为~$2$。
          \Sage{} 给出如下结果。
\begin{lstlisting}
sage: M = matrix(RDF, [[3,2], [-1,0]])
sage: v = vector(RDF, [1, 2])
sage: v = eigen(M,v)
sage: (M*v).dot_product(v)/v.dot_product(v)
2.00097741083
\end{lstlisting}
     \end{exparts}
    
\end{ans}
\begin{ans}{2}
       \begin{exparts}
         \partsitem 下面是 \Sage{} 的结果。
 \begin{lstlisting}
sage: def eigen_by_iter(M, v, toler=0.01):
....:     dex = 0
....:     diff = 10
....:     while abs(diff)>toler:
....:         dex = dex+1
....:         v_next = M*v
....:         v_normalized = (1/v.norm())*v
....:         v_next_normalized = (1/v_next.norm())*v_next
....:         diff = (v_next_normalized-v_normalized).norm()
....:         v_prior = v_normalized
....:         v = v_next_normalized
....:     return v, v_prior, dex
....:
sage: M = matrix(RDF, [[1,5], [0,4]])
sage: v = vector(RDF, [1, 2])
sage: v,v_prior,dex = eigen_by_iter(M,v)
sage: (M*v).norm()/v.norm()
4.00604111686
sage: dex
4
 \end{lstlisting}
      \partsitem 对这道题 \Sage{} 要多迭代几次。
        这里用到了上一小题中定义的过程。
\begin{lstlisting}
sage: M = matrix(RDF, [[3,2], [-1,0]])
sage: v = vector(RDF, [1, 2])
sage: v,v_prior,dex = eigen_by_iter(M,v)
sage: (M*v).norm()/v.norm()
2.01585174302
sage: dex
6
\end{lstlisting}
       \end{exparts}
     
\end{ans}
\begin{ans}{3}
     \begin{exparts}
       \partsitem 最大特征值是 $3$。
         \Sage{} 给出如下结果。
\begin{lstlisting}
sage: M = matrix(RDF, [[4,0,1], [-2,1,0], [-2,0,1]])
sage: v = vector(RDF, [1, 2, 3])
sage: v = eigen(M,v)
sage: (M*v).dot_product(v)/v.dot_product(v)
3.02362112326
\end{lstlisting}
       \partsitem 最大特征值是 $-3$。
\begin{lstlisting}
sage: M = matrix(RDF, [[-1,2,2], [2,2,2], [-3,-6,-6]])
sage: v = vector(RDF, [1, 2, 3])
sage: v = eigen(M,v)
sage: (M*v).dot_product(v)/v.dot_product(v)
-3.00941127145
\end{lstlisting}
     \end{exparts}
    
\end{ans}
\begin{ans}{4}
       \begin{exparts}
         \partsitem
           \Sage{} 得到如下结果，其中 \lstinline!eigen_by_iter!
           已在上面定义。
\begin{lstlisting}
sage: M = matrix(RDF, [[4,0,1], [-2,1,0], [-2,0,1]])
sage: v = vector(RDF, [1, 2, 3])
sage: v,v_prior,dex = eigen_by_iter(M,v)
sage: (M*v).dot_product(v)/v.dot_product(v)
3.05460392934
sage: dex
8
\end{lstlisting}
         \partsitem 采用如下设置，
\begin{lstlisting}
sage: M = matrix(RDF, [[-1,2,2], [2,2,2], [-3,-6,-6]])
sage: v = vector(RDF, [1, 2, 3])
\end{lstlisting}
            \Sage{} 不返回结果（可用 \lstinline!<Ctrl>-c! 中断
            计算）。
            在该例程中加入一些错误检查代码
\begin{lstlisting}
def eigen_by_iter(M, v, toler=0.01):
    dex = 0
    diff = 10
    while abs(diff)>toler:
        dex = dex+1
        if dex>1000:
            print "oops! probably in some loop: \nv=",v,"\nv_next=",v_next
        v_next = M*v
        if (v.norm()==0):
            print "oops! v is zero"
^^I    return None
        if (v_next.norm()==0):
            print "oops! v_next is zero"
            return None
        v_normalized = (1/v.norm())*v
        v_next_normalized = (1/v_next.norm())*v_next
        diff = (v_next_normalized-v_normalized).norm()
        v_prior = v_normalized
        v = v_next_normalized
    return v, v_prior, dex
\end{lstlisting}
           得到如下输出。
\begin{lstlisting}
oops! probably in some loop:
  v= (0.707106781187, -1.48029736617e-16, -0.707106781187)
  v_next= (2.12132034356, -4.4408920985e-16, -2.12132034356)
oops! probably in some loop:
  v= (-0.707106781187, 1.48029736617e-16, 0.707106781187)
  v_next= (-2.12132034356, 4.4408920985e-16, 2.12132034356)
oops! probably in some loop:
  v= (0.707106781187, -1.48029736617e-16, -0.707106781187)
  v_next= (2.12132034356, -4.4408920985e-16, -2.12132034356)
\end{lstlisting}
          可见它在循环打转。
       \end{exparts}
     
\end{ans}
\begin{ans}{5}
       理论上，这一方法将给出 $\lambda_2$。
       然而在实际计算中，舍入误差会引入
       $\vec{v}_1$ 方向上的分量，因此该方法
       仍会给出 $\lambda_1$，只是与选取较为幸运的初始向量
       相比，可能要多花些时间。
     
\end{ans}
\begin{ans}{6}
      不使用 $\vec{v}_k=T\vec{v}_{k-1}$，
      而使用 $T^{-1}\vec{v}_k=\vec{v}_{k-1}$。
    
\end{ans}
\protect \topic {稳定种群}
\begin{ans}{1}
      方程
      \begin{equation*}
         0.89I-T
         =
         \begin{mat}
           0.89 &0 \\
           0 &0.89
         \end{mat}
         -
         \begin{mat}
           .90  &.01  \\
           .10  &.99
         \end{mat}
         =
         \begin{mat}
           -.01  &-.01  \\
           -.10  &-.10
         \end{mat}
      \end{equation*}
      导出如下方程组。
      \begin{equation*}
        \begin{mat}
          -.01  &-.01  \\
          -.10  &-.10
        \end{mat}
        \colvec{p \\ r}
        =
        \colvec{0  \\ 0}
      \end{equation*}\
      于是特征向量满足 $p=-r$。
    
\end{ans}
\begin{ans}{2}
      \Sage{} 求得如下结果。
\begin{lstlisting}
sage: M = matrix(RDF, [[0.90,0.01], [0.10,0.99]])
sage: v = vector(RDF, [10000, 100000])
sage: for y in range(10):
....:     v[1] = v[1]*(1+.01)^y
....:     print "pop vector year",y," is",v
....:     v = M*v
....:
pop vector year 0  is (10000.0, 100000.0)
pop vector year 1  is (10000.0, 101000.0)
pop vector year 2  is (10010.0, 103019.899)
pop vector year 3  is (10039.19899, 106111.421211)
pop vector year 4  is (10096.3933031, 110360.453787)
pop vector year 5  is (10190.3585107, 115891.187687)
pop vector year 6  is (10330.2345365, 122872.349786)
pop vector year 7  is (10525.9345807, 131525.973067)
pop vector year 8  is (10788.6008533, 142139.351965)
pop vector year 9  is (11131.1342876, 155081.09214)
\end{lstlisting}
       于是园内数量增长约百分之十一，
       而园外数量增长约百分之五十五。
    
\end{ans}
\begin{ans}{3}
      我们要求明年公园的数量
      由留下的动物与从外面进入的动物组合而成，
      即 $p_{n+1}=0.95p_n+0.01r_n$。
      而明年世界其余地区的数量
      由离开公园的动物与留在世界其余地区的动物组合而成，
      即 $r_{n+1}=0.05p_n+0.99r_n$。
      矩阵方程
      \begin{equation*}
        \colvec{p_{n+1}  \\ r_{n+1}}
        =
        \begin{mat}
          0.95  &0.01  \\
          0.05  &0.99
        \end{mat}
        \colvec{p_n  \\ r_n}
      \end{equation*}
      表明，为求特征值我们要解
      \begin{equation*}
        0
        =
        \begin{vmat}
          \lambda-0.95  &0.01  \\
          0.10          &\lambda-0.99
        \end{vmat}
        =\lambda^2-1.94\lambda-0.9405
      \end{equation*}
      \Sage{} 给出如下结果。
\begin{lstlisting}
sage: a,b,c = var('a,b,c')
sage: qe = (x^2 - 1.94*x -.9405 == 0)
sage: print solve(qe, x)
[
x == -1/100*sqrt(18814) + 97/100,
x == 1/100*sqrt(18814) + 97/100
]
sage: n(-1/100*sqrt(18814) + 97/100)
-0.401641352540816
sage: n(1/100*sqrt(18814) + 97/100)
2.34164135254082
\end{lstlisting}
      于是，数量要动态地稳定，唯一的途径是公园的数量每年增长 $234\%$。
    
\end{ans}
\begin{ans}{4}
      \begin{exparts}
        \partsitem 递推关系如下。
          \begin{equation*}
            \colvec{c_{n+1} \\ u_{n+1} \\ m_{n+1}}
            =
            \begin{mat}
              .95  &.06  &0  \\
              .04  &.90  &.10  \\
              .01  &.04  &.90
            \end{mat}
            \colvec{c_n \\ u_n  \\ m_n}
          \end{equation*}
        \partsitem 方程组
          \begin{equation*}
            \begin{linsys}{3}
              .05c &- &.06u &  &     &= &0  \\
             -.04c &+ &.10u &- &.10m &= &0  \\
             -.01c &- &.04u &+ &.10m &= &0
            \end{linsys}
          \end{equation*}
          有无穷多解。
\begin{lstlisting}
sage: var('c,u,m')
(c, u, m)
sage: eqns = [ .05*c-.06*u  == 0,
....:         -.04*c+.10*u-.10*m == 0,
....:         -.01*c-.04*u+.10*m == 0 ]
sage: solve(eqns, c,u,m)
[[c == 30/13*r1, u == 25/13*r1, m == r1]]
\end{lstlisting}
      \end{exparts}
    
\end{ans}
\protect \topic {网页排名}
\begin{ans}{1}
      第~$j$ 列的元素之和为
      $\sum_i \alpha h_{i,j}+(1-\alpha)s_{i,j}
       =\sum_i \alpha h_{i,j}+\sum_i (1-\alpha)s_{i,j}
       =\alpha\sum_i \alpha h_{i,j} +(1-\alpha)\sum_i s_{i,j}
       =\alpha\cdot 1+(1-\alpha)\cdot 1$，
      其结果为一。
    
\end{ans}
\begin{ans}{2}
      下面的 \textit{Sage} 会话给出了相等的值。
\begin{lstlisting}
sage: H=matrix(QQ,[[0,0,0,1], [1,0,0,0], [0,1,0,0], [0,0,1,0]])
sage: S=matrix(QQ,[[1/4,1/4,1/4,1/4], [1/4,1/4,1/4,1/4], [1/4,1/4,1/4,1/4], [1/4,1/4,1/4,1/4]])
sage: alpha=0.85
sage: G=alpha*H+(1-alpha)*S
sage: I=matrix(QQ,[[1,0,0,0], [0,1,0,0], [0,0,1,0], [0,0,0,1]])
sage: N=G-I
sage: 1200*N
[-1155.00000000000  45.0000000000000  45.0000000000000  1065.00000000000]
[ 1065.00000000000 -1155.00000000000  45.0000000000000  45.0000000000000]
[ 45.0000000000000  1065.00000000000 -1155.00000000000  45.0000000000000]
[ 45.0000000000000  45.0000000000000  1065.00000000000 -1155.00000000000]
sage: M=matrix(QQ,[[-1155,45,45,1065], [1065,-1155,45,45], [45,1065,-1155,45], [45,45,1065,-1155]])
sage: M.echelon_form()
[ 1  0  0 -1]
[ 0  1  0 -1]
[ 0  0  1 -1]
[ 0  0  0  0]
sage: v=vector([1,1,1,1])
sage: (v/v.norm()).n()
(0.500000000000000, 0.500000000000000, 0.500000000000000, 0.500000000000000)
\end{lstlisting}
    
\end{ans}
\begin{ans}{3}
      我们有下式。
      \begin{equation*}
          H=\begin{mat}
               0     &0    &1    &1/2   \\
               1/3   &0    &0    &0   \\
               1/3   &1/2  &0    &1/2 \\
               1/3   &1/2  &0    &0
          \end{mat}
      \end{equation*}
      \begin{exparts}
        \item 下面的 \Sage{} 会话给出了答案。
\begin{lstlisting}
sage: H=matrix(QQ,[[0,0,1,1/2], [1/3,0,0,0], [1/3,1/2,0,1/2], [1/3,1/2,0,0]])
sage: S=matrix(QQ,[[1/4,1/4,1/4,1/4], [1/4,1/4,1/4,1/4], [1/4,1/4,1/4,1/4], [1/4,1/4,1/4,1/4]])
sage: I=matrix(QQ,[[1,0,0,0], [0,1,0,0], [0,0,1,0], [0,0,0,1]])
sage: alpha=0.85
sage: G=alpha*H+(1-alpha)*S
sage: N=G-I
sage: 1200*N
[-1155.00000000000  45.0000000000000  1065.00000000000  555.000000000000]
[ 385.000000000000 -1155.00000000000  45.0000000000000  45.0000000000000]
[ 385.000000000000  555.000000000000 -1155.00000000000  555.000000000000]
[ 385.000000000000  555.000000000000  45.0000000000000 -1155.00000000000]
sage: M=matrix(QQ,[[-1155,45,1065,555], [385,-1155,45,45], [385,555,-1155,555],
....:      [385,555,45,-1155]])
sage: M.echelon_form()
[            1             0             0 -106613/58520]
[            0             1             0        -40/57]
[            0             0             1        -57/40]
[            0             0             0             0]
sage: v=vector([106613/58520,40/57,57/40,1])
sage: (v/v.norm()).n()
(0.696483066294572, 0.268280959381099, 0.544778023143244, 0.382300367118066)
\end{lstlisting}
        \item 继续该 \Sage{} 会话得到下式。
\begin{lstlisting}
 sage: alpha=0.95
sage: G=alpha*H+(1-alpha)*S
sage: N=G-I
sage: 1200*N
[-1185.00000000000  15.0000000000000  1155.00000000000  585.000000000000]
[ 395.000000000000 -1185.00000000000  15.0000000000000  15.0000000000000]
[ 395.000000000000  585.000000000000 -1185.00000000000  585.000000000000]
[ 395.000000000000  585.000000000000  15.0000000000000 -1185.00000000000]
sage: M=matrix(QQ,[[-1185,15,1155,585], [395,-1185,15,15], [395,585,-1185,585],
....:      [395,585,15,-1185]])
sage: M.echelon_form()
[             1              0              0 -361677/186440]
[             0              1              0         -40/59]
[             0              0              1         -59/40]
[             0              0              0              0]
sage: v=vector([361677/186440,40/59,59/40,1])
sage: (v/v.norm()).n()
(0.713196892748114, 0.249250262646952, 0.542275102671275, 0.367644137404254)
\end{lstlisting}
       \item 网页 $p_3$ 是重要的，但它的重要度只传递给
         一个网页，即 $p_1$。
         于是该网页得到很大的提升。
      \end{exparts}
    
\end{ans}
\protect \topic {线性递推}
\begin{ans}{1}
     我们使用公式
     \begin{equation*}
       F(n)=\frac{1}{\sqrt{5}}\left[\left(\frac{1+\sqrt{5}}{2}\right)^{n}
                               -\left(\frac{1-\sqrt{5}}{2}\right)^{n}\right]
     \end{equation*}
     如前面所观察到的，$(1+\sqrt{5})/2$ 大于一，
     而 $(1+\sqrt{5})/2$ 的绝对值小于一。
\begin{lstlisting}
sage: phi = (1+5^(0.5))/2
sage: psi = (1-5^(0.5))/2
sage: phi
1.61803398874989
sage: psi
-0.618033988749895
\end{lstlisting}
     所以该表达式的值由第一项主导。
     解 $1000=(1/\sqrt{5})\cdot ((1+\sqrt{5})/2)^n$ 给出
     下式。
\begin{lstlisting}
sage: a = ln(1000*5^(0.5))/ln(phi)
sage: a
16.0271918385296
sage: psi^(17)
-0.000280033582072583
\end{lstlisting}
   于是到十七次幂时，第二项的贡献已不足以
   改变舍入的结果。
   对于一万和一百万的计算，情况
   更加极端。
\begin{lstlisting}
sage: b = ln(10000*5^(0.5))/ln(phi)
sage: b
20.8121638053112
sage: c = ln(1000000*5^(0.5))/ln(phi)
sage: c
30.3821077388746
\end{lstlisting}
   在这两种情形下，答案是 $21$ 和 $31$。
   
\end{ans}
\begin{ans}{2}
    \begin{exparts}
      \partsitem
        我们把该关系写成矩阵形式。
        \begin{equation*}
          \colvec{f(n) \\ f(n-1)}
          =
          \begin{mat}[r]
            5  &-6  \\
            1  &0
          \end{mat}
          \colvec{f(n-1) \\ f(n-2)}
        \end{equation*}
        这个矩阵的特征方程
        \begin{equation*}
          \begin{vmat}
            5-\lambda &-6       \\
            1         &-\lambda
          \end{vmat}
          =\lambda^2-5\lambda+6
        \end{equation*}
        有根 $2$ 和 $3$。
        任何形如
        $f(n)=c_12^n+c_23^n$
        的函数都满足该递推。
      \partsitem
        该关系的矩阵表示是
        \begin{equation*}
          \colvec{f(n) \\ f(n-1)}
          =
          \begin{mat}[r]
            0  &4  \\
            1  &0
          \end{mat}
          \colvec{f(n-1) \\ f(n-2)}
        \end{equation*}
        而特征方程
        \begin{equation*}
          \begin{vmat}
            \lambda^2-2
          \end{vmat}
          =(\lambda-2)(\lambda+2)
        \end{equation*}
        有两个根 $2$ 和 $-2$。
        任何形如
        $f(n)=c_12^n+c_2(-2)^n$
        的函数都满足这个递推。
      \partsitem
        用矩阵形式，该关系
        \begin{equation*}
          \colvec{f(n) \\ f(n-1) \\ f(n-2)}
          =
          \begin{mat}[r]
            5  &-2  &-8  \\
            1  &0  &0  \\
            0  &1  &0
          \end{mat}
          \colvec{f(n-1) \\ f(n-2) \\ f(n-3)}
        \end{equation*}
        的特征方程有根 $-1$、$2$ 和 $4$。
        任何形如
        $c_1(-1)^n+c_22^n+c_34^n$ 的组合都是该递推的解。
    \end{exparts}
   
\end{ans}
\begin{ans}{3}
      \begin{exparts}
        \partsitem 该齐次递推的解是
          $f(n)=c_12^n+c_23^n$。
          代入 $f(0)=1$ 与 $f(1)=1$ 给出这个线性方程组。
          \begin{equation*}
            \begin{linsys}{2}
              c_1 &+ &c_2   &= &1 \\
             2c_1 &+ &3c_2  &= &1
            \end{linsys}
          \end{equation*}
          一眼可见 $c_1=2$，$c_2=-1$。
        \partsitem
          该齐次递推的解是
          $c_12^n+c_2(-2)^n$。
          初始条件给出这个线性方程组。
          \begin{equation*}
            \begin{linsys}{2}
              c_1 &+ &c_2  &= &0  \\
             2c_1 &- &2c_2 &= &1
            \end{linsys}
          \end{equation*}
          解为 $c_1=1/4$，$c_2=-1/4$。
        \partsitem
          该齐次递推有解
          $f(n)=c_1(-1)^n+c_22^n+c_34^n$。
          结合初始条件，我们得到这个线性方程组。
          \begin{equation*}
            \begin{linsys}{3}
              c_1 &+  &c_2   &+  &c_3  &= &1 \\
             -c_1 &+  &2c_2  &+  &4c_3  &= &1 \\
              c_1 &+  &4c_2  &+  &16c_3  &= &3
            \end{linsys}
          \end{equation*}
          其解为 $c_1=1/3$，$c_2=2/3$，$c_3=0$。
      \end{exparts}
    
\end{ans}
\begin{ans}{4}
      取定一个~$k$ 阶线性齐次递推。
      设 $S$ 是满足该递推的
      函数 $\map{f}{\N}{\C}$ 的集合。
      考虑如下给出的函数 $\map{\Phi}{S}{\C^k}$。
      \begin{equation*}
        f\mapsunder{\Phi}\colvec{f(0) \\ f(1) \\ \vdots \\ f(k-1)}
      \end{equation*}
      这就证明了线性性。
      \begin{multline*}
        \Phi(a\cdot f_1+b\cdot f_2)
        =
        \colvec{af_1(0)+bf_2(0) \\ \vdots \\ af_1(k-1)+bf_2(k-1)}     \\
        =
        a\colvec{f_1(0) \\ \vdots \\ f_1(k-1)}
        +
        b\colvec{f_2(0) \\ \vdots \\ f_2(k-1)}
        =a\cdot \Phi(f_1)+b\cdot \Phi(f_2)
      \end{multline*}
    
\end{ans}
\begin{ans}{5}
      我们利用提示来证明这一点。
      \begin{align*}
        0 &=
        \begin{vmat}
          a_{n-1}-\lambda &a_{n-2}  &a_{n-3}  &\ldots  &a_{n-k+1} &a_{n-k}  \\
          1              &-\lambda &0        &\ldots  &0         &0        \\
          0              &1        &-\lambda                               \\
          0              &0        &1                                      \\
          \vdots         &\vdots &         &\ddots   &           &\vdots  \\
          0             &0        &0        &\ldots   &1          &-\lambda
        \end{vmat}                                                           \\
        &=(-1)^{k-1}(-\lambda^k+a_{n-1}\lambda^{k-1}+a_{n-2}\lambda^{k-2}
               +\dots+a_{n-k+1}\lambda+a_{n-k})                      \\
        &=\pm(-\lambda^k+a_{n-1}\lambda^{k-1}+a_{n-2}\lambda^{k-2}
               +\dots+a_{n-k+1}\lambda+a_{n-k})
      \end{align*}
      基础步骤是平凡的。
      对于归纳步骤，
      沿最后一列展开给出两个非零项。
      \begin{equation*}
        (-1)^{k-1}a_{n-k}\cdot 1
        -\lambda\cdot
        \begin{vmat}
          a_{n-1}-\lambda &a_{n-2}  &a_{n-3}  &\ldots  &a_{n-k+1}  \\
          1              &-\lambda &0        &\ldots  &0       \\
          0              &1        &-\lambda                               \\
          0              &0        &1                                      \\
          \vdots         &\vdots &         &\ddots   &        \\
          0             &0        &0        &\ldots   &1
        \end{vmat}
      \end{equation*}
      （该矩阵是方阵，所以 $-\lambda$
      前面的符号是 $-1^{\text{偶}}$）。
      应用归纳假设就得到想要的结果。
      \begin{multline*}
        =(-1)^{k-1}a_{n-k}\cdot 1                     \\
        -\lambda\cdot
         (-1)^{k-2}(-\lambda^{k-1}+a_{n-1}\lambda^{k-2}+a_{n-2}\lambda^{k-3}
               +\dots+a_{n-k+1}\lambda^0)
      \end{multline*}
    
\end{ans}
\begin{ans}{6}
       这是对 $n$ 的直接归纳。
    
\end{ans}
\begin{ans}{7}
      \Sage{} 表明我们是安全的。
\begin{lstlisting}
sage: T64 = 18446744073709551615
sage: T64_days = T64/(60*60*24)
sage: T64_days
1229782938247303441/5760
sage: T64_years = T64_days/365.25
sage: T64_years
5.84542046090626e11
sage: age_of_universe = 13.8e9
sage: T64_years/age_of_universe
42.3581192819294
\end{lstlisting}
    
\end{ans}
\protect \topic {耦合振子}
\begin{ans}{1}
    余弦函数的和角公式为
    $\cos(\alpha+\beta)=cos(\alpha)\cos(\beta)-\sin(\alpha)\sin(\beta)$。
    把 $A\cos(\omega t+\phi)$ 展开为 $A\cdot[cos(\omega
    t)\cos(\phi)-\sin(\omega t)\sin(\phi)]$。由于 $\cos(\phi)$ 与
    $\sin(\phi)$ 不随~$t$ 变化，所以我们得到通解
    $x(t)=B\cos(\omega t)+C\sin(\omega t)$。
  
\end{ans}
\begin{ans}{2}
    我们对第一个根进行计算；第二个类似。
    我们有如下方程。
    \begin{equation*}
      \begin{mat}
        \omega_0^2-\omega^2  &\epsilon/2m \\
        \epsilon/2I         &\omega_0^2-\omega^2
      \end{mat}
      \colvec{A_1 \\ A_2}
      =\colvec{0 \\ 0}
    \end{equation*}
    代入第一个根~$\omega^2=\omega_0^2+\epsilon/2\sqrt{mI}$。
    两个方程中有一个是多余的，所以我们只考虑第一个。
    \begin{equation*}
      -(\epsilon/2\sqrt{mI})\cdot A_1+(\epsilon/2m)\cdot A_2=0
    \end{equation*}
    这是平面中一条过原点的直线，因此要确定它，
    我们只需求出第一个变量与第二个变量之比。
    \begin{equation*}
      \frac{A_2}{A_1}=\frac{(\epsilon/2\sqrt{mI})}{(\epsilon/2m)}
                     =\sqrt{\frac{m}{I}}
    \end{equation*}
    所以这就是与第一个特征值相对应的特征向量空间。
    \begin{equation*}
      \set{\colvec{A_1 \\ \sqrt{m/I}\cdot A_1}\suchthat A_1\in\C}
    \end{equation*}
  
\end{ans}
\begin{ans}{3}
    注意到 $\cos(\omega t+(\phi+\pi))=-\cos(\omega t+\phi)$。
    取导数
    \begin{equation*}
      \frac{d\,x(t)}{dt}=-\omega\cdot\sin(\omega t+\phi)
      \qquad
      \frac{d^2\,x(t)}{dt^2}=-\omega^2\cdot\cos(\omega t+\phi)
    \end{equation*}
    以及
    \begin{equation*}
      \frac{d\,\theta(t)}{dt}=\omega\cdot\sin(\omega t+\phi)
      \qquad
      \frac{d^2\,\theta(t)}{dt^2}=\omega^2\cdot\cos(\omega t+\phi)
    \end{equation*}
    并代入运动方程组~($**$)。
    \begin{align*}
      mA_1\cdot(-\omega^2\cos(\omega t+\phi))
          +kA_1\cdot (\cos(\omega t+\phi))
          +\frac{\epsilon}{2}A_2\cdot (-\cos(\omega t+\phi))
          &=0  \\
      IA_2\cdot(\omega^2\cos(\omega t+\phi))
          +\kappa A_2\cdot (-\cos(\omega t+\phi))
          +\frac{\epsilon}{2}A_1\cdot (\cos(\omega t+\phi))
          &=0
    \end{align*}
    提出 $\cos(\omega t+\phi)$
    \begin{align*}
     \big(m(-\omega^2)\big)\cdot A_1
         +kA_1
         -\frac{\epsilon}{2}\cdot A_2
         &=0  \\
     \big(I\omega^2\big)\cdot A_2
         -\kappa A_2
         +\frac{\epsilon}{2}\cdot A_1
         &=0
   \end{align*}
   再用~$m$ 和~$I$ 去除。
   \begin{align*}
    \big(\frac{k}{m}-\omega^2\big)\cdot A_1
        -\frac{\epsilon}{2m}\cdot A_2
        &=0  \\
    \frac{\epsilon}{2I}\cdot A_1
    -\big(\frac{\kappa}{I}-\omega^2\big)\cdot A_2
        &=0
  \end{align*}
   我们假定 $k/m=\omega_{x}^2$
   与 $\kappa/I=\omega_{\theta}^2$ 相等，
   并把该数记为 $\omega_0^2$。
   作此替换，并将其改写为矩阵方程。
   \begin{equation*}
     \begin{mat}
       \omega_0^2-\omega^2    &-\epsilon/2m  \\
       \epsilon/2I &\omega_0^2-\omega^2
     \end{mat}
     \colvec{
       A_1 \\
       A_2}
     =
     \colvec{
       0 \\
       0}
   \end{equation*}
   我们要求使该方程组有非平凡
   解的频率~$\omega$。
   \begin{equation*}
      \begin{vmatrix}
       \omega_0^2-\omega^2    &-\epsilon/2m  \\
       \epsilon/2I           &\omega^2-\omega_0^2
      \end{vmatrix}
      =0
    \end{equation*}
    这与专题正文中算过的行列式相同，
    只是第二列乘了 $-1$。
    用一个标量乘某一行或某一列，就相当于用该标量乘
    整个行列式。
    所以这个行列式是专题正文中那个行列式的相反数。
    但我们要令它为零，所以这无关紧要。
    根与专题正文中相同。
  
\end{ans}
\begin{ans}{4}
    参见 \cite{Mewes}。
  
\end{ans}
\protect \topic {内积}
\begin{ans}{1}
本专题开头一段指出点积对两个输入都是线性的，
因此第一个条件得到满足。
条件~(2) 涉及共轭，
而共轭对实数没有影响，
所以这个条件平凡地得到满足。
条件~(3) 要求长度（$n$ 维实空间中的通常 Euclid 长度）
为零当且仅当向量是~$\zero$。
若向量有分量 $v_1$，\ldots~$v_n$，则
$v_1^2+\cdots v_n^2=0$ 当且仅当所有分量全为零。
\end{ans}
\begin{ans}{2}
专题正文的推导证明了~(1) 与~(2)。
对于~(3)，设
~$\C^n$ 的两个成员 $\vec{v},\vec{w}$
的第 $i$ 个分量分别为 $v_i$ 与 $w_i$，于是
$\innerprod{\vec{v}}{\vec{w}}=v_1\compconj{w}_1+\cdots+v_n\compconj{w}_n$。
那么 $\innerprod{\vec{v}}{\vec{v}}=v_1\compconj{v}_1+\cdots+v_n\compconj{v}_n$
等于 $\absval{\vec{v}_1}^2+\cdots \absval{\vec{v}_n}^2$，
显然它等于~$0$ 当且仅当每个分量都是~$0$。
\end{ans}
\begin{ans}{3}
条件~(1)，即对第一个变量的线性性，是容易的。
\begin{multline*}
  \innerprod{
    \begin{mat}
    a_1  &b_1  \\
    c_1  &d_1
    \end{mat}
    +
    \begin{mat}
    a_2  &b_2  \\
    c_2  &d_2
    \end{mat}
    }{
      \begin{mat}
       e  &f  \\
       g  &h
      \end{mat}
    }                                                  \\
  \begin{split}
    &=(a_1+a_2)e+(b_1+b_2)f+(c_1+c_2)g+(d_1+d_2)h   \\
    &=(a_1e+b_1f+c_1g+d_1h)+(a_2e+b_2f+c_2g+d_2h)   \\
    &=
  \innerprod{
    \begin{mat}
    a_1  &b_1  \\
    c_1  &d_1
    \end{mat}
    }{
      \begin{mat}
       e  &f  \\
       g  &h
      \end{mat}
    }
  +
  \innerprod{
    \begin{mat}
    a_2  &b_2  \\
    c_2  &d_2
    \end{mat}
    }{
      \begin{mat}
       e  &f  \\
       g  &h
      \end{mat}
    }
  \end{split}
\end{multline*}
标量乘法的验证同样常规。
条件~(2) 涉及共轭，
而共轭对实数没有影响，
所以这个条件平凡地得到满足。
对于~(3)，当然有 $a^2+b^2+c^2+d^2\geq 0$，而等号成立当且仅当
每个元素都是零。
\end{ans}
\begin{ans}{4}
条件~(1) 是它对第一个输入线性。
验证这个运算保持加法是例行的
\begin{multline*}
  \innerprod{(a_2x^2+a_1x+a_0)+(c_2x^2+c_1x+c_0)}{b_2x^2+b_1x+b_0}  \\
  \begin{split}
    &=(a_2+c_2)b_2+(a_1+c_1)b_1+(a_0+c_0)b_0             \\
    &=(a_2b_2+a_1b_1+a_0b_0)+(c_2b_2+c_1b_1+c_0b_0)             \\
    &=\innerprod{(a_2x^2+a_1x+a_0)}{b_2x^2+b_1x+b_0}      \\
      &\qquad+\innerprod{(c_2x^2+c_1x+c_0)}{b_2x^2+b_1x+b_0}
  \end{split}
\end{multline*}
验证它保持标量乘法亦然。
因为这是实向量空间，而共轭对实数标量
没有影响，所以它满足条件~(2)。
对于条件~(3)，令 $\vec{v}=a_2x^2+a_1x+a_0$。
那么 $\innerprod{\vec{v}}{\vec{v}}=a_2^2+a_1^2+a_0^2$ 显然非负，
且等于~$0$ 当且仅当 $\vec{v}$ 是零多项式。

至于 $x^2-1$ 的范数，我们有
\begin{equation*}
 \norm{x^2-1}=\sqrt{\innerprod{x^2-1}{x^2-1}}=\sqrt{1^2+0^2+(-1)^2}=\sqrt{2}
\end{equation*}
\end{ans}
\begin{ans}{5}
向量空间定义中的十个条件
都是微积分诸结果的推论。
例如，第一个条件是对加法封闭，而在微积分中
我们有：两个连续函数的和
是连续函数。
（不过，学生往往要到后面的课程
才见到完全严格的证明。）
类似地，两个函数的加法是可交换的。
当然，零向量是函数 $f(x)=0$，它是连续的。
\end{ans}
\begin{ans}{6}
假设
$\innerprod{\vec{v}}{\vec{v}}$ 对所有
复向量~$\vec{v}$ 都大于或等于~$0$。
那么对两个输入的线性性给出
$0\leq \innerprod{i\vec{v}}{i\vec{v}}
  =i\innerprod{\vec{v}}{i\vec{v}}
  =i^2\innerprod{\vec{v}}{\vec{v}}\leq 0$。
\end{ans}
